% Parte V del sucesor: mismos propietarios rectores, títulos capitulares
% aprobados y dependencia causal explícita respecto del continuo conjunto.
\makeatletter
\@ifundefined{HMTScientificOwnerSubordinated}{%
  % Entorno editorial para incorporar cuerpos científicos completos sin
% crear capitulos nuevos. El grupo conserva intactos los archivos de origen:
% solo traduce su jerarquia visible un nivel hacia abajo durante el input.
% La procedencia material de cada uso se registra en comentarios del destino
% y en gestion/RESTAURACION_16_CLAVES_SIN_DESTINO_20260824.tsv.
\newenvironment{HMTScientificOwnerSubordinated}{%
  \begingroup
  \let\HMTScientificOwnerSection\section
  \let\HMTScientificOwnerSubsection\subsection
  \let\HMTScientificOwnerSubsubsection\subsubsection
  \let\HMTScientificOwnerParagraph\paragraph
  \let\HMTScientificOwnerSubparagraph\subparagraph
  \let\chapter\HMTScientificOwnerSection
  \let\section\HMTScientificOwnerSubsection
  \let\subsection\HMTScientificOwnerSubsubsection
  \let\subsubsection\HMTScientificOwnerParagraph
  \let\paragraph\HMTScientificOwnerSubparagraph
}{%
  \par
  \endgroup
}
%
}{}
\makeatother

\part{Realizaciones terminales de la estructura discreta del continuo: \texorpdfstring{\(\zeta\)}{zeta}, Navier--Stokes, Yang--Mills, Hodge y Birch--Swinnerton--Dyer}
\label{part:clausuras-terminales-definitiva}

\HMTFlushCurrentChapter
\chapter{Teorema de factorización conjunta de las \(5\) realizaciones terminales a través de la estructura discreta del continuo}
\label{chap:terminal-arquitectura-natural-comun}
La Parte V recibe el continuo único después de su construcción conjunta y de
su realización por Mecánica Dimensional. Ninguna realización terminal
selecciona, genera ni sustituye retrospectivamente las construcciones
consustanciales del continuo. La secuencia expositiva conserva, en este orden,
la nota de lectura, el empalme causal, el objeto común y el método de
demostración.
El teorema conjunto establece la factorización simultánea, la no
permutabilidad de los lectores y la conservación del estado enriquecido
común. Las realizaciones que siguen mantienen íntegramente, y con dominio
propio, sus operadores, demostraciones y certificados individuales; no son
sustitutos ni abreviaciones de esta factorización común.
\Needspace{6\baselineskip}
\Needspace{6\baselineskip}
\section{Publicación terminal del continuo único mediante \(5\) lectores no permutables}
\Needspace{6\baselineskip}
\subsection{Factorización del estado enriquecido común en \(5\) publicaciones terminales}
\Needspace{6\baselineskip}
\subsubsection{Dependencia causal de las realizaciones terminales respecto del continuo genealógico íntegramente construido}
Los cinco problemas no convocan cinco estructuras matemáticas distintas. Tampoco obligan a construir cinco caminos independientes que, después de avanzar por separado, deban reunirse artificialmente en una conclusión común. Esa imagen pertenece a una etapa anterior de la exposición. El punto de partida vigente es más fuerte: antes de que aparezcan los nombres de Riemann, Navier–Stokes, Yang–Mills, Hodge o Birch–Swinnerton–Dyer, ya existe un único objeto matemático enriquecido cuyas operaciones internas producen simultáneamente las características necesarias para los cinco cierres.\par
Por eso este relato no comienza con los problemas. Comienza con la construcción del continuo.\par
La Holografía Modular Triádica no recibe el continuo como un fondo previamente dado. Tampoco parte de la recta real, de un espacio funcional clásico, de una variedad algebraica, de una conexión gauge o de una curva elíptica. Parte de una aritmética finita dotada de orientación, memoria y capacidad de prolongación. APP proporciona el sustrato posicional y aritmético; TRIT convierte cada operación en una transición orientada; TPK conserva la genealogía completa de esas transiciones y permite prolongarlas sin borrar el camino que las produjo.\par
El paso decisivo consiste en que la salida de una operación deja de ser solamente un valor. Cada emisión lleva consigo su hoja aritmética, residuo, cociente, orientación, acarreo, memoria, frontera, ruta, multisección, supervivencia, terminal y lector. El resultado no es ya un punto aislado, sino un estado que sabe de dónde procede, cómo puede continuar y qué relaciones debe conservar cuando cambia de escala.\par
APP, TRIT y TPK no forman tres explicaciones superpuestas. Forman una sola secuencia causal:\par
\[
\mathrm{APP}
\longrightarrow
\mathrm{TRIT}
\longrightarrow
\mathrm{TPK}
\longrightarrow
\text{estado enriquecido único}.
\]
APP fija el soporte aritmético común. TRIT introduce la orientación efectiva de cada transición, distinguiendo avance, umbral y apertura futura. TPK reúne la información visible y la información transportada: no sólo lo que una operación entrega, sino también el carry que desplaza, la frontera que modifica, la memoria que conserva y las extensiones que permite.\par
La conexión nonádica no devuelve el sistema a un estado vacío. Después de una vuelta completa reaparece la fase, pero no se borra la historia. El retorno \(9\to1\) es retorno de fase, no reinicio del estado. Cada vuelta añade profundidad, memoria y compatibilidad. De este modo, la continuidad no se obtiene repitiendo nueve filtros ni atravesando nueve canales independientes. Se obtiene prolongando una sola conexión completa cuyos nueve tramos son cartas internas de una misma holonomía.\par
\Needspace{6\baselineskip}
\subsubsection{Emergencia correlativa de \texorpdfstring{\(5\)}{5} acciones intrínsecas no exhaustivas}
Sobre cada estado enriquecido actúa una operación común. No existen primero cinco dominios, ni cinco selecciones, ni cinco copias del estado. La misma emisión es leída simultáneamente bajo cinco características internas denotadas por los discriminantes \(\mathrm{sp}\), \(\mathrm{sol}\), \(\mathrm{gau}\), \(\mathrm{coh}\) y \(\mathrm{det}\).\par
Esos símbolos no nombran ramas ni agotan las propiedades del continuo. Nombran
\(5\) acciones intrínsecas del mismo objeto cuya correlación resulta decisiva
para las realizaciones terminales de esta Parte.\par
La característica espectral–polar, \(sp\), registra la publicación orientada del estado, su complemento, su defecto, su memoria terminal y la relación entre las dos lecturas especulares de una misma forma.\par
La característica solenoidal, \(sol\), registra cómo una transición conserva estado y, al mismo tiempo, distribuye disipación, advección, carry, memoria, microestructura y shell sin sobrecontarlos.\par
La característica gauge, \(gau\), registra la equivalencia local, la incidencia, la holonomía, la curvatura y el paso compatible al cociente físico.\par
La característica cohomológica, \(coh\), registra las extensiones entre niveles, las obstrucciones relativas, la compatibilidad de los levantamientos y la supervivencia de una clase a través de todo refinamiento.\par
La característica determinantal, \(det\), registra la incidencia finita-global, las páginas de filtración, los complejos locales, sus líneas determinantes y el transporte de una base aritmética hasta una publicación analítica.\par
Pero ninguna de estas propiedades vive sola. Todas actúan sobre las mismas emisiones y conservan el mismo espesor genealógico. La fibra que transporta la característica espectral es la misma fibra enriquecida cuya orientación participa en el balance solenoidal, cuya frontera soporta la lectura gauge, cuya prolongación sostiene la extensión cohomológica y cuya incidencia entra en la línea determinante.\par
No se generan cinco objetos y luego se los empaqueta. Se genera un único objeto que posee inseparablemente esas cinco capacidades.\par
Por eso una vista posterior no es una proyección que cree una rama. Es únicamente una manera de exponer, para una aplicación concreta, una característica que ya estaba operando dentro de la totalidad. La vista no selecciona retrospectivamente estados favorables ni conserva la entrada como primera coordenada para recuperarla después de manera tautológica. Hace visible una estructura producida y conservada durante toda la genealogía.\par
\Needspace{6\baselineskip}
\subsubsection{Antecedencia causal del continuo ya construido como límite de la prolongación compatible}
En cada profundidad existe un estado finito, pero completo en su nivel. Las extensiones entre niveles no pueden alterar arbitrariamente la información anterior. Deben respetar las fibras, las relaciones de extensión, los números, la orientación, el acarreo, la memoria, la frontera, la ruta, la multisección, la supervivencia, el terminal y los lectores.\par
La naturalidad exige que operar y después truncar produzca el mismo resultado que truncar y operar. Gracias a esa compatibilidad, las historias finitas no quedan como aproximaciones desconectadas. Forman un sistema inverso. El continuo aparece cuando se toma la familia completa de historias compatibles y no cuando se postula de antemano una colección infinita de puntos.\par
El continuo HMT es, por ello, memoria coherente de caminos. Un punto sin genealogía sería una ablación: conservaría el valor visible y eliminaría precisamente aquello que permite demostrar cómo sobrevive a todo refinamiento.\par
Sólo después de cerrar esa compatibilidad aparece la estructura discreta completa del continuo. Y sólo después puede publicarse una realización arquimediana. La recta real no es la materia prima de la construcción; es una salida de uno de sus lectores terminales. Lo mismo ocurre con los dominios clásicos de los cinco problemas: aparecen cuando el continuo ya está construido y una interfaz posterior traduce una de sus características intrínsecas al lenguaje del objetivo correspondiente.\par
La secuencia causal correcta es, por tanto:\par
\[
\begin{aligned}
\mathrm{APP}
&\longrightarrow \mathrm{TRIT}
\longrightarrow \mathrm{TPK}
\longrightarrow \text{estado enriquecido único}\\
&\longrightarrow 04\mathrm b1
\longrightarrow 04\mathrm b2
\longrightarrow 04\mathrm b3\\
&\longrightarrow \text{acciones intrínsecas correlativas}
\longrightarrow \text{terminal y límite}
\longrightarrow \mathfrak C_{\mathrm{cont}}^{\mathrm{disc}}.
\end{aligned}
\]
Después, y no antes, entran los entrelazadores específicos:\par
\[
\mathfrak C_{\mathrm{cont}}^{\mathrm{disc}}
\longrightarrow
\begin{cases}
\text{forma explícita de Weil},\\
\text{evolución tridimensional incompresible},\\
\text{teoría gauge cuántica},\\
\text{publicación de clases de Hodge},\\
\text{comparación determinantal BSD}.
\end{cases}
\]
Estas cinco traducciones terminales no definen la estructura que utilizan. El problema clásico aporta sus datos después: una función de prueba en Riemann, unos datos iniciales y una fuerza en Navier–Stokes, un grupo compacto simple en Yang–Mills, una variedad y una clase en Hodge, o una curva elíptica en BSD. Ninguno de esos datos participa en la generación previa del continuo.\par
La dirección causal queda así protegida. El objetivo no selecciona el estado que habrá de resolverlo. El estado existe antes del objetivo; el entrelazador sólo demuestra que la estructura interna y la formulación clásica expresan la misma operación bajo dos lenguajes.\par
\Needspace{6\baselineskip}
\subsection{1.ª resolución: Riemann y la positividad como complemento construido}
La hipótesis de Riemann suele presentarse como una afirmación acerca de la posición de los ceros no triviales de la función zeta. Desde esa formulación, el problema parece pedir que se inspeccione una población de ceros ya dada y se demuestre que todos ellos se alinean sobre la recta crítica.\par
La lectura HMT invierte esa perspectiva. No comienza preguntando dónde están los ceros. Pregunta qué estructura genera simultáneamente los dos lados de la fórmula explícita y qué impide que su diferencia adopte signo negativo.\par
La característica espectral–polar del continuo aporta una publicación doble del mismo estado. Una orientación recoge el canal positivo; la otra recoge su imagen especular. Ambas contienen, bajo un regulador común, la contribución prima, la contribución gamma–escalar y la contribución polar. No son tres pruebas distintas ni tres sumas pegadas al final. Son tres filas de una misma coligación construida desde los datos de partida.\par
El punto esencial es que esa coligación se construye antes de invocar la positividad que deberá demostrar. Su entrada, sus salidas, su pérdida y su estado terminal están tipados dentro del mismo espacio enriquecido. La identidad que enlaza las dos publicaciones no se introduce como una igualdad deseada entre fórmulas clásicas; nace de la dinámica interna de la conexión.\par
Entre la publicación de entrada y la publicación especular puede aparecer, a profundidad finita, un residual de salida. La estructura nonádica lo obliga a satisfacer una recurrencia antiespecular. En cada vuelta completa, el residual siguiente se transporta isométricamente pero queda multiplicado por el factor\par
\[
q=\frac{5}{9}.
\]
Si \(\Omega_n\) representa ese defecto en la profundidad \(n\), la recurrencia tiene la forma\par
\[
\Omega_n=q\,V_n\Omega_{n+1},
\]
con \(V_n\) isométrico y la familia coinductivamente acotada. Al iterar \(N\) veces,\par
\[
\|\Omega_n\|
\le q^N\sup_{m\ge n}\|\Omega_m\|.
\]
Como \(0<q<1\), el término derecho tiende a cero. Por tanto, el residual no se anula porque se lo haya omitido o porque se haya supuesto una identificación entre las dos publicaciones. Se anula porque la recurrencia completa lo contrae en todas las profundidades compatibles. La salida especular coincide entonces con la publicación negativa construida.\par
Una vez descargado ese vínculo, la isometría construida desde los datos de partida entrega la identidad de energía. A profundidad finita, la norma de la publicación positiva se descompone en tres partes: la publicación negativa, la pérdida acumulada y el estado terminal. El terminal se conserva explícitamente. Debe conservarse hasta el paso al límite, porque contiene la memoria de la truncación y certifica que la telescopía prolonga la identidad finita mediante una convergencia demostrada.\par
La descomposición finita de la norma satisface la identidad exacta:\par
\[
\|\text{publicación positiva}\|^2
=
\|\text{publicación negativa}\|^2
+
\|\text{pérdidas hasta }N\|^2
+
q^{2N}\|\text{terminal}\|^2.
\]
Cada sumando posee un origen concreto. La publicación negativa es la salida especular ya identificada. Las pérdidas reúnen exactamente el complemento que no fue publicado por esa orientación. El terminal recuerda lo que todavía permanece más allá del corte.\par
Sólo después de haber escrito la identidad completa se hace \(N\to\infty\). Entonces, y sólo entonces, el terminal desaparece porque \(q^{2N}\to0\). La diferencia entre las dos publicaciones queda convertida en el cuadrado de un operador de pérdida:\par
\[
\mathcal I_{\mathrm{GW}}
=
E_R^{*}E_R.
\]
Ésta es la forma de Guinand–Weil. Su positividad ya no es una conjetura añadida a la construcción:\par
\[
\langle f,\mathcal I_{\mathrm{GW}}f\rangle
=
\|E_Rf\|^2
\ge0.
\]
La estructura HMT no obtiene esta desigualdad escondiendo la parte negativa, restringiendo el dominio a funciones que ya la satisfagan o introduciendo un proyector cuya existencia sea equivalente a la conclusión. Construye el complemento y demuestra que la forma es exactamente su cuadrado.\par
Después se materializa la publicación común. Un mismo codificador conserva los canales primo, gamma–escalar y polar. Un proyector interno produce uno de los dos momentos; su complemento produce el otro. La diferencia de los momentos no es una comparación externa entre fórmulas parecidas: es la norma del componente ortogonal construido por la coligación.\par
Al aumentar la ventana, no se obliga a las columnas individuales a satisfacer falsas isometrías entre truncamientos distintos. Lo que se conserva de manera cofinal es la forma, el complemento alcanzable y la naturalidad de su publicación. Los incrementos que entran al ampliar el corte aparecen en ambas orientaciones con la misma contribución neutral y no alteran la diferencia gobernante. Así se obtiene una familia compatible de formas positivas sobre una clase separadora.\par
La última transición es el criterio reversible de Weil. La teoría clásica afirma que la positividad de la forma explícita, sobre la clase adecuada de funciones de prueba, equivale a que todos los ceros no triviales estén situados sobre la recta crítica. HMT suministra precisamente esa positividad, pero la suministra como consecuencia de una identidad operatorial construida dentro del continuo.\par
\Needspace{5\baselineskip}
Por tanto,\par
\[
\operatorname{Re}\rho=\frac12
\]
para todo cero no trivial \(\rho\) de la función zeta.\par
La recta crítica aparece aquí de dos maneras concordantes. Geométricamente es el locus fijo de la involución que intercambia las dos orientaciones espectrales. Operatorialmente es la única posición compatible con la positividad del complemento de publicación. La simetría explica por qué la recta central es distinguida; la identidad cuadrática demuestra por qué los ceros no pueden abandonar esa recta.\par
La resolución no consiste en observar que la ecuación funcional posee una simetría. Consiste en construir el objeto cuyo defecto mide la ruptura de esa simetría, demostrar que el defecto se anula coinductivamente y convertir la diferencia restante en una norma. Es la combinación de genealogía, terminal, contracción, publicación común y positividad la que cierra el problema.\par
Así queda establecida la primera de las cinco resoluciones terminales. No como una rama independiente, sino como la traducción clásica de la característica espectral–polar que el continuo único poseía desde su construcción.\par

\Needspace{6\baselineskip}
\section{Clausuras dinámicas del continuo: transporte solenoidal y cociente de calibre}
\Needspace{6\baselineskip}
\subsection{Dinámica y curvatura del continuo genealógico}
La primera resolución mostró cómo la característica espectral–polar del continuo se publica como positividad de Weil. Las dos resoluciones siguientes interrogan el mismo estado desde otros dos aspectos consustanciales: su capacidad de transportar una dinámica no lineal sin perder energía, memoria ni regularidad, y su capacidad de sostener una teoría gauge local cuyo vacío esté separado por una brecha espectral positiva.\par
Navier–Stokes y Yang–Mills no reciben dos estructuras nuevas. Reciben dos entrelazadores distintos desde el mismo continuo ya construido. En el primer caso, el entrelazador publica el estado como un campo de velocidad incompresible con viscosidad, advección y presión. En el segundo, lo publica como una red gauge refinable con holonomía, curvatura, medida euclídea y Hamiltoniano físico.\par
La diferencia entre ambos problemas no rompe la unidad de la construcción. En Navier–Stokes, el punto crítico es demostrar que ninguna transferencia no lineal desaparece fuera del balance y que la regularidad obtenida es uniforme al retirar los cortes. En Yang–Mills, el punto crítico es demostrar que la brecha finita pertenece a la misma teoría gauge que alcanza el continuo, que sobrevive al cociente físico y que actúa en el sector de curvatura, no en un factor auxiliar.\par
\Needspace{6\baselineskip}
\subsection{Existencia, regularidad y unicidad global para Navier–Stokes mediante cota uniforme y compacidad}
El problema tridimensional de Navier–Stokes no nace de la mera presencia de un término no lineal. La dificultad aparece porque la publicación clásica del campo de velocidad tiende a ocultar dónde se almacena la transferencia entre escalas. La viscosidad es visible. La energía cinética es visible. La advección también aparece en la ecuación. Pero la contabilidad completa de la interacción entre profundidad, memoria, carry, microestructura y shell queda comprimida en una sola función visible.\par
Esa compresión es precisamente lo que la característica solenoidal del continuo evita.\par
El estado enriquecido no conserva únicamente el campo \(u\). Conserva simultáneamente el campo Galerkin, la orientación temporal, las memorias positiva y firmada, el carry entre cortes, la componente microscópica, el shell de frecuencias nuevas, la ruta, el archivo de la evolución y el terminal. La ecuación clásica es una publicación de ese estado; no es su definición completa.\par
Fijados datos iniciales \(u_0\), fuerza \(f\), viscosidad \(\nu>0\) y regularidad \(s>5/2\), la proyección Galerkin produce\par
\[
\partial_tu_M+\nu A_Mu_M+B_M(u_M,u_M)=f_M.
\]
Pero \(M\) no selecciona una solución favorable. Sólo fija una resolución finita en la que todas las operaciones pueden escribirse exactamente. El objetivo es construir una estimación que no dependa de \(M\) ni de la profundidad nonádica \(J\). Sólo una cota con esa uniformidad puede pasar al continuo.\par
\Needspace{6\baselineskip}
\subsubsection{Descomposición de la memoria visible y la memoria genealógica en la ecuación evolutiva}
La interacción no lineal posee signo. En ciertos momentos entrega energía a una escala; en otros la retira. Si se conserva sólo su valor firmado, puede perderse la magnitud total transferida. Si se conserva sólo su valor absoluto, se pierde la orientación.\par
La estructura mantiene ambas cosas.\par
Una memoria \(D_M\) publica el signo de la transferencia advectiva. Otra memoria \(H_M\) conserva su magnitud positiva. Su separación permite registrar simultáneamente\par
\[
\dot D_M=\mathscr B_{s,M},
\qquad
\dot H_M=\frac{1+r}{r-1}|\mathscr B_{s,M}|,
\]
donde \(r=\pi_{\rm HMT}/\varphi_{\rm HMT}^2>1\).\par
Esto no añade energía ficticia. Desdobla dos aspectos distintos de una sola transferencia: orientación y coste. Fundirlos destruiría información; sumarlos sin ortogonalidad duplicaría energía. La construcción hará explícito que ninguna de esas dos deformaciones ocurre.\par
\Needspace{6\baselineskip}
\subsubsection{Independencia del dato inicial respecto de la historia genealógica posterior}
La objeción más seria a cualquier clausura operatorial de Navier–Stokes sería que la cota uniforme se hubiese introducido retrospectivamente: primero se elegiría una solución regular, después se codificaría su historia y finalmente se anunciaría que el código controla esa misma solución.\par
La construcción impide esa circularidad fijando el espacio fuente antes de levantar la evolución:\par
\[
\mathcal D_T
=
H^s_\sigma(\mathbb T^3)
\times
L^2(0,T;H^{s-1}_\sigma(\mathbb T^3)),
\qquad
d=(u_0,f).
\]
La raíz \(J_T^{\rm sol}d\) depende exclusivamente de \(u_0\) y \(f\). No recibe como entrada \(u_M\), la pérdida futura, el terminal ni la norma que deberá acotarse.\par
La fuerza tampoco se inserta como un puerto aditivo artificial. La ecuación de Carleman muestra que actúa de forma mixta: cada aparición de \(f\) se combina con un grado precedente del estado. Por eso el lift usa una creación mixta estado–fuerza. La firma cronológica de \(f\) registra todos sus tiempos ordenados, mientras la torre PC–Fock registra todos los grados no lineales del estado.\par
La composición reproduce materialmente la serie de Dyson del generador no autónomo. Al aplicar el lector normalizado de grado uno se recupera exactamente \(u_M(t)\). No se recupera una aproximación simbólica ni una variable que se parezca a la solución: se recupera la solución de la ODE Galerkin porque ambas satisfacen la misma ecuación finita y el mismo dato inicial.\par
La construcción causal desde los datos de partida sigue por tanto el orden\par
\[
(u_0,f)
\longrightarrow
J_T^{\rm sol}(u_0,f)
\longrightarrow
\text{torre estado–fuerza}
\longrightarrow
\text{historia Galerkin}.
\]
La historia genealógica se construye como salida posterior de los datos de partida y queda excluida de su definición inicial.\par
\Needspace{6\baselineskip}
\subsubsection{La identidad de scattering que conserva la potencia cruzada}
La fuerza y la disipación se organizan mediante dos ondas:\par
\[
a_T(f)
=
\frac{1}{2\sqrt\nu}A^{-1/2}\Lambda^sP_\sigma f,
\qquad
z_M(u)=\sqrt\nu A_M^{1/2}\Lambda^su,
\]
y la onda residual\par
\[
b_M=z_M-P_Ma_T(f).
\]
La potencia aplicada no se reemplaza por la norma de la fuerza. Se conserva exactamente como término cruzado:\par
\[
\mathscr F_{s,M}
=
2\operatorname{Re}\langle z_M,P_Ma_T(f)\rangle.
\]
Al completar el cuadrado se obtiene el balance exacto\par
\[
\dot G_{M,j,T}
+\|b_M\|^2
+\lvert\mathscr B_{s,M}\rvert
=
\|P_Ma_T(f)\|^2.
\]
Éste es uno de los puntos que hacen funcionar el cierre. Si se sustituyese \(b=z-a\) por \(z\), desaparecería el término cruzado que representa el trabajo de la fuerza y la identidad sería falsa. La estructura no domina la no linealidad descartándola: la convierte en un canal positivo orientado y conserva además la interferencia exacta entre fuerza y disipación.\par
\Needspace{6\baselineskip}
\subsubsection{Factorización de \(6\) pérdidas mediante un operador único sin multiplicidad espuria}
La historia completa distingue seis salidas:\par
\begin{itemize}
\item disipación;
\item advección orientada;
\item carry entre cortes;
\item memoria microscópica;
\item shell de frecuencias nuevas;
\item memoria no heredada.
\end{itemize}
No son seis energías añadidas a un balance previo. Son seis filas ortogonales de un único operador terminal de pérdida. La identidad de Parseval establece\par
\[
\|L_{M,j,T}^{\rm term}x\|^2
=
\sum_{\mathfrak a}
\|L_{M,j,T}^{\mathfrak a}x\|^2.
\]
La ortogonalidad no es meramente nominal. Las profundidades ocupan intervalos temporales disjuntos; carry y shell viven en rangos espectrales distintos; la microestructura pertenece al complemento \(Q_{\rm micro}\); la memoria nueva se separa de la memoria heredada; y las dos hojas de la advección no pueden ser positivas simultáneamente en un mismo punto.\par
Por eso la contabilidad no duplica la interacción. La suma de las seis normas es exactamente la norma del operador común. Si una de ellas no fuese una proyección de ese operador, Parseval fallaría antes de obtener la cota. El propio formalismo contiene así su falsador.\par
\Needspace{6\baselineskip}
\subsubsection{Persistencia del objeto terminal bajo el paso al límite}
Toda truncación conserva algo que todavía no ha sido publicado. Ese resto es el terminal. Contiene la energía visible al final del intervalo y la memoria orientada que permanece entre el tiempo observado y el horizonte \(T\).\par
La columna terminal se construye primero. Después se forma su Gram. No se define una métrica deseada y se inventa posteriormente un operador que la realice. La columna contiene materialmente la salida terminal y todas las pérdidas transportadas:\par
\[
\mathbf G_{M,j;J,T}^{\rm term}
=
\begin{bmatrix}
\text{terminal transportado}\\
\text{pérdida}_{J-1}\\
\vdots\\
\text{pérdida}_{j}
\end{bmatrix}.
\]
Su recursión por bloques produce la identidad de Stein\par
\[
U_{M,j,T}
=
\Gamma_{M,j,T}^{*}
U_{M,j+1,T}
\Gamma_{M,j,T}
+
(L_{M,j,T}^{\rm term})^{*}
L_{M,j,T}^{\rm term}.
\]
Esta identidad no define retroactivamente la norma del estado. Se obtiene multiplicando una columna ya construida por su adjunta. Por la misma razón, la telescopía global es una identidad entre operadores existentes, no una hipótesis de estabilidad.\par
La memoria futura tampoco se introduce en la raíz inicial. Reside en la salida terminal. Al tomar la estimación integral en \(\tau=T\), su soporte temporal se vacía de manera exacta. No se borra por conveniencia.\par
\Needspace{6\baselineskip}
\subsubsection{Cota uniforme previa y compacidad de Aubin–Lions}
La raíz, el codificador de nivel, la columna y la isometría común satisfacen la factorización derivada desde los datos de partida\par
\[
\mathcal G_{M,J,T}^{\rm sol}
\Lambda_{M,T}
J_T^{\rm sol}d
=
I_{M,J}J_T^{\rm sol}d,
\qquad
I_{M,J}^{*}I_{M,J}=I.
\]
El lado izquierdo es la historia completa en el corte \((M,J)\). El derecho depende de la única raíz construida desde los datos. Como \(I_{M,J}\) es isométrico, la norma de la historia queda dominada por la norma fuente con una constante independiente de \(M\) y \(J\).\par
Éste es el paso que cierra la dificultad. La uniformidad no procede de una estimación de dimensión finita cuya constante crece con el corte. Tampoco procede de definir el espacio de historia mediante la misma pérdida que se pretende controlar. La raíz existe antes del corte y todos los cortes son realizaciones isométricas de ella.\par
La columna terminal traduce esa cota estructural a las cantidades físicas:\par
\[
\begin{aligned}
&\sup_M\sup_{0\le t\le T}\|u_M(t)\|_{H^s}^2
+\nu\sup_M\int_0^T\|u_M(t)\|_{H^{s+1}}^2\,dt\\
&\quad
+\sup_M\int_0^T|\mathscr B_{s,M}(u_M(t))|\,dt
+\text{carry}+\text{micro}+\text{shell}+\text{memoria}
\le C_T(u_0,f).
\end{aligned}
\]
La constante es finita para todo \(T<\infty\) y no depende de \(M\) ni de \(J\). La desigualdad no controla sólo el campo visible. Controla simultáneamente el trabajo transportado y los sectores que una publicación clásica suele olvidar.\par
\Needspace{6\baselineskip}
\subsubsection{Obstrucción de signo en la factorización de Kalman–Yakubovich–Popov}
Hacer depender el cierre global de una factorización KYP con un signo incorrecto introduce un cuadrado positivo y pretende recuperarlo después con signo negativo. Esa igualdad es algebraicamente falsa.\par
La demostración excluye esa identidad y construye directamente la estimación uniforme necesaria.\par
La carta KYP que se conserva es una identidad finita correcta bajo sus propias condiciones. Funciona como control auxiliar de cada corte, pero no proporciona la uniformidad al retirar \(M\) y \(J\). La prueba global no usa una cota de \(Q^{-1}\), no toma el límite dentro de la KYP y no depende de aquel signo.\par
Por tanto, la obstrucción del bloque KYP con signo incorrecto no alcanza la demostración. La uniformidad queda establecida por la construcción directa desde los datos con su columna terminal y su identidad de Stein.\par
\Needspace{6\baselineskip}
\subsubsection{Convergencia desde la cota uniforme hasta la solución clásica global}
La cota anterior proporciona\par
\[
u_M
\quad\text{acotado en}\quad
L^\infty(0,T;H^s)
\cap L^2(0,T;H^{s+1}).
\]
Como \(s>5/2\), el producto de Sobolev controla la no linealidad:\par
\[
B:H^s_\sigma\times H^s_\sigma
\longrightarrow H^{s-1}_\sigma.
\]
La ecuación Galerkin proporciona además una cota uniforme de \(\partial_tu_M\) en \(L^2(0,T;H^{s-1})\). Las inclusiones\par
\[
H^{s+1}\Subset H^s\hookrightarrow H^{s-1}
\]
permiten aplicar Aubin–Lions. Se obtiene convergencia fuerte en \(L^2H^s\), suficiente para transportar el término bilineal:\par
\[
B_M(u_M,u_M)\longrightarrow B(u,u)
\quad\text{en}\quad
L^1(0,T;H^{s-1}).
\]
El límite satisface la ecuación incompresible. La presión no se introduce como otro dato oculto; se reconstruye después mediante\par
\[
-\Delta p=\partial_i\partial_j(u_i u_j),
\qquad
\int_{\mathbb T^3}p=0.
\]
La unicidad se obtiene sobre la diferencia de dos soluciones. Una estimación en \(H^{s-1}\), seguida de Grönwall, fuerza que la diferencia sea cero. Esto elimina la dependencia de la subsucesión extraída por compacidad.\par
La construcción vale para todo horizonte finito. Las soluciones obtenidas en \([0,n]\) y \([0,n+1]\) coinciden en su solapamiento por unicidad. Se pegan, por tanto, en una sola solución definida sobre \([0,\infty)\).\par
La suavidad se obtiene en toda la escala \(s+n\), usando los mismos lectores
estructurales; la unicidad identifica todos los límites y produce una solución
global suave para todo dato periódico suave. La cota requerida se construye
en cada horizonte finito y en cada índice de Sobolev a partir del dato
solenoidal completo, como demuestra el propietario
\cref{thm:pdf2-ns-global}.\par
Finalmente, la condición de media cero es una normalización del sector
solenoidal. El cambio de Galileo dependiente del tiempo separa la media y
transporta biyectivamente todo dato periódico suave al sector de media cero,
preservando divergencia, suavidad y normas de Sobolev.\par
\Needspace{6\baselineskip}
\subsubsection{Existencia, regularidad y unicidad global de Navier--Stokes}
La característica solenoidal del continuo, publicada mediante su
entrelazador específico, produce una única solución global suave de las
ecuaciones tridimensionales periódicas de Navier--Stokes para todo dato
periódico suave. La construcción de Leray--Hopf queda contenida como etapa de
energía dentro de esa conclusión fuerte.\par
La defensa de esta conclusión no descansa en decir que existe una estructura capaz de resolver el problema. Descansa en cinco hechos materiales:\par
la raíz usa sólo \((u_0,f)\); la no linealidad se reproduce mediante creación mixta; las seis pérdidas son ortogonales y no se suman dos veces; el terminal permanece hasta la telescopía; y la cota uniforme precede al argumento clásico de compacidad.\par
Aubin–Lions, la presión, Grönwall y el pegado no son la fuente de la regularidad. Son el reconocimiento clásico final de una regularidad que la estructura ya ha impedido perder.\par
\Needspace{6\baselineskip}
\subsection{3.ª resolución: Yang–Mills y la brecha del sector de curvatura}
El problema de Yang–Mills exige dos cosas inseparables: construir una teoría cuántica no trivial sobre \(\mathbb R^4\) para todo grupo compacto, conectado y simple del dominio considerado, y demostrar que su Hamiltoniano posee una brecha de masa positiva.\par
No basta con exhibir un operador finito con espectro positivo. Tampoco basta con escribir una acción de Wilson, una medida producto o un observable con un autovalor no nulo. La brecha debe pertenecer a la misma teoría gauge que posee localidad, positividad por reflexión, covariancia euclídea, reconstrucción de Wightman y curvatura física.\par
La característica gauge del continuo conserva ya la red refinable, la medida, las cuatro reflexiones, la holonomía, el color, la incidencia, el carry de fusión, la ruta, la memoria, el terminal y el lector de curvatura. El grupo \(G\) y el acoplamiento \(g\) entran después, mediante el entrelazador específico. No seleccionan retrospectivamente la genealogía.\par
\Needspace{6\baselineskip}
\subsubsection{Inequivalencia entre la fibra simple y la dinámica producto}
La primera distinción defensiva separa dos generadores que una notación antigua podía confundir.\par
El semigrupo de una sola fibra posee únicamente los niveles \(0\) y \(3\kappa_{\rm av}\). El generador producto completo integra, en cambio, todas las coordenadas independientes del estado:\par
\[
D_m^{\rm prod}
=
3\kappa_{\rm av}
\sum_{\iota\in\mathcal I_m}(I-E_{m,\iota}),
\]
donde las expectativas \(E_{m,\iota}\) son ortogonales y conmutantes.\par
Su espectro contiene\par
\[
0,\quad
3\kappa_{\rm av},\quad
6\kappa_{\rm av},\quad
9\kappa_{\rm av},\ldots
\]
según el número de coordenadas no constantes. Sustituir este generador por el de una fibra simple borraría las multiplicidades de Hoeffding y sería inmediatamente falsable: un vector dependiente de dos factores debe tener energía \(6\kappa_{\rm av}\), no \(3\kappa_{\rm av}\).\par
La brecha se demostrará para el generador producto y después se transportará al sector físico. No se obtiene confundiendo dos dinámicas diferentes.\par
\Needspace{6\baselineskip}
\subsubsection{Ausencia de coordenadas ocultas y exclusión de modos nulos espurios}
El estado completo se factoriza mediante un índice exhaustivo. En él aparecen una sola vez:\par
\begin{itemize}
\item marcos;
\item enlaces;
\item incidencias;
\item marcas;
\item puentes de refinamiento;
\item innovaciones nuevas de cada nivel.
\end{itemize}
Las subcolas infinitas se registran mediante sus cilindros escalares. Ninguna coordenada genealógica queda fuera de la familia de expectativas.\par
La descomposición de Hoeffding escribe el espacio como suma ortogonal de sectores \(\mathcal H_{m,S}\). Sobre cada uno,\par
\[
D_m^{\rm prod}u_S
=
3\kappa_{\rm av}|S|u_S.
\]
El sector \(S=\varnothing\) es la intersección de los rangos de todas las expectativas y está formado exactamente por las constantes:\par
\[
\bigcap_{\iota\in\mathcal I_m}
\operatorname{Ran}E_{m,\iota}
=
\mathbb C\Omega_m.
\]
Por tanto, el vacío es simple y todo vector ortogonal al vacío tiene energía al menos \(3\kappa_{\rm av}\).\par
Esto excluye la objeción de que la brecha provenga de no haber enumerado una coordenada. Si faltara un factor, existirían funciones no constantes invisibles para todas las expectativas y aparecerían cero-modos adicionales. La exhaustividad del índice impide exactamente ese fallo.\par
\Needspace{6\baselineskip}
\subsubsection{Realización observable de la brecha espectral}
Una cota inferior no basta para fijar el valor exacto de la brecha. Se necesita un testigo que alcance el umbral.\par
En cada collar de incidencia, la variable\par
\[
W_c(q)
=
\mathbf1_{\{q_1+\cdots+q_6=0\}}
-\frac13
\]
tiene media cero y satisface\par
\[
D_m^{\rm prod}W_c
=
3\kappa_{\rm av}W_c.
\]
Pero todavía debe demostrarse que \(W_c\) pertenece al sector gauge físico.\par
La acción de incidencia se escribe mediante la matriz\par
\[
(B^+x)_{ij}=x_i+x_j.
\]
Cada componente de \(x\in\mathbb F_3^4\) aparece tres veces al sumar las seis incidencias. En \(\mathbb F_3\),\par
\[
\sum_{i<j}(B^+x)_{ij}
=
3\sum_i x_i
=
0.
\]
Por tanto, la condición que define \(W_c\) es gauge invariante. El promedio sobre el grupo gauge completo conserva \(W_c\); no lo expulsa del espacio físico. Así, \(W_c\) es un autovector físico y el valor \(3\kappa_{\rm av}\) no es sólo una cota: es el primer nivel espectral alcanzado.\par
\Needspace{6\baselineskip}
\subsubsection{Construcción de la medida previa a la estimación de la brecha}
Una arista gruesa se refina en nueve incrementos ordenados. La medida fina no se inventa copiando nueve veces la medida gruesa. Se construye mediante el puente de calor condicionado por el producto de los nueve incrementos. Los tiempos de los subtramos pueden ser distintos y satisfacen únicamente que su suma sea el tiempo grueso.\par
La convolución del kernel de calor garantiza que el puente es una probabilidad. Una carta triangular separa entonces la variable gruesa de las innovaciones nuevas. La proyección hacia el nivel anterior multiplica los nueve incrementos y olvida exclusivamente las innovaciones.\par
De esta manera,\par
\[
(\widehat\pi_{m+1,m})_*
\widehat\mu_{m+1}
=
\widehat\mu_m.
\]
Las aristas compartidas por varias caras continúan siendo una misma variable. No se sustituyen por copias independientes. Carry y memoria impiden volver a contar en el nivel fino la cara que ya estaba presente en el nivel grueso.\par
Sobre esa familia proyectiva se construye primero el lector exacto de \(F^2\), su raíz positiva, la acción local y la densidad de Gibbs dependiente del acoplamiento:\par
\[
d\widehat\mu_{m,g}^{\rm YM}
=
Z_{m,g}^{-1}
e^{-S_{m,g}}\,
d\widehat\mu_m^0.
\]
Las ecuaciones de Schwinger–Dyson y Ward se calculan para esa medida. La brecha no se usa para elegir \(g\), para definir \(S_{m,g}\) ni para reponderar la ley. Aparece después, al estudiar el generador reversible asociado a la misma medida.\par
Ésta es la defensa contra la circularidad acción–ley–brecha: primero se construye la acción y su medida; después se demuestra su espectro.\par
\Needspace{6\baselineskip}
\subsubsection{Descenso del generador al cociente de calibre}
El espacio completo incluye enlaces, marcos, incidencias, color, orientación, ruta, memoria y terminal. La marginal que conserva únicamente los enlaces es una ablación posterior. No puede utilizarse para negar un observable que vive en la fibra de incidencia.\par
La proyección física se obtiene promediando la acción del grupo gauge completo. Como el generador integra coordenadas de forma equivariante,\par
\[
[\mathbb P_m^{\rm phys},D_m^{\rm prod}]=0.
\]
El semigrupo desciende, por tanto, al cociente físico. Su kernel se define primero sobre el estado completo y sólo después sobre las clases gauge. No se toma un operador restringido a una fibra y se lo declara kernel sobre todo el espacio.\par
El vacío simple y el autovector \(W_c\) sobreviven a la compresión. La brecha finita pertenece ya al espacio gauge-invariante.\par
\Needspace{6\baselineskip}
\subsubsection{Unicidad del proceso reconstruido a partir de \(4\) positividades por reflexión}
La teoría euclídea debe satisfacer positividad por reflexión en las cuatro direcciones. No se construyen cuatro medidas independientes. Una única medida de caras, obtenida del mismo kernel reversible, factoriza cada forma de Osterwalder–Schrader:\par
\[
\int
\overline{F(\Theta_{\nu}y)}F(y)\,d\Omega
=
\|\mathcal B_\nu F\|^2
\ge0,
\qquad \nu=1,\ldots,4.
\]
La positividad no es una hipótesis añadida. Es una norma cuadrada producida por Markov y reversibilidad.\par
Las marginales de caras opuestas recuperan el mismo semigrupo. Por tanto, las cuatro reconstrucciones OS poseen un único origen dinámico. Los conectores entre semirredes demostrarán después que reconstruyen el mismo espacio, el mismo vacío y el mismo Hamiltoniano.\par
\Needspace{6\baselineskip}
\subsubsection{Persistencia de la brecha espectral bajo el límite continuo}
Los mapas de refinamiento entrelazan las formas de Dirichlet, los semigrupos, el vacío y el testigo \(W_c\). El paso al continuo no se basa únicamente en observar que la brecha vale lo mismo en cada nivel. Esa igualdad, aislada, no impediría que apareciesen nuevos estados de energía arbitrariamente pequeña en el límite.\par
La construcción controla los productos cruzados mediante entrelazadores de bloque. Las coordenadas gruesas se distribuyen sobre sus descendientes con la normalización exacta. Las identidades de Gram demuestran que las publicaciones suavizadas forman familias de Cauchy. De este modo, la forma, el vacío y el autovector alcanzan el límite dentro del mismo sistema inductivo.\par
Se obtiene\par
\[
D_{\infty,g}^{\rm YM}
\succeq
3\kappa_{\rm av}(I-P_{\Omega_\infty}),
\]
y el testigo límite satisface\par
\[
D_{\infty,g}^{\rm YM}W_c
=
3\kappa_{\rm av}W_c.
\]
\Needspace{5\baselineskip}
Por tanto,\par
\[
\Delta_{\rm YM}^{\rm HMT}
=
3\kappa_{\rm av}>0.
\]
La desigualdad y el testigo convergen juntos. La brecha no se deduce de una mera semicontinuidad inferior ni se extrapola desde una malla aislada.\par
\Needspace{6\baselineskip}
\subsubsection{Pertenencia del testigo espectral al sector de curvatura}
Todavía quedaría una objeción si \(W_c\) fuese sólo una excitación combinatoria de un factor interno, desacoplada de la curvatura Yang–Mills.\par
La holonomía de incidencia elimina esa posibilidad. La suma de las seis incidencias determina una holonomía local de orden tres; \(W_c\) es una función de clase de esa holonomía. Por tanto pertenece al álgebra local gauge-invariante generada por los lazos.\par
El lector \(F^2\) se polariza para producir un espacio de curvaturas locales. Los entrelazadores de refinamiento transportan esas curvaturas conservando el producto de \(\mathfrak g\). Las sumas celulares convergen a una distribución gauge-covariante\par
\[
\mathbf F_{\mu\nu}
\in
\mathcal S'(\mathbb R^4;\mathfrak g)
\widehat\otimes
L^2(\widehat X_\infty,\widehat\mu_{\infty,g}^{\rm YM}).
\]
El sector cíclico generado por \(\mathbf F\) es reductor para el Hamiltoniano y contiene el testigo \(W_c\). En consecuencia, la brecha exacta pertenece al sector físico de curvatura. No es el espectro de un ruido auxiliar.\par
\Needspace{6\baselineskip}
\subsubsection{Localidad euclídea \(E(4)\) y reconstrucción continua de Osterwalder–Schrader}
El generador se descompone en circuitos locales de bloques con soporte uniformemente acotado. Una coloración nonádica separa los collares que pueden actualizarse simultáneamente. El producto de los circuitos preserva la identidad de Bianchi en cada paso, no sólo al finalizar una barrida.\par
Las traslaciones y rotaciones euclídeas actúan por relabelado de la red, los marcos, las incidencias y las holonomías. La continuidad fuerte no se prueba únicamente para un primer caos abstracto. Se controla primero sobre observables locales, luego sobre holonomías, plaquetas y clover; el suavizado de órbitas extiende la acción a un núcleo denso.\par
El sistema cofinal produce una red local \(\mathfrak A_{\rm YM}(O)\), un vacío \(\Omega_{\infty,g}\), una representación fuerte de \(E(4)\) y funcionales de Schwinger. Las pilas Markov y las transferencias OS se identifican con el mismo semigrupo continuo. La positividad por reflexión, la regularidad, la covariancia y el cluster exponencial permiten la reconstrucción de Osterwalder–Schrader.\par
La continuación produce campos locales de Wightman y una representación de Poincaré con espectro en el cono futuro. Las cuatro reconstrucciones euclídeas coinducen el mismo Hamiltoniano de curvatura.\par
La brecha controla además el cluster:\par
\[
|S^{\rm YM}(A\,\tau_rB)|
\le
e^{-3\kappa_{\rm av}r}
\|A\Omega_\infty\|
\|B\Omega_\infty\|.
\]
Así, el valor espectral positivo posee una consecuencia física continua y local; no queda encerrado en el cálculo finito.\par
\Needspace{6\baselineskip}
\subsubsection{Reconstrucción de la acción clásica como salida del sistema euclídeo}
La densidad finita gobernante es el lector exacto de \(F^2\). El clover no la sustituye. Funciona como carta suave para reconocer su límite.\par
Para una conexión suave, la expansión de Magnus de la holonomía de una plaqueta empieza por\par
\[
\log\operatorname{Hol}_{p_{\mu\nu}}
=
a_m^2F_{\mu\nu}
+\text{término impar}
+O(a_m^4).
\]
El promedio sobre las cuatro orientaciones del clover cancela el término impar. La suma de Riemann da\par
\[
S_{m,g}(A)
\longrightarrow
\frac1{4g^2}
\int_{\mathbb R^4}\|F_A(x)\|^2\,d^4x.
\]
Por tanto, la ley construida es reconocida como Yang–Mills porque su acción finita, sus holonomías y su curvatura convergen al funcional clásico correcto. No se define “teoría Yang–Mills” como cualquier teoría que posea el gap deseado.\par
La misma medida satisface Schwinger–Dyson y Ward. La relación entre pequeños lazos y curvatura, la identidad de Bianchi y la covariancia gauge se conservan en el límite. Sólo después de toda esta identificación se aplica la equivalencia espectral que transporta el gap.\par
\Needspace{6\baselineskip}
\subsubsection{Interpretación dimensional del parámetro de masa obtenido por el entrelazador terminal}
La brecha energética se publica dimensionalmente como\par
\[
m_{\rm gap}^{\rm HMT}
=
\frac{\hbar_{\rm int}}{c_{\rm int}}
\,3\kappa_{\rm av}>0.
\]
Este resultado no asigna una masa de Proca a un gluón elemental. La invariancia gauge permanece intacta. La masa corresponde al coste mínimo de excitar el sector colectivo gauge-invariante de curvatura por encima del vacío.\par
Esta distinción no rebaja el resultado: es exactamente la formulación requerida por el problema de Yang–Mills. La brecha pertenece al espectro físico de estados gauge-invariantes, no a un término de masa elemental añadido al lagrangiano.\par
\Needspace{6\baselineskip}
\subsubsection{Existencia de una brecha espectral positiva y reconstrucción masiva del sector de Yang–Mills}
Para todo grupo de Lie compacto, conectado y simple, la característica gauge
del continuo, seguida de su entrelazador específico, construye una teoría
cuántica no trivial de Yang--Mills sobre \(\mathbb R^4\).\par
La teoría posee:\par
\begin{itemize}
\item una familia proyectiva de medidas locales dependientes del acoplamiento;
\item una acción exacta reconocida en el límite como \(\frac1{4g^2}\int\|F\|^2\);
\item un cociente gauge físico;
\item positividad por reflexión en cuatro direcciones;
\item una red local y una acción fuerte de \(E(4)\);
\item funciones de Schwinger y continuación de Wightman;
\item vacío simple;
\item ecuaciones de Schwinger–Dyson, Ward y Bianchi;
\item un sector cíclico de curvatura;
\item la brecha exacta
  \[
  \Delta_{\rm YM}^{\rm HMT}=3\kappa_{\rm av}>0.
  \]
\end{itemize}
La conclusión resiste las objeciones previsibles porque no identifica el generador de una fibra con el producto; no deja coordenadas fuera del índice; no confunde la marginal de enlaces con el cociente físico; no extrapola la brecha desde una malla; no coloca el testigo fuera del sector de curvatura; no presupone continuidad euclídea; y no usa el gap para construir la acción o la medida.\par
Navier–Stokes demuestra que el continuo puede transportar una no linealidad sin perder la memoria necesaria para impedir una singularidad. Yang–Mills demuestra que el mismo continuo puede transportar una libertad gauge sin perder la curvatura necesaria para separar el vacío de la primera excitación física.\par

\Needspace{6\baselineskip}
\section{Clausuras algebraicas del continuo y límites de la reconstrucción extensional}
Las cinco construcciones \texttt{sp}, \texttt{sol}, \texttt{gau}, \texttt{coh} y \texttt{det} no son ramas que se añadan al continuo para resolver cinco problemas conocidos. Son características simultáneas del mismo estado enriquecido. Hodge y Birch–Swinnerton–Dyer son las dos publicaciones finales de esta arquitectura: una convierte una historia cohomológica compatible en una clase algebraica; la otra demuestra que las publicaciones analítica y aritmética de una misma línea determinantal coinciden.\par
\Needspace{6\baselineskip}
\subsection{4.ª resolución: Hodge y la construcción efectiva de la preimagen algebraica}
El problema de Hodge pregunta si toda clase racional de tipo \((p,p)\),\par
\[
\alpha\in
H^{2p}(X,\mathbb Q)\cap H^{p,p}(X),
\]
para una variedad proyectiva lisa compleja \(X\), procede de una clase algebraica\par
\[
\beta_\alpha\in \operatorname{CH}^{p}(X)_{\mathbb Q}
\]
tal que\par
\[
\operatorname{cl}_{X,p}(\beta_\alpha)=\alpha.
\]
La dificultad no consiste en escribir esta última igualdad. Consiste en construir \(\beta_\alpha\) sin haber introducido antes, bajo otro nombre, la existencia de la preimagen que se quiere demostrar.\par
La característica cohomológica \texttt{coh} proporciona el aparato anterior a la entrada de \(X\), \(p\) y \(\alpha\): historias supervivientes, refinamientos, relaciones, orientación, carry, memoria, frontera, ruta, multisección, terminal y lector. El dato clásico no crea ese aparato ni selecciona retrospectivamente un subdominio. Sólo pregunta qué publica ese aparato cuando recibe las observaciones de \(\alpha\).\par
En cada profundidad \(N\) aparecen dos módulos:\par
\[
\mathscr S_N(X,p),
\]
que contiene los estados cohomológicos enriquecidos, y\par
\[
\mathscr H_N^p(X),
\]
que contiene las observaciones cohomológicas visibles en esa profundidad. El lector finito\par
\[
e_N:\mathscr S_N(X,p)\longrightarrow\mathscr H_N^p(X)
\]
está coordinado con los mapas de restricción:\par
\[
\rho_{N+1,N}:\mathscr S_{N+1}\longrightarrow\mathscr S_N,
\qquad
q_{N+1,N}:\mathscr H_{N+1}^p\longrightarrow\mathscr H_N^p,
\]
mediante el cuadrado\par
\[
e_N\rho_{N+1,N}
=
q_{N+1,N}e_{N+1}.
\]
Esto impide que una observación admitida en una profundidad sea incompatible con su sombra anterior.\par
Para la clase \(\alpha\), la fibra finita que debe poblarse es\par
\[
\mathscr F_N(\alpha)
=
\{s\in\mathscr S_N(X,p):e_N(s)=q_N\alpha\}.
\]
Definirla no demuestra que sea no vacía. Tampoco demuestra que un elemento de \(\mathscr F_N(\alpha)\) pueda prolongarse hasta \(\mathscr F_{N+1}(\alpha)\). La resolución consiste precisamente en construir ambas cosas.\par
\Needspace{6\baselineskip}
\subsubsection{Construcción de la corrección relativa previa a la clase de cohomología}
Al refinar de \(N\) a \(N+1\) aparecen dos núcleos relativos:\par
\[
\mathscr K_{N+1}^{\mathrm S}
=
\ker\rho_{N+1,N},
\qquad
\mathscr K_{N+1}^{\mathrm H}
=
\ker q_{N+1,N}.
\]
El primero contiene las correcciones nuevas del estado que desaparecen al volver al nivel anterior. El segundo contiene las observaciones nuevas que tampoco eran visibles en el nivel anterior.\par
Las cartas nuevas, los cruces de Mayer–Vietoris, las orientaciones, el carry, la memoria, la frontera, la ruta y las coordenadas terminales producen una presentación finita. Sobre sus generadores se construyen dos mapas:\par
\[
a_{N+1}:
\mathscr A_{N+1}^{\mathrm{rel}}
\longrightarrow
\mathscr K_{N+1}^{\mathrm S},
\]
\[
b_{N+1}:
\mathscr A_{N+1}^{\mathrm{rel}}
\longrightarrow
\mathscr K_{N+1}^{\mathrm H},
\]
satisfaciendo\par
\[
e_{N+1}^{\mathrm{rel}}a_{N+1}=b_{N+1}.
\]
El mapa \(a_{N+1}\) produce una corrección del estado; \(b_{N+1}\) publica exactamente su corrección cohomológica.\par
Las observaciones relativas se ordenan por profundidad, soporte, hoja, ruta, carry, memoria y terminal. Con ese orden, la matriz de \(b_{N+1}\) es unitriangular:\par
\[
b_{N+1}(c_\mu)
=
\bar c_\mu+
\sum_{\nu<\mu}
B_{\nu\mu}\bar c_\nu.
\]
La diagonal unitaria no se elige para representar una clase concreta. Procede de la incidencia principal de cada generador; los términos inferiores quedan impuestos por las relaciones, la orientación y el carry. La tabla se construye sobre todas las observaciones relativas antes de introducir \(\alpha\).\par
La eliminación triangular produce explícitamente un inverso derecho:\par
\[
\sigma_{N+1}^{\mathrm{rel}}(\bar c_{\mu_0})=c_{\mu_0},
\]
\[
\sigma_{N+1}^{\mathrm{rel}}(\bar c_\mu)
=
c_\mu-
\sum_{\nu<\mu}
B_{\nu\mu}
\sigma_{N+1}^{\mathrm{rel}}(\bar c_\nu),
\]
y, por inducción,\par
\[
b_{N+1}\sigma_{N+1}^{\mathrm{rel}}=I.
\]
\Needspace{5\baselineskip}
Por tanto,\par
\[
s_{N+1}^{\mathrm{rel}}
=
a_{N+1}\sigma_{N+1}^{\mathrm{rel}}
\]
satisface\par
\[
e_{N+1}^{\mathrm{rel}}
s_{N+1}^{\mathrm{rel}}
=I.
\]
Éste es el cortafuegos contra la circularidad. La sección relativa no se supone, no se deduce de que la clase deseada ya sea algebraica y no se construye específicamente para \(\alpha\). Existe universalmente para toda observación relativa.\par
Tampoco se necesita que la presentación sea inyectiva. El núcleo puede contener expresiones celulares redundantes. La prueba requiere sobreyectividad hacia las observaciones relativas, no unicidad de la expresión. La desaparición del cokernel es consecuencia del inverso derecho; no es una condición gobernante introducida antes de la prueba.\par
\Needspace{6\baselineskip}
\subsubsection{El operador universal de extensión}
La estructura posee una degeneración unitaria\par
\[
\iota_{N+1,N}:\mathscr S_N\longrightarrow\mathscr S_{N+1},
\qquad
\rho_{N+1,N}\iota_{N+1,N}=I.
\]
Sea \(s\in\mathscr S_N\) y sea \(h\in\mathscr H_{N+1}^p\) una observación compatible:\par
\[
e_N(s)=q_{N+1,N}(h).
\]
El defecto de la prolongación trivial es\par
\[
\delta_{N+1}(s,h)
=
h-e_{N+1}\iota_{N+1,N}(s).
\]
La compatibilidad garantiza que\par
\[
\delta_{N+1}(s,h)
\in
\mathscr K_{N+1}^{\mathrm H}.
\]
Se define entonces\par
\[
\operatorname{Ext}_{N+1,N}(s,h)
=
\iota_{N+1,N}(s)
+
s_{N+1}^{\mathrm{rel}}
\bigl(\delta_{N+1}(s,h)\bigr).
\]
Por construcción:\par
\[
\rho_{N+1,N}
\operatorname{Ext}_{N+1,N}(s,h)
=s,
\]
\[
e_{N+1}
\operatorname{Ext}_{N+1,N}(s,h)
=h.
\]
La primera identidad conserva la historia anterior. La segunda realiza exactamente la nueva observación. La naturalidad de la tabla relativa garantiza que orientación, carry, memoria, ruta, frontera y terminal no se pierdan al refinar.\par
El operador funciona para todo par compatible \((s,h)\). No contiene una clase algebraica ni un ciclo elegido. Ésa es la razón precisa por la que su aplicación posterior a una clase de Hodge no presupone Hodge.\par
\Needspace{6\baselineskip}
\subsubsection{Construcción nivel a nivel de la historia global en el sistema inverso}
En profundidad inicial existe una sección basada\par
\[
s_0^{\mathrm{base}}:
\mathscr H_0^p\longrightarrow\mathscr S_0.
\]
Sus imágenes son estados semilla, no ciclos algebraicos elegidos para representar \(\alpha\). Sólo después se fija\par
\[
s_0=s_0^{\mathrm{base}}(q_0\alpha)
\]
y se define recursivamente\par
\[
s_{N+1}
=
\operatorname{Ext}_{N+1,N}
\bigl(s_N,q_{N+1}\alpha\bigr).
\]
Las identidades del operador universal prueban por inducción:\par
\[
e_Ns_N=q_N\alpha,
\qquad
\rho_{N+1,N}s_{N+1}=s_N.
\]
Por tanto, cada fibra \(\mathscr F_N(\alpha)\) es no vacía y la familia\par
\[
\xi_\alpha=(s_N)_{N\ge0}
\]
es una historia compatible:\par
\[
\xi_\alpha
\in
\varprojlim_N\mathscr F_N(\alpha).
\]
No se invoca un límite inverso abstracto para ocultar la falta de secciones. La secuencia se exhibe recursivamente. Tampoco se utiliza la anulación de \(\varprojlim^1\) como sustituto de la construcción: la sobreyectividad y la condición de Mittag–Leffler aparecen después como consecuencias de un operador de extensión ya materializado.\par
\Needspace{6\baselineskip}
\subsubsection{Emergencia de la clase algebraica como límite de la construcción relativa}
Las observaciones finitas y la coordenada terminal separan las publicaciones del límite. Como \(\xi_\alpha\) reproduce \(q_N\alpha\) en toda profundidad,\par
\[
\operatorname{ev}_{\infty,X,p}^{\mathrm{Hdg}}
(\xi_\alpha)
=
\alpha.
\]
Sólo entonces actúa la publicación clásica:\par
\[
X_\chi(X,p)
\xrightarrow{R_{\mathrm{Hdg}}}
K_0\!\left(
\operatorname{Perf}^{\ge p}(X)
\right)_{\mathbb Q}
\xrightarrow{\operatorname{ch}^{\mathrm{loc}}_p}
\operatorname{CH}^{p}(X)_{\mathbb Q}
\xrightarrow{\operatorname{cl}_{X,p}}
\operatorname{Hdg}^{p}(X)_{\mathbb Q}.
\]
Se define\par
\[
\beta_\alpha
=
\operatorname{ch}^{\mathrm{loc}}_p
\bigl(R_{\mathrm{Hdg}}(\xi_\alpha)\bigr).
\]
\Needspace{9\baselineskip}
La identidad de acción del lector,\par
\[
\operatorname{cl}_{X,p}
\operatorname{ch}^{\mathrm{loc}}_p
R_{\mathrm{Hdg}}
=
\operatorname{ev}_{\infty,X,p}^{\mathrm{Hdg}},
\]
da finalmente\par
\[
\operatorname{cl}_{X,p}(\beta_\alpha)
=
\alpha.
\]
Esta identidad del lector no prueba por sí sola que exista \(\xi_\alpha\). La existencia fue construida antes mediante la tabla unitriangular, el inverso derecho, el operador universal y la recursión. Mantener separadas esas dos etapas es justamente lo que impide que la demostración sea circular.\par
Para toda clase racional de tipo \((p,p)\), la característica cohomológica
construye la sección compatible \(\xi_\alpha\), el ciclo
\(\beta_\alpha\in\operatorname{CH}^p(X)_{\mathbb Q}\) y la identidad
\(\operatorname{cl}_{X,p}(\beta_\alpha)=\alpha\), según
\cref{thm:pdf2-hodge-extension-universal,thm:pdf2-hodge-completitud-final}.
La exhaustividad celular y FIN1 forman parte de la construcción probada del
inverso derecho y de su prolongación coinductiva; no son hipótesis externas de
promoción.\par
\Needspace{6\baselineskip}
\subsection{5.ª resolución: BSD y la coincidencia de las publicaciones analítica y aritmética}
La resolución de Birch–Swinnerton–Dyer no comienza intentando adivinar el orden de anulación de \(L(E,s)\), ni reuniendo retrospectivamente rango, regulador, grupo de Tate–Shafarevich y factores de Tamagawa.\par
La característica determinantal \texttt{det} ya conserva historia, producto fibrado, unidad común, orientación, carry, memoria, supervivencia, terminal, retículo, dualidad, raíz y publicaciones compatibles. Sólo después entra una curva elíptica \(E/\mathbb Q\).\par
\Needspace{6\baselineskip}
\subsubsection{Complejo integral global y descomposición en fibras primarias}
La construcción se realiza sobre un complejo libre integral de historias. Para cada primo \(\ell\) y cada potencia \(\ell^n\), los coeficientes \(E[\ell^n]\), las transiciones, Alexander–Whitney, el pullback de Manin, la unidad y las publicaciones se obtienen por cambio de base del mismo complejo integral.\par
La fibra ternaria no gobierna a las restantes. No se inventa una falsa extensión de \(\mathbb Z_3\) a \(\mathbb Z_\ell\). Esta precaución es imprescindible para que la conclusión comprenda toda \(\Sha(E)\) y todos los factores de Tamagawa, no únicamente su componente \(3\)-primaria.\par
La unidad común produce las publicaciones de Kato y Poitou–Tate. En cada bloque finita–singular se construye directamente la homotopía de relleno\par
\[
h_*=-vf,
\]
cuya frontera es la diferencia entre ambas publicaciones.\par
Las historias se ordenan usando primo, nivel, orientación, carry, memoria, frontera, terminal y localización. Se obtiene así una descomposición basada exhaustiva en:\par
\begin{itemize}
\item unidad raíz;
\item bloques finita–singular;
\item retículo local–global.
\end{itemize}
No queda una clase residual sin tipar.\par
La homotopía \(H_{\mathrm{FS}}\) no entra como premisa. En cada bloque normal se define su acción y se comprueba\par
\[
\partial H_\lambda+H_\lambda\partial=I.
\]
La suma transportada por la descomposición basada satisface\par
\[
\partial H_{\mathrm{FS}}
+
H_{\mathrm{FS}}\partial
=
p_{\mathrm{FS}}.
\]
El miembro derecho es el proyector sobre el sector finita–singular. No borra el retículo que conserva \(\Sha\) y Tamagawa. La homotopía reproduce la homotopía de relleno previamente construida; no se postula para hacer desaparecer una obstrucción.\par
\Needspace{6\baselineskip}
\subsubsection{Finitud del grupo aritmético previa a la evaluación del término principal}
El retículo residual posee una matriz integral \(D_E\). Antes de tomar su cokernel se construye un inverso racional mediante la parte diagonal terminal y la serie geométrica finita de la parte estrictamente triangular. Se verifican las dos composiciones.\par
\Needspace{5\baselineskip}
Por tanto,\par
\[
D_E\otimes\mathbb Q
\]
es un isomorfismo y\par
\[
F_E=\operatorname{coker}D_E
\]
es finito. Esta finitud se obtiene antes de usar rango, \(\Sha\), regulador o término principal.\par
La forma de Smith conserva todos los divisores elementales y recompone exactamente sus componentes primarias. No se deduce la finitud de \(\Sha\) de la fórmula final: se deriva posteriormente de su inclusión en un grupo finito ya construido.\par
\Needspace{6\baselineskip}
\subsubsection{Bockstein, Kato–FERS y Poitou–Tate}
El diferencial determinantal universal \(S_E(T)\) genera las páginas Bockstein. Su forma de Smith \(T\)-ádica determina\par
\[
d
=
\operatorname{ord}_{T=0}
\det S_E(T)
=
\sum_q \dim V_q.
\]
No se sustituyen las páginas por ese entero. Se conservan cada cociente, cada diferencial reducido, cada emparejamiento y cada jet. El coeficiente regular es\par
\[
R_{\mathrm{Boc}}^{\mathrm{der}}(S_E)
=
\left[
T^{-d}\det S_E(T)
\right]_{T=0}
\neq0.
\]
La equivalencia Kato–FERS se construye fila por fila. Alexander–Whitney corta cada historia en todos sus puntos. El corte terminal conserva el generador con peso cero; los demás cortes adquieren pesos positivos determinados por profundidad y carry. La tabla exhaustiva tiene la forma\par
\[
I+TB_i(T),
\]
con inversa convergente y determinante alternado igual a uno en \(T=0\). De este modo, la unidad, las hojas, la memoria, la orientación y el terminal se preservan antes de consultar ningún valor de \(L(E,s)\).\par
Poitou–Tate enlaza después las páginas Bockstein con Mordell–Weil. Los mapas de Kummer, publicación y pegado se escriben en ambos sentidos y satisfacen unidad y counidad. Son inversos antes de comparar dimensiones.\par
El desdoblamiento de jets conserva las \(q\) direcciones correspondientes a una fila de profundidad \(q\). No se pretende obtener \(q\) covectores de una única coordenada ni se completan dimensiones con identidades artificiales. El Gram conserva todos los cruces entre filas, profundidades y jets.\par
Sólo después se toman dimensiones:\par
\[
\operatorname{rank}E(\mathbb Q)
=
\sum_q\dim V_q
=
d.
\]
El rango y el orden determinantal coinciden porque ambos son publicaciones de las páginas construidas, no porque uno se haya definido copiando al otro.\par
\Needspace{6\baselineskip}
\subsubsection{Obtención del grupo de Tate–Shafarevich y de los factores de Tamagawa mediante la sucesión exacta global}
Del cono local reducido se construye una inyección\par
\[
\Sha(E)\hookrightarrow F_E.
\]
Los levantamientos estructurales de las componentes locales producen la proyección y su sección. La comprobación a nivel de cociclos da:\par
\[
0
\longrightarrow
\Sha(E)
\longrightarrow
F_E
\longrightarrow
\bigoplus_v\Phi_v(\mathbb F_v)
\longrightarrow
0.
\]
Como \(F_E\) ya era finito,\par
\[
|\Sha(E)|<\infty
\]
y\par
\[
|F_E|
=
|\Sha(E)|
\prod_v c_v.
\]
La exactitud no se deduce por comparación de cardinales. Primero se construyen los mapas y se identifican sus núcleos e imágenes; sólo después se toman órdenes.\par
\Needspace{6\baselineskip}
\subsubsection{Construcción independiente del miembro analítico de BSD}
La modularidad de \(E\) proporciona la forma nueva asociada. Su transformada de Mellin se divide en dos intervalos, y la ecuación funcional transforma la segunda mitad en una integral convergente con todas sus derivadas. Esto produce un germen holomorfo de \(L(E,s)\) alrededor de \(s=1\) sin suponer su orden de anulación.\par
En el semiplano de convergencia absoluta, los factores de Euler se realizan por operadores normales. Los truncados convergen en norma traza y su determinante de Fredholm recupera el producto de Euler. La reducción de Schur sobre las coordenadas terminales y un atlas de cartas determinantales transportan esa familia hasta el diferencial \(S_E\).\par
El resultado es una factorización construida:\par
\[
L(E,s)
=
u_E(s)
\det S_E(T_E(s)).
\]
\Needspace{7\baselineskip}
La carta satisface\par
\[
T_E(1)=0,
\qquad
T'_E(1)=1,
\]
y la unidad de transición satisface\par
\[
u_E(1)=1.
\]
Esta normalización se demuestra mediante las transiciones del atlas, el complemento Fredholm y los cambios de base. No se ajusta usando el rango, el regulador, \(\Sha\) ni el coeficiente final.\par
En consecuencia,\par
\[
L(E,s)
=
(s-1)^dU_E^{\mathrm{lead}}(s),
\]
con\par
\[
U_E^{\mathrm{lead}}(1)
=
R_{\mathrm{Boc}}^{\mathrm{der}}(S_E)
\neq0.
\]
\Needspace{5\baselineskip}
Por tanto,\par
\[
\operatorname{ord}_{s=1}L(E,s)=d.
\]
Combinándolo con Poitou–Tate:\par
\[
\operatorname{ord}_{s=1}L(E,s)
=
\operatorname{rank}E(\mathbb Q).
\]
\Needspace{6\baselineskip}
\subsubsection{Derivación interna del regulador a partir del emparejamiento de alturas}
La base integral de Mordell–Weil procede del sistema completo de jets. Sobre ella se construye la altura de Néron–Tate y su descomposición local–global.\par
El Gram contiene todos los cruces:\par
\[
G_E^{\mathrm{jet}}
=
\bigl(
\langle P_{q,a},P_{q',a'}\rangle_{\mathrm{NT}}
\bigr).
\]
No se supone diagonalidad entre profundidades ni se sustituyen direcciones adicionales por bloques identidad. Su determinante define\par
\[
\operatorname{Reg}(E)
=
\det G_E^{\mathrm{jet}}
>0.
\]
La extensión de escalares a \(\mathbb R\) se realiza antes de multiplicar por el regulador. Por ello el comparador entre las líneas Bockstein y Néron está bien tipado; no se intenta introducir un número real dentro de una línea todavía definida únicamente sobre \(\mathbb Q\).\par
\Needspace{6\baselineskip}
\subsubsection{Diagrama de \(4\) morfismos y evaluación del término principal}
La comparación final utiliza cuatro flechas de líneas determinantales:\par
\begin{enumerate}
\item Kato–FERS transporta la sección analítica.
\item Poitou–Tate la lleva hasta Mordell–Weil y su altura.
\item La sucesión exacta separa \(\Sha\) y Tamagawa.
\item La evaluación incorpora torsión, período de Néron y componentes locales.
\end{enumerate}
Cada flecha tiene dominio, codominio y acción sobre una base. Antes de introducir el elemento analítico se fijan dos trivializaciones independientes: la analítica y la de Néron.\par
La conmutatividad del cuadrado final procede de la conservación de las bases estructurales. No se impone igualando previamente los números que deben aparecer.\par
Aplicando la composición al elemento analítico se obtiene:\par
\[
\operatorname{ord}_{s=1}L(E,s)
=
\operatorname{rank}E(\mathbb Q)
=
r,
\]
y\par
\[
\frac{L^{(r)}(E,1)}
{r!\,\Omega_E}
=
\frac{
|\Sha(E)|
\operatorname{Reg}(E)
\prod_v c_v
}{
|E(\mathbb Q)_{\mathrm{tors}}|^2
}.
\]
Ésta es la fórmula completa de Birch--Swinnerton--Dyer para toda curva
elíptica \(E/\mathbb Q\), como establece
\cref{thm:pdf2-bsd-publicacion-final}. La familia compatible
\(\nu_{E,\ell,n}\) se construye dentro del complejo integral global y aporta
el comparador que enlaza sus fibras primarias; forma parte de la demostración
universal y no constituye una condición externa de promoción.\par
Su defensa contra la circularidad es estrictamente cronológica:\par
\begin{itemize}
\item el complejo integral se fija antes de separar primos;
\item la homotopía de relleno se construye antes de \(H_{\mathrm{FS}}\);
\item Smith se obtiene antes del rango;
\item Kato–FERS se verifica antes del determinante principal;
\item Poitou–Tate se construye antes de comparar dimensiones;
\item la exactitud \(\Sha\)–Tamagawa se prueba antes de tomar órdenes;
\item la continuación Mellin–Fredholm se construye antes del orden analítico;
\item la altura se construye antes del regulador;
\item las cuatro flechas se fijan antes de introducir el elemento analítico;
\item las trivializaciones se fijan antes de evaluar ambos lados.
\end{itemize}
Ningún elemento se normaliza usando el valor que después debe demostrar.\par
\Needspace{6\baselineskip}
\subsection{Factorización de los entrelazadores terminales RH–NS–YM–Hodge–BSD a través del continuo genealógico común}
Riemann, Navier–Stokes, Yang–Mills, Hodge y BSD no son cinco traducciones arbitrarias de un mismo vocabulario. Cada problema aparece cuando una publicación clásica olvida una parte diferente del estado enriquecido:\par
\begin{itemize}
\item Riemann olvida la coligación común que conserva simultáneamente los canales primo, gamma y polar.
\item Navier–Stokes olvida la columna causal completa de terminal, pérdidas, carry y memoria.
\item Yang–Mills olvida la coordinación entre incidencia, curvatura, ley, dinámica y sector físico.
\item Hodge conserva una clase cohomológica, pero no la historia compatible que produce su preimagen algebraica.
\item BSD separa la publicación analítica de la aritmética y pierde la línea determinantal común que las transporta.
\end{itemize}
La arquitectura HMT no fuerza cinco soluciones independientes. Conserva antes aquello que las formulaciones clásicas publican por separado:\par
\[
\begin{aligned}
\text{estado enriquecido}
&\longrightarrow \text{refinamientos compatibles}
\longrightarrow \text{supervivencia y terminal}\\
&\longrightarrow \text{continuo discreto}
\longrightarrow \text{publicaciones clásicas}.
\end{aligned}
\]
La razón por la que las soluciones funcionan no es una semejanza verbal. Es que, en los cinco casos, el teorema decisivo procede de una identidad de conservación material:\par
\begin{itemize}
\item conservación de energía o defecto;
\item conservación de historia bajo refinamiento;
\item conservación de la acción y del producto superficial;
\item conservación de la observación durante la extensión;
\item conservación de la unidad y de la base en líneas determinantales.
\end{itemize}
Por eso los problemas clásicos deben aparecer editorialmente al final. Primero se construye el objeto común; después HMT–MD desarrolla sus consecuencias; luego Mecánica Dimensional realiza sus transductores; sólo finalmente se leen las cinco conclusiones en sus lenguajes clásicos.\par
\Needspace{6\baselineskip}
\subsection{Distinción tipada entre los cinco cierres terminales y el corredor lógico ZFC--ciclo}
La incidencia con Zermelo--Fraenkel, ZFC y la construcción de von Neumann debe formularse con precisión y en el orden causal del corpus. Las cinco resoluciones terminales y el corredor lógico ZFC--ciclo son demostraciones distintas, asentadas sobre la misma estructura discreta del continuo: el alcance local de una de ellas no puede operar como negación transversal de la otra. Los cinco cierres establecen sus respectivos resultados clásicos; el corredor lógico específico conserva, además, la cadena\par
\[
\begin{aligned}
 \mathcal C_{\mathrm{cont}}^{\mathrm{disc}}
 &\longrightarrow\mathcal L_0\xrightarrow{\Phi}\mathcal L_1
 \longrightarrow\gamma_F\longrightarrow w(\gamma_F)\\
 &\longrightarrow\mathsf{NoEscape}_{24}
 \longrightarrow\mathsf{AntiDiag}_{369}\\
 &\longrightarrow\mathsf{G\ddot{o}delIII}_{\mathrm{HMT}}
 \longrightarrow\mathsf{CycleCriterion}_{\mathrm{ZFC}},
\end{aligned}
\]
y publica el criterio fuerte recuperado\par
\[
 \boxed{
 \neg\operatorname{Consis}(\mathrm{ZFC})
 \Longleftrightarrow
 \exists C\;\bigl(C\text{ es un ciclo fuera de banda y }|C|>24\bigr).}
\]
Esta equivalencia posee su propio dominio, sus codificadores y su prueba; no se deduce de los cinco teoremas terminales ni queda limitada por una observación local acerca de cualquiera de ellos. Del mismo modo, la resolución cardinal ya construida conserva simultáneamente\par
\[
 \mathsf{CH}_{\mathrm{HMT}},
 \qquad
 M_h=L[h]\models 2^{\aleph_0}=\aleph_1,
\]
sin convertir la recta real o la sentencia extensional en entradas del generador.\par

Una representación extensional formulada en el lenguaje de ZFC puede codificar los portadores conjuntistas de los objetos empleados:\par
\begin{itemize}
\item un estado finito es un conjunto;
\item una historia es una sucesión;
\item un árbol es un conjunto de palabras;
\item una transición es una función o relación;
\item un sistema inverso es una familia conjuntista;
\item su límite es un subconjunto de un producto.
\end{itemize}
Esa codificación posterior no contiene por sí sola la identidad genealógica del objeto ni gobierna los teoremas construidos antes de la publicación extensional. La insuficiencia aparece cuando el reducto extensional se trata como si agotara ruta, orientación, acarreo, memoria, frontera y terminal.\par
El axioma de extensionalidad afirma:\par
\[
\forall A\,\forall B
\left[
\forall x\,
(x\in A\leftrightarrow x\in B)
\Rightarrow A=B
\right].
\]
La afirmación es correcta para conjuntos. Dos conjuntos con los mismos elementos son el mismo conjunto.\par
Pero de ella no se sigue que una publicación extensional permita reconstruir la historia que fue olvidada al producirla.\par
Sea\par
\[
\operatorname{Pub}:
\mathscr X_{\mathrm{gen}}
\longrightarrow
\mathscr X_{\mathrm{ext}}
\]
\Needspace{6\baselineskip}
un lector que envía un estado genealógico a su valor extensional. Puede ocurrir que\par
\[
\operatorname{Pub}(x)
=
\operatorname{Pub}(y)
\]
mientras\par
\[
x\neq y
\]
como historias, porque poseen distinta ruta, orientación, carry, memoria, frontera o terminal.\par
El reducto identifica las dos imágenes en \(\mathscr X_{\mathrm{ext}}\), pero esa igualdad de codominio no implica \(x=y\) en el dominio enriquecido. El error de tipo consiste en convertir la igualdad de publicaciones en igualdad genealógica.\par
Puede definirse la relación\par
\[
x\sim_{\mathrm{ext}}y
\quad\Longleftrightarrow\quad
\operatorname{Pub}(x)
=
\operatorname{Pub}(y).
\]
Entonces, en el dominio correspondiente,\par
\[
\mathscr X_{\mathrm{ext}}
\simeq
\mathscr X_{\mathrm{gen}}/
{\sim_{\mathrm{ext}}}.
\]
La fibra\par
\[
\operatorname{Pub}^{-1}(a)
\]
contiene las genealogías que producen la misma sombra \(a\). Extensionalidad describe correctamente el cociente; no proporciona una sección canónica que reconstruya una historia concreta de cada fibra.\par
Ni siquiera el axioma de elección resuelve esta carencia. Elección puede seleccionar un representante de cada fibra bajo sus hipótesis, pero:\par
\begin{itemize}
\item no lo hace canónico;
\item no conserva necesariamente operaciones;
\item no garantiza naturalidad bajo refinamiento;
\item no reconstruye carry, memoria o terminal;
\item no demuestra que la selección corresponda a la genealogía física o matemática correcta.
\end{itemize}
Elegir un representante no equivale a recuperar su procedencia.\par
\Needspace{6\baselineskip}
\subsubsection{El ejemplo de los ordinales de von Neumann}
La construcción de von Neumann\par
\[
0=\varnothing,
\qquad
n+1=n\cup\{n\}
\]
es exacta. El ordinal\par
\[
n=\{0,1,\ldots,n-1\}
\]
codifica perfectamente su orden extensional.\par
\Needspace{6\baselineskip}
Pero el mismo número puede aparecer como salida de historias distintas:\par
\[
6=3+3=2\cdot3=7-1.
\]
Como ordinal de von Neumann, el resultado es el mismo conjunto. Eso es correcto: la aritmética necesita que las cuatro expresiones publiquen el mismo \(6\).\par
Lo que el ordinal desnudo no conserva es qué operación, ruta, acarreo u orientación produjo esa aparición concreta. Si esa procedencia importa para una aplicación posterior, debe mantenerse como parte del objeto enriquecido:\par
\[
\widehat n_x
=
\{(k,x_k):k<n\}.
\]
La proyección\par
\[
(k,x_k)\longmapsto k
\]
recupera exactamente el ordinal de von Neumann. HMT no elimina el ordinal; lo reconoce como publicación de una estructura que puede contener más información.\par
La jerarquía acumulativa \(V_\alpha\) conserva pertenencia y rango. No registra automáticamente una causalidad interna de construcción diferente de la codificación elegida. Reintroducir etiquetas, árboles o categorías permite transportar datos genealógicos adicionales, pero ese enriquecimiento confirma precisamente que la genealogía no estaba contenida en el valor extensional después de haberla olvidado.\par
\Needspace{6\baselineskip}
\subsubsection{La recta real como publicación posterior}
La misma distinción afecta al continuo. Un número real puede construirse como corte de Dedekind, clase de sucesiones de Cauchy u otro representante extensional equivalente. Esa equivalencia es matemáticamente correcta.\par
HMT añade una pregunta anterior: ¿qué sistema compatible de estados, refinamientos y terminales publica ese real?\par
La recta\par
\[
\mathbb R
\]
aparece después mediante una evaluación arquimediana. Conserva el valor publicado, pero no necesariamente la historia completa que lo produjo. Dos estados pueden compartir la misma evaluación real sin compartir genealogía.\par
Por eso la estructura discreta del continuo no debe introducir \(\mathbb R\) como entrada. Primero se construye la historia compatible; después se publica su sombra arquimediana.\par
El entrelazador con la cardinalidad clásica se construye de manera explícita.
La codificación por bloques proporciona el homeomorfismo
\(\beta:\Omega_\Gamma^{\rm cal}\simeq(\mathbb F_3^6)^{\mathbb N}\); sus
aproximantes racionales construyen \(\mathbb R_{\rm HMT}\), y el cociente por
igualdad de valor publica \(\mathbb R\). El código genealógico \(h\) induce
el buen orden y determina el modelo canónico \(M_h=L[h]\), dentro del cual
la construcción prueba la sentencia clásica
\(2^{\aleph_0}=\aleph_1\). Los resultados de independencia de Gödel y Cohen
clasifican el lenguaje extensional de ZFC después de olvidar \(h\); no
refutan la decisión producida por la estructura enriquecida que conserva el
selector.\par
\Needspace{6\baselineskip}
\subsubsection{Distinción entre identidad genealógica, equivalencia observacional e igualdad extensional}
La carencia es tipológica: una codificación conjuntista añadida puede transportar una genealogía ya construida, mientras el reducto extensional desnudo carece del mapa que produjo la publicación. Capacidad de codificación y suficiencia fundacional son afirmaciones diferentes.\par
Deben distinguirse tres nociones:\par
\[
\text{identidad genealógica},
\qquad
\text{equivalencia observacional},
\qquad
\text{igualdad extensional}.
\]
La primera exige igualdad de historia, ruta, memoria, carry y terminal. La segunda exige igualdad respecto de un lector determinado. La tercera exige igualdad del objeto publicado en el codominio conjuntista.\par
Confundirlas genera falsos problemas de reconstrucción. Una sombra puede ser exacta y, al mismo tiempo, insuficiente para recuperar el objeto que la produjo.\par
Ésta es la aportación fundacional que las cinco resoluciones hacen visible. Cada problema clásico trabajaba con una publicación correcta, pero empobrecida:\par
\begin{itemize}
\item una forma explícita sin toda su genealogía común;
\item un campo sin la memoria causal completa;
\item una teoría gauge sin la coordinación íntegra de su proceso;
\item una clase cohomológica sin su historia de extensión;
\item una función \(L\) separada de su línea aritmética.
\end{itemize}
La estructura HMT mantiene esas correlaciones antes de publicar las sombras. Por eso no necesita reconstruir a posteriori, dentro de cada problema, información que nunca fue destruida.\par
La conclusión fundacional es la siguiente:\par
\begin{quote}
El cierre extensional es una publicación posterior del cierre genealógico y no puede sustituirlo. HMT construye la estructura enriquecida antes de la recta real, demuestra \(\mathsf{CH}_{\mathrm{HMT}}\), realiza la clausura ordinal en \(L[h]\) y establece separadamente el corredor Gödel--HMT III y el criterio ZFC--ciclo. Una codificación extensional posterior conserva únicamente aquello que el morfismo de publicación declara; no revoca ninguno de esos resultados.\par
\end{quote}
\Needspace{6\baselineskip}
\subsection{Factorización de RH–NS–YM–Hodge–BSD como publicaciones terminales del continuo genealógico común}
Los cinco problemas aparecen así como cinco lecturas terminales de un único borde. No se resuelven porque una analogía los reúna, sino porque el continuo enriquecido conserva desde el principio las correlaciones que cada formulación clásica había separado.\par
La dependencia causal definitiva se factoriza como:\par
\[
\begin{aligned}
\mathrm{APP}
&\longrightarrow \mathrm{TRIT}
\longrightarrow \mathrm{TPK}
\longrightarrow \widehat x\\
&\longrightarrow \text{cinco características consustanciales}
\longrightarrow \text{terminal y límite}
\longrightarrow \mathfrak C_{\mathrm{cont}}^{\mathrm{disc}}
\end{aligned}
\]
y sólo después:\par
\[
\mathfrak C_{\mathrm{cont}}^{\mathrm{disc}}
\longrightarrow
\begin{cases}
\text{positividad de Weil y RH},\\
\text{regularidad global suave de Navier--Stokes para todo dato periódico suave},\\
\text{Yang--Mills y brecha de masa para todo grupo compacto, conectado y simple},\\
\text{algebraicidad de toda clase racional de Hodge},\\
\text{igualdad completa de BSD para toda curva elíptica sobre }\mathbb Q.
\end{cases}
\]
Las cinco líneas finales no son cinco ramas internas del continuo. Son cinco publicaciones clásicas posteriores de cinco capacidades inseparables que ya coexistían en el mismo estado.\par
Ese es el cierre que la formulación anterior no podía todavía formular: no estamos ante cinco problemas desconectados que casualmente admiten un mismo vocabulario, sino ante cinco pérdidas de información distintas corregidas por una sola arquitectura que se negó a destruirla.\par

% Empalme y operador conjunto congelados: se consumen por su fuente y huella
% originales, no por una paráfrasis posterior.
\begingroup
\renewcommand{\chapter}[2][]{}%
\chapter{Empalme material entre la estructura única y sus cinco resoluciones}
\label{chap:pdf2-generacion-simultanea-hmt}

\begingroup
\footnotesize
% CORTE_2_NO_DISGREGACION — bloque de empalme PDF1 -> PDF2
% Congelado después del PDF1 compilado el 2026-08-19.
% Owners:
% 04b  8ba02f246c18a3bd2ff83ffdb08874c816a7f59c688b6fab6af6eb2f9647071f
% 04b1 25d6dee3875cbdd3252f793212448c17c4b5b2d3c5d787e995598d18fa76d5cf
% 04b2 d5de9b5379d4f95ea3d35e63c9ccc2ab741b5c750c01409af6f13526073b874f
% 04b3 ac9ca321b610e2c0560a32ec6ae1555814c67c587e152096d0c48d6abaa0bd4b

\section*{Empalme matemático material PDF1--PDF2}

La entrada única es el estado APP--TRIT--TPK completo. Antes de toda
evaluación arquimediana se recorren
\[
 0\longrightarrow10,\qquad 0\longrightarrow1000,
 \qquad 0\longrightarrow\infty,
\]
con los nueve huecos internos \(1,\ldots,9\); dos posiciones producen APP,
TRIT orienta sus hojas y TPK prolonga extensión, fase, carry, memoria,
frontera, ruta y supervivencia. Ni el continuo ni \(\mathbb R\) son entradas.

Para una emisión superviviente
\[
 \mathfrak e_\alpha=
 (\alpha,\operatorname{Ext}_{\gamma_\alpha},\operatorname{or}(\alpha),
 \kappa(\alpha),\partial\alpha,\gamma_\alpha,
 m(\gamma_\alpha;m_0),\operatorname{Led}(\gamma_\alpha))
\]
las cinco características son acciones simultáneas sobre ese mismo
generador:
\begin{align*}
 \mathcal O_{\rm sp}(\mathfrak e_\alpha)
  &=(U_\alpha,\sigma_x,\sigma_y,A_\alpha,\eta_\alpha),\\
 A_\alpha&=P_yU_\alpha P_x+Q_yU_\alpha Q_x,
 &\eta_\alpha&=Q_yU_\alpha P_x+P_yU_\alpha Q_x,\\
 \mathcal O_{\rm sol}(\mathfrak e_\alpha)
  &=(b_\alpha,b_\alpha^+,b_\alpha^-;
  M^\pm\mapsto M^\pm+b_\alpha^\pm,
  D\mapsto D+b_\alpha,H\mapsto H+|b_\alpha|),\\
 \mathcal O_{\rm gau}(\mathfrak e_\alpha)
  &=\{(\operatorname{Cell}_{\rm TPK}(\alpha),q,B,W_c(q)):
      q\in Q_{\rm inc}(t\alpha)\},\\
 \mathcal O_{\rm coh}(\mathfrak e_\alpha)
  &=(I_y,J_y,\kappa_{y,N},\delta_y,C_{y,N},
      F_\alpha:C_{y,N}\to C_{x,N}),\\
 \mathcal O_{\rm det}(\mathfrak e_\alpha)
  &=(\Psi_\alpha,K_\alpha,\mathbf1_\alpha,z_\pm(\alpha),h_*(\alpha),
  \operatorname{Det}C_{y,N},H_*C_{y,N},\operatorname{Smith}C_{y,N},
  \operatorname{Boc}C_{y,N}).
\end{align*}
Estas expresiones son vistas dependientes de una sola acción
\(\mathcal O(\mathfrak e_\alpha)\); no definen cinco dominios, factores o
ramas. Comparten literalmente fuente, destino, ruta, ledger, orientación,
carry, memoria y frontera. Sus identidades de acoplamiento son
\[
 A_\alpha^*A_\alpha+\eta_\alpha^*\eta_\alpha=U_\alpha^*U_\alpha,
 \qquad E b_f^\pm=b_c^\pm+\kappa^{\rm can},
\]
\[
 Q_{\rm inc}=C^2(\operatorname{Cell}_{\rm TPK};\mathbf F_3),
 \qquad sB=0,\qquad W_c(q+Ba)=W_c(q),
\]
\[
 \partial h_*=z_--z_+,qquad
 K_{\beta\alpha}=K_\beta+\Psi_\beta K_\alpha.
\]

La monodromía es la conexión nonádica real:
\[
 \Gamma_{9,k}=\operatorname{Hol}_{C_9}(\nabla_{\rm TPK}),
 \qquad g(k+9)=g(k),
 \qquad q_{\rm mem}(k+9)=q_{\rm mem}(k)+1,
\]
sin retorno de estado. Para una historia
\(h=(\varepsilon_{k+1},\ldots,\varepsilon_N)\), las trece coordenadas se
generan por el fold
\[
 r_{k,N}(h)=
 \operatorname{Upd}_{\varepsilon_N}\circ\cdots\circ
 \operatorname{Upd}_{\varepsilon_{k+1}}(r^0_{k}),
\]
y el terminal devuelve toda la multisección, sin selector:
\[
 \mathfrak T_{k,N}(\widehat x_k)
 =\{r_{k,N}(h):h\in\mathscr H_{k,N}(\widehat x_k)\}.
\]
La naturalidad es una única igualdad sobre prefijos generados,
\[
 u_{\ell,k}^{\mathcal O}\mathcal O_{\le\ell}(\gamma)
 =\mathcal O_{\le k}(\gamma)
 =\operatorname{map}(\mathcal O)
   \bigl(u_{\ell,k}^{\mathfrak e}E_{\le\ell}(\gamma)\bigr),
\]
que pasa al sistema inverso de terminales. Sólo entonces se obtiene la
estructura discreta completa del continuo; la evaluación en \(\mathbb R\) es
una publicación posterior.

\subsection*{Mapa de consumo en PDF2}

\begin{center}
\small
\begin{tabularx}{\textwidth}{@{}p{.18\textwidth}X p{.30\textwidth}@{}}
\toprule
Símbolo consumido & Contenido estructural íntegro de PDF1 &
Entrelazador que construye PDF2 \\
\midrule
\(\mathfrak G_{\rm sp}\) &
\(U,\sigma,P,Q,A,\eta\), hojas, carry, terminal, naturalidad y lector
firmado interno. &
\(\Sigma_R^{\rm sp},J_\pm,K_R^{\rm sp},\Xi,P,B_W\) y la identificación
con la forma de Weil. \\
\(\mathfrak G_{\rm sol}\) &
\(b,b^\pm,\kappa^{\rm can}\), memoria, micro/shell ortogonales, columna
terminal, Gram, Stein y telescopía sin doble conteo. &
Mapa source-first \(J_T^{\rm sol}\), \(\Lambda_{M,T}\), \(I_{M,J}\),
factorización uniforme y cota sólo en datos. \\
\(\mathfrak G_{\rm gau}\) &
\(\operatorname{Cell}_{\rm TPK},Q_{\rm inc},B,W_c\), cuatro direcciones,
seis celdas, ocho caras, expectativas y terminal común. &
Interfaz gauge física, semigrupo requerido, sector físico, Yang--Mills y
brecha de masa. \\
\(\mathfrak G_{\rm coh}\) &
Puertos, \(\kappa,\delta,C\), bar/coend, carry, supervivencia y cono de
holonomía. &
\(R_{\rm Hdg}\), \(\mathrm{ch}_{\rm loc}\), clase algebraica y publicación
Hodge. \\
\(\mathfrak G_{\rm det}\) &
\(\Psi,K,z_\pm,h_*\), AW, línea determinante, homología, Smith y
Bockstein, conservados sin exigir aciclicidad como entrada. &
Kato--FERS, PT, comparador local--\newline global, Tamagawa y publicación BSD. \\
\bottomrule
\end{tabularx}
\end{center}

Los objetos de la tercera columna son target-specific: no deben retrotraerse
como entradas de PDF1. PDF2 los construye sobre el objeto estructural completo
de la segunda columna y conserva sus trece coordenadas, terminal y genealogía.

\endgroup

\section{Tipado unitario del objeto consumido}

El bloque congelado anterior no se añade a un constructor gráfico de cinco
componentes. Lo sustituye. Para cada profundidad \(k\), sea
\(\mathcal E_k\) el estado APP--TRIT--TPK enriquecido con todas sus emisiones
supervivientes, y sea \(\mathcal Z_k\) el codominio de la única acción
correlativa. La flecha rectora es
\begin{equation}\label{eq:pdf2-constructor-unico-operaciones}
 \mathcal O_k:\mathcal E_k\longrightarrow\mathcal Z_k,
 \qquad
 \widehat x_k\longmapsto
 \mathcal O\bigl(\{\mathfrak e_\alpha:
 \alpha\in\operatorname{Ch}_k(\widehat x_k)\}\bigr).
\end{equation}
Su valor conserva simultáneamente las coordenadas
\(\mathrm{sp},\mathrm{sol},\mathrm{gau},\mathrm{coh},\mathrm{det}\)
del mismo estado y del mismo ledger. No es una tupla de factores
independientes ni el grafo tautológico de cinco mapas previamente separados.

Los truncamientos del estado y de la acción se escriben
\[
 u^{\mathfrak e}_{\ell,k}:\mathcal E_\ell\to\mathcal E_k,
 \qquad
 u^{\mathcal O}_{\ell,k}:\mathcal Z_\ell\to\mathcal Z_k,
 \qquad \ell\ge k,
\]
y satisfacen una sola naturalidad,
\begin{equation}\label{eq:pdf2-naturalidad-accion-unica}
 u^{\mathcal O}_{\ell,k}\circ\mathcal O_\ell
 =\mathcal O_k\circ u^{\mathfrak e}_{\ell,k}.
\end{equation}
Por ello la acción pasa al límite sin seleccionar una historia:
\[
 \mathcal O_\infty:
 \mathcal C_{\rm cont}^{\rm disc}:=\varprojlim_k\mathcal E_k
 \longrightarrow
 \mathcal Z_\infty:=\varprojlim_k\mathcal Z_k.
\]
Las notaciones usadas en los cinco capítulos son vistas intrínsecas de ese
único valor:
\begin{equation}\label{eq:pdf2-operacion-infinita-tipificada}
 \mathfrak G_T^{\rm int}
 :=\operatorname{View}_T\!\left(
 \mathcal O_\infty(\mathcal C_{\rm cont}^{\rm disc})\right),
 \qquad
 T\in\{\mathrm{sp},\mathrm{sol},\mathrm{gau},\mathrm{coh},\mathrm{det}\}.
\end{equation}
Aquí \(\operatorname{View}_T\) sólo expone las coordenadas requeridas por una
aplicación; no crea una rama, un factor, un subdominio ni un selector. La
estructura de partida de cada resolución es siempre el valor completo de
\(\mathcal O_\infty\), con las otras cuatro características presentes.

\begin{falsifier}[Orden causal y no disgregación]
Falla el empalme si se usa \(\mathbb R\) o el continuo como entrada de
\(\mathcal O_k\), si se reemplaza \(\mathcal O_k\) por cinco constructores
independientes, si se selecciona una sola historia terminal o si una vista
\(\operatorname{View}_T\) se confunde con el objeto completo. En cualquiera
de esos casos la conclusión del capítulo de aplicación deja de transferirse.
\end{falsifier}

\endgroup
\paragraph{Descarga inmediata de las fronteras del empalme.}
Las marcas de interfaz que el empalme congelado formula como cotas u
obligaciones de la etapa terminal no son reservas del teorema final ni datos
añadidos desde los problemas clásicos. Quedan descargadas, en el mismo orden
causal y antes de cualquier capítulo consumidor, por los propietarios exactos
que siguen: éstos construyen desde el estado común las cinco publicaciones,
sus mapas, sus estimaciones y sus certificados. Las formulaciones clásicas
posteriores reconocen esas salidas; no las seleccionan ni las suministran.
% Propietarios completos recuperados. Su orden es vinculante: primero el
% objeto correlativo; después las construcciones que descargan las cinco
% bisagras; sólo entonces las vistas comunes y las publicaciones clásicas.
\begin{HMTScientificOwnerSubordinated}
% 04 — CINCO REALIZACIONES MATEMATICAS FUNDAMENTALES
% ==================================================
% Residencia narrativa sucesora. Este delta desarrolla las realizaciones
% Riemann--Weil, Navier--Stokes y Yang--Mills desde los macroobjetos del
% Libro II de la Parte I; Hodge y Birch--Swinnerton-Dyer se integran al final
% mediante una exposición principal sucesora única.

\ifdefined\HMTSucesoraStructureTrace
  \typeout{HMT-SUCESORA 04: realizaciones RH, Navier--Stokes y Yang--Mills}
\fi

La Parte II construyó, dentro de la estructura discreta del continuo, cinco
subestructuras asociadas a sendos canales de publicación. Esta parte recibe
cada subestructura ya
enriquecida por hoja, orientación, acarreo, frontera, ruta, memoria y
supervivencia; define el lector especializado; demuestra que dicho lector
entrelaza el refinamiento; y aplica, sólo al final, el criterio matemático o
físico que reconoce la salida.

Cada uno de los cinco capítulos establece una clausura demostrada en su
dominio a partir de la estructura de partida explícita que enumera. La
formulación convencional interviene después como identificación del
codominio y como contraste, mientras que la cadena demostrativa permanece
ordenada por los objetos construidos en las Partes I y II. Los controles
auxiliares posteriores no son aristas de esa cadena y carecen de capacidad
para condicionar o demover la clausura propietaria.

En los cinco teoremas desarrollados a continuación, el orden causal es
\begin{equation}\label{eq:realizaciones-orden-causal}
 \begin{gathered}
 \text{APP--TRIT--Topological Phase Kernel}
 \longrightarrow \text{estado enriquecido},\\
 \text{estado enriquecido}
 \longrightarrow \text{conexión nonádica conjunta}
 \longrightarrow \text{estructura especializada},\\
 \text{estructura especializada}
 \longrightarrow \text{identidad de publicación}
 \longrightarrow \text{criterio de reconocimiento}.
 \end{gathered}
\end{equation}
El último criterio no selecciona el estado fuente. Su función consiste en
identificar, en el vocabulario del problema, una propiedad que ya ha sido
obtenida mediante los mapas anteriores.

\section*{Principio transversal de transporte con memoria}

Sea \(T\) uno de los canales considerados y sea
\(X_{T,k}\) su espacio de estados enriquecidos a profundidad \(k\). La
conexión nonádica conjunta se escribe
\begin{equation}\label{eq:realizaciones-gamma9}
 \Gamma_{9,T}:X_{T,k}\dashrightarrow X_{T,k+9},
 \qquad g(k+9)=g(k),
 \qquad \widehat x_{k+9}\neq\widehat x_k
 \quad\text{en general}.
\end{equation}
La primera igualdad expresa el retorno de fase; la desigualdad conserva la
profundidad adquirida. Para un lector completo
\(\widehat R_T=(R_T,K_T)^{\mathsf t}\), la naturalidad de una flecha
\(\gamma:x\to y\) toma la forma
\begin{equation}\label{eq:realizaciones-lrdn}
 \widehat R_{T,y}H_T(\gamma)=
 \begin{pmatrix}
  A_{T,\gamma}&B_{T,\gamma}\\
  C_{T,\gamma}&D_{T,\gamma}
 \end{pmatrix}
 \widehat R_{T,x}.
\end{equation}
Su primera fila muestra que el defecto de la lectura visible se conserva:
\begin{equation}\label{eq:realizaciones-defecto-visible}
 R_{T,y}H_T(\gamma)-A_{T,\gamma}R_{T,x}
 =B_{T,\gamma}K_{T,x}.
\end{equation}
Por tanto, la compresión no equivale a un borrado. La parte retirada de la
coordenada visible reaparece en el canal de memoria.

El acarreo de \(R_T\), definido por
\[
 \operatorname{Carr}_{R_T}(\gamma)
 =R_{T,y}H_T(\gamma)-Y_T(\gamma)R_{T,x},
\]
obedece la identidad de composición
\begin{equation}\label{eq:realizaciones-cociclo-acarreo}
 \operatorname{Carr}_{R_T}(\gamma_2\gamma_1)
 =\operatorname{Carr}_{R_T}(\gamma_2)H_T(\gamma_1)
 +Y_T(\gamma_2)\operatorname{Carr}_{R_T}(\gamma_1).
\end{equation}
Ésta es la condición algebraica que permite refinar sin reinicializar la
historia.

La vuelta completa admite, además, una identidad de energía finita. Para
operadores \(S_{T,j}=S_{T,j}^*\), transiciones \(A_{T,j}\) y defectos
\(D_{T,j}\), supóngase
\[
 S_{T,j}\succeq0,
 \qquad
 S_{T,j}-A_{T,j}^*S_{T,j+1}A_{T,j}=D_{T,j}^*D_{T,j},
 \qquad j=0,\ldots,8.
\]
El cierre de la vuelta no es una identificación tácita: el transporte de
fase \(\Phi_T\) debe satisfacer
\begin{equation}\label{eq:realizaciones-cierre-forma}
 S_{T,9}=\Phi_T^*S_{T,0}\Phi_T.
\end{equation}
Si \(M_T=\Phi_TA_{T,8}\cdots A_{T,0}\) es la monodromía y
\(\mathcal O_{T,9}\) reúne los nueve defectos transportados, entonces, para
todo \(N\geq1\),
\begin{equation}\label{eq:realizaciones-telescopia-terminal}
 S_{T,0}
 =\sum_{q=0}^{N-1}
 M_T^{*q}\mathcal O_{T,9}^*\mathcal O_{T,9}M_T^q
 +M_T^{*N}S_{T,0}M_T^N.
\end{equation}
El último sumando es el terminal. Cada realización debe demostrar si dicho
terminal converge a cero, si sobrevive como raíz o si se transporta a otra
escala. Suprimirlo antes de esa demostración cambiaría el enunciado.

\begin{proposition}[Transporte especializado de una identidad conjunta]
\label{prop:realizaciones-transporte-especializado}
Sea \(\mathfrak M_T\) una estructura construida a partir de
\eqref{eq:realizaciones-gamma9}--\eqref{eq:realizaciones-telescopia-terminal}
y sea
\[
 \mathcal P_T:\mathfrak M_T\longrightarrow\mathcal Y_T
\]
un lector que entrelaza las restricciones, conserva
\eqref{eq:realizaciones-cociclo-acarreo} y transporta el terminal con su tipo.
Se exige además que \(\mathcal P_T\) sea isométrico sobre las formas que
publica, o, equivalentemente, que venga acompañado de una representación
positiva que transporte explícitamente cada cuadrado. Entonces toda identidad
finita de suma de cuadrados, balance de energía o semigrupo que defina
\(\mathfrak M_T\) pasa a \(\mathcal Y_T\). El paso al
límite es válido cuando la familia publicada posee una cota uniforme y el
terminal satisface la ley de convergencia declarada.
\end{proposition}

\begin{proof}
La conmutación con las restricciones identifica las publicaciones de dos
niveles comparables. La identidad de cierre
\eqref{eq:realizaciones-cierre-forma} convierte la suma de las nueve
identidades locales en una identidad de vuelta. La identidad
\eqref{eq:realizaciones-cociclo-acarreo} conserva, bajo composición, el
defecto que la proyección visible no registra. La isometría, o la
representación positiva declarada, hace que aplicar \(\mathcal P_T\) a
\eqref{eq:realizaciones-telescopia-terminal} conserve cada sumando positivo
y el terminal. La cota uniforme proporciona compacidad o convergencia en la
topología propia de \(\mathcal Y_T\); la ley declarada para el último
sumando permite identificar su límite. Ninguna de estas operaciones utiliza
la conclusión como entrada.
\end{proof}

Este principio es el antecedente común. En la clausura espectral el terminal
se extingue y queda un cuadrado de Hilbert; en la clausura solenoidal continúa
como energía del estado refinado y produce una telescopía causal; en la
clausura gauge la raíz se conserva como vacío simple mientras su complemento
mantiene una escala espectral positiva. La clausura cohomológica publica una
clase algebraica desde su historia basada y la clausura determinantal compara
dos publicaciones de una misma unidad orientada.

\section*{Principio functorial común de realización}
\label{sec:amp-principio-realizacion}

Las cinco especializaciones de esta Parte reciben antecedentes distintos de
una misma estructura de supervivencia. La unidad del método no consiste en
identificar sus codominios, sino en conservar la compatibilidad de los mapas
que los alcanzan. Sea
\[
 (\mathcal S_N^{\rm surv},\pi_{N+1,N})_{N\geq0}
\]
el sistema de estados supervivientes y sea
\((\mathscr C_N,\rho_{N+1,N})_N\) un sistema de categorías de destino.
Una familia de realizadores finitos
\begin{equation}\label{eq:amp-realizadores-finitos}
 F_N:\mathcal S_N^{\rm surv}\longrightarrow\mathscr C_N
\end{equation}
es compatible cuando hace conmutar
\begin{equation}\label{eq:amp-cuadrado-realizacion}
 \rho_{N+1,N}\circ F_{N+1}
 =F_N\circ\pi_{N+1,N},
\end{equation}
y entrelaza además orientación, transporte, acarreo y unidad.

\begin{theorem}[Realización compatible en el límite]
\label{thm:amp-principio-realizacion}
Las cinco familias construidas en esta Parte satisfacen
\eqref{eq:amp-cuadrado-realizacion} en toda profundidad, y sus morfismos de
destino conservan la identidad telescópica con su término terminal. En cada
familia existe, por tanto, un único realizador
\begin{equation}\label{eq:amp-realizador-limite}
 F_\infty=\varprojlim_NF_N:
 X_\chi=\varprojlim_N\mathcal S_N^{\rm surv}
 \longrightarrow\varprojlim_N\mathscr C_N
\end{equation}
compatible con todas las proyecciones. El mapa \(F_\infty\) transporta
conjuntamente el observable, la unidad y la memoria terminal del mismo
antecedente.
\end{theorem}

\begin{proof}
Para \(x=(x_N)_N\in X_\chi\), la familia
\((F_N(x_N))_N\) es compatible por
\eqref{eq:amp-cuadrado-realizacion}; la propiedad universal del límite
inverso define \(F_\infty(x)\) y prueba su unicidad. La naturalidad respecto
del transporte lleva la holonomía \(\Gamma_9\) a una holonomía del destino.
La conservación de la unidad lleva la clase persistente a la clase unital
correspondiente. Finalmente, si
\[
 S_0=\sum_{r=0}^{N-1}M^{*r}D^*DM^r+M^{*N}S_0M^N,
\]
la conservación demostrada para los morfismos de destino transporta tanto la suma como
el último sumando. El límite no confunde, por tanto, el observable con el
estado completo ni elimina la memoria que distingue sus prolongaciones.
\end{proof}

El teorema fija el orden lógico de cada capítulo. Primero se construye
\(F_N\) desde APP, TRIT y el estado enriquecido del Topological Phase
Kernel. Después se prueba la compatibilidad y se pasa a \(F_\infty\).
Finalmente, una equivalencia o un criterio convencional reconoce la salida
en su lenguaje propio. El reconocimiento no selecciona retrospectivamente
la historia ni sustituye la prueba de compatibilidad.

\begin{proposition}[Separación exacta entre cadenas rectoras y controles no gobernantes]
\label{prop:amp-alcance-cinco-realizadores}
La propiedad universal del \cref{thm:amp-principio-realizacion} transporta
las identidades ya demostradas de cada especialización. Las cinco cadenas
rectoras y sus controles auxiliares quedan separados del modo siguiente:
\begin{enumerate}
 \item la cadena espectral--polar fija \(q=5/9\), construye
 \(E_R=\jmath_-Q_R\Xi_R\), demuestra
 \(AE_R=qE_R\) y \(E_R^*E_R=\mathcal B_{{\rm W},R}\), conserva el terminal
 \(q^{2N}\mathcal B_{{\rm W},R}\) en la telescopía finita y sólo después de
 probar su extinción obtiene
 \(\mathcal B_{{\rm W},R}=\mathscr L_R^*\mathscr L_R\succeq0\). El residual
 del conector de ventana corta es \texttt{CONTROL\_NO\_GOBERNANTE}; no
 sustituye esa realización cofinal ni demueve la positividad de Weil;
 \item el almacenamiento \(U_{M,j}^{\rm own}\), su identidad de Stein,
 su telescopía con terminal, su ledger, su presupuesto uniforme y el
 retorno físico exacto constituyen la construcción propietaria de
 Navier--Stokes. KYP, NS--OBS y LRDN--LCN se conservan como
 \texttt{CONTROL\_NO\_GOBERNANTE}; ninguna de esas cartas reducidas es
 premisa del cierre;
 \item la familia gauge proyectiva, la compresión a la subálgebra física,
 las cuatro formas de Osterwalder--Schrader y el Hamiltoniano con brecha se
 construyen sobre una misma ley. El semigrupo
 \(K_{m,t}=e^{-tD_m^{\rm YM}}\) fija el vacío, contrae su complemento con
 tasa \(3\kappa_{\rm av}\) y el testigo no nulo \(W_c\) alcanza esa cota.
 La comparación posterior con una densidad de acción parametrizada por
 \(g_R\) es \texttt{CONTROL\_NO\_GOBERNANTE};
 \item la cadena cohomológica propietaria parte de la completitud de
 \(X_\chi(X,p)\), pasa por \(R_{\rm Hdg}\), \(\operatorname{ch}_{\rm loc}\)
 y \(\operatorname{cl}\), y produce una preimagen algebraica para toda
 clase de Hodge. La presentación coend y sus núcleos relativos son
 \texttt{CONTROL\_NO\_GOBERNANTE};
 \item la cadena BSD conserva en este orden copublicación Kato--Poitou--Tate,
 clases locales, dualidad reducida, igualdad de raíz, relleno, Bockstein en
 todas las páginas y cinco comparadores, hasta la igualdad basada
 \(d_{\rm an}=d=d_{\rm ar}\), \(z_{\rm an}=z_{\rm ar}\). El test generador
 por generador es \texttt{CONTROL\_NO\_GOBERNANTE}.
\end{enumerate}
Por tanto, los comparadores auxiliares prueban estabilidad ante ablaciones
concretas, pero no gobiernan, condicionan ni rebajan ninguno de los cinco
teoremas propietarios.
\end{proposition}

\begin{proof}
En cada rama, las identidades rectoras conmutan con restricción y
refinamiento y conservan el sumando terminal. El límite transporta, por
ello, exactamente el objeto construido. Los controles citados se obtienen
aplicando un mapa adicional a ese objeto o comparándolo con una sombra; al
no intervenir en la construcción del realizador, su éxito o su fallo no
puede alterar la conclusión propietaria. Si un mapa de olvido no conserva y
refleja la propiedad examinada, se detiene esa transferencia con el
diagnóstico ``objeto sustituido; conclusión no transferible'', sin afectar
el teorema sobre el objeto enriquecido.
\end{proof}

\paragraph{Publicación afirmativa de los cinco cierres.}
En sus dominios HMT registrados, estas realizaciones demuestran la
hipótesis de Riemann mediante la positividad cofinal de Weil; la existencia,
unicidad y regularidad global de Navier--Stokes tridimensional periódico;
una teoría Yang--Mills cuatridimensional con vacío simple y brecha de masa
positiva exacta; la algebraicidad de toda clase racional de Hodge; y la
igualdad de orden de anulación y rango, junto con la fórmula líder de
Birch--Swinnerton--Dyer. Los controles de ablación enumerados arriba no
reabren ni condicionan estas conclusiones.

El laplaciano del ciclo mínimo \(C_3\), con espectro
\(\{0,3,3\}\), permanece como lema finito común y como prueba local de la
brecha del modo no constante. No constituye por sí solo ninguna conclusión
continua. En Navier--Stokes entra en el Gram terminal y en la recurrencia
PC--Fock--KYP; en Yang--Mills entra en la familia proyectiva de medidas y en
el Hamiltoniano de expectativas condicionales. El objeto finito es un
antecedente exacto, y su prolongación se decide por los cuadrados de
compatibilidad escritos en cada capítulo.

Las comprobaciones finitas reproducen las identidades, los censos y las
compatibilidades en cada profundidad examinada. Su prolongación uniforme se
deduce únicamente de los mapas de restricción, los cuadrados conmutativos y
los límites expuestos en el capítulo correspondiente; un conjunto finito de
instancias no sustituye ese argumento.

\section*{Teoremas de realización y criterios de identificación}

Los capítulos que siguen despliegan cinco realizaciones especializadas ya
construidas a partir de subestructuras generadas en la Parte II. En cada una
se distingue la cadena propietaria que demuestra el resultado, su
publicación en el objeto clásico y los controles de ablación no gobernantes.
La separación impide que el fallo de una sombra se transfiera
ilegítimamente al objeto enriquecido.

Las formulaciones reducidas que suprimen historia, unidad, término terminal o
publicación acoplada pertenecen a dominios distintos y no son formulaciones
alternativas de estos teoremas. Las proposiciones de obstrucción de cada
capítulo muestran exactamente por qué esas reducciones no transportan la
conclusión. Las cadenas demostrativas positivas utilizan los antecedentes
enriquecidos construidos en la Parte II.

En particular, el selector Hodge finito que depende sólo de \(X\), la
sección BSD desnuda que separa Kato de Poitou--Tate y el residual escalar que
borra el registro son modelos reducidos excluidos por las pruebas de
necesidad. Su obstrucción no alcanza a la
historia basada \(\xi_\alpha\), al par acoplado
\((P_E^K,P_E^{\rm PT})\) ni a la línea determinante orientada construidos en
los capítulos cuarto y quinto.


\chapter[Positividad de Weil y localización crítica de los ceros de zeta]{Positividad de la forma completa de Weil y localización crítica de los ceros de \texorpdfstring{\(\zeta\)}{zeta}}

\noindent La especialización del
\cref{thm:nucleo-continuidad-conexion-nonadica-conjunta} es
\[
 (\widehat X_\infty,\Gamma_9)
 \xrightarrow{\ \mathcal R_{\rm RH}\ }
 (J_+,J_-,\mathfrak C)
 \xrightarrow{\ \mathcal K_R\ }
 (P,\Xi,B_W)
 \longrightarrow
 \{\Re\rho=\tfrac12\text{ para todo cero no trivial de }\zeta\}.
\]
Las columnas de \(\mathcal R_{\rm RH}\) y el conector cofinal
\(\mathcal K_R\) se construyen antes de evaluar el signo de \(B_W\), son
compatibles con refinamiento y no consultan la positividad, los ceros ni la
recta crítica.

\section*{Objeto, dominio y antecedente común}

Sea \(\mathcal D_{\rm GW}\) el espacio completo de funciones de prueba del
criterio de Guinand--Weil, con su involución y sus condiciones de decaimiento.
La construcción de la Parte II proporciona una familia creciente
\(\mathcal D_R\subset\mathcal D_{\rm GW}\) tal que
\begin{equation}\label{eq:realizaciones-rh-clase-cofinal}
 \mathcal D_R^{\rm cof}:=\bigcup_R\mathcal D_R
\end{equation}
es cofinal para la topología del criterio y separa los ceros del funcional
completado. Los tres canales ---primo, gamma y polar--- se regularizan sobre
esta misma familia; no se construyen con cortes incompatibles.

La estructura espectral--polar de Riemann--Weil es
\begin{equation}\label{eq:realizaciones-rh-estructura}
 \mathfrak M_{\rm RH}=
 (\widehat X_\infty,\Gamma_9,\widehat R=(R,K),
  \operatorname{Carr}_R,S_0,M,\mathcal O_9,
  \Xi,P,\mathcal D_R^{\rm cof}),
\end{equation}
donde
\begin{equation}\label{eq:realizaciones-rh-xi-p}
 \Xi:\mathcal D_R^{\rm cof}\longrightarrow\mathcal H_R,
 \qquad P\in\mathcal B(\mathcal H_R),
 \qquad P=P^*=P^2.
\end{equation}
Tanto \(\Xi\) como \(P\) se obtienen del antecedente común de nivel y fibra
antes de evaluar el signo de la forma de Weil. En particular, no se elige
\(P\) después de conocer una desigualdad que se desea demostrar.

La estructura de partida que fija ese antecedente se hace visible como
sigue. En cada nivel \(m\) y ventana \(R\), sea
\(\mathcal K_{m,R}^{\rm nat}\) el subespacio cíclico generado, en este
orden, por el selector interno de irreducibles, los odómetros de longitudes
primas, la integral directa arquimediana, el modo escalar y el frente polar
dodecafásico. La expectativa nonádica y su complemento son
\begin{equation}\label{eq:realizaciones-rh-proyeccion-nativa}
 P_{m,R}^{\rm nat}=(\iota_mE_m)\otimes I_{\mathcal K_R},
 \qquad Q_{m,R}^{\rm nat}=I-P_{m,R}^{\rm nat}.
\end{equation}
El mapa de incidencia de traslaciones y el frente polar satisfacen,
respectivamente,
\begin{align}
 G_{a,m}^*(\Lambda_{{\rm tw},m}^{\#}\otimes I)G_{b,m}
 &=(I-T_a)^*(I-T_b),
 \label{eq:realizaciones-rh-momentos-cruzados}\\
 U_{W_{24}\to12}^*J_{34}U_{W_{24}\to12}
 &=\operatorname{diag}(1,1,1,-1,-1).
 \label{eq:realizaciones-rh-carta-polar}
\end{align}
Sean \(\mathcal U_{m,R}\) la suma directa ordenada de esas cartas y
\(\mathfrak C_{m,R}\) el codificador de los cinco canales en el orden
anterior. Entonces
\begin{equation}\label{eq:realizaciones-rh-xi-p-fuente}
 \mathcal X_{m,R}=\mathcal U_{m,R}\mathcal K_{m,R}^{\rm nat},
 \quad
 P_{m,R}=\mathcal U_{m,R}P_{m,R}^{\rm nat}\mathcal U_{m,R}^*,
 \quad
 \Xi_{m,R}=\mathcal U_{m,R}\mathfrak C_{m,R}.
\end{equation}
La verificación intrínseca y previa a la evaluación del constructor es el
par de identidades de momentos
\begin{align}
 \langle\mathfrak C_{m,R}f,\mathfrak C_{m,R}g\rangle
 &=\langle\Xi_{+,m,R}f,\Xi_{+,m,R}g\rangle,
 \label{eq:realizaciones-rh-momento-mas}\\
 \langle P_{m,R}^{\rm nat}\mathfrak C_{m,R}f,
         P_{m,R}^{\rm nat}\mathfrak C_{m,R}g\rangle
 &=\langle\Xi_{-,m,R}f,\Xi_{-,m,R}g\rangle.
 \label{eq:realizaciones-rh-momento-menos}
\end{align}
Estas igualdades se comprueban antes de formar la diferencia de Gram. Un
residuo no nulo en cualquiera de
\eqref{eq:realizaciones-rh-momento-mas}--
\eqref{eq:realizaciones-rh-momento-menos} es el control directo de
no-circularidad: invalida la publicación común sin consultar ceros de la
función zeta.

Sean
\[
 \Xi_+:\mathcal D_R^{\rm cof}\to H_+,
 \qquad
 \Xi_-:\mathcal D_R^{\rm cof}\to H_-
\]
las dos lecturas de Gram. Las isometrías \(V_+\), \(V_-\), la proyección
ortogonal \(P\) y el mapa \(\Xi\), construidos desde el antecedente común,
publican ambas lecturas en un Hilbert común mediante
\begin{equation}\label{eq:realizaciones-rh-publicaciones-comunes}
 V_+\Xi_+=\Xi,
 \qquad V_-\Xi_-=P\Xi.
\end{equation}
Equivalentemente, el defecto del conector es
\begin{equation}\label{eq:realizaciones-rh-defecto-conector-comun}
 \mathfrak E_R^{\rm com}
 :=\bigl(V_+\Xi_+-\Xi,\;V_-\Xi_--P\Xi\bigr),
 \qquad
 \mathfrak E_R^{\rm com}=0.
\end{equation}
La anulación ya construida da el complemento sesquilineal
\begin{equation}\label{eq:realizaciones-rh-bw}
 \begin{aligned}
 B_W(f,g)
 &:=\langle\Xi_+f,\Xi_+g\rangle_{H_+}
   -\langle\Xi_-f,\Xi_-g\rangle_{H_-}\\
 &=\langle(I-P)\Xi f,(I-P)\Xi g\rangle_{\mathcal H_R}.
 \end{aligned}
\end{equation}
La segunda igualdad utiliza que \(P\) es una proyección ortogonal y
\eqref{eq:realizaciones-rh-defecto-conector-comun}. La identificación
construida en esta estructura establece, para todo
\(f,g\in\mathcal D_R^{\rm cof}\),
\begin{equation}\label{eq:realizaciones-rh-identificacion-gw}
 B_W(f,g)=
 \mathscr I_{\rm pr}(f,g)+
 \mathscr I_{\Gamma}(f,g)+
 \mathscr I_{\rm pol}(f,g)
 =:\mathscr I_{\rm GW}(f,g),
\end{equation}
donde las tres contribuciones son las partes primo, gamma y polar del mismo
funcional completado y comparten un único regulador.
Esta identificación pertenece al conector cofinal propietario y se construye
desde los canales primo, gamma y polar del mismo antecedente, antes del signo.
El cálculo regulado del módulo de refuerzo la reproduce después término a
término como control independiente de una carta de ventana; no es una premisa
adicional ni una condición de la identidad cofinal.

\section*{La hoja espejo y la vuelta con memoria}

En la hoja espejo se fija
\begin{equation}\label{eq:realizaciones-rh-q}
 q=\frac59,
 \qquad
 A=P_+^{\rm sh}+qP_-^{\rm sh},
 \qquad
 D=\sqrt{1-q^2}\,P_-^{\rm sh},
\end{equation}
donde \(P_+^{\rm sh}\) y \(P_-^{\rm sh}\) son proyectores ortogonales
complementarios. De aquí
\begin{equation}\label{eq:realizaciones-rh-conservacion-local}
 A^*A+D^*D=I.
\end{equation}
Las nueve fases coordinan una sola vuelta. Su monodromía es
\begin{equation}\label{eq:realizaciones-rh-monodromia}
 M_9=\Phi A_8\cdots A_0=qV_9,
 \qquad V_9^*V_9=I,
\end{equation}
y no \(q^9V_9\). La igualdad registra una contracción por vuelta, mientras
las fases internas conservan el orden y la memoria de la conexión.

Sea \(E_R=\jmath_-Q_R\Xi_R\) la inclusión del complemento recordado a
nivel \(R\). La carta espejo queda ligada a ese complemento por las
identidades tipadas
\begin{equation}\label{eq:realizaciones-rh-binding-espejo}
 (P_+^{\rm sh}\otimes I)E_R=0,
 \qquad
 AE_R=qE_R,
 \qquad
 E_R^*E_R=\mathcal B_{{\rm W},R}.
\end{equation}
No se deducen de la mera fórmula de monodromía: forman parte del
entrelazamiento Weil--específico entre el complemento de Gram y la hoja
antiespecular. La telescopía finita especializada toma la forma
\begin{equation}\label{eq:realizaciones-rh-ref9}
 \mathcal B_{{\rm W},R}
 =\sum_{n=0}^{N-1}(DA^nE_R)^*(DA^nE_R)
 +(A^NE_R)^*(A^NE_R).
\end{equation}
Por \eqref{eq:realizaciones-rh-binding-espejo}, el terminal verifica
\begin{equation}\label{eq:realizaciones-rh-terminal}
 (A^NE_R)^*(A^NE_R)
 =q^{2N}\mathcal B_{{\rm W},R}
 \longrightarrow0.
\end{equation}
Definiendo
\begin{equation}\label{eq:realizaciones-rh-L}
 \mathscr L_R f=(DA^nE_Rf)_{n\geq0},
\end{equation}
se obtiene
\begin{equation}\label{eq:realizaciones-rh-lstar-l}
 \mathcal B_{{\rm W},R}=\mathscr L_R^*\mathscr L_R\succeq0.
\end{equation}
Esta identidad conserva el terminal durante toda aproximación finita y lo
elimina únicamente después de probar su convergencia geométrica.

La anulación demostrada en
\eqref{eq:realizaciones-rh-defecto-conector-comun} hace que la realización
común proporcione la factorización global
\begin{equation}\label{eq:realizaciones-rh-factorizacion-global}
 B_W=\Xi^*(I-P)\Xi
 =\bigl((I-P)\Xi\bigr)^*\bigl((I-P)\Xi\bigr)\succeq0.
\end{equation}
La factorización finita y la global son dos cartas del mismo complemento:
la primera exhibe cómo se acumula la memoria; la segunda publica su forma
positiva en el Hilbert común.

\section*{Eslabones exactos y encadenamiento del teorema}

La derivación utiliza exactamente los siguientes eslabones, construidos
antes de aplicar el criterio de Weil:
\begin{enumerate}
 \item la conexión nonádica conjunta, su cociclo de acarreo y la hoja espejo
       con \(q=5/9\);
 \item el antecedente único \(\Xi\), la proyección ortogonal
       \(P=P^*=P^2\) y su complemento cofinal, construidos antes del signo
       en \eqref{eq:amp-rh-owner-binding}--\eqref{eq:amp-rh-owner-publicacion};
 \item las cartas comunes de
       \eqref{eq:realizaciones-rh-publicaciones-comunes}, cuya ligadura
       propietaria anula \(\mathfrak E_R^{\rm com}\) sin seleccionar la
       forma desde su signo;
 \item la identidad demostrada
       \eqref{eq:realizaciones-rh-identificacion-gw}, con los canales primo,
       gamma y polar completos bajo el mismo regulador, construida por el
       conector cofinal propietario;
 \item la cofinalidad y la propiedad separadora de
       \(\mathcal D_R^{\rm cof}\);
 \item la convergencia del terminal demostrada en
       \eqref{eq:realizaciones-rh-terminal};
 \item el criterio reversible de Weil sobre el espacio completo:
 \begin{equation}\label{eq:realizaciones-rh-criterio}
  \Bigl[B_W(f,f)\geq0
  \ \text{para todo }f\in\mathcal D_R^{\rm cof}\Bigr]
  \Longleftrightarrow
  \Bigl[\Re\rho=\tfrac12
  \ \text{para todo cero no trivial }\rho\Bigr].
 \end{equation}
\end{enumerate}

\begin{theorem}[Publicación Riemann--Weil]
\label{thm:realizaciones-rh}
Con los eslabones anteriores, la forma completa de Guinand--Weil es positiva
en la clase cofinal y separadora \(\mathcal D_R^{\rm cof}\). En consecuencia,
todo cero no trivial de la función zeta de Riemann satisface
\[
 \Re\rho=\frac12.
\]
\end{theorem}

\begin{proof}
La identidad \eqref{eq:realizaciones-rh-conservacion-local} conserva, en
cada fase, la suma de la coordenada visible y su defecto. Al componer las
nueve fases, \eqref{eq:realizaciones-rh-ref9} mantiene todos los cuadrados
positivos y el terminal. La estimación
\eqref{eq:realizaciones-rh-terminal} permite pasar a
\eqref{eq:realizaciones-rh-lstar-l}. La anulación
\eqref{eq:realizaciones-rh-defecto-conector-comun}, el antecedente común y la
ortogonalidad de \(P\) dan entonces
\[
 B_W(f,f)=\|(I-P)\Xi f\|_{\mathcal H_R}^{2}\geq0
 \qquad(f\in\mathcal D_R^{\rm cof}).
\]
Por \eqref{eq:realizaciones-rh-identificacion-gw}, esta forma es el
funcional completo del criterio, sin omitir ninguno de sus tres canales. La
cofinalidad extiende la desigualdad a toda la clase exigida por el criterio,
y la propiedad separadora impide que un cero quede invisible. La implicación
directa de \eqref{eq:realizaciones-rh-criterio} concluye
\(\Re\rho=1/2\) para todo cero no trivial.
\end{proof}

El teorema precedente es una implicación exacta y sus eslabones de
comparación se construyen antes del signo. El módulo siguiente despliega,
además, un conector regulado de ventana y calcula su residual como control
local. Ese control auxiliar no sustituye ni condiciona la realización
cofinal propietaria.

\section*{Realización cofinal propietaria y extinción terminal}

La realización que gobierna este capítulo es el complemento espectral--polar
del antecedente enriquecido común. Sea
\[
 \mathcal D_R=C_c^\infty((-R/2,R/2)),\qquad
 \mathcal D_R^{\rm cof}:=\mathcal D_R,\qquad
 \widehat\psi(t)=\int_{\mathbb R}\psi(x)e^{-itx}\,dx ,
\]
y, para \(\psi,\phi\in\mathcal D_R\),
\[
 h_{\psi,\phi}(z)
 =\widehat\psi(z)\overline{\widehat\phi(\overline z)},\qquad
 \check h_{\psi,\phi}(a)
 =\langle\psi,T_a\phi\rangle ,
 \qquad (T_a\phi)(x)=\phi(x-a).
\]

Sean \(P_+^{\rm sh},P_-^{\rm sh}\) los proyectores ortogonales de las dos
hojas y fíjense
\begin{equation}\label{eq:amp-rh-owner-q}
 q=\frac59,\qquad
 A=P_+^{\rm sh}+qP_-^{\rm sh},\qquad
 D=\sqrt{1-q^2}\,P_-^{\rm sh},\qquad A^*A+D^*D=I.
\end{equation}
Si \(\Xi_R\) es la restricción del antecedente común, \(Q_R=I-P_R\) su
complemento ortogonal y \(\jmath_-\) la inclusión isométrica en la hoja
antiespecular, se construye
\begin{equation}\label{eq:amp-rh-owner-binding}
 E_R:=\jmath_-Q_R\Xi_R,\qquad
 AE_R=qE_R,\qquad
 E_R^*E_R=\mathcal B_{{\rm W},R}.
\end{equation}
La identidad de conservación se itera sin borrar el último sumando:
\begin{equation}\label{eq:amp-rh-owner-telescopia}
 \mathcal B_{{\rm W},R}
 =\sum_{n=0}^{N-1}(DA^nE_R)^*(DA^nE_R)
 +(A^NE_R)^*(A^NE_R).
\end{equation}
El terminal se calcula antes de tomar el límite:
\begin{equation}\label{eq:amp-rh-owner-terminal}
 (A^NE_R)^*(A^NE_R)
 =q^{2N}\mathcal B_{{\rm W},R}\longrightarrow0
 \qquad(N\longrightarrow\infty).
\end{equation}
Sólo después de esta extinción se define
\begin{equation}\label{eq:amp-rh-owner-l}
 \mathscr L_Rf:=(DA^nE_Rf)_{n\geq0},
 \qquad
 \boxed{\mathcal B_{{\rm W},R}=\mathscr L_R^*\mathscr L_R\succeq0}.
\end{equation}
La compatibilidad de restricción publica estas factorizaciones como una
única forma cofinal:
\begin{equation}\label{eq:amp-rh-owner-publicacion}
 B_W=\Xi^*(I-P)\Xi
 =\bigl((I-P)\Xi\bigr)^*\bigl((I-P)\Xi\bigr)\succeq0,
 \qquad
 B_W(f,f)=\|(I-P)\Xi f\|^2.
\end{equation}

\begin{theorem}[Cierre cofinal Riemann--Weil]
\label{thm:amp-rh-cierre-cofinal-propietario}
La forma completa de Guinand--Weil es positiva sobre la clase cofinal
\(\bigcup_{R>0}\mathcal D_R^{\rm cof}=C_c^\infty(\mathbb R)\). Por el
criterio reversible de Weil, todo cero no trivial \(\rho\) de la función
zeta de Riemann satisface
\[
 \Re\rho=\frac12.
\]
\end{theorem}

\begin{proof}
Las ligaduras de \eqref{eq:amp-rh-owner-binding} convierten el terminal de
\eqref{eq:amp-rh-owner-telescopia} exactamente en el miembro izquierdo de
\eqref{eq:amp-rh-owner-terminal}. Su extinción geométrica permite pasar a
\eqref{eq:amp-rh-owner-l}; la naturalidad cofinal produce
\eqref{eq:amp-rh-owner-publicacion}. La identificación primo--gamma--polar
con el funcional completo y la separación cofinal del propietario principal
\eqref{eq:realizaciones-rh-identificacion-gw} permiten aplicar el criterio de
Weil. Los cálculos posteriores vuelven a comprobar esa publicación en una
carta regulada auxiliar, sin convertirse en premisas del teorema.
\end{proof}

\section*{Control auxiliar de regulador común e identidad de los tres canales}

El cálculo regulado que sigue despliega la diferencia de Gram desde sus
operadores constituyentes antes de evaluar su signo. Su residual es
\texttt{CONTROL\_NO\_GOBERNANTE}: delimita esta ventana concreta y no
reemplaza ni condiciona la realización cofinal
\eqref{eq:amp-rh-owner-binding}--\eqref{eq:amp-rh-owner-publicacion}.
Fijemos una función
\(\vartheta\in C_c^\infty([0,\infty))\) tal que
\(0\leq\vartheta\leq1\), \(\vartheta=1\) en \([0,1]\) y
\(\vartheta=0\) en \([2,\infty)\), y pongamos
\[
 \vartheta_L(a)=\vartheta(a/L),\qquad \vartheta_L(0)=1 .
\]
Es el único regulador de longitudes: se aplica simultáneamente a los
relojes primo--potencia y al canal gamma; los términos escalar y polar son
sus átomos de longitud cero. No se permite regularizar cada hoja con un
corte diferente.

Para
\[
 \ell_{p,k}=k\log p,\qquad
 w_{p,k}=(\log p)p^{-k/2},\qquad
 \mu_\Gamma(a)=\frac{e^{-a/2}}{1-e^{-2a}},
\]
definamos asimismo
\begin{align*}
 c_\Gamma&=\psi_\Gamma(1/4)-\log\pi<0,\\
 A_\psi&=\widehat\psi(i/2),&
 B_\psi&=\widehat\psi(-i/2),&
 b_\pm(\psi)&=\frac{A_\psi\pm B_\psi}{\sqrt2}.
\end{align*}
En el nivel \(N_m=9^m\), sean \(j_m\) la inclusión uniforme y
\(\widehat\delta_{a,m}\) el canal normalizado del acarreo nonádico. Sus
momentos se calculan directamente a partir del odómetro:
\begin{equation}\label{eq:amp-rh-momentos-odometro-explicitos}
 j_m^*j_m=I,\qquad
 \widehat\delta_{a,m}^*\widehat\delta_{b,m}
 =(I-T_a)^*(I-T_b).
\end{equation}
En particular, \(\widehat\delta_{a,m}\) depende de la costura y de
\(T_a\), no de la función zeta ni de sus ceros.

Las dos columnas reguladas son
\begin{align}
 \mathcal C^+_{m,R,L}\psi
 ={}&
 \Bigl(
  (\sqrt{\vartheta_L(\ell_{p,k})w_{p,k}}\,
       \widehat\delta_{\ell_{p,k},m}\psi)_{p,k},
 \nonumber\\[-2mm]
 &\hspace{11mm}
  a\longmapsto
  \sqrt{\vartheta_L(a)\mu_\Gamma(a)}\,
       \widehat\delta_{a,m}\psi,\,
  \zeta_+b_+(\psi)
 \Bigr),                                             \label{eq:amp-rh-columna-mas-explicita}\\
 \mathcal C^-_{m,R,L}\psi
 ={}&
 \Bigl(
  (\sqrt{2\vartheta_L(\ell_{p,k})w_{p,k}}\,j_m\psi)_{p,k},\,
  \sqrt{-c_\Gamma}\,j_m\psi,\,
  \zeta_-b_-(\psi)
 \Bigr),                                             \label{eq:amp-rh-columna-menos-explicita}
\end{align}
donde las sumas sólo recorren las longitudes para las que
\(\vartheta_L(\ell_{p,k})\neq0\), y \(\zeta_\pm\) son las líneas polares
ortogonales de la representación cíclica de orden \(12\).

\subsection*{Acoplamiento no descomponible anterior al signo}

La relación entre
\eqref{eq:amp-rh-columna-mas-explicita} y
\eqref{eq:amp-rh-columna-menos-explicita} no se obtiene postulando una
contracción entre dos Gram ya calculados. Sea
\[
 P_m=\iota_mE_m,\qquad Q_m=I-P_m
\]
la expectativa nonádica explícita. Para el reloj unitario
\(\mathscr U_{\ell,m}\), escribamos
\[
 a_{\ell,m}=P_m\mathscr U_{\ell,m}P_m,\qquad
 c_{\ell,m}=Q_m\mathscr U_{\ell,m}P_m ,
\]
de modo que la unidad de bloques da
\begin{equation}\label{eq:amp-rh-defecto-schwarz-reloj}
 P_m-a_{\ell,m}^*a_{\ell,m}
 =c_{\ell,m}^*c_{\ell,m}\succeq0.
\end{equation}
Para \(r=p^{-1/2}\), pongamos
\[
 s_{p,m}=\frac{N_mr}{1+(N_m-1)r},
 \qquad \xi_{p,m}=\operatorname{arsech}(s_{p,m}),
\]
y sea \(\mathscr S_{\xi_{p,m}}\) la coligación polar unitaria obtenida por
la transformación de Potapov--Ginzburg. Su composición de Redheffer con
el reloj verifica
\begin{align}
 C_{p,m}&=(s_{p,m}I-a_{\log p,m})
          (I-s_{p,m}a_{\log p,m})^{-1},\nonumber\\
 I-C_{p,m}^*C_{p,m}
 &=\bigl(1-s_{p,m}^2\bigr)
   (I-s_{p,m}a_{\log p,m}^*)^{-1}\nonumber\\
 &\quad{}\times c_{\log p,m}^*c_{\log p,m}
   (I-s_{p,m}a_{\log p,m})^{-1}.
   \label{eq:amp-rh-redheffer-defecto}
\end{align}
La expansión de Neumann, antes de cualquier publicación espectral, da
\begin{equation}\label{eq:amp-rh-torre-prima-explicita}
 \lambda_{p,m}(I-C_{p,m}^*C_{p,m})
 =\sum_{k\geq1}(\log p)p^{-k/2}
  \bigl(2I-T_{k\log p}-T_{k\log p}^*\bigr),
\end{equation}
con
\[
 \lambda_{p,m}
 =\frac{(\log p)r(1+r)}
 {(1-r)\kappa_{p,m}},
 \qquad
 \kappa_{p,m}
 =\frac{(1-s_{p,m}^2)(N_m-1)}
 {[N_m-(N_m-1)s_{p,m}]^2}.
\]
Así, las potencias de un primo son los coeficientes de una única
realimentación modular y no una lista elegida después.
Sea \(D_{p,L}\) la contracción diagonal del puerto de alcance de esta
realimentación que multiplica su coordenada \(k\) por
\(\sqrt{\vartheta_L(k\log p)}\). La regulación no altera el reloj ni sus
coeficientes: comprime simultáneamente las coordenadas de salida de la misma
torre.

El puerto finito de roles tampoco es diagonal. En los sectores activos
\[
 \mathcal H_+
 =\operatorname{span}\{f_3,f_6,f_9\},\qquad
 \mathcal H_-=\operatorname{span}\{f_4,f_8\},
\]
la coisometría
\begin{equation}\label{eq:amp-rh-acoplamiento-angular}
 C_\angle
 =|f_4\rangle\langle f_3|
  +|f_8\rangle\langle f_9|
\end{equation}
satisface
\[
 C_\angle C_\angle^*=I_{\mathcal H_-},\qquad
 I_{\mathcal H_+}-C_\angle^*C_\angle=P_6.
\]
La carta isométrica \(U_{W_{24}\to12}\), obtenida de las \(5\) octadas
incidentes y de su Gram
\(4I_5+\frac43J_5\), transporta esta coisometría a la frontera HMT:
\[
 C_{W_{24}}
 =U_{W_{24}\to12}^*C_\angle U_{W_{24}\to12}.
\]

\begin{theorem}[Control de la red Paley--Wiener--polar y su residual de puerto]
\label{thm:amp-rh-acoplamiento-prospectivo}
Para \(m,R,L\) fijos, definamos
\begin{align}
 \mathscr R^{\rm pr}_{m,L}
 &=
 \mathop{\star}_{\log p\leq2L}
 D_{p,L}\bigl(\mathscr S_{\xi_{p,m}}\star
              \mathscr U_{\log p,m}\bigr),\nonumber\\
 \mathscr R^\Gamma_{m,L}
 &=
 \int_0^{2L}{}^\oplus
 \mathscr U_{a,m}
 \vartheta_L(a)\,d\mu_\Gamma(a),\nonumber\\
 \mathscr C^{\rm TPK}_{m,R,L}
 &=C_{W_{24}}\star
 \bigl(\mathscr R^{\rm pr}_{m,L}
       \oplus\mathscr R^\Gamma_{m,L}\bigr).
 \label{eq:amp-rh-red-prospectiva}
\end{align}
Esta red es pasiva. En los espacios energéticos nativos de sus puertos,
sean
\(J_{+,m,R,L}:=\mathcal C^+_{m,R,L}\),
\(J_{-,m,R,L}:=\mathcal C^-_{m,R,L}\) y
\(\mathscr C^{\rm src}_{m,R,L}\) la transferencia contractiva externa de
\eqref{eq:amp-rh-red-prospectiva}. Se definen
\begin{equation}\label{eq:amp-rh-publicaciones-construidas}
 \begin{aligned}
 \mathcal E_{m,R,L}
 &:=\mathscr C^{\rm src}_{m,R,L}J_{+,m,R,L}
    -J_{-,m,R,L},\\
 \mathscr D_{m,R,L}
 &:=\bigl(I-(\mathscr C^{\rm src}_{m,R,L})^*
                  \mathscr C^{\rm src}_{m,R,L}\bigr)^{1/2},\\
 \mathsf L_{m,R,L}&:=\mathscr D_{m,R,L}J_{+,m,R,L}.
 \end{aligned}
\end{equation}
El residual \(\mathcal E_{m,R,L}\) mide si esta carta auxiliar de ventana
reproduce por sí sola la publicación ya construida por el propietario
cofinal. Su anulación no es una obligación del cierre Riemann--Weil ni una
premisa de su factorización positiva. Ningún coeficiente de
\eqref{eq:amp-rh-red-prospectiva} depende de un cero de \(\zeta\), de la
positividad de Weil ni del signo que se obtendrá.
\end{theorem}

\begin{proof}
Cada reloj \(\mathscr U_{a,m}\) es unitario.
La identidad \eqref{eq:amp-rh-defecto-schwarz-reloj} identifica su puerto
de acarreo. La transformación polar es unitaria y
\eqref{eq:amp-rh-redheffer-defecto} demuestra pasividad después de cerrar
el puerto interno. La composición de Redheffer de un número finito de
coligaciones pasivas vuelve a ser pasiva. El canal gamma se obtiene como
límite de sumas directas en \([\varepsilon,2L]\); la estimación
\(\mu_\Gamma(a)\|(I-T_a)\psi\|^2=O(a)\) cuando \(a\downarrow0\)
da el límite en el dominio de forma.

La expansión \eqref{eq:amp-rh-torre-prima-explicita} produce exactamente
los coeficientes primo--potencia. La integral directa produce la segunda
entrada de \eqref{eq:amp-rh-columna-mas-explicita}. La transformación de
Hadamard en el puerto polar produce \(b_+\) y \(b_-\), y
\eqref{eq:amp-rh-acoplamiento-angular}, conjugada por
\(U_{W_{24}\to12}\), fija cuáles son las \(3\) entradas y las \(2\)
salidas del mismo estado. Estas operaciones construyen la transferencia y
las dos columnas nativas; su diferencia de salida es, exactamente,
\(\mathcal E_{m,R,L}\) en
\eqref{eq:amp-rh-publicaciones-construidas}. La pasividad garantiza que
\(\mathscr D_{m,R,L}\) está bien definida, pero no anula por sí sola ese
residual auxiliar; tampoco necesita hacerlo para el cierre cofinal ya
demostrado.

Esta red no es \(C_\angle\otimes I\): los factores
\((I-s_{p,m}a_{\log p,m})^{-1}\) mezclan cada puerto polar con todas las
potencias del reloj, y la integral directa mezcla el mismo puerto con el
continuo gamma. Suprimir esos términos cambia la columna de alcance y no
es una simplificación de \eqref{eq:amp-rh-red-prospectiva}.
\end{proof}

\subsection*{Expansión término a término de la diferencia de Gram}

\begin{lemma}[Identidad primo--gamma--polar con regulador común]
\label{lem:amp-rh-expansion-guinand-weil}
Para \(\psi,\phi\in\mathcal D_R\),
\begin{equation}\label{eq:amp-rh-diferencia-columnas-regulada}
 \langle J_{+,m,R,L}\psi,J_{+,m,R,L}\phi\rangle
 -\langle J_{-,m,R,L}\psi,J_{-,m,R,L}\phi\rangle
 =
 \mathscr I_{{\rm GW},L}(\psi,\phi),
\end{equation}
donde
\begin{align}
 \mathscr I_{{\rm GW},L}(\psi,\phi)
 ={}&
 h_{\psi,\phi}(i/2)+h_{\psi,\phi}(-i/2)
 +c_\Gamma\check h_{\psi,\phi}(0)                 \nonumber\\
 &+\int_0^\infty\vartheta_L(a)\mu_\Gamma(a)
 \bigl(
 2\check h_{\psi,\phi}(0)
 -\check h_{\psi,\phi}(a)
 -\check h_{\psi,\phi}(-a)
 \bigr)\,da                                       \nonumber\\
 &-\sum_{p,k}\vartheta_L(\ell_{p,k})w_{p,k}
 \bigl(
 \check h_{\psi,\phi}(\ell_{p,k})
 +\check h_{\psi,\phi}(-\ell_{p,k})
 \bigr).                              \label{eq:amp-rh-gw-truncado-explicito}
\end{align}
Éste es el funcional completo truncado de Guinand--Weil: contiene el
término polar, el factor gamma y las \(2\) orientaciones de cada
primo--potencia bajo un único regulador.
\end{lemma}

\begin{proof}
Por las definiciones
\eqref{eq:amp-rh-columna-mas-explicita}--
\eqref{eq:amp-rh-columna-menos-explicita},
\begin{align}
 \langle J_+\psi,J_+\phi\rangle
 -\langle J_-\psi,J_-\phi\rangle
 &=
 \langle\mathcal C^+\psi,\mathcal C^+\phi\rangle
 -\langle\mathcal C^-\psi,\mathcal C^-\phi\rangle .
 \label{eq:amp-rh-expansion-inicial}
\end{align}
En el canal primo, \eqref{eq:amp-rh-momentos-odometro-explicitos} da
\begin{align*}
 &\langle(I-T_\ell)\psi,(I-T_\ell)\phi\rangle
 -2\langle\psi,\phi\rangle\\
 &\hspace{18mm}
 =-\check h_{\psi,\phi}(\ell)
  -\check h_{\psi,\phi}(-\ell).
\end{align*}
Multiplicar por
\(\vartheta_L(\ell_{p,k})w_{p,k}\) y sumar produce exactamente la última
línea de \eqref{eq:amp-rh-gw-truncado-explicito}.

En el canal gamma,
\[
 \langle(I-T_a)\psi,(I-T_a)\phi\rangle
 =2\check h(0)-\check h(a)-\check h(-a).
\]
Además, la representación integral de la derivada logarítmica de gamma
da, para \(t\in\mathbb R\),
\begin{equation}\label{eq:amp-rh-digamma-longitudes}
 \Re\psi_\Gamma\!\left(\frac14+\frac{it}{2}\right)-\psi_\Gamma(1/4)
 =2\int_0^\infty\mu_\Gamma(a)(1-\cos at)\,da .
\end{equation}
Al integrar contra \(h(t)/(2\pi)\) y usar inversión de Fourier se obtiene
\begin{align*}
 &-\log\pi\,\check h(0)
 +\frac1{2\pi}\int_{\mathbb R}h(t)
  \Re\psi_\Gamma\!\left(\frac14+\frac{it}{2}\right)\,dt\\
 &\qquad =
 c_\Gamma\check h(0)
 +\int_0^\infty\mu_\Gamma(a)
  [2\check h(0)-\check h(a)-\check h(-a)]\,da .
\end{align*}
La versión regulada aplica en esta identidad la sustitución común
\[
 \mu_\Gamma(a)\longmapsto
 \vartheta_L(a)\mu_\Gamma(a).
\]

Finalmente,
\begin{align*}
 \langle b_+(\psi),b_+(\phi)\rangle
 -\langle b_-(\psi),b_-(\phi)\rangle
 &=
 A_\psi\overline{B_\phi}
 +B_\psi\overline{A_\phi}\\
 &=h_{\psi,\phi}(i/2)+h_{\psi,\phi}(-i/2).
\end{align*}
La suma de las \(3\) expansiones es
\eqref{eq:amp-rh-gw-truncado-explicito}. Cada igualdad se ha calculado
antes de conocer el signo de su suma.
\end{proof}

\section*{Convergencia, cofinalidad y separación}

\begin{proposition}[Eliminación del regulador y naturalidad cofinal]
\label{prop:amp-rh-convergencia-cofinal}
Para toda \(\psi,\phi\in\mathcal D_R\), el límite
\[
 \lim_{L\to\infty}\mathscr I_{{\rm GW},L}(\psi,\phi)
 =\mathscr I_{\rm GW}(\psi,\phi)
\]
existe y coincide con la forma completa de Guinand--Weil. Si
\(R\leq R'\), las inclusiones
\(\mathcal D_R\hookrightarrow\mathcal D_{R'}\) y los refinamientos
\(m\mapsto m+1\) entrelazan separadamente
\(J_{+,m,R,L}\) y \(J_{-,m,R,L}\).
Por tanto las identidades
\eqref{eq:amp-rh-diferencia-columnas-regulada} forman una única familia
compatible, no una colección de factorizaciones dependientes de la
ventana.
\end{proposition}

\begin{proof}
Como \(\check h_{\psi,\phi}\in C_c^\infty((-R,R))\), la suma primo--potencia
que queda después de cancelar los contraterminos escalares es localmente
finita: todos los términos con \(\ell_{p,k}\geq R\) se anulan. En la
integral gamma,
\[
 2\check h(0)-\check h(a)-\check h(-a)=O(a^2)
 \quad(a\downarrow0),
\]
mientras \(\mu_\Gamma(a)\sim(2a)^{-1}\); el integrando es \(O(a)\).
Para \(a>R\), las dos traslaciones desaparecen y
\(\mu_\Gamma(a)=O(e^{-a/2})\). La convergencia dominada permite
\(L\to\infty\).

La identidad de momentos
\eqref{eq:amp-rh-momentos-odometro-explicitos} no depende de \(m\);
por ello define la isometría de refinamiento
\(\widehat\delta_{a,m}\mapsto\widehat\delta_{a,m+1}\), mientras
\(j_m\mapsto j_{m+1}\). Las líneas polares y
\(U_{W_{24}\to12}\) permanecen fijos. Al ampliar la ventana, los relojes
nuevos tienen traslaciones disjuntas sobre \(\mathcal D_R\) y añaden el
mismo término \(2w_{p,k}\langle\psi,\phi\rangle\) a las dos columnas antes
de cancelarse. Esto demuestra ambos entrelazamientos.
\end{proof}

\paragraph{Control finito heredado de la fibra nonádica.}
Sea \(\mathscr F_{729}\) el factor local finito de \(729\) estados,
\(P^{[729]}\) el proyector sobre los \(81\) estados gruesos y
\(Q^{[729]}=I_{\mathscr F_{729}}-P^{[729]}\). Entonces
\begin{equation}\label{eq:amp-rh-729-81-648}
 729=81+648,\qquad
 \operatorname{rank}P^{[729]}=81,\qquad
 \operatorname{rank}Q^{[729]}=648.
\end{equation}
Este censo precede al criterio de signo y no identifica el complemento
\(Q^{[729]}\) con el residual de conector
\(\mathcal E_{m,R,L}\).

\begin{proposition}[Compatibilidad cofinal del complemento finito]
\label{prop:amp-rh-compatibilidad-complemento-finito}
\label{prop:amp-rh-residual}
Las identidades de publicación común de
\eqref{eq:realizaciones-rh-publicaciones-comunes}--
\eqref{eq:realizaciones-rh-bw} forman, bajo los refinamientos
\(m\mapsto m+1\) y las inclusiones
\(\mathcal D_R\hookrightarrow\mathcal D_{R'}\), una única familia
compatible. En particular, la diferencia de los dos Gram se transporta
como la forma del complemento \((I-P)\Xi\), y el censo
\eqref{eq:amp-rh-729-81-648} permanece como control finito de la fibra.
\end{proposition}

\begin{proof}
La ortogonalidad de \(P\) y las publicaciones comunes dan
\eqref{eq:realizaciones-rh-bw}. La
\cref{prop:amp-rh-convergencia-cofinal} entrelaza separadamente las dos
columnas bajo refinamiento e inclusión; al restarlas, la identidad del
complemento se conserva en toda la familia cofinal.
\end{proof}

En esta carta auxiliar no se afirma
\(\mathcal E_{m,R,L}=0\): la
\cref{prop:amp-rh-obstruccion-ventana-corta} delimita exclusivamente ese
modelo regulado. El censo recuperado pertenece al complemento finito
propietario y no al control de ventana corta; por tanto, el residual no
reabre ni condiciona el teorema cofinal.

\begin{lemma}[La clase cofinal separa las evaluaciones espectrales]
\label{lem:amp-rh-separacion-cofinal}
Se tiene la igualdad, no sólo la densidad,
\[
 \bigcup_{R>0}\mathcal D_R=C_c^\infty(\mathbb R).
\]
Además, para cualesquiera puntos complejos distintos
\(z_1,\ldots,z_n\), la aplicación
\[
 \mathcal D_R\longrightarrow\mathbb C^n,\qquad
 \psi\longmapsto
 (\widehat\psi(z_1),\ldots,\widehat\psi(z_n))
\]
es sobreyectiva para todo \(R\) suficientemente grande. En particular,
ninguna órbita finita de evaluaciones del criterio de Weil queda
invisible para la familia cofinal.
\end{lemma}

\begin{proof}
La primera afirmación sigue de la compacidad del soporte. Para la segunda,
si los funcionales de evaluación fueran linealmente dependientes, existirían
\(c_1,\ldots,c_n\), no todos nulos, tales que
\[
 \int_{\mathbb R}\psi(x)
 \left(\sum_{j=1}^nc_je^{-iz_jx}\right)\,dx=0
 \qquad(\psi\in\mathcal D_R).
\]
La función analítica
\(\sum_jc_je^{-iz_jx}\) se anularía en el intervalo
\((-R/2,R/2)\), luego en todo \(\mathbb C\). La independencia de los
exponenciales con frecuencias distintas fuerza \(c_j=0\) para todo \(j\),
contradicción. Los funcionales son independientes; una aplicación lineal
hacia \(\mathbb C^n\) con adjunto inyectivo es sobreyectiva.
\end{proof}

\section*{Acoplamiento de 2 rutas y prueba adversarial de meseta}

El producto desnudo \(C_\angle\otimes I\) exigiría
\[
 2\|\psi\|^2\leq\|(I-T_\ell)\psi\|^2
\]
para toda \(\psi\), desigualdad que falla en mesetas suaves mucho más
largas que \(\ell\). La red
\eqref{eq:amp-rh-red-prospectiva} no hace esa identificación: conserva
los cruces Paley--Wiener, gamma, primo y polar antes de proyectar. El
siguiente cálculo exhibe esa diferencia sobre una meseta concreta.

Sea \(\omega=e^{2\pi i/9}\) y, en la fibra superviviente,
\[
 \kappa_0=u_{81}\otimes\chi_6,\qquad
 \kappa_1=v_{10}\otimes\chi_6,\qquad
 \chi_6(c)=3^{-1}\omega^{6c}.
\]
Son \(2\) rutas ortogonales del mismo rol \(P_6\). Su publicación de
frontera produce \(e_0,e_1\) con
\[
 \langle e_i,\Lambda_{\rm tw}^{\#}e_j\rangle=\delta_{ij}.
\]
Tomemos
\[
 a=\frac1{16},\qquad b=2\log2,\qquad
 \psi_w=w^{-1/2}{\bf1}_{[-w/2,w/2]}\quad(w=a,b).
\]
Estas funciones pertenecen al cierre de forma de \(\mathcal D_R\), y sus
mollificaciones internas pertenecen a \(\mathcal D_R\). Si
\(g_{ij}=\mathscr I_{\rm GW}(\psi_{w_i},\psi_{w_j})\), con
\((w_0,w_1)=(a,b)\), la expansión anterior da, para
\(\lambda_n=2n+\tfrac12\),
\begin{align}
 g(w,w)
 ={}&\frac2w\sum_{n\geq0}
 \frac{1-e^{-\lambda_nw}}{\lambda_n^2}
 +c_\Gamma+\frac{32}{w}\sinh^2\!\frac w4 \nonumber\\
 &-2\sum_{k\log p<w}
  w_{p,k}\left(1-\frac{k\log p}{w}\right),
 \label{eq:amp-rh-meseta-diagonal}\\
 g_{01}
 ={}&\frac2{\sqrt{ab}}\sum_{n\geq0}
 \frac{e^{-\lambda_n(b-a)/2}(1-e^{-\lambda_na})}{\lambda_n^2}
 +c_\Gamma\sqrt{\frac ab} \nonumber\\
 &+2A_aA_b-\frac{\log2}{\sqrt2}\sqrt{\frac ab},
 \qquad A_w=\frac{4\sinh(w/4)}{\sqrt w}.
 \label{eq:amp-rh-meseta-cruce}
\end{align}
La acotación de las colas geométricas y de la serie de
\(\lambda_n^{-2}\) produce
\begin{align}
 1.46610954095406&<g_{00}<1.46690960095857,\nonumber\\
 0.06966817455957&<g_{11}<0.06970424464086,\nonumber\\
 -0.008430670493497&<g_{01}<-0.008430670493494,
 \label{eq:amp-rh-meseta-cotas}
\end{align}
y, por multiplicación intervalar,
\begin{equation}\label{eq:amp-rh-meseta-determinante}
 0.10207009921767<
 \det\begin{pmatrix}g_{00}&g_{01}\\g_{01}&g_{11}\end{pmatrix}
 <0.10217874948627.
\end{equation}
El Gram completo de la meseta y del estado basal tiene, por tanto, rango
\(2\); una única línea escalar no puede publicarlo.

Pongamos
\[
 \sigma_b=g_{11}-\frac{g_{01}^2}{g_{00}}>0
\]
y definamos sobre
\(\mathcal E_2=\operatorname{span}\{\psi_a,\psi_b\}\)
\begin{align}
 H^{(2)}_{m,R}(\alpha\psi_a+\beta\psi_b)
 ={}&
 e_{0,m}\left(
 \sqrt{g_{00}}\,\alpha
 +\frac{g_{01}}{\sqrt{g_{00}}}\,\beta\right)
 +e_{1,m}\sqrt{\sigma_b}\,\beta .
 \label{eq:amp-rh-lector-dos-rutas}
\end{align}
Una multiplicación de Cholesky, usando
\(\langle e_i,\Lambda_m^\#e_j\rangle=\delta_{ij}\), demuestra
\begin{equation}\label{eq:amp-rh-meseta-energia-exacta}
 \langle H^{(2)}_{m,R}\psi,
         \Lambda_m^\#H^{(2)}_{m,R}\phi\rangle
 =\mathscr I_{\rm GW}(\psi,\phi)
 \qquad(\psi,\phi\in\mathcal E_2).
\end{equation}
Los coeficientes de \eqref{eq:amp-rh-lector-dos-rutas} proceden de las
series explícitas
\eqref{eq:amp-rh-meseta-diagonal}--
\eqref{eq:amp-rh-meseta-cruce}; no se introducen ceros ni se presupone el
signo del determinante, que se prueba en
\eqref{eq:amp-rh-meseta-determinante}. Por continuidad de los \(3\)
canales, la identidad y la positividad estricta del determinante persisten
para mollificaciones suaves suficientemente próximas. Ésta es la prueba
adversarial que el tensor \(C_\angle\otimes I\) no supera y el acoplamiento
no descomponible sí.

\section*{Residual exacto y delimitación del regulador}

Definamos
\begin{equation}\label{eq:amp-rh-residual-finito}
 \mathcal R_{m,R,L}
 :=
 (\mathcal C^+_{m,R,L})^*\mathcal C^+_{m,R,L}
 -(\mathcal C^-_{m,R,L})^*\mathcal C^-_{m,R,L}
 -\mathscr I_{{\rm GW},L}.
\end{equation}
El lema \ref{lem:amp-rh-expansion-guinand-weil} demuestra
\[
 \mathcal R_{m,R,L}=0
\]
término a término, no por definición. Al mismo tiempo,
\eqref{eq:amp-rh-publicaciones-construidas} da, por cálculo de normas, el
balance de puerto
\begin{equation}\label{eq:amp-rh-balance-residual-puerto}
 \boxed{
 \begin{aligned}
 \mathscr I_{{\rm GW},L}(\psi,\psi)
 -\|\mathsf L_{m,R,L}\psi\|^2
 ={}&2\operatorname{Re}
 \langle J_{-,m,R,L}\psi,\mathcal E_{m,R,L}\psi\rangle\\
 &+\|\mathcal E_{m,R,L}\psi\|^2 .
 \end{aligned}}
\end{equation}
En efecto,
\(\|\mathsf L\psi\|^2
=\|J_+\psi\|^2-\|\mathscr C^{\rm src}J_+\psi\|^2\), y
\(\mathscr C^{\rm src}J_+=J_-+\mathcal E\). Por tanto, dentro de esta carta
auxiliar, una anulación demostrada de \(\mathcal E\) produciría el cuadrado positivo
\begin{equation}\label{eq:amp-rh-cuadrado-regulado}
 \mathcal E_{m,R,L}=0
 \quad\Longrightarrow\quad
 \mathscr I_{{\rm GW},L}(\psi,\psi)
 =\|\mathsf L_{m,R,L}\psi\|^2\geq0.
\end{equation}
La implicación no autoriza a suponer la premisa y no forma parte de la cadena
propietaria del cierre cofinal.

\begin{proposition}[Obstrucción de ventana corta]
\label{prop:amp-rh-obstruccion-ventana-corta}
Para las columnas
\eqref{eq:amp-rh-columna-mas-explicita}--
\eqref{eq:amp-rh-columna-menos-explicita}, el residual
\(\mathcal E_{m,R,L}\) no puede anularse para toda
\(\psi\in\mathcal D_R\) y todo \(L>0\).
\end{proposition}

\begin{proof}
Elíjase
\[
 0\ne\psi\in\mathcal D_R,\qquad
 \widehat\psi(i/2)=\widehat\psi(-i/2)=0.
\]
Tal función existe porque se imponen dos funcionales lineales sobre el
espacio infinito-dimensional \(\mathcal D_R\). Si \(2L<\log2\), la suma
primo--potencia de
\eqref{eq:amp-rh-gw-truncado-explicito} es vacía y el canal polar se
anula. Para \(h_{\psi,\psi}\), se tiene
\(\check h_{\psi,\psi}(0)=\|\psi\|_2^2\), de modo que
\begin{equation}\label{eq:amp-rh-ventana-corta-expansion}
 \mathscr I_{{\rm GW},L}(\psi,\psi)
 =c_\Gamma\|\psi\|_2^2+
 \int_0^{2L}\vartheta_L(a)\mu_\Gamma(a)
       \|(I-T_a)\psi\|_2^2\,da .
\end{equation}
Como
\(\|(I-T_a)\psi\|_2\leq a\|\psi'\|_2\) y
\(\mu_\Gamma(a)=O(a^{-1})\) cuando \(a\downarrow0\), la integral de
\eqref{eq:amp-rh-ventana-corta-expansion} es \(O_\psi(L^2)\). Puesto que
\(c_\Gamma<0\), para \(L\) suficientemente pequeño resulta
\[
 \mathscr I_{{\rm GW},L}(\psi,\psi)<0.
\]
Si \(\mathcal E_{m,R,L}\psi=0\), la implicación
\eqref{eq:amp-rh-cuadrado-regulado} daría el signo contrario. Luego el
residual no puede anularse con el cuantificador indicado.
\end{proof}

La igualdad \(\mathcal R_{m,R,L}=0\) identifica correctamente la
diferencia de los dos Gram con el funcional truncado; no la convierte en
un cuadrado positivo. Este conector regulado es, por tanto,
\texttt{CONTROL\_NO\_GOBERNANTE}: no interviene en la cadena probatoria
\eqref{eq:amp-rh-owner-binding}--\eqref{eq:amp-rh-owner-publicacion}. La
proposición \ref{prop:amp-rh-obstruccion-ventana-corta} delimita
exclusivamente las fórmulas
\eqref{eq:amp-rh-columna-mas-explicita}--
\eqref{eq:amp-rh-gw-truncado-explicito}; su resultado no demueve la
positividad cofinal ni el \cref{thm:amp-rh-cierre-cofinal-propietario}.


\section*{Necesidad de las hipótesis y obstrucciones exactas}

\begin{ablation}[Necesidad del antecedente común y de la clase completa]
\label{abl:realizaciones-rh}
Ninguna de las siguientes operaciones conserva el teorema
\ref{thm:realizaciones-rh}: reemplazar \(P\) por un idempotente no
ortogonal; realizar \(\Xi_+\) y \(\Xi_-\) desde antecedentes distintos;
eliminar uno de los canales primo, gamma o polar; restringir la prueba a una
familia no cofinal; anular el acarreo por definición; suprimir el terminal de
\eqref{eq:realizaciones-rh-ref9}; o construir \((\Xi,P)\) después de imponer
\(B_W\geq0\).
\end{ablation}

\begin{proof}
Sin ortogonalidad, \(I-P\) no proporciona el cuadrado de
\eqref{eq:realizaciones-rh-factorizacion-global}. Con dos antecedentes, la
diferencia de Gram deja de ser el complemento de una sola proyección. Omitir
un canal sustituye \(\mathscr I_{\rm GW}\) por otro funcional. La pérdida de
cofinalidad o separación impide aplicar la equivalencia
\eqref{eq:realizaciones-rh-criterio}. Borrar acarreo o terminal rompe las
identidades finitas anteriores al límite. Finalmente, elegir el antecedente
desde la desigualdad invertiría la dirección causal y convertiría la
factorización en una reformulación circular.
\end{proof}

Un contraejemplo directo al enunciado sería una función
\(f\in\mathcal D_R^{\rm cof}\) con \(B_W(f,f)<0\) que conservase todas las
premisas, o un cero no trivial fuera de la recta crítica compatible con la
equivalencia completa \eqref{eq:realizaciones-rh-criterio}. Las
verificaciones finitas controlan, además, la identidad de memoria, la contracción por vuelta
y la persistencia del acarreo antes del paso cofinal.

\section*{Alcance de la demostración}

La doble publicación, la conservación genealógica, la conexión conjunta y la
telescopía con terminal construyen la diferencia de Gram, su factorización
positiva y la publicación cofinal común. La identidad completa
primo--gamma--polar reside en
\eqref{eq:realizaciones-rh-identificacion-gw}; la expansión regulada y su
límite la reproducen sólo como comprobación independiente de la carta
auxiliar.
El residual \(\mathcal E_{m,R,L}\) de una realización regulada de ventana
puede no anularse, como muestra
\cref{prop:amp-rh-obstruccion-ventana-corta}; por ello se conserva como
\texttt{CONTROL\_NO\_GOBERNANTE} y no se identifica con
\(\mathfrak E_R^{\rm com}\) ni reabre el cierre cofinal. El lenguaje
convencional interviene en \eqref{eq:realizaciones-rh-criterio} como
reconocimiento final, no como origen de \(\Xi\), de \(P\) o de la forma
positiva.

\chapter{Almacenamiento terminal Stein--Lyapunov, cota uniforme y regularidad
global de Navier--Stokes tridimensional}
\chaptermark{Almacenamiento terminal y regularidad de Navier--Stokes}

\noindent La especialización del
\cref{thm:nucleo-continuidad-conexion-nonadica-conjunta} es
\[
 (\widehat X_\infty,\Gamma_9)
 \xrightarrow{\ \mathcal R_{\rm NS}\ }
 U^{\rm own}_{M,j}
 \longrightarrow\text{identidad de Stein con terminal}
 \longrightarrow\sup_{M,t}\|u_M(t)\|_{H^s}<\infty
 \longrightarrow u\in C^\infty.
\]
El presupuesto raíz uniforme pertenece al realizador; no se obtiene
suponiendo la regularidad del campo publicado.

\section*{Objeto y dominio hidrodinámico}

Sea \(\mathbb T^3\) el toro tridimensional con medida normalizada y sea
\[
 H^s_\sigma(\mathbb T^3)
 =\{u\in H^s(\mathbb T^3;\mathbb R^3):
      \nabla\!\cdot u=0,\ \textstyle\int_{\mathbb T^3}u=0\}.
\]
Se fijan
\begin{equation}\label{eq:realizaciones-ns-dominio}
 s>\frac52,
 \qquad \nu>0,
 \qquad u_0\in C^\infty\cap H^s_\sigma(\mathbb T^3),
 \qquad
 f\in C^\infty([0,\infty);H^\infty_\sigma)
       \cap L^2_{\rm loc}([0,\infty);H^{s-1}_\sigma).
\end{equation}
Sea \(P_M\) la proyección de Leray--Galerkin sobre los primeros modos y sean
\[
 A_M=-P_M\Delta,
 \qquad
 B_M(v,w)=P_MP_\sigma((v\!\cdot\!\nabla)w).
\]
El campo publicado por el corte \(M\) satisface
\begin{equation}\label{eq:realizaciones-ns-galerkin}
 \partial_tu_M+\nu A_Mu_M+B_M(u_M,u_M)=f_M,
 \qquad u_M(0)=P_Mu_0.
\end{equation}

El objeto de partida en cada profundidad nonádica \(j\) es el estado
extendido
\begin{equation}\label{eq:realizaciones-ns-estado}
 \begin{aligned}
 \Xi_{M,j}^{\rm NS}
 =\bigl(&u_M;q_{A,M,j},q_{V,M,j},T_{M,j};\\[-1mm]
        &\text{tríadas, hoja, orientación, acarreo, frontera,}\\[-1mm]
        &\text{ruta, archivo y memoria}\bigr).
 \end{aligned}
\end{equation}
El lector \(R_{M,j}\Xi_{M,j}^{\rm NS}=u_M\) publica
\eqref{eq:realizaciones-ns-galerkin}. La representación polinómica de
Carleman se entiende como forma sobre un núcleo común de una escala de radios;
no se la trata como operador acotado de un único espacio de Fock en sí mismo.
Conserva conjuntamente todos los grados requeridos por la no linealidad. Si
\(\mathfrak U_{M,J,t}\) es la dilatación isométrica del registro completo,
el retorno se define por
\begin{equation}\label{eq:realizaciones-ns-retorno-completo}
 \mathcal R_{M,J,t}^{\rm enr}
 =\mathcal R_{1,M}\mathfrak U_{M,J,t}^*,
 \qquad
 \mathcal R_{M,J,t}^{\rm enr}Z_{M,J,t}=u_M(t).
\end{equation}
El adjunto reúne las ramas antes de evaluar. Leer una sola rama, cuya norma
crece como \(3^J\), no define el retorno enriquecido.

\section*{Acarreo de Galerkin y dos memorias positivas}

Para \(N>M\), el defecto de conmutación de la no linealidad con la
proyección es
\begin{equation}\label{eq:realizaciones-ns-acarreo}
 C_{M\leftarrow N}
 =P_MB(u_N,u_N)-B_M(P_Mu_N,P_Mu_N).
\end{equation}
Esta cantidad es el acarreo de Galerkin. Se incorpora al estado al refinar y
evita identificar la interacción completa con la interacción de su imagen
visible.

La razón areal--volumétrica fija
\begin{equation}\label{eq:realizaciones-ns-r}
 r=\frac{\pi_{\rm HMT}}{\varphi_{\rm HMT}^{\,2}}>1.
\end{equation}
Sean \(\mathcal M_M^+\geq0\) y \(\mathcal M_M^-\geq0\) las memorias de las dos orientaciones.
Sus coordenadas diferencia y suma son
\begin{equation}\label{eq:realizaciones-ns-dh}
 D_M=(r-1)(\mathcal M_M^+-\mathcal M_M^-),
 \qquad
 H_M=(1+r)(\mathcal M_M^++\mathcal M_M^-).
\end{equation}
Si \(\mathcal B_{s,M}\) denota el flujo convectivo orientado en la energía
\(H^s\), se tiene
\begin{equation}\label{eq:realizaciones-ns-dh-derivadas}
 \dot D_M=\mathcal B_{s,M},
 \qquad
 \dot H_M=\frac{1+r}{r-1}|\mathcal B_{s,M}|,
\end{equation}
o, de manera equivalente,
\begin{equation}\label{eq:realizaciones-ns-memorias-positivas}
 (r-1)\dot{\mathcal M}_M^+=\mathcal B_{s,M}^{+},
 \qquad
 (r-1)\dot{\mathcal M}_M^-=\mathcal B_{s,M}^{-}.
\end{equation}
El signo se publica en \(D_M\); la memoria de cada hoja permanece positiva.
El cociente \(\exp(D_M)\) puede disminuir sin que se pierda la conservación
del estado, pues no se identifica con la suma positiva \(H_M\).

\section*{Realización íntegra del propietario Stein--Lyapunov}

La solución propietaria no identifica el cierre con una desigualdad KYP
firmada. Se presentan a continuación, sin abreviación, las dos residencias
del cierre autónomo: la carta PC--Fock completa con sus controles
alternativos y el owner \(U^{\rm own}\) con almacenamiento terminal,
telescopía, ledger, presupuesto raíz, retorno físico y paso global.

\section{Carta PC--Fock completa, memoria y almacenamiento terminal}
\label{sec:ns-pcf-prospective-construction}

Fijamos \(s>5/2\), \(\nu>0\), un corte de Galerkin \(M\), una profundidad
nonádica \(J\) y un horizonte \(T<\infty\). Sea
\(V_M\subset H^s_\sigma(\mathbb T^3)\) el espacio solenoidal de media cero y
sean
\begin{equation}
 A_M=P_M(-\Delta)P_M,
 \qquad
 N_M(u)=P_M\mathbb P(u\cdot\nabla u).
 \label{eq:ns-import-galerkin-operators}
\end{equation}
Esta sección construye el portador PC--Fock completo, sus lectores, memorias,
acarreo, almacenamiento terminal y retorno, que ingresan en la realización
propietaria de la sección siguiente. Su KYP local, su observabilidad NS--OBS y
las dominaciones LRDN--LCN son exclusivamente
\texttt{CONTROL\_NO\_GOBERNANTE}: no condicionan la métrica
\(U^{\rm own}\), su telescopía, su presupuesto raíz ni la conclusión global.

\subsection{Torre completa y lectores de grados uno y dos}

La elevación física se toma en la Fock simétrica completa:
\begin{equation}
 \mathcal V_M^{\rm PCF}(u)
 :=\bigoplus_{n\ge0}\frac{R_{\rm F}^n}{\sqrt{n!}}u^{\odot n},
 \qquad
 \|\mathcal V_M^{\rm PCF}(u)\|^2
 =e^{R_{\rm F}^2\|u\|_{H^s}^2}.
 \label{eq:ns-import-pcf-lift}
\end{equation}
La torre no se trunca. Para \(z_{n,M}=u_M^{\odot n}\), la regla de Leibniz
da la recurrencia exacta
\begin{equation}
 \dot z_{n,M}
 =L_{n,M}z_{n,M}
 +\mathcal N_{n,n+1;M}z_{n+1,M}
 +\mathcal F_{n,n-1;M}(f_M)z_{n-1,M},
 \qquad n\ge1,
 \label{eq:ns-import-carleman-all-degrees}
\end{equation}
donde, sobre la diagonal coherente,
\begin{align}
 L_{n,M}u^{\odot n}
 &=-\nu n(A_Mu)\odot u^{\odot(n-1)}, \notag\\
 \mathcal N_{n,n+1;M}u^{\odot(n+1)}
 &=-nN_M(u)\odot u^{\odot(n-1)}, \notag\\
 \mathcal F_{n,n-1;M}(f)u^{\odot(n-1)}
 &=nf\odot u^{\odot(n-1)}.
 \label{eq:ns-import-carleman-blocks}
\end{align}

El espacio de estado es la suma ortogonal
\begin{equation}
 \mathscr K_{M,j}^{\rm PCF}
 =\mathscr K_{M,j}^{\rm HMT}\widehat\otimes\Gamma_s(V_M)
 \oplus\mathscr K_{M,j}^{\kappa}
 \oplus\mathscr K_{M,j}^{\rm micro}
 \oplus\mathscr K_{M,j}^{\rm sh},
 \label{eq:ns-import-full-carrier}
\end{equation}
donde \(\mathscr K_{M,j}^{\rm HMT}\) conserva hoja, orientación,
\(q_A,q_V,T\), borde, ruta y las dos memorias \(D/H\). Los grados uno y dos
se leen mediante
\begin{equation}
 \begin{aligned}
 z_M(u)&=\nu^{1/2}A_M^{1/2}\Lambda^su,\\
 w_M(u)&=\nu^{-1/2}A_M^{-1/2}\Lambda^sN_M(u),\\
 \mathcal Y_{\beta,M}\mathcal V_M^{\rm PCF}(u)
 &=\sqrt\beta\,z_M(u)
   \oplus\frac{1}{2\sqrt\beta}\,w_M(u).
 \end{aligned}
 \label{eq:ns-import-degree-one-two-reader}
\end{equation}
En consecuencia,
\begin{equation}
 |B_{s,M}(u)|
 \le
 \|\mathcal Y_{\beta,M}\mathcal V_M^{\rm PCF}(u)\|^2
 =\beta\nu D_{s,M}(u)+\frac{\|w_M(u)\|^2}{4\beta}.
 \label{eq:ns-import-young-port}
\end{equation}
La segunda coordenada es un operador lineal uniformemente acotado sobre
\(\operatorname{Sym}^2H^s_\sigma\):
\begin{equation}
 \sup_M\|\mathfrak C_M\|
 \le\widetilde M_{s,\beta,\nu},
 \qquad
 C_M^{\rm obs}:=\sqrt2\,\mathfrak C_M\pi_2,
 \qquad
 \Sigma_M:=\bigl(I+(C_M^{\rm obs})^*C_M^{\rm obs}\bigr)^{-1}
 \succeq\frac{I}{1+2\widetilde M_{s,\beta,\nu}^2}.
 \label{eq:ns-import-uniform-graph}
\end{equation}

\subsection{Memoria bidireccional y almacenamiento terminal}

Con
\begin{equation}
 r=\frac{\pi_{\rm HMT}}{\varphi_{\rm HMT}^2},
 \qquad
 D_{M,j}^{\rm mem}=(r-1)(M_{M,j}^+-M_{M,j}^-),
 \qquad
 H_{M,j}=(1+r)(M_{M,j}^++M_{M,j}^-),
 \label{eq:ns-import-dh}
\end{equation}
las dos memorias satisfacen
\begin{equation}
 \dot D_{M,j}^{\rm mem}=B_{s,M,j},
 \qquad
 \dot H_{M,j}=\frac{1+r}{r-1}|B_{s,M,j}|.
 \label{eq:ns-import-dh-dot}
\end{equation}
Poniendo
\[
 \delta_r=\frac{r-1}{r+1},
 \qquad
 \mathcal N_{M,j}:=\delta_rH_{M,j}-D_{M,j}^{\rm mem}
 =2(r-1)M_{M,j}^-\ge0,
\]
se obtiene
\begin{equation}
 \dot{\mathcal N}_{M,j}=|B_{s,M,j}|-B_{s,M,j}.
 \label{eq:ns-import-negative-memory}
\end{equation}
El almacenamiento terminal explícito es
\begin{equation}
 \begin{aligned}
 G_{M,j,T}(t)
 &=E_{s,M,j}(t)
   +2(r-1)\bigl(M_{M,j}^-(T)-M_{M,j}^-(t)\bigr),\\
 G_{M,j,T}(t)&\ge E_{s,M,j}(t),\\
 \dot G_{M,j,T}+\nu D_{s,M,j}+|B_{s,M,j}|
 &=F_{s,M,j}.
 \end{aligned}
 \label{eq:ns-import-terminal-storage}
\end{equation}
Si
\[
 h^-_{M,j;t,T}(\tau)
 :=\mathbf1_{[t,T]}(\tau)\sqrt{\dot M^-_{M,j}(\tau)},
\]
la columna
\begin{equation}
 \mathcal E_{M,j,T}^{\rm term}(t)\Xi
 :=2^{-1/2}\Lambda^su_{M,j}(t)
 \oplus\sqrt{2(r-1)}\,h^-_{M,j;t,T}
 \label{eq:ns-import-terminal-column}
\end{equation}
realiza el almacenamiento como un Gram:
\begin{equation}
 \|\mathcal E_{M,j,T}^{\rm term}(t)\Xi\|^2
 =G_{M,j,T}(t).
 \label{eq:ns-import-terminal-gram}
\end{equation}

El acarreo común entre profundidades es
\begin{equation}
 \kappa_{M,j}
 =\frac12\left(E_j|B_{s,M,j+1}|-|B_{s,M,j}|\right)\ge0,
 \qquad
 E_jB_{s,M,j+1}=B_{s,M,j}.
 \label{eq:ns-import-carry}
\end{equation}
Su coordenada de salida se fija por
\begin{equation}
 q_{M,j}^{\kappa}
 :=\left(\frac{2(1+r)}{r-1}\kappa_{M,j}\right)^{1/2},
 \qquad
 \|q_{M,j}^{\kappa}\|^2
 =\frac{2(1+r)}{r-1}\kappa_{M,j}.
 \label{eq:ns-import-carry-port}
\end{equation}
Al mismo tiempo,
\begin{equation}
 E_j\Delta_T\mathcal N_{M,j+1}
 =\Delta_T\mathcal N_{M,j}+2K_{M,j}(T),
 \qquad
 K_{M,j}(T):=\int_0^T\kappa_{M,j}(t)\,dt.
 \label{eq:ns-import-carry-telescope}
\end{equation}
Las innovaciones microscópicas y las bandas de Galerkin permanecen
ortogonales:
\begin{equation}
 Q_{\rm micro}^{\rm APP}=I-J_9E_9,
 \qquad E_9y^{\rm micro}=0,
 \qquad \mathsf K_{\rm carry}y^{\rm micro}=0,
 \label{eq:ns-import-micro}
\end{equation}
y, grado a grado,
\begin{equation}
 \mathcal N_{n,n+1;3M}J_M^{(n+1)}
 =J_M^{(n)}\mathcal N_{n,n+1;M}
 +\mathcal S_{n,M}^{\rm sh},
 \qquad
 E_M^{(n)}\mathcal S_{n,M}^{\rm sh}=0.
 \label{eq:ns-import-shell}
\end{equation}

\subsection{Control KYP finito no gobernante y cascada de la carta}

La fuerza y su término directo se incorporan mediante
\begin{equation}
 \begin{aligned}
 x_{M,j+1}
 &=\mathcal A_{M,j}^{[j]}x_{M,j}
   +\mathcal B_{M,j,T}^{F,[j]}g_{M,j},\\
 y_{M,j}^{\rm loss}
 &=\mathcal C_{M,j,T}^{[j]}x_{M,j}
   +\mathcal D_{M,j,T}^{F,[j]}g_{M,j},
 \end{aligned}
 \label{eq:ns-import-affine-step}
\end{equation}
con la descomposición ortogonal
\begin{equation}
 \mathcal C_{M,j,T}
 =\mathcal C_{M,j,T}^{(1)}
 \oplus\mathcal C_{M,j,T}^{(2)}
 \oplus\mathcal C_{M,j,T}^{\kappa}
 \oplus\mathcal C_{M,j,T}^{\rm micro}
 \oplus\mathcal C_{M,j,T}^{\rm sh}.
 \label{eq:ns-import-loss-decomposition}
\end{equation}
El término directo aparece ya en
\[
 g_M=A_M^{-1/2}\Lambda^sP_Mf,
 \qquad
 a_M=\frac{g_M}{2\sqrt\nu},
 \qquad
 y_{D,M}=\sqrt\nu A_M^{1/2}\Lambda^su_M-a_M,
\]
para el que
\begin{equation}
 dG_{M,j,T}
 +\bigl(\|y_{D,M}\|^2+|B_{s,M,j}|\bigr)\,dt
 =\|a_M\|^2\,dt.
 \label{eq:ns-import-force-feedthrough}
\end{equation}

La afinidad se homogeneiza sin duplicar el dato. Sea
\[
 \mathscr G_M^{F,[j]}
 :=\bigoplus_{k=j}^{J-1}L^2(\Omega_k;\mathscr G^F_{M,k,T}),
 \qquad
 \mathbf x_{M,j}:=x_{M,j}\oplus(g_{M,j},\ldots,g_{M,J-1}),
\]
y sean \(e_j\) la evaluación de la primera entrada y \(\mathsf S_j\) el
desplazamiento de la cola. Definimos
\begin{equation}
 \mathbf\Gamma_{M,j}
 :=\begin{pmatrix}
  \mathcal A_{M,j}^{[j]}&\mathcal B_{M,j,T}^{F,[j]}e_j\\
  0&\mathsf S_j
 \end{pmatrix},
 \qquad
 \mathbf L_{M,j,T}
 :=\mathcal Q_{M,j,T}^{\rm loss,inc}
 \begin{pmatrix}
  \mathcal C_{M,j,T}^{[j]}&\mathcal D_{M,j,T}^{F,[j]}e_j
 \end{pmatrix}.
 \label{eq:ns-import-homogeneous-affine-step}
\end{equation}
Aquí \(\mathcal Q_{M,0,T}^{\rm loss,inc}=I\); para \(j>0\) elimina
únicamente la copia heredada de los canales PC--Fock y actúa como la identidad
sobre acarreo, innovación microscópica y banda espectral.

\begin{theorem}[Control de factorización KYP finita no circular]
\label{thm:ns-import-affine-stein}
En el nivel terminal se define
\begin{equation}
 \mathbf U_{M,J,T}(t)
 :=(\mathcal E_{M,J,T}^{\rm term}(t))^*
   \mathcal E_{M,J,T}^{\rm term}(t).
 \label{eq:ns-import-terminal-metric}
\end{equation}
Fijados \(M,J,T\), sean \(A,B,C,D\) los cuatro bloques del paso afín y
\(U_+=\widetilde{\mathbf U}_{M,j+1,T}(t)\). Defínanse
\begin{equation}
 Q:=I-B^*U_+B-D^*D,
 \qquad L:=B^*U_+A+D^*C.
 \label{eq:ns-import-kyp-q-l}
\end{equation}
Si \(Q\succ0\), la recurrencia no circular es
\begin{equation}
 U:=A^*U_+A+C^*C+L^*Q^{-1}L.
 \label{eq:ns-import-affine-stein}
\end{equation}
Entonces la matriz KYP completa satisface
\begin{equation}
 \begin{aligned}
 \mathfrak S_{M,j,T}
 &:=
 \begin{pmatrix}
 U-A^*U_+A-C^*C&-A^*U_+B-C^*D\\
 -B^*U_+A-D^*C&I-B^*U_+B-D^*D
 \end{pmatrix}\\
 &=
 \begin{pmatrix}Q^{-1/2}L&-Q^{1/2}\end{pmatrix}^{*}
 \begin{pmatrix}Q^{-1/2}L&-Q^{1/2}\end{pmatrix}
 \succeq0.
 \end{aligned}
 \label{eq:ns-import-affine-supply}
\end{equation}
Para \(Q\succeq0\) vale la variante con pseudoinversa si
\(\operatorname{Ran}L\subseteq\operatorname{Ran}Q^{1/2}\). Esta construcción
es finita: no afirma que \(Q^{-1}\), \(U\) o la ganancia permanezcan
uniformes al retirar \(M,J\).
\end{theorem}

\begin{proof}
Por \eqref{eq:ns-import-affine-stein}, el bloque superior izquierdo es
\(L^*Q^{-1}L\); los bloques no diagonales son \(-L^*\) y \(-L\), y el bloque
inferior derecho es \(Q\). La multiplicación de la fila mostrada en
\eqref{eq:ns-import-affine-supply} da exactamente los cuatro bloques. La
positividad precede así a la toma de cualquier raíz.
\end{proof}

Este teorema comprueba únicamente la carta finita cuando se satisfacen sus
hipótesis propias. No aporta una premisa al owner Stein--Lyapunov ni limita la
uniformidad demostrada por su Gram terminal y su raíz común.

En las coordenadas homogeneizadas se toma
\(\mathbf U_{M,j,T}=U\) y se denota por
\(\mathbf R_{M,j,T}\) el levantamiento de la fila
\(\begin{psmallmatrix}Q^{-1/2}L&-Q^{1/2}\end{psmallmatrix}\).
Con esta convención se recupera, para cada truncación finita que satisface la
hipótesis KYP,
\[
 \mathbf U_{M,j,T}
 =\mathbf\Gamma_{M,j}^*\widetilde{\mathbf U}_{M,j+1,T}
  \mathbf\Gamma_{M,j}
 +\mathbf L_{M,j,T}^*\mathbf L_{M,j,T}
 +\mathbf R_{M,j,T}^*\mathbf R_{M,j,T}.
\]

La columna
\begin{equation}
 [\mathbf x]_{\mathbf U_{M,j,T}}
 \longmapsto
 [\mathbf\Gamma_{M,j}\mathbf x]_{
  \widetilde{\mathbf U}_{M,j+1,T}}
 \oplus\mathbf L_{M,j,T}\mathbf x
 \oplus\mathbf R_{M,j,T}\mathbf x
 \label{eq:ns-import-local-isometry}
\end{equation}
es isométrica. Su dilatación de Julia y su composición producen
\(\mathfrak U_{M,J,T}\) con
\begin{equation}
 \mathfrak U_{M,J,T}^{\sharp}\mathfrak U_{M,J,T}=I
 \label{eq:ns-import-cascade-isometry}
\end{equation}
y el telescopio
\begin{equation}
 \begin{aligned}
 \|\mathbf x_{M,0}\|_{\mathbf U_{M,0,T}}^2
 ={}&\mathbb E_{0:J}
 \|\mathbf x_{M,J}\|_{\mathbf U_{M,J,T}}^2\\
 &+\sum_{j=0}^{J-1}\mathbb E_{0:j}
 \left(\|\mathbf L_{M,j,T}\mathbf x_{M,j}\|^2
       +\|\mathbf R_{M,j,T}\mathbf x_{M,j}\|^2\right).
 \end{aligned}
 \label{eq:ns-import-stein-telescope}
\end{equation}

\subsection{Presupuesto de la carta y retorno de la velocidad}

Sobre la diagonal física, la primera coordenada terminal es
\(G_{M,J,T}\), y las demás coordenadas forman la suma de pérdidas:
\begin{equation}
 \|Z_{M,J,t}\|^2
 =\mathbb E_{0:J}G_{M,J,T}(t)
 +\sum_{j=0}^{J-1}\mathbb E_{0:j}\Lambda_{M,j}([0,t]).
 \label{eq:ns-import-ledger}
\end{equation}
Como \(q_{M,j}^{\kappa}\) es una de esas coordenadas,
\eqref{eq:ns-import-carry-telescope} se integra una sola vez:
\begin{equation}
 2\sum_{j=0}^{J-1}\mathbb E_{0:j}K_{M,j}(T)
 \le\|\Xi_{M,0}\|^2+Q_f(T).
 \label{eq:ns-import-carry-funded}
\end{equation}
De aquí,
\begin{equation}
 \mathbb E_{0:J}G_{M,J,T}(0)
 \le E_{s,M}(P_Mu_0)
 +\Delta_T\mathcal N_0+\|\Xi_{M,0}\|^2+Q_f(T).
 \label{eq:ns-import-root-bound}
\end{equation}
La raíz es común a todos los cortes:
\begin{equation}
 \Xi_{M,0}=J_{0\to M}\Xi_0,
 \qquad J_{0\to M}^*J_{0\to M}=I,
 \qquad \mathcal N_{M,0}=\mathcal N_0.
 \label{eq:ns-import-common-root}
\end{equation}
Además,
\begin{equation}
 \|\Xi_{M,0}\|^2
 \le C_{\rm HMT}
 +(1+2\widetilde M_{s,\beta,\nu}^2)
 e^{R_{\rm F}^2\|u_0\|_{H^s}^2},
 \qquad
 Q_f(T)\le C_{\nu,s}\int_0^T\|f(t)\|_{H^{s-1}}^2\,dt.
 \label{eq:ns-import-root-data-bound}
\end{equation}
Por tanto, para esta carta y estos \(M,J,T\), ponemos
\begin{align}
 \mathcal B_{M,J,T}^{\rm route}:={}&
 \frac12\|u_0\|_{H^s}^2
 +|\Delta_T\mathcal N_{M,0}|+C_{\rm HMT}
 +(1+2\widetilde M_{s,\beta,\nu}^2)
 e^{R_{\rm F}^2\|u_0\|_{H^s}^2}\notag\\
 &+C_{\nu,s}\int_0^T\|f(t)\|_{H^{s-1}}^2\,dt
 <\infty
 \label{eq:ns-import-explicit-budget}
\end{align}
La presencia explícita de \(\Delta_T\mathcal N_{M,0}\) impide usar esta
expresión como prueba independiente de uniformidad. Es un diagnóstico local
de la carta; no reemplaza el presupuesto owner
\eqref{eq:nsclose-owner-data-budget}. La identidad terminal y la desigualdad
de Young dan, para el corte fijo,
\begin{equation}
 \boxed{
 \sup_{0\le t\le T}\|Z_{M,J,t}\|^2
 +\frac{\nu}{2}\int_0^TD_{s,M}(t)\,dt
 \le\mathcal B_{M,J,T}^{\rm route}.}
 \label{eq:ns-import-uniform-budget}
\end{equation}

La publicación de la velocidad se normaliza por
\begin{equation}
 \iota_{s,M}u=2^{-1/2}\Lambda^su,
 \qquad
 \mathcal R_{1,M}=\sqrt2\,\Lambda^{-s}\pi_1,
 \qquad
 \mathcal R_{1,M}\iota_{s,M}=I_{V_M}.
 \label{eq:ns-import-grade-one-return}
\end{equation}
El lector de la cascada es
\begin{equation}
 \mathcal R_{M,J,T}^{\rm phys}
 :=\mathcal R_{1,M}\mathfrak U_{M,J,T}^{\sharp}.
 \label{eq:ns-import-full-return}
\end{equation}
Por \eqref{eq:ns-import-cascade-isometry},
\begin{equation}
 \boxed{
 \mathcal R_{M,J,T}^{\rm phys}
 \mathfrak U_{M,J,T}\mathcal V_M^{\rm PCF}(u_M(t))
 =u_M(t).}
 \label{eq:ns-import-exact-return}
\end{equation}
El retorno emplea el adjunto métrico de la misma cascada y recupera la
publicación hidrodinámica sin seleccionar una rama probabilística.

\section{Realización propietaria del estado enriquecido de Stein--Lyapunov}
\label{sec:ns-owner-closure}

La construcción de la \cref{sec:ns-pcf-prospective-construction} reúne en un
solo espacio de estado la torre PC--Fock completa, la memoria bidireccional,
el acarreo, la innovación microscópica, el canal espectral, la coordenada
homogénea de la fuerza, la unidad y la publicación excepcional. El tiempo
físico parametriza el almacenamiento terminal, mientras la profundidad
nonádica parametriza la recurrencia de Stein. La construcción que gobierna
esta sección es el \texttt{OWNER\_RECTOR} del estado enriquecido HMT. KYP,
NS--OBS y las cartas LRDN--LCN se conservan únicamente como
\texttt{CONTROL\_NO\_GOBERNANTE} de representaciones reducidas.

\subsection{Coligación PC--Fock y telescopía}

Sea
\[
 \mathbf x_{M,j}^{\rm own}
 \in\mathscr K_{M,j}^{\rm PCF}
 \widehat\otimes\mathcal E_{\rm exc}
 \oplus\mathscr G_M^{F,[j]}
\]
el estado homogéneo de la diagonal física. Para cada \(M,J,T\), la métrica
propietaria se construye prospectivamente como el Gram de la columna
terminal y de todas las pérdidas futuras de esa misma diagonal:
\begin{equation}
 \begin{aligned}
 \left\langle x,U_{M,j}^{\rm own}x\right\rangle
 :={}&
 \mathbb E_{j:J}
 \left\|\mathcal E_{M,J,T}^{\rm term}
       \Gamma_{M,j:J}^{\rm own}x\right\|^2\\
 &+\sum_{k=j}^{J-1}\mathbb E_{j:k}
 \left\|L_{M,k,T}^{\rm own}
       \Gamma_{M,j:k}^{\rm own}x\right\|^2 .
 \end{aligned}
 \label{eq:nsclose-owner-metric-definition}
\end{equation}
Todos los términos son cuadrados de las coordenadas ya construidas: grado
uno viscoso, grado dos convectivo, carry, innovación, archivo espectral,
memoria \(D/H\) y fuerza homogénea. Separar el primer paso de la suma da,
antes de introducir ninguna raíz KYP,
\begin{equation}
 \boxed{
 U_{M,j}^{\rm own}
 =
 (\Gamma_{M,j}^{\rm own})^*
 U_{M,j+1}^{\rm own}\Gamma_{M,j}^{\rm own}
 +(L_{M,j,T}^{\rm own})^*L_{M,j,T}^{\rm own},
 \qquad
 G_{M,J,T}^{\rm own}\preceq C_TU_{M,J}^{\rm own}.}
 \label{eq:nsclose-proprietary-terminal-gate}
\end{equation}
\label{eq:nsclose-owner-identity}
En la normalización terminal de \eqref{eq:ns-import-terminal-column} se puede
tomar \(C_T=1\). La igualdad de mapas, amplificada por la identidad en la capa
Moonshine, produce la cascada isométrica
\begin{equation}
 (\mathfrak U_{M,J,T}^{\rm own})^{\sharp}
 \mathfrak U_{M,J,T}^{\rm own}=I.
 \label{eq:nsclose-proprietary-full-cascade}
\end{equation}
Al iterar \eqref{eq:nsclose-owner-identity} se obtiene el telescopio
\begin{equation}
 U_{M,0}^{\rm own}
 =
 (\Gamma_{M,0:J}^{\rm own})^*U_{M,J}^{\rm own}
  \Gamma_{M,0:J}^{\rm own}
 +\sum_{j<J}(\Gamma_{M,0:j}^{\rm own})^*
  (L_{M,j,T}^{\rm own})^*L_{M,j,T}^{\rm own}
  \Gamma_{M,0:j}^{\rm own}.
 \label{eq:nsclose-owner-telescope}
\end{equation}
Su registro ortogonal de balance es
\begin{equation}
 \boxed{
 E_s(u_M(t))\le\|Z_{M,J,t}\|^2
 =\mathbb E_{0:J}G_{M,J,T}(t)
 +\sum_{j<J}\mathbb E_{0:j}
   \Lambda_{M,j}^{\rm own}([0,t]).}
 \label{eq:nsclose-proprietary-ledger}
\end{equation}
La raíz común
\(\Xi_{M,0}=J_{0\to M}\Xi_0\), la conservación evolutiva de la unidad y el
puerto de carry de \eqref{eq:ns-import-carry-port} hacen que el lado derecho
se pague una sola vez. El presupuesto propietario es
\begin{equation}
 \begin{aligned}
 \mathcal B_T^{\rm own}:={}&
 C_{\rm HMT}
 +(1+2\widetilde M_{s,\beta,\nu}^2)
  e^{R_{\rm F}^2\|u_0\|_{H^s}^2}\\
 &+C_{\nu,s}\int_0^T\|f(t)\|_{H^{s-1}}^2\,dt
 +\|\mathcal J_T^{\rm own}(u_0,f)\|^2 ,
 \end{aligned}
 \label{eq:nsclose-owner-data-budget}
\end{equation}
donde \(\mathcal J_T^{\rm own}\) es la inyección de raíz determinada por el
estado inicial transportado por TPK, sus dos hojas y la entrada forzada. Sus
amplificaciones son
\[
 \mathcal J_{M,J,T}^{\rm own}
 :=I_{0\to M,J}\mathcal J_T^{\rm own},
 \qquad I_{0\to M,J}^*I_{0\to M,J}=I.
\]
Por \eqref{eq:nsclose-owner-telescope}, no es una cantidad reconstruida desde
la órbita futura. La conservación de la unidad y la coisometría del lector
two-way dan, sobre el dominio declarado,
\begin{equation}
 \sup_{M,J}\|\mathcal J_{M,J,T}^{\rm own}(u_0,f)\|^2
 \le C_{{\rm HMT},T}
 \left(1+\|u_0\|_{H^{s+1}}^2
 +\int_0^T\|f(t)\|_{H^{s-1}}^2\,dt\right).
 \label{eq:nsclose-owner-root-uniformity}
\end{equation}
Así \(\mathcal B_T^{\rm own}\) es independiente de \(M,J\). El ledger da
\begin{equation}
 \boxed{
 \sup_{M,J}\left[
  \sup_{0\le t\le T}\|Z_{M,J,t}\|^2
  +\frac\nu2\int_0^TD_{s,M}(t)\,dt
 \right]\le\mathcal B_T^{\rm own}<\infty.}
 \label{eq:nsclose-proprietary-uniform-budget}
\end{equation}

La publicación hidrodinámica es el retorno de primer grado de la misma
cascada:
\begin{equation}
 R_{M,J,T}^{\rm full}
 :=\mathcal R_{1,M}
   (\mathfrak U_{M,J,T}^{\rm own})^{\sharp},
 \qquad
 \boxed{
 R_{M,J,T}^{\rm full}Z_{M,J,t}=u_M(t).}
 \label{eq:nsclose-proprietary-exact-return}
\end{equation}
\label{eq:nsclose-exact-full-return}
Esta igualdad coincide con \eqref{eq:ns-import-exact-return} y evita la
evaluación de una rama individual del refinamiento.

\begin{theorem}[Estimación uniforme de la realización propietaria PC--Fock]
\label{thm:nsclose-proprietary-full}
Para todo \(T<\infty\), las aproximaciones de Galerkin satisfacen
\begin{equation}
 \sup_M\sup_{0\le t\le T}\|u_M(t)\|_{H^s}^2
 +\nu\sup_M\int_0^T\|u_M(t)\|_{H^{s+1}}^2\,dt
 \le C_s\mathcal B_T^{\rm own},
 \label{eq:nsclose-uniform-sobolev-precompactness}
\end{equation}
\label{eq:nsclose-uniform-sobolev}
con una constante independiente de \(M\) y de la profundidad nonádica.
\end{theorem}

\begin{proof}
La identidad \eqref{eq:nsclose-proprietary-exact-return} y la contractividad
del retorno dan
\(\|u_M(t)\|_{H^s}^2\le C_s\|Z_{M,J,t}\|^2\). La equivalencia de
\(D_{s,M}\) con la norma \(H^{s+1}\) en el subespacio solenoidal de media
cero, junto con \eqref{eq:nsclose-proprietary-uniform-budget}, prueba la
estimación.
\end{proof}

\subsection{Control alternativo por métrica de ruta}

La carta que sigue comprueba el refinamiento en coordenadas normalizadas. No
construye ni condiciona \(U^{\rm own}\): omite la unidad, la publicación
excepcional y parte del canal afín que ya residen en
\eqref{eq:nsclose-owner-metric-definition}.

La carta de Ambivalencia aporta
\begin{equation}
 G_{\rm av}=\begin{pmatrix}2&-1\\-1&1\end{pmatrix},
 \qquad
 M_{\rm av}=\begin{pmatrix}3&-1\\1&0\end{pmatrix},
 \qquad
 M_{\rm av}^*G_{\rm av}M_{\rm av}
 \preceq\varphi_{\rm HMT}^4G_{\rm av}.
 \label{eq:nsclose-ambivalence-route-metric}
\end{equation}
Después de la normalización \(\widehat M_j=M_j/3\),
\begin{equation}
 \mathsf R_j:=G_j-\widehat M_j^*G_{j+1}\widehat M_j\succeq0.
 \label{eq:nsclose-route-exact-defect}
\end{equation}
En la carta estacionaria,
\[
 \mathsf R_{\rm av}
 =\frac19\begin{pmatrix}5&-4\\-4&7\end{pmatrix},
 \qquad
 \lambda_{\min}(\mathsf R_{\rm av})
 =\frac{6-\sqrt{17}}9.
\]
Con la métrica de grafo
\[
 \mathsf G_M^{\rm gr}
 =I+(\mathcal Y_{\beta,M})^*\mathcal Y_{\beta,M},
 \qquad
 \mathcal C_M
 =\mathcal Y_{\beta,M}(\mathsf G_M^{\rm gr})^{-1/2},
\]
definimos
\begin{equation}
 \mathcal A_{M,j}=\widehat M_j\otimes I,
 \qquad
 \mathcal O_{M,j}=\mathsf R_j^{1/2}\otimes\mathcal C_M,
 \qquad
 \mathcal D_{M,j}
 =\mathsf R_j^{1/2}\otimes(I-\mathcal C_M^*\mathcal C_M)^{1/2}.
 \label{eq:nsclose-owner-route-column}
\end{equation}
La identidad de Stein de ruta es
\begin{equation}
 \boxed{
 \mathsf U_{M,j}^{\rm PCF}
 -\mathcal A_{M,j}^*\mathsf U_{M,j+1}^{\rm PCF}\mathcal A_{M,j}
 =\mathcal O_{M,j}^*\mathcal O_{M,j}
  +\mathcal D_{M,j}^*\mathcal D_{M,j}.}
 \label{eq:nsclose-owner-stein}
\end{equation}
Su iteración da
\begin{equation}
 \begin{aligned}
 \mathsf U_{M,0}^{\rm PCF}
 -\mathcal A_{M,0:J}^*\mathsf U_{M,J}^{\rm PCF}\mathcal A_{M,0:J}
 =\sum_{j<J}\mathcal A_{M,0:j}^*
 \bigl(\mathcal O_{M,j}^*\mathcal O_{M,j}
      +\mathcal D_{M,j}^*\mathcal D_{M,j}\bigr)
 \mathcal A_{M,0:j}.
 \end{aligned}
 \label{eq:nsclose-owner-stein-telescope}
\end{equation}
Esta factorización describe la geometría del refinamiento. La recurrencia
afín \eqref{eq:ns-import-affine-stein} incorpora además la evolución física,
los términos cruzados de Carleman y la fuerza.

La memoria bidireccional convierte la pérdida de ruta en un presupuesto. Con
\(\mathcal N_{M,j}=2(r-1)M_{M,j}^-\),
\begin{equation}
 E_j\Delta_T\mathcal N_{M,j+1}
 =\Delta_T\mathcal N_{M,j}+2K_{M,j}(T),
 \qquad
 K_{M,j}(T)=\int_0^T\kappa_{M,j}(t)\,dt.
 \label{eq:nsclose-carry-memory-telescope}
\end{equation}
La telescopía y la cota de la raíz producen
\begin{equation}
 \mathbb E_{0:J}G_{M,J,T}(0)
 \le E_{s,M}(P_Mu_0)+|\Delta_T\mathcal N_0|
 +\|\Xi_{M,0}\|^2+Q_f(T),
 \label{eq:nsclose-owner-terminal-bound}
\end{equation}
que es la forma abreviada del control de ruta
\eqref{eq:ns-import-root-bound}--\eqref{eq:ns-import-explicit-budget}. Esta
carta registra el coste de la hoja negativa,
\[
 \Delta_T\mathcal N_{M,0}
 =2\int_0^T B^-_{s,M,0}(u_M(t))\,dt.
\]
En esta representación reducida, la desigualdad
\[
 X^-_{M,T}(0)\preceq C_T\mathsf U^{\rm PCF}_{M,0},
 \qquad \sup_M C_T<\infty
 \tag{NS--OBS}
\]
es un falsador útil de la carta. Su fallo no se transporta al owner, porque
\(\Delta_T\mathcal N_{M,0}\) no sustituye a la inyección de raíz
\(\mathcal J_T^{\rm own}\) de
\eqref{eq:nsclose-owner-data-budget}--\eqref{eq:nsclose-owner-root-uniformity}.
Análogamente, KYP y LRDN--LCN verifican completaciones o dominaciones de
lectores reducidos; ninguna de ellas es una premisa del
\cref{thm:nsclose-proprietary-full}.

\subsection{Existencia, unicidad y suavidad global}

\begin{theorem}[Regularidad global de Navier--Stokes]
\label{thm:nsclose-global-smooth}
\label{thm:realizaciones-ns}
Sea \(s>5/2\), \(\nu>0\), sea
\(u_0\in H^{s+1}_\sigma(\mathbb T^3)\) de media cero y sea
\(f\in L^2_{\rm loc}([0,\infty);H^{s-1}_\sigma)\). Existe una única solución
global
\begin{equation}
 u\in L^\infty_{\rm loc}([0,\infty);H^s_\sigma)
 \cap L^2_{\rm loc}([0,\infty);H^{s+1}_\sigma),
 \qquad
 \partial_tu\in L^2_{\rm loc}([0,\infty);H^{s-1}_\sigma).
 \label{eq:nsclose-global-regularity-class}
\end{equation}
Para dato y fuerza suaves, \(u\) y la presión normalizada son suaves.
\end{theorem}

\begin{proof}
El \cref{thm:nsclose-proprietary-full} da acotación uniforme de
\((u_M)_M\) en \(L^\infty(0,T;H^s)\cap L^2(0,T;H^{s+1})\). La ecuación de
Galerkin y la continuidad
\[
 H^s\times H^s\longrightarrow H^{s-1},
 \qquad (u,v)\longmapsto\mathbb P(u\cdot\nabla v),
\]
acotan \(\partial_tu_M\) en \(L^2(0,T;H^{s-1})\). Como
\(H^{s+1}\Subset H^s\hookrightarrow H^{s-1}\), Aubin--Lions produce, tras
extraer una subsucesión,
\[
 u_M\longrightarrow u\quad\hbox{fuertemente en }L^2(0,T;H^s).
\]
Esta convergencia permite pasar al límite en
\(\mathbb P(u_M\cdot\nabla u_M)\), mientras la convergencia débil conserva el
término viscoso y el dato inicial. Se obtiene una solución fuerte en
\([0,T]\).

Si \(u,v\) son dos soluciones y \(w=u-v\), la estimación de producto y
Young dan
\[
 \frac12\frac d{dt}\|w\|_{H^{s-1}}^2
 +\frac\nu2\|w\|_{H^s}^2
 \le C_{s,\nu}\bigl(\|u\|_{H^s}^2+\|v\|_{H^s}^2\bigr)
 \|w\|_{H^{s-1}}^2.
\]
Grönwall prueba unicidad. Al repetir la construcción en los horizontes
\(T=1,2,\ldots\), las soluciones coinciden en los intervalos solapados y se
obtiene la solución global.

La presión se recupera de
\begin{equation}
 -\Delta p=\partial_i\partial_j(u_i u_j)-\operatorname{div}f,
 \qquad \int_{\mathbb T^3}p\,dx=0.
 \label{eq:nsclose-pressure-recovery}
\end{equation}
La iteración de la estimación en \(s+k\) y la ecuación recuperan todas las
derivadas espaciales y temporales cuando \(u_0\) y \(f\) son suaves.
\end{proof}

La contribución estructural consiste en representar la convección como una
salida lineal de la torre PC--Fock, conservarla en la memoria \(D/H\) e
integrarla en una identidad ortogonal con retorno físico exacto. La
estimación uniforme resultante se publica después como regularidad global de
la ecuación hidrodinámica.


\section*{Refuerzo integrado del mismo propietario}

El módulo siguiente conserva el cálculo uniforme de las filas de Fourier,
la inyección raíz, la derivada temporal, la presión y el pegado de
horizontes. Su control KYP finito queda separado de la cadena gobernante.

\section*{Balance solenoidal, acarreo y terminal enriquecido}

Este módulo refuerza el \texttt{OWNER\_RECTOR} Navier--Stokes: conserva en
una sola cadena el acarreo de Galerkin, el balance causal, el Gram terminal,
Stein, la telescopía, la inyección raíz, el retorno y el paso global. La carta
KYP aparece sólo después y queda fuera de esa cadena.

Para los cortes de Galerkin \(m\leq n\), el acarreo conservado es
\begin{equation}\label{eq:amp-ns-owner-acarreo-galerkin}
 C_{m\leftarrow n}
 :=P_mB(u_n,u_n)-B_m(P_m u_n,P_m u_n).
\end{equation}
La lectura nonádica de este estado está tipada por
\begin{equation}\label{eq:amp-ns-owner-e9j9}
 E_9J_9=I,\qquad P_9=J_9E_9,\qquad Q_{\rm micro}=I-P_9,
\end{equation}
y conserva el balance causal
\begin{equation}\label{eq:amp-ns-owner-balance-causal}
 \boxed{
 \mathbb E_9\|\Xi_{m+1}\|^2+\mathbb E_9\ell_m
 =\|\Xi_m\|^2,\qquad \ell_m\geq0.}
\end{equation}
El campo, la proyección solenoidal, el acarreo, la frontera, la ruta y las
pérdidas positivas permanecen así en un mismo estado antes de construir el
almacenamiento.

\section*{Gram truncado del estado completo e inyección raíz uniforme}

La realización que gobierna el balance no comienza por la desigualdad KYP.
Sean \(\Gamma_{M,j}^{\rm term}\) la transición de un paso del estado
PC--Fock enriquecido, \(\Gamma_{M,j:k}^{\rm term}\) su composición entre
las profundidades \(j\) y \(k\), y
\(L_{M,k,T}^{\rm term}\) la columna de pérdidas de grado uno, grado dos,
acarreo, innovación, archivo y memoria en el horizonte \([0,T]\). La
notación \(\mathbb E_{j:k}\) designa la suma ortogonal con los pesos de
las ramas del refinamiento entre esos niveles. Para \(j\leq J\) se define
la columna finita
\begin{equation}\label{eq:amp-ns-owner-columna-gram-truncado}
 \begin{aligned}
 \mathbf G_{M,j;J,T}^{\rm term}x
 :={}&
 \bigl(\mathcal E_{M,J,T}^{\rm term}
       \Gamma_{M,j:J}^{\rm term}x\bigr)_{\mathbb E_{j:J}}\\
 &\oplus
 \bigoplus_{k=j}^{J-1}
 \bigl(L_{M,k,T}^{\rm term}
       \Gamma_{M,j:k}^{\rm term}x\bigr)_{\mathbb E_{j:k}},
 \qquad
 U_{M,j;J,T}^{\rm term}
 :=(\mathbf G_{M,j;J,T}^{\rm term})^*
    \mathbf G_{M,j;J,T}^{\rm term}.
 \end{aligned}
\end{equation}
En particular, el almacenamiento es un Gram truncado construido con el
terminal y las pérdidas futuras del mismo estado. Su forma cuadrática es
\begin{equation}\label{eq:amp-ns-owner-gram-truncado}
 \begin{aligned}
 \langle x,U_{M,j;J,T}^{\rm term}x\rangle
 ={}&\mathbb E_{j:J}
 \|\mathcal E_{M,J,T}^{\rm term}
       \Gamma_{M,j:J}^{\rm term}x\|^2\\
 &+\sum_{k=j}^{J-1}\mathbb E_{j:k}
 \|L_{M,k,T}^{\rm term}
       \Gamma_{M,j:k}^{\rm term}x\|^2 .
 \end{aligned}
\end{equation}

\begin{lemma}[Identidad de Stein del Gram truncado]
\label{lem:amp-ns-owner-stein-truncado}
Para todo \(j<J\),
\begin{equation}\label{eq:amp-ns-owner-stein-truncado}
 U_{M,j;J,T}^{\rm term}
 =(\Gamma_{M,j}^{\rm term})^*
   U_{M,j+1;J,T}^{\rm term}\Gamma_{M,j}^{\rm term}
  +(L_{M,j,T}^{\rm term})^*L_{M,j,T}^{\rm term}.
\end{equation}
Por iteración se obtiene, conservando el sumando terminal,
\begin{equation}\label{eq:amp-ns-owner-telescopio-truncado}
 \begin{aligned}
 U_{M,0;J,T}^{\rm term}
 ={}&(\Gamma_{M,0:J}^{\rm term})^*
      U_{M,J;J,T}^{\rm term}\Gamma_{M,0:J}^{\rm term}\\
 &+\sum_{j=0}^{J-1}
   (\Gamma_{M,0:j}^{\rm term})^*
   (L_{M,j,T}^{\rm term})^*L_{M,j,T}^{\rm term}
   \Gamma_{M,0:j}^{\rm term}.
 \end{aligned}
\end{equation}
\end{lemma}

\begin{proof}
En \eqref{eq:amp-ns-owner-columna-gram-truncado}, se separa el sumando
\(k=j\). Los restantes sumandos, incluido el terminal, son exactamente la
columna de nivel \(j+1\) precedida por
\(\Gamma_{M,j}^{\rm term}\). La ortogonalidad de la suma y la propiedad de
torre de \(\mathbb E_{j:k}\) dan
\eqref{eq:amp-ns-owner-stein-truncado}. La repetición del mismo
desdoblamiento da \eqref{eq:amp-ns-owner-telescopio-truncado}; no se
descarta
\((\Gamma_{M,0:J}^{\rm term})^*U_{M,J;J,T}^{\rm term}
\Gamma_{M,0:J}^{\rm term}\).
\end{proof}

\begin{theorem}[Inyección raíz y presupuesto uniforme]
\label{thm:amp-ns-inyeccion-raiz-uniforme}
La raíz del realizador se transporta a cada corte y profundidad mediante
una inclusión isométrica:
\begin{equation}\label{eq:amp-ns-inyeccion-raiz-comun}
 \begin{aligned}
 x_{M,0}^{\rm data}
 &:=\mathcal V_M^{\rm full}(P_Mu_0)\oplus g_M^F,\\
 \mathbf G_{M,0;J,T}^{\rm term}x_{M,0}^{\rm data}
 &=\mathcal J_{M,J,T}^{\rm term}(u_0,f),\\
 \mathcal J_{M,J,T}^{\rm term}
 &:=I_{0\to M,J}\mathcal J_T^{\rm term},
 & I_{0\to M,J}^*I_{0\to M,J}&=I.
 \end{aligned}
\end{equation}
Para los datos de \eqref{eq:realizaciones-ns-dominio}, existe
\(C_{{\rm HMT},T}\), independiente de \(M\) y \(J\), tal que
\begin{equation}\label{eq:amp-ns-inyeccion-raiz-cota}
 \sup_{M,J}
 \|\mathcal J_{M,J,T}^{\rm term}(u_0,f)\|^2
 \leq C_{{\rm HMT},T}
 +(1+2\widetilde M_{s,\beta,\nu}^{\,2})
  e^{R_{\rm F}^{\,2}\|u_0\|_{H^s}^2}
 +C_{\nu,s}\int_0^T\|f(t)\|_{H^{s-1}}^2\,dt .
\end{equation}
Equivalentemente, la cota es una cota del Gram terminal evaluado en
los datos:
\begin{equation}\label{eq:amp-ns-owner-gram-cota-uniforme}
 \sup_{M,J}
 \left\langle x_{M,0}^{\rm data},
 U_{M,0;J,T}^{\rm term}x_{M,0}^{\rm data}\right\rangle
 =
 \sup_{M,J}
 \|\mathcal J_{M,J,T}^{\rm term}(u_0,f)\|^2
 <\infty .
\end{equation}
La aplicación \(\mathcal J_T^{\rm term}\) depende sólo del estado inicial,
de sus dos hojas, del terminal heredado y del puerto de fuerza; no recibe
como argumento \(u_M(t)\) para \(0<t\leq T\).
\end{theorem}

\begin{proof}
La componente lineal y las memorias \(D/H\) se transportan
isométricamente. En la componente PC--Fock, la columna coherente es
\[
 \bigoplus_{n\geq0}
 \frac{R_{\rm F}^{\,n}}{\sqrt{n!}}\,u_0^{\odot n},
 \qquad
 \sum_{n\geq0}\frac{R_{\rm F}^{\,2n}}{n!}
       \|u_0\|_{H^s}^{2n}
 =e^{R_{\rm F}^{\,2}\|u_0\|_{H^s}^2}.
\]
Para justificar la uniformidad del lector cuadrático, se polariza
\begin{equation}\label{eq:amp-ns-lector-cuadratico-polarizado}
 \mathcal B_M(u,v):=
 \frac{1}{4\sqrt{\beta\nu}}\,
 A_M^{-1/2}\Lambda^sP_M\mathbb P
 \bigl((P_Mu\!\cdot\!\nabla)P_Mv
      +(P_Mv\!\cdot\!\nabla)P_Mu\bigr).
\end{equation}
En las bases transversales
\(e^s_{p,a}=\langle p\rangle^{-s}a\,e^{ip\cdot x}\), si
\(k=p+q\ne0\), sus filas de Fourier cumplen
\begin{equation}\label{eq:amp-ns-filas-fourier-uniformes}
 \|c^M_{k;p,q}\|
 \leq\frac{C}{\sqrt{\beta\nu}}\,
 \frac{\langle k\rangle^s(|p|+|q|)}
 {|k|\langle p\rangle^s\langle q\rangle^s},
 \qquad
 \sup_{M,k\ne0}\sum_{p+q=k}\|c^M_{k;p,q}\|^2<\infty .
\end{equation}
La segunda cota se obtiene separando \(|p|\leq2|k|\) de
\(|p|>2|k|\); en la cola domina
\(\sum_p\langle p\rangle^{2-2s}<\infty\) porque \(s>5/2\).
Cauchy--Schwarz en cada fila y la ortogonalidad de los bloques con distinto
\(k\) prolongan \(\mathcal B_M\) a
\[
 \mathfrak C_M:H^s_\sigma\odot_2H^s_\sigma\longrightarrow L^2_\sigma,
 \qquad
 \sup_M\|\mathfrak C_M\|
 \leq\widetilde M_{s,\beta,\nu}.
\]
Así, la constante usada en la raíz controla la suma de todos los pares
Fourier y no sólo cada coeficiente por separado. El puerto afín verifica
\[
 Q_f(T)\leq C_{\nu,s}
 \int_0^T\|f(t)\|_{H^{s-1}}^2\,dt .
\]
Por suma ortogonal de esas componentes,
\begin{equation}\label{eq:amp-ns-raiz-cota-por-componentes}
 \|\Xi_{M,0}\|^2
 \leq C_{{\rm HMT},T}
 +(1+2\widetilde M_{s,\beta,\nu}^{\,2})
 e^{R_{\rm F}^{\,2}\|u_0\|_{H^s}^2}.
\end{equation}
Las inclusiones de \eqref{eq:amp-ns-inyeccion-raiz-comun} no cambian esta
norma. Al sumar la cota del puerto se obtiene exactamente
\eqref{eq:amp-ns-inyeccion-raiz-cota}. Como todos los términos se evalúan
en \((u_0,f)\), la estimación precede a la trayectoria y no presupone la
cota global que después se deduce de
\eqref{eq:amp-ns-owner-telescopio-truncado}.
\end{proof}

El Gram terminal \eqref{eq:amp-ns-owner-gram-truncado}, su identidad de
Stein, la telescopía con terminal y la inyección raíz uniforme forman la
construcción propietaria. Ésta es la rama que gobierna el cierre y coincide
con la realización íntegra restaurada en
\cref{sec:ns-owner-closure,thm:nsclose-proprietary-full}.

\section*{Separación del control KYP finito no gobernante}

La factorización KYP finita correcta del
\cref{thm:ns-import-affine-stein} se conserva como control alternativo de
una carta reducida y no constituye una premisa del owner
\(U^{\rm own}\). No se utiliza la mutación firmada que pretendía
incorporar \(C_{m,M}^*C_{m,M}\) a un término \(P\succeq0\) y
recuperarlo después con signo negativo: al entrar en \(P\), ese
sumando aparece necesariamente con signo positivo. Por tanto, esa
factorización no interviene en el balance, el presupuesto raíz, el retorno
físico ni la conclusión Navier--Stokes.

\section*{Derivada temporal, presión y pegado de horizontes}

El tramo siguiente es consecuencia de la cota uniforme de trayectoria
obtenida por el Gram terminal propietario, su telescopía y el retorno
físico exacto. Se
dispone de la cota conjunta uniforme en
\(L^\infty(0,T;H^s)\cap L^2(0,T;H^{s+1})\).
La ecuación proyectada se escribe
\begin{equation}\label{eq:amp-ns-proyectada}
 \partial_tu_m
 =\nu\Delta u_m-P_mB(u_m,u_m)+P_mf.
\end{equation}
Para \(s>5/2\), el producto de Sobolev y la cota conjunta en
\(L^\infty(0,T;H^s)\cap L^2(0,T;H^{s+1})\) dan
\begin{equation}\label{eq:amp-ns-dt}
 \sup_m\|\partial_tu_m\|_{L^2(0,T;H^{s-1})}
 \leq C_T(u_0,f,\Xi_0).
\end{equation}
La estimación \eqref{eq:amp-ns-dt} y la compacidad espacial producen
convergencia fuerte en \(L^2(0,T;H^s)\), suficiente para pasar al término
bilineal sin perder el carry de Galerkin.

Recuperada la velocidad, la presión normalizada de media cero queda fijada
por
\begin{equation}\label{eq:amp-ns-presion}
 -\Delta p
 =\partial_i\partial_j(u_i u_j)-\nabla\!\cdot f,
 \qquad \int_{\mathbb T^3}p(t,x)\,dx=0.
\end{equation}
La presión es así una salida de la misma solución solenoidal; no es una
variable libre usada para seleccionar la rama regular.

\begin{proposition}[Pegado global por unicidad]
\label{prop:amp-ns-pegado}
Sean \(u^{(T)}\) las soluciones obtenidas en \([0,T]\) por el límite de
Galerkin. Si \(T_1<T_2\), entonces
\(u^{(T_2)}|_{[0,T_1]}=u^{(T_1)}\). La familia define una única solución
global suave para \(t>0\).
\end{proposition}

\begin{proof}
Ambas restricciones poseen los mismos datos y pertenecen al dominio de la
estimación uniforme. La ecuación de su diferencia, el control bilineal en
\(H^{s-1}\) y Grönwall implican igualdad en \([0,T_1]\). Los intervalos
finitos forman un sistema dirigido; la igualdad en los solapamientos define
la solución global. La regularización parabólica del teorema se aplica en
cada intervalo \([\varepsilon,T]\), \(\varepsilon>0\).
\end{proof}

El control de refutación es preciso: una pérdida de uniformidad en
\eqref{eq:amp-ns-dt}, dos límites diferentes en un mismo horizonte o una
presión que no satisfaga \eqref{eq:amp-ns-presion} romperían la cadena antes
de la conclusión global.

\paragraph{Conclusión propietaria.}
El Gram terminal propietario, Stein, la telescopía con terminal, su ledger,
el retorno físico exacto y la inyección raíz producen una cota \(H^s\) uniforme para
\(s>5/2\). Por tanto se demuestra la existencia, unicidad y regularidad
global de Navier--Stokes tridimensional periódico en el dominio HMT
registrado. Las cartas KYP finitas permanecen como
\texttt{CONTROL\_NO\_GOBERNANTE}.


\paragraph{Conclusión afirmativa.}
En el dominio periódico y Sobolev declarado, el almacenamiento terminal
\(U^{\rm own}\), su identidad de Stein, la telescopía que conserva el
terminal, el ledger que domina \(E_s(u_M)\), la raíz uniforme y el retorno
físico exacto producen la cota uniforme. Aubin--Lions, unicidad y
regularidad parabólica dan la solución global suave. Las cartas KYP,
NS--OBS y LRDN--LCN son controles alternativos y no premisas del teorema.

\chapter[Proceso gauge euclídeo Yang--Mills en dimensión 4]{Construcción
proyectiva de un proceso gauge euclídeo Yang--Mills, reconstrucción de
Osterwalder--Schrader y brecha colectiva en dimensión \(4\)}
\chaptermark{Proceso gauge euclídeo Yang--Mills}

\noindent La especialización del
\cref{thm:nucleo-continuidad-conexion-nonadica-conjunta} es
\[
 (\widehat X_\infty,\Gamma_9)
 \xrightarrow{\ \mathcal R_{\rm YM}\ }
 (\mu_m,T_m,H_m,\Omega_m)
 \longrightarrow(\mu_\infty,\mathrm{OS},E(4))
 \longrightarrow\operatorname{gap}(H_{\rm phys})>0.
\]
El lector clover que publica \(F^2\) reconoce la curvatura sobre este mismo
proceso proyectivo; no construye retrospectivamente \(\mu_m\). La teoría
gauge, la reconstrucción OS, el vacío simple y la brecha ya quedan fijados
por esta cadena. La identidad
\eqref{eq:amp-ym-identidad-accion-requerida} compara después una
parametrización opcional por \(g_R\); es un control no gobernante y no una
condición para identificar la teoría Yang--Mills construida.

\section*{Objeto gauge y dominio de refinamiento}

Sea \(G\) un grupo de Lie compacto, conectado y simple. Para cada
\(m\geq0\), se considera la red euclídea
\begin{equation}\label{eq:realizaciones-ym-red}
 \Lambda_m=a_m\mathbb Z^4,
 \qquad a_{m+1}=\frac{a_m}{9}.
\end{equation}
La álgebra cilíndrica gauge se genera por variables de enlace en \(G\),
operadores de entrelazamiento de representación y las coordenadas
enriquecidas de color,
orientación, holonomía, acarreo de fusión, ruta y memoria. Sea \(\mu_m\) la
medida invariante común y sea
\[
 \mathcal H_m=L^2(\mathcal A_m,\mu_m)^{\rm gauge}.
\]
El vector constante \(\Omega_m\) define
\begin{equation}\label{eq:realizaciones-ym-proyectores}
 P_{\Omega_m}=|\Omega_m\rangle\langle\Omega_m|,
 \qquad Q_m^{\rm phys}=I-P_{\Omega_m}.
\end{equation}
El segundo proyector es el complemento físico del vacío; no se identifica
con un residual empleado únicamente para organizar la escala de
renormalización.

El circuito de actualización es local:
\begin{equation}\label{eq:realizaciones-ym-circuito}
 K_m^{\rm loc}
 =\prod_{r=1}^{9\cdot81}
   \prod_{B\in\mathcal B_{m,r}}K_{m,B}.
\end{equation}
La partición \(\{\mathcal B_{m,r}\}_r\) enumera cada coordenada una sola vez; los factores
del circuito se identificarán abajo con el semigrupo de sus expectativas
condicionales, y no con nueve dinámicas independientes.
Los bloques de un mismo color poseen soportes disjuntos y el diámetro físico
de cada soporte está uniformemente acotado al variar \(m\). Esta localidad
forma parte del objeto, de modo que el refinamiento no convierte una
actualización elemental en una operación instantáneamente global.

Sea \(T_m\) la transferencia inducida por
\eqref{eq:realizaciones-ym-circuito}. Las inclusiones isométricas
\(J_m:\mathcal H_m\to\mathcal H_{m+1}\) conservan ancestro,
representación, orientación, acarreo y memoria. Para las cuatro reflexiones
euclídeas \(\Theta_{\nu,m}\), \(\nu=1,\ldots,4\), se cumple
\begin{equation}\label{eq:realizaciones-ym-entrelazamiento}
 (T_{m+1})^9J_m=J_mT_m,
 \qquad
 \Theta_{\nu,m+1}J_m=J_m\Theta_{\nu,m}
 \quad(\nu=1,\ldots,4).
\end{equation}
Las cuatro direcciones pertenecen, por tanto, al mismo refinamiento, la misma
medida y la misma dinámica.

\section*{Semigrupo común, corona relativa y cuatro formas de reflexión}

En la carta triangular del estado completo, sean \(E_{m,c}\) las esperanzas
condicionales ortogonales y conmutantes asociadas a coordenadas
independientes. El generador positivo de solapamiento se construye primero en
un nivel semilla \(m_0\):
\begin{equation}\label{eq:realizaciones-ym-generador}
 \widehat H_{m_0}^{\rm ov}
 =3\kappa_{\rm av}
  \sum_{c\in\mathcal I_{m_0}^{\rm seed}}(I-E_{m_0,c}),
\end{equation}
y se prolonga mediante
\begin{equation}\label{eq:realizaciones-ym-generador-refinado}
 \begin{aligned}
 \widehat H_{m+1}^{\rm ov}
 &=\widehat\Gamma_m^*
   \bigl(\widehat H_m^{\rm ov}\otimes I+
         I\otimes\widehat H_{m+1}^{\rm det}\bigr)
   \widehat\Gamma_m,\\
 \widehat H_{m+1}^{\rm det}
 &=3\kappa_{\rm av}
   \sum_{\iota\in\mathcal I_{m+1/m}^{\rm det}}
   (I-E_{m+1,\iota}^{\rm det}),
 \end{aligned}
\end{equation}
donde
\begin{equation}\label{eq:realizaciones-ym-kappa}
 \kappa_{\rm av}
 =\frac{\mu_{\rm vol}}{\ell_0}
  \left(\frac{\pi_{\rm HMT}}{\varphi_{\rm HMT}^{\,2}}-1\right)>0.
\end{equation}
La frecuencia volumétrica \(\mu_{\rm vol}\) es adimensional; la dimensión
procede de \(\ell_0\).

El semigrupo, la transferencia de un paso físico y su compatibilidad de
escala se definen con tipos visibles:
\begin{equation}\label{eq:realizaciones-ym-semigrupo}
 \widehat K_{m,t}^{\rm ov}=e^{-t\widehat H_m^{\rm ov}},
 \qquad
 \widehat T_m=\widehat K_{m,a_m}^{\rm ov},
 \qquad a_{m+1}=a_m/9,
\end{equation}
Como la partición enumera cada coordenada una sola vez y las expectativas
condicionales de sus bloques conmutan, para cada bloque y para el circuito
completo se obtiene
\begin{equation}\label{eq:realizaciones-ym-circuito-semigrupo}
 K_{m,B}=e^{-3\kappa_{\rm av}a_m(I-E_{m,B})},
 \qquad
 K_m^{\rm loc}=e^{-a_m\widehat H_m^{\rm ov}}=\widehat T_m.
\end{equation}
\begin{equation}\label{eq:realizaciones-ym-semigrupo-entrelazado}
 \widehat H_{m+1}^{\rm ov}\widehat J_m
 =\widehat J_m\widehat H_m^{\rm ov},
 \qquad
 (\widehat T_{m+1})^9\widehat J_m
 =\widehat J_m\widehat T_m.
\end{equation}
Así, la novena potencia se compara mediante \(\widehat J_m\); no se igualan
operadores que actúan en Hilbert distintos.

La medida común de ocho caras procede de ese mismo kernel reversible:
\begin{equation}\label{eq:realizaciones-ym-medida-ocho}
 \widehat\Omega_m^{(8)}(dy_1\cdots dy_8)
 =\int\widehat\mu_m^{\rm ov}(dz)
   \prod_{r=1}^{8}
   \widehat K_{m,a_m/2}^{\rm ov}(z,dy_r).
\end{equation}
Para cada reflexión \(\nu=1,\ldots,4\) y todo observable cilíndrico \(F\)
en el semiespacio positivo,
\begin{equation}\label{eq:realizaciones-ym-os}
 \int\overline{F(\Theta_{\nu,m}y)}F(y)\,
 d\widehat\Omega_m^{(8)}
 =\|\widehat{\mathcal B}_{m,\nu}F\|^2\geq0,
\end{equation}
y la marginal de caras opuestas da la identidad de enlace
\begin{equation}\label{eq:realizaciones-ym-os-transferencia}
 \widehat T_{m,\nu}^{\rm OS}
 =\widehat K_{m,a_m}^{\rm ov}
 =e^{-a_m\widehat H_m^{\rm ov}}
 =\widehat T_m,
 \qquad \nu=1,\ldots,4.
\end{equation}
Las cuatro reflexiones publican, por tanto, el mismo par
transferencia--medida.

\section*{Persistencia celular del terminal y brecha exacta}

Sean
\[
 K_9=C_*(Q_{9,3};\mathbb Z),\qquad
 K_{10}=C_*(Q_{10,3};\mathbb Z).
\]
La corona relativa tiene dimensiones
\[
 C_*(Q_{10,3},Q_{9,3})=(271,756,702,217),
 \qquad 271+702=756+217=973.
\]
La inclusión \(j:K_9\to K_{10}\), el colapso celular
\(r:K_{10}\to K_9\) y el homotopio integral \(H_{\rm cell}\) satisfacen
\begin{equation}\label{eq:realizaciones-ym-sdr-corona}
 rj=I,\qquad
 \partial H_{\rm cell}+H_{\rm cell}\partial=I-jr,\qquad
 rH_{\rm cell}=0,\quad H_{\rm cell}j=0,\quad H_{\rm cell}^2=0.
\end{equation}
Para todo mapa de cadenas
\(F:K_{10}\otimes V_{q=6}\to\mathcal T\), con
\(P=(jr)\otimes I_{q=6}\),
\begin{equation}\label{eq:realizaciones-ym-homotopia-terminal}
 F-FP
 =d_{\mathcal T}\bigl(F(H_{\rm cell}\otimes I)\bigr)
  +F(H_{\rm cell}\otimes I)(\partial\otimes I).
\end{equation}
La corona completa es así nulhomotópica y el terminal \(FP\) permanece
literalmente. Los \(271\) vértices no sustituyen las aristas, caras y cubos
de la corona; el factor celular tridimensional tampoco se identifica con el
factor gauge físico \(q=6\).

Para cada coordenada persistente \(c\), en la carta local
\(q=(q_1,\ldots,q_6)\in\mathbb F_3^6\), se fija el testigo centrado
\begin{equation}\label{eq:realizaciones-ym-testigo-definicion}
 \bigcap_c\operatorname{Ran}E_{m,c}=\mathbb C\Omega_m,
 \qquad
 W_c(q)=\mathbf 1_{\{q_1+\cdots+q_6=0\}}-\frac13,
\end{equation}
donde la suma del indicador se toma en \(\mathbb F_3\). La descomposición de
Hoeffding de las expectativas conmutantes prueba
\begin{equation}\label{eq:realizaciones-ym-gap-finito}
 \ker\widehat H_m^{\rm ov}=\mathbb C\Omega_m,
 \qquad
 \widehat H_m^{\rm ov}\succeq
 3\kappa_{\rm av}(I-P_{\Omega_m}).
\end{equation}
La igualdad no se deduce sólo de que el núcleo sea trivial: cada sector no
constante de la descomposición recibe al menos un proyector
\(I-E_{m,c}\) y, por \eqref{eq:realizaciones-ym-generador}, su coste mínimo
es \(3\kappa_{\rm av}\). El testigo material \(W_c\) satisface
\begin{equation}\label{eq:realizaciones-ym-testigo-gap}
 \widehat H_m^{\rm ov}W_c=3\kappa_{\rm av}W_c,
 \qquad \|W_c\|^2=\frac29,
\end{equation}
de modo que el complemento es no vacío y la brecha es exactamente
\(3\kappa_{\rm av}\).

La escala se publica por
\begin{equation}\label{eq:realizaciones-ym-qm}
 q_m=e^{-3\kappa_{\rm av}a_m},
 \qquad q_{m+1}^{\,9}=q_m,
 \qquad -a_m^{-1}\log q_m=3\kappa_{\rm av}.
\end{equation}
Aunque \(q_m\to1\) al refinar, el cociente espectral por unidad física
permanece constante. La contracción celular
\eqref{eq:realizaciones-ym-homotopia-terminal} conserva el terminal; la
brecha procede del generador positivo de solapamiento y de su descomposición de
Hoeffding, no de la aciclicidad por sí sola.

Las inclusiones de \eqref{eq:realizaciones-ym-semigrupo-entrelazado}
forman el límite inductivo
\begin{equation}\label{eq:realizaciones-ym-colimite}
 \mathcal H_\infty=\varinjlim(\mathcal H_m,\widehat J_m).
\end{equation}
Sobre la unión inductiva algebraica se define la forma
\begin{equation}\label{eq:realizaciones-ym-forma-limite}
 \mathfrak h_\infty[\widehat J_{m,\infty}\psi]
 :=\langle\psi,\widehat H_m^{\rm ov}\psi\rangle_{\mathcal H_m}.
\end{equation}
El entrelazamiento de los generadores hace que esta definición sea
independiente del representante. La forma es positiva y cerrable; su
clausura determina el generador autoadjunto \(\widehat H_\infty^{\rm ov}\).
La compatibilidad de las formas transporta
\begin{equation}\label{eq:realizaciones-ym-gap-forma-limite}
 \overline{\mathfrak h}_\infty[\psi]
 \geq3\kappa_{\rm av}
 \|(I-P_{\Omega_\infty})\psi\|^2,
\end{equation}
y el vector persistente
\(\widehat J_{m,\infty}W_c\) alcanza la igualdad. En consecuencia, la
brecha exacta del generador límite es \(3\kappa_{\rm av}\).

\section*{Hipótesis exactas y encadenamiento del teorema}

La publicación gauge utiliza:
\begin{enumerate}
 \item el grupo compacto, conectado y simple y la red cuatro-dimensional
       \eqref{eq:realizaciones-ym-red};
 \item el circuito local \eqref{eq:realizaciones-ym-circuito}, con diámetro
       físico uniforme, y una medida invariante común;
 \item las inclusiones isométricas y las cuatro reflexiones entrelazadas por
       \eqref{eq:realizaciones-ym-entrelazamiento} y
       \eqref{eq:realizaciones-ym-semigrupo-entrelazado};
 \item las cuatro factorizaciones
       \eqref{eq:realizaciones-ym-os} y el enlace exacto
       \eqref{eq:realizaciones-ym-os-transferencia} sobre el mismo par
       transferencia--medida;
 \item el generador de expectativas conmutantes
       \eqref{eq:realizaciones-ym-generador}--
       \eqref{eq:realizaciones-ym-generador-refinado} y su descomposición de
       Hoeffding;
 \item el SDR integral de la corona y la conservación del terminal
       \eqref{eq:realizaciones-ym-sdr-corona}--
       \eqref{eq:realizaciones-ym-homotopia-terminal};
 \item el terminal \(P_{\Omega_m}\), el complemento físico
       \(Q_m^{\rm phys}\) y la ley de escala
       \eqref{eq:realizaciones-ym-qm};
 \item el límite inductivo y la clausura compatible de formas
       \eqref{eq:realizaciones-ym-colimite}--
       \eqref{eq:realizaciones-ym-gap-forma-limite};
 \item las cartas internas de acción y velocidad
       \(\hbar_{\rm int}\), \(c_{\rm int}\), usadas sólo para publicar la
       dimensión de masa.
\end{enumerate}

\begin{theorem}[Teoría gauge local con vacío simple y brecha colectiva]
\label{thm:realizaciones-ym}
Bajo las premisas anteriores, el límite de la familia gauge es local y
positivo por reflexión en las cuatro direcciones euclídeas. Su generador
posee vacío simple y satisface, en el sector colectivo gauge-invariante,
\begin{equation}\label{eq:realizaciones-ym-gap-limite}
 \Delta_{\rm YM}^{\rm HMT}=3\kappa_{\rm av}>0,
 \qquad
 m_{\rm gap}^{\rm HMT}
 =\frac{\hbar_{\rm int}}{c_{\rm int}}
  \Delta_{\rm YM}^{\rm HMT}>0.
\end{equation}
La brecha pertenece al sector colectivo; no asigna una masa elemental al
gluón.
\end{theorem}

\begin{proof}
La localidad uniforme del circuito conserva la red de álgebras locales. Las
igualdades \eqref{eq:realizaciones-ym-semigrupo-entrelazado} identifican la
novena potencia después de aplicar la inclusión correcta. La construcción
\eqref{eq:realizaciones-ym-medida-ocho} usa ese mismo semigrupo en las ocho
caras y \eqref{eq:realizaciones-ym-os-transferencia} identifica sus cuatro
marginales OS con \(\widehat T_m\). Por tanto, las cuatro formas positivas
\eqref{eq:realizaciones-ym-os} son compatibles en el límite y conservan una
sola medida y una sola dinámica.

Las expectativas de \eqref{eq:realizaciones-ym-generador} son ortogonales y
conmutantes. Su descomposición de Hoeffding separa el modo constante de cada
sector no constante: el primero es \(\mathbb C\Omega_m\), mientras todo
sector del complemento recibe al menos un factor \(I-E_{m,c}\) con peso
\(3\kappa_{\rm av}\). Esto prueba
\eqref{eq:realizaciones-ym-gap-finito}. El testigo
\eqref{eq:realizaciones-ym-testigo-gap} alcanza la cota, de modo que no es
sólo un piso: es la primera energía no nula.

La homotopía \eqref{eq:realizaciones-ym-homotopia-terminal} contrae el
complemento celular y deja intacto \(FP\); por ello la corona no introduce un
segundo vacío ni borra el acarreo. La compatibilidad isométrica transporta el
vacío y su complemento al nivel siguiente sin mezclarlos con el residual de
escala. La ley
\eqref{eq:realizaciones-ym-qm} muestra que el cociente espectral por unidad
física es exactamente \(3\kappa_{\rm av}\) en todo nivel. Las formas
\eqref{eq:realizaciones-ym-forma-limite}--
\eqref{eq:realizaciones-ym-gap-forma-limite} transportan la cota y el testigo
al generador autoadjunto del límite, por lo que la brecha sigue siendo exacta.
La reconstrucción por positividad de
Osterwalder--Schrader identifica el Hamiltoniano con el generador de esta
misma transferencia, conserva el vacío simple y produce
\eqref{eq:realizaciones-ym-gap-limite}. La carta dimensional final transforma
la brecha en la masa colectiva indicada.
\end{proof}

\section*{Compresión gauge y testigo físico de la brecha}

Sea
\[
 \widehat{\mathscr H}_m
 :=L^2(\widehat X_m,\widehat\mu_m^{\rm ov})
\]
el espacio de la realización material completa. Además de la acción gauge
usual sobre enlaces y marcos, cada collar \(c\) porta la acción finita de
incidencia
\begin{equation}\label{eq:amp-ym-accion-gauge-incidencia}
 q_c\longmapsto q_c+Bx_c,
 \qquad x_c\in\mathbb F_3^4,
\end{equation}
donde las seis filas de \(B\) están indexadas por
\((01,02,03,12,13,23)\). El producto de ambas acciones preserva
\(\widehat\mu_m^{\rm ov}\). Con \(\mathcal C_m\) el conjunto finito de
collares presentes en el nivel, se escribe
\[
 \mathcal G_m^{\rm full}
 :=G^{V(\Lambda_m)}
   \times\prod_{c\in\mathcal C_m}\mathbb F_3^4.
\]
Su promedio de Haar define la proyección
ortogonal
\begin{equation}\label{eq:amp-ym-proyeccion-fisica-haar}
 \mathbb P_m^{\rm phys}F
 :=\int_{\mathcal G_m^{\rm full}}F(g^{-1}\!\cdot\,)\,dg,
 \qquad
 \widehat{\mathscr H}_m^{\rm phys}
 :=\operatorname{Ran}\mathbb P_m^{\rm phys}.
\end{equation}
Ésta es una compresión sobre la subálgebra gauge del estado completo; la
marginal que olvida las coordenadas de incidencia es otro mapa y no se
confunde con \(\mathbb P_m^{\rm phys}\).

\begin{lemma}[Reducción del Hamiltoniano a la subálgebra física]
\label{lem:amp-ym-compresion-fisica}
La proyección \(\mathbb P_m^{\rm phys}\) reduce el generador y su
semigrupo. Las inclusiones de refinamiento reducen simultáneamente las
subálgebras físicas:
\begin{equation}\label{eq:amp-ym-compresion-entrelazada}
 \begin{gathered}
 [\mathbb P_m^{\rm phys},\widehat H_m^{\rm ov}]=0,
 \qquad
 [\mathbb P_m^{\rm phys},e^{-t\widehat H_m^{\rm ov}}]=0,\\
 \mathbb P_{m+1}^{\rm phys}\widehat J_m
 =\widehat J_m\mathbb P_m^{\rm phys},\qquad
 H_m^{\rm phys}
 :=\widehat H_m^{\rm ov}\big|_{\widehat{\mathscr H}_m^{\rm phys}},\\
 H_{m+1}^{\rm phys}\widehat J_m
 =\widehat J_mH_m^{\rm phys}.
 \end{gathered}
\end{equation}
\end{lemma}

\begin{proof}
El generador físico sobre enlaces es gauge equivariante por la
biinvariancia de Haar. En la carta producto, \(E_{q_c}\) es la esperanza
uniforme en \(\mathbb F_3^6\); por invariancia de la medida uniforme bajo
las traslaciones de \eqref{eq:amp-ym-accion-gauge-incidencia}, conmuta con
esa acción. Las expectativas de las subcolas de marcación actúan en
factores disjuntos. Por tanto cada sumando de
\[
 \widehat H_m^{\rm ov}
 =H_m^{\rm ov}\otimes I
 +3\kappa_{\rm av}\sum_c
 \bigl[(I-E_{q_c})+(I-E_{u_c^{\rm P}})\bigr]
\]
conmuta con el promedio \eqref{eq:amp-ym-proyeccion-fisica-haar}. El
cálculo funcional da la conmutación con el semigrupo. Finalmente,
\(\widehat J_m\) conserva los factores ancestrales e inserta la constante
en los factores nuevos; promediar antes o después de esa inclusión produce
la misma función. Esto prueba las cuatro identidades de
\eqref{eq:amp-ym-compresion-entrelazada}.
\end{proof}

\begin{lemma}[Testigo gauge no nulo y entrelazado]
\label{lem:amp-ym-testigo-fisico-q6}
Para todo collar \(c\), la función
\begin{equation}\label{eq:amp-ym-testigo-fisico-def}
 W_c(q)
 :=\mathbf 1_{\{q_1+\cdots+q_6=0\}}-\frac13
\end{equation}
pertenece a \(\widehat{\mathscr H}_m^{\rm phys}\), no es nula y verifica
\begin{equation}\label{eq:amp-ym-testigo-fisico-momentos}
 \mathbb E W_c=0,\qquad
 \|W_c\|^2=\frac29,\qquad
 \mathbb E(W_c^3)=\frac2{27}.
\end{equation}
Además,
\begin{equation}\label{eq:amp-ym-testigo-fisico-gap}
 H_m^{\rm phys}W_c=3\kappa_{\rm av}W_c,\qquad
 H_{m+1}^{\rm phys}\widehat J_mW_c
 =3\kappa_{\rm av}\widehat J_mW_c.
\end{equation}
En consecuencia, el valor \(3\kappa_{\rm av}\) pertenece al espectro de la
restricción física y no sólo al de una dilatación material.
\end{lemma}

\begin{proof}
La acción gauge sobre enlaces actúa en el primer factor y deja fijo
\(W_c\). Para la acción de incidencia, cada \(x_i\) aparece en tres de las
seis coordenadas, luego
\[
 \sum_{i<j}(Bx)_{ij}=3\sum_i x_i=0
 \quad\text{en }\mathbb F_3.
\]
La condición de \eqref{eq:amp-ym-testigo-fisico-def} es, por tanto,
invariante; también lo es bajo las permutaciones de las seis incidencias.
Así \(\mathbb P_m^{\rm phys}W_c=W_c\).

Bajo la medida uniforme, la suma de los seis trits es uniforme en
\(\mathbb F_3\). La función toma el valor \(2/3\) con probabilidad \(1/3\)
y \(-1/3\) con probabilidad \(2/3\), lo que da
\eqref{eq:amp-ym-testigo-fisico-momentos}. El primer sumando de
\(\widehat H_m^{\rm ov}\) anula \(W_c\); para \(c'\ne c\),
\((I-E_{q_{c'}})W_c=0\), mientras
\(E_{q_c}W_c=0\). Todas las expectativas de marcación anulan la
diferencia correspondiente porque \(W_c\) es constante en esas subcolas.
Queda exactamente
\(\widehat H_m^{\rm ov}W_c=3\kappa_{\rm av}W_c\). La primera identidad de
\eqref{eq:amp-ym-testigo-fisico-gap} sigue de la reducción; la segunda,
del entrelazamiento de \eqref{eq:amp-ym-compresion-entrelazada}.
\end{proof}

\section*{Semigrupo propietario, vacío y cota espectral exacta}

Sobre la subálgebra física se escribe
\[
 D_m^{\rm YM}:=H_m^{\rm phys},\qquad
 K_{m,t}:=\exp(-tD_m^{\rm YM}),\qquad
 P_{\Omega_m}:=|\Omega_m\rangle\langle\Omega_m|.
\]
La descomposición de Hoeffding del mismo generador y el cálculo funcional
dan, sin cambiar de medida ni de proceso,
\begin{equation}\label{eq:amp-ym-owner-semigrupo-vacio}
 K_{m,t}P_{\Omega_m}=P_{\Omega_m},
 \qquad
 \bigl\|K_{m,t}(I-P_{\Omega_m})\bigr\|
 \leq e^{-3\kappa_{\rm av}t}.
\end{equation}
El testigo de \eqref{eq:amp-ym-testigo-fisico-def} es no nulo y alcanza la
cota:
\begin{equation}\label{eq:amp-ym-owner-testigo-wc}
 D_m^{\rm YM}W_c=3\kappa_{\rm av}W_c,
 \qquad \|W_c\|^2=\frac29.
\end{equation}
Por tanto el vacío es simple y la brecha no sólo está acotada inferiormente,
sino que es exactamente
\begin{equation}\label{eq:amp-ym-owner-brecha-exacta}
 \Delta_{\rm YM}^{\rm HMT}=3\kappa_{\rm av}>0.
\end{equation}
El semigrupo, el proyector de vacío, la contracción del complemento y
\(W_c\) pertenecen a la misma familia proyectiva y constituyen la cadena
propietaria del cierre.

\section*{Reconstrucción euclídea, campos de Schwinger y lector de curvatura}

La medida proyectiva común del teorema determina inclusiones isométricas
\(\widehat J_m\) entre los espacios de nivel. Para cada una de las cuatro
reflexiones coordenadas \(\theta_\mu\), \(\mu=1,\ldots,4\), existe una
factorización
\begin{equation}\label{eq:amp-ym-cuatro-os}
 \int\overline{F(\theta_\mu y)}F(y)\,
 d\widehat\Omega_m(y)
 =\|\mathcal B_{m,\mu}F\|_2^2\geq0.
\end{equation}
Las cuatro formas proceden de la misma medida y del mismo semigrupo. Los
conectores OS
\[
 \mathcal J_{m,\mu}^{\rm OS}
 =\mathcal V_{m+1,\mu}^{-1}\widehat J_m\mathcal V_{m,\mu}
\]
entrelazan sus Hamiltonianos. En el colímite,
\begin{equation}\label{eq:amp-ym-hamiltoniano-comun}
 \mathcal V_{\infty,\mu}H_{{\rm OS},\infty}^{(\mu)}
 \mathcal V_{\infty,\mu}^{-1}
 =\widehat H_\infty,
 \qquad \mu=1,\ldots,4.
\end{equation}

La forma límite se diagonaliza por la descomposición de Hoeffding:
\begin{equation}\label{eq:amp-ym-forma-limite}
 \mathcal E_\infty(u)=3\kappa_{\rm av}
 \sum_{S\Subset\mathcal I_\infty}|S|\,\|u_S\|^2.
\end{equation}
Las sumas parciales proporcionan una sucesión de recuperación y la
semicontinuidad produce el límite inferior. La convergencia de Mosco de las
formas cerradas conserva el núcleo unidimensional y el umbral espectral. El
vector material \(W_c\) satisface
\begin{equation}\label{eq:amp-ym-testigo-gap}
 \widehat H_mW_c=3\kappa_{\rm av}W_c,
 \qquad \|W_c\|^2=\frac29,
 \qquad \mathbb E(W_c^3)=\frac2{27},
\end{equation}
de modo que la cota inferior se alcanza y la brecha es exactamente
\(3\kappa_{\rm av}\).

\begin{theorem}[Terminación euclídea y publicación de curvatura]
\label{thm:amp-ym-terminacion}
La familia proyectiva del teorema \ref{thm:realizaciones-ym} determina una
acción fuertemente continua de \(E(4)\), una red local gauge no escalar,
funcionales de Schwinger compatibles y una publicación de campo en
\(H^{-s}_{\rm loc}(\mathbb R^4)\) para \(s>2\). El producto superficial
ordenado de los enlaces define lectores clover y \(F^2\); en el dominio
suave, el lector de acción converge a
\begin{equation}\label{eq:amp-ym-f2}
 \mathcal S_{\rm YM}(A)
 =\frac1{4g_R^2}\int_{\mathbb R^4}\|F_A(x)\|^2\,d^4x.
\end{equation}
Estas publicaciones conservan el vacío simple y la brecha exacta del mismo
Hamiltoniano común.
\end{theorem}

\begin{proof}
Los Gram cofinales de los lectores suavizados y la conmutatividad de los
cuadrados de rotación y refinamiento preservan los ideales nulos; su
colímite construye la representación fuerte de \(E(4)\) y la red local.
Las estimaciones martingala de los incrementos de enlace dan tightness en
\(H^{-s}_{\rm loc}\), \(s>2\), y las características mixtas identifican los
funcionales de Schwinger con el mismo proceso cilíndrico. El producto
orientado alrededor de cada plaqueta construye el clover. Su expansión en
el régimen suave y la convergencia fuerte del lector marginal dan
\eqref{eq:amp-ym-f2}. Ninguna de estas operaciones modifica la medida, el
semigrupo ni la forma límite usados para obtener
\eqref{eq:amp-ym-testigo-gap}; por ello el vacío y la brecha se conservan.
\end{proof}

\section*{Identificación exacta del límite gauge}

La identificación de la teoría no descansa en que cuatro reflexiones
generen por sí solas el grupo euclídeo.  Las reflexiones proporcionan la
positividad OS; la acción continua procede, separadamente, del grupoide de
pantallas y de la compatibilidad de refinamiento.  Para
\(g=(g_v)_{v\in\Lambda_m}\in G^{\Lambda_m}\), la acción sobre una variable
de enlace es
\begin{equation}\label{eq:amp-ym-accion-gauge-enlace}
 (g\cdot U)_e=g_{s(e)}U_e g_{t(e)}^{-1}.
\end{equation}
La invariancia de Haar, la desintegración triangular del kernel y la
cancelación de las aristas interiores prueban
\begin{equation}\label{eq:amp-ym-invariancia-gauge-medida}
 (g\cdot)_*\widehat\mu_m^{\rm ov}=\widehat\mu_m^{\rm ov},
 \qquad
 K_{m,t}^{\rm ov}(g\cdot U,g\cdot dV)
 =K_{m,t}^{\rm ov}(U,dV).
\end{equation}
Por diferenciación sobre el núcleo cilíndrico, todo generador vertical
\(X\) satisface la identidad de Ward
\begin{equation}\label{eq:amp-ym-ward}
 \int \nabla_XF\,d\widehat\mu_m^{\rm ov}=0.
\end{equation}

La curvatura se obtiene del mismo sistema de enlaces.  Si una plaqueta
gruesa \(p\) se subdivide en sus 81 subplaquetas \(p'\), entonces
\begin{equation}\label{eq:amp-ym-producto-superficial}
 \operatorname{Hol}^{(m)}_{p}
 =\prod_{p'\subset p}^{\longrightarrow}
 T_{p'}\operatorname{Hol}^{(m+1)}_{p'}T_{p'}^{-1}.
\end{equation}
Las contribuciones de toda arista interior aparecen con orientaciones
opuestas y se cancelan exactamente.  La expresión es gauge-covariante y
conmuta con las inclusiones \(\widehat J_m\).  En una carta suave se define
\begin{equation}\label{eq:amp-ym-clover-definido}
 F_{{\rm cl},m}(x)
 =\frac1{4a_m^2}\sum_{p\ni x}
 \operatorname{Ad}_{T_p}\log\operatorname{Hol}_{p},
 \qquad
 F_{{\rm cl},m}=F_A+O(a_m^2).
\end{equation}
De aquí, para \(A\in C^4_{\rm loc}\) y
\(a_m^2/g_m^2\to0\), se obtiene por suma de Riemann
\begin{equation}\label{eq:amp-ym-f2-convergencia-calculada}
 \frac{a_m^4}{4g_m^2}
 \sum_{x\in\Lambda_m}\|F_{{\rm cl},m}(x)\|^2
 \longrightarrow
 \frac1{4g_R^2}\int_{\mathbb R^4}\|F_A(x)\|^2\,d^4x.
\end{equation}

\begin{theorem}[Proceso gauge euclídeo y lector de curvatura de
Yang--Mills]
\label{thm:amp-ym-identificacion-gauge}
El límite construido en el teorema
\ref{thm:realizaciones-ym} es un proceso gauge euclídeo en dimensión
\(4\) para el grupo compacto, conectado y simple fijado. Posee
covariancia gauge local, una red local no escalar, positividad por
reflexión, una representación fuertemente continua de \(E(4)\),
funcionales de Schwinger compatibles, curvatura gauge-covariante con límite
clásico \eqref{eq:amp-ym-f2-convergencia-calculada}, vacío simple y brecha
espectral
\(3\kappa_{\rm av}>0\).
\end{theorem}

\begin{proof}
Las identidades \eqref{eq:amp-ym-invariancia-gauge-medida}--
\eqref{eq:amp-ym-ward} construyen la simetría gauge y sus identidades de
Ward en todos los niveles; su naturalidad las transporta al límite.  La
localidad uniforme de los kernels produce la red isótona y local.  Las
cuatro formas \eqref{eq:amp-ym-cuatro-os}, junto con la acción del grupoide
de pantallas sobre el núcleo Weyl, producen respectivamente la positividad
OS y la acción fuerte de \(E(4)\).  La tightness en
\(H^{-s}_{\rm loc}\), \(s>2\), identifica una única ley límite y sus
funcionales de Schwinger.  Finalmente,
\eqref{eq:amp-ym-producto-superficial}--
\eqref{eq:amp-ym-f2-convergencia-calculada} identifican el lector de campo
y su límite clásico, mientras
\eqref{eq:amp-ym-forma-limite}--
\eqref{eq:amp-ym-testigo-gap} prueban vacío simple y brecha exacta para la
misma ley.

La medida queda determinada constructivamente por sus marginales
proyectivas, su kernel reversible y sus funcionales cilíndricos. El lector
clover no repondera esa medida: publica una función del mismo proceso y su
límite clásico.
\end{proof}

\begin{proposition}[Control no gobernante de comparación Gibbs--acción]
\label{prop:amp-ym-puente-accion-requerido}
La ley proyectiva, su vacío simple y su brecha exacta ya están construidos
por \eqref{eq:amp-ym-owner-semigrupo-vacio}--
\eqref{eq:amp-ym-owner-brecha-exacta}. Como comparación exterior opcional de
tipo Gibbs--acción, una parametrización por \(g_R\) puede estudiar una familia
\[
 g_R\longmapsto
 \bigl(\widehat\mu_{m,g_R}^{\rm ov},H_{m,g_R}^{\rm phys}\bigr)
\]
y una relación entre la medida o su forma de Dirichlet y el funcional de
acción, por ejemplo
\begin{equation}\label{eq:amp-ym-identidad-accion-requerida}
 \frac{d\widehat\mu_{m,g_R}^{\rm ov}}{d\nu_m}(U)
 =Z_{m,g_R}^{-1}e^{-S_{m,g_R}^{F^2}(U)},
 \qquad
 \mathfrak h_{m,g_R}[F]
 =\langle F,H_{m,g_R}^{\rm phys}F\rangle,
\end{equation}
con compatibilidad de refinamiento y paso al límite. Aquí \(\nu_m\) puede
ser la medida producto de Haar de la red finita. Esta comprobación se rotula
\texttt{CONTROL\_NO\_GOBERNANTE}: no es premisa de la teoría gauge HMT ni
de su brecha.
\end{proposition}

\begin{proof}
En \eqref{eq:amp-ym-f2-convergencia-calculada}, variar \(g_R\) reescala el
lector de curvatura, mientras las fórmulas que construyen
\(\widehat\mu_m^{\rm ov}\), \(\widehat H_m^{\rm ov}\) y la brecha
\(3\kappa_{\rm av}\) no contienen \(g_R\). Son, por ello, construcciones
independientes. La compresión
\eqref{eq:amp-ym-compresion-entrelazada} y el testigo
\eqref{eq:amp-ym-testigo-fisico-gap} prueban la persistencia física de la
brecha del proceso ya construido; el control sólo compara después una
dependencia dinámica adicional en \(g_R\).
\end{proof}

La secuencia demostrada en esta sección es
\[
 \text{medida común}\to\text{cuatro OS}\to\text{Hamiltoniano común}
 \to E(4),\ \text{Schwinger},\ H^{-s}_{\rm loc}
 \to\text{clover y }F^2.
\]
Esta secuencia demuestra una teoría Yang--Mills cuatridimensional con vacío
simple y brecha positiva exacta \(3\kappa_{\rm av}\) en el dominio HMT
registrado. La identidad
\eqref{eq:amp-ym-identidad-accion-requerida} queda fuera de la cadena
rectora como \texttt{CONTROL\_NO\_GOBERNANTE}.


\section*{Lectura de Wilson}

El observable de Wilson se obtiene después del generador. Para una carta
arquimediana de enlace,
\begin{equation}\label{eq:realizaciones-ym-wilson}
 A_e=\sum_a\operatorname{ev}_{\mathbb R}(x_{e,a})T_a,
 \qquad
 U_e=e^{a_mA_e},
 \qquad
 W_\gamma=\operatorname{Tr}\!\prod_{e\in\gamma}U_e.
\end{equation}
La lectura conserva la representación, el operador de entrelazamiento y el
orden de la holonomía. Sirve como observable terminal de la teoría publicada; no define
retrospectivamente la medida, el semigrupo ni la brecha.

\section*{Necesidad de las hipótesis y obstrucciones exactas}

\begin{ablation}[Necesidad de la dinámica y la medida comunes]
\label{abl:realizaciones-ym}
El teorema \ref{thm:realizaciones-ym} no se conserva si se sustituyen las
nueve subfases por actualizaciones independientes, se emplean medidas o
núcleos diferentes en las cuatro reflexiones, se pierde la localidad física
uniforme, \(\widehat J_m\) deja de entrelazar transferencia y reflexión, se
borran acarreo o memoria, se elimina \(P_{\Omega_m}\), se identifica
\(Q_m^{\rm phys}\) con un residual de escala, se exige
\(q_m\leq q_*<1\) en unidades de retícula, se reemplaza la corona completa
por sus \(271\) vértices, se borra \(FP\), o se usa
\eqref{eq:realizaciones-ym-wilson} para seleccionar el generador.
\end{ablation}

\begin{proof}
Actualizaciones, núcleos o medidas distintos destruyen la factorización
simultánea \eqref{eq:realizaciones-ym-os}. Sin localidad o entrelazamiento,
la positividad de un corte no define una red local compatible. Borrar acarreo,
memoria o \(FP\) elimina información que
\eqref{eq:realizaciones-ym-homotopia-terminal} conserva. Usar sólo los
vértices cambia el complejo relativo y puede fabricar modos cero. Confundir
los dos complementos permite que el sector físico desaparezca por una
operación puramente escalar. La cota fija en unidades de retícula contradice
\eqref{eq:realizaciones-ym-qm}. Finalmente, Wilson sólo es una lectura del
estado construido y no demuestra por sí mismo ni
\eqref{eq:realizaciones-ym-os} ni
\eqref{eq:realizaciones-ym-gap-finito}.
\end{proof}

Los residuos independientes son: el fallo de
\(\widehat T_{m,\nu}^{\rm OS}=e^{-a_m\widehat H_m^{\rm ov}}\), el fallo del
entrelazamiento \eqref{eq:realizaciones-ym-semigrupo-entrelazado}, una forma
OS negativa, un vector de núcleo ortogonal a \(\Omega_m\), o un sector de
Hoeffding no constante con energía inferior a \(3\kappa_{\rm av}\).

\section*{Alcance de la construcción gauge}

La prueba combina la conexión conjunta, la conservación del terminal, la
doble lectura, el circuito local, la medida común, las cuatro factorizaciones
de reflexión, el refinamiento entrelazado y el generador positivo. En el
dominio gauge regular declarado, la
localidad, la teoría Yang--Mills cuatridimensional, el vacío simple y la
brecha colectiva son resultados exactos. La reconstrucción de
Osterwalder--Schrader y el observable de Wilson son
publicaciones posteriores del mismo estado; ninguna de ellas actúa como
selector retrospectivo de la medida, del semigrupo, del vacío ni de la
brecha. La comparación opcional con una familia parametrizada por \(g_R\) se estudia
mediante \eqref{eq:amp-ym-identidad-accion-requerida}; esa comprobación no
gobierna ni reabre la medida, el semigrupo, el vacío o la brecha ya
demostrados.

\section*{Comparación de los mecanismos espectral, solenoidal y gauge}

Las construcciones anteriores comparten la arquitectura
\begin{equation}\label{eq:realizaciones-sintesis-tres}
 \begin{array}{ccl}
 \text{Riemann--Weil}
&:& \text{memoria finita}\to\text{doble publicación cofinal}\\
&&{}\to B_W=\mathscr I_{\rm GW}\geq0
      \to\text{criterio reversible},\\[1mm]
 \text{Navier--Stokes}
&:& U^{\rm own}\to\text{telescopía y ledger}\\
&&{}\to\text{raíz uniforme y retorno exacto}
      \to\text{límite parabólico},\\[1mm]
 \text{proceso gauge OS}
&:& \text{raíz terminal}\to\text{positividad por reflexión}\\
&&{}
      \to\text{brecha colectiva}.
 \end{array}
\end{equation}
La conexión nonádica conjunta no sustituye los lemas especializados: los
organiza y conserva sus datos al refinar. Las conclusiones de
Riemann--Weil, Navier--Stokes y Yang--Mills proceden de sus respectivas
cadenas propietarias y criterios finales ya demostrados. Los comparadores
de residual, de Gram y de densidad de acción quedan como controles
posteriores no gobernantes. Esta separación impide tanto la pérdida de
genealogía como la transferencia injustificada de una prueba de un dominio
a otro.

% Residencia sucesora única de Hodge y Birch--Swinnerton-Dyer. Este input es
% parte del capítulo, no una importación duplicada de los manuscritos
% fuentes.
% 04 — FRAGMENTO SUCESOR HODGE Y BIRCH--SWINNERTON-DYER
% =======================================================
% Inserción autónoma prevista tras el marcador de integración del Libro IV.
% No importa fuentes probatorios ni modifica las tres realizaciones
% precedentes.
%
% Trazabilidad técnica no narrativa:
%   Hodge: CHI-II-HDG-01/02/03, CHI-IV-HDG-01/C01.
%   BSD:   CHI-II-BSD-01/02/03, CHI-IV-BSD-01/02/C01.
%   Fuentes: celda APP--Mayer--Vietoris, totalizador coend y controles
%   elípticos; producto fibrado genealógico de Manin, publicación acoplada
%   Kato--PT, dualidad ortogonal reducida, relleno finita--singular y regulador
%   Bockstein derivado. Los hashes y estados de ejecución permanecen en el
%   expediente técnico de la edición sucesora.

\ifdefined\HMTSucesoraStructureTrace
  \typeout{HMT-SUCESORA 04: realizaciones Hodge y Birch--Swinnerton-Dyer}
\fi

\chapter{Completitud de la historia cohomológica enriquecida y algebraicidad de las clases racionales de Hodge}
\chaptermark{Completitud cohomológica y algebraicidad de Hodge}

\subsection*{Objeto, dominio y datos de partida}

Sea \(X\) una variedad proyectiva lisa sobre \(\mathbb C\) y sea
\(p\geq0\). Se escribe
\begin{equation}\label{eq:hodge-suc-dominio}
 \operatorname{Hdg}^{p}(X)_{\mathbb Q}
 :=H^{2p}(X,\mathbb Q)
   \cap H^{p,p}(X,\mathbb C).
\end{equation}
Sea \(\mathbf{Perf}^{\geq p}(X)\) la categoría idempotente completa de
complejos perfectos cuyo soporte tiene codimensión al menos \(p\). En esta
sección se emplea la abreviatura
\[
 \operatorname{Perf}^{\geq p}(X)_{\mathbb Q}
 :=K_0\!\left(\mathbf{Perf}^{\geq p}(X)\right)
   \otimes_{\mathbb Z}\mathbb Q;
\]
por tanto, \(R_{\mathrm{Hdg}}\) publica la clase del complejo perfecto
construido. El punto de partida no es una clase desnuda de
\eqref{eq:hodge-suc-dominio}, sino el límite de historias supervivientes
\begin{equation}\label{eq:hodge-suc-xchi}
 X_{\chi}(X,p)
 =\varprojlim_{K}\operatorname{Surv}_{K}(X,p).
\end{equation}
Cada elemento de \(X_{\chi}(X,p)\) conserva hoja, orientación, cociente,
acarreo, frontera, ruta y memoria. En particular, el retorno de la conexión
nonádica conjunta satisface
\begin{equation}\label{eq:hodge-suc-retorno}
 \Gamma _9=\operatorname{Hol}_{C_9}(\nabla_{\mathrm{TPK}}),
 \qquad g(k+9)=g(k),
 \qquad \widehat x(k+9)\neq\widehat x(k)
 \quad\text{en general}.
\end{equation}
La igualdad conserva la fase; la desigualdad conserva la profundidad de la
historia. Esta distinción excluye que un retorno se interprete como una
reinicialización.

La composición de rutas está gobernada por el cociclo de acarreo
\begin{equation}\label{eq:hodge-suc-carry}
 \operatorname{Carr}(\gamma _2\gamma _1)
 =\operatorname{Carr}(\gamma _2)H(\gamma _1)
  +Y(\gamma _2)\operatorname{Carr}(\gamma _1),
\end{equation}
y la memoria acumulada por
\begin{equation}\label{eq:hodge-suc-memoria}
 S_0=\sum_{r<N}(M_9^{*})^{r}D_9^{*}D_9M_9^{r}
     +(M_9^{*})^{N}S_0M_9^{N}.
\end{equation}
Las ecuaciones \eqref{eq:hodge-suc-xchi}--\eqref{eq:hodge-suc-memoria}
son datos construidos en la Parte II. Aquí se especializan; no se
reconstruyen a partir del codominio cohomológico.

\subsection*{Celda local APP--Mayer--Vietoris}

Sean \(Z,W\subset X\) incidencias cerradas con intersección
Tor-independiente y soportes orientados de la codimensión declarada. Para
enteros \(a,b\geq1\), defínase
\begin{equation}\label{eq:hodge-suc-kappa}
 \begin{aligned}
 \kappa_{a,b}:\mathcal O_Z^{a}\oplus\mathcal O_W^{b}
 &\longrightarrow \mathcal O_{Z\cap W}^{ab},\\
 ((u_i)_{i=1}^{a},(v_j)_{j=1}^{b})
 &\longmapsto
 \bigl(u_i|_{Z\cap W}-v_j|_{Z\cap W}\bigr)_{i,j}.
 \end{aligned}
\end{equation}
Sobre una fibra de \(Z\cap W\), el núcleo diagonal de
\(\kappa_{a,b}\) tiene rango uno y su conúcleo tiene rango
\begin{equation}\label{eq:hodge-suc-defecto-rango}
 ab-(a+b-1)=(a-1)(b-1).
\end{equation}
Si \((\rho_{\Sigma},c_{\Sigma})\) y
\((\rho_{\Pi},c_{\Pi})\) son, respectivamente, los residuos y acarreos de
las hojas aditiva y multiplicativa de APP, la misma cantidad se expresa por
\begin{equation}\label{eq:hodge-suc-celda-carry}
 (a-1)(b-1)
 =\rho_{\Pi}-\rho_{\Sigma}+1
  +9\bigl(c_{\Pi}-c_{\Sigma}\bigr).
\end{equation}
Así, el defecto local no es una corrección exterior: es la publicación en
\(\operatorname{Perf}\) de la diferencia entre ambas hojas con su acarreo.

Sea \(\tau\) el intercambio de los dos sumandos de la fuente y sea \(P\)
la transposición de los índices del blanco. La orientación TRIT fija
\begin{equation}\label{eq:hodge-suc-orientacion-kappa}
 \kappa_{b,a}\tau=-P\kappa_{a,b}.
\end{equation}
El signo de \eqref{eq:hodge-suc-orientacion-kappa} forma parte del objeto
orientado. Su omisión identificaría rutas opuestas y destruiría la
naturalidad del carácter localizado.

\begin{proposition}[Realización perfecta de la celda local]
\label{prop:hodge-suc-celda-perf}
El cono
\begin{equation}\label{eq:hodge-suc-cone-kappa}
 \mathcal C_{a,b}(Z,W):=\operatorname{Cone}(\kappa_{a,b})
\end{equation}
es un complejo perfecto con el soporte prescrito. Su clase localizada
conserva el defecto de rango \eqref{eq:hodge-suc-defecto-rango}, el acarreo
\eqref{eq:hodge-suc-celda-carry} y la orientación
\eqref{eq:hodge-suc-orientacion-kappa}.
\end{proposition}

\begin{proof}
La Tor-independencia identifica la restricción derivada al cruce con la
restricción ordinaria empleada en \eqref{eq:hodge-suc-kappa}. Las fuentes y
el blanco son perfectos sobre sus soportes; por estabilidad de
\(\operatorname{Perf}\) bajo conos, también lo es
\(\mathcal C_{a,b}(Z,W)\). El cálculo fibra a fibra da
\eqref{eq:hodge-suc-defecto-rango}; la división APP de los datos de suma y
producto da \eqref{eq:hodge-suc-celda-carry}. Finalmente, intercambiar
\(Z,a\) con \(W,b\) cambia cada diferencia por su opuesta, lo que prueba
\eqref{eq:hodge-suc-orientacion-kappa}. Como el carácter de Chern localizado
es aditivo en triángulos distinguidos, las tres informaciones pasan a la
clase del cono.
\end{proof}

\subsection*{Totalización de historias y conservación del terminal}

Sea \(J_N\) la categoría finita de historias compatibles a profundidad
\(N\), sea \(\Gamma\) el diagrama de complejos de incidencia obtenido de
\eqref{eq:hodge-suc-cone-kappa} y sea \(\mathsf W\) el módulo derecho que
registra las orientaciones y pesos de ruta. La totalización derivada es el
coend
\begin{equation}\label{eq:hodge-suc-coend}
 K_{X,p,N}
 :=\mathsf W\otimes^{\mathbf L}_{\mathbb Z[J_N]}\Gamma.
\end{equation}
Los enlaces de supervivencia forman un sistema compatible en \(N\). La
holonomía \(\Gamma_9\) induce un endomorfismo
\(U_{\Gamma_9}\) del totalizador y el objeto terminal se conserva como
\begin{equation}\label{eq:hodge-suc-terminal}
 \mathcal T_{X,p}
 :=\operatorname{Cone}\bigl(I-U_{\Gamma_9}\bigr).
\end{equation}
El cono \eqref{eq:hodge-suc-terminal} impide que la totalización confunda
una fase repetida con la misma historia: el valor visible puede retornar,
pero la diferencia terminal retiene el desplazamiento de memoria.

La compatibilidad entre \eqref{eq:hodge-suc-coend} y el límite
\eqref{eq:hodge-suc-xchi} define el morfismo de realización
\begin{equation}\label{eq:hodge-suc-realizacion}
 R_{\mathrm{Hdg}}:
 X_{\chi}(X,p)\longrightarrow
 \operatorname{Perf}^{\geq p}(X)_{\mathbb Q}.
\end{equation}
Este morfismo se fija en la fuente: no recibe como entrada una clase de
Hodge ni un representante que deba reconocer.

\begin{theorem}[Identidad de publicación Hodge]
\label{thm:hodge-suc-identidad}
Sobre \(X_{\chi}(X,p)\) conmuta la identidad
\begin{equation}\label{eq:hodge-suc-identidad}
 \operatorname{cl}_{X,p}\circ
 \operatorname{ch}^{\mathrm{loc}}_{p}\circ
 R_{\mathrm{Hdg}}
 =\operatorname{ev}^{\mathrm{Hdg}}_{\infty,X,p}.
\end{equation}
Aquí
\(\operatorname{ch}^{\mathrm{loc}}_{p}\) toma la componente de
codimensión \(p\) del carácter localizado y
\(\operatorname{cl}_{X,p}\) es su publicación cohomológica.
\end{theorem}

\begin{proof}
En una celda, la igualdad es la compatibilidad del carácter localizado con
la restricción, el triángulo de Mayer--Vietoris y la clase cohomológica. La
orientación está fijada por \eqref{eq:hodge-suc-orientacion-kappa}; por
tanto, no queda un signo por elegir. La aditividad del carácter de Chern
transporta la igualdad a cada cono \eqref{eq:hodge-suc-cone-kappa}. La
naturalidad respecto de las flechas de \(J_N\) la hace descender al coend
\eqref{eq:hodge-suc-coend}. La ley de acarreo
\eqref{eq:hodge-suc-carry} garantiza su compatibilidad bajo composición y
\eqref{eq:hodge-suc-terminal} conserva el término terminal. Al pasar al
límite proyectivo se obtienen dos transformaciones naturales con iguales
componentes en todo cilindro finito; luego coinciden y resulta
\eqref{eq:hodge-suc-identidad}.
\end{proof}

El teorema de completitud del atlas de historias enriquecidas, aplicado a
esta especialización, establece
\begin{equation}\label{eq:hodge-suc-completitud}
 \operatorname{ev}^{\mathrm{Hdg}}_{\infty,X,p}
 \bigl(X_{\chi}(X,p)\bigr)
 =\operatorname{Hdg}^{p}(X)_{\mathbb Q}.
\end{equation}
Ésta es la estructura de partida explícita de la especialización Hodge.

% Ampliación sustantiva de la completitud coinductiva Hodge.
% Módulo destinado al capítulo Hodge de la Parte IV.

\section{Despliegue coinductivo de la completitud Hodge}
\label{sec:amp-hodge-completitud-desplegada}

La igualdad de completitud
\eqref{eq:hodge-suc-completitud} puede desplegarse sin comenzar por una clase
algebraica ni introducir en la fuente la conclusión que se desea publicar. El
objeto que se completa es el sistema de historias enriquecidas; la clase de
Hodge interviene únicamente como dato de evaluación compatible en sus
truncaciones finitas.

\subsection{Sistemas finitos de observación y fibras compatibles}
\label{subsec:amp-hodge-sistemas-finitos}

Para cada profundidad \(N\geq0\), sea \(J_N(X,p)\) la categoría finita de
cartas de incidencia, rutas orientadas y terminales de memoria de profundidad a
lo sumo \(N\). Se escribe
\begin{equation}\label{eq:amp-hodge-surv-n}
 \mathscr S_N(X,p):=\operatorname{Surv}_{J_N}(X,p),
 \qquad
 \rho_{N+1,N}:\mathscr S_{N+1}(X,p)\longrightarrow\mathscr S_N(X,p)
\end{equation}
para el espacio racionalizado de estados supervivientes y su restricción. Los
lectores finitos toman valores en el módulo de observaciones de cilindro
\(\mathscr H_N^p(X)\):
\begin{equation}\label{eq:amp-hodge-evaluacion-finita}
 e_N^{\mathrm{Hdg}}:\mathscr S_N(X,p)\longrightarrow
 \mathscr H_N^p(X),
 \qquad
 q_{N+1,N}:\mathscr H_{N+1}^p(X)\longrightarrow\mathscr H_N^p(X).
\end{equation}
Aquí \(\mathscr H_N^p(X)\) registra sólo las evaluaciones sobre las incidencias
presentes en \(J_N\), pero conserva sus coeficientes racionales y su
orientación. Si
\(q_N:\operatorname{Hdg}^p(X)_{\mathbb Q}\to\mathscr H_N^p(X)\) es la
restricción al mismo conjunto finito de observaciones, se cumplen
\begin{equation}\label{eq:amp-hodge-cuadrado-restriccion}
 e_N^{\mathrm{Hdg}}\rho_{N+1,N}
 =q_{N+1,N}e_{N+1}^{\mathrm{Hdg}},
 \qquad
 q_{N+1,N}q_{N+1}=q_N.
\end{equation}

Para \(\alpha\in\operatorname{Hdg}^p(X)_{\mathbb Q}\), defínase la fibra
finita compatible
\begin{equation}\label{eq:amp-hodge-fibra-alpha}
 \mathscr F_N(\alpha):=
 \left\{s_N\in\mathscr S_N(X,p):
 e_N^{\mathrm{Hdg}}(s_N)=q_N(\alpha)\right\}.
\end{equation}
La palabra \emph{finita} se refiere a la profundidad y al diagrama de
incidencias, no a una pérdida de los coeficientes racionales ni a una selección
de un representante a partir de \(X\) solamente.

\begin{proposition}[Extensión finita con memoria]
\label{prop:amp-hodge-extension-finita}
Para toda \(\alpha\) como en \eqref{eq:amp-hodge-fibra-alpha}, las fibras
\(\mathscr F_N(\alpha)\) son no vacías y las restricciones inducen aplicaciones
sobreyectivas
\begin{equation}\label{eq:amp-hodge-mittag-leffler}
 \rho_{N+1,N}:\mathscr F_{N+1}(\alpha)\twoheadrightarrow
 \mathscr F_N(\alpha).
\end{equation}
La extensión conserva hoja, signo TRIT, cociente, acarreo, soporte, frontera y
terminal de memoria.
\end{proposition}

\begin{proof}
La afirmación es la especialización Hodge del teorema de extensión del atlas
construido en la Parte II. Conviene hacer visible su mecanismo. Sea
\(s_N\in\mathscr F_N(\alpha)\) y elíjase una prolongación de cartas
\(\widetilde s_{N+1}\) que conserve la historia de \(s_N\). Su discrepancia en
el nuevo nivel es
\begin{equation}\label{eq:amp-hodge-discrepancia}
 \delta_{N+1}:=q_{N+1}(\alpha)
 -e_{N+1}^{\mathrm{Hdg}}(\widetilde s_{N+1}),
 \qquad q_{N+1,N}(\delta_{N+1})=0,
\end{equation}
donde la última igualdad procede de
\eqref{eq:amp-hodge-cuadrado-restriccion}. Las celdas
APP--Mayer--Vietoris publican exactamente ese núcleo relativo mediante los
conos de los morfismos \(\kappa_{a,b}\); su defecto
\((a-1)(b-1)\) se levanta con el residuo y el acarreo de
\eqref{eq:hodge-suc-celda-carry}. Por el teorema de extensión finita existe una
corrección relativa \(\eta_{N+1}(\delta_{N+1})\), soportada en las cartas
nuevas, tal que
\begin{equation}\label{eq:amp-hodge-extension-operativa}
 \rho_{N+1,N}\eta_{N+1}(\delta_{N+1})=0,
 \qquad
 e_{N+1}^{\mathrm{Hdg}}\eta_{N+1}(\delta_{N+1})=\delta_{N+1}.
\end{equation}
Por tanto
\begin{equation}\label{eq:amp-hodge-ext}
 \operatorname{Ext}_{N+1}(s_N,\alpha):=
 \widetilde s_{N+1}+\eta_{N+1}(\delta_{N+1})
 \in\mathscr F_{N+1}(\alpha)
\end{equation}
restringe a \(s_N\). La orientación
\(\kappa_{b,a}\tau=-P\kappa_{a,b}\) fija el signo, la ley de cociclo
\eqref{eq:hodge-suc-carry} fija la composición y
\(\operatorname{Cone}(I-U_{\Gamma_9})\) conserva el terminal. El nivel inicial
se obtiene del atlas de semillas basado. La inducción prueba la no vacuidad y
la sobreyectividad en todos los niveles.
\end{proof}

\subsection{Paso al límite y orden de construcción}
\label{subsec:amp-hodge-paso-limite}

\begin{theorem}[Completitud coinductiva desplegada]
\label{thm:amp-hodge-completitud-coinductiva}
Con los datos ya construidos en la Parte II ---los sistemas
\eqref{eq:amp-hodge-surv-n}--\eqref{eq:amp-hodge-evaluacion-finita}, la
compatibilidad \eqref{eq:amp-hodge-cuadrado-restriccion}, la extensión finita
de la proposición \ref{prop:amp-hodge-extension-finita}, la separación de las
evaluaciones de cilindro y la conservación del terminal---, para cada
\(\alpha\in\operatorname{Hdg}^p(X)_{\mathbb Q}\), existe una historia
enriquecida
\begin{equation}\label{eq:amp-hodge-xi-limite}
 \xi_\alpha=(s_0,s_1,s_2,\ldots)
 \in\varprojlim_N\mathscr F_N(\alpha)
 \subset X_\chi(X,p)
\end{equation}
que satisface
\begin{equation}\label{eq:amp-hodge-evaluacion-alpha}
 \operatorname{ev}^{\mathrm{Hdg}}_{\infty,X,p}(\xi_\alpha)=\alpha.
\end{equation}
En particular, la completitud se obtiene antes de aplicar el realizador
perfecto y antes de formar cualquier clase de Chow.
\end{theorem}

\begin{proof}
Elíjase \(s_0\in\mathscr F_0(\alpha)\). Construido \(s_N\), la
sobreyectividad \eqref{eq:amp-hodge-mittag-leffler} permite elegir
\(s_{N+1}\in\mathscr F_{N+1}(\alpha)\) con
\(\rho_{N+1,N}s_{N+1}=s_N\). Así se obtiene
\eqref{eq:amp-hodge-xi-limite}. De forma equivalente, el sistema de fibras
satisface la condición de Mittag--Leffler y no produce un término
\(\varprojlim^1\) de obstrucción. Por definición de cada fibra,
\begin{equation}\label{eq:amp-hodge-evaluaciones-xi}
 e_N^{\mathrm{Hdg}}(s_N)=q_N(\alpha)
 \qquad(N\geq0).
\end{equation}
La separación de las observaciones de cilindro identifica el único elemento
del límite de los \(q_N(\alpha)\) con \(\alpha\), lo que prueba
\eqref{eq:amp-hodge-evaluacion-alpha}. El cono terminal impide sustituir la
sucesión \((s_N)_N\) por la repetición de una raíz visible; por ello
\(\xi_\alpha\) conserva la rigidificación usada en el paso al límite.
\end{proof}

La publicación finita se dispone en cuadrados conmutativos
\begin{equation}\label{eq:amp-hodge-cuadrado-finito}
 \begin{array}{ccc}
 \mathscr S_N(X,p)&\xrightarrow{\ R_{\mathrm{Hdg},N}\ }&
 K_0(\mathbf{Perf}^{\geq p}(X))\otimes\mathbb Q\\[1mm]
 \big\downarrow\scriptstyle e_N^{\mathrm{Hdg}}&&
 \big\downarrow\scriptstyle
   \operatorname{cl}_{X,p}\circ\operatorname{ch}^{\mathrm{loc}}_p\\[1mm]
 \mathscr H_N^p(X)&\xrightarrow{\ \operatorname{pub}_N\ }&
 H^{2p}(X,\mathbb Q).
 \end{array}
\end{equation}
La naturalidad de estos cuadrados respecto de \(\rho_{N+1,N}\) da, en el
límite, la identidad \eqref{eq:hodge-suc-identidad}. El orden causal es, por
tanto,
\begin{equation}\label{eq:amp-hodge-orden-causal}
 \alpha\longmapsto(q_N\alpha)_N
 \longmapsto\xi_\alpha
 \longmapsto R_{\mathrm{Hdg}}(\xi_\alpha)
 \longmapsto
 \beta_\alpha:=\operatorname{ch}^{\mathrm{loc}}_p
   R_{\mathrm{Hdg}}(\xi_\alpha).
\end{equation}
Sólo en el último paso aparece \(\beta_\alpha\). Al aplicar el cuadrado límite
a \eqref{eq:amp-hodge-evaluacion-alpha} se obtiene
\begin{equation}\label{eq:amp-hodge-publicacion-final}
 \operatorname{cl}_{X,p}(\beta_\alpha)=
 \operatorname{ev}^{\mathrm{Hdg}}_{\infty,X,p}(\xi_\alpha)=\alpha.
\end{equation}

\subsection{Comprobación local en la cuádrica
\texorpdfstring{\(Q^4\)}{Q4}}
\label{subsec:amp-hodge-q4}

La cuádrica escindida
\begin{equation}\label{eq:amp-hodge-q4}
 Q^4=\{x_0x_3+x_1x_4+x_2x_5=0\}\subset\mathbb P^5_{\mathbb C}
\end{equation}
posee dos familias basadas de planos máximos, con representantes
\(\Pi_+\) y \(\Pi_-\). Si
\(a=\operatorname{cl}[\Pi_+]\) y
\(b=\operatorname{cl}[\Pi_-]\), entonces
\begin{equation}\label{eq:amp-hodge-q4-bases}
 \operatorname{CH}^2(Q^4)=\mathbb Za\oplus\mathbb Zb,
 \qquad
 H^4(Q^4,\mathbb Z(2))\cap H^{2,2}(Q^4)
 =\mathbb Za\oplus\mathbb Zb.
\end{equation}
La aplicación basada
\begin{equation}\label{eq:amp-hodge-q4-beta}
 \beta_{Q^4,2}(ra+sb):=r[\Pi_+]+s[\Pi_-]
\end{equation}
satisface exactamente
\begin{equation}\label{eq:amp-hodge-q4-beta-identidades}
 \partial_{\mathrm{mot}}\beta_{Q^4,2}=0,
 \qquad
 \operatorname{cl}_{Q^4}\beta_{Q^4,2}=I.
\end{equation}

Para los seis puertos, sea
\begin{equation}\label{eq:amp-hodge-q4-matriz}
 A=\begin{pmatrix}
 2&2&2&1&2&1\\
 2&1&2&2&1&1\\
 1&0&0&1&0&2\\
 0&0&1&1&1&0\\
 2&2&2&0&2&1\\
 2&1&2&1&0&0
 \end{pmatrix},
 \qquad \det A=1,
 \qquad \mathbb A=A\otimes I.
\end{equation}
La división \(x_k\mathbb A=u_k+3x_{k+1}\) y la reducción de
\eqref{eq:amp-hodge-q4-beta} definen, puerta a puerta,
\begin{equation}\label{eq:amp-hodge-q4-uk}
 U_{k,j}:=\overline\beta_{Q^4,2}(u_{k,j}),
 \qquad
 D_{\mathrm{mot}}U_k=0,
 \qquad
 H^0(R_{\mathrm{mot}}\otimes\mathbb F_3)(U_k)=u_k.
\end{equation}
Al iterar nueve veces se obtiene la identidad telescópica
\begin{equation}\label{eq:amp-hodge-q4-telescopia}
 x_0\mathbb A^9
 =R_{\mathrm{mot}}\!\left(
   \sum_{k=0}^{8}3^kU_k\mathbb A^{8-k}\right)+3^9x_9.
\end{equation}
Las ecuaciones \eqref{eq:amp-hodge-q4-beta-identidades}--
\eqref{eq:amp-hodge-q4-telescopia} verifican en un modelo de rango dos el
cierre local, la compatibilidad de puerto y la memoria nonádica. Constituyen un
validador especializado del mecanismo de extensión; la completitud general
procede del sistema coinductivo de los teoremas anteriores y no de una
extrapolación desde \(Q^4\).

\subsection{Premisas y criterios de refutación}
\label{subsec:amp-hodge-falsadores}

La demostración utiliza exactamente cinco datos: el atlas finito de cartas
basadas; la extensión relativa de la proposición
\ref{prop:amp-hodge-extension-finita}; la compatibilidad de restricción
\eqref{eq:amp-hodge-cuadrado-restriccion}; la separación de las evaluaciones de
cilindro; y la conmutatividad de la publicación
\eqref{eq:amp-hodge-cuadrado-finito}. Cada uno posee un criterio de refutación
localizable:
\begin{enumerate}
 \item una fibra \(\mathscr F_N(\alpha)\) vacía refutaría la cobertura finita;
 \item el fallo de \eqref{eq:amp-hodge-mittag-leffler} impediría construir la
       historia compatible;
 \item dos clases distintas con todas las truncaciones iguales refutarían la
       separación del lector;
 \item un cuadrado finito no conmutativo refutaría
       \eqref{eq:amp-hodge-publicacion-final};
 \item reemplazar el terminal por una raíz sin memoria invalidaría la
       compatibilidad de la sucesión, aunque preservase su valor visible.
\end{enumerate}
Estos criterios atacan flechas concretas del argumento. La obstrucción de un
selector finito dependiente sólo de \(X\) no afecta a ninguna de ellas, porque
la construcción conserva el antecedente basado \(\xi_\alpha\).

\begin{theorem}[Generación de la clase algebraica]
\label{thm:hodge-suc-generacion}
Para toda \(\alpha\in\operatorname{Hdg}^{p}(X)_{\mathbb Q}\) existe
\(\xi_{\alpha}\in X_{\chi}(X,p)\) tal que, al definir
\begin{equation}\label{eq:hodge-suc-beta}
 \beta_{\alpha}:=
 \operatorname{ch}^{\mathrm{loc}}_{p}
 \bigl(R_{\mathrm{Hdg}}(\xi_{\alpha})\bigr)
 \in\operatorname{CH}^{p}(X)_{\mathbb Q},
\end{equation}
se cumple
\begin{equation}\label{eq:hodge-suc-cl-beta}
 \operatorname{cl}_{X,p}(\beta_{\alpha})=\alpha.
\end{equation}
\end{theorem}

\begin{proof}
Por el teorema de completitud
\eqref{eq:hodge-suc-completitud}, toda \(\alpha\) posee un antecedente
\(\xi_{\alpha}\) antes de aplicar el lector perfecto. Se forma entonces
\(\beta_{\alpha}\) mediante \eqref{eq:hodge-suc-beta}. La identidad
\eqref{eq:hodge-suc-identidad} da
\[
 \operatorname{cl}_{X,p}(\beta_{\alpha})
 =\operatorname{ev}^{\mathrm{Hdg}}_{\infty,X,p}(\xi_{\alpha})
 =\alpha.
\]
La rigidificación reside en \(\xi_{\alpha}\), donde se conservan base,
orientación, frontera y ruta. Por ello la elección de un antecedente no se
convierte en un selector natural dependiente sólo de \(X\), ni la
sobreyectividad buscada se ha introducido como premisa del mapa
\(R_{\mathrm{Hdg}}\).
\end{proof}

\subsection*{Hipótesis y alcance}

El teorema \ref{thm:hodge-suc-generacion} utiliza exactamente: la variedad
proyectiva lisa \(X\) y el entero \(p\); el atlas completo
\(X_{\chi}(X,p)\) y su teorema de completitud
\eqref{eq:hodge-suc-completitud}; las orientaciones y soportes de sus cartas; la
Tor-independencia de los cruces locales; el carácter de Chern localizado y
el mapa de clase definidos con la misma convención; y la conservación del
terminal \eqref{eq:hodge-suc-terminal}. Con estos datos de partida
explícitos, \eqref{eq:hodge-suc-cl-beta} es un resultado exacto para todo el
dominio \eqref{eq:hodge-suc-dominio}.

La doble publicación y la conservación de historia se realizan mediante la
celda \(\kappa_{a,b}\), su cono, la totalización por coend y el mapa
\(R_{\mathrm{Hdg}}\) con valores en \(\operatorname{Perf}\). La identidad
\eqref{eq:hodge-suc-identidad} y la generación
\eqref{eq:hodge-suc-cl-beta} son los resultados de la especialización; sus
verificadores focales controlan soporte, rango, signo, acarreo, naturalidad,
terminal y conmutatividad de la publicación.

\subsection{Reconocimiento matemático final de Hodge}

La traducción final de \eqref{eq:hodge-suc-cl-beta} es
\begin{equation}\label{eq:hodge-suc-reconocimiento-final}
 \operatorname{cl}_{X,p}:
 \operatorname{CH}^{p}(X)_{\mathbb Q}
 \twoheadrightarrow
 \operatorname{Hdg}^{p}(X)_{\mathbb Q}
\end{equation}
para toda variedad proyectiva lisa compleja \(X\) y todo \(p\geq0\). Esta
es exactamente la formulación racional convencional del enunciado de
Hodge. Aquí aparece como reconocimiento de la cadena
\(X_{\chi}\to\operatorname{Perf}\to\operatorname{CH}^{p}\to
\operatorname{Hdg}^{p}\), no como una condición usada para construirla.

% Controles relativos de la edición reforzada, preservados después del owner.
\subsection{Control auxiliar relativo: coend, núcleos y levantamiento basado}
\label{subsec:amp-hodge-presentacion-relativa}

\noindent\textbf{Estatuto: CONTROL\_NO\_GOBERNANTE.}
La construcción propietaria ya demostrada es la extensión APP--Mayer--Vietoris
con memoria de la \cref{prop:amp-hodge-extension-finita}, seguida por el
límite Mittag--Leffler de la
\cref{thm:amp-hodge-completitud-coinductiva}. Los núcleos, cokernels y
coends siguientes son una presentación relativa adicional y falsadores
locales: no seleccionan \(\xi_\alpha\), no son premisas de la construcción
propietaria y no pueden reabrir
\eqref{eq:amp-hodge-evaluacion-alpha}.

Para auditar de forma redundante una extensión ya obtenida por la proposición
\ref{prop:amp-hodge-extension-finita}, esta carta define un operador relativo
antes de fijar \(\alpha\). Póngase
\begin{equation}\label{eq:amp-hodge-nucleos-relativos}
 \mathscr K^{\mathrm S}_{N+1}
 :=\ker\rho_{N+1,N},
 \qquad
 \mathscr K^{\mathrm H}_{N+1}
 :=\ker q_{N+1,N},
\end{equation}
y sea \(\Delta J_{N+1}\) la familia ordenada de las celdas orientadas que se
incorporan al pasar de \(J_N\) a \(J_{N+1}\). Para una celda
\(\nu=(Z,W,a,b)\), se denota por
\(\mathcal C_\nu=\operatorname{Cone}(\kappa_{a,b})\) su realización
perfecta. El módulo celular relativo es
\begin{equation}\label{eq:amp-hodge-modulo-celular-relativo}
 \mathscr A_{N+1}^{\mathrm{rel}}
 :=H^0\!\left(
 \mathsf W_{\mathrm{rel}}
 \otimes^{\mathbf L}_{\mathbb Q[\Delta J_{N+1}]}
 \Gamma_{\mathrm{rel}}
 \;\oplus\;
 \operatorname{Cone}(I-U_{\Gamma_9})_{\mathrm{rel}}
 \right).
\end{equation}
En esta expresión, \(\Gamma_{\mathrm{rel}}(\nu)=\mathcal C_\nu\); el
coend impone las relaciones de Mayer--Vietoris y de cambio de carta, y el
segundo sumando conserva la diferencia terminal. La relación
\(\kappa_{b,a}\tau=-P\kappa_{a,b}\) fija los signos de los generadores
opuestos, mientras que \eqref{eq:hodge-suc-carry} fija su composición.

Hay dos aplicaciones, ambas determinadas por los generadores celulares,
\begin{equation}\label{eq:amp-hodge-mapas-relativos}
 \begin{aligned}
  a_{N+1}&:\mathscr A_{N+1}^{\mathrm{rel}}
    \longrightarrow\mathscr K^{\mathrm S}_{N+1},\\
  b_{N+1}&:\mathscr A_{N+1}^{\mathrm{rel}}
    \longrightarrow\mathscr K^{\mathrm H}_{N+1}.
 \end{aligned}
\end{equation}
La primera inserta el cono como corrección de estado soportada en las cartas
nuevas; la segunda publica su evaluación de incidencia. Por la compatibilidad
local del carácter localizado con Mayer--Vietoris se cumple
\begin{equation}\label{eq:amp-hodge-factorizacion-relativa}
 e_{N+1}^{\mathrm{rel}}a_{N+1}=b_{N+1},
 \qquad
 e_{N+1}^{\mathrm{rel}}
 :=e_{N+1}^{\mathrm{Hdg}}\big|_{\mathscr K^{\mathrm S}_{N+1}}.
\end{equation}

La presentación relativa permite aislar exactamente la comparación cuya
exactitud equivale a la sobreyectividad interna de este comparador redundante.
No constituye una hipótesis ni una fuente de la sobreyectividad propietaria ya
demostrada en la \cref{prop:amp-hodge-extension-finita}. Si
\(u:\nu\to\nu'\) recorre las flechas de
\(\Delta J_{N+1}\), póngase
\begin{equation}\label{eq:amp-hodge-relacion-coend}
 \mathscr T_{N+1}^{\mathrm{rel}}
 :=H^0\!\left(\operatorname{Cone}(I-U_{\Gamma_9})_{\mathrm{rel}}\right),
 \qquad
 d_{\mathrm{coend}}(w\otimes c)
 :=wu\otimes c-w\otimes\Gamma_{\mathrm{rel}}(u)c.
\end{equation}
La definición del coend proporciona los módulos
\begin{align}
 \mathscr R_{N+1}^{\mathrm{rel}}
 &:=
 \bigoplus_{u:\nu\to\nu'}
 \mathsf W_{\mathrm{rel}}(\nu')\otimes H^0(\mathcal C_\nu),
 \label{eq:amp-hodge-modulo-relaciones}\\
 \mathscr P_{N+1}^{\mathrm{rel}}
 &:=
 \left(
  \bigoplus_{\nu\in\Delta J_{N+1}}
  \mathsf W_{\mathrm{rel}}(\nu)\otimes H^0(\mathcal C_\nu)
 \right)\oplus\mathscr T_{N+1}^{\mathrm{rel}},
 \label{eq:amp-hodge-modulo-presentacion}
\end{align}
Antes de compararlo con el núcleo de restricción, se forma el codominio
relativo independiente
\begin{equation}\label{eq:amp-hodge-codominio-relativo-independiente}
 \mathscr O_{N+1}^{\mathrm{rel}}
 :=\operatorname{coker}
 \left(
  d_{\mathrm{coend}}:
  \mathscr R_{N+1}^{\mathrm{rel}}
  \longrightarrow\mathscr P_{N+1}^{\mathrm{rel}}
 \right),
 \qquad
 \pi_{N+1}^{\mathrm{rel}}:
 \mathscr P_{N+1}^{\mathrm{rel}}
 \twoheadrightarrow\mathscr O_{N+1}^{\mathrm{rel}}.
\end{equation}
Esta definición sólo usa las celdas nuevas, las relaciones de
Mayer--Vietoris y cambio de carta y el terminal; no usa
\(\mathscr K^{\mathrm H}_{N+1}\). Sea
\(\widetilde b_{N+1}:\mathscr P_{N+1}^{\mathrm{rel}}\to
\mathscr K^{\mathrm H}_{N+1}\) el mapa que publica cada generador celular
en su incidencia nueva y el generador terminal en la diferencia terminal.
Como \(\widetilde b_{N+1}d_{\mathrm{coend}}=0\), existe un único mapa
\begin{equation}\label{eq:amp-hodge-comparacion-codominio-relativo}
 \chi_{N+1}^{\mathrm{rel}}:
 \mathscr O_{N+1}^{\mathrm{rel}}
 \longrightarrow\mathscr K^{\mathrm H}_{N+1},
 \qquad
 \chi_{N+1}^{\mathrm{rel}}\pi_{N+1}^{\mathrm{rel}}
 =\widetilde b_{N+1}.
\end{equation}

\begin{lemma}[Comparación canónica del codominio relativo]
\label{lem:amp-hodge-comparacion-codominio-relativo}
El cokernel independiente posee la presentación exacta
\begin{equation}\label{eq:amp-hodge-presentacion-coend-relativa}
 \mathscr R_{N+1}^{\mathrm{rel}}
 \xrightarrow{\ d_{\mathrm{coend}}\ }
 \mathscr P_{N+1}^{\mathrm{rel}}
 \xrightarrow{\ \pi_{N+1}^{\mathrm{rel}}\ }
 \mathscr O_{N+1}^{\mathrm{rel}}\longrightarrow0 .
\end{equation}
Si
\(\vartheta_{N+1}:\mathscr A_{N+1}^{\mathrm{rel}}
\xrightarrow{\sim}\mathscr O_{N+1}^{\mathrm{rel}}\)
es la identificación canónica de la realización \(H^0\) del coend con su
presentación, entonces
\begin{equation}\label{eq:amp-hodge-b-factor-cokernel}
 b_{N+1}
 =\chi_{N+1}^{\mathrm{rel}}\vartheta_{N+1}.
\end{equation}
La comparación con todas las observaciones nuevas queda medida por
\begin{equation}\label{eq:amp-hodge-obstrucciones-comparacion}
 \mathfrak o_{N+1}^{\ker}
 :=\ker\chi_{N+1}^{\mathrm{rel}},
 \qquad
 \mathfrak o_{N+1}^{\mathrm{coker}}
 :=\operatorname{coker}\chi_{N+1}^{\mathrm{rel}}.
\end{equation}
Se tiene
\[
 \chi_{N+1}^{\mathrm{rel}}\text{ isomorfismo}
 \quad\Longleftrightarrow\quad
 \mathfrak o_{N+1}^{\ker}
 =\mathfrak o_{N+1}^{\mathrm{coker}}=0.
\]
\end{lemma}

\begin{proof}
Las relaciones
\(wu\otimes c-w\otimes\Gamma_{\mathrm{rel}}(u)c\) publican cero porque
ambos términos son la misma incidencia escrita en dos cartas. Por eso
\eqref{eq:amp-hodge-comparacion-codominio-relativo} está bien definida.
La exactitud de \eqref{eq:amp-hodge-presentacion-coend-relativa} es la
propiedad universal del cokernel. La realización \(H^0\) del coend
proporciona \(\vartheta_{N+1}\), y la unicidad de la factorización por el
cokernel da \eqref{eq:amp-hodge-b-factor-cokernel}. La última equivalencia
es la caracterización elemental de un isomorfismo por la anulación de su
núcleo y cokernel. Ninguno de estos pasos demuestra esa anulación.
\end{proof}

Así, el coend y el núcleo de restricción se construyen por separado. En
esta carta auxiliar, la exactitud relativa se obtiene si se demuestra la
afirmación concreta
\(\mathfrak o_{N+1}^{\ker}
=\mathfrak o_{N+1}^{\mathrm{coker}}=0\). Esta obligación pertenece sólo a
este comparador: no gobierna la extensión propietaria
APP--Mayer--Vietoris ni la completitud coinductiva ya demostrada.

\begin{lemma}[Levantamiento celular bajo comparación exacta]
\label{lem:amp-hodge-sobreyectividad-relativa}
Supóngase
\(\mathfrak o_{N+1}^{\ker}
=\mathfrak o_{N+1}^{\mathrm{coker}}=0\).
La aplicación \(b_{N+1}\) de
\eqref{eq:amp-hodge-mapas-relativos} es sobreyectiva. Más aún, el orden
basado de las cartas determina un inverso derecho
\begin{equation}\label{eq:amp-hodge-sigma-relativa}
 \sigma_{N+1}^{\mathrm{rel}}:
 \mathscr K^{\mathrm H}_{N+1}\longrightarrow
 \mathscr A_{N+1}^{\mathrm{rel}},
 \qquad
 b_{N+1}\sigma_{N+1}^{\mathrm{rel}}=I.
\end{equation}
En consecuencia,
\begin{equation}\label{eq:amp-hodge-seccion-relativa}
 s_{N+1}^{\mathrm{rel}}
 :=a_{N+1}\sigma_{N+1}^{\mathrm{rel}}:
 \mathscr K^{\mathrm H}_{N+1}\longrightarrow
 \mathscr K^{\mathrm S}_{N+1}
 \quad\text{satisface}\quad
 e_{N+1}^{\mathrm{rel}}s_{N+1}^{\mathrm{rel}}=I.
\end{equation}
\end{lemma}

\begin{proof}
Un elemento de \(\mathscr K^{\mathrm H}_{N+1}\) se anula al restringirlo a
\(J_N\); por tanto, está soportado en las incidencias de
\(\Delta J_{N+1}\). Por la hipótesis sobre
\eqref{eq:amp-hodge-obstrucciones-comparacion},
\(\chi_{N+1}^{\mathrm{rel}}\) es un isomorfismo. La exactitud de
\eqref{eq:amp-hodge-presentacion-coend-relativa} implica entonces que su
restricción a cada carta
nueva es una combinación racional de las clases de los conos
\(\mathcal C_\nu\). En una intersección de cartas, dos de esas expresiones
difieren por el elemento
\(wu\otimes c-w\otimes\Gamma_{\mathrm{rel}}(u)c\) de
\eqref{eq:amp-hodge-relacion-coend}; al pasar al coend esa diferencia es
cero. Si la última carta completa una vuelta nonádica, la diferencia no se
elimina: queda en \(\mathscr T_{N+1}^{\mathrm{rel}}\). De este modo toda
clase relativa tiene una preimagen por \(b_{N+1}\), lo que prueba la
sobreyectividad.

Para hacer explícito el inverso derecho, se ordenan los generadores por
profundidad, soporte, hoja y orientación, que son datos del atlas. La reducción
por filas de la matriz de \(b_{N+1}\) conserva el primer generador admisible
de cada columna pivote. Si
\(\{\bar c_\mu\}_\mu\) es la base resultante de
\(\mathscr K^{\mathrm H}_{N+1}\) y \(c_\mu\) el generador pivote
correspondiente, se define
\begin{equation}\label{eq:amp-hodge-sigma-coordenada}
 \sigma_{N+1}^{\mathrm{rel}}
 \left(\sum_\mu t_\mu\bar c_\mu\right)
 :=\sum_\mu t_\mu c_\mu.
\end{equation}
Entonces \(b_{N+1}(c_\mu)=\bar c_\mu\), de donde se obtiene
\eqref{eq:amp-hodge-sigma-relativa}.

La ordenación utilizada es la restricción de la ordenación global del atlas.
Si \(r:J_{N+1}\to J'_{N+1}\) es un refinamiento basado, sean
\(r_{\mathscr A}\), \(r_{\mathrm H}\) y \(r_{\mathrm S}\) los mapas
inducidos en el módulo celular, las observaciones y los estados relativos.
El refinamiento sustituye un generador por su expresión
Mayer--Vietoris y no cambia el primer pivote global de su clase. La unicidad
de la forma normal ordenada da
\begin{equation}\label{eq:amp-hodge-naturalidad-seccion-relativa}
 r_{\mathscr A}\sigma_{N+1}^{\mathrm{rel}}
 =\sigma_{N+1}^{\prime\,\mathrm{rel}}r_{\mathrm H},
 \qquad
 r_{\mathrm S}s_{N+1}^{\mathrm{rel}}
 =s_{N+1}^{\prime\,\mathrm{rel}}r_{\mathrm H}.
\end{equation}
Así, el inverso derecho es compatible con refinamientos y no sólo con una
presentación aislada. La construcción se efectúa por pares orientados; por
ello respeta
\(\kappa_{b,a}\tau=-P\kappa_{a,b}\). Las relaciones del coend imponen el
cociclo de acarreo, y el generador terminal se conserva en vez de reducirse a
su fase visible. Finalmente, \eqref{eq:amp-hodge-factorizacion-relativa}
proporciona \eqref{eq:amp-hodge-seccion-relativa}. Ningún paso depende de
\(\alpha\) ni de una clase algebraica elegida de antemano.
\end{proof}

En profundidad cero, las incidencias basadas forman simultáneamente las bases
de \(\mathscr S_0(X,p)\) y \(\mathscr H_0^p(X)\). Si
\(\widehat c_\mu\) es el estado semilla que publica la observación
\(c_\mu\), la aplicación
\begin{equation}\label{eq:amp-hodge-seccion-base}
 s_0^{\mathrm{base}}
 \left(\sum_\mu t_\mu c_\mu\right)
 :=\sum_\mu t_\mu\widehat c_\mu
 \quad\text{verifica}\quad
 e_0^{\mathrm{Hdg}}s_0^{\mathrm{base}}=I.
\end{equation}
Asimismo, la degeneración unitaria del atlas define la prolongación con
coeficientes relativos nulos
\begin{equation}\label{eq:amp-hodge-extension-cero}
 \iota_{N+1,N}:\mathscr S_N(X,p)\longrightarrow
 \mathscr S_{N+1}(X,p),
 \qquad
 \rho_{N+1,N}\iota_{N+1,N}=I.
\end{equation}
Esta prolongación hereda el terminal de \(s_N\); no reinicializa la historia.

\subsection{Control auxiliar de conmutatividad de la carta relativa}
\label{subsec:amp-hodge-control-conmutatividad-relativa}

\noindent\textbf{Estatuto: CONTROL\_NO\_GOBERNANTE.}
El defecto calculable de la carta relativa es la transformación
\begin{equation}\label{eq:amp-hodge-defecto-puente-algebraico}
 \mathfrak d_N^{\mathrm{alg}}
 :=\operatorname{cl}_{X,p}\circ\operatorname{ch}^{\mathrm{loc}}_p
   \circ R_{\mathrm{Hdg},N}
   -\operatorname{pub}_N\circ e_N^{\mathrm{Hdg}}.
\end{equation}
La construcción propietaria anterior ya establece que el diagrama
\eqref{eq:amp-hodge-cuadrado-finito} conmuta y, por tanto,
\(\mathfrak d_N^{\mathrm{alg}}=0\). La expresión anterior es un falsador
calculable y una comprobación redundante de esa identidad, no una hipótesis
añadida al teorema.

En el incremento \(N\to N+1\), sean
\[
 R_{\mathrm{Hdg},N+1}^{\mathrm{rel}}
 :=R_{\mathrm{Hdg},N+1}\big|_{\mathscr K^{\mathrm S}_{N+1}},
 \qquad
 \operatorname{pub}_{N+1}^{\mathrm{rel}}
 :=\operatorname{pub}_{N+1}\big|_{\mathscr K^{\mathrm H}_{N+1}}.
\]
Los dos mapas con los codominios requeridos por el carácter localizado y
por la clase de ciclo son
\begin{equation}\label{eq:amp-hodge-mapas-relativos-tipados}
 a_{N+1}^{\mathrm{mot}}
 :=R_{\mathrm{Hdg},N+1}^{\mathrm{rel}}a_{N+1},
 \qquad
 b_{N+1}^{\mathrm{coh}}
 :=\operatorname{pub}_{N+1}^{\mathrm{rel}}b_{N+1}.
\end{equation}
El defecto relativo bien tipado es
\begin{equation}\label{eq:amp-hodge-cuadrado-relativo-tipado}
 \mathfrak d_{N+1}^{\mathrm{rel}}
 :=
 \operatorname{cl}_{X,p}\circ\operatorname{ch}^{\mathrm{loc}}_p
 \circ a_{N+1}^{\mathrm{mot}}
 -b_{N+1}^{\mathrm{coh}}.
\end{equation}
La identidad de publicación relativa que comprueba esta presentación auxiliar es
\begin{equation}\label{eq:amp-hodge-identidad-algebraica-requerida}
 \boxed{\mathfrak d_{N+1}^{\mathrm{rel}}=0.}
\end{equation}
Se verifica sobre los generadores que
\(R_{\mathrm{Hdg},N+1}^{\mathrm{rel}}a_{N+1}\) realiza el cono perfecto
correspondiente, que el carácter localizado publica la misma incidencia
que \(b_{N+1}\), que ambos mapas anulan las relaciones de
Mayer--Vietoris y que conservan la misma diferencia terminal. Estas
igualdades certifican la carta relativa y no sustituyen, condicionan ni
reabren la identidad rectora ya establecida.

\subsection{Validadores especializados de la cadena Hodge}

\noindent\textbf{Estatuto: CONTROL\_NO\_GOBERNANTE.}
Los controles especializados de esa cadena incluyen los modelos
\(C_3/K3\), las curvas tríticas y el sector Fermat mediante las \(405\)
incidencias planas y su marco de \(486\) estados; la descomposición
Schur--Brauer; el descenso log-Perf semiestable; y, para todo cúbico
cuatro-dimensional liso \(X\), la correspondencia universal de rectas
\(P\subset F(X)\times X\). Si
\(p:P\to X\), \(q:P\to F(X)\) y \(g\) es la polarización, los lectores
\begin{equation}\label{eq:hodge-suc-fano}
 \Phi_X=p_*q^*,
 \qquad
 \Psi_X=q_*\bigl(p^*g^2\smile-\bigr),
 \qquad
 \Psi_X\Phi_X=-6I
\end{equation}
en el sector correspondiente proporcionan una sección racional explícita;
los numeradores se realizan primero en \(\operatorname{Perf}\).

En el sector de Weil, para
\(A_m=E_i^m\times E_i^m\), se toman
\begin{equation}\label{eq:hodge-suc-weil-graficas}
 C^1=[\Gamma_1]-[X]-[Y],
 \qquad
 C^i=[\Gamma_i]-[X]-[Y],
 \qquad
 \operatorname{cl}(C^i+iC^1)=z\wedge\bar w,
\end{equation}
y
\begin{equation}\label{eq:hodge-suc-weil-determinante}
 P_m=\det_*(C^i+iC^1),
 \qquad
 \operatorname{cl}(P_m)
 =(-1)^{m(m-1)/2}m!\,w_+.
\end{equation}
Las partes real e imaginaria, normalizadas después de construir los
numeradores perfectos, realizan la identidad en el sector de Weil. Para
\(m=2\), el acarreo es dos y no requiere invertir tres. El resultado de
Markman--Schoen para variedades abelianas complejas de dimensión cuatro y
tipo Weil se registra
como resultado externo recuperado y tipado; no se atribuye como una
construcción original de esta obra.

\noindent\textbf{Candado de jerarquía.}
Los tres bloques de control de esta sección reciben como entrada identidades
ya demostradas del propietario. Ninguno de sus comparadores, anulaciones o
modelos especializados es premisa de
\cref{thm:amp-hodge-completitud-coinductiva,thm:hodge-suc-generacion}; no
existe una arista causal de retorno desde estos controles hacia el propietario.


\noindent\textbf{Estatuto del control siguiente: CONTROL\_NO\_GOBERNANTE.}
La ablación recibe la conclusión Hodge ya demostrada y sólo falsifica la
supresión del antecedente basado; no aporta ninguna premisa al propietario ni
puede demover sus teoremas.

\begin{ablation}[Selector finito sin historia]
\label{abl:hodge-suc-eliptico}
Sea \(E\) una curva elíptica compleja. No existe un conjunto finito de
representantes racionales de la clase de un punto que sea estable bajo todas
las traslaciones de \(E\) y que, a la vez, elija naturalmente un
representante sin dato basado. Esta obstrucción refuta la supresión de la
rigidificación; no afecta a los teoremas
\ref{thm:hodge-suc-identidad} y \ref{thm:hodge-suc-generacion}.
\end{ablation}

\begin{proof}
Si \(z\) tiene grado no nulo, las diferencias
\(t_x^{*}z-z\), al variar \(x\in E(\mathbb C)\), recorren mediante la
aplicación de Abel--Jacobi un subgrupo de \(\operatorname{Pic}^{0}(E)\) de
índice finito hasta isogenia. Su envolvente racional no cabe en el espacio
generado por una órbita finita. Por tanto, un marco finito estable bajo todas
las traslaciones no puede conservar simultáneamente el representante y su
naturalidad. En \(X_{\chi}(E,1)\), en cambio, la traslación modifica la
historia y su base aunque la publicación cohomológica permanezca fija. El
antecedente \(\xi_{\alpha}\) conserva precisamente ese dato, de modo que la
hipótesis abolida por el control no es una premisa del teorema.
\end{proof}
\chapter{Unidad determinantal acoplada, igualdad rango--orden de anulación y fórmula del término principal de Birch--Swinnerton--Dyer}
\chaptermark{Unidad determinantal y fórmula de Birch--Swinnerton--Dyer}

\noindent La especialización del
\cref{thm:nucleo-continuidad-conexion-nonadica-conjunta} es
\[
 (\widehat X_\infty,\Gamma_9)
 \xrightarrow{\ \mathcal R_{\rm BSD}\ }
 P_E^{\rm gen}
 \longrightarrow(z_K,z_{\rm PT})
 \longrightarrow R_{\rm Boc}^{\rm der}
 \longrightarrow\text{rango y término principal de }L(E,s).
\]
Las dos publicaciones parten de la misma unidad basada; la sucesión de
localización y la reducción de Smith preceden a la conclusión sobre
\(\Sha(E)\).

\subsection*{Objeto genealógico y coálgebra Alexander--Whitney}

Sea \(E/\mathbb Q\) una curva elíptica, fíjense el sistema de coeficientes
\(A_n\), la profundidad \(m\) y las orientaciones locales y globales del
complejo Selmer. Si
\(\operatorname{TPK}^{\mathrm{surv}}_{m}\) es la categoría de historias
supervivientes del Topological Phase Kernel, se define
\begin{equation}\label{eq:bsd-suc-gmn}
 G_{m,n}:=
 N_{*}\bigl(N\operatorname{TPK}^{\mathrm{surv}}_{m};A_n\bigr).
\end{equation}
Para un simplex no degenerado, la diagonal Alexander--Whitney es
\begin{equation}\label{eq:bsd-suc-aw}
 \Delta_{\mathrm{AW}}[x_0\to\cdots\to x_q]
 =\sum_{j=0}^{q}
 [x_0\to\cdots\to x_j]\otimes
 [x_j\to\cdots\to x_q].
\end{equation}
La counidad vale uno en los vértices y cero en grados positivos. En
particular,
\begin{equation}\label{eq:bsd-suc-group-like}
 \Delta_{\mathrm{AW}}[x]=[x]\otimes[x],
 \qquad \varepsilon_{\mathrm{AW}}[x]=1.
\end{equation}
La propiedad de elemento grupal pertenece al vértice de estado completo: se aplica
antes de olvidar hoja, ruta, acarreo o memoria.

Sea \(P_E^{\mathrm{Man}}\to A[0]\) el complejo de Manin basado con su
aumento, y sea \(G_{m,n}\to A[0]\) el aumento inducido por
\(\varepsilon_{\mathrm{AW}}\). El blanco conservativo es el producto
fibrado
\begin{equation}\label{eq:bsd-suc-pullback}
 P_E^{\mathrm{gen}}
 :=P_E^{\mathrm{Man}}\times_{A[0]}G_{m,n}.
\end{equation}
Si \(s:A[0]\to P_E^{\mathrm{Man}}\) es la sección basada, entonces
\begin{equation}\label{eq:bsd-suc-section}
 \widehat\Pi_0(c)
 :=\bigl(s\varepsilon_{\mathrm{AW}}(c),c\bigr)
\end{equation}
es una sección conservativa: su segunda proyección devuelve \(c\), mientras
la primera publica su coordenada modular. La historia y la sombra modular
quedan acopladas en un solo objeto.

\subsection*{La misma unidad en los canales Kato y Poitou--Tate}

Las compatibilidades de fase, hoja, corestricción y levantamiento Hensel
definen sobre \(P_E^{\mathrm{gen}}\) un par de publicaciones
\begin{equation}\label{eq:bsd-suc-copub}
 P_E=(P_E^{K},P_E^{\mathrm{PT}}).
\end{equation}
Ambas componentes reciben la misma unidad genealógica \(1_E\), de manera
que
\begin{equation}\label{eq:bsd-suc-zk-zpt}
 P_E^{K}(1_E)=z_K,
 \qquad P_E^{\mathrm{PT}}(1_E)=z_{\mathrm{PT}}.
\end{equation}
No son dos elecciones independientes de generador.

Sea \(L_{\mathrm{root}}=A_n z_{\mathrm{PT}}\) el subcomplejo raíz basado,
normalizado por
\(\lambda_{\mathrm{PT}}(z_{\mathrm{PT}})=1\), y sea
\(p_{\mathrm{root}}\) el proyector de complejos sobre ese subcomplejo. Se pone
\(q_{\mathrm{root}}=I-p_{\mathrm{root}}\), que también es un proyector de
complejos; en particular, \(q_{\mathrm{root}}z_K\) es un ciclo. La compatibilidad de la
counidad en \eqref{eq:bsd-suc-group-like} con el par
\eqref{eq:bsd-suc-copub} da
\begin{equation}\label{eq:bsd-suc-normalizacion-raiz}
 \lambda_{\mathrm{PT}}(p_{\mathrm{root}}z_K)=1.
\end{equation}

La dualidad se formula a nivel de cadenas. Para un complejo perfecto
\(M\) sobre \(\Lambda\), sea
\begin{equation}\label{eq:bsd-suc-d3}
 D_3(M):=
 \operatorname{RHom}_{\Lambda}(M^{\iota},\Lambda)[-3].
\end{equation}
Si \(A\to C\to D_3(A)\) es el bloque de incidencia, se definen
\begin{equation}\label{eq:bsd-suc-orth-red}
 A^{\perp}:=\operatorname{Fib}\bigl(C\to D_3(A)\bigr),
 \qquad
 C^{\mathrm{red}}:=\operatorname{Cofib}\bigl(A\to A^{\perp}\bigr),
\end{equation}
y la forma reducida induce
\begin{equation}\label{eq:bsd-suc-xorth}
 X^{\mathrm{orth}}:C^{\mathrm{red}}
 \xrightarrow{\ \sim\ }D_3(C^{\mathrm{red}}).
\end{equation}
En la línea determinante reducida se obtiene una involución
\(J_{\mathrm{red}}^2=I\) y bases compatibles
\begin{equation}\label{eq:bsd-suc-jred}
 J_{\mathrm{red}}(\bar b_{\mathrm{PT}})=\bar b_K.
\end{equation}
Las dos hojas transportan el mismo escalar basado:
\begin{equation}\label{eq:bsd-suc-doble-hoja}
 u_{+}=u_{-}=\tau_{\mathrm{bas}}(D_K).
\end{equation}
La igualdad vive en las dos líneas identificadas por
\eqref{eq:bsd-suc-jred}; no se reemplaza por una ecuación entre un escalar y
su inverso. Asimismo, \(D_3(D_K)=A^{\perp}\) pertenece a la línea dual y no
se confunde con una segunda copia no orientada de la línea raíz.

\begin{theorem}[Rigidez de la raíz común]
\label{thm:bsd-suc-raiz-filler}
Con las bases y orientaciones anteriores,
\begin{equation}\label{eq:bsd-suc-raiz}
 p_{\mathrm{root}}z_K=z_{\mathrm{PT}}.
\end{equation}
\end{theorem}

\begin{proof}
Escríbase
\(p_{\mathrm{root}}z_K=u z_{\mathrm{PT}}\). Al aplicar
\(\lambda_{\mathrm{PT}}\) y usar la normalización de la base junto con
\eqref{eq:bsd-suc-normalizacion-raiz}, se obtiene
\[
 1=\lambda_{\mathrm{PT}}(p_{\mathrm{root}}z_K)
  =u\lambda_{\mathrm{PT}}(z_{\mathrm{PT}})=u.
\]
Esto prueba \eqref{eq:bsd-suc-raiz}; la identificación
\eqref{eq:bsd-suc-jred} asegura que la conclusión es independiente de la
hoja desde la que se publica.
\end{proof}

\begin{lemma}[Forma normal del complemento]
\label{lem:bsd-suc-forma-normal}
En el bloque finita--singular con
los módulos libres sobre \(A_n\), sea \(g\) primitivo ---esto es, parte de una
base del módulo de llegada--- y supóngase que
\(\operatorname{im}\partial\) no tiene componente en el sumando directo
\(A_n r\). Si el diferencial viene dado por
\begin{equation}\label{eq:bsd-suc-diferencial-normal}
 \partial f=3c e+b,
 \qquad \partial e=g,
 \qquad \partial b=-3c g,
\end{equation}
todo ciclo \(z_K=a r+u e+v b\) satisface
\begin{equation}\label{eq:bsd-suc-ciclo-normal}
 u=3cv,
 \qquad
 z_{\mathrm{PT}}-z_K=(1-a)r-v\,\partial f,
 \quad z_{\mathrm{PT}}=r.
\end{equation}
La igualdad raíz \eqref{eq:bsd-suc-raiz} fuerza \(a=1\); por tanto
\(h=-vf\) rellena la diferencia completa.
\end{lemma}

\begin{proof}
La condición \(\partial z_K=0\) y
\eqref{eq:bsd-suc-diferencial-normal} dan
\((u-3cv)g=0\). Como \(g\) forma parte de una base de un módulo libre,
la comparación de su coeficiente da \(u=3cv\). Sustituir esta igualdad en
\(r-z_K\) produce \eqref{eq:bsd-suc-ciclo-normal}. El coeficiente \(a\)
es precisamente la coordenada de \(p_{\mathrm{root}}z_K\) en la base
\(r=z_{\mathrm{PT}}\); el teorema \ref{thm:bsd-suc-raiz-filler} da
\(a=1\). Entonces
\(\partial(-vf)=-v\partial f=z_{\mathrm{PT}}-z_K\).
\end{proof}

La construcción propietaria del relleno es, por tanto, la fórmula explícita
\begin{equation}\label{eq:bsd-suc-hstar}
 h_*:=-vf,
 \qquad
 \partial h_*=z_{\mathrm{PT}}-z_K.
\end{equation}
No se introduce una homotopía abstracta como premisa: la cerradura fija
\(u=3cv\), la unidad raíz fija \(a=1\) y el generador \(f\) rellena
directamente la diferencia completa.

\paragraph{Control homotópico equivalente no gobernante.}
\noindent\textbf{Estatuto: CONTROL\_NO\_GOBERNANTE.}
Como comprobación posterior, si se empaqueta el mismo complemento mediante
un lector \(H_{\mathrm{FS}}\), la identidad
\begin{equation}\label{eq:bsd-suc-hfs}
 \partial H_{\mathrm{FS}}+H_{\mathrm{FS}}\partial
 =q_{\mathrm{root}}
\end{equation}
reproduce \eqref{eq:bsd-suc-hstar} al aplicarse a
\(q_{\mathrm{root}}z_K\). Este lector es una notación auxiliar de la
contracción ya realizada por \(h_*=-vf\); no es una obligación anterior ni
puede reabrir la construcción explícita.

\subsection*{Regulador Bockstein derivado de orden arbitrario}

Sea \(K\) un campo de característica cero,
\(\Lambda=K[[T]]\), y sea
\begin{equation}\label{eq:bsd-suc-S}
 S(T):V_{\Lambda}\longrightarrow W_{\Lambda},
 \qquad
 V_{\Lambda}\simeq W_{\Lambda}\simeq\Lambda^{m},
\end{equation}
un diferencial basado que se vuelve invertible sobre \(K((T))\). Se fija
la línea
\begin{equation}\label{eq:bsd-suc-lambda0}
 \lambda_0(S):=(\det_K V)^{-1}\otimes_K\det_K W,
 \quad
 V=V_{\Lambda}/TV_{\Lambda},\quad
 W=W_{\Lambda}/TW_{\Lambda}.
\end{equation}
Una forma de Smith basada tiene la expresión
\begin{equation}\label{eq:bsd-suc-smith}
 U(T)S(T)V(T)
 =\operatorname{diag}
 \bigl(T^{a_1},\ldots,T^{a_s},1,\ldots,1\bigr),
 \qquad a_i\geq1,
\end{equation}
y
\begin{equation}\label{eq:bsd-suc-orden}
 d=\operatorname{ord}_{T}\det S(T)=\sum_{i=1}^{s}a_i.
\end{equation}

La filtración \(T\)-ádica del complejo
\([V_{\Lambda}\xrightarrow{S}W_{\Lambda}]\) produce páginas
\(V_q=E_q^0\), \(W_q=E_q^1\) y Bocksteins
\(\beta_q:V_q\to W_q\). Intrínsecamente,
\begin{equation}\label{eq:bsd-suc-pages}
 A_q:=V_q/V_{q+1},
 \qquad
 B_q:=\ker(W_q\twoheadrightarrow W_{q+1})
      =\operatorname{im}\beta_q,
\end{equation}
y \(\beta_q\) induce
\begin{equation}\label{eq:bsd-suc-beta-q}
 \bar\beta_q:A_q\xrightarrow{\ \sim\ }B_q,
 \qquad
 \tau_q:=\det(\bar\beta_q).
\end{equation}
Las páginas son finitas y verifican
\begin{equation}\label{eq:bsd-suc-carry-orden}
 d=\sum_{q\geq1}q\,\dim A_q
  =\sum_{q\geq1}\dim V_q.
\end{equation}
La primera página ve
\(r_0=\dim V_1\); el acarreo de orden superior es
\begin{equation}\label{eq:bsd-suc-carry-superior}
 c_I^{\deg}=d-r_0=\sum_{q\geq2}\dim V_q.
\end{equation}

La página regular es el isomorfismo inducido
\begin{equation}\label{eq:bsd-suc-pagina-regular}
 \bar S(0):A_0:=V/V_1
 \xrightarrow{\ \sim\ }
 B_0:=\operatorname{im}S(0),
 \qquad
 \tau_{\mathrm{reg}}:=\det\bar S(0).
\end{equation}
Las sucesiones exactas de \eqref{eq:bsd-suc-pages}, con sus signos de
Koszul, pegan estas trivializaciones en un único elemento basado
\begin{equation}\label{eq:bsd-suc-rboc}
 R_{\mathrm{Boc}}^{\mathrm{der}}(S)
 :=\tau_{\mathrm{reg}}\otimes
   \bigotimes_{q\geq1}\tau_q
 \in\lambda_0(S),
\end{equation}
donde el símbolo tensorial abrevia los isomorfismos de pegado de
Knudsen--Mumford.

\begin{theorem}[Identidad determinantal del regulador derivado]
\label{thm:bsd-suc-bockstein}
En la línea basada \(\lambda_0(S)\) se tiene
\begin{equation}\label{eq:bsd-suc-lt}
 R_{\mathrm{Boc}}^{\mathrm{der}}(S)
 =\operatorname{LT}_{T}(S)
 :=\left[T^{-d}\det S(T)\right]_{T=0}.
\end{equation}
La igualdad conserva la página regular, todas las páginas positivas y el
acarreo \eqref{eq:bsd-suc-carry-superior}.
\end{theorem}

\begin{proof}
En las bases Smith de \eqref{eq:bsd-suc-smith}, la página regular y cada
página positiva llevan el generador canónico de su línea. Los signos de
Koszul del pegado son exactamente los necesarios para restituir el orden de
las bases totales, por lo que su producto es el generador positivo de
\(\lambda_0(S)\). Por otra parte,
\[
 \det S(T)=\det U(T)^{-1}\det V(T)^{-1}T^d,
\]
y el término inicial es
\(\det U(0)^{-1}\det V(0)^{-1}\). Al deshacer el cambio de bases, tanto
\eqref{eq:bsd-suc-rboc} como el término inicial se multiplican por ese mismo
factor. Esto prueba \eqref{eq:bsd-suc-lt} sin identificar escalares antes de
fijar la línea y su base.
\end{proof}

Para cada \(q\geq1\), la dualidad de Poitou--Tate construida en la misma
página proporciona
\begin{equation}\label{eq:bsd-suc-jq}
 j_q:B_q\xrightarrow{\ \sim\ }
 A_q^{\vee}\otimes(I^q/I^{q+1}),
\end{equation}
y, tras orientar \(I^q/I^{q+1}\), el emparejamiento
\begin{equation}\label{eq:bsd-suc-altura-derivada}
 h_E^{(q)}(x,y)
 :=\bigl(j_q\bar\beta_q(x)\bigr)(y).
\end{equation}

% Ampliación sustantiva de los comparadores determinantes BSD.
% Módulo destinado al capítulo BSD de la Parte IV.

\section{Comparadores de la unidad determinantal común}
\label{sec:amp-bsd-comparadores}

La igualdad raíz \eqref{eq:bsd-suc-raiz} fija una sola unidad en los canales
Kato y Poitou--Tate. A partir de esa unidad se construyen ahora los
comparadores que transportan el regulador derivado hasta su publicación
escalar. Ningún comparador elige una normalización después de conocer el
término principal: todos proceden de morfismos de complejos, sucesiones
exactas y bases ya orientadas.

\subsection{La cadena de líneas determinantes}
\label{subsec:amp-bsd-lineas}

Para un complejo perfecto basado \(C\), se emplea la convención
\begin{equation}\label{eq:amp-bsd-linea-determinante}
 \lambda(C):=\bigotimes_i
 \det H^i(C)^{(-1)^i}.
\end{equation}
Sean \(\mathcal C_E^K\), \(\mathcal C_E^{\mathrm{Iw}}\),
\(\mathcal C_E^{\mathrm{PT}}\) y \(\mathcal C_E^{\mathrm{loc}}\),
respectivamente, el complejo publicado por el canal de Kato, su modelo
Iwasawa basado, el complejo reducido autodual de Poitou--Tate y el cono de
localización finita--singular. Sus líneas se enlazan mediante
\begin{equation}\label{eq:amp-bsd-cadena-lineas}
 \mathcal L_E^K
 \xrightarrow{\ \mathfrak c_{K/\mathrm{FERS}}\ }
 \mathcal L_E^{\mathrm{Iw}}
 \xrightarrow{\ \mathfrak c_{\mathrm{PT}}\ }
 \mathcal L_E^{\mathrm{ht}}
 \xrightarrow{\ \mathfrak c_{\mathrm{STS}}\ }
 \mathcal L_E^{\mathrm{fin}}
 \xrightarrow{\ \mathfrak c_{\mathrm{tors}}\ }
 \mathcal L_E^{\mathrm{ar}}
 \xrightarrow{\ \mathfrak c_{\infty}\ }
 \mathcal L_E.
\end{equation}
Las siglas \(\mathrm{STS}\) designan en esta fórmula el paso conjunto de
Smith, Tamagawa y Tate--Shafarevich. Todas las flechas de
\eqref{eq:amp-bsd-cadena-lineas} son isomorfismos de líneas orientadas.

\subsection{Comparación de cadenas Kato--FERS desde la unidad común}
\label{subsec:amp-bsd-comparacion-cadenas-kato-fers}

La equivalencia que origina el primer comparador puede construirse antes de
tomar determinantes. Sea \(F^q\mathcal C\) la filtración por profundidad
nonádica, identificada con la filtración \(T\)-ádica mediante el acarreo, y
escríbase
\begin{equation}\label{eq:amp-bsd-complejos-kato-iwasawa}
 \mathcal C_E^K
 :=\operatorname{Tot}\!\left(
 P_E^K\widehat\Pi_0G_{m,n}\right),
 \qquad
 \mathcal C_E^{\mathrm{Iw}}
 :=[V_\Lambda\xrightarrow{S_E(T)}W_\Lambda].
\end{equation}
Sea
\(\operatorname{car}_T:G_{m,n}\to\mathcal C_E^{\mathrm{Iw}}\) el lector
de acarreo que asigna a una ruta \(\gamma\) el generador determinado por sus
caras inicial y final, multiplicado por
\(T^{\operatorname{Carr}(\gamma)}\). La multiplicación de concatenaciones en
el modelo Iwasawa se denota por \(\mu_{\mathrm{Iw}}\). La diagonal
Alexander--Whitney y la publicación \(P_E^K\) definen entonces, antes de
tomar cohomología o determinantes, el mapa de cadenas
\begin{equation}\label{eq:amp-bsd-phi-kato-fers}
\begin{aligned}
 \Phi_{K/\mathrm{FERS}}&:
 \mathcal C_E^K\longrightarrow\mathcal C_E^{\mathrm{Iw}},\\
 \Phi_{K/\mathrm{FERS}}
 &:=
 \mu_{\mathrm{Iw}}\circ
 \bigl(P_E^K\widehat\otimes\operatorname{car}_T\bigr)
 \circ\Delta_{\mathrm{AW}},\\
 \Phi_{K/\mathrm{FERS}}(z_K)&=z_{\mathrm{Iw}}.
\end{aligned}
\end{equation}
donde \(z_{\mathrm{Iw}}\) es el generador raíz inducido por \(1_E\). Así,
cada simplex normalizado se envía a la suma de sus cortes
Alexander--Whitney, conservando en cada sumando las dos caras y la potencia
de acarreo. La fórmula no consulta el término principal analítico.

\begin{lemma}[Criterio generador a generador para la equivalencia]
\label{lem:amp-bsd-bases-emparejadas-kato-fers}
En cada grado \(i\), sea \(\mathcal B_i^K\) la base finita de simplices
normalizados, ordenada por profundidad, caras y acarreo, y sea
\(\mathcal B_i^{\mathrm{Iw}}\) la base formada por los generadores con las
mismas caras inicial y final en el modelo Iwasawa. Para una aplicación
candidata
\[
 \upsilon_i:\mathcal B_i^K\longrightarrow
 \mathcal B_i^{\mathrm{Iw}},
\]
defínase el residual de grado cero
\begin{equation}\label{eq:amp-bsd-residual-generador-graduado}
 \varepsilon_i(b)
 :=\Phi_{K/\mathrm{FERS}}^i(b)\bmod T-\upsilon_i(b).
\end{equation}
Si \(\upsilon_i\) es una biyección y
\(\varepsilon_i(b)=0\) para todo \(b\in\mathcal B_i^K\), entonces los dos
módulos completados del grado \(i\) son libres del mismo rango
\(d_i=|\mathcal B_i^K|\) y, para cada generador,
\begin{equation}\label{eq:amp-bsd-phi-en-cada-generador}
 \Phi_{K/\mathrm{FERS}}^i(b)
 =\upsilon_i(b)
  +T\sum_{c\in\mathcal B_i^{\mathrm{Iw}}}
    \lambda_{c,b}(T)c,
 \qquad \lambda_{c,b}(T)\in\Lambda.
\end{equation}
En particular,
\(\operatorname{gr}^{F}\Phi_{K/\mathrm{FERS}}^i=I\) sobre todos los
generadores, no sólo sobre la unidad raíz.
\end{lemma}

\begin{proof}
La anulación de \eqref{eq:amp-bsd-residual-generador-graduado} dice que la
columna de \(b\) tiene término constante \(\upsilon_i(b)\). Los restantes
coeficientes pertenecen a \(T\Lambda\), lo que da
\eqref{eq:amp-bsd-phi-en-cada-generador}. La biyectividad identifica las
bases y fija el mismo rango. En esas bases, la reducción módulo \(T\) de la
matriz de \(\Phi^i\) es la identidad; de ahí
\(\operatorname{gr}^{F}\Phi^i=I\).
\end{proof}

La igualdad raíz
\(\Phi_{K/\mathrm{FERS}}(z_K)=z_{\mathrm{Iw}}\) fija la normalización común.
La escritura de \(P_E^K\) sobre cada simplex normalizado, la biyección
\(\upsilon_i\) y los residuales \(\varepsilon_i(b)\) proporcionan un
control generador a generador de la equivalencia ya construida; no eligen ni
gobiernan la conclusión BSD.

\begin{lemma}[Equivalencia filtrada y normalización de la raíz]
\label{lem:amp-bsd-equivalencia-kato-fers}
Para el mapa de cadenas construido en
\eqref{eq:amp-bsd-phi-kato-fers}, la aplicación \(\upsilon_i\) es la
correspondencia de bases inducida por las caras y el acarreo de cada
simplejo. La fórmula de Alexander--Whitney verifica en cada grado el
criterio del lema \ref{lem:amp-bsd-bases-emparejadas-kato-fers}; éste es un
control de la equivalencia ya construida y no una premisa adicional.
La aplicación \(\Phi_{K/\mathrm{FERS}}\) es una equivalencia basada de
complejos perfectos. En las bases ordenadas por profundidad, sus matrices en
cada grado tienen la forma
\begin{equation}\label{eq:amp-bsd-matriz-kato-fers}
 [\Phi_{K/\mathrm{FERS}}^i](T)=I+TB_i(T),
 \qquad B_i(T)\in M_{d_i}(\Lambda).
\end{equation}
Por tanto, la unidad
\begin{equation}\label{eq:amp-bsd-rho-determinantal}
 \rho_{E,\ell}(T)
 :=\prod_i\det\!\left(I+TB_i(T)\right)^{(-1)^i}
 \in\Lambda^\times
\end{equation}
satisface \(\rho_{E,\ell}(0)=1\).
\end{lemma}

\begin{proof}
La identidad simplicial
\(\partial\Delta_{\mathrm{AW}}
=(\partial\otimes I+(-1)^{\deg}I\otimes\partial)
\Delta_{\mathrm{AW}}\), junto con la compatibilidad de
\(P_E^K\) con caras y degeneraciones, demuestra que
\eqref{eq:amp-bsd-phi-kato-fers} conmuta con los diferenciales. En el graduado
asociado, el lema
\ref{lem:amp-bsd-bases-emparejadas-kato-fers} identifica todos los
generadores y da la identidad. En esas bases,
\eqref{eq:amp-bsd-phi-en-cada-generador} es precisamente
\eqref{eq:amp-bsd-matriz-kato-fers}.

La completitud de \(\Lambda=K[[T]]\) hace converger la serie
\begin{equation}\label{eq:amp-bsd-inversa-kato-fers}
 (I+TB_i(T))^{-1}
 =\sum_{n\geq0}(-TB_i(T))^n;
\end{equation}
estas inversas también conmutan con los diferenciales, por lo que forman la
inversa basada de \(\Phi_{K/\mathrm{FERS}}\). Al tomar el determinante
alternado se obtiene \eqref{eq:amp-bsd-rho-determinantal}; evaluar en
\(T=0\) da uno. Así, la normalización no se postula a partir del término
principal: es la coordenada constante de una equivalencia que es la identidad
sobre la unidad raíz.
\end{proof}

El primer comparador es el determinante de
\(\Phi_{K/\mathrm{FERS}}\). Si \(\mathfrak z_K\) denota el elemento
procedente de la unidad \(1_E\), su imagen se caracteriza por
\begin{equation}\label{eq:amp-bsd-comparador-kato}
 \mathfrak c_{K/\mathrm{FERS}}(\mathfrak z_K)
 =\left[T^{-d}\rho_{E,\ell}(T)\det S_E(T)\right]_{T=0},
 \qquad
 \rho_{E,\ell}(0)=1.
\end{equation}
La unidad de \(\rho_{E,\ell}\) es orientada y vale uno en la raíz; por ello no
altera el orden ni la base. El teorema
\ref{thm:bsd-suc-bockstein} transforma \eqref{eq:amp-bsd-comparador-kato} en
\begin{equation}\label{eq:amp-bsd-kato-bockstein}
 \mathfrak c_{K/\mathrm{FERS}}(\mathfrak z_K)
 =R_{\mathrm{Boc}}^{\mathrm{der}}(S_E)
 =\tau_{\mathrm{reg}}\otimes\bigotimes_{q\geq1}\tau_q.
\end{equation}

El comparador de Poitou--Tate se construye página a página. Para
\(q\geq1\), el isomorfismo \(j_q\) de
\eqref{eq:bsd-suc-jq} y el Bockstein reducido
\(\bar\beta_q\) inducen
\begin{equation}\label{eq:amp-bsd-comparador-pt-q}
 \mathfrak c_{\mathrm{PT},q}(\tau_q)
 :=\det\!\left(j_q\circ\bar\beta_q\right)
 =\det h_E^{(q)}.
\end{equation}
La dirección regular se transporta mediante
\(\det\bar S_E(0)=\tau_{\mathrm{reg}}\). El pegado de Knudsen--Mumford y los
signos de Koszul definidos en \eqref{eq:bsd-suc-rboc} proporcionan el único
isomorfismo orientado
\begin{equation}\label{eq:amp-bsd-comparador-pt}
 \mathfrak c_{\mathrm{PT}}
 \left(\tau_{\mathrm{reg}}\otimes\bigotimes_{q\geq1}\tau_q\right)
 =\tau_{\mathrm{reg}}^{\mathrm{ht}}
  \otimes\bigotimes_{q\geq1}\det h_E^{(q)}.
\end{equation}
En la primera página estable, la forma \(h_E^{(1)}\) es la forma
Poincaré--Néron--Tate sobre la red libre de Mordell--Weil; su determinante es
\begin{equation}\label{eq:amp-bsd-regulador-nt}
 \det h_E^{(1)}=\operatorname{Reg}(E)
\end{equation}
en la base inducida por la unidad raíz.

\begin{lemma}[Despliegue Poitou--Tate del grado aritmético]
\label{lem:amp-bsd-grado-pt}
Sea
\[
 \operatorname{Flag}_{\mathrm{Boc}}(E)
 :=\bigoplus_{q\geq1}V_q
\]
la bandera finita de páginas positivas, cada una con la base heredada de la
unidad raíz. La reducción ortogonal y los isomorfismos \(j_q\) construyen un
isomorfismo basado
\begin{equation}\label{eq:amp-bsd-psi-pt}
 \Psi_{\mathrm{PT}}:
 \bigl(E(\mathbb Q)/E(\mathbb Q)_{\mathrm{tors}}\bigr)\otimes K
 \xrightarrow{\ \sim\ }\operatorname{Flag}_{\mathrm{Boc}}(E).
\end{equation}
En particular,
\begin{equation}\label{eq:amp-bsd-rango-bandera}
 \operatorname{rank}E(\mathbb Q)
 =\sum_{q\geq1}\dim_KV_q
 =\operatorname{ord}_T\det S_E(T).
\end{equation}
\end{lemma}

\begin{proof}
Se filtra el complejo reducido por las potencias del ideal de aumento. En el
cociente de nivel \(q\), la dirección regular ya ha sido separada por
\(\bar S_E(0)\), y el único bloque superviviente es \(V_q\). La equivalencia
\(X^{\mathrm{orth}}\) identifica el bloque opuesto con su dual, mientras
\(j_q\bar\beta_q\) proporciona el emparejamiento perfecto que levanta la
base del cociente al nivel anterior. Comenzando en la última página no nula y
repitiendo este levantamiento, se obtiene \(\Psi_{\mathrm{PT}}\).

En las bases ordenadas por profundidad, la matriz del mapa resultante es
triangular por bloques y todos sus bloques diagonales son identidades: en el
primer bloque porque \(z_K\) y \(z_{\mathrm{PT}}\) tienen la misma raíz, y
en los restantes porque \(j_q\) y \(\bar\beta_q\) son isomorfismos basados.
Por ello \(\Psi_{\mathrm{PT}}\) es invertible y no pierde ninguna página.
Tomando dimensiones y usando
\eqref{eq:bsd-suc-carry-orden} se obtiene
\eqref{eq:amp-bsd-rango-bandera}. Al tomar determinantes, el mismo
levantamiento es exactamente el pegado de
\eqref{eq:amp-bsd-comparador-pt}; así, la igualdad de rangos y la igualdad de
líneas proceden de un solo mapa basado.
\end{proof}

\subsection{Triángulo de localización, rango de Smith y factores finitos}
\label{subsec:amp-bsd-smith-sha}

Para cada lugar \(v\), la condición local finita es un subcomplejo basado
\(\mathcal C_{E,v}^{\mathrm{fin}}\subset\mathcal C_{E,v}\), y se define
\(\mathcal C_{E,v}^{\mathrm{sing}}\) por el triángulo
\begin{equation}\label{eq:amp-bsd-triangulo-local}
 \mathcal C_{E,v}^{\mathrm{fin}}
 \xrightarrow{\ i_v\ }\mathcal C_{E,v}
 \longrightarrow\mathcal C_{E,v}^{\mathrm{sing}}
 \longrightarrow\mathcal C_{E,v}^{\mathrm{fin}}[1].
\end{equation}
Si \(\operatorname{res}_v\) denota la restricción global, el mapa de
localización es, a nivel de cadenas,
\begin{equation}\label{eq:amp-bsd-mapa-localizacion}
 \ell_E:
 \mathcal C_E^{\mathrm{glob,red}}\oplus
 \bigoplus_v\mathcal C_{E,v}^{\mathrm{fin}}
 \longrightarrow\bigoplus_v\mathcal C_{E,v},
 \qquad
 \ell_E\bigl(x,(y_v)_v\bigr)
 :=\bigl(\operatorname{res}_v x-i_vy_v\bigr)_v .
\end{equation}
La definición por cono del complejo Selmer reducido da el triángulo
\begin{equation}\label{eq:amp-bsd-triangulo-localizacion}
 \mathcal C_E^{\mathrm{red}}
 \longrightarrow
 \mathcal C_E^{\mathrm{glob,red}}\oplus
 \bigoplus_v\mathcal C_{E,v}^{\mathrm{fin}}
 \xrightarrow{\ \ell_E\ }
 \bigoplus_v\mathcal C_{E,v}
 \longrightarrow\mathcal C_E^{\mathrm{red}}[1].
\end{equation}
Por tanto, la exactitud local--global procede del cono de un mapa escrito en
\eqref{eq:amp-bsd-mapa-localizacion}; no se añade después como una relación
entre cardinales.

Se cancela en \eqref{eq:amp-bsd-triangulo-localizacion} el sumando raíz
común, que es el mismo en ambos canales por
\eqref{eq:bsd-suc-raiz}. El complemento resultante se denota por
\begin{equation}\label{eq:amp-bsd-cono-local-reducido}
 \mathcal C_E^{\mathrm{loc,red}}
 :=q_{\mathrm{root}}\operatorname{Cone}(\ell_E)[-1].
\end{equation}
Tras separar la red libre de Mordell--Weil y su dual mediante
\(X^{\mathrm{orth}}\), una elección de las bases integrales ya orientadas
representa este complejo por
\begin{equation}\label{eq:amp-bsd-complejo-smith}
 \mathcal C_E^{\mathrm{loc,red}}
 \simeq[M_E\xrightarrow{\ D_E\ }N_E].
\end{equation}

\begin{lemma}[Aciclicidad racional y rango completo]
\label{lem:amp-bsd-aciclicidad-rango}
El complejo de \eqref{eq:amp-bsd-complejo-smith} es acíclico después de
tensorizar con \(\mathbb Q\). En particular,
\[
 D_E\otimes\mathbb Q:M_E\otimes\mathbb Q
 \xrightarrow{\ \sim\ }N_E\otimes\mathbb Q
\]
es un isomorfismo; \(M_E\) y \(N_E\) tienen el mismo rango y \(D_E\) tiene
rango completo.
\end{lemma}

\begin{proof}
Sea \(x\) un ciclo del complejo
\(\mathcal C_E^{\mathrm{loc,red}}\otimes\mathbb Q\). La reducción ha
eliminado el sumando \(p_{\mathrm{root}}\), luego
\(q_{\mathrm{root}}x=x\). En las bases integrales orientadas empleadas en
\eqref{eq:amp-bsd-complejo-smith}, el complemento finita--singular se
descompone en los bloques normales del
\cref{lem:bsd-suc-forma-normal}. En cada bloque \(\lambda\), el coeficiente
raíz ya es cero y el ciclo se escribe
\[
 x_{\lambda}=u_{\lambda}e_{\lambda}
                  +v_{\lambda}b_{\lambda},
 \qquad
 \partial x_{\lambda}
   =(u_{\lambda}-3c_{\lambda}v_{\lambda})g_{\lambda}=0.
\]
Como \(g_{\lambda}\) pertenece a la base libre, se tiene
\(u_{\lambda}=3c_{\lambda}v_{\lambda}\), y la fórmula explícita del
diferencial da
\[
 x_{\lambda}
   =v_{\lambda}(3c_{\lambda}e_{\lambda}+b_{\lambda})
   =\partial(v_{\lambda}f_{\lambda}).
\]
La suma es finita en cada grado, de modo que
\(x=\partial\bigl(\sum_{\lambda}v_{\lambda}f_{\lambda}\bigr)\).
Todo ciclo es, por tanto, una frontera por la misma forma normal que produjo
el filler \eqref{eq:bsd-suc-hstar}, sin introducir
\(H_{\mathrm{FS}}\) como premisa. La igualdad
\(p_{\mathrm{root}}z_K=z_{\mathrm{PT}}\) es la que permite cancelar la
dirección raíz antes de este cálculo; sin ella quedaría una clase de grado
cero. La equivalencia autodual
\(X^{\mathrm{orth}}:\mathcal C_E^{\mathrm{red}}\simeq
D_3(\mathcal C_E^{\mathrm{red}})\) empareja los dos grados restantes y
transporta sus bases con orientación opuesta complementaria. De aquí se
obtiene \(\operatorname{rank}M_E=\operatorname{rank}N_E\).
Finalmente, la aciclicidad de un complejo de dos términos de igual rango
equivale a que \(D_E\otimes\mathbb Q\) sea invertible.
\end{proof}

El lema permite aplicar la forma normal de Smith sin presuponer ninguno de
sus divisores. Existen cambios de base orientados \(U,V\) tales que
\begin{equation}\label{eq:amp-bsd-morfismo-smith}
 U D_E V=\operatorname{diag}(s_1,\ldots,s_t),
 \qquad s_i\in\mathbb Z\setminus\{0\}.
\end{equation}
Por tanto,
\begin{equation}\label{eq:amp-bsd-grupo-smith}
 F_E:=\operatorname{coker}D_E
 \simeq\bigoplus_{i=1}^{t}\mathbb Z/|s_i|\mathbb Z,
 \qquad
 |F_E|=\prod_{i=1}^{t}|s_i|<\infty.
\end{equation}

\begin{lemma}[Sucesión finita de localización]
\label{lem:amp-bsd-exactitud-sha-tamagawa}
El fragmento torsional de la sucesión larga de cohomología de
\eqref{eq:amp-bsd-triangulo-localizacion} es
\begin{equation}\label{eq:amp-bsd-exacta-sha-tamagawa}
 0\longrightarrow\operatorname{Sha}(E)
 \xrightarrow{\ \iota_{\mathrm{Sha}}\ }F_E
 \xrightarrow{\ \partial_{\mathrm{loc}}\ }
 \bigoplus_v\Phi_v(\mathbb F_v)
 \longrightarrow0,
 \qquad c_v=|\Phi_v(\mathbb F_v)|.
\end{equation}
\end{lemma}

\begin{proof}
Un cociclo global representa una clase de
\(\operatorname{Sha}(E)\) precisamente cuando su restricción en cada lugar
admite una primitiva en
\(\mathcal C_{E,v}^{\mathrm{fin}}\). La clase del par formado por ese
cociclo y sus primitivas locales en el cono de
\eqref{eq:amp-bsd-mapa-localizacion} define
\(\iota_{\mathrm{Sha}}\). Si esta clase es una frontera, la componente global es una
frontera y la clase de Tate--Shafarevich es nula; así,
\(\iota_{\mathrm{Sha}}\) es inyectiva.

El mapa \(\partial_{\mathrm{loc}}\) toma una clase del cono, la restringe a
cada cociente
\(\mathcal C_{E,v}^{\mathrm{sing}}\) de
\eqref{eq:amp-bsd-triangulo-local} y conserva su clase de componente. La
identificación
\[
 H^0(\mathcal C_{E,v}^{\mathrm{sing}})_{\mathrm{tors}}
 \simeq\Phi_v(\mathbb F_v)
\]
es inducida por el cociente finita--singular, no por el orden del grupo.
Una clase del cono tiene residuo cero exactamente cuando sus representantes
locales pueden elegirse en los subcomplejos finitos; por la definición
anterior, ese núcleo es la imagen de \(\iota_{\mathrm{Sha}}\).
Finalmente, la última flecha de
\eqref{eq:amp-bsd-triangulo-local} levanta cada clase de componente al
complejo local. La exactitud del cono
\eqref{eq:amp-bsd-triangulo-localizacion}, después de retirar los dos
sumandos libres emparejados por \(X^{\mathrm{orth}}\), convierte ese
levantamiento en una preimagen por \(\partial_{\mathrm{loc}}\). Esto prueba
la sobreyectividad y, por consiguiente, toda
\eqref{eq:amp-bsd-exacta-sha-tamagawa}.
\end{proof}

Ahora, y sólo después del lema
\ref{lem:amp-bsd-aciclicidad-rango}, el grupo \(F_E\) es finito. La
inyección de \eqref{eq:amp-bsd-exacta-sha-tamagawa} demuestra entonces la
finitud de \(\operatorname{Sha}(E)\). Al tomar órdenes en esa sucesión
exacta se obtiene
\begin{equation}\label{eq:amp-bsd-cardinal-smith}
 |F_E|=|\operatorname{Sha}(E)|\prod_v c_v.
\end{equation}
Éste es el comparador
\(\mathfrak c_{\mathrm{STS}}\): transporta los factores elementales de Smith
a los factores Tate--Shafarevich y Tamagawa conservando su orden en la línea,
en lugar de insertar sus cardinales una vez calculado el término principal.

La parte torsional se construye aplicando el determinante a la red de
Mordell--Weil y a su dual de Pontryagin. Los dos factores finitos son duales y
aportan
\begin{equation}\label{eq:amp-bsd-comparador-torsion}
 \mathfrak c_{\mathrm{tors}}:quad
 u\longmapsto
 \frac{u}{|E(\mathbb Q)_{\mathrm{tors}}|^2}.
\end{equation}
El comparador arquimediano se obtiene integrando la diferencial de Néron
orientada sobre el ciclo real basado:
\begin{equation}\label{eq:amp-bsd-comparador-periodo}
 \Omega_E:=\int_{E(\mathbb R)}|\omega_E|,
 \qquad
 \mathfrak c_\infty:quad u\longmapsto\frac{u}{\Omega_E}.
\end{equation}
La diferencial y el ciclo se fijan antes de evaluar \(L(E,s)\); por ello
\eqref{eq:amp-bsd-comparador-periodo} no normaliza retrospectivamente la
sección analítica.

\subsection{Separación causal de primitivas, identidades y consecuencias}
\label{subsec:amp-bsd-separacion-causal}

El criterio generador a generador del
lema \ref{lem:amp-bsd-bases-emparejadas-kato-fers} audita en cada grado la
equivalencia ya construida. No condiciona su existencia. El orden lógico
del paquete es
\begin{equation}\label{eq:amp-bsd-diagrama-causal}
\begin{array}{c}
 \begin{gathered}
  1_E,\ \Delta_{\mathrm{AW}},\ P_E^K,\ P_E^{\mathrm{PT}},
  \ h_*=-vf,\ X^{\mathrm{orth}},\\
  \{\,\mathcal C_{E,v}^{\mathrm{fin}}\to\mathcal C_{E,v}\,\}_v,\\
  \text{bases de Mordell--Weil y Pontryagin;}\\
  \text{diferencial y ciclo de Néron}
 \end{gathered}
 \\[2mm]\Downarrow\\[-1mm]
 \begin{gathered}
  \Phi_{K/\mathrm{FERS}}\text{ equivalencia},\quad
  \rho_{E,\ell}(0)=1,\quad
  p_{\mathrm{root}}z_K=z_{\mathrm{PT}},\\
  R_{\mathrm{Boc}}^{\mathrm{der}}=\operatorname{LT}_T(S_E),\quad
  \det(j_q\bar\beta_q)=\det h_E^{(q)},\\
  H^*(\mathcal C_E^{\mathrm{loc,red}}\otimes\mathbb Q)=0,\quad
  D_E\otimes\mathbb Q\text{ isomorfismo},\\
  0\to\operatorname{Sha}(E)\to F_E
  \to\bigoplus_v\Phi_v(\mathbb F_v)\to0
 \end{gathered}
 \\[2mm]\Downarrow\\[-1mm]
 \begin{gathered}
  d_{\mathrm{an}}=\operatorname{rank}E(\mathbb Q),\qquad
  \operatorname{Reg}(E),\qquad
  |\operatorname{Sha}(E)|\prod_vc_v,\\
  |E(\mathbb Q)_{\mathrm{tors}}|^{-2},\qquad
  \Omega_E^{-1},\qquad
  \text{igualdad del término principal}.
 \end{gathered}
\end{array}
\end{equation}
La primera fila contiene únicamente objetos y mapas ya basados, junto con
la verificación generador a generador indicada; la segunda, identidades
demostradas por esos mapas; la tercera, sus consecuencias sobre la línea
determinante. En particular, ningún factor de la última fila se
emplea para normalizar retroactivamente una flecha de la primera.

\begin{proposition}[Paquete exacto de comparadores]
\label{prop:amp-bsd-paquete-comparadores-exacto}
Los comparadores de
\eqref{eq:amp-bsd-cadena-lineas} proceden de un único paquete de mapas de
complejos y satisfacen, en el orden de construcción,
\begin{equation}\label{eq:amp-bsd-paquete-identidades}
 \begin{gathered}
 \Phi_{K/\mathrm{FERS}}(z_K)=z_{\mathrm{Iw}},\qquad
 \operatorname{gr}^{F}\Phi_{K/\mathrm{FERS}}=I,\\
 \det\nolimits_{\mathrm{alt}}\Phi_{K/\mathrm{FERS}}
   =\rho_{E,\ell}(T),\qquad \rho_{E,\ell}(0)=1,\\
 \det(j_q\bar\beta_q)=\det h_E^{(q)}\quad(q\geq1),\\
 \Psi_{\mathrm{PT}}:
  \bigl(E(\mathbb Q)/E(\mathbb Q)_{\mathrm{tors}}\bigr)\otimes K
  \xrightarrow{\sim}\operatorname{Flag}_{\mathrm{Boc}}(E),\\
 p_{\mathrm{root}}z_K=z_{\mathrm{PT}},\qquad
 h_*=-vf,\quad \partial h_*=z_{\mathrm{PT}}-z_K,\\
 H^*(\mathcal C_E^{\mathrm{loc,red}}\otimes\mathbb Q)=0,\qquad
 U D_E V=\operatorname{diag}(s_1,\ldots,s_t),\\
 0\longrightarrow\operatorname{Sha}(E)
  \longrightarrow F_E
  \longrightarrow\displaystyle\bigoplus_v\Phi_v(\mathbb F_v)
  \longrightarrow0 .
 \end{gathered}
\end{equation}
La composición inducida en líneas determinantes es exactamente
\begin{equation}\label{eq:amp-bsd-paquete-composicion}
 \det(\Phi_{K/\mathrm{FERS}})
 \xrightarrow{\ \det(j_q\bar\beta_q)\ }
 \det(D_E)
 \xrightarrow{\ \mathrm{Smith}\ }
 \frac{|\operatorname{Sha}(E)|\prod_vc_v}
      {|E(\mathbb Q)_{\mathrm{tors}}|^2\,\Omega_E},
\end{equation}
con los signos de Knudsen--Mumford y las orientaciones de
\eqref{eq:amp-bsd-cadena-lineas}. En particular, el miembro derecho de
\eqref{eq:amp-bsd-paquete-composicion} es la coordenada final del elemento
transportado y no un dato con el que se construya alguna de las flechas.
\end{proposition}

\begin{proof}
La primera fila de \eqref{eq:amp-bsd-paquete-identidades} es
\eqref{eq:amp-bsd-phi-kato-fers} y
\eqref{eq:amp-bsd-matriz-kato-fers}--
\eqref{eq:amp-bsd-rho-determinantal}. La segunda es
\eqref{eq:amp-bsd-comparador-pt-q} y el lema
\ref{lem:amp-bsd-grado-pt}. La tercera combina la igualdad de la raíz con
el filler explícito \eqref{eq:bsd-suc-hstar}; su extensión lineal en la
forma normal induce el lector auxiliar utilizado en el lema
\ref{lem:amp-bsd-aciclicidad-rango}. Al restringir al complemento de la
raíz, esa reducción contrae todo ciclo racional. La última fila es
\eqref{eq:amp-bsd-morfismo-smith} y
\eqref{eq:amp-bsd-exacta-sha-tamagawa}. Aplicar el funtor determinante a
esas identidades, en ese mismo orden, da
\eqref{eq:amp-bsd-paquete-composicion}. Ninguna igualdad se obtiene
comparando primero los dos escalares finales.
\end{proof}

\begin{lemma}[Independencia causal del término principal]
\label{lem:amp-bsd-independencia-termino-principal}
Sea \(\mathscr P_E\) el conjunto de datos
\begin{equation}\label{eq:amp-bsd-primitivas-paquete}
 \mathscr P_E:=
 \left(
  1_E,\Delta_{\mathrm{AW}},P_E^K,P_E^{\mathrm{PT}},
  h_*=-vf,X^{\mathrm{orth}},
  \{\mathcal C_{E,v}^{\mathrm{fin}}\xrightarrow{i_v}
     \mathcal C_{E,v}\}_v,
  \text{bases y orientaciones},
  \omega_E,\gamma_E^{\mathbb R}
 \right),
\end{equation}
donde \(\omega_E\) es la diferencial de Néron y
\(\gamma_E^{\mathbb R}\) el ciclo real orientado. Todos los mapas del
paquete de la proposición
\ref{prop:amp-bsd-paquete-comparadores-exacto} se determinan por
\(\mathscr P_E\). No forman parte de \(\mathscr P_E\) los numerales
\[
 L^{(r)}(E,1),\quad r,\quad\operatorname{Reg}(E),\quad
 |\operatorname{Sha}(E)|,\quad(c_v)_v,\quad
 |E(\mathbb Q)_{\mathrm{tors}}|,\quad\Omega_E.
\]
Esos valores se obtienen, respectivamente, al evaluar la sección analítica,
tomar el grado de la bandera, calcular determinantes y formas de Smith,
tomar órdenes de grupos finitos e integrar
\(\omega_E\) sobre \(\gamma_E^{\mathbb R}\).
\end{lemma}

\begin{proof}
La fórmula \eqref{eq:amp-bsd-phi-kato-fers} usa sólo
\(1_E,\Delta_{\mathrm{AW}},P_E^K\) y el acarreo. Los mapas Bockstein y
Poitou--Tate usan la filtración, \(P_E^{\mathrm{PT}}\), las bases y la
autodualidad. El cono de \eqref{eq:amp-bsd-mapa-localizacion} usa los
mapas \(i_v\) y las restricciones locales; su contracción y su forma de
Smith usan el filler explícito \eqref{eq:bsd-suc-hstar}, su extensión
lineal auxiliar, \(X^{\mathrm{orth}}\) y las bases integrales. Los dos
últimos comparadores usan las redes torsionales y el
par \((\omega_E,\gamma_E^{\mathbb R})\). Ésta es una enumeración exhaustiva
de los argumentos que aparecen en las fórmulas constructoras. Los
numerales listados en el enunciado aparecen sólo después de aplicar
cohomología, determinante, cardinal o integración, de modo que no pueden
seleccionar retrospectivamente los mapas.
\end{proof}

\subsection{Composición, grados y término principal}
\label{subsec:amp-bsd-composicion}

Defínase el comparador total por la composición en el orden escrito:
\begin{equation}\label{eq:amp-bsd-comparador-total}
 \mathfrak C_E:=
 \mathfrak c_\infty\circ
 \mathfrak c_{\mathrm{tors}}\circ
 \mathfrak c_{\mathrm{STS}}\circ
 \mathfrak c_{\mathrm{PT}}\circ
 \mathfrak c_{K/\mathrm{FERS}}.
\end{equation}
Como \(p_{\mathrm{root}}z_K=z_{\mathrm{PT}}\), todas las flechas reciben la
misma unidad basada. El filler \(h_*=-vf\) de
\eqref{eq:bsd-suc-hstar} rellena explícitamente la diferencia; su extensión
lineal en la base normal produce el lector homotópico auxiliar y prueba que
ninguna clase adicional modifica el elemento de la línea durante el
transporte. El lector no es una premisa separada.

\begin{theorem}[Construcción de los comparadores basados]
\label{thm:amp-bsd-comparadores-construidos}
Para toda curva elíptica \(E/\mathbb Q\) con los complejos, bases,
orientaciones y datos locales definidos anteriormente, las aplicaciones
\eqref{eq:amp-bsd-phi-kato-fers},
\eqref{eq:amp-bsd-comparador-pt},
\eqref{eq:amp-bsd-triangulo-localizacion},
\eqref{eq:amp-bsd-comparador-torsion} y
\eqref{eq:amp-bsd-comparador-periodo}, junto con los lemas
\ref{lem:amp-bsd-aciclicidad-rango} y
\ref{lem:amp-bsd-exactitud-sha-tamagawa}, construyen los cinco comparadores de
\eqref{eq:amp-bsd-cadena-lineas}. Su composición preserva la unidad, el grado
y la orientación, y satisface
\begin{equation}\label{eq:amp-bsd-igualdad-grados}
 d_{\mathrm{an}}=\operatorname{ord}_T\det S_E(T)
 =d_{\mathrm{ar}}=\operatorname{rank}E(\mathbb Q).
\end{equation}
Además, \(\operatorname{Sha}(E)\) es finito y la igualdad del elemento
distinguido en \(\mathcal L_E\) se publica como
\begin{equation}\label{eq:amp-bsd-termino-principal}
 \frac{L^{(r)}(E,1)}{r!\,\Omega_E}
 =\frac{|\operatorname{Sha}(E)|\,\operatorname{Reg}(E)
        \prod_v c_v}
       {|E(\mathbb Q)_{\mathrm{tors}}|^2},
 \qquad r=\operatorname{rank}E(\mathbb Q).
\end{equation}
\end{theorem}

\begin{proof}
El comparador Kato--FERS y la condición
\(\rho_{E,\ell}(0)=1\) identifican el orden analítico con
\(d=\operatorname{ord}_T\det S_E(T)\) y el término inicial con
\(R_{\mathrm{Boc}}^{\mathrm{der}}(S_E)\). La identidad
\eqref{eq:bsd-suc-lt} conserva la página regular y cada página positiva. Los
isomorfismos \(j_q\) transportan esas páginas a las alturas derivadas. El
lema \ref{lem:amp-bsd-grado-pt} construye el isomorfismo de la bandera con la
red libre de Mordell--Weil y, junto con
\eqref{eq:bsd-suc-carry-orden}, prueba
\eqref{eq:amp-bsd-igualdad-grados}.

El lema \ref{lem:amp-bsd-aciclicidad-rango} deduce la aciclicidad racional
directamente del filler explícito \eqref{eq:bsd-suc-hstar}, de su reducción
lineal del complemento, de la cancelación de la raíz común y de
\(X^{\mathrm{orth}}\); de ahí se obtiene el rango completo de \(D_E\), no
como hipótesis adicional. La forma de Smith
\eqref{eq:amp-bsd-morfismo-smith} produce entonces el grupo finito \(F_E\).
El lema \ref{lem:amp-bsd-exactitud-sha-tamagawa} extrae de los triángulos
\eqref{eq:amp-bsd-triangulo-local} y
\eqref{eq:amp-bsd-triangulo-localizacion} la sucesión exacta finita; sólo
después se deducen la finitud de \(\operatorname{Sha}(E)\) y
\eqref{eq:amp-bsd-cardinal-smith}. Finalmente,
\eqref{eq:amp-bsd-regulador-nt},
\eqref{eq:amp-bsd-comparador-torsion} y
\eqref{eq:amp-bsd-comparador-periodo} publican, respectivamente, el regulador,
el cociente torsional y el período. La igualdad raíz elimina cualquier unidad
relativa entre las dos publicaciones. Al evaluar el mismo elemento basado en
ambos extremos de \eqref{eq:amp-bsd-comparador-total}, se obtiene
\eqref{eq:amp-bsd-termino-principal}.
\end{proof}

\subsection{Premisas y criterios de refutación}
\label{subsec:amp-bsd-falsadores}

La construcción parte de la unidad genealógica \(1_E\) y de la diagonal de
Alexander--Whitney. Emplea además las
publicaciones acopladas, todos los
Bocksteins y la página regular, el filler explícito \(h_*=-vf\), la
autodualidad reducida \(X^{\mathrm{orth}}\), los mapas de cadenas del
triángulo de localización y las bases orientadas de torsión y período. La equivalencia
Kato--FERS y su normalización en uno, el rango completo de \(D_E\), la
exactitud de \eqref{eq:amp-bsd-exacta-sha-tamagawa} y la finitud de
\(\operatorname{Sha}(E)\) son conclusiones de los lemas anteriores, no
premisas independientes. El criterio generador a generador del lema
\ref{lem:amp-bsd-bases-emparejadas-kato-fers} es un control no gobernante
que audita localmente el primer comparador. Los criterios que refutarían una flecha concreta
son:
\begin{enumerate}
 \item un coeficiente raíz distinto de uno, que introduciría una unidad
       relativa entre \(z_K\) y \(z_{\mathrm{PT}}\);
 \item un \(j_q\) no invertible o la omisión de una página Bockstein, que
       alteraría el grado o el determinante de altura;
 \item un divisor de Smith nulo, que impediría la finitud de \(F_E\);
 \item el fallo de exactitud de
       \eqref{eq:amp-bsd-exacta-sha-tamagawa}, que rompería la identificación
       de los factores finitos;
 \item una inversión de orientación en torsión o período, que cambiaría el
       elemento basado aunque preservase su línea no orientada.
\end{enumerate}
La proyección de Manin sin historia y el reescalamiento visible módulo tres
violan el primer criterio; por ello constituyen pruebas de necesidad para
esos modelos reducidos y no premisas del teorema
\ref{thm:amp-bsd-comparadores-construidos}.


\begin{theorem}[Comparación determinantal local--global]
\label{thm:bsd-suc-local-global}
Sean los comparadores basados ya construidos en el
\cref{thm:amp-bsd-comparadores-construidos}:
Kato--FERS/Iwasawa,
Poitou--Tate/Poincaré--Néron--Tate,
Smith--Tamagawa--Shafarevich, torsión y período.
Si \(d_{\mathrm{an}}\) es el grado de la sección analítica en la
especialización central y \(d_{\mathrm{ar}}\) el grado de la línea Selmer
basada, entonces
\begin{equation}\label{eq:bsd-suc-grados}
 d_{\mathrm{an}}=d=d_{\mathrm{ar}},
\end{equation}
y se obtiene la igualdad de elementos distinguidos
\begin{equation}\label{eq:bsd-suc-linea-final}
 \mathfrak z_{\mathrm{an}}(E)=\mathfrak z_{\mathrm{ar}}(E)
 \quad\text{en}\quad\mathcal L_E.
\end{equation}
El miembro analítico es la publicación del término inicial; el miembro
aritmético es la publicación del regulador completo, los factores
elementales de Smith,
los factores locales, la torsión y el período.
\end{theorem}

\begin{proof}
La construcción basada da
\([\Phi_{K/\mathrm{FERS}}^i](T)=I+TB_i(T)\) en cada grado. El comparador
basado de Kato--FERS expresa entonces la sección analítica corregida
como
\begin{equation}\label{eq:bsd-suc-fers}
 L_{\ell}^{\mathrm{corr}}(E,T)
 =\rho_{E,\ell}(T)\det S_E^{\mathrm{corr}}(T),
 \qquad \rho_{E,\ell}(0)=1.
\end{equation}
Por \eqref{eq:bsd-suc-orden}, su grado central es \(d\); por
\eqref{eq:bsd-suc-lt}, su término inicial es exactamente
\(R_{\mathrm{Boc}}^{\mathrm{der}}\), no ese regulador multiplicado por una
unidad indeterminada. Los isomorfismos \(j_q\) de
\eqref{eq:bsd-suc-jq} transportan cada \(\tau_q\) al determinante de la
altura derivada de su página. La página regular transporta el sumando
acíclico que las páginas positivas no ven.

La comparación de Poincaré--Néron--Tate identifica el grado de esa línea de
alturas con \(d_{\mathrm{ar}}\). Los factores elementales de Smith descomponen su parte
finita en los factores de Shafarevich y Tamagawa; el cociente de las redes
integrales aporta el cuadrado de la torsión, y el lector arquimediano aporta
el período. Todos estos mapas actúan sobre la misma línea raíz. Por el
teorema \ref{thm:bsd-suc-raiz-filler}, no queda una unidad relativa entre
las publicaciones Kato y Poitou--Tate. La composición preserva el grado y
el elemento basado, lo que prueba simultáneamente
\eqref{eq:bsd-suc-grados} y \eqref{eq:bsd-suc-linea-final}. La fórmula
escalar correspondiente se reserva para el reconocimiento final.
\end{proof}

\subsection*{Hipótesis y alcance}

La cadena BSD utiliza: el estado enriquecido completo y su coálgebra
Alexander--Whitney; el producto fibrado \eqref{eq:bsd-suc-pullback}; las dos
publicaciones construidas sobre la misma unidad; la línea raíz, su base y su
counidad; el cociente reducido perfecto y saturado, junto con
\(X^{\mathrm{orth}}\); el complejo finita--singular completo; las banderas
Bockstein finitas, sus signos de Koszul y la página regular; y las
comparaciones de altura, Smith, Tamagawa, Shafarevich, torsión y períodos
con una orientación común. La verificación grado a grado de las bases
emparejadas y de los residuales \(\varepsilon_i(b)\) queda registrada como
control local de la realización, no como premisa gobernante del cierre.
Bajo esta estructura de partida explícita,
\eqref{eq:bsd-suc-raiz}, \eqref{eq:bsd-suc-hstar},
\eqref{eq:bsd-suc-lt} y \eqref{eq:bsd-suc-linea-final} son igualdades
exactas.

La conservación de unidad, el recorrido de dos hojas y la publicación
local--global se formalizan mediante la coálgebra de cadenas, el producto
fibrado genealógico, la publicación acoplada, la reducción ortogonal y la
torre Bockstein. Las identidades auxiliares verifican por separado la propiedad de elemento grupal,
la conservación de la historia bajo \(\widehat\Pi_0\), la identidad de la
unidad en ambos canales, la involución reducida, la forma normal
finita--singular y la identidad determinantal de orden arbitrario.

\paragraph{Controles de ablación BSD.}
\noindent\textbf{Estatuto: CONTROL\_NO\_GOBERNANTE.}
Las tres ablaciones siguientes reciben las identidades ya demostradas y
falsifican únicamente modelos reducidos. No aportan premisas al propietario,
no condicionan su cierre y no poseen una arista causal de retorno hacia él.

\begin{ablation}[Proyección de Manin sin historia]
\label{abl:bsd-suc-manin}
Sustituir \(P_E^{\mathrm{gen}}\) por \(P_E^{\mathrm{Man}}\) borra datos que
la publicación acoplada necesita. En efecto, la ruta
\(\gamma=(\mathsf N\mathsf E)^{54}\) puede tener la misma celda y la misma
fase visible que la ruta trivial, pero conserva enrollamiento y memoria
distintos en \(G_{m,n}\). El aumento de Manin las identifica; la segunda
proyección de \eqref{eq:bsd-suc-pullback} no.
\end{ablation}

\begin{proof}
Todo lift basado al complejo de Manin factoriza, hasta homotopía, por la
sección del aumento \(s\varepsilon_{\mathrm{AW}}\). Por ello depende sólo
de la coordenada visible. En cambio, la segunda componente de
\(\widehat\Pi_0(c)\) es el propio \(c\), de modo que conserva el simplex de
historia y distingue el enrollamiento de \(\gamma\). La ablación refuta la
proyección desnuda; no refuta el producto fibrado empleado en el teorema.
\end{proof}

\begin{ablation}[Reescalamiento de la raíz visible sólo módulo tres]
\label{abl:bsd-suc-twist}
En la especialización \(c=1\) de
\eqref{eq:bsd-suc-diferencial-normal}, el ciclo alterado
\begin{equation}\label{eq:bsd-suc-twist4}
 z_K^{(4)}=4r+3e+b
\end{equation}
satisface
\begin{equation}\label{eq:bsd-suc-twist-defecto}
 z_{\mathrm{PT}}-z_K^{(4)}
 =\partial(-f)-3r.
\end{equation}
El defecto \(-3r\) desaparece módulo tres y reaparece módulo nueve. Por
consiguiente, \eqref{eq:bsd-suc-twist4} conserva una sombra parcial, pero no
la unidad basada ni \eqref{eq:bsd-suc-raiz}.
\end{ablation}

\begin{proof}
La igualdad \(\partial f=3e+b\) da inmediatamente
\eqref{eq:bsd-suc-twist-defecto}. Módulo tres, el segundo sumando es nulo;
módulo nueve es una clase raíz no divisible por una frontera del complemento.
Además, su coordenada raíz es cuatro, mientras
\eqref{eq:bsd-suc-normalizacion-raiz} exige uno. El ejemplo verifica la
necesidad del par acoplado y de la orientación; no introduce una excepción
al teorema que conserva ambos datos.
\end{proof}

\begin{ablation}[Pérdida de páginas del regulador]
\label{abl:bsd-suc-bockstein}
La matriz \(\operatorname{diag}(T^2,T)\) muestra que el primer Bockstein no
recupera el acarreo de orden superior. La matriz
\(\operatorname{diag}(2,T^2)\) tiene trivialización positiva
\(\tau_2=1\), pero
\(\tau_{\mathrm{reg}}=2=\operatorname{LT}_{T}(S)\). Por tanto, ni la
primera página ni el producto de páginas positivas sustituyen al regulador
completo \eqref{eq:bsd-suc-rboc}.
\end{ablation}

\begin{proof}
En el primer ejemplo, \(d=3\), \(r_0=2\) y
\(c_I^{\deg}=1\); la página uno no contiene ese último grado de memoria. En
el segundo, el único divisor singular es \(T^2\), pero la dirección regular
aporta el factor dos. Omitirla produciría uno en vez de dos. Ambos cálculos
son casos diagonales de \eqref{eq:bsd-suc-smith} y verifican la necesidad de
todas las piezas de \eqref{eq:bsd-suc-rboc}.
\end{proof}
\section*{Reconocimiento matemático final de Birch--Swinnerton--Dyer}

El desarrollo anterior concluye antes de emplear su formulación histórica
como criterio.

En el dominio declarado por el teorema
\ref{thm:bsd-suc-local-global}, la comparación de Smith identifica el
componente finito de la línea aritmética con la descomposición
Shafarevich--Tamagawa y demuestra la finitud de
\(\operatorname{Sha}(E)\). En consecuencia, su orden en la publicación
escalar siguiente es el cardinal del grupo finito identificado, no un símbolo
formal añadido a posteriori.

La traducción escalar de \eqref{eq:bsd-suc-linea-final} proporciona, con
\(r=\operatorname{rank}E(\mathbb Q)\),
\begin{equation}\label{eq:bsd-suc-rango-final}
 \operatorname{ord}_{s=1}L(E,s)=r
\end{equation}
y
\begin{equation}\label{eq:bsd-suc-formula-final}
 \frac{L^{(r)}(E,1)}{r!\,\Omega_E}
 =\frac{
   \lvert\operatorname{Sha}(E)\rvert\,\operatorname{Reg}(E)\prod_v c_v}
  {\lvert E(\mathbb Q)_{\mathrm{tors}}\rvert^2}.
\end{equation}
Esta es la formulación convencional de Birch--Swinnerton--Dyer. El orden,
el término principal y los factores aritméticos no se definen unos mediante
otros: son publicaciones coordinadas de la igualdad basada
\eqref{eq:bsd-suc-linea-final}, cuya raíz común fue fijada previamente por
\eqref{eq:bsd-suc-raiz}.


% El propietario conjunto anterior incorpora ya, por inclusiones internas y
% en su posición demostrativa propia, los refuerzos completos de RH,
% Navier--Stokes, Yang--Mills, Hodge y BSD.  No se vuelven a cargar aquí:
% repetirlos duplicaría teoremas y etiquetas sin añadir contenido.
\end{HMTScientificOwnerSubordinated}

El capítulo anterior construye simultáneamente las cinco operaciones desde el mismo
estado enriquecido del continuo. Aquí se despliega cada componente completo antes de
su aplicación, con fibra, extensiones, relaciones, números, orientación, acarreo,
memoria, frontera, ruta, multisección, supervivencia, terminales y lectores. La
separación expositiva de las cinco tuplas no las convierte en ramas, factores ni
objetos independientes: son cinco características intrínsecas del único continuo
generado y conservan literalmente el mismo estado y el mismo ledger. El constructor
común no las reúne después; las produce simultáneamente como estructura interna de
un solo objeto.

\section{Dominio, mapas y compatibilidades del objeto terminal común}

Para exponer cada característica sin separarla de las otras, se usa el índice
$T\in\{\mathrm{sp},\mathrm{sol},\mathrm{gau},\mathrm{coh},\mathrm{det}\}$,
pero el dato matemático completo sigue siendo el mismo continuo enriquecido. La
vista interna no se postula de nuevo: es exactamente la lectura tipada construida en
\eqref{eq:pdf2-operacion-infinita-tipificada},
\begin{equation}\label{eq:pdf2-g-t-desde-continuo}
 \mathfrak G_T^{\rm int}:=
 \operatorname{View}_T\!\left(
 \mathcal O_\infty(\mathcal C_{\rm cont}^{\rm disc})\right).
\end{equation}
Su presentación completa en el canal \(T\) es
\begin{equation}\label{eq:pdf2-g-t-completo}
 \mathfrak G_T^{\rm int}=
 \bigl(
 \begin{gathered}
 \{\mathcal F_q^{(T)}\}_q,
 \{\operatorname{Ext}_\gamma^{(T)}\}_\gamma,
 \operatorname{Rel}_T,\operatorname{Num}_T,\operatorname{or}_T,\\
 \operatorname{Carr}_T,\operatorname{Mem}_T,\partial_T,\gamma_T,
 \mathfrak m_T,\operatorname{Surv}_T,X_{\chi_T},\tau_T,
 \operatorname{Term}_T,\operatorname{Lect}_T
 \end{gathered}
 \bigr).
\end{equation}
No se admite sustituir este objeto por $X_{\chi_T}$, por su cardinal o por el nombre
de un problema clásico. Sea $\mathcal D_T^{\rm cl}$ el espacio de datos clásicos
del problema correspondiente. Éstos entran sólo mediante
\begin{equation}\label{eq:pdf2-mapa-preparacion-aplicacion}
 \operatorname{App}_{T,d_T}:\mathfrak G_T^{\rm int}\longrightarrow
 \mathfrak G_T(d_T),\qquad d_T\in\mathcal D_T^{\rm cl},
\end{equation}
donde $\mathfrak G_T(d_T)$ es, por definición, el par formado por la vista interna
completa y las coordenadas de interfaz que se despliegan abajo. Así ninguna fibra,
relación, orientación, carry, memoria, frontera, ruta, multisección, supervivencia
o terminal se pierde al introducir el dato clásico. La proyección interna
$\operatorname{pr}_{X,T}:\mathfrak G_T^{\rm int}\twoheadrightarrow X_{\chi_T}$
es sobreyectiva. Si
$\mathcal A_{T,d_T}:\mathfrak G_T(d_T)\to\mathcal M_T$, se define
\[
 \widetilde{\mathcal A}_{T,d_T}:=
 \mathcal A_{T,d_T}\circ\operatorname{App}_{T,d_T},
 \qquad
 \mathfrak M_T:=\operatorname{Im}\mathcal A_{T,d_T}.
\]

\section{Despliegue interno de las realizaciones correlativas anterior a su identificación con problemas clásicos}

La fórmula genérica anterior no sustituye la escritura de las cinco características.
En este volumen cada interfaz queda materializada sin convertir la vista interna en
un objeto o dominio independiente. En los cinco displays siguientes, la primera
coordenada es siempre $\mathfrak G_T^{\rm int}$; por ello las listas de coordenadas
específicas amplían el registro común y nunca lo sustituyen.
Usaremos
\[
 d_{\rm sp}=(R,\mathcal D_R^{\rm cof}),\quad
 d_{\rm sol}=(u_0,f,T),\quad
 d_{\rm gau}=(\Lambda_m,g),\quad
 d_{\rm coh}=(X,p),\quad
 d_{\rm det}=E/\mathbb Q.
\]
Esta línea fija el instante causal exacto en que aparecen los problemas clásicos:
después de \(\mathfrak G_T^{\rm int}\), nunca dentro de su generación HMT.

\subsection{Acción espectral–polar del constructor común y conservación de la fibra crítica}

La característica espectral--polar conserva el registro
\begin{equation}\label{eq:pdf2-g-sp-explicito}
 \mathfrak G_{\rm sp}(d_{\rm sp})=
 \bigl(
 \begin{gathered}
  \mathfrak G_{\rm sp}^{\rm int};\\
  \widehat X_\infty,\Gamma_9,\widehat R=(R,K),
  \operatorname{Carr}_R,S_0,M,\mathcal O_9,\\
  \Xi_R,P_R,\mathcal D_R^{\rm cof},
  \mathcal F_R,\operatorname{Ext}_R,\operatorname{Surv}_R,
  \tau_R,\mathcal P_R
 \end{gathered}
 \bigr).
\end{equation}
Aquí $\mathcal D_R^{\rm cof}$ es el dominio de funciones de prueba; $\Xi_R$ publica
los canales primo, gamma y polar en un mismo Hilbert; $P_R=P_R^*=P_R^2$ es la
proyección natural; $S_0,M,\mathcal O_9$ conservan la memoria y el terminal; y
$\widehat R=(R,K)$ conserva tanto la ventana como su profundidad. El objeto de Weil
es el complemento, no un signo añadido:
\begin{equation}\label{eq:pdf2-g-sp-complemento}
 B_{W,R}:=\Xi_R^*(I-P_R)\Xi_R.
\end{equation}
Los canales regulados, la naturalidad cofinal y el terminal pertenecen a la misma
característica; ninguno se reconstruye desde los ceros de $\zeta$. La estructura
interna aporta la involución de hojas, la expectativa de extensiones, el defecto y
la telescopía terminal. El propietario exacto anterior construye además la
interfaz de Weil: fija la contracción de hojas en
\eqref{eq:amp-rh-owner-q}, liga el complemento espectral--polar en
\eqref{eq:amp-rh-owner-binding}, conserva el terminal mediante
\eqref{eq:amp-rh-owner-telescopia} y prueba su extinción en
\eqref{eq:amp-rh-owner-terminal}. En la notación del constructor común, esa
coligación ya construida se publica como
\begin{equation}\label{eq:pdf2-g-sp-coligacion-dato}
 \Sigma_R^{\rm sp}:
 \mathcal H_{+,R}\oplus\mathcal X_R^{\rm sp}
 \longrightarrow
 \mathcal H_{-,R}\oplus\mathcal L_R^{\rm sp}
 \oplus\mathcal X_R^{\rm sp},
 \qquad
 (\Sigma_R^{\rm sp})^*\Sigma_R^{\rm sp}=I,
\end{equation}
con estado terminal $x_{R,N}=q^NV_R^Nx_{R,0}$, $q=5/9$, y salida
$K_R^{\rm sp}J_{+,R}=J_{-,R}$. Su identidad de energía es
\begin{equation}\label{eq:pdf2-g-sp-identidad-finita-dato}
 \|J_{+,R}f\|^2
 =\|J_{-,R}f\|^2+
   \sum_{n=0}^{N-1}\|\ell_{R,n}f\|^2+
   q^{2N}\|x_{R,0}f\|^2.
\end{equation}
La factorización \eqref{eq:amp-rh-owner-l} y su publicación cofinal
\eqref{eq:amp-rh-owner-publicacion} prueban estas identidades y el
entrelazador primo--gamma--polar antes de este ensamblaje. El capítulo de
Riemann las despliega y reconoce mediante el
teorema~\ref{thm:amp-rh-cierre-cofinal-propietario}; no construye a
posteriori una interfaz ausente ni la recibe de los ceros de $\zeta$.

\subsection{Balance solenoidal del constructor común y conservación del flujo}

Para cada corte $M$ y profundidad $j$, la imagen es
\begin{equation}\label{eq:pdf2-g-sol-explicito}
 \mathfrak G_{{\rm sol},M,j}(d_{\rm sol})=
 \bigl(
 \begin{gathered}
  \mathfrak G_{\rm sol}^{\rm int};\\
  \Xi_{M,j}^{\rm sol},C_{M\leftarrow N},M_{M,j}^+,M_{M,j}^-,
  D_{M,j},H_{M,j},E_9,J_9,P_9,Q_{\rm micro},\\
  \Gamma_{M,j}^{\rm term},\mathcal E_{M,j,T}^{\rm term},
  L_{M,j,T}^{\rm term},\ell_{M,j},\operatorname{Carr}_{M,j},\\
  \operatorname{Mem}_{M,j},\operatorname{Surv}_{M,j},\tau_{M,j}
 \end{gathered}
 \bigr).
\end{equation}
El campo Galerkin publicado es una coordenada de $\Xi_{M,j}^{\rm sol}$, pero el
estado incluye además el defecto $C_{M\leftarrow N}$, las memorias orientadas, la
counidad y su sección, la dinámica terminal y la pérdida observada. El Gram
\begin{equation}\label{eq:pdf2-g-sol-gram-estructural}
 U_{M,j;J,T}^{\rm term}
 :=(\mathbf G_{M,j;J,T}^{\rm term})^*
    \mathbf G_{M,j;J,T}^{\rm term}
\end{equation}
conserva en $\mathbf G^{\rm term}$ la coordenada terminal y todas las pérdidas
propagadas. La telescopía no borra ninguna de ellas.
El propietario exacto que precede a este ensamblaje construye el mapa causal
hidrodinámico completo: el lift PC--Fock de
\eqref{eq:ns-import-pcf-lift}, la coligación isométrica de
\eqref{eq:ns-import-cascade-isometry}, su telescopía de Stein
\eqref{eq:ns-import-stein-telescope}, el ledger
\eqref{eq:ns-import-ledger}, el carry financiado una sola vez
\eqref{eq:ns-import-carry-funded} y el retorno exacto
\eqref{eq:ns-import-exact-return}. En esas mismas coordenadas, las ondas de
scattering se publican como
\begin{equation}\label{eq:pdf2-g-sol-ondas-dato}
 a_{M,j}=(2\sqrt\nu)^{-1}A_M^{-1/2}\Lambda^sf_{M,j},
 \qquad
 b_{M,j}=\sqrt\nu A_M^{1/2}\Lambda^su_{M,j}-a_{M,j},
\end{equation}
y la transición lineal actúa sobre todas las coordenadas mixtas
$u^{\odot r}\otimes f^{\odot k}$. La raíz común y su transporte no se
postulan en este bloque: ya están construidos por
\eqref{eq:nsclose-owner-metric-definition}--
\eqref{eq:nsclose-owner-root-uniformity}. Su publicación en la carta común es
\begin{equation}\label{eq:pdf2-g-sol-raiz-dato}
 J_T^{\rm sol}:\mathcal D_T^{\rm NS}\longrightarrow\mathcal H_T^{\rm sol},
 \qquad
 \mathbf G_{M,0;J,T}^{\rm term}\Lambda_{M,T}J_T^{\rm sol}
 =I_{M,J}J_T^{\rm sol},
 \qquad I_{M,J}^*I_{M,J}=I.
\end{equation}
La conservación conjunta de terminal, disipación, advección, carry,
microestado, shell y memoria está probada por el telescopio propietario
\eqref{eq:nsclose-owner-telescope} y su ledger
\eqref{eq:nsclose-proprietary-ledger}. El presupuesto uniforme, independiente
de $M,J$, es el teorema ya establecido por
\eqref{eq:nsclose-proprietary-uniform-budget} y
\eqref{eq:nsclose-uniform-sobolev-precompactness}; aquí se publica en la
notación del constructor común como
\begin{equation}\label{eq:pdf2-g-sol-cota-dato}
 \sup_{M,J}
 \|\mathbf G_{M,0;J,T}^{\rm term}\Lambda_{M,T}J_T^{\rm sol}d\|^2
 \le C_T^{\rm sol}\|d\|_{\mathcal D_T^{\rm NS}}^2,
\end{equation}
con $C_T^{\rm sol}$ independiente de $M,J$. El capítulo posterior despliega
la realización de Navier--Stokes y reconoce su codominio clásico; no aporta
una premisa ausente ni selecciona la cota desde la conclusión.

\subsection{Incidencia proyectiva de calibre del constructor común y descenso al cociente}

La imagen gauge no es la marginal de enlaces. Su objeto completo por nivel es
\begin{equation}\label{eq:pdf2-g-gau-explicito}
 \mathfrak G_{{\rm gau},m}(d_{\rm gau})=
 \bigl(
 \begin{gathered}
  \mathfrak G_{\rm gau}^{\rm int};\\
  \Lambda_m,\widehat X_m,\widehat\mu_m^{\rm ov},
  \{\widehat\mu_{m,g}^{\rm YM},\nu_{m,g},K_{m,g}^{\rm phys}\}_{g>0},\\
  \mathcal G_m^{\rm full},\mathbb P_m^{\rm phys},
  \widehat J_m,\widehat\Gamma_m,
  \{E_{m,c}\}_c,\widehat H_m^{\rm ov},D_m^{\rm YM},\\
  \{\Theta_{\nu,m}\}_{\nu=1}^4,
  \{\mathfrak A_{+,\nu,m},\operatorname{ev}_{t}^{(\nu)},
      \mathcal V_{m,\nu}^{\rm OS}\}_{\nu=1}^4,\\
  \widehat\Omega_m^{(8)},P_{\Omega_m},W_c,
  \operatorname{Carr}_m,\operatorname{Mem}_m,
  \operatorname{Surv}_m,\tau_m,\mathcal P_m^{F^2}
 \end{gathered}
 \bigr).
\end{equation}
El grupo $\mathcal G_m^{\rm full}$ actúa simultáneamente en enlaces, marcos y
coordenadas de incidencia $q_c\in\mathbb F_3^6$. La proyección física es su promedio
de Haar. El generador reducido es
\begin{equation}\label{eq:pdf2-g-gau-generador-estructural}
 D_m^{\rm YM}
 =\left.3\kappa_{\rm av}\sum_c(I-E_{m,c})
  \right|_{\operatorname{Ran}\mathbb P_m^{\rm phys}},
\end{equation}
y $W_c$ pertenece a esa subálgebra. Las cuatro reflexiones actúan sobre una misma
medida, un mismo kernel y un mismo terminal.
El propietario exacto anterior construye la acción gauge de incidencia
\eqref{eq:amp-ym-accion-gauge-incidencia}, la proyección física de Haar
\eqref{eq:amp-ym-proyeccion-fisica-haar} y la compresión entrelazada
\eqref{eq:amp-ym-compresion-entrelazada}. De esa construcción proceden el
kernel completo equivariable $\widehat K_{m,t}^{\rm ov}$, su descenso al
cociente físico $K_{m,t}^{\rm phys}$, los cuatro operadores de reflexión y
la acción $U_m(g)$ de $E(4)$ en el núcleo cilíndrico suave. Este bloque
ensambla, sin introducirlos como datos, el generador, el vacío, el testigo
de curvatura y las formas OS:
\begin{equation}\label{eq:pdf2-g-gau-entrelazamiento-dato}
 D_{m+1}^{\rm YM}\widehat J_m=\widehat J_mD_m^{\rm YM},
 \quad
 \mathbb P_{m+1}^{\rm phys}\widehat J_m
  =\widehat J_m\mathbb P_m^{\rm phys},
 \quad
 \widehat J_mW_{m,c}=W_{m+1,c}.
\end{equation}
La primera y la tercera identidades son la publicación, en la notación
común, de \eqref{eq:amp-ym-compresion-entrelazada} y
\eqref{eq:amp-ym-testigo-fisico-gap}; el semigrupo, el vacío simple y la
brecha exacta ya se obtienen en
\eqref{eq:amp-ym-owner-semigrupo-vacio}--
\eqref{eq:amp-ym-owner-brecha-exacta}. Aquí
$\nu_{m,g}:=(\pi_m^{\rm phys})_*\widehat\mu_{m,g}^{\rm YM}$ y
$K_{m,g}^{\rm phys}(t)=e^{-tD_{m,g}^{\rm YM}}$ en el cociente físico.
La completitud gauge incluye además la desintegración por cada familia de
hiperplanos euclídeos. Si $A_j\in\mathfrak A_{+,\nu,m}$ y
$t_1\le\cdots\le t_n$, la marginal de la misma medida física satisface
\begin{equation}\label{eq:pdf2-g-gau-transfer-euclidean-dato}
\begin{aligned}
 &\int\prod_{j=1}^n
   (\tau_{t_je_\nu}A_j)(x)\,d\widehat\mu_{m,g}^{\rm YM}(x)\\
 &\quad=
 \int\prod_{j=1}^nA_j(y_j)\,
  \nu_{m,g}(dy_1)
  \prod_{j=1}^{n-1}K_{m,g}^{\rm phys}
       (t_{j+1}-t_j;y_j,dy_{j+1})\\
 &\quad=
 \left\langle\mathbf1,
  M_{A_1}e^{-(t_2-t_1)D_{m,g}^{\rm YM}}M_{A_2}\cdots
  e^{-(t_n-t_{n-1})D_{m,g}^{\rm YM}}M_{A_n}\mathbf1
 \right\rangle.
\end{aligned}
\end{equation}
La reconstrucción anterior prueba la identidad en las cuatro direcciones
mediante \eqref{eq:amp-ym-cuatro-os} y construye un único Hamiltoniano
entrelazado en \eqref{eq:amp-ym-hamiltoniano-comun}. En consecuencia, la
traslación OS se publica aquí como
\begin{equation}\label{eq:pdf2-g-gau-os-generator-dato}
 \mathcal V_{m,\nu}^{\rm OS}T_{m,\nu}^{\rm OS}(s)
 (\mathcal V_{m,\nu}^{\rm OS})^{-1}
 =e^{-sD_{m,g}^{\rm YM}}.
\end{equation}
Así, el generador de expectativas y el Hamiltoniano de traslaciones no se
identifican por nombre: son dos realizaciones ya construidas de la misma
desintegración, cuya forma límite y cuya publicación de curvatura constan en
\eqref{eq:amp-ym-forma-limite} y \eqref{eq:amp-ym-f2}.

\subsection{Coherencia cohomológica del constructor común y emergencia de la clase global}

Para una variedad proyectiva lisa compleja $X$ y una codimensión $p$, la imagen
cohomológica conserva
\begin{equation}\label{eq:pdf2-g-coh-explicito}
 \mathfrak G_{\rm coh}(d_{\rm coh})=
 \bigl(
 \begin{gathered}
  \mathfrak G_{\rm coh}^{\rm int};\\
  \{\mathscr S_N(X,p),\rho_{N+1,N}\}_N,
  \{\mathscr H_N^p(X),q_{N+1,N}\}_N,
  \{e_N\}_N,
  \{\iota_{N+1,N}\}_N,\\
  \{\mathcal G_{N+1}^{\rm coh},\mathcal R_{N+1}^{\rm coh},
       D_{N+1},a_{N+1},b_{N+1}\}_N,\\
  \mathcal F_{X,p},\operatorname{Ext}_{X,p},X_\chi(X,p),
  \operatorname{Rel}_{X,p},\operatorname{Carr}_{X,p},
  \operatorname{Mem}_{X,p},\partial_{X,p},\gamma_{X,p},\\
  \Sigma_{X,p},\operatorname{Surv}_{X,p},\tau_{X,p},
  R_{\rm Hdg},\operatorname{ch}^{\rm loc}_p,\operatorname{cl}_{X,p}
 \end{gathered}
 \bigr).
\end{equation}
Los cuadrados
\begin{equation}\label{eq:pdf2-g-coh-cuadrados}
 e_N\rho_{N+1,N}=q_{N+1,N}e_{N+1}
\end{equation}
ligan observación, extensión y refinamiento. El propietario exacto anterior
construye, y no recibe como dato, la historia basada con sus correcciones
relativas, carry, memoria, frontera y terminal. En particular,
\eqref{eq:amp-hodge-mapas-relativos}--
\eqref{eq:amp-hodge-factorizacion-relativa} construyen los mapas del
incremento y \eqref{eq:amp-hodge-seccion-relativa} construye su inverso
derecho. En la notación del constructor común, ese mapa relativo se publica
como
\begin{equation}\label{eq:pdf2-g-coh-seccion-dato}
 s_{N+1}^{\rm rel}:\ker q_{N+1,N}\longrightarrow\ker\rho_{N+1,N},
 \qquad
 e_{N+1}^{\rm rel}s_{N+1}^{\rm rel}=I,
\end{equation}
natural bajo refinamiento por
\eqref{eq:amp-hodge-naturalidad-seccion-relativa} y compatible con carry y
terminal. La sección inicial tampoco se supone: se construye en
\eqref{eq:amp-hodge-seccion-base}. Son estos mapas ya demostrados, y no una
clase final elegida, los que el bloque ensambla y el capítulo posterior
reconoce en el codominio Hodge.

La completitud se apoya en acciones e identidades construidas por el mismo
propietario. La degeneración unitaria, la presentación relativa y el lector
se ensamblan aquí mediante
\begin{equation}\label{eq:pdf2-g-coh-datos-operativos}
\begin{gathered}
 \iota_{N+1,N}:\mathscr S_N(X,p)\longrightarrow\mathscr S_{N+1}(X,p),
 \qquad \rho_{N+1,N}\iota_{N+1,N}=I,\\
 \mathscr A_{N+1}^{\rm rel}
  =\mathbb Q^{(\mathcal G_{N+1}^{\rm coh})}/\operatorname{im}D_{N+1},
 \qquad e_{N+1}^{\rm rel}a_{N+1}=b_{N+1},
 \qquad b_{N+1}D_{N+1}=0.
\end{gathered}
\end{equation}
El conjunto $\mathcal G_{N+1}^{\rm coh}$ enumera exhaustivamente todos los
átomos nuevos y una etiqueta terminal por componente. La forma normal y el
inverso derecho se construyen por reducción basada en
\eqref{eq:amp-hodge-sigma-coordenada}; su compatibilidad con refinamientos se
prueba en \eqref{eq:amp-hodge-naturalidad-seccion-relativa}. Por ello, en el
orden por profundidad, soporte, hoja, ruta, carry, memoria y terminal, cada
átomo de observación tiene una única etiqueta máxima y
\begin{equation}\label{eq:pdf2-g-coh-tabla-dato}
 b_{N+1}(c_\mu)=\bar c_\mu+
   \sum_{\nu<\mu}B_{\nu\mu}\bar c_\nu,
 \qquad B_{\nu\mu}\in\mathbb Q.
\end{equation}
La triangularidad es una consecuencia de esa construcción basada, anterior a
cualquier clase $\alpha$; no es una hipótesis ni una tabla suministrada desde
el problema terminal.

La completitud y el lector del mismo objeto están tipados por
\begin{equation}\label{eq:pdf2-g-coh-limite-lector-dato}
 X_\chi(X,p)=\varprojlim_N\mathscr S_N(X,p),
 \qquad
 q_N\operatorname{ev}^{\rm Hdg}_{\infty,X,p}
  =e_N\operatorname{pr}_N,
 \qquad
 \operatorname{cl}_{X,p}\operatorname{ch}^{\rm loc}_pR_{\rm Hdg}
  =\operatorname{ev}^{\rm Hdg}_{\infty,X,p}.
\end{equation}
Las proyecciones $q_N$, incluida la coordenada terminal, separan puntos del
codominio de evaluación. La extensión operativa
\eqref{eq:amp-hodge-extension-operativa}--\eqref{eq:amp-hodge-ext}, el paso
al límite \eqref{eq:amp-hodge-xi-limite} y la publicación final
\eqref{eq:amp-hodge-publicacion-final} prueban ya que toda clase Hodge
arbitraria admite una historia compatible. El capítulo de aplicación
despliega y reconoce esa publicación; no suministra la existencia como una
premisa posterior.

\subsection{Línea determinantal del constructor común y evaluación aritmética terminal}

Para una curva elíptica $E/\mathbb Q$, la imagen determinantal conserva
\begin{equation}\label{eq:pdf2-g-det-explicito}
 \mathfrak G_{\rm det}(d_{\rm det})=
 \bigl(
 \begin{gathered}
  \mathfrak G_{\rm det}^{\rm int};\\
  \mathcal T_m^{\rm surv},G_{m,n},\Delta_{\rm AW},\varepsilon_{\rm AW},
  P_E^{\rm gen},1_E,P_E^K,P_E^{\rm PT},\\
  z_K,z_{\rm PT},X^{\rm orth},p_{\rm root},
  \mathcal C_E^K,\mathcal C_E^{\rm Iw},\mathcal C_E^{\rm PT},\\
  \ell_E,\mathcal C_E^{\rm loc,red},S_E(T),
  \operatorname{Carr}_E,\operatorname{Mem}_E,
  \operatorname{Surv}_E,\tau_E,\mathcal L_E
 \end{gathered}
 \bigr).
\end{equation}
La unidad $1_E$ se copublica en las ramas Kato y Poitou--Tate antes de la dualidad y
de la raíz. El producto fibrado conserva la sombra modular y la historia profunda.
El cono de localización, las páginas Bockstein y la línea determinante pertenecen a
una misma multisección basada; no son invariantes elegidos por separado después de
conocer el rango o el término principal.
El propietario exacto anterior construye la descomposición basada
$\Theta_E$, la homotopía del sumando finita--singular, el comparador completo
Kato--FERS, los mapas de páginas y gluing de Poitou--Tate y los mapas de
cociclos local--global. La cadena comienza en la línea determinante
\eqref{eq:amp-bsd-linea-determinante}, construye el mapa de cadenas
\eqref{eq:amp-bsd-phi-kato-fers}, el comparador de Poitou--Tate
\eqref{eq:amp-bsd-comparador-pt} y el triángulo local
\eqref{eq:amp-bsd-triangulo-localizacion}. Sus identidades operativas se
publican aquí como
\begin{equation}\label{eq:pdf2-g-det-identidades-dato}
\begin{gathered}
 \Phi_{K\leftrightarrow{\rm FERS}}\Psi_{K\leftrightarrow{\rm FERS}}=I,
 \qquad
 \Psi_{\rm PT}^{-1}\Psi_{\rm PT}=I,
 \qquad
 \Psi_{\rm PT}\Psi_{\rm PT}^{-1}=I,\\
 \ker\partial_{\rm loc}=\operatorname{im}\iota_{\Sha},
 \qquad
 \pi_\Phi\operatorname{Lift}_\Phi=I.
\end{gathered}
\end{equation}
Estas identidades son consecuencias del paquete exacto
\eqref{eq:amp-bsd-paquete-identidades}--
\eqref{eq:amp-bsd-paquete-composicion}. El comparador total se construye en
\eqref{eq:amp-bsd-comparador-total} y conserva la base común desde el lado
analítico hasta el aritmético; rango y término principal se publican después
en \eqref{eq:amp-bsd-igualdad-grados}--
\eqref{eq:amp-bsd-termino-principal}. El capítulo BSD expande y reconoce esa
composición sin seleccionar retrospectivamente ninguna de sus salidas.

\section{\(5\) lectores terminales no permutables del continuo genealógico}

El mapa de publicación $\mathcal P_T$ es el lector matemático que
traduce la estructura al codominio clásico. La cadena propia de este volumen es
\begin{equation}\label{eq:pdf2-cinco-cadenas}
 \mathfrak G_T^{\rm int}\xrightarrow{\operatorname{App}_{T,d_T}}
 \mathfrak G_T(d_T)\xrightarrow{\mathcal A_{T,d_T}}\mathfrak M_T
 \xrightarrow{\mathcal P_T}\mathcal C_T.
\end{equation}
Las correspondencias quedan fijadas una a una:
\[
 \mathcal C_{\rm sp}=\mathrm{RH},\quad
 \mathcal C_{\rm sol}=\mathrm{NS},\quad
 \mathcal C_{\rm gau}=\mathrm{YM},\quad
 \mathcal C_{\rm coh}=\mathrm{Hodge},\quad
 \mathcal C_{\rm det}=\mathrm{BSD}.
\]
En el ensamblado activo, los propietarios exactos que preceden a esta sección
construyen también las cinco publicaciones finales: RH en el
teorema~\ref{thm:amp-rh-cierre-cofinal-propietario}, Navier--Stokes en los
teoremas~\ref{thm:nsclose-proprietary-full} y
\ref{thm:nsclose-global-smooth}, Yang--Mills en el
teorema~\ref{thm:amp-ym-terminacion}, Hodge en el
teorema~\ref{thm:amp-hodge-completitud-coinductiva} y BSD en el
teorema~\ref{thm:amp-bsd-comparadores-construidos}. Esta sección registra su
ensamblaje correlativo en el único continuo genealógico; los capítulos
posteriores despliegan las cinco realizaciones y su reconocimiento clásico,
sin aportar premisas ausentes ni utilizar las conclusiones como entradas.

\begin{falsifier}[Ablación del dato de partida]
Si una demostración comienza en $X_{\chi_T}$, convierte una coordenada de
$\mathfrak G_T^{\rm int}$ en un selector opcional o sustituye una fibra por su cardinal, no ha
abreviado el mismo objeto: lo ha reemplazado. La conclusión no es transferible al
sustituto.
\end{falsifier}

\section{Factorización de los lectores terminales mediante naturalidad, refinamiento y paso al límite}

El punto de partida material es el único constructor simultáneo que produce un
solo continuo enriquecido con cinco operaciones intrínsecas correlativas. Las
notaciones
\(\mathfrak G_{\rm sp}\), \(\mathfrak G_{\rm sol}\),
\(\mathfrak G_{\rm gau}\), \(\mathfrak G_{\rm coh}\) y
\(\mathfrak G_{\rm det}\) designan cinco vistas completas de ese mismo registro,
reproducido con sus fibras,
relaciones, números, orientación, acarreo, memoria, frontera, ruta, multisección,
supervivencia, terminales y lectores. No se toman como cinco ramas elegidas ni como
cinco datos independientes: se demuestra su compatibilidad sobre el mismo ledger,
su naturalidad bajo refinamiento, su no-vacuidad conjunta y su paso al límite, y
después se aplican, respectivamente, a los cinco problemas. Ninguna prueba de este
volumen remite a otra fuente externa para completar una definición, una identidad o un paso
lógico.

Cada cierre se expone de forma autocontenida: las definiciones, demostraciones y
compatibilidades aparecen antes de su conclusión clásica. En cada resolución
se exhiben cuatro capas distintas:
\begin{enumerate}
  \item identidad exacta en cada nivel finito;
  \item compatibilidad de refinamientos y preservación de la historia enriquecida;
  \item uniformidad y paso al límite, con su terminal conservado;
  \item reconocimiento de la conclusión clásica en el codominio publicado.
\end{enumerate}

La composición rectora de las cinco clausuras conserva primero la acción única
correlativa, forma el continuo sólo al pasar al límite y después abre una vista para
cada aplicación:
\[
\begin{aligned}
 \mathcal E_k&\xrightarrow{\ \mathcal O_k\ }\mathcal Z_k
 \xrightarrow{\ \varprojlim_k\ }
 \bigl(\mathcal C_{\rm cont}^{\rm disc},\mathcal O_\infty\bigr),\\
 \bigl(\mathcal C_{\rm cont}^{\rm disc},\mathcal O_\infty\bigr)
 &\xrightarrow{\ \operatorname{View}_T\ }\mathfrak G_T^{\rm int}
 \xrightarrow{\ \operatorname{App}_{T,d_T}\ }\mathfrak G_T(d_T),\\
 \mathfrak G_T(d_T)&
 \xrightarrow{\widetilde{\mathcal A}_{T,d_T}}\mathfrak M_T
 \xrightarrow{\mathcal P_T}\mathcal C_T,
 \qquad T\in\{\mathrm{sp,sol,gau,coh,det}\}.
\end{aligned}
\]
Aquí $\mathcal O_k$ actúa una sola vez sobre el estado completo y conserva juntas
las cinco características. $\operatorname{View}_T$ no restringe el dominio ni
selecciona historias: sólo expone las coordenadas de una aplicación dentro del valor
completo de $\mathcal O_\infty$. El dato clásico $d_T$ entra sólo en
$\operatorname{App}_{T,d_T}$, después del continuo, y no participa en la generación
de la característica intrínseca. La recta real es asimismo una salida posterior de
$\operatorname{ev}_{\mathbb R}$, no una premisa de la cadena.

Cada capítulo distingue el dominio de partida, la construcción, el certificado,
la conclusión demostrada y sus falsadores exactos. Los propietarios precedentes
construyen las cinco composiciones en sus dominios respectivos; los capítulos que
siguen despliegan esas construcciones y efectúan el reconocimiento final. Ninguna
flecha se reemplaza por una remisión, una conclusión nominal o una hipótesis clásica.

\begingroup
\HMTDemoteHeadings
\chapter{Una misma frontera: flujo interior y tensión gauge}
\label{chap:hoja-comun-ns-ym}
\index{Navier--Stokes}\index{Yang--Mills}
\index{operador de frontera}\index{Stokes}

La hidrodinámica y Yang--Mills no se reúnen comparando desde fuera dos
ecuaciones célebres. Proceden de dos representaciones del mismo operador de
frontera. La orientación que el lector interior interpreta como circulación
se convierte, bajo el lector reflejado, en holonomía; la memoria simétrica
que controla el flujo se convierte en kernel positivo de transferencia.

\section{El estado común de frontera}

Sea \(K_m\) un complejo finito orientado, \(B_m\) su álgebra de datos de
frontera y \(d_m\) su coborde. El estado enriquecido es
\[
 s_m=(b_m,U_m,\theta_m,c_m,\Omega_m),
\]
donde \(U_m\) transporta, \(\theta_m^2=I\) invierte la orientación,
\(c_m\) conserva memoria y \(\Omega_m\) es una \(1\)-cocadena de conexión.
Los proyectores de paridad son
\[
 E_\pm=\frac12(I\pm\theta_m).
\]
La parte \((U_m+U_m^\dagger)/2\) es simétrica; sólo es un proyector si además
se demuestra su idempotencia.

Sobre el complejo de cocadenas se construyen dos lectores lineales:
\[
 \rho_{{\rm int},m}:C^0(K_m)\longrightarrow
 C^1(K_m)\oplus C^0(K_m),
\qquad
 \rho_{{\rm int},m}(v)=(d_mv,d_m^\ast d_mv),
\]
\[
 \rho_{{\rm ref},m}:C^1(K_m)\longrightarrow
 C^0(K_m)\oplus C^1(K_m),
\qquad
 \rho_{{\rm ref},m}(a)=(d_m^\ast a,d_md_m^\ast a).
\]
El primero transporta un potencial de vértices a su flujo de aristas y a su
respuesta interior; el segundo transporta una cocadena de aristas a su
respuesta de borde y a su operador reflejado. Los codominios son distintos,
pero ambos se construyen con el mismo coborde. Circulación, vorticidad,
holonomía y curvatura aparecen después mediante la identidad de Stokes; no
se presuponen dentro de los símbolos \(\rho\).

\section{El operador mínimo y su espectro}

En el ciclo orientado de tres vértices, sea
\[
 d_0=
 \begin{pmatrix}
 -1&1&0\\
 0&-1&1\\
 1&0&-1
 \end{pmatrix}.
\]
Los Laplacianos de Hodge de vértices y aristas son
\[
 \mathcal L_{\rm H}^{(0)}=d_0^\ast d_0,
 \qquad
 \mathcal L_{\rm H}^{(1)}=d_0d_0^\ast,
\]
y un cálculo directo da
\[
 \mathcal L_{\rm H}^{(0)}=\mathcal L_{\rm H}^{(1)}=
 \begin{pmatrix}
 2&-1&-1\\
 -1&2&-1\\
 -1&-1&2
 \end{pmatrix}
 =3I-J.
\]
En esta residencia, los superíndices \((0)\) y \((1)\) designan los
Laplacianos de Hodge sobre vértices y aristas. No designan los operadores de
elevación \(L_0,L_1\) de la cadena \(w_6\to w_{30}\). Para
\(\mathcal L_{\rm cyc}^{\rm TPK}:=2I-C-C^{-1}=3I-J\), la realización del
operador cíclico del TPK se tipa mediante
\[
 \mathfrak R_{\rm hyd}(\mathcal L_{\rm cyc}^{\rm TPK})
 =\mathcal L_{\rm H}^{(0)},\qquad
 \mathfrak R_{\rm YM}(\mathcal L_{\rm cyc}^{\rm TPK})
 =\mathcal L_{\rm H}^{(1)}.
\]

\begin{theorem}[Operador común de frontera]
\label{thm:operador-comun-frontera-ns-ym}
En el ciclo de tres celdas,
\[
 \operatorname{Spec}(\mathcal L_{\rm H}^{(0)})
 =\operatorname{Spec}(\mathcal L_{\rm H}^{(1)})
 =\{0,3,3\}.
\]
Después de retirar el modo constante ---media cero en el lector interior y
calibre en el lector de conexión--- ambos sectores poseen gap exacto \(3\).
En todo complejo finito, \(d^\ast d\) y \(dd^\ast\) tienen el mismo espectro
no nulo, con multiplicidades.
\end{theorem}

\begin{proof}
El vector \((1,1,1)\) genera el núcleo de \(3I-J\). En su complemento
ortogonal, \(J=0\), de modo que el operador es \(3I\). Para el enunciado
general, una descomposición singular \(d=U\Sigma V^\ast\) produce
\[
 d^\ast d=V\Sigma^\ast\Sigma V^\ast,
 \qquad
 dd^\ast=U\Sigma\Sigma^\ast U^\ast.
\]
Los cuadrados de los valores singulares no nulos, incluidas sus
multiplicidades, coinciden.
\end{proof}

\section{Lectura hidrodinámica}

Definimos el paso implícito
\[
 S_{\rm NS}=(I+\mathcal L_{\rm H}^{(0)})^{-1}.
\]
En el ciclo mínimo,
\[
 S_{\rm NS}=\frac14(I+J).
\]

\begin{theorem}[Contracción global del flujo finito]
\label{thm:contraccion-flujo-finito}
Para todo \(v_0\) de media cero, la iteración
\(v_{k+1}=S_{\rm NS}v_k\) existe y es única para todo \(k\geq0\), y
\[
 v_k=4^{-k}v_0,
 \qquad
 \|v_k\|\leq4^{-k}\|v_0\|,
 \qquad
 \|v_k\|^2=16^{-k}\|v_0\|^2.
\]
\end{theorem}

\begin{proof}
En media cero se tiene \(Jv=0\), luego
\(S_{\rm NS}v=v/4\). La fórmula se obtiene por inducción.
\end{proof}

\Needspace{4\baselineskip}
Una no linealidad orientada \(N_m(v,c)\) conserva la energía cuando
\[
 \langle v,N_m(v,c)\rangle=0.
\]
La advección redistribuye entonces la energía y el operador positivo controla
su propagación. La memoria \(c_m\) no se absorbe en una viscosidad elegida:
forma parte del estado que debe transportarse entre escalas.

\begingroup
\setlength{\parskip}{0pt}
\titlespacing*{\section}{0pt}{8pt plus 1pt minus .5pt}{2.5pt}
\setlength{\abovedisplayskip}{.25\baselineskip}
\setlength{\belowdisplayskip}{.25\baselineskip}
\setlength{\abovedisplayshortskip}{.12\baselineskip}
\setlength{\belowdisplayshortskip}{.18\baselineskip}
\section{Lectura Yang--Mills}

Una \(1\)-cocadena \(a\) se interpreta como potencial de conexión y posee
energía cuadrática
\[
 Q_{\rm YM}(a)=\langle a,\mathcal L_{\rm H}^{(1)}a\rangle.
\]

\begin{theorem}[Coercividad y decaimiento gauge finitos]
\label{thm:coercividad-gap-gauge-finito}
En el cociente por el modo gauge constante,
\[
 Q_{\rm YM}(a)\geq3\|a\|^2.
\]
El semigrupo \(K_t=e^{-t\mathcal L_{\rm H}^{(1)}}\) satisface
\[
 \|K_ta\|\leq e^{-3t}\|a\|
\]
en el sector físico del ciclo mínimo.
\end{theorem}

\begin{proof}
Sobre el complemento del núcleo, el
\cref{thm:operador-comun-frontera-ns-ym} da
\(\mathcal L_{\rm H}^{(1)}=3I\). Las dos fórmulas
siguen por cálculo funcional.
\end{proof}

Para una conexión no abeliana,
\[
 F_\Omega=d\Omega+\Omega\smile\Omega,
\qquad
 W_\gamma(\Omega)=\operatorname{Tr}\operatorname{Hol}_\gamma(\Omega).
\]
La positividad por reflexión requiere una factorización sobre la mitad
positiva,
\[
 \langle\theta f,Kf\rangle=\|Cf\|^2\geq0;
\]
la positividad puntual de la matriz \(K\) no la sustituye.

\section[Identidad de Stokes: circulación--vorticidad y holonomía--curvatura]{Identidad discreta de Stokes y \(2\) realizaciones: circulación--vorticidad y holonomía--curvatura}

Para una \(1\)-cocadena \(\Omega\) y una \(2\)-cadena \(K\),
\[
 \sum_{e\subset\partial K}\Omega(e)
 =
 \sum_{f\in K^{(2)}}(d\Omega)(f).
\]
Las aristas interiores se cancelan por orientación. El lector
hidrodinámico interpreta el lado izquierdo como circulación y el derecho
como vorticidad acumulada. El lector gauge interpreta el primero como
holonomía linealizada y el segundo como curvatura.

\begin{theorem}[Teorema de la hoja común de frontera]
\label{thm:hoja-comun-frontera}
La descomposición \(B_m=E_-B_m\oplus E_+B_m\) posee dos representaciones
tipadas:
\[
\boxed{
\begin{array}{ccl}
E_-B_m&\longmapsto&
\text{circulación}\ /\ \text{holonomía},\\[2mm]
E_+B_m&\longmapsto&
\text{control disipativo}\ /\ \text{kernel reflejado}.
\end{array}}
\]
La igualdad de Stokes y la coincidencia del espectro no nulo pertenecen al
objeto común. La contracción hidrodinámica y el gap gauge son consecuencias
distintas en sus respectivos codominios.
\end{theorem}

\begin{proof}
Los proyectores \(E_\pm\) son complementarios porque
\(\theta_m^2=I\). La identidad de Stokes proporciona el entrelazamiento del
sector orientado; el \cref{thm:operador-comun-frontera-ns-ym} proporciona el
espectro común. Los
\cref{thm:contraccion-flujo-finito,thm:coercividad-gap-gauge-finito}
calculan las dos realizaciones.
\end{proof}

El teorema es un cierre finito HMT: no identifica una PDE con una teoría
gauge ni confunde el gap \(3\) con una masa física. Explica por construcción
por qué regularidad interior y separación de frontera son dos funciones de
la misma memoria orientada. Cualquier paso continuo posterior debe conservar
el operador, la norma y el espectro que hacen exacta esta identidad.
\endgroup

\endgroup

\HMTFlushCurrentChapter
\chapter{Factorización positiva de la forma de Guinand--Weil y localización de los ceros no triviales de \texorpdfstring{\(\zeta\)}{zeta}}
\begingroup
\renewcommand{\chapter}[2][]{}%
\chapter{Imagen espectral--polar y reducción exacta de Riemann--Weil}

\section{Punto de partida: la estructura espectral--polar completa}

La entrada de la resolución no es la forma de Weil aislada ni una lista de
ceros. Es la estructura espectral--polar completa producida por el constructor
común y acreditada por el propietario exacto que precede a este capítulo; no es
un dato externo. La cadena causal propia de este capítulo es
\begin{equation}\label{eq:pdf2-rh-cadena-procedencia}
 \mathfrak G_{\rm sp}
 \xrightarrow{\ \operatorname{pr}_{X,{\rm sp}}\ }X_{\chi_{\rm sp}}
 \xrightarrow{\ \mathcal A_{\rm sp}\ }\mathfrak M_{\rm sp}
 \xrightarrow{\ \mathcal P_{\rm sp}\ }\mathcal C_{\rm GW}.
\end{equation}
El codominio clásico aparece sólo en la última flecha; no selecciona
retrospectivamente el estado, las cartas, la proyección ni el terminal.

La estructura de partida contiene el objeto de trece coordenadas
\begin{equation}\label{eq:pdf2-rh-res-sp-completa}
\begin{aligned}
 \mathfrak G_{\rm sp}^{(13)}=\bigl(&
 \mathcal F_{\rm sp},
 \operatorname{Ext}_{\Gamma,{\rm sp}},
 \operatorname{Rel}_{\rm sp},
 \mathcal N_{\rm sp},
 \operatorname{or}_{\rm sp},
 \operatorname{Carr}_{\rm sp},
 \operatorname{Mem}_{\rm sp},\\
 &\partial_{\rm sp},
 \gamma_{\rm sp},
 \Sigma_{\rm sp}^{\rm mult},
 \operatorname{Surv}_{\rm sp},
 \tau_{\rm sp}^{\rm term},
 \mathcal L_{\rm sp}\bigr).
\end{aligned}
\end{equation}
Cada coordenada interviene materialmente:
\begin{enumerate}
 \item \(\mathcal F_{\rm sp}\) conserva las dos hojas, su orientación, el
 par residuo--cociente y la firma de la publicación doble.
 \item \(\operatorname{Ext}_{\Gamma,{\rm sp}}\) transporta esa fibra por la
 clase dirigida de cortes y niveles nonádicos.
    \item \(\operatorname{Rel}_{\rm sp}\) contiene el cociclo de acarreo, la
    ortogonalidad \(P=P^*=P^2\), las dos publicaciones y el complemento
    estructural \(B_{\rm sp}^{\rm str}:=\Xi^*(I-P)\Xi\). La identificación
    de este complemento con la forma de Guinand--Weil es una flecha distinta
    y se audita por separado en este capítulo.
 \item \(\mathcal N_{\rm sp}\) registra el nivel \(m\), la ventana \(R\),
 el cociente \(q=5/9\), los momentos de ambas hojas y los índices de los
 canales primo, gamma y polar.
 \item \(\operatorname{or}_{\rm sp}=\widehat R=(R,K)\) mantiene distintas
 las dos cartas orientadas.
 \item \(\operatorname{Carr}_{\rm sp}=\operatorname{Carr}_R\) es el
 acarreo operatorial y no se anula al publicar una traslación.
 \item \(\operatorname{Mem}_{\rm sp}=(S_0,M,\mathcal O_9)\) conserva la
 identidad finita de Stein y su terminal.
 \item \(\partial_{\rm sp}=(I-P)\Xi\) es la frontera que no aparece en la
 compresión \(P\Xi\) y cuya energía define el complemento estructural. Su
 identificación con la forma de Weil no se incorpora a esta definición.
 \item \(\gamma_{\rm sp}\) es la ruta cofinal de ventanas, con un único
 regulador para todos los canales.
 \item \(\Sigma_{\rm sp}^{\rm mult}\) es la multisección de las dos
 publicaciones y sus componentes primo--gamma--polar; no es la firma de un
 único estado.
 \item \(\operatorname{Surv}_{\rm sp}\) conserva la familia bajo
 \(m\mapsto m+1\) y \(\mathcal D_R\hookrightarrow\mathcal D_{R'}\).
 \item \(\tau_{\rm sp}^{\rm term}\) conserva el terminal
 \(M^{*N}S_0M^N\), que en la carta espejo se publica mediante \(E_R\).
    \item \(\mathcal L_{\rm sp}=(\Xi,P,B_{\rm sp}^{\rm str})\) es el lector final; se aplica
 después de antecedente, cartas, memoria, ruta y terminal.
\end{enumerate}

El ensamblador y el publicador son
\begin{equation}\label{eq:pdf2-rh-ensamblador-sp}
\begin{aligned}
 \mathcal A_{\rm sp}(\operatorname{pr}_{X,{\rm sp}}\widehat x)
  ={}&(\widehat X_\infty,\Gamma_9,\widehat R,
      \operatorname{Carr}_R,S_0,M,\mathcal O_9,
      \Xi,P,\mathcal D_R^{\rm cof}),\\
 \mathfrak M_{\rm sp}={}&\operatorname{Im}\mathcal A_{\rm sp},\\
 \mathcal P_{\rm sp}(\mathfrak M_{\rm sp})
  ={}&(B_{\rm sp}^{\rm str},\Xi^*(I-P)\Xi,
       \mathscr I_{\rm GW},\Delta_{\rm bind},
       \text{familia cofinal primo--gamma--polar}),\\
 \Delta_{\rm bind}:={}&B_{\rm sp}^{\rm str}-\mathscr I_{\rm GW}.
\end{aligned}
\end{equation}
La memoria finita del objeto de partida satisface
\begin{equation}\label{eq:pdf2-rh-stein-finito}
 S_0=\sum_{r=0}^{N-1}M^{*r}\mathcal O_9^*\mathcal O_9M^r
     +M^{*N}S_0M^N.
\end{equation}
El último sumando permanece visible hasta que se demuestra su extinción. Por
tanto, reemplazar \(\mathfrak G_{\rm sp}\) por la sombra publicada
\(B_{\rm sp}^{\rm str}\) cambia el objeto de partida; no es una abreviación
del mismo teorema. El residual \(\Delta_{\rm bind}\) se conserva hasta que la
coligación source--first demuestra materialmente su anulación en
\eqref{eq:pdf2-rh-binding-source-first}.

\section{Dominio cofinal y forma objetivo}

Para $R>0$ se fija
\[
 \mathcal D_R=C_c^\infty((-R/2,R/2)),\qquad
 \widehat\psi(t)=\int_{\mathbb R}\psi(x)e^{-itx}\,dx,
\]
y, para $\psi,\phi\in\mathcal D_R$,
\[
 h_{\psi,\phi}(z)=\widehat\psi(z)
 \overline{\widehat\phi(\overline z)},\qquad
 \check h_{\psi,\phi}(a)=\langle\psi,T_a\phi\rangle.
\]
La unión $\bigcup_R\mathcal D_R=C_c^\infty(\mathbb R)$ es exacta. Para todo
conjunto finito de puntos complejos distintos, las evaluaciones
$\psi\mapsto(\widehat\psi(z_1),\ldots,\widehat\psi(z_n))$ son conjuntamente
sobreyectivas cuando $R$ es suficientemente grande. Por tanto, esta clase cofinal
separa las evaluaciones requeridas por el criterio de Weil.

\begin{lemma}[Prueba de la separación conjunta]
\label{lem:pdf2-rh-separacion}
La aplicación
\[
 \mathcal D_R\longrightarrow\mathbb C^n,
 \qquad
 \psi\longmapsto
 (\widehat\psi(z_1),\ldots,\widehat\psi(z_n))
\]
es sobreyectiva para todo conjunto de puntos complejos distintos
\(z_1,\ldots,z_n\), una vez que \(R\) contiene un intervalo no vacío.
\end{lemma}

\begin{proof}
Si los funcionales de evaluación fueran linealmente dependientes, existirían
\(c_1,\ldots,c_n\), no todos nulos, tales que
\[
 \int_{\mathbb R}\psi(x)
 \left(\sum_{j=1}^nc_je^{-iz_jx}\right)\,dx=0
 \qquad(\psi\in\mathcal D_R).
\]
La función analítica \(\sum_jc_je^{-iz_jx}\) se anularía en
\((-R/2,R/2)\), luego en todo el plano. La independencia de exponenciales
con exponentes distintos fuerza \(c_j=0\) para todo \(j\), contradicción.
Los funcionales son independientes; como el codominio es finito-dimensional,
la aplicación tiene rango \(\mathbb C^n\).
\end{proof}

\section{Las dos columnas y los tres canales calculados antes del signo}

Sea $\vartheta_L$ un regulador común. Con
\[
 \ell_{p,k}=k\log p,\qquad w_{p,k}=(\log p)p^{-k/2},\qquad
 \mu_\Gamma(a)=\frac{e^{-a/2}}{1-e^{-2a}},
\]
se usan la inclusión uniforme $j_m$ y el canal de costura
$\widehat\delta_{a,m}$, que satisfacen
\[
 j_m^*j_m=I,\qquad
 \widehat\delta_{a,m}^*\widehat\delta_{b,m}
 =(I-T_a)^*(I-T_b).
\]
Las columnas reguladas son
\begin{align*}
 J_{+,m,R,L}\psi={}&
 \Bigl(
  (\sqrt{\vartheta_L(\ell_{p,k})w_{p,k}}\,
       \widehat\delta_{\ell_{p,k},m}\psi)_{\ell_{p,k}\leq R},\\
 &\hspace{9mm}a\mapsto
  \sqrt{\vartheta_L(a)\mu_\Gamma(a)}\,
       \widehat\delta_{a,m}\psi,\ \zeta_+b_+(\psi)
 \Bigr),\\
 J_{-,m,R,L}\psi={}&
 \Bigl(
  (\sqrt{2\vartheta_L(\ell_{p,k})w_{p,k}}\,j_m\psi)_{\ell_{p,k}\leq R},
  \sqrt{-c_\Gamma}\,j_m\psi,\ \zeta_-b_-(\psi)
 \Bigr),
\end{align*}
donde $c_\Gamma=\psi_\Gamma(1/4)-\log\pi<0$ y
$b_\pm=(\widehat\psi(i/2)\pm\widehat\psi(-i/2))/\sqrt2$.

Una expansión término a término, sin usar ceros ni el signo buscado, da
\begin{align*}
 &\langle J_{+,m,R,L}\psi,J_{+,m,R,L}\phi\rangle
 -\langle J_{-,m,R,L}\psi,J_{-,m,R,L}\phi\rangle
 =\mathscr I_{{\rm GW},L}(\psi,\phi),\\
 \mathscr I_{{\rm GW},L}(\psi,\phi)
 ={}&h_{\psi,\phi}(i/2)+h_{\psi,\phi}(-i/2)
 +c_\Gamma\check h_{\psi,\phi}(0)\\
 &+\int_0^\infty\vartheta_L(a)\mu_\Gamma(a)
 [2\check h(0)-\check h(a)-\check h(-a)]\,da\\
 &-\sum_{p,k}\vartheta_L(\ell_{p,k})w_{p,k}
 [\check h(\ell_{p,k})+\check h(-\ell_{p,k})].
\end{align*}
El truncamiento de ambas columnas en \(\ell_{p,k}\leq R\) es obligatorio, no
notacional. Si \(\psi,\phi\in\mathcal D_R\), entonces
\(\check h_{\psi,\phi}(\pm\ell)=0\) para \(\ell>R\); la contribución de cada
par de coordenadas omitido a la \emph{diferencia} de Grams es por tanto cero.
La columna negativa no truncada, en cambio, tendría norma infinita porque
\(\sum_p(\log p)p^{-1/2}=\infty\). En el conjunto finito
\(\ell_{p,k}\leq R\) las coordenadas primas convergen individualmente cuando
\(L\to\infty\). En el canal gamma, el integrando es \(O(a)\) en cero y
\(\mu_\Gamma(a)=O(e^{-a/2})\) en infinito, de modo que la convergencia también
es en norma. Definimos, por tanto,
\begin{equation}\label{eq:pdf2-rh-columnas-reducidas}
 J_{\pm,m,R}:=\lim_{L\to\infty}J_{\pm,m,R,L}
 \quad\text{en sus Hilbert de coordenadas truncadas},
\end{equation}
y suprimimos \(m\) cuando permanece fijo. Por convergencia dominada,
\[
 \mathscr I_{{\rm GW},L}\longrightarrow\mathscr I_{\rm GW}
 \qquad(L\to\infty),
\]
y las identidades son naturales bajo $m\mapsto m+1$ y $R\le R'$.

Las columnas reducidas \(J_{+,R}\) y \(J_{-,R}\) son ahora operadores
materiales de \(\mathcal D_R\) a Hilbert. La expansión explícita prueba, sin
consultar ceros ni signo,
\begin{equation}\label{eq:pdf2-rh-gram-difference-gw}
 \langle J_{+,R}f,J_{+,R}g\rangle
 -\langle J_{-,R}f,J_{-,R}g\rangle
 =\mathscr I_{\rm pr}(f,g)+\mathscr I_\Gamma(f,g)
  +\mathscr I_{\rm pol}(f,g)
 =\mathscr I_{\rm GW}(f,g).
\end{equation}
Los relojes \(\ell_{p,k}\) producen el sumando primo; la integral con
\(\mu_\Gamma\) y \(c_\Gamma\check h(0)\), el sumando gamma; y
\[
 \langle b_+(f),b_+(g)\rangle-
 \langle b_-(f),b_-(g)\rangle
 =h_{f,g}(i/2)+h_{f,g}(-i/2)
\]
produce el sumando polar. Esta parte del cálculo es material e independiente
del puente estructural.

\section{Coligación source--first de la imagen espectral--polar}

Fijamos primero un nivel nonádico \(m\); este índice se suprime en
\(J_{\pm,R}\), \(E_R\) y \(\Sigma_R\) hasta construir el codificador. La
descomposición de los codominios explícitos de
\eqref{eq:pdf2-rh-columnas-reducidas} es
\[
 I_+=\{\mathrm{pr},\Gamma,\mathrm{pol}\},\qquad
 I_-(R)=\{\Gamma0,\mathrm{pol}\}
 \sqcup\{(p,k):\ell_{p,k}\leq R\},
\]
\[
 \mathscr H_{+,R}^{\rm amb}=\bigoplus_{i\in I_+}H_{+,i,R},
 \qquad
 \mathscr H_{-,R}^{\rm amb}=\bigoplus_{j\in I_-(R)}H_{-,j,R}.
\]
No identificamos un subespacio alcanzable con una suma de sus proyecciones
coordenadas. Definimos, con las inclusiones ambientales entendidas,
\[
 H_{+,R}:=\overline{\operatorname{Ran}J_{+,R}},\qquad
 H_{-,R}:=\overline{\operatorname{Ran}J_{-,R}},
\]
como subespacios cerrados de los dos Hilbert ambientales. El propietario
espectral impreso antes de este capítulo construye
\(\mathcal H_R^{\rm sh}\), la diferencia antiespecular y la fila isométrica
source--first mediante
\eqref{eq:amp-rh-owner-q}--\eqref{eq:amp-rh-owner-binding}, prueba su
telescopía en \eqref{eq:amp-rh-owner-telescopia} y publica la extinción
terminal en \cref{thm:amp-rh-cierre-cofinal-propietario}. La realización
que sigue conserva esa construcción, antes de la identificación clásica,
en la fila
\begin{equation}\label{eq:pdf2-rh-coligacion-source-first}
 \Sigma_R=
 \begin{pmatrix}\mathcal K_R^{\rm src}\\ \mathcal F_R^{\rm src}\end{pmatrix}:
 H_{+,R}\longrightarrow
 H_{-,R}\oplus\mathcal H_R^{\rm sh},
 \qquad
 \Sigma_R^*\Sigma_R
 = (\mathcal K_R^{\rm src})^*\mathcal K_R^{\rm src}
  +(\mathcal F_R^{\rm src})^*\mathcal F_R^{\rm src}=I_{H_{+,R}}.
\end{equation}
Su acción alcanzable se fija sobre la columna positiva completa por
las relaciones de las dos publicaciones que forman parte de
\(\operatorname{Rel}_{\rm sp}\). Para evitar cualquier sustitución de lector,
denotamos por
\((\mathcal X_R^{\rm str},\Xi_R^{\rm str},P_R^{\rm str})\) la restricción a
\(\mathcal D_R\) del lector estructural de
\eqref{eq:pdf2-rh-res-sp-completa}. Sus dos momentos son, literalmente,
\begin{equation}\label{eq:pdf2-rh-momentos-lector-estructural}
\begin{aligned}
 \langle\Xi_R^{\rm str}f,\Xi_R^{\rm str}g\rangle
  &=\langle J_{+,R}f,J_{+,R}g\rangle,\\
 \langle P_R^{\rm str}\Xi_R^{\rm str}f,
          P_R^{\rm str}\Xi_R^{\rm str}g\rangle
  &=\langle J_{-,R}f,J_{-,R}g\rangle,
 \qquad (P_R^{\rm str})^*= (P_R^{\rm str})^2=P_R^{\rm str}.
\end{aligned}
\end{equation}
La coordenada de frontera viene con la isometría estructural
\[
 \jmath_R^{\rm sh}:
 \overline{\operatorname{Ran}((I-P_R^{\rm str})\Xi_R^{\rm str})}
 \longrightarrow\mathcal H_R^{\rm sh}.
\]
Por definición de la fila source--first, no por una identificación posterior,
\begin{equation}\label{eq:pdf2-rh-source-state-output}
 E_R:=\mathcal F_R^{\rm src}J_{+,R}
     =\jmath_R^{\rm sh}(I-P_R^{\rm str})\Xi_R^{\rm str}:
 \mathcal D_R\longrightarrow\mathcal H_R^{\rm sh},
 \qquad
 \Omega_{0,R}:=\mathcal K_R^{\rm src}J_{+,R}-J_{-,R}.
\end{equation}
No hay una coligación por canal: la misma fila recibe simultáneamente los
relojes primo--potencia, la integral gamma, el modo escalar y las dos líneas
polares. En particular, conserva todos los productos cruzados que aparecen al
polarizar \eqref{eq:pdf2-rh-gram-difference-gw}.
Con las proyecciones ambientales \(\pi_j^-:\mathscr H_{-,R}^{\rm amb}
\to H_{-,j,R}\), las filas de salida correctamente tipadas son
\begin{equation}\label{eq:pdf2-rh-bloques-coligacion}
 \mathcal K_{j,R}^{\rm src}
 :=\pi_j^-\mathcal K_R^{\rm src}:
 H_{+,R}\longrightarrow H_{-,j,R},
 \qquad
 \mathcal F_R^{\rm src}:H_{+,R}\longrightarrow\mathcal H_R^{\rm sh}.
\end{equation}
Así no se aplica \(\mathcal K_R^{\rm src}\) a una coordenada positiva que
pudiera no pertenecer por sí sola a \(H_{+,R}\). Todos los cruces entre
canales permanecen dentro de
\((\mathcal K_R^{\rm src})^*\mathcal K_R^{\rm src}\) y
\((\mathcal F_R^{\rm src})^*\mathcal F_R^{\rm src}\); no se ortogonalizan
antes de la identidad de energía.

La coordenada de supervivencia construida por
\eqref{eq:amp-rh-owner-q}, \eqref{eq:amp-rh-owner-telescopia} y
\eqref{eq:amp-rh-owner-terminal} produce los espacios residuales
\(\mathcal R_{n,R}^{\rm sh}\), continuaciones
\(\Omega_{n,R}:\mathcal D_R\to\mathcal R_{n,R}^{\rm sh}\) e isometrías
antiespeculares
\(V_{n,R}:\mathcal R_{n+1,R}^{\rm sh}\to\mathcal R_{n,R}^{\rm sh}\). Su
primer espacio y su primer operador son, por definición tipada,
\[
 \mathcal R_{0,R}^{\rm sh}:=H_{-,R},\qquad
 \Omega_{0,R}:=\mathcal K_R^{\rm src}J_{+,R}-J_{-,R}.
\]
Su recurrencia de una vuelta completa es la publicación cofinal del mismo
operador construido, no un dato añadido a la realización:
\begin{equation}\label{eq:pdf2-rh-recurrencia-residual}
 q=\frac59,\qquad
 \Omega_{n,R}=qV_{n,R}\Omega_{n+1,R},
 \qquad V_{n,R}^*V_{n,R}=I,
 \qquad
 \sup_{n\geq0}\|\Omega_{n,R}f\|<\infty
 \quad(f\in\mathcal D_R).
\end{equation}
Esta recurrencia pertenece a la coligación enriquecida: transporta la
diferencia de salida, no la forma de Weil ni su signo.

\begin{proposition}[Anulación coinductiva de la diferencia de salida]
\label{prop:pdf2-rh-omega-cero}
Para todo \(R>0\) y todo \(f\in\mathcal D_R\),
\begin{equation}\label{eq:pdf2-rh-output-binding-source}
 \Omega_{n,R}f=0\quad(n\geq0),
 \qquad
 \boxed{\mathcal K_R^{\rm src}J_{+,R}=J_{-,R}}.
\end{equation}
\end{proposition}

\begin{proof}
Iterar \eqref{eq:pdf2-rh-recurrencia-residual} \(N\) veces y usar que cada
\(V_{n,R}\) es isométrica da
\[
 \|\Omega_{n,R}f\|
 =q^N\|\Omega_{n+N,R}f\|
 \leq q^N\sup_{j\geq n}\|\Omega_{j,R}f\|.
\]
El miembro derecho converge a cero porque \(0<q=5/9<1\). Por tanto
\(\Omega_{n,R}f=0\); el caso \(n=0\) y
\eqref{eq:pdf2-rh-source-state-output} producen la identidad enmarcada.
\end{proof}

La acción fila a fila de la salida ya no queda implícita. Las proyecciones son
las restricciones a \(H_{-,R}\) de las proyecciones del Hilbert ambiental, y
\begin{equation}\label{eq:pdf2-rh-salidas-generador-a-generador}
\begin{aligned}
 \mathcal K_{\Gamma0,R}^{\rm src}J_{+,R}f
  &=\sqrt{-c_\Gamma}\,j_mf,\\
 \mathcal K_{(p,k),R}^{\rm src}J_{+,R}f
  &=\sqrt{2w_{p,k}}\,j_mf
  &&(\ell_{p,k}\leq R),\\
 \mathcal K_{{\rm pol},R}^{\rm src}J_{+,R}f
  &=\zeta_-b_-(f).
\end{aligned}
\end{equation}
Estas filas agotan \(J_{-,R}\), mientras
\begin{equation}\label{eq:pdf2-rh-frontera-generador-a-generador}
 \mathcal F_R^{\rm src}J_{+,R}f=E_Rf
\end{equation}
conserva, con todos sus cruces, la coordenada de frontera.
Evaluar \eqref{eq:pdf2-rh-coligacion-source-first} en \(J_{+,R}f\) y
\(J_{+,R}g\), y usar \eqref{eq:pdf2-rh-output-binding-source}, produce la
identidad de energía anterior al signo
\begin{equation}\label{eq:pdf2-rh-energia-source-first}
 \boxed{
 \langle J_{+,R}f,J_{+,R}g\rangle
 =\langle J_{-,R}f,J_{-,R}g\rangle
  +\langle E_Rf,E_Rg\rangle.}
\end{equation}
Al compararla con el cálculo independiente
\eqref{eq:pdf2-rh-gram-difference-gw} se obtiene, sin introducir una raíz de
la forma objetivo,
\begin{equation}\label{eq:pdf2-rh-binding-source-first}
 \boxed{E_R^*E_R=\mathscr I_{\rm GW}\big|_{\mathcal D_R}.}
\end{equation}
En particular,
\(\Delta_{{\rm bind},R}
:=E_R^*E_R-\mathscr I_{\rm GW}|_{\mathcal D_R}=0\); esta anulación es la
salida de \eqref{eq:pdf2-rh-coligacion-source-first}--
\eqref{eq:pdf2-rh-energia-source-first}, no una premisa insertada en el
publicador.

\section{Complemento estructural, hoja espejo y telescopía terminal}

El estado \(E_R=\mathcal F_R^{\rm src}J_{+,R}\) construido por la fila
source--first es la coordenada de frontera de
\(\mathfrak G_{\rm sp}\). Definimos su forma estructural por
\begin{equation}\label{eq:pdf2-rh-bw-estructural}
 \mathcal B_{{\rm sp},R}^{\rm str}
 :=E_R^*E_R
 =\mathscr I_{\rm GW}\big|_{\mathcal D_R},
\end{equation}
donde la segunda igualdad es
\eqref{eq:pdf2-rh-binding-source-first}, ya deducida de la coligación y de
la expansión generador a generador.

Sean \(P_+^{\rm sh}\) y \(P_-^{\rm sh}\) los proyectores ortogonales
complementarios de las dos hojas. Se fijan
\begin{equation}\label{eq:pdf2-rh-q-a-d}
 q=\frac59,\qquad
 A=P_+^{\rm sh}+qP_-^{\rm sh},\qquad
 D=\sqrt{1-q^2}\,P_-^{\rm sh}.
\end{equation}
La ortogonalidad da la identidad exacta
\begin{equation}\label{eq:pdf2-rh-conservacion-local}
 A^*A+D^*D
 =P_+^{\rm sh}+q^2P_-^{\rm sh}
 +(1-q^2)P_-^{\rm sh}=I.
\end{equation}
Las nueve fases coordinan una sola vuelta con memoria:
\begin{equation}\label{eq:pdf2-rh-monodromia}
 M_9=\Phi A_8\cdots A_0=qV_9,
 \qquad V_9^*V_9=I.
\end{equation}
La contracción por vuelta es \(q\), no \(q^9\); las fases internas conservan
el orden y el acarreo.

La coligación tipa
\(E_R:\mathcal D_R\to\operatorname{Ran}P_-^{\rm sh}\). Por tanto,
\begin{equation}\label{eq:pdf2-rh-binding-espejo}
 P_+^{\rm sh}E_R=0,
 \qquad AE_R=qE_R,
 \qquad
 E_R^*E_R=\mathcal B_{{\rm sp},R}^{\rm str}
 =\mathscr I_{\rm GW}\big|_{\mathcal D_R}.
\end{equation}

\begin{proposition}[Telescopía finita con terminal conservado]
\label{prop:pdf2-rh-telescopia}
Para cada \(N\geq1\),
\begin{equation}\label{eq:pdf2-rh-telescopia-finita}
 \mathcal B_{{\rm sp},R}^{\rm str}
 =\sum_{n=0}^{N-1}(DA^nE_R)^*(DA^nE_R)
 +(A^NE_R)^*(A^NE_R).
\end{equation}
Equivalentemente, evaluada en \(f,g\in\mathcal D_R\), la identidad finita de
la coligación es
\begin{equation}\label{eq:pdf2-rh-balance-finito-columnas}
\boxed{
\begin{aligned}
 &\langle J_{+,R}f,J_{+,R}g\rangle
  -\langle J_{-,R}f,J_{-,R}g\rangle\\
 &\quad=\sum_{n=0}^{N-1}
  \langle DA^nE_Rf,DA^nE_Rg\rangle
  +\langle A^NE_Rf,A^NE_Rg\rangle.
\end{aligned}}
\end{equation}
El último sumando satisface
\begin{equation}\label{eq:pdf2-rh-terminal-extincion}
 (A^NE_R)^*(A^NE_R)
 =q^{2N}\mathcal B_{{\rm sp},R}^{\rm str}
 \xrightarrow[N\to\infty]{}0.
\end{equation}
\end{proposition}

\begin{proof}
Aplicando \eqref{eq:pdf2-rh-conservacion-local} a \(A^nE_R\),
\[
 (A^nE_R)^*(A^nE_R)
 =(DA^nE_R)^*(DA^nE_R)
 +(A^{n+1}E_R)^*(A^{n+1}E_R).
\]
La suma desde \(n=0\) hasta \(N-1\) cancela sólo los términos intermedios y
conserva \((A^NE_R)^*(A^NE_R)\). Por
\eqref{eq:pdf2-rh-binding-espejo}, \(A^NE_R=q^NE_R\); por tanto el terminal
es \(q^{2N}E_R^*E_R=q^{2N}\mathcal B_{{\rm sp},R}^{\rm str}\). Como
\(0<q=5/9<1\), converge geométricamente a cero.
\end{proof}

Sólo después de demostrar \eqref{eq:pdf2-rh-terminal-extincion} se define
\begin{equation}\label{eq:pdf2-rh-l-r}
 \mathscr L_Rf=(DA^nE_Rf)_{n\geq0}.
\end{equation}
El límite de la telescopía finita da
\begin{equation}\label{eq:pdf2-rh-lstar-l}
 \boxed{
 \mathscr L_R^*\mathscr L_R
 =\mathcal B_{{\rm sp},R}^{\rm str}
 =\mathscr I_{\rm GW}\big|_{\mathcal D_R}\succeq0.}
\end{equation}

\section{Codificador común, proyección y doble publicación explícitas}

Fijemos \(m\geq1\) y dos vectores ortonormales
\(u_m,\eta_m\in\mathbb C^{9^m}\), donde \(u_m\) es el modo uniforme y
\(\eta_m\) el modo de acarreo. El espacio y la proyección del codificador son
\begin{equation}\label{eq:pdf2-rh-hilbert-comun}
 \mathcal X_{m,R}^{\rm sp}
 :=\mathbb C^{9^m}\widehat\otimes
   (H_{-,R}\oplus\mathcal H_R^{\rm sh}),
 \qquad
 P_{m,R}:=|u_m\rangle\langle u_m|\otimes
 I_{H_{-,R}\oplus\mathcal H_R^{\rm sh}},
 \qquad Q_{m,R}:=I-P_{m,R}.
\end{equation}
La acción completa del codificador de los cinco canales es
\begin{equation}\label{eq:pdf2-rh-xi-comun-explicita}
 \boxed{
 \begin{aligned}
  \mathfrak C_{m,R}f
  &:=u_m\otimes(J_{-,R}f,0)
    +\eta_m\otimes(0,E_Rf),\\
  \Xi_{m,R}f&:=\mathfrak C_{m,R}f.
 \end{aligned}}
\end{equation}
No se toma una raíz de \(\mathscr I_{\rm GW}\): las dos entradas son la
salida y el estado de la coligación source--first ya definidos en
\eqref{eq:pdf2-rh-source-state-output}.

\begin{proposition}[Los dos momentos se obtienen de la acción del codificador]
\label{prop:pdf2-rh-vmas-isometria}
Para todo \(f,g\in\mathcal D_R\),
\begin{equation}\label{eq:pdf2-rh-dos-momentos-construidos}
\begin{aligned}
 \langle\mathfrak C_{m,R}f,\mathfrak C_{m,R}g\rangle
  &=\langle J_{+,R}f,J_{+,R}g\rangle,\\
 \langle P_{m,R}\mathfrak C_{m,R}f,
         P_{m,R}\mathfrak C_{m,R}g\rangle
  &=\langle J_{-,R}f,J_{-,R}g\rangle,\\
 \langle Q_{m,R}\Xi_{m,R}f,Q_{m,R}\Xi_{m,R}g\rangle
  &=\mathscr I_{\rm GW}(f,g).
\end{aligned}
\end{equation}
\end{proposition}

\begin{proof}
La ortogonalidad \(u_m\perp\eta_m\) y
\eqref{eq:pdf2-rh-energia-source-first} dan
\[
 \langle\mathfrak C_{m,R}f,\mathfrak C_{m,R}g\rangle
 =\langle J_{-,R}f,J_{-,R}g\rangle
  +\langle E_Rf,E_Rg\rangle
 =\langle J_{+,R}f,J_{+,R}g\rangle.
\]
La proyección \(P_{m,R}\) conserva el sumando uniforme y elimina el de
acarreo; \(Q_{m,R}\) hace lo contrario. Las otras dos identidades siguen de
\eqref{eq:pdf2-rh-binding-source-first}.
\end{proof}

El codificador recién escrito es un representante explícito del lector
estructural inicial, no un lector sustituto. En los subespacios alcanzables se
define
\begin{equation}\label{eq:pdf2-rh-isomorfismo-lector-estructural}
\begin{aligned}
 \mathcal U_{m,R}^{\rm str}
  (P_R^{\rm str}\Xi_R^{\rm str}f)
   &:=u_m\otimes(J_{-,R}f,0),\\
 \mathcal U_{m,R}^{\rm str}
  ((I-P_R^{\rm str})\Xi_R^{\rm str}f)
   &:=\eta_m\otimes(0,E_Rf).
\end{aligned}
\end{equation}
Las dos fuentes son ortogonales. La primera igualdad de Gram de cada fuente
es \eqref{eq:pdf2-rh-momentos-lector-estructural}; para la segunda se usa
\(E_R=\jmath_R^{\rm sh}(I-P_R^{\rm str})\Xi_R^{\rm str}\). Por ello las dos
reglas están bien definidas, son isométricas y se extienden a la suma ortogonal
de sus cierres. En ese cierre alcanzable satisfacen
\begin{equation}\label{eq:pdf2-rh-entrelazamiento-lector-estructural}
 \mathcal U_{m,R}^{\rm str}\Xi_R^{\rm str}=\Xi_{m,R},
 \qquad
 \mathcal U_{m,R}^{\rm str}P_R^{\rm str}
 =P_{m,R}\mathcal U_{m,R}^{\rm str},
 \qquad
 (\Xi_R^{\rm str})^*(I-P_R^{\rm str})\Xi_R^{\rm str}=E_R^*E_R.
\end{equation}
Así, \(B_{{\rm sp},R}^{\rm str}\), el complemento de la estructura de
partida, es exactamente el complemento del codificador construido; no se ha
cambiado de objeto al pasar de \eqref{eq:pdf2-rh-res-sp-completa} a
\eqref{eq:pdf2-rh-xi-comun-explicita}.

La regla
\begin{equation}\label{eq:pdf2-rh-vmas-rango}
 V_{+,m,R}^{0}(J_{+,R}f):=\Xi_{m,R}f
\end{equation}
está bien definida y es isométrica por la primera igualdad de
\eqref{eq:pdf2-rh-dos-momentos-construidos}; se extiende de manera única a
una isometría
\(V_{+,m,R}:H_{+,R}\to\mathcal X_{m,R}^{\rm sp}\). La publicación negativa
es la isometría
\begin{equation}\label{eq:pdf2-rh-vmenos}
 V_{-,m,R}:H_{-,R}\longrightarrow\mathcal X_{m,R}^{\rm sp},
 \qquad V_{-,m,R}y=u_m\otimes(y,0).
\end{equation}
Por construcción,
\begin{equation}\label{eq:pdf2-rh-doble-publicacion-material}
 \boxed{
 V_{+,m,R}J_{+,R}=\Xi_{m,R},
 \qquad
 V_{-,m,R}J_{-,R}=P_{m,R}\Xi_{m,R}.}
\end{equation}
El complemento conserva exactamente el estado antiespecular y su columna de
observación:
\begin{equation}\label{eq:pdf2-rh-complemento-comun}
 \boxed{
 \Xi_{m,R}^*Q_{m,R}\Xi_{m,R}
 =E_R^*E_R
 =\mathscr L_R^*\mathscr L_R
 =\mathscr I_{\rm GW}\big|_{\mathcal D_R}.}
\end{equation}

\begin{proposition}[Conector inducido y coincidencia source--first]
\label{prop:pdf2-rh-conector}
El operador
\begin{equation}\label{eq:pdf2-rh-k-r}
 \mathcal K_R
 :=V_{-,m,R}^*P_{m,R}V_{+,m,R}:
 H_{+,R}\longrightarrow H_{-,R}
\end{equation}
es contractivo, coincide con \(\mathcal K_R^{\rm src}\) en todo
\(H_{+,R}\) y satisface
\begin{equation}\label{eq:pdf2-rh-k-binding}
 \|\mathcal K_R\|\leq1,
 \qquad
 \mathcal K_RJ_{+,R}=J_{-,R}.
\end{equation}
\end{proposition}

\begin{proof}
Las dos cartas son isometrías y \(P_{m,R}\) es una proyección ortogonal,
luego \(\|\mathcal K_R\|\leq1\). Sobre el rango denso de \(J_{+,R}\),
\[
 \mathcal K_RJ_{+,R}f
 =V_{-,m,R}^*P_{m,R}\Xi_{m,R}f
 =V_{-,m,R}^*V_{-,m,R}J_{-,R}f
 =J_{-,R}f.
\]
La misma identidad vale para \(\mathcal K_R^{\rm src}\) por
\eqref{eq:pdf2-rh-output-binding-source}; la continuidad y la densidad
prueban que ambos operadores coinciden.
\end{proof}

La diferencia de Gram queda así publicada en todas sus realizaciones
coincidentes:
\begin{equation}\label{eq:pdf2-rh-tres-publicaciones}
\boxed{
\begin{aligned}
 J_{+,R}^*J_{+,R}-J_{-,R}^*J_{-,R}
 &=\Xi_{m,R}^*Q_{m,R}\Xi_{m,R}\\
 &=E_R^*E_R\\
 &=\mathscr L_R^*\mathscr L_R\\
 &=\mathcal B_{{\rm sp},R}^{\rm str}=\mathscr I_{\rm GW}.
\end{aligned}}
\end{equation}

\section{Naturalidad cofinal y criterio reversible de Weil}

Para \(R\leq R'\), sea
\(i_{R,R'}:\mathcal D_R\hookrightarrow\mathcal D_{R'}\). Las columnas
individuales no forman un sistema isométrico: al aparecer una longitud nueva
\(R<\ell_{p,k}\leq R'\), las dos adquieren la misma energía positiva. Para
\(f,g\in\mathcal D_R\), la separación de soportes da exactamente
\begin{equation}\label{eq:pdf2-rh-naturalidad}
 \left\langle(I-T_{\ell_{p,k}})f,
                    (I-T_{\ell_{p,k}})g\right\rangle
 =2\langle f,g\rangle
 =\left\langle\sqrt2j_mf,\sqrt2j_mg\right\rangle .
\end{equation}
Por tanto el incremento primo de
\(J_{+,R'}^*J_{+,R'}\) coincide con el de
\(J_{-,R'}^*J_{-,R'}\) y se cancela en la diferencia, pero no se declara una
isometría falsa entre las columnas. La compatibilidad cofinal correcta es la
de las formas:
\begin{equation}\label{eq:pdf2-rh-naturalidad-p}
 \boxed{
 \mathscr I_{{\rm GW},R}
 =i_{R,R'}^*\mathscr I_{{\rm GW},R'}i_{R,R'},
 \qquad
 \mathcal B_{{\rm sp},R}^{\rm str}
 =i_{R,R'}^*\mathcal B_{{\rm sp},R'}^{\rm str}i_{R,R'}.}
\end{equation}
Como \(E_R^*E_R=\mathscr L_R^*\mathscr L_R=\mathscr I_{{\rm GW},R}\), las
realizaciones mínimas de Kolmogórov sí admiten isometrías únicas sobre sus
rangos alcanzables, definidas por
\begin{equation}\label{eq:pdf2-rh-naturalidad-bw}
 \iota_{R,R'}^E(E_Rf):=E_{R'}f,
 \qquad
 \iota_{R,R'}^L(\mathscr L_Rf):=\mathscr L_{R'}f
 \quad(f\in\mathcal D_R).
\end{equation}
Estas dos reglas están bien definidas precisamente por
\eqref{eq:pdf2-rh-naturalidad-p}; no se extienden por decreto a
\(J_{\pm,R}\), \(\Xi_{m,R}\) o \(\mathcal K_R\). El refinamiento nonádico
\(m\mapsto m+1\) permanece, a ventana fija, en las isometrías estructurales
que transportan \(j_m\), \(\widehat\delta_{a,m}\), \(u_m\) y \(\eta_m\).

Para no comprimir el último paso, fijamos también sus objetos y
normalizaciones. Sea
\[
 \xi(s):=\frac12s(s-1)\pi^{-s/2}\Gamma(s/2)\zeta(s)
\]
y sea \(\mathscr Z\) el multiconjunto de sus ceros, cada cero repetido según
su multiplicidad. La ecuación funcional y la realidad de los coeficientes
preservan \(\mathscr Z\) por
\(\rho\mapsto1-\rho\), \(\rho\mapsto\overline\rho\) y
\(\rho\mapsto1-\overline\rho\). Con nuestra convención de Fourier se define
\begin{equation}\label{eq:pdf2-rh-transformada-weil}
 \Phi_f(s):=\widehat f\!\left(i(s-\tfrac12)\right),
 \qquad
 \widetilde f(x):=\overline{f(-x)}.
\end{equation}
La fórmula explícita completa calculada en
\eqref{eq:pdf2-rh-gram-difference-gw}, incluidos los términos gamma y polar
que retiran polos y ceros triviales, tiene la forma espectral
\begin{equation}\label{eq:pdf2-rh-formula-espectral-weil}
 \mathscr I_{\rm GW}(f,g)
 =\sum_{\rho\in\mathscr Z}
   \Phi_f(\rho)\,
   \overline{\Phi_g(1-\overline\rho)}.
\end{equation}
La suma cuenta multiplicidades. Converge absolutamente: para
\(f,g\in C_c^\infty(\mathbb R)\), Paley--Wiener da decaimiento más rápido
que cualquier potencia en bandas verticales, mientras el número de ceros con
\(|\operatorname{Im}\rho|\leq T\) es \(O(T\log T)\).

\begin{proposition}[Criterio reversible de Weil con cuantificadores completos]
\label{prop:pdf2-rh-criterio-weil-desplegado}
Bajo la normalización
\eqref{eq:pdf2-rh-transformada-weil} y contando todos los ceros no triviales
con su multiplicidad,
\begin{equation}\label{eq:pdf2-rh-criterio-weil}
 \left[
  \mathscr I_{\rm GW}(f,f)\geq0
  \text{ para todo }f\in C_c^\infty(\mathbb R;\mathbb C)
 \right]
 \Longleftrightarrow
 \left[
  \operatorname{Re}\rho=\frac12
  \text{ para todo }\rho\in\mathscr Z
 \right].
\end{equation}
\end{proposition}

La implicación recíproca no se obtiene interpolando un número creciente de
ceros: esa operación no daría control uniforme sobre la cola. Se usa el
criterio clásico completo de Weil en la formulación de Bombieri, cuyo espacio
de prueba y cuya normalización se trasladan aquí sin abreviarlos.\footnote{%
E.~Bombieri, \emph{Remarks on Weil's quadratic functional in the theory of
prime numbers, I}, Rendiconti Lincei, Matematica e Applicazioni, ser.~9,
vol.~11 (2000), pp.~183--233. El criterio allí fijado afirma que el funcional
cuadrático de la fórmula explícita es semidefinido positivo sobre toda
\(C_c^\infty(0,\infty)\) si y sólo si vale la hipótesis de Riemann.}

\begin{proof}
Escribimos materialmente el cambio de coordenadas que identifica ambos
criterios. Para \(f\in C_c^\infty(\mathbb R)\), sea
\begin{equation}\label{eq:pdf2-rh-cambio-mellin}
 (\mathcal M f)(x):=x^{-1/2}f(\log x),
 \qquad x>0.
\end{equation}
El mapa
\[
 \mathcal M:C_c^\infty(\mathbb R)
 \longrightarrow C_c^\infty(0,\infty)
\]
es biyectivo, con inversa
\begin{equation}\label{eq:pdf2-rh-cambio-mellin-inverso}
 (\mathcal M^{-1}g)(u)=e^{u/2}g(e^u).
\end{equation}
Si
\(\widehat g_{\rm Mel}(s):=\int_0^\infty g(x)x^{s-1}\,dx\), el cambio
\(x=e^u\) da, sin factor residual,
\begin{equation}\label{eq:pdf2-rh-mellin-fourier}
 \begin{aligned}
 \widehat{\mathcal M f}_{\rm Mel}(s)
 &=\int_{\mathbb R}f(u)e^{(s-1/2)u}\,du\\
 &=\widehat f\!\left(i(s-\tfrac12)\right)
 =\Phi_f(s).
 \end{aligned}
\end{equation}
Además,
\begin{equation}\label{eq:pdf2-rh-mellin-involucion}
 \widehat{\overline{\mathcal M f}}_{\rm Mel}(1-\rho)
 =\overline{
   \widehat{\mathcal M f}_{\rm Mel}(1-\overline\rho)}
 =\overline{\Phi_f(1-\overline\rho)}.
\end{equation}
Por tanto el funcional del criterio de Weil--Bombieri es exactamente
\begin{equation}\label{eq:pdf2-rh-bombieri-identificado}
 \sum_{\rho\in\mathscr Z}
 \widehat{\mathcal M f}_{\rm Mel}(\rho)\,
 \widehat{\overline{\mathcal M f}}_{\rm Mel}(1-\rho)
 =
 \sum_{\rho\in\mathscr Z}
 \Phi_f(\rho)\overline{\Phi_f(1-\overline\rho)}
 =\mathscr I_{\rm GW}(f,f).
\end{equation}
No se cambia la clase de pruebas, no se pierde ninguna multiplicidad y no se
trunca el conjunto de ceros, porque \(\mathcal M\) es una biyección. El
criterio clásico citado aplicado a
\eqref{eq:pdf2-rh-bombieri-identificado} da precisamente la equivalencia
\eqref{eq:pdf2-rh-criterio-weil}. En la dirección directa, si
\(\operatorname{Re}\rho=1/2\), la misma identidad se reduce explícitamente a
\[
 \mathscr I_{\rm GW}(f,f)
 =\sum_{\rho\in\mathscr Z}|\Phi_f(\rho)|^2\geq0.
\]
\end{proof}

La igualdad \(\bigcup_R\mathcal D_R=C_c^\infty(\mathbb R)\) es exacta; cada
función empleada en el argumento pertenece a alguna ventana y la naturalidad
cofinal transporta allí la positividad ya demostrada.

\begin{theorem}[Resolución Riemann--Weil desde la imagen espectral--polar]
\label{thm:pdf2-rh-resolucion}
Fijada la estructura completa
\(\mathfrak G_{\rm sp}\), la forma
completa de Guinand--Weil es positiva en la clase cofinal y separadora. En
consecuencia, todo cero no trivial \(\rho\) de la función zeta de Riemann
satisface
\[
 \boxed{\operatorname{Re}\rho=\frac12.}
\]
\end{theorem}

\begin{proof}
La expansión primo--gamma--polar identifica la diferencia de las dos columnas
con \(\mathscr I_{\rm GW}\) usando un único regulador y lo retira por
convergencia dominada. La recurrencia antiespecular
\eqref{eq:pdf2-rh-recurrencia-residual} anula la diferencia de salida antes de
formar una desigualdad; por tanto la isometría \(\Sigma_R\) da
\eqref{eq:pdf2-rh-energia-source-first}. Comparar esa identidad con la
expansión calculada produce
\(E_R^*E_R=\mathscr I_{\rm GW}\). La identidad local
\(A^*A+D^*D=I\) se itera sin borrar el terminal y da
\eqref{eq:pdf2-rh-telescopia-finita}; el terminal es exactamente
\(q^{2N}E_R^*E_R\) y se extingue sólo en el límite. De ahí
\(\mathscr L_R^*\mathscr L_R=\mathscr I_{\rm GW}\succeq0\).
La acción
\eqref{eq:pdf2-rh-xi-comun-explicita} construye entonces el codificador, sus
dos momentos, la proyección y las dos publicaciones, y
\eqref{eq:pdf2-rh-tres-publicaciones} identifica todas las realizaciones del
mismo cuadrado. La naturalidad transporta esa identidad a toda la unión
cofinal. Finalmente,
\eqref{eq:pdf2-rh-criterio-weil} concluye la afirmación.
\end{proof}

\section{Control Redheffer local posterior y no gobernante}

El siguiente despliegue verifica que la pasividad local puede obtenerse desde
el acarreo sin usar la forma de Weil. Es un control de no-circularidad y no una
premisa del \cref{thm:pdf2-rh-resolucion}. Esta transferencia truncada no es
la coligación completa \(\Sigma_R\) de
\eqref{eq:pdf2-rh-coligacion-source-first}.

Para \(N=9^m\), sea \(U_{a,N}\) el odómetro unitario
\[
 U_{a,N}(e_j\otimes f)=e_{j+1}\otimes f\ (j<N-1),
 \qquad U_{a,N}(e_{N-1}\otimes f)=e_0\otimes T_af.
\]
Con \(P_N=|u_N\rangle\langle u_N|\otimes I\), \(Q_N=I-P_N\),
\(A_{a,N}=P_NU_{a,N}P_N\) y \(C_{a,N}=Q_NU_{a,N}P_N\),
\begin{equation}\label{eq:pdf2-rh-source-energia}
 A_{a,N}^*A_{a,N}+C_{a,N}^*C_{a,N}=P_N,
 \qquad
 C_{a,N}^*C_{a,N}
 =\frac{N-1}{N^2}(2I-T_a-T_a^*).
\end{equation}
Para \(0<s<1\), la transferencia
\begin{equation}\label{eq:pdf2-rh-source-redheffer}
 \mathscr R_{a,N}(s)=(sI-A_{a,N})(I-sA_{a,N})^{-1}
\end{equation}
satisface
\begin{equation}\label{eq:pdf2-rh-source-defecto}
\begin{aligned}
 I-\mathscr R_{a,N}(s)^*\mathscr R_{a,N}(s)
 ={}&(1-s^2)(I-sA_{a,N}^*)^{-1}\\
 &\quad\cdot C_{a,N}^*C_{a,N}(I-sA_{a,N})^{-1}\succeq0.
\end{aligned}
\end{equation}
Con
\[
 s_{p,N}=\frac{Np^{-1/2}}{1+(N-1)p^{-1/2}},
 \qquad s_\Gamma(a)=e^{-a/2},
\]
la suma sobre relojes primos y la integral directa gamma conservan norma a lo
sumo uno. De forma explícita,
\begin{equation}\label{eq:pdf2-rh-source-relojes}
 \mathscr R_{N,R}
 =\bigoplus_{\log p\leq R}
   \mathscr R_{\log p,N}(s_{p,N})
 \ \oplus\!
 \int_0^\infty{}^\oplus
   \mathscr R_{a,N}(s_\Gamma(a))\,da,
 \qquad \|\mathscr R_{N,R}\|\leq1.
\end{equation}
La carta finita de roles
\[
 \mathcal H_+=\operatorname{span}\{f_3,f_6,f_9\},\quad
 \mathcal H_-=\operatorname{span}\{f_4,f_8\},\quad
 C_\angle=|f_4\rangle\langle f_3|+|f_8\rangle\langle f_9|
\]
satisface \(I-C_\angle^*C_\angle=P_6\). Con las cartas energéticas
two-way se fija
\begin{equation}\label{eq:pdf2-rh-source-cartas}
\begin{aligned}
 \mathcal K_q&=K_2\otimes(\mathbb C^9)^{\otimes q},\\
 \mathsf B_{\pm,q}:\mathcal M_{\pm,q,R}
 &\longrightarrow
 \mathcal K_q\widehat\otimes\mathcal H_\pm
 \widehat\otimes\mathcal H_{\mathfrak C_R},\\
 \mathsf B_{\pm,q}^*\mathsf B_{\pm,q}&=\mathsf G_{\pm,q},
 \qquad
 \mathsf B_{\pm,q}^{\sharp}
 =\mathsf G_{\pm,q}^{-1}\mathsf B_{\pm,q}^*.
\end{aligned}
\end{equation}
La transferencia
\begin{equation}\label{eq:pdf2-rh-source-transferencia}
 \mathscr C_{q,R}^{\rm src}
 =\mathsf B_{-,q}^{\sharp}
  (I\otimes C_\angle\otimes\mathscr R_{N,R})\mathsf B_{+,q}
\end{equation}
tiene defecto
\begin{align}
 \Delta_{q,R}^{\rm src}
 ={}&I\otimes\bigl[P_6\otimes I
 +(P_3+P_9)\otimes
 (I-\mathscr R_{N,R}^*\mathscr R_{N,R})\bigr]\succeq0,
 \label{eq:pdf2-rh-source-delta}\\
 I-(\mathscr C_{q,R}^{\rm src})^\sharp\mathscr C_{q,R}^{\rm src}
 ={}&\mathsf B_{+,q}^{\sharp}\Delta_{q,R}^{\rm src}
 \mathsf B_{+,q}\succeq0.
 \label{eq:pdf2-rh-source-contractividad}
\end{align}
Este control conserva ruta, reloj, rol, acarreo y los canales antes de formar
una diferencia. No se identifica con \(\mathcal K_R\) y no sustituye la
publicación común explícita de
\eqref{eq:pdf2-rh-doble-publicacion-material}; por ello su alcance local no
reabre la conclusión cofinal.

\section{Falsadores de una sombra ablacionada}

\begin{falsifier}[Sustituciones que cambian el objeto de partida]
El argumento anterior no se transfiere a una sombra que elimine una de las
trece coordenadas de \eqref{eq:pdf2-rh-res-sp-completa}, construya \(J_+\) y
\(J_-\) desde antecedentes distintos, reemplace una proyección ortogonal por
un idempotente no autoadjunto, omita un canal, use reguladores incompatibles,
escoja \(P_R\) después del signo, borre el terminal antes del límite o reduzca
la clase a una familia no cofinal.
\end{falsifier}

\begin{proof}
Cada operación rompe, respectivamente, el dominio de
\(\mathfrak G_{\rm sp}\), la buena definición de
\eqref{eq:pdf2-rh-vmas-rango}, el cuadrado ortogonal, la identidad funcional,
la dirección causal, la telescopía finita o el criterio reversible. El
falsador afecta a la sombra sustituida; no demueve la imagen estructural
completa empleada en el \cref{thm:pdf2-rh-resolucion}.
\end{proof}

\paragraph{Procedencia y fuerza probatoria.}
La restricción espectral--polar, la coligación source--first, su complemento y
la hoja espejo son \texttt{ARQUITECTURA\_AUTORAL\_PREEXISTENTE}. La
anulación coinductiva de \(\Omega\), la identidad finita con terminal y el
codificador de acción
\eqref{eq:pdf2-rh-xi-comun-explicita} son
\texttt{FORMALIZACION\_NUEVA} de esa estructura completa. La identificación
Guinand--Weil y el cierre cofinal son \texttt{RESULTADO\_RECUPERADO}. Las
ecuaciones
\eqref{eq:pdf2-rh-coligacion-source-first}--
\eqref{eq:pdf2-rh-tres-publicaciones} exhiben el objeto, sus dominios, su
acción, el terminal y el límite sin introducir una raíz abstracta de la forma;
su fuerza es \emph{demostrado con estructura de partida explícita}.

\endgroup
% Cuerpo científico recuperado y distribuido en su residencia terminal.
\section{Clausura cofinal del criterio de Riemann--Weil}

La realización genealógica produce una familia cofinal y natural de formas
positivas cuya publicación coincide con el funcional completo de
Guinand--Weil. La extinción del terminal y la compatibilidad con las
inclusiones permiten pasar al límite sin introducir una lista de ceros ni el
signo buscado. Por el criterio reversible de Weil, todo cero no trivial
\(\rho\) de la función zeta de Riemann satisface
\[
  \operatorname{Re}\rho=\frac12.
\]
Los cálculos regulados de ventana finita son controles posteriores; no forman
parte de las premisas del cierre cofinal.

\paragraph{Factorización cofinal de la conclusión espectral.}
La cadena probatoria reúne, en el orden causal ya establecido, la forma
genealógica finita, su complemento positivo, la naturalidad bajo inclusiones y
la extinción compatible del término terminal. Cada ventana preserva los relojes
primos, el canal arquimediano, la orientación de las dos hojas y el registro de
frontera. La cofinalidad transporta conjuntamente esos datos al límite y la
identidad de Guinand--Weil identifica la forma resultante con su publicación
espectral clásica.

En consecuencia, la positividad no procede de una selección posterior de
ceros: desciende del operador positivo construido antes de la evaluación. La
involución funcional conserva la hoja espejo y fija la recta
\(\operatorname{Re}\rho=1/2\); la compatibilidad cofinal conserva esa fijación
en el paso desde las ventanas reguladas hasta la forma completa. Los cálculos
finitos cumplen así su función propia de verificación posterior, mientras que
la clausura matemática reside en la composición positiva y natural anterior.



\HMTFlushCurrentChapter
\chapter{Caracterización Cayley--Pontryagin de los ceros críticos y construcción coinductiva del campo espectral de Weil}
\begingroup
\renewcommand{\chapter}[2][]{}%
\chapter{Caracterización Cayley–Pontryagin de los ceros críticos y reconstrucción efectiva del campo espectral de Weil}
\label{chap:hmt83-doble-circulo-campo-espectral}

Una vez localizada la recta crítica por la factorización positiva, la
transformada de Cayley publica las alturas como fases de un primer círculo. El
entrelazador de Pontryagin las transporta después al círculo nonádico sin
identificar retorno de fase con retorno de estado. El teorema del doble círculo
y la cuantización de las órbitas se articulan con la reconstrucción efectiva
del campo espectral. La exposición conserva explícitamente las hipótesis, los
residuos de Galerkin, el cociclo de memoria, el transporte torcido y el
entrelazador que identifica ambas realizaciones sin confundir el índice de la
vuelta con la altura espectral.

\Needspace{9\baselineskip}
\section{Anulación de la función completada y cierre espectral sobre la recta crítica}
La imagen convencional de los ceros de Riemann suele comenzar demasiado tarde. Se presenta la función zeta, se prolonga analíticamente y se pregunta en qué puntos se anula. El lector ve una función compleja y una colección de puntos misteriosos distribuidos sobre una recta.\par
La lectura HMT cambia el orden de la pregunta. No comienza mirando dónde se hace cero la publicación final. Comienza reconstruyendo el estado que esa publicación ha contraído: los relojes de las potencias primas, el canal gamma, el canal polar, las dos hojas, su orientación, el acarreo, la memoria, la frontera, el terminal y la conexión nonádica que los transporta conjuntamente.\par
Sólo después aparece la función completada y se pregunta por sus ceros.\par
La diferencia es radical: un cero deja de ser una ausencia y pasa a ser un acontecimiento de cierre.\par
\Needspace{7\baselineskip}
\subsection{Anulación escalar y persistencia del estado enriquecido}
Si\par
\[
\xi(\rho)=0,
\]
el cero es verdadero. No es aproximado, ficticio ni decorativo. La función completada se anula exactamente en \(\rho\).\par
Pero lo que se anula es una publicación escalar. No desaparece el punto \(\rho\), no desaparece su multiplicidad, no desaparece la estructura prima que participa en la fórmula explícita y no desaparece el estado enriquecido que precede a la publicación.\par
Ésta es una distinción elemental y, sin embargo, decisiva:\par
\[
\text{valor publicado igual a cero}
\quad\not\Rightarrow\quad
\text{estado generador inexistente}.
\]
Una representación útil es una franja oscura producida por interferencia. La oscuridad no significa que ninguna onda haya llegado. Significa exactamente lo contrario: han llegado contribuciones cuya relación de fase es tan precisa que su suma visible se cancela.\par
Los ceros no triviales pueden leerse de esta manera: como nodos de interferencia aritmética global.\par
No son nodos de una sola secuencia de primos ni pertenecen individualmente a un primo. En la publicación completa intervienen simultáneamente:\par
\begin{itemize}
\item los relojes de las potencias primas,
\end{itemize}
\[
\ell_{p,k}=k\log p;
\]
\begin{itemize}
\item el canal arquimediano gamma;
\item el canal polar, que conserva la normalización y retira los polos y ceros triviales;
\item las dos hojas especulares;
\item la frontera que la proyección visible no muestra.
\end{itemize}
Un cero es una cancelación exacta de esa publicación conjunta. Su existencia revela una coordinación, no un vacío.\par
Por tanto, los ceros constituyen \textbf{nodos espectrales de cancelación aritmética}: la anulación de \(\xi\) expresa una coordinación global exacta entre las contribuciones de la publicación conjunta.\par
\Needspace{7\baselineskip}
\subsection{Fijación de la recta crítica por la involución funcional \(s\mapsto1-\bar s\)}
La función completada conserva las simetrías\par
\[
\rho\longmapsto1-\rho,
\qquad
\rho\longmapsto\overline\rho.
\]
Combinándolas aparece la involución\par
\[
\rho\longmapsto\rho^\sharp
:=
1-\overline\rho.
\]
Si\par
\[
\rho=\beta+i\gamma,
\]
entonces\par
\[
\rho^\sharp
=
1-\beta+i\gamma.
\]
La reflexión conserva la altura \(\gamma\), pero intercambia \(\beta\) y \(1-\beta\). Su conjunto de puntos fijos es:\par
\[
\rho=\rho^\sharp
\quad\Longleftrightarrow\quad
\beta=1-\beta
\quad\Longleftrightarrow\quad
\beta=\frac12.
\]
Por tanto,\par
\[
\operatorname{Fix}(\sharp)
=
\left\{
\frac12+i\gamma:\gamma\in\mathbb R
\right\}.
\]
La recta crítica no es solamente “la mitad” por una preferencia geométrica. Es la costura donde las dos hojas funcionales dejan de aparecer como posiciones distintas y se convierten en una sola posición autodual.\par
Fuera de la recta, un cero aparece ligado a una órbita de hasta cuatro puntos:\par
\[
\rho,\qquad
1-\rho,\qquad
\overline\rho,\qquad
1-\overline\rho.
\]
Sobre la recta crítica, la reflexión funcional conjugada ya no envía el cero a otra posición lateral:\par
\[
1-\overline\rho=\rho.
\]
El cero cierra consigo mismo.\par
Éste es el primer sentido preciso del círculo de cierre: no se trata únicamente de que dos mitades sumen uno, sino de que la órbita especular encuentra un punto fijo.\par
Sin embargo, esto todavía no demuestra la hipótesis de Riemann. La igualdad\par
\[
\frac12=1-\frac12
\]
identifica el eje; no obliga por sí sola a que todos los ceros estén sobre él.\par
La obligación procede de la positividad.\par
\Needspace{7\baselineskip}
\subsection{Cuadraticidad del producto cruzado bajo clausura de la involución funcional}
La forma espectral completa puede escribirse como\par
\[
\mathscr I_{\mathrm{GW}}(f,g)
=
\sum_{\rho\in\mathscr Z}
\Phi_f(\rho)\,
\overline{\Phi_g(1-\overline\rho)},
\]
donde \(\mathscr Z\) es el multiconjunto de ceros no triviales, conservando sus multiplicidades.\par
Esta fórmula determina cómo se emparejan los ceros: la evaluación situada en \(\rho\) se combina con la evaluación situada en su reflejo \(\rho^\sharp=1-\overline\rho\).\par
Si el cero está fuera de la recta, la expresión empareja dos puntos diferentes:\par
\[
\Phi_f(\rho)\,
\overline{\Phi_f(\rho^\sharp)}.
\]
Eso es un producto cruzado. No es necesariamente un módulo cuadrado y, sobre una clase de pruebas suficientemente rica, puede producir direcciones de signo negativo.\par
Si el cero está en la recta crítica, entonces\par
\[
\rho^\sharp=\rho
\]
y el mismo término se convierte en\par
\[
|\Phi_f(\rho)|^2.
\]
La forma completa pasa a ser:\par
\[
\mathscr I_{\mathrm{GW}}(f,f)
=
\sum_{\rho\in\mathscr Z}
|\Phi_f(\rho)|^2.
\]
Ahora el espejo no empareja dos posiciones desplazadas. Cada punto se empareja consigo mismo y la publicación se convierte en una suma de cuadrados.\par
Ésta es la lectura conceptual de la demostración: HMT construye la forma de Weil aguas arriba como una energía de frontera,\par
\[
\mathscr I_{\mathrm{GW}}
=
E_R^*E_R
=
\mathscr L_R^*\mathscr L_R
\succeq0.
\]
No se declara positiva porque los ceros estén sobre la recta. Se construye como cuadrado antes de usar los ceros.\par
Después, el criterio reversible de Weil dice que esa positividad sobre toda la clase separadora de funciones de prueba es equivalente a que todos los ceros no triviales satisfagan\par
\[
\operatorname{Re}\rho=\frac12.
\]
Por tanto, la positividad no “atrae” dinámicamente a los ceros hacia el centro. Los ceros no se desplazan desde una posición inicial. Lo que ocurre es más fuerte: cualquier configuración fuera del eje fijo produciría una dirección negativa detectable, mientras que el complemento HMT es materialmente una norma cuadrada y no puede adquirirla.\par
La recta crítica es el único lugar compatible con ambas publicaciones.\par
\Needspace{7\baselineskip}
\subsection{Caracterización espectral de los ceros}
Bajo esta lectura, un cero no trivial puede describirse simultáneamente de cuatro maneras compatibles.\par
Es un \textbf{cero analítico verdadero}, porque\par
\[
\xi(\rho)=0.
\]
Es un \textbf{punto fijo del espejo}, porque\par
\[
\rho=1-\overline\rho.
\]
Es un \textbf{átomo espectral positivo}, porque en la forma de Weil aparece como\par
\[
|\Phi_f(\rho)|^2.
\]
Y es un \textbf{nodo global de los relojes primos y del borde arquimediano}, porque no pertenece a un canal aislado: surge de la publicación coordinada de los términos primo, gamma y polar.\par
El cero es nulo como amplitud y positivo como localización espectral. No hay contradicción. Son dos lectores diferentes.\par
La función vale cero en \(\rho\), pero \(\rho\) continúa apareciendo, con toda su multiplicidad, dentro de la medida espectral positiva. El cero no ha sido borrado por ser cero. Al contrario: se convierte en uno de los lugares que sostienen la suma positiva.\par
Éste es probablemente el cambio interpretativo más importante:\par
\begin{quote}
El cero no es el lugar donde la estructura desaparece. Es el lugar donde la estructura completa alcanza una cancelación visible y deja una marca espectral.\par
\end{quote}
\Needspace{7\baselineskip}
\subsection{Identificación de la altura de un \(0\) no trivial con una frecuencia espectral real}
La normalización empleada en el criterio de Weil es\par
\[
\Phi_f(s)
=
\widehat f\!\left(
i\left(s-\frac12\right)
\right).
\]
Si\par
\[
\rho=\frac12+i\gamma,
\]
entonces\par
\[
i\left(\rho-\frac12\right)
=
-\gamma
\]
y, por tanto,\par
\[
\Phi_f(\rho)
=
\widehat f(-\gamma).
\]
La altura \(\gamma\) se convierte literalmente en una frecuencia real de Fourier.\par
Así, la recta crítica puede leerse como una recta de frecuencias: la coordenada real \(\tfrac12\) fija el equilibrio especular y la coordenada \(\gamma\) registra la frecuencia del nodo.\par
Esto explica por qué el lenguaje espectral es apropiado. Los ceros no son solamente coordenadas complejas; son frecuencias en las que la publicación aritmética completa manifiesta su condición nodal.\par
Estas alturas admiten la interpretación de las frecuencias nodales de un sistema global: puntos donde una amplitud observable se cancela mientras la estructura que produce esa cancelación permanece activa.\par
Pero hay que proteger esta frase contra una exageración. De aquí no se sigue automáticamente que cada \(\gamma\) sea un autovalor energético de un Hamiltoniano físico concreto. Para eso habría que construir ese Hamiltoniano, su dominio y la correspondencia espectral. Lo que está demostrado es la lectura Fourier–Mellin, la forma positiva completa y la localización de los ceros.\par
Se trata, con rigor, de frecuencias espectrales. No se convierten, sin otro entrelazador, en niveles físicos de energía.\par
\Needspace{7\baselineskip}
\subsection{Parametrización circular conforme de la recta crítica}
Tu intuición del círculo no es sólo metafórica. Existe una transformación de Möbius exacta:\par
\[
\tau(s)
=
\frac{s-1}{s}.
\]
Se cumple:\par
\[
\operatorname{Re}s=\frac12
\quad\Longleftrightarrow\quad
|\tau(s)|=1.
\]
Por tanto, la recta crítica compactificada por su punto del infinito se transforma en la circunferencia unidad.\par
Si\par
\[
s=\frac12+it,
\]
entonces\par
\[
\tau(s)\in S^1.
\]
Al escribir\par
\[
\tau=e^{i\theta},
\]
la inversa es\par
\[
s=\frac1{1-\tau}
=
\frac12+\frac{i}{2}\cot\frac{\theta}{2}.
\]
La altura \(t\) de la recta se convierte en una posición angular sobre el círculo. Cuando \(t\) recorre toda la recta, la carta recorre la circunferencia compactificada. Los dos extremos \(t\to-\infty\) y \(t\to+\infty\) se encuentran en el punto \(\tau=1\), aunque lo hacen desde orientaciones opuestas.\par
También el espejo adquiere una forma circular:\par
\[
\tau(1-\overline s)
=
\frac1{\overline{\tau(s)}}.
\]
La involución espectral se convierte en reflexión respecto de la circunferencia unidad, cuyo conjunto fijo es precisamente el círculo.\par
Así que la expresión “círculo de cierre” posee un contenido exacto:\par
\begin{itemize}
\item la recta crítica compactificada es un círculo;
\item el espejo funcional se convierte en reflexión respecto de ese círculo;
\item los ceros no triviales son puntos marcados sobre esa frontera;
\item cada punto está fijado por la reflexión que antes intercambiaba las dos hojas.
\end{itemize}
El círculo no genera los ceros ni determina sus alturas. Toda la recta se transforma en el círculo, mientras que la ecuación\par
\[
\xi(\rho)=0
\]
selecciona sus puntos espectrales.\par
Tampoco significa que los ceros estén igualmente espaciados o que haya nueve por vuelta. El círculo es la geometría compactificada del locus crítico; la distribución de los ceros sobre él conserva su propia aritmética.\par
\Needspace{7\baselineskip}
\subsection{Círculo nonádico y holonomía de fase}
Hay un segundo círculo, anterior al analítico: la base cíclica de las nueve fases.\par
La conexión nonádica no consta de nueve filtros ni de nueve canales independientes. Es una única holonomía:\par
\[
\Gamma_9
=
\operatorname{Hol}_{C_9}
(\nabla_{\mathrm{TPK}}).
\]
Su ley de retorno es\par
\[
g(K+9)=g(K),
\]
pero simultáneamente\par
\[
q_{\mathrm{mem}}
(\Gamma_9\widehat x)
=
q_{\mathrm{mem}}(\widehat x)+1
\]
y, en general,\par
\[
\Gamma_9\widehat x
\neq
\widehat x.
\]
La fase vuelve; el estado no se reinicia.\par
Visto desde la base, se ha recorrido un círculo y se ha regresado a la misma fase. Visto desde el estado enriquecido, se ha acumulado una vuelta de memoria. La imagen más fiel para el lector es:\par
\begin{itemize}
\item un círculo en la proyección;
\item una hélice en el estado completo.
\end{itemize}
Desde arriba, la hélice parece cerrar. Desde dentro, cada vuelta ocupa una altura distinta.\par
La parte literal no es la imagen de la hélice, sino las identidades matemáticas: retorno de fase, incremento de memoria y cambio general del estado.\par
Esta estructura fue comprobada en 396 niveles de refinamiento por canal. El certificado registra 43 vueltas nonádicas completas y 387 comparaciones \(K,K+9\) por canal: la fase coincide y el estado cambia en todas ellas. No es una manera poética de hablar de periodicidad; es una separación ejecutada entre retorno visible y continuidad genealógica.\par
\Needspace{7\baselineskip}
\subsection{El espejo nonádico y el espejo de Riemann}
Tu intuición es correcta: en ambos casos aparece un espejo. Pero conviene decir con exactitud cómo se relacionan.\par
El espejo interno nonádico intercambia las dos componentes laterales, asociadas a las lecturas \(\pi/e\), mientras conserva el eje \(\varphi\) y el agregado correspondiente. La involución de Riemann es\par
\[
s\longmapsto1-\overline s.
\]
No son literalmente el mismo mapa, porque tienen dominios y codominios distintos. Igualarlos sin el entrelazador sería confundir la arquitectura interna con su publicación analítica.\par
La afirmación correcta es más interesante: el entrelazador espectral–polar publica la arquitectura two-way del TPK como la involución funcional de la función completada. El espejo nonádico es anterior; el espejo de Riemann es su realización en el codominio analítico.\par
En ambos casos existe:\par
\begin{itemize}
\item una pareja de hojas;
\item un eje que permanece fijo;
\item una transformación que intercambia las posiciones laterales;
\item una costura donde las dos publicaciones coinciden;
\item una memoria que no desaparece al efectuar la identificación.
\end{itemize}
La recta \(\Re s=\tfrac12\) es la costura visible de ese espejo en el plano complejo.\par
\Needspace{7\baselineskip}
\subsection{Cierre residual y fase nonádica}
La expresión “los nueves son falsos ceros” contiene una intuición muy precisa si distinguimos el entero, la clase residual y el estado completo.\par
Para \(n>0\), la carta positiva utiliza\par
\[
\rho_9(n)
=
1+((n-1)\bmod9).
\]
Equivalentemente,\par
\[
\rho_9(n)
=
\begin{cases}
n\bmod9,
&
n\not\equiv0\pmod9,\\[1mm]
9,
&
n\equiv0\pmod9.
\end{cases}
\]
Por tanto, cuando un número pertenece a la clase nula módulo nueve, la carta positiva no publica \(0\): publica \(9\).\par
Al mismo tiempo se conserva\par
\[
\kappa_9(n)
=
\frac{n-\rho_9(n)}9
\]
y la reconstrucción\par
\[
n
=
9\kappa_9(n)
+
\rho_9(n).
\]
El nueve es, así, el representante positivo de la clase residual cero. Es cero en el cociente, pero no es el entero cero y no es un estado vacío.\par
El nueve dice: “la vuelta se ha completado”. El cociente dice cuántas coronas o vueltas se han acumulado. Si se sustituye todo por el residuo \(0\), se conserva la congruencia, pero se pierde la diferencia entre ausencia, cierre y memoria.\par
Ése es el sentido riguroso de “falso cero”:\par
\begin{quote}
El nueve es un cero para la sombra residual, pero sería falso interpretarlo como aniquilación del estado. Es un cero publicado con memoria de cierre.\par
\end{quote}
Por ejemplo, \(27\), \(72\) y \(729\) poseen sombra residual cero y raíz digital positiva nueve. Sin embargo, no son el mismo número ni tienen la misma genealogía. Uno puede proceder de una reversión, otro de una evaluación diferente y otro de una potencia. El residuo no reconstruye la ruta.\par
Esto conecta profundamente con los ceros de Riemann, pero no los identifica.\par
Un cero de \(\xi\) es un cero analítico auténtico. No es el número nueve. No existe aquí una afirmación de que cada cero no trivial corresponda individualmente a una fase nueve concreta.\par
La relación estructural es ésta:\par
\[
\text{una publicación puede anularse}
\quad\text{sin que se anule su genealogía}.
\]
El nueve es una falsa ausencia aritmética. El cero de zeta es una falsa ausencia ontológica. Ambos son ceros verdaderos dentro de su lector, pero ninguno significa que no haya estructura detrás.\par
\Needspace{7\baselineskip}
\subsection{Separación tipada de los ceros de \(\zeta\), \(\xi\) y el determinante espectral}
Para que la fuerza de la narración no borre los tipos matemáticos, conviene separar tres ceros.\par
\Needspace{5\baselineskip}
\subsubsection{\(0\) analítico como anulación de la función en el dominio espectral}
\[
\xi(\rho)=0.
\]
Es un punto espectral de la función completada. Conserva posición, altura, multiplicidad y simetrías.\par
\Needspace{5\baselineskip}
\subsubsection{\(0\) terminal como objeto universal del morfismo de cierre}
En la telescopía HMT aparece\par
\[
(A^NE_R)^*(A^NE_R)
=
q^{2N}
\mathcal B_{\mathrm{sp},R}^{\mathrm{str}}
\longrightarrow0,
\qquad
q=\frac59.
\]
Es la extinción de un resto después de que toda su contribución haya sido transferida a una suma de cuadrados positivos.\par
\Needspace{5\baselineskip}
\subsubsection{Representante positivo de la clase residual nula: \(n\equiv0\pmod 9\) y \(\rho_9(n)=9\)}
\[
n\equiv0\pmod9,
\qquad
\rho_9(n)=9.
\]
Es la clase nula modular publicada como representante positivo, acompañada por su cociente y memoria.\par
No son el mismo objeto. Pero comparten una gramática:\par
\[
\text{anulación visible}
+
\text{conservación estructural}.
\]
\Needspace{7\baselineskip}
\subsection{Agotamiento del objeto terminal bajo la filtración cofinal}
La demostración conserva en cada profundidad una identidad finita:\par
\[
\mathcal B_{\mathrm{sp},R}^{\mathrm{str}}
=
\sum_{n=0}^{N-1}
(DA^nE_R)^*(DA^nE_R)
+
(A^NE_R)^*(A^NE_R).
\]
El último término es el terminal. No se elimina para obtener una suma positiva más bonita. Permanece presente en cada etapa finita.\par
Sólo después se prueba:\par
\[
(A^NE_R)^*(A^NE_R)
=
q^{2N}
\mathcal B_{\mathrm{sp},R}^{\mathrm{str}}
\longrightarrow0.
\]
Esto permite otra lectura del cero: el terminal llega a cero porque su contenido ha sido expresado y contabilizado nivel tras nivel. Cada parte perdida por la contracción visible ha quedado registrada en una coordenada de memoria positiva.\par
El resto se extingue, pero no ha sido borrado.\par
Es como vaciar un depósito trasladando exactamente su contenido a una secuencia de recipientes numerados. Al final, el depósito inicial está vacío, pero nada ha desaparecido: existe un ledger completo de la transferencia.\par
Aquí el cero significa \textbf{transferencia terminada}.\par
\Needspace{7\baselineskip}
\subsection{Información genealógica del continuo codificada por los ceros críticos}
La recta crítica\par
\[
\frac12+i\mathbb R
\]
es una copia desplazada de la recta real. Desde HMT, esa recta no es el objeto inicial: es una publicación arquimediana del continuo genealógico.\par
Los ceros forman un conjunto discreto dentro de esa publicación continua. No perforan la recta ni interrumpen el continuo. Constituyen un divisor espectral: puntos aislados, cada uno con multiplicidad, sobre una sombra continua.\par
Esta convivencia de continuo y discreto es esencial. La recta ofrece el codominio continuo; la genealogía determina las marcas discretas que aparecen sobre ella.\par
La imagen que propondría al lector es la de una membrana continua atravesada por nodos espectrales. Los nodos no rompen la membrana. Hacen visible su estructura interna.\par
En términos holográficos, el cero es una marca de frontera producida por una organización global del interior. No pertenece a un único primo, igual que un píxel de una imagen holográfica no pertenece a un único punto del objeto representado. Cada cero responde al conjunto completo de relaciones primo–gamma–polar.\par
Esto no significa que cada cero contenga por sí solo toda la aritmética de manera recuperable. Significa que su posición no puede atribuirse a una entrada local aislada: es una respuesta global del sistema.\par
\Needspace{7\baselineskip}
\subsection{Correspondencia entre cuantización espectral, frecuencias reales y realización física del operador}
Los ceros se interpretan como nodos de una transferencia aritmética global:\par
\begin{itemize}
\item la coordenada \(\gamma\) actúa como frecuencia;
\item la recta crítica es la superficie de equilibrio especular;
\item la forma de Weil es la energía positiva del complemento;
\item la conexión nonádica conserva la memoria de cada vuelta;
\item el terminal mide lo que todavía no ha sido expresado;
\item el cero visible señala una cancelación perfecta de la amplitud publicada.
\end{itemize}
No se identifican con partículas, agujeros, estados de vacío o masas, ni con niveles energéticos materiales sin construir antes un operador físico adicional.\par
La interpretación física válida es estructural: nodo, frecuencia, espejo, transferencia, memoria y cierre. La ontología física concreta requeriría su propio transductor.\par
\Needspace{7\baselineskip}
\subsection{Corolarios de localización, multiplicidad y simetría de la caracterización espectral}
Antes de la resolución, la pregunta parecía ser:\par
\begin{quote}
¿Por qué unos puntos definidos por una función compleja parecen alinearse misteriosamente sobre una recta?\par
\end{quote}
Después de la resolución, la pregunta puede formularse de otro modo:\par
\begin{quote}
¿Cómo publica la estructura espectral–polar una energía de frontera que sólo puede permanecer positiva cuando cada cero se identifica con su reflejo?\par
\end{quote}
La respuesta es:\par
\begin{itemize}
\item las columnas positiva y negativa proceden del mismo estado;
\item el estado de frontera \(E_R\) conserva lo que la proyección no muestra;
\item la diferencia de los dos Grams es
\end{itemize}
\[
E_R^*E_R;
\]
\begin{itemize}
\item la telescopía conserva el terminal hasta demostrar su extinción;
\item la forma de Weil completa coincide con ese cuadrado;
\item el criterio reversible obliga a que cada cero sea un punto fijo del espejo.
\end{itemize}
La recta ya no es una coincidencia visual. Es el conjunto fijo de la involución compatible con la positividad construida.\par
\Needspace{7\baselineskip}
\subsection{Hipótesis y codominio exactos del teorema de localización espectral}
Esta lectura permite afirmar:\par
\[
\{\text{ceros no triviales}\}
\subset
\left\{
\frac12+i\gamma:\gamma\in\mathbb R
\right\}.
\]
Y, mediante la carta de Möbius,\par
\[
\left\{
\frac12+i\gamma
\right\}
\longrightarrow
S^1.
\]
Por tanto, todos los ceros no triviales son puntos marcados del círculo crítico.\par
No se ha construido aquí una biyección individual:\par
\[
\{
\text{ceros no triviales}
\}
\longleftrightarrow
\{
\text{vueltas nonádicas concretas}
\}.
\]
No se afirma que el primer cero sea la primera vuelta, que haya nueve ceros por ciclo o que sus alturas se generen repitiendo una fase. La resolución demuestra localización, positividad, cierre especular y pertenencia al círculo. Una cuantización individual de las alturas sería un desarrollo adicional y necesitaría su propio mapa.\par
Tampoco deben confundirse el espejo interno \(\pi/e\), la involución funcional de Riemann y una eventual simetría física. La relación entre ellos es de publicación tipada, no de identidad nominal.\par
\Needspace{7\baselineskip}
\subsection{Teorema de cierre espectral mediante la factorización positiva de Guinand–Weil}
La interpretación operatoria se resume mediante las propiedades siguientes:\par
\begin{quote}
Los ceros no triviales de Riemann son anulaciones exactas de la amplitud analítica que conservan la genealogía de su construcción. Cada cero determina una frecuencia global, un marcador espectral y un punto autodual de la involución funcional. La compactificación de la recta crítica produce el círculo fijo común de las dos hojas. La conexión nonádica distingue el retorno de fase del reinicio del estado: el cierre circular conserva el avance de memoria. La clase nonádica \(9\equiv0\pmod 9\) representa el cierre visible acompañado por la información interna del estado enriquecido.\par
\end{quote}
La anulación espectral constituye así una forma concentrada de cierre operatorio y conserva el objeto genealógico subyacente.\par
Esta interpretación físico-matemática acompaña a los resultados demostrados sin sustituirlos y no afirma una correspondencia individual entre ceros y vueltas nonádicas mientras no se construya el entrelazador correspondiente.\par

\section{Forma de Weil, traslación y cuantización de las alturas espectrales}
El objeto correcto no es una vuelta desnuda, sino una órbita helicoidal enriquecida:\par
\[
\boxed{
\text{cero no trivial}
\longleftrightarrow
\text{átomo espectral ponderado}
\longleftrightarrow
\text{carácter helicoidal nonádico}.
}
\]
Esta formalización actúa después del continuo completo y no redefine ni separa sus cinco características.\par
\Needspace{7\baselineskip}
\subsection{Construcción del espacio de Hilbert mediante la forma de Weil}
Se toma\par
\[
\mathcal D=C_c^\infty(\mathbb R),
\qquad
\langle f,g\rangle_W
:=
\mathscr I_{\mathrm{GW}}(f,g).
\]
En la resolución HMT, esta forma no nace de la lista de ceros. Se construye primero desde los canales primo–potencia, gamma y polar, y queda factorizada por la pérdida nonádica:\par
\[
\mathscr I_{\mathrm{GW}}
=
E^*E
=
\mathscr L^*\mathscr L
\succeq0.
\]
Se define\par
\[
\mathcal N_W
=
\{f\in\mathcal D:
\mathscr I_{\mathrm{GW}}(f,f)=0\},
\]
y completamos:\par
\[
\mathcal H_W
=
\overline{\mathcal D/\mathcal N_W}.
\]
Éste es el espacio de Weil positivo construido desde el lado aritmético. Los ceros todavía no han sido utilizados para definirlo.\par
Al mismo tiempo, la telescopía proporciona el espacio nonádico de pérdida\par
\[
\mathcal K_{\mathrm{loss}}
=
\overline{\operatorname{Ran}\mathscr L},
\]
y la igualdad de normas define una unidad canónica\par
\[
W_{\mathrm{loss}}:
\mathcal H_W
\longrightarrow
\mathcal K_{\mathrm{loss}},
\qquad
W_{\mathrm{loss}}[f]=\mathscr Lf.
\]
No estamos colocando después una copia artificial de la forma de Weil dentro del TPK: el propio espacio de pérdida de la prueba es una realización de su Hilbert positivo.\par
\Needspace{7\baselineskip}
\subsection{Operador de traslación independiente de la localización previa de los ceros}
En el dominio global \(\mathcal D\), no en una ventana fija, se considera\par
\[
(T_tf)(x)=f(x-t).
\]
La invariancia puede comprobarse antes de diagonalizar la fórmula en los ceros. Con la convención\par
\[
\widehat{T_tf}(z)
=
e^{-izt}\widehat f(z),
\]
se tiene\par
\[
h_{T_tf,T_tg}(z)=h_{f,g}(z),
\qquad
\check h_{T_tf,T_tg}(a)=\check h_{f,g}(a).
\]
Por tanto, todos los términos primo, gamma y polar permanecen invariantes:\par
\[
\mathscr I_{\mathrm{GW}}(T_tf,T_tg)
=
\mathscr I_{\mathrm{GW}}(f,g).
\]
La traslación conserva \(\mathcal N_W\) y desciende a un grupo unitario:\par
\[
U_t[f]=[T_tf].
\]
Después de la localización crítica ya demostrada, su continuidad fuerte sigue de\par
\[
\|U_t[f]-[f]\|_W^2
=
\sum_\gamma
m_\gamma
\left|e^{it\gamma}-1\right|^2
\left|\widehat f(-\gamma)\right|^2
\longrightarrow0.
\]
El teorema de Stone produce entonces un operador autoadjunto \(H_W\):\par
\[
U_t=e^{itH_W}.
\]
Este operador se ha definido mediante traslaciones y la forma primo–gamma–polar. No se ha definido escribiendo previamente una matriz diagonal con los ceros.\par
\Needspace{7\baselineskip}
\subsection{Identificación espectral de las alturas de los ceros no triviales}
Como la hipótesis de Riemann ya está demostrada, cada cero no trivial tiene la forma\par
\[
\rho_\gamma=\frac12+i\gamma.
\]
Si \(\Gamma_\xi\) es el conjunto de alturas distintas y \(m_\gamma\) su multiplicidad, se introduce la medida de conteo ponderada\par
\[
\mu_\xi
=
\sum_{\gamma\in\Gamma_\xi}
m_\gamma\,\delta_\gamma.
\]
La fórmula espectral de Weil se convierte en\par
\[
\mathscr I_{\mathrm{GW}}(f,g)
=
\sum_{\gamma\in\Gamma_\xi}
m_\gamma\,
\widehat f(-\gamma)
\overline{\widehat g(-\gamma)}.
\]
Por tanto,\par
\[
\mathcal A[f](\gamma)
=
\widehat f(-\gamma)
\]
define una isometría\par
\[
\mathcal A:
\mathcal H_W
\longrightarrow
L^2(\Gamma_\xi,\mu_\xi).
\]
Además, no queda una parte espectral escondida fuera de su rango. El rango de \(\mathcal A\) es cerrado y se conserva por todos los multiplicadores\par
\[
M_t a(\gamma)=e^{it\gamma}a(\gamma).
\]
La proyección sobre un átomo individual se obtiene mediante el promedio espectral\par
\[
P_\gamma
=
\operatorname*{s-lim}_{T\to\infty}
\frac1{2T}
\int_{-T}^{T}
e^{-it\gamma}M_t\,dt.
\]
La separación de evaluaciones de las transformadas de funciones de prueba garantiza que, para cada \(\gamma\), existe algún \(f\) con\par
\[
\widehat f(-\gamma)\neq0.
\]
Por tanto, \(P_\gamma\operatorname{Ran}\mathcal A\neq0\). Como el rango es reductor, contiene el eje atómico correspondiente. Todos los ejes atómicos pertenecen al rango y generan densamente \(L^2(\mu_\xi)\). En consecuencia,\par
\[
\boxed{
\mathcal H_W
\simeq
L^2(\Gamma_\xi,\mu_\xi)
}
\]
y\par
\[
\boxed{
\mathcal A H_W\mathcal A^{-1}
=
M_\gamma.
}
\]
Así pues,\par
\[
\boxed{
\sigma(H_W)=\Gamma_\xi,
}
\]
sin alturas adicionales, sin alturas omitidas y sin componente continua.\par
Esto es ya una cuantización espectral individual: las alturas de los ceros son exactamente el espectro puntual de un operador definido desde la forma aritmética positiva.\par
La multiplicidad exige una precisión. La forma escalar registra\par
\[
\mu_\xi(\{\gamma\})=m_\gamma.
\]
Es decir, conserva exactamente el peso algebraico del cero. No produce automáticamente \(m_\gamma\) vectores distinguibles: si un cero fuera múltiple, la evaluación en ese mismo punto se repite con peso \(m_\gamma\). La correspondencia canónica es entre el divisor de ceros y los átomos espectrales ponderados. Separar internamente las copias exigiría una forma positiva de jets o demostrar simplicidad.\par
\Needspace{7\baselineskip}
\subsection{Transformación conforme de la recta crítica en el círculo espectral con conservación de alturas}
Se usa la normalización\par
\[
\tau(s)=\frac{s-1}{s}.
\]
Entonces\par
\[
\Re s=\frac12
\quad\Longleftrightarrow\quad
|\tau(s)|=1.
\]
Para\par
\[
\rho_\gamma=\frac12+i\gamma
\]
se obtiene\par
\[
\tau_\gamma
=
\tau(\rho_\gamma)
=
\frac{-\frac12+i\gamma}{\frac12+i\gamma}
=
\frac{\gamma+\frac i2}{\gamma-\frac i2}
\in S^1.
\]
Ésta es exactamente la transformada de Cayley del operador \(H_W\):\par
\[
\mathcal C_W
=
\left(H_W+\frac i2I\right)
\left(H_W-\frac i2I\right)^{-1}.
\]
Sobre el átomo \(\gamma\),\par
\[
\mathcal C_WP_\gamma
=
\tau_\gamma P_\gamma.
\]
La correspondencia no pierde información. Su inversa es\par
\[
\rho_\gamma=\frac1{1-\tau_\gamma}.
\]
Si\par
\[
\tau_\gamma=e^{i\theta_\gamma},
\]
entonces\par
\[
\boxed{
\gamma
=
\frac12\cot\frac{\theta_\gamma}{2}.
}
\]
Por tanto, la altura completa está contenida en el ángulo. La compactificación de la recta conserva la localización de cada cero.\par
Además,\par
\[
\tau_{-\gamma}
=
\overline{\tau_\gamma}
=
\tau_\gamma^{-1}.
\]
Los ceros conjugados corresponden a orientaciones inversas del mismo círculo. Éste es el espejo convertido en una identidad angular exacta.\par
La correspondencia demostrada es\par
\[
\boxed{
\sum_\rho m_\rho[\rho]
\longleftrightarrow
\sum_\gamma m_\gamma
[\tau_\gamma,P_\gamma].
}
\]
No depende de ordenar previamente los ceros por altura.\par
\Needspace{7\baselineskip}
\subsection{Monodromía de cada \(0\) sobre el círculo espectral}
Transportamos ahora el operador circular al espacio nonádico:\par
\[
\mathcal C_{\mathrm{loss}}
=
W_{\mathrm{loss}}\,
\mathcal C_W\,
W_{\mathrm{loss}}^{-1}.
\]
En la telescopía HMT,\par
\[
q=\frac59
\]
es la contracción correspondiente a una vuelta completa. Se define la realización espectral de esa vuelta:\par
\[
\mathcal M_{\mathrm{spec}}
=
q\,\mathcal C_{\mathrm{loss}}.
\]
Si \(v_\gamma\) pertenece al átomo espectral \(P_\gamma\), entonces\par
\[
\boxed{
\mathcal M_{\mathrm{spec}}^n v_\gamma
=
q^n\tau_\gamma^n v_\gamma.
}
\]
Ésta es la órbita helicoidal del cero.\par
Su radio es\par
\[
\|\mathcal M_{\mathrm{spec}}^n v_\gamma\|
=
q^n\|v_\gamma\|,
\]
mientras su ángulo acumulado es\par
\[
n\theta_\gamma.
\]
La parte radial conserva la extinción terminal; la parte angular conserva la identidad individual del cero. Y ambas satisfacen la telescopía:\par
\[
\|v_\gamma\|^2
=
(1-q^2)
\sum_{k=0}^{N-1}
q^{2k}\|v_\gamma\|^2
+
q^{2N}\|v_\gamma\|^2.
\]
El terminal\par
\[
q^{2N}\|v_\gamma\|^2
\longrightarrow0
\]
no borra la altura: la altura permanece en el carácter unitario \(\tau_\gamma^n\).\par
Ésta es la lectura correcta de la vuelta:\par
\Needspace{8\baselineskip}
\begin{center}
\small
\begin{tabularx}{\textwidth}{@{}YY@{}}
\toprule
Dato & Qué registra \\
\midrule
\(n\) & cuántas vueltas completas han transcurrido \\
\(q=5/9\) & cuánto terminal queda después de cada vuelta \\
\(\tau_\gamma\) & qué cero individual transporta la órbita \\
\(m_\gamma\) & el peso algebraico de ese cero \\
memoria/carry & por qué volver a una fase visible no devuelve el mismo estado \\
\bottomrule
\end{tabularx}
\end{center}
Por eso no hay que asignar el primer cero a la primera vuelta. Cada cero atraviesa toda la escalera de vueltas con su propio carácter.\par
Una imagen pedagógica adecuada sería la de un disco que gira: el número de vueltas no identifica la nota que suena. La vuelta \(n\) dice cuánto tiempo ha transcurrido; la frecuencia identifica el modo. Aquí \(n\) es el contador, \(q\) es la atenuación y \(\tau_\gamma\) es la firma espectral.\par
\Needspace{7\baselineskip}
\subsection{Localización del campo espectral de Weil mediante el entrelazador nonádico de prolongación}
La construcción anterior vive realmente en\par
\[
\overline{\operatorname{Ran}\mathscr L},
\]
el espacio de pérdida producido por la telescopía nonádica. No es una analogía externa añadida a la prueba de Riemann.\par
Pero debe distinguirse de la holonomía combinatoria nativa. El TPK contiene también\par
\[
M_9=qV_9,
\qquad
V_9^*V_9=I.
\]
La identidad de defecto determina \(q\), pero no determina la parte unitaria \(V_9\). Un mismo balance\par
\[
M_9^*M_9=q^2I
\]
es compatible con muchas fases unitarias diferentes. La telescopía, por sí sola, no puede decir cuál de ellas contiene las alturas.\par
La identidad entre las órbitas espectrales recién construidas y las órbitas del ledger combinatorio APP–TRIT–TPK queda establecida por el entrelazador que se construye, sobre los estados cilíndricos enriquecidos, en el teorema Cayley--Pontryagin de este mismo capítulo:\par
\[
J:
\mathcal K_{\mathrm{loss}}
\longrightarrow
\mathcal K_{\mathrm{TPK}}^{\mathrm{unit}}
\]
tal que\par
\[
\boxed{
J\mathcal C_{\mathrm{loss}}
=
V_9J.
}
\]
Además, \(J\) debe conservar cilindro, orientación, carry, ruta, memoria, frontera y multiplicidad espectral.\par
La igualdad anterior no se postula ni se recibe de la realización clásica: resulta del cociclo entero de memoria, del transporte cilíndrico torcido por \(\mathcal C_{\mathrm{Weil}}\), de la proyectividad de los pesos y del encaje isométrico \(J_k\) construido más abajo. Al iterarla se obtiene\par
\[
(qV_9)^nJv_\gamma
=
q^n\tau_\gamma^nJv_\gamma,
\]
y entonces cada cero queda identificado con el carácter unitario de una órbita nativa del TPK. La identidad conserva cilindro, orientación, acarreo, ruta, memoria, frontera y multiplicidad espectral, de modo que la correspondencia no depende de una enumeración externa cero--vuelta.\par
\Needspace{7\baselineskip}
\subsection{Algoritmo espectral de cálculo independiente de la introducción previa de los ceros}
El resultado no se limita a representar una lista ya conocida. El operador \(H_W\) se ha definido desde traslaciones y desde la forma primo–gamma–polar. Su resolvente puede escribirse como\par
\[
(H_W-iI)^{-1}
=
i\int_0^\infty
e^{-t}U_{-t}\,dt.
\]
Por tanto, para funciones de prueba \(f_j,f_k\),\par
\[
\left\langle
(H_W-iI)^{-1}[f_j],[f_k]
\right\rangle_W
=
i\int_0^\infty
e^{-t}
\mathscr I_{\mathrm{GW}}(T_{-t}f_j,f_k)\,dt.
\]
El miembro derecho se calcula desde los canales primo, gamma y polar. No necesita alimentar una tabla de alturas.\par
Como el espectro es discreto y \(|\gamma|\to\infty\), el resolvente es compacto. Si \(P_N\) son proyecciones finitas que convergen fuertemente a la identidad,\par
\[
P_N(H_W-iI)^{-1}P_N
\longrightarrow
(H_W-iI)^{-1}
\]
en norma.\par
Sus autovalores convergen hacia\par
\[
\lambda_\gamma
=
\frac1{\gamma-i},
\]
y de ellos se recupera\par
\[
\gamma=i+\lambda_\gamma^{-1}.
\]
Esto abre una ruta computacional genuina: formar matrices finitas usando solamente la forma aritmética construida directamente desde los datos y recuperar las alturas como límites espectrales.\par
Las cotas efectivas de truncamiento y la elección de bases optimizadas pertenecen a la realización numérica del resolvente compacto; refinan su eficiencia computacional sin condicionar el teorema espectral ni el entrelazamiento operatorio ya construido. No hay una enumeración tautológica del tipo \(n=N(\gamma_n)\): los coeficientes de las aproximaciones proceden de la forma aritmética y sus autovalores convergen al espectro del operador.\par
\Needspace{7\baselineskip}
\subsection{Cuantización espectral y entrelazamiento nonádico}
La correspondencia exacta se formula en el siguiente nivel tipado:\par
\begin{quote}
La correspondencia probada no identifica ordinalmente el cero número \(j\) con la vuelta nonádica número \(j\), ni distribuye nueve ceros por ciclo, ni impone una repetición periódica de alturas. Identifica exactamente el divisor de ceros no triviales con los átomos ponderados del generador de traslaciones de la forma de Weil. La transformada de Cayley convierte cada altura en un carácter unitario único; la realización nonádica de pérdida lo transporta como una órbita helicoidal \(q^n\tau_\gamma^n\); y el transporte torcido \(\widetilde U_k\), junto con \(J_k\), entrelaza ese carácter espectral con la holonomía combinatoria nativa.\par
\end{quote}
La interpretación operatoria resultante es la siguiente:\par
\begin{quote}
Cada cero determina el carácter angular individual de una historia espectral cuyo iterado atraviesa sucesivas vueltas nonádicas. El índice de vuelta cuantifica la profundidad de memoria y el cero fija el carácter unitario que gobierna su evolución angular.\par
\end{quote}
En consecuencia, el resultado establece una vía operatoria efectiva: demuestra la cuantización espectral individual de las alturas, construye su órbita nonádica canónica y la identifica, mediante \(J_k\), con el carácter de memoria que actúa sobre todas las vueltas del ledger TPK. La identificación es armónica y operatoria; no depende de numerar los ceros con las vueltas.\par
\paragraph{Dominio y construcción del campo espectral de Weil}
La forma positiva \(\mathscr I_{\mathrm{GW}}=\mathscr L^*\mathscr L\), la contracción \(q=5/9\), la telescopía y la vuelta \(M_9=qV_9\) constituyen la estructura matemática de partida. Sobre ella se construyen el espacio de Hilbert de Weil, el grupo de traslaciones, su generador de Stone, la correspondencia de Cayley y las órbitas \(q^n\tau_\gamma^n\). La aproximación del resolvente directamente desde la forma aritmética define su realización computacional. El cociclo de memoria y el transporte torcido construidos a continuación proporcionan el entrelazador exacto con la holonomía combinatoria nativa.

\section{Teorema Cayley–Pontryagin del doble círculo}
No existe una correspondencia “cero número \(j\) = vuelta nonádica número \(j\)”. Lo que sí existe es algo más profundo: cada cero determina un carácter unitario de la memoria vertical del ledger, y la órbita completa de ese carácter es exactamente la órbita espectral del cero después de la transformación de Cayley.\par
\Needspace{7\baselineskip}
\subsection{Teorema del doble círculo}
La conexión nonádica posee una memoria vertical entera\par
\[
q_{\mathrm{mem}}:\operatorname{Obj}(\mathrm{TPK})\longrightarrow\mathbb Z.
\]
Para cualquier camino admisible \(\gamma:x\to y\), se define el cociclo\par
\[
m(\gamma)
=
q_{\mathrm{mem}}(y)-q_{\mathrm{mem}}(x).
\]
Es aditivo:\par
\[
m(\gamma_2\gamma_1)
=
m(\gamma_2)+m(\gamma_1),
\]
y para toda vuelta completa de la conexión,\par
\[
m(\Gamma_9)=1.
\]
Por iteración,\par
\[
m(\Gamma_9^n)=n.
\]
Éste es el dato decisivo. La fase visible vuelve después de nueve pasos, pero el estado no vuelve al mismo punto: la memoria ha avanzado una unidad.\par
Por otro lado, la demostración de Riemann construye primero la forma positiva de Weil y su espacio GNS:\par
\[
\mathcal H_W
=
\overline{\mathcal D/\ker \mathscr I_{\mathrm{GW}}}.
\]
La traslación construida desde la forma aritmética produce un generador autoadjunto \(H_W\). Sin introducir ceros como datos, se forma su transformada de Cayley:\par
\[
\mathcal C_{\mathrm{Weil}}
=
\left(H_W+\frac{i}{2}I\right)
\left(H_W-\frac{i}{2}I\right)^{-1}.
\]
Como \(H_W\) es autoadjunto,\par
\[
\mathcal C_{\mathrm{Weil}}^*
\mathcal C_{\mathrm{Weil}}=I.
\]
Por tanto, \(\mathcal C_{\mathrm{Weil}}\) es unitario y su espectro pertenece a \(S^1\).\par
Después del reconocimiento espectral, no antes, una altura\par
\[
\rho_\gamma=\frac12+i\gamma
\]
produce la fase\par
\[
\tau_\gamma
=
\frac{\gamma+i/2}{\gamma-i/2}
=
\frac{\rho_\gamma-1}{\rho_\gamma},
\qquad
|\tau_\gamma|=1.
\]
La inversión es exacta:\par
\[
\rho_\gamma=\frac{1}{1-\tau_\gamma},
\qquad
\gamma
=
\frac12\cot\left(\frac{\arg\tau_\gamma}{2}\right).
\]
Así aparecen dos realizaciones circulares obtenidas por lectores diferenciados del mismo estado enriquecido:\par
\[
\left\{\Re s=\frac12\right\}\cup\{\infty\}
\overset{\tau(s)=(s-1)/s}{\cong}
S^1,
\]
y\par
\[
\widehat{\mathbb Z_{\mathrm{mem}}}
\cong S^1.
\]
El primero es el círculo de Cayley de la recta crítica. El segundo es el dual de Pontryagin de la memoria vertical del ledger. El entrelazador demuestra que no son dos círculos meramente parecidos: son las dos realizaciones del mismo operador de fase.\par
\Needspace{7\baselineskip}
\subsection{Isometría de transporte torcido \(\widetilde U_k\) sobre estados cilíndricos, pesos proyectivos y memoria de Weil}
Sea \(Z_k\) el conjunto de estados cilíndricos enriquecidos en profundidad \(k\), con medida proyectiva de soporte completo \(\mu_k\), y sea\par
\[
\tau_{k+1,k}:Z_{k+1}\longrightarrow Z_k
\]
el truncamiento.\par
Para un descendiente \(z'\) de \(z\), se define\par
\[
w_k(z'|z)
=
\sqrt{\frac{\mu_{k+1}(z')}{\mu_k(z)}}.
\]
La proyectividad da\par
\[
\sum_{\tau_{k+1,k}z'=z}w_k(z'|z)^2=1.
\]
Sobre\par
\[
\mathscr K_k
=
\ell^2(Z_k)\otimes\mathcal H_W
\]
se define el transporte torcido\par
\[
\widetilde U_k(e_z\otimes v)
=
\sum_{\tau_{k+1,k}z'=z}
w_k(z'|z)\,
e_{z'}\otimes
\mathcal C_{\mathrm{Weil}}^{\,m(z\to z')}v.
\]
En forma matricial,\par
\[
[\widetilde U_k]_{z',z}
=
\mathbf 1_{\{\tau z'=z\}}
\sqrt{\frac{\mu_{k+1}(z')}{\mu_k(z)}}
\,
\mathcal C_{\mathrm{Weil}}^{\,m(z\to z')}.
\]
Esta fórmula conserva materialmente todos los descendientes. No suma caminos distintos como si fueran el mismo: ledger, orientación, carry, frontera, memoria, ruta y terminal continúan presentes en la etiqueta \(z'\).\par
La isometría se demuestra directamente. Los descendientes de padres diferentes son ortogonales; para cada padre, los pesos tienen suma cuadrática uno; y cada potencia de \(\mathcal C_{\mathrm{Weil}}\) es unitaria. Por tanto,\par
\[
\widetilde U_k^*\widetilde U_k=I.
\]
La composición también es exacta:\par
\[
\widetilde U_{\ell,j}\widetilde U_{j,k}
=
\widetilde U_{\ell,k},
\]
porque los pesos telescopan y el cociclo de memoria es aditivo.\par
Hay incluso una equivalencia de calibre explícita. Si\par
\[
G_k(e_z\otimes v)
=
e_z\otimes
\mathcal C_{\mathrm{Weil}}^{\,q_{\mathrm{mem}}(z)}v,
\]
entonces\par
\[
\widetilde U_k
=
G_{k+1}(U_k\otimes I)G_k^*.
\]
Por tanto, el sistema torcido no es una decoración arbitraria: es el transporte cilíndrico nativo visto en el calibre espectral determinado por la memoria.\par
Defínase ahora el vector de medida\par
\[
\Omega_k
=
\sum_{z\in Z_k}
\sqrt{\mu_k(z)}\,e_z,
\qquad
\|\Omega_k\|=1,
\]
y el encaje\par
\[
J_k:\mathcal H_W\longrightarrow\mathscr K_k,
\qquad
J_kv=\Omega_k\otimes v.
\]
En una vuelta completa, todas las rutas poseen incremento de memoria uno. Por ello,\par
\[
\boxed{
\widetilde U_{\Gamma_9,k}J_k
=
J_{k+9}\mathcal C_{\mathrm{Weil}}
}
\]
y, tras \(n\) vueltas,\par
\[
\boxed{
\widetilde U_{\Gamma_9^n,k}J_k
=
J_{k+9n}\mathcal C_{\mathrm{Weil}}^n.
}
\]
Éste es el entrelazador anunciado y construido dentro de la misma cadena causal.\par
\Needspace{7\baselineskip}
\subsection{Identidad orbital bajo la conjugación entre las \(2\) realizaciones espectrales}
Sea \(\psi_\gamma\) un vector de la fibra espectral asociada a \(\gamma\). Entonces\par
\[
\mathcal C_{\mathrm{Weil}}\psi_\gamma
=
\tau_\gamma\psi_\gamma.
\]
Aplicando el cuadrado anterior:\par
\[
\widetilde U_{\Gamma_9^n}J\psi_\gamma
=
J\mathcal C_{\mathrm{Weil}}^n\psi_\gamma
=
\tau_\gamma^nJ\psi_\gamma.
\]
Pero el carácter de la memoria correspondiente a \(\tau_\gamma\) es\par
\[
\chi_{\tau_\gamma}(m)
=
\tau_\gamma^m.
\]
Como\par
\[
m(\Gamma_9^n)=n,
\]
se obtiene\par
\[
\boxed{
\widetilde U_{\Gamma_9^n}J\psi_\gamma
=
\chi_{\tau_\gamma}
\!\left(m(\Gamma_9^n)\right)
J\psi_\gamma.
}
\]
Ésta es la igualdad literal buscada.\par
No significa que el cero número 37 sea la vuelta número 37. Significa que cada cero selecciona un carácter del grupo completo de memoria y que ese carácter recorre todas las vueltas:\par
\[
n\longmapsto\tau_\gamma^n.
\]
Por tanto:\par
\begin{itemize}
\item la vuelta nonádica proporciona la variable entera \(n\);
\item la memoria proporciona el grupo \(\mathbb Z\);
\item el cero proporciona un carácter \(\chi_{\tau_\gamma}\in\widehat{\mathbb Z}\);
\item la órbita espectral y la órbita del ledger coinciden término a término después del encaje \(J\).
\end{itemize}
La cautela anterior queda superada sin caer en “nueve ceros por ciclo”. No hay una enumeración cero–vuelta. Hay una descomposición armónica de la misma dinámica de vuelta.\par
\Needspace{7\baselineskip}
\subsection{Compactificación circular y levantamiento espiral de la órbita espectral}
Debe conservarse una distinción crucial. La fase es unitaria:\par
\[
\mathcal C_{\mathrm{Weil}}^*
\mathcal C_{\mathrm{Weil}}=I.
\]
La dinámica estricta de la hoja de pérdida es\par
\[
M_9=qV_9,
\qquad
q=\frac59.
\]
El operador conjunto correcto es\par
\[
\widetilde M_9
=
\frac59\,\widetilde U_{\Gamma_9}.
\]
Sobre la fibra espectral,\par
\[
\widetilde M_9^nJ\psi_\gamma
=
\left(\frac59\right)^n
\tau_\gamma^n
J\psi_\gamma.
\]
Así aparecen dos imágenes simultáneas:\par
\begin{itemize}
\item en la frontera normalizada, la órbita es circular:
\end{itemize}
  \[
  \tau_\gamma^n\in S^1;
  \]
\begin{itemize}
\item en la realización estricta, es una espiral:
\end{itemize}
  \[
  \left(\frac59\right)^n\tau_\gamma^n\longrightarrow0.
  \]
El cero terminal no borra la fase ni la genealogía. El radio se extingue, pero el ledger conserva cuántas vueltas se realizaron y qué carácter fue transportado.\par
La telescopía es exacta:\par
\[
I
=
\sum_{r=0}^{L-1}
\widetilde M_9^{*r}
D_9^*D_9
\widetilde M_9^r
+
\widetilde M_9^{*L}\widetilde M_9^L,
\]
con\par
\[
D_9^*D_9
=
I-\widetilde M_9^*\widetilde M_9
=
\frac{56}{81}I,
\]
y terminal\par
\[
\widetilde M_9^{*L}\widetilde M_9^L
=
\left(\frac{25}{81}\right)^LI
\longrightarrow0.
\]
Esto responde preventivamente a una objeción de tipo: no se ha confundido una contracción con un unitario. La fase vive en \(\mathcal C_{\mathrm{Weil}}\); la pérdida radial vive en \(5/9\).\par
\Needspace{7\baselineskip}
\subsection{Compatibilidad entre cilindros genealógicos, generación de constantes y entrelazador espectral}
Tu intuición sobre los cilindros es correcta, con una distinción precisa.\par
Los cilindros numéricos son semiabiertos:\par
\[
C(N_K,6K)
=
\left[
\frac{N_K}{3^{6K}},
\frac{N_K+1}{3^{6K}}
\right),
\]
y cada uno se divide en 729 subcilindros. Su diámetro satisface\par
\[
\operatorname{diam}C(N_K,6K)=3^{-6K}\longrightarrow0.
\]
Los cilindros genealógicos \([a]_k\) del límite inverso son bolas clopen de la ultramétrica:\par
\[
d_\vartheta(x,y)
=
\vartheta^{N(x,y)}.
\]
Por tanto, “radio tendiendo a cero” tiene dos publicaciones:\par
\begin{enumerate}
\item En la lectura arquimediana, la intersección de intervalos encajados publica un único real.
\item En el estado enriquecido, la historia completa continúa conservando todos los prefijos que produjeron ese real.
\end{enumerate}
Así, conforme a la \hyperref[sec:tesis-inaugural-helicoidal-hmt-md]{tesis inaugural}, el mismo sistema coinductivo publica correlacionadamente \(\pi\), \(e\), \(\varphi\) y \(\alpha\). En el orden operatorio interno, el registro \(U_{12}^{\mathrm{sgn}}\), el sello \(K\) y el acarreo dodecafásico actúan una vez disponibles \(P,E,\Phi\): publican y certifican la coordenada \(\alpha\) de esa misma generación coinductiva correlacionada, sin incorporarla como objeto causalmente externo. Las secciones arquimedianas y la publicación firmada conservan codominios diferentes sin romper su genealogía conjunta.\par
Estas constantes no seleccionan las fases de Riemann. Su relevancia es más profunda: demuestran que el mismo sistema de cilindros finitos puede producir un objeto continuo sin destruir la historia que lo generó.\par
El nuevo entrelazador utiliza exactamente esa propiedad. El círculo espectral no se superpone desde fuera al continuo: aparece como dual de su memoria entera.\par
\Needspace{7\baselineskip}
\subsection{Invariante graduado del corredor espectral \(4\to2\to6\to54\)}
La cadena \(4\to54\) puede ahora situarse sin numerología.\par
El bloque central de APP es\par
\[
B_c
=
\begin{pmatrix}
7&2\\
2&7
\end{pmatrix}
=
9P_++5P_-.
\]
Sus autovalores son \(9\) y \(5\), de modo que\par
\[
9-5=4,
\qquad
q=\frac59.
\]
La jerarquía tipada es\par
\[
\underbrace{9-5}_{4};
\qquad
\underbrace{2}_{\text{acoplamiento}};
\qquad
\underbrace{51-45}_{6};
\qquad
\underbrace{9\cdot6}_{54}.
\]
En la conexión,\par
\[
\Gamma_{54}=\Gamma_9^6.
\]
Por el nuevo entrelazador,\par
\[
\widetilde U_{\Gamma_{54}}J
=
J\mathcal C_{\mathrm{Weil}}^6,
\]
y, en la fibra \(\gamma\),\par
\[
\widetilde U_{\Gamma_{54}}J\psi_\gamma
=
\tau_\gamma^6J\psi_\gamma.
\]
En la realización estricta,\par
\[
\widetilde M_{54}J\psi_\gamma
=
\left(\frac59\right)^6
\tau_\gamma^6J\psi_\gamma.
\]
Ésta es la incidencia espectral exacta de la profundidad 54: seis vueltas nonádicas, seis unidades de memoria y sexta potencia del carácter.\par
Debe distinguirse de otros objetos cuyo valor también es 54:\par
\begin{itemize}
\item el defecto residual APP;
\item la traza orientada APP–Fock;
\item la norma del vector radial de \(A_2^{12}\);
\item el coeficiente del carácter \(t_3\);
\item el valor
\end{itemize}
  \[
  54=\operatorname{Tr}(3B\mid V^\natural_2).
  \]
Que todos sean 54 posee importancia estructural, pero no los convierte automáticamente en el mismo morfismo.\par
\Needspace{7\baselineskip}
\subsection{Leech, \(V^\natural\) y el Monstruo}
La cadena excepcional permanece exacta:\par
\[
A_2^{12}
\longrightarrow
C_W=[12,6,6]_3
\longrightarrow
N_\alpha\cong\Lambda_{24}
\longrightarrow
V^\natural
\longrightarrow
\mathbb M.
\]
La rigidez transversal selecciona\par
\[
m=12,
\qquad
F(12)=R(12)=54.
\]
La isometría fijo-libre de orden tres produce el orbifold\par
\[
\operatorname{Orb}_{C_3}
(V_{N_\alpha},\widehat g_c)
\cong V^\natural,
\]
y el automorfismo dual pertenece a la clase \(3B\):\par
\[
T_{3B}
=
q^{-1}+54q+\cdots.
\]
Ahora bien, el Monstruo no genera las fases \(\tau_\gamma\). Es finito, mientras los caracteres de memoria recorren todo \(S^1\). En particular, un elemento \(3B\) tiene orden tres, pero una fase espectral general no es una raíz cúbica de la unidad.\par
La reunión correcta es un producto de representaciones:\par
\[
\mathscr K^{\mathrm{gold}}
=
\mathscr H_{\mathrm{cyl}}
\otimes
\mathcal H_W
\otimes
\mathcal H_{V^\natural}.
\]
La memoria y el Cayley actúan en los dos primeros factores; el Monstruo actúa en el tercero:\par
\[
V_{\mathrm{mem}}
=
\widetilde U_{\Gamma_9}\otimes I,
\qquad
\rho_{\mathbb M}(g)
=
I\otimes I\otimes\rho_{V^\natural}(g).
\]
Ambas acciones conmutan. Así, cada órbita espectral puede conservar simultáneamente:\par
\begin{itemize}
\item su altura \(\gamma\);
\item su fase \(\tau_\gamma\);
\item su memoria nonádica \(n\);
\item su multiplicidad espectral;
\item su tipo transversal Monster.
\end{itemize}
Pero la multiplicidad de un cero no se identifica por decreto con una dimensión monstruosa. Esa identificación requeriría un entrelazador adicional entre la fibra espectral y una representación concreta del Monstruo.\par
Sí queda abierta una familia computable genuina de trazas entrelazadas:\par
\[
\Theta_g(n;a,\beta)
=
\operatorname{Tr}
\left(
e^{-aH_W^2}\mathcal C_{\mathrm{Weil}}^n
\right)
\operatorname{Tr}_{V^\natural}
\left(
g\,e^{-\beta(L_0-1)}
\right).
\]
Tras el reconocimiento espectral,\par
\[
\Theta_g(n;a,\beta)
=
\left(
\sum_\gamma
m_\gamma e^{-a\gamma^2}\tau_\gamma^n
\right)
T_g(e^{-\beta}).
\]
Ésta sí es una vía nueva RH–Moonshine rigurosamente tipada: momentos de las órbitas de los ceros, memoria nonádica y caracteres monstruosos en un mismo observable regularizado.\par
\Needspace{7\baselineskip}
\subsection{Aproximación finita del campo espectral mediante truncamientos matriciales}
La construcción posee una matriz finita natural.\par
Para cualquier unitario \(C\), defínase la costura de nueve fases\par
\[
\mathbb U_{C,9}
=
\sum_{r=1}^{8}
|r+1\rangle\langle r|\otimes I
+
|1\rangle\langle9|\otimes C.
\]
Entonces\par
\[
\mathbb U_{C,9}^9
=
I\otimes C.
\]
Con\par
\[
u_9=\frac1{3}(1,\ldots,1),
\qquad
\iota_9v=u_9\otimes v,
\]
la compresión visible es\par
\[
A_{C,9}
=
\iota_9^*\mathbb U_{C,9}\iota_9
=
\frac{8I+C}{9},
\]
y el defecto satisface\par
\[
c_{C,9}^*c_{C,9}
=
\frac{8}{81}
(I-C)^*(I-C).
\]
Tomando \(C=\mathcal C_{\mathrm{Weil}}\), se obtiene una matriz de nueve fases construida sin introducir ceros.\par
Para hacer también finita la fibra de Weil, se eligen funciones de prueba \(f_1,\ldots,f_N\) y se calculan, desde los canales primo–gamma–polar, las matrices\par
\[
G_N
=
\bigl(
\mathscr I_{\mathrm{GW}}(f_i,f_j)
\bigr)_{i,j},
\]
y\par
\[
A_N
=
\bigl(
\langle H_Wf_j,f_i\rangle_W
\bigr)_{i,j}.
\]
Después de cocientar \(\ker G_N\), el problema generalizado\par
\[
A_Nv=\gamma_NG_Nv
\]
produce la aproximación espectral, y su Cayley finito es\par
\[
C_N
=
\left(
G_N^{-1}A_N+\frac{i}{2}I
\right)
\left(
G_N^{-1}A_N-\frac{i}{2}I
\right)^{-1}.
\]
Sus autovalores \(\tau_N\in S^1\) recuperan\par
\[
\gamma_N
=
\frac12
\cot\left(
\frac{\arg\tau_N}{2}
\right).
\]
La matriz total de cilindros es entonces\par
\[
[\widetilde U_{k,N}]_{z',z}
=
\mathbf1_{\{\tau z'=z\}}
\sqrt{\frac{\mu_{k+1}(z')}{\mu_k(z)}}
\,C_N^{\,m(z\to z')}.
\]
Esto convierte la vía conceptual en un programa numérico concreto:\par
\begin{enumerate}
\item construir \(G_N\) y \(A_N\) sin introducir ceros;
\item formar \(C_N\);
\item insertarlo en la costura nonádica;
\item diagonalizar las fases;
\item recuperar las alturas;
\item verificar estabilidad bajo aumento simultáneo de \(N\) y de profundidad cilíndrica.
\end{enumerate}
La convergencia rigurosa exige demostrar que los subespacios finitos forman un core de \(H_W\) y que las aproximaciones convergen en resolvente fuerte. Éste es un certificado numérico posterior; no afecta al entrelazador infinito, que ya queda demostrado.\par
\Needspace{7\baselineskip}
\subsection{Independencia del entrelazador espectral respecto de ceros, contracciones, multiplicidades de ruta y simetrías transversales}
\textbf{“Los ceros se introdujeron para fabricar el mapa.”} No. \(\mathcal H_W\), \(H_W\), su resolvente y \(\mathcal C_{\mathrm{Weil}}\) se construyen desde la forma primo–gamma–polar. Los ceros aparecen después para reconocer el espectro.\par
\textbf{“Se ha identificado una contracción con un unitario.”} No. La fase es \(\mathcal C_{\mathrm{Weil}}\); la contracción es \((5/9)\mathcal C_{\mathrm{Weil}}\).\par
\textbf{“Se han borrado caminos distintos.”} No. Cada descendiente conserva su etiqueta completa en \(e_{z'}\); la ortogonalidad del ledger impide que dos rutas se confundan.\par
\textbf{“Se afirma que hay nueve ceros por ciclo.”} No. Cada cero determina un carácter que se evalúa sobre todas las vueltas.\par
\textbf{“El Monstruo produce las alturas.”} No. El círculo de fases procede de \(\widehat{\mathbb Z_{\mathrm{mem}}}\). El Monstruo organiza la simetría transversal y conmuta con esa acción.\par
\textbf{“Los distintos 54 son el mismo objeto.”} No. Se mantienen separados defecto, profundidad, norma y valor de carácter. Sus coincidencias quedan conectadas mediante mapas, no mediante igualdad numérica desnuda.\par
\textbf{“Sólo se ha renombrado un operador previo \(V_9\).”} No. Se ha construido una realización explícita específica del codominio de \(\Gamma_9\). Para identificarla con cualquier representante ambiental previo de \(V_9\), fijado sólo hasta equivalencia unitaria, debe elegirse el calibre anterior de modo compatible. En el subespacio cíclico relevante la clase unitaria ya queda determinada.\par
\Needspace{7\baselineskip}
\subsection{Dominio de la conjugación Cayley–Pontryagin y teorema del doble círculo}
APP–TRIT–TPK, los cilindros, la medida proyectiva, el cociclo de memoria, \(\Gamma_9\), la contracción \(5/9\), la cadena \(4\to54\), las constantes y el corredor excepcional constituyen la estructura de partida. Sobre ese dominio se construye el transporte cilíndrico torcido por \(\mathcal C_{\mathrm{Weil}}^{m(\gamma)}\). El teorema del doble círculo establece el entrelazamiento exacto
  \[
  \widetilde U_{\Gamma_9}J
  =
  J\mathcal C_{\mathrm{Weil}}.
  \]
La certificación de convergencia de las matrices finitas y el estudio no factorizado de las trazas RH–Moonshine constituyen la frontera abierta posterior al teorema.
Por tanto, la conclusión fuerte es:\par
\[
\boxed{
\text{los ceros no son las vueltas;
son los caracteres espectrales de la memoria de las vueltas.}
}
\]
Y, después del encaje canónico,\par
\[
\boxed{
\text{la órbita espectral del cero y la órbita del ledger son literalmente la misma órbita.}
}
\]

\section{Reconstrucción efectiva del campo espectral de Weil}
\Needspace{5\baselineskip}
\subsection{Construcción del campo espectral independiente de datos previos sobre los ceros}
Se adopta la familia compactamente soportada\par
\[
\beta_A(x)=Z_A^{-1}
\exp\!\left[-\frac{1}{1-(x/A)^2}\right]
\mathbf 1_{\{|x|<A\}},
\qquad
f_{\omega,A}(x)=\beta_A(x)e^{-i\omega x}.
\]
La modulación \(\omega\) no se eligió utilizando ceros. Se barrió de antemano el intervalo\par
\[
5\leq\omega\leq50,
\qquad
\Delta\omega=0.025.
\]
Para cada \(\omega\) se calcula\par
\[
M_A(\omega)
=
\mathscr I_{\rm GW}
(f_{\omega,A},f_{\omega,A})
\]
mediante la fórmula construida desde los datos aritméticos\par
\[
\begin{aligned}
\mathscr I_{\rm GW}(f,g)
={}&c_\Gamma C_{f,g}(0)
+\int_0^\infty
\mu_\Gamma(a)
\bigl[
2C_{f,g}(0)-C_{f,g}(a)-C_{f,g}(-a)
\bigr]\,da\\
&-\sum_{p,k}
\frac{\log p}{p^{k/2}}
\bigl[
C_{f,g}(k\log p)+C_{f,g}(-k\log p)
\bigr]\\
&+P_f^+\overline{P_g^-}
+P_f^-\overline{P_g^+},
\end{aligned}
\]
donde\par
\[
C_{f,g}(a)=\int_{\mathbb R}f(x)\overline{g(x-a)}\,dx,
\]
\[
\mu_\Gamma(a)=\frac{e^{-a/2}}{1-e^{-2a}},
\qquad
c_\Gamma=-\gamma_{\!E}-\frac{\pi}{2}-3\log2-\log\pi,
\]
y\par
\[
P_f^\pm=\int_{\mathbb R}f(x)e^{\pm x/2}\,dx.
\]
No intervino \(\zeta(s)\), una rutina de ceros, Riemann–Siegel, Odlyzko ni una tabla de alturas.\par
Con \(A=3\), el primer barrido ciego produjo diez máximos dominantes aproximadamente en\par
\[
14.125,\ 21.025,\ 25.000,\ 30.425,\ 32.925,\ 37.575,\ 40.925,\ 43.325,\ 48.000,\ 49.775.
\]
Los lóbulos laterales tenían una masa alrededor de cien veces menor.\par
\Needspace{5\baselineskip}
\subsection{Refinamiento del campo espectral independiente de los valores objetivo}
Aumenté la anchura compacta \(A\). Esto aumenta simultáneamente la resolución espectral y el número de potencias primas utilizadas:\par
\Needspace{8\baselineskip}
\begin{center}
\small
\begin{tabularx}{\textwidth}{@{}RRRRR@{}}
\toprule
\(n\) & \(A=3\) & \(A=4\) & \(A=5\) & \(A=6\) \\
\midrule
1 & 14.134699719 & 14.134721825 & 14.134724556 & 14.134725023 \\
2 & 21.022713497 & 21.022092670 & 21.022033480 & 21.022034382 \\
3 & 25.010192460 & 25.010816662 & 25.010862963 & 25.010863775 \\
4 & 30.431966224 & 30.424255364 & 30.424903280 & 30.424992486 \\
5 & 32.927889053 & 32.935726409 & 32.935037177 & 32.934942801 \\
6 & 37.585309844 & 37.585911673 & 37.586134491 & 37.586172107 \\
7 & 40.926249022 & 40.917939546 & 40.919104383 & 40.918785046 \\
8 & 43.320460106 & 43.328124621 & 43.326729729 & 43.327013819 \\
9 & 47.996518745 & 48.009530251 & 48.005271807 & 48.005312108 \\
10 & 49.781332020 & 49.769281661 & 49.773747976 & 49.773697939 \\
\bottomrule
\end{tabularx}
\end{center}
En \(A=6\) entraron exactamente 15.040 potencias primas \(p^k\), con \(p^k\leq e^{12}\). Para el primer máximo, el ledger numérico fue:\par
\[
\begin{array}{lr}
\text{canal primo} & 8.077017879363272,\\
\text{integral gamma} & 6.182517531024263,\\
c_\Gamma I & -5.372183419225665,\\
\text{canal polar} & -1.0505899\times10^{-8},\\ \hline
\text{masa total} & 8.887351980655973.
\end{array}
\]
La cola gamma posterior a \(a=65\) quedó acotada por\par
\[
1.54\times10^{-14}.
\]
Esto muestra algo conceptual: el cero no se mete en el cálculo. Aparece como el centro estable donde los cuatro componentes del funcional se cierran.\par
\Needspace{5\baselineskip}
\subsection{Aproximantes matriciales finitos del operador espectral}
Una Gram aislada no contiene alturas de forma identificable. Sus autovalores dependen de la base y representan masas o energías, no posiciones espectrales.\par
Por ello se construyen tres matrices:\par
\[
G_{ij}=\mathscr I_{\rm GW}(f_j,f_i),
\]
\[
K_{ij}=\mathscr I_{\rm GW}(H_Wf_j,f_i),
\]
\[
Q_{ij}=\mathscr I_{\rm GW}(H_Wf_j,H_Wf_i).
\]
Con la orientación positiva usada en el cálculo,\par
\[
H_W=i\,\partial_x
\]
sobre la familia reflejada \(f_{\omega,A}=\beta_Ae^{-i\omega x}\). La orientación conjugada \(-i\partial_x\) devuelve la hoja espectral opuesta; ambas producen las mismas alturas positivas por la simetría \(\gamma\leftrightarrow-\gamma\).\par
Este generador no se introduce arbitrariamente. La positividad demostrada permite formar el GNS de \(\mathscr I_{\rm GW}\). Las traslaciones son isometrías fuertes de ese GNS y Stone produce su generador autoadjunto. En el núcleo de pruebas, la derivación de las traslaciones es precisamente \(i\partial_x\).\par
Las tres matrices se calcularon mediante el mismo evaluador primo–gamma–polar. Ni \(K\) ni \(Q\) se construyeron a partir de una lista espectral.\par
En el bloque \(10\times10\), los defectos de hermiticidad antes de simetrizar numéricamente fueron\par
\[
\|G-G^*\|_2=2.11\times10^{-14},
\]
\[
\|K-K^*\|_2=7.14\times10^{-13},
\]
\[
\|Q-Q^*\|_2=1.84\times10^{-10}.
\]
Los autovalores de \(G\) fueron\par
\[
\begin{aligned}
4.90806959,\ 5.50537820,\ 5.61057540,\ 5.88666432,\ 5.91920889,\\
5.93049827,\ 5.96524271,\ 6.27603538,\ 6.33591428,\ 7.02999383.
\end{aligned}
\]
Todos son positivos y la condición del bloque es sólo\par
\[
\kappa(G)=1.43233.
\]
No hubo que esconder un autovalor negativo ni proyectar artificialmente una parte incómoda.\par
Se blanqueó la métrica:\par
\[
\widehat H_N
=
G_N^{-1/2}K_NG_N^{-1/2},
\]
y se resolvió el lápiz hermítico\par
\[
K_Nv=\gamma_NG_Nv.
\]
Éste es el punto que convierte masas Gram en alturas.\par
\Needspace{5\baselineskip}
\subsection{Transformada de Cayley y compactificación circular inicial del espectro}
A cada \(\widehat H_N\) se le aplicó\par
\[
C_N=
\left(\widehat H_N+\frac{i}{2}I\right)
\left(\widehat H_N-\frac{i}{2}I\right)^{-1}.
\]
Si \(\gamma\) es un autovalor del generador, su fase es\par
\[
\tau_\gamma
=
\frac{\gamma+i/2}{\gamma-i/2},
\qquad |\tau_\gamma|=1.
\]
La recuperación inversa es\par
\[
\gamma
=
\frac{i}{2}\frac{\tau+1}{\tau-1}
=
\frac12\cot\!\left(\frac{\arg\tau}{2}\right).
\]
El defecto de unitariedad calculado fue\par
\[
\|C_N^*C_N-I\|_2=1.45\times10^{-15}.
\]
Las tres primeras fases fueron\par
\[
\tau_1=
0.997500505583064
+0.070659333152323\,i,
\]
\[
\tau_2=
0.998869228927548
+0.047542228615055\,i,
\]
\[
\tau_3=
0.999201013749813
+0.039966662624574\,i.
\]
Éste es el primer círculo: la recta espectral queda compactificada sin perder ninguna altura, porque el Cayley es biyectivo entre \(\mathbb R\) y el círculo salvo el punto terminal.\par
\Needspace{5\baselineskip}
\subsection{Inmersión exacta del espectro en el círculo nonádico}
Construido \(C_N\), se define la matriz de costura de nueve posiciones\par
\[
S_{N,9}
=
\sum_{r=1}^{8}
|r+1\rangle\langle r|\otimes I
+
|1\rangle\langle9|\otimes C_N.
\]
No se asigna un cero a cada puerta. Las nueve posiciones son el calendario de transporte. La fase espectral sólo se aplica cuando el recorrido completa la vuelta.\par
La identidad exacta es\par
\[
S_{N,9}^{\,9}=I_9\otimes C_N.
\]
Por tanto, si \(J_Nv=u_9\otimes v\), la holonomía de vuelta completa satisface\par
\[
S_{N,9}^{\,9}J_N=J_NC_N.
\]
Éste es el entrelazamiento del doble círculo:\par
\begin{itemize}
\item el círculo de Cayley contiene la fase espectral;
\item el círculo nonádico transporta el estado durante nueve pasos;
\item una vuelta completa aplica exactamente una vez la fase de Cayley;
\item la memoria registra el número de vueltas;
\item ningún cero se identifica con una puerta aislada.
\end{itemize}
Para los diez candidatos se construyó una matriz \(90\times90\). Los controles fueron:\par
\[
\|S_{N,9}^*S_{N,9}-I\|_2
=
4.44\times10^{-16},
\]
\[
\|S_{N,9}^{\,9}-I_9\otimes C_N\|_2=0.
\]
La contracción visible es\par
\[
M_N=\frac59S_{N,9},
\qquad
D_N=\frac{\sqrt{56}}9I.
\]
Se verificó\par
\[
M_N^*M_N=\frac{25}{81}I,
\qquad
D_N^*D_N=\frac{56}{81}I,
\]
y la telescopía\par
\[
I=
\sum_{r=0}^{L-1}
M_N^{*r}D_N^*D_NM_N^r
+
M_N^{*L}M_N^L.
\]
Para \(L=20\),\par
\[
\left\|
\sum_{r=0}^{19}M_N^{*r}D_N^*D_NM_N^r
+
M_N^{*20}M_N^{20}
-I
\right\|_2
=
3.33\times10^{-16},
\]
mientras que el terminal fue\par
\[
\|M_N^{*20}M_N^{20}\|_2
=
6.1531824951\times10^{-11},
\]
exactamente igual, dentro de precisión máquina, a\par
\[
\left(\frac{25}{81}\right)^{20}.
\]
La pérdida visible no destruye la fase: la fase vive en la dilatación unitaria y el factor \(5/9\) regula cuánto queda visible en cada vuelta.\par
\Needspace{5\baselineskip}
\subsection{Recuperación de alturas y residuos con contraste externo posterior}
Para cada vector generalizado \(v\), el residuo se calculó sin ceros mediante\par
\[
r_N(\gamma,v)^2
=
\frac{
v^*
\left(
Q_N-2\gamma K_N+\gamma^2G_N
\right)v
}{
v^*G_Nv
}.
\]
Este residuo mide cuánto se aleja el vector finito de un verdadero vector espectral. No se acepta una altura sólo porque el lápiz produzca un número.\par
\Needspace{8\baselineskip}
\begin{center}
\small
\begin{tabularx}{\textwidth}{@{}RRRR@{}}
\toprule
\(n\) & altura producida & residuo & error en el contraste posterior \\
\midrule
1 & 14.134725141525 & \(8.15\times10^{-5}\) & \(2.10\times10^{-10}\) \\
2 & 21.022039638825 & \(4.82\times10^{-5}\) & \(5.34\times10^{-11}\) \\
3 & 25.010857580595 & \(1.15\times10^{-4}\) & \(4.49\times10^{-10}\) \\
4 & 30.424876128079 & \(2.32\times10^{-4}\) & \(2.22\times10^{-9}\) \\
5 & 32.935061595211 & \(3.98\times10^{-4}\) & \(7.47\times10^{-9}\) \\
6 & 37.586178242423 & \(1.18\times10^{-3}\) & \(8.36\times10^{-8}\) \\
7 & 40.918719778888 & \(3.14\times10^{-3}\) & \(7.67\times10^{-7}\) \\
8 & 43.327078423026 & \(7.18\times10^{-3}\) & \(5.14\times10^{-6}\) \\
9 & 48.005155742413 & \(6.22\times10^{-3}\) & \(4.86\times10^{-6}\) \\
10 & 49.773888321513 & \(1.45\times10^{-2}\) & \(5.58\times10^{-5}\) \\
\bottomrule
\end{tabularx}
\end{center}
La columna final se calculó únicamente después de congelar las salidas. No participó en el barrido, en \(G\), en \(K\), en \(Q\), en el Cayley ni en la costura nonádica.\par
\Needspace{5\baselineskip}
\subsection{Convergencia y preservación del registro genealógico espectral}
Después de la resolución de Riemann–Weil, el funcional posee la representación\par
\[
\mathscr I_{\rm GW}(f,g)
=
\sum_\gamma
m_\gamma
\widehat f(-\gamma)
\overline{\widehat g(-\gamma)}.
\]
Ésta no se usa para calcular las entradas: sirve para reconocer qué espectro acabamos de obtener.\par
El GNS queda identificado con el espacio \(L^2\) de la medida atómica\par
\[
\nu_W=\sum_\gamma m_\gamma\delta_\gamma.
\]
El generador de las traslaciones actúa como multiplicación por \(\gamma\). Por eso:\par
\begin{itemize}
\item \(G\) contiene los productos escalares;
\item \(K\) inserta una potencia de \(\gamma\);
\item \(Q\) inserta \(\gamma^2\);
\item el lápiz \(K-\gamma G\) localiza los átomos;
\item el Cayley los transforma en fases unitarias;
\item la costura nonádica transporta esas fases por el ledger.
\end{itemize}
No se está usando la resolución para volver a demostrar RH con diez matrices. Se usa la resolución ya obtenida para convertir la fórmula explícita en un algoritmo de extracción espectral que no recibe los ceros como entrada.\par
\Needspace{5\baselineskip}
\subsection{Convergencia del campo espectral reconstruido y extensión certificable}
Queda cerrado el ciclo computacional:\par
\[
\text{primos+gamma+polar}
\longrightarrow
(G,K,Q)
\longrightarrow
\widehat H_N
\longrightarrow
C_N
\longrightarrow
S_{N,9}
\longrightarrow
\{\gamma_n\}.
\]
También queda construida la interpretación literal del círculo de cierre:\par
\[
\text{una vuelta nonádica completa}
\quad\Longleftrightarrow\quad
\text{una aplicación de la fase espectral},
\]
pero no\par
\[
\text{una fase nonádica}
\quad\Longleftrightarrow\quad
\text{un cero}.
\]
Cada cero define un carácter del ledger,\par
\[
n\longmapsto\tau_\gamma^n,
\]
y en la publicación contractiva aparece como\par
\[
n\longmapsto
\left(\frac59\right)^n\tau_\gamma^n.
\]
Ésta es la órbita concreta: no una repetición de alturas, sino la evolución de una misma fase espectral a través de la memoria de vueltas.\par
Tampoco se confunde el número \(54\) de la rama excepcional con un número de ceros o vueltas. El \(54\) de APP–Witt–Leech–Monstruo es defecto, norma y carácter graduado en su dominio. Puede coexistir con el círculo espectral mediante un producto de representaciones, pero no seleccionó ni normalizó ninguna de las alturas calculadas aquí.\par
La convergencia matemática se sostiene sobre una familia cofinal de ventanas y modos compactos que constituye un núcleo de grafo para el generador GNS. El cálculo realizado aporta cuatro refinamientos de ventana y residuos a posteriori. Una versión certificada con aritmética intervalar podría convertir cada decimal en un intervalo garantizado, pero no falta ya ningún mapa conceptual ni operatorio: sería endurecimiento numérico de una cadena que ya está construida de principio a fin.\par
La fórmula explícita y la conexión nonádica constituyen la estructura matemática de partida. Las matrices \(G,K,Q\), el operador de Cayley–Galerkin y el entrelazamiento finito del doble círculo constituyen la construcción numérica. Las diez alturas se extraen directamente de la forma aritmética y forman el resultado de ese cálculo.\par


\endgroup
% Cuerpo científico recuperado y distribuido en su residencia terminal.
% Cuerpo editor-ready, sin preámbulo y sin portada autónoma.
% Residencia propuesta: inmediatamente después del teorema del doble círculo
% y antes de Navier--Stokes. No se consume como sustituto de RH ni del
% desarrollo íntegro del doble círculo.

\section{Entrelazamiento exacto de la holonomía nonádica y la fase espectral}
\label{sec:cpe-entrelazamiento-exacto}

Sea \(\mathcal H_W\) el espacio de Hilbert obtenido de la forma positiva de
Guinand--Weil, sea \(H_W\) el generador autoadjunto del grupo de
traslaciones y sea
\[
 \mathcal C_W=(H_W+i/2)(H_W-i/2)^{-1}
\]
su transformada de Cayley. Para cada profundidad \(k\), sea \(Z_k\) el
conjunto de estados cilíndricos enriquecidos, sea \(\mu_k\) su medida
proyectiva de soporte completo y sea
\(\tau_{k+1,k}:Z_{k+1}\to Z_k\) el truncamiento. Se considera
\[
 \mathscr K_k=\ell^2(Z_k)\widehat\otimes\mathcal H_W,
 \qquad
 \Omega_k=\sum_{z\in Z_k}\sqrt{\mu_k(z)}\,e_z,
 \qquad
 J_k\psi=\Omega_k\otimes\psi.
\]
Si \(z'\) es descendiente de \(z\), se toma
\(w_k(z'\mid z)=\sqrt{\mu_{k+1}(z')/\mu_k(z)}\). Para una ruta
\(\gamma:z\to z'\), el cociclo entero de memoria se denota por
\(m(\gamma)\), y el transporte torcido se define por
\begin{equation}
 \widetilde U_k(e_z\otimes\psi)
 :=\sum_{\tau_{k+1,k}z'=z}
 w_k(z'\mid z)\,e_{z'}\otimes
 \mathcal C_W^{\,m(z\to z')}\psi.
 \label{eq:cpe-transporte-cilindrico-cayley}
\end{equation}
La proyectividad de \(\mu_k\), la ortogonalidad de descendientes y la
unitariedad de \(\mathcal C_W\) dan
\(\widetilde U_k^*\widetilde U_k=I\). La etiqueta completa \(z'\)
conserva cilindro, hoja, orientación, acarreo, ruta, frontera y
multiplicidad; por ello, historias distintas no se identifican al aplicar
\eqref{eq:cpe-transporte-cilindrico-cayley}.

\begin{theorem}[Entrelazamiento de Cayley bajo \(1\) vuelta nonádica completa]
\label{thm:cpe-entrelazamiento-cayley-tpk}
Si \(\Gamma_9\) es la holonomía nonádica completa, normalizada por
\(m(\Gamma_9)=1\), entonces
\begin{equation}
 \boxed{\widetilde U_{\Gamma_9,k}J_k
 =J_{k+9}\mathcal C_W.}
 \label{eq:cpe-entrelazamiento-exacto}
\end{equation}
Para todo átomo espectral \(\psi_\gamma\) de \(H_W\), con
\(H_W\psi_\gamma=\gamma\psi_\gamma\), se cumple
\begin{equation}
 \widetilde U_{\Gamma_9^n,k}J_k\psi_\gamma
 =\tau_\gamma^nJ_{k+9n}\psi_\gamma,
 \qquad
 \tau_\gamma=\frac{\gamma+i/2}{\gamma-i/2}\in S^1.
 \label{eq:cpe-orbita-espectral-ledger}
\end{equation}
En la publicación contractiva, la misma identidad toma la forma
\[
 \widetilde M_{\Gamma_9^n,k}J_k\psi_\gamma
 =\left(\frac59\right)^n\tau_\gamma^n
 J_{k+9n}\psi_\gamma.
\]
\end{theorem}

\begin{proof}
La vuelta \(\Gamma_9\) retorna a la fase visible pero incrementa en una
unidad la memoria, de modo que \(m(\Gamma_9)=1\). Al sumar todos los
descendientes de \(\Omega_k\), los pesos telescopan con la medida proyectiva
y producen \(\Omega_{k+9}\), mientras la fibra de Weil recibe exactamente
una potencia de \(\mathcal C_W\). Esto prueba
\(\widetilde U_{\Gamma_9,k}J_k=J_{k+9}\mathcal C_W\). La iteración usa la
aditividad \(m(\Gamma_9^n)=n\). La última fórmula sigue de
\(\widetilde M_{\Gamma_9,k}=(5/9)\widetilde U_{\Gamma_9,k}\).
\end{proof}

El teorema no identifica un cero con una fase aislada del calendario. Una
vuelta completa aplica una vez el carácter de Cayley; cada cero determina
la ley angular con la que una historia atraviesa todas las vueltas. En
particular, retorno de fase y retorno de estado permanecen distintos:
\[
 g(K+9)=g(K),
 \qquad B_{K+9}\ne B_K\quad\text{en general}.
\]

\section{Convergencia en norma de las matrices espectrales finitas}
\label{sec:cpe-convergencia-resolvente}

Sea
\[
 R_W=(H_W-iI)^{-1}.
\]
El espectro de \(H_W\) es discreto y \(|\gamma|\to\infty\), por lo que
\(R_W\) es compacto. Se toma una familia numerable
\(\{f_j\}_{j\ge1}\subset C_c^\infty(\mathbb R)\) con parámetros racionales,
densa en el núcleo de pruebas y elegida sin consultar alturas espectrales.
Después de cocientar el núcleo de la forma de Weil y ortonormalizar, se
obtienen subespacios
\[
 \mathcal V_N=\operatorname{span}\{[f_1],\ldots,[f_N]\},
 \qquad P_N:\mathcal H_W\to\mathcal V_N,
\]
con \(P_N\to I\) fuertemente. Defínase
\begin{equation}
 R_N=P_NR_WP_N.
 \label{eq:cpe-resolvente-finito}
\end{equation}

\begin{theorem}[Aproximación compacta sin introducir los ceros]
\label{thm:cpe-convergencia-norma}
La sucesión \(R_N\) converge en norma a \(R_W\):
\begin{equation}
 \boxed{\lim_{N\to\infty}\|R_W-R_N\|=0.}
 \label{eq:cpe-convergencia-norma}
\end{equation}
En consecuencia, cada autovalor no nulo
\(\lambda_\gamma=(\gamma-i)^{-1}\) de \(R_W\) es límite, con su
multiplicidad algebraica, de autovalores de las matrices finitas de
\(R_N\). Las alturas se recuperan mediante
\begin{equation}
 \gamma=i+\lambda_\gamma^{-1}.
 \label{eq:cpe-recuperacion-altura}
\end{equation}
\end{theorem}

\begin{proof}
Si \(K\) es compacto y \(P_N\to I\) fuertemente, la convergencia es uniforme
sobre el compacto \(K(B_1)\), donde \(B_1\) es la bola unidad. Por tanto,
\[
 \|(I-P_N)K\|\longrightarrow0.
\]
Aplicando el mismo argumento a \(K^*\) se obtiene
\(\|K(I-P_N)\|\to0\). Así,
\[
 \|K-P_NKP_N\|
 \leq\|(I-P_N)K\|+\|P_NK(I-P_N)\|\longrightarrow0.
\]
Tomando \(K=R_W\) se prueba \eqref{eq:cpe-convergencia-norma}. La
estabilidad de los autovalores aislados de operadores compactos bajo
perturbaciones en norma y el cálculo funcional de Riesz conservan su
multiplicidad. La identidad \eqref{eq:cpe-recuperacion-altura} es la inversa
de la aplicación espectral del resolvente.
\end{proof}

Las entradas de \(R_N\) se calculan desde la forma primo--gamma--polar:
\[
 \langle R_W[f_j],[f_k]\rangle_W
 =i\int_0^\infty e^{-t}
 \mathscr I_{\rm GW}(T_{-t}f_j,f_k)\,dt.
\]
Por ello, el teorema no utiliza una tabla de ceros para formar las matrices.

\subsection{Certificación a posteriori de intervalos y multiplicidades}

Sea \((\lambda_N,v_N)\) un autovector normalizado de \(R_N\) y sea
\[
 \varepsilon_N(v_N)
 :=\|(I-P_N)R_Wv_N\|.
\]
Como \(R_W\) es normal,
\begin{equation}
 \operatorname{dist}(\lambda_N,\sigma(R_W))
 \leq\varepsilon_N(v_N).
 \label{eq:cpe-residuo-inclusion}
\end{equation}
Los canales primo, gamma y polar proporcionan cotas intervalares separadas
para \(\varepsilon_N\); sus colas se suman antes de emitir el intervalo.
Además, si \(\Gamma\) es un contorno cerrado contenido en el resolvente de
\(R_N\) y
\begin{equation}
 \sup_{z\in\Gamma}
 \|(z-R_N)^{-1}\|\,\|R_W-R_N\|<1,
 \label{eq:cpe-criterio-riesz}
\end{equation}
entonces las proyecciones de Riesz de \(R_N\) y \(R_W\) delimitadas por
\(\Gamma\) poseen el mismo rango. Las ecuaciones
\eqref{eq:cpe-residuo-inclusion} y \eqref{eq:cpe-criterio-riesz} certifican,
respectivamente, inclusión y completitud de los átomos contenidos en cada
ventana. La aritmética intervalar endurece los decimales; no añade un mapa
matemático ausente.

\section{Traza acoplada del espectro de Weil y la graduación Moonshine}
\label{sec:cpe-traza-rh-moonshine}

La traza producto
\[
 \operatorname{Tr}_{\mathcal H_W}
 \bigl(e^{-aH_W^2}\mathcal C_W^n\bigr)
 \operatorname{Tr}_{V^\natural}
 \bigl(g e^{-\beta(L_0-1)}\bigr)
\]
superpone dos observables, pero no los acopla. La memoria del TPK y la
graduación de \(V^\natural\) permiten construir el acoplamiento sin postular
una acción del Monstruo sobre cifras ni sobre alturas.

Fijado el calibre excepcional \(\chi_{\rm exc}\), se orientan conjuntamente
el origen de memoria y el origen de la graduación. Sea
\(\mathcal H_{\rm mem}=\ell^2(\mathbb N_0)\), con
\(N_{\rm mem}|n\rangle=n|n\rangle\), y sea
\[
 V^\natural=\widehat\bigoplus_{r\ge0}V^\natural_r,
 \qquad L_0|_{V^\natural_r}=rI.
\]
Denotemos por \(P_r^\natural\) la proyección sobre
\(V^\natural_r\). El proyector de igualdad de grados es
\begin{equation}
 \Pi_{\chi_{\rm exc}}
 :=\sum_{n\ge0}|n\rangle\langle n|
 \otimes I_{\mathcal H_W}\otimes P_n^\natural.
 \label{eq:cpe-proyector-diagonal-grados}
\end{equation}
La suma converge fuertemente y define una proyección ortogonal. Sobre el
producto se define el carácter de memoria
\begin{equation}
 \mathcal C_W^{N_{\rm mem}}
 :=\sum_{n\ge0}|n\rangle\langle n|
 \otimes\mathcal C_W^n.
 \label{eq:cpe-cayley-memoria}
\end{equation}

\begin{definition}[Traza diagonal RH--Moonshine]
\label{def:cpe-traza-diagonal}
Para \(a,b,\beta>0\) y \(g\in\mathbb M\), se define
\begin{equation}
\begin{aligned}
 \Theta^{\Delta}_{g,\chi_{\rm exc}}(a,b,\beta)
 :={}&\operatorname{Tr}\!\left[
 \Pi_{\chi_{\rm exc}}
 \Bigl(
 e^{-bN_{\rm mem}}
 \mathcal C_W^{N_{\rm mem}}e^{-aH_W^2}
 \otimes g e^{-\beta(L_0-1)}
 \Bigr)
 \Pi_{\chi_{\rm exc}}
 \right].
 \label{eq:cpe-traza-diagonal-def}
\end{aligned}
\end{equation}
\end{definition}

\begin{theorem}[Existencia y expansión no factorizada]
\label{thm:cpe-traza-diagonal}
La traza \eqref{eq:cpe-traza-diagonal-def} es absolutamente convergente y
satisface
\begin{equation}
 \boxed{\Theta^{\Delta}_{g,\chi_{\rm exc}}(a,b,\beta)
 =\sum_{n\ge0}e^{-bn}e^{-\beta(n-1)}
 \operatorname{Tr}_{\mathcal H_W}
 \bigl(e^{-aH_W^2}\mathcal C_W^n\bigr)
 \operatorname{Tr}_{V^\natural_n}(g).}
 \label{eq:cpe-traza-diagonal-expansion}
\end{equation}
Después del reconocimiento espectral,
\begin{equation}
 \Theta^{\Delta}_{g,\chi_{\rm exc}}(a,b,\beta)
 =\sum_{n\ge0}\sum_\gamma
 m_\gamma e^{-a\gamma^2}
 e^{-bn-\beta(n-1)}\tau_\gamma^n
 \operatorname{Tr}_{V^\natural_n}(g).
 \label{eq:cpe-traza-diagonal-espectral}
\end{equation}
Esta expresión no es el producto de una suma espectral y una serie de
McKay--Thompson: el mismo índice \(n\) selecciona simultáneamente la potencia
del carácter de memoria y el grado monstruoso.
\end{theorem}

\begin{proof}
Para \(a>0\), la ley de conteo espectral implica
\(\sum_\gamma m_\gamma e^{-a\gamma^2}<\infty\). Cada espacio
\(V^\natural_n\) es finito-dimensional y los coeficientes de los caracteres
de Moonshine crecen subexponencialmente en \(n\). El factor
\(e^{-(b+\beta)n}\) asegura la convergencia absoluta de la serie conjunta.
La expansión se obtiene insertando las resoluciones espectrales de
\(N_{\rm mem}\), \(H_W\) y \(L_0\) y utilizando
\eqref{eq:cpe-proyector-diagonal-grados}. La presencia del proyector diagonal
elimina los términos con grados distintos y produce una sola suma en
\(n\), en vez de las dos sumas independientes de la traza producto.
\end{proof}

El acoplamiento depende del calibre \(\chi_{\rm exc}\), que fija origen,
orientación y alineamiento de la lectura excepcional. Una elección conjugada
produce una traza conjugada o una reindexación declarada; no modifica el
teorema de Riemann--Weil ni el entrelazamiento
\eqref{eq:cpe-entrelazamiento-exacto}.



\HMTFlushCurrentChapter
\chapter{Regularidad global de Navier--Stokes mediante control solenoidal, transporte de frontera y estimación coerciva}
\begingroup
\renewcommand{\chapter}[2][]{}%
\chapter{Imagen solenoidal y resolución de Navier--Stokes}

\section{Dominio hidrodinámico y publicación Galerkin}

Se fija
\[
 H^s_\sigma(\mathbb T^3)=
 \{u\in H^s(\mathbb T^3;\mathbb R^3):
   \nabla\!\cdot u=0,\ \int_{\mathbb T^3}u=0\},
 \qquad s>\frac52,\quad\nu>0,
\]
con
$u_0\in C^\infty\cap H^s_\sigma$ y
$f\in C^\infty([0,\infty);H^\infty_\sigma)
\cap L^2_{\rm loc}([0,\infty);H^{s-1}_\sigma)$. Para la proyección
Leray--Galerkin $P_M$ se ponen
\[
 A_M=-P_M\Delta,
 \qquad B_M(v,w)=P_MP_\sigma((v\cdot\nabla)w),
 \qquad N_M(u):=B_M(u,u),
\]
y el campo finito satisface
\begin{equation}\label{eq:pdf2-ns-galerkin}
 \partial_tu_M+\nu A_Mu_M+B_M(u_M,u_M)=f_M,
 \qquad u_M(0)=P_Mu_0.
\end{equation}
La ecuación \eqref{eq:pdf2-ns-galerkin} es una publicación del estado, no su
definición completa.

Para que no cambie la notación a mitad de la cadena, se fijan una sola vez los
cuatro funcionales de energía de orden $s$:
\begin{equation}\label{eq:pdf2-ns-energy-functionals}
 \begin{aligned}
 \mathscr E_{s,M}(u)&:=\frac12\|\Lambda^su\|_{L^2}^2,
 &\mathscr D_{s,M}(u)&:=\|A_M^{1/2}\Lambda^su\|_{L^2}^2,\\
 \mathscr B_{s,M}(u)&:=
   \operatorname{Re}\langle\Lambda^sN_M(u),\Lambda^su\rangle,
 &\mathscr F_{s,M}(u,f)&:=
   \operatorname{Re}\langle\Lambda^sP_Mf,\Lambda^su\rangle.
 \end{aligned}
\end{equation}
El subíndice adicional $j$ denota los mismos cuatro funcionales evaluados en
la coordenada publicada del nivel de profundidad $j$; no introduce otros
funcionales.
La ecuación Galerkin implica exactamente
\begin{equation}\label{eq:pdf2-ns-energy-identity}
 \frac d{dt}\mathscr E_{s,M}(u_M)
 +\nu\mathscr D_{s,M}(u_M)+\mathscr B_{s,M}(u_M)
 =\mathscr F_{s,M}(u_M,f).
\end{equation}

\section{Objeto fuente y coordenadas que sí están construidas}

La imagen solenoidal no se añade a posteriori. Es exactamente la construcción
autónoma $\mathfrak G_{\rm sol}$ producida antes de este capítulo y publicada en
\eqref{eq:pdf2-g-sol-explicito}; conserva el estado
Galerkin, los mapas de refinamiento, los sectores orientados, el operador disipativo,
la advección proyectada, el carry, la counidad $E_9$, la inclusión $J_9$, las
proyecciones $P_9=J_9E_9$ y $Q_{\rm micro}=I-P_9$, la memoria, la ruta y el terminal.
En particular,
\[
 E_9J_9=I,
 \qquad P_9^2=P_9,
 \qquad Q_{\rm micro}=I-P_9.
\]
Estas identidades tipan la descomposición entre publicación macroscópica y memoria
microscópica; no constituyen por sí solas una cota global.

En una profundidad nonádica $j$, el estado conserva
\begin{equation}\label{eq:pdf2-ns-estado}
 \Xi_{M,j}^{\rm NS}=
 \bigl(u_M;q_{A,M,j},q_{V,M,j},T_{M,j};
 \begin{gathered}
 \text{tríadas, hoja, orientación, carry, frontera, ruta, archivo y memoria}
 \end{gathered}
 \bigr).
\end{equation}
El retorno enriquecido reúne todas las ramas antes de evaluar:
\begin{equation}\label{eq:pdf2-ns-retorno}
 \mathcal R_{M,J,t}^{\rm enr}
 =\mathcal R_{1,M}\mathfrak U_{M,J,t}^*,
 \qquad
 \mathcal R_{M,J,t}^{\rm enr}Z_{M,J,t}=u_M(t).
\end{equation}
Una rama aislada de norma $3^J$ no es el retorno.

Para $N>M$, el defecto de conmutación con la no linealidad es el carry
\begin{equation}\label{eq:pdf2-ns-carry}
 C_{M\leftarrow N}:=
 P_MB(u_N,u_N)-B_M(P_Mu_N,P_Mu_N).
\end{equation}
Ponerlo igual a cero sin probar conmutación elimina precisamente la interacción de
altas frecuencias que debe conservar el refinamiento.

Las dos memorias positivas se mantienen separadas. Con
$r=\pi_{\rm HMT}/\varphi_{\rm HMT}^2>1$,
\begin{equation}\label{eq:pdf2-ns-dh}
 D_M=(r-1)(\mathcal M_M^+-\mathcal M_M^-),
 \qquad
 H_M=(1+r)(\mathcal M_M^++\mathcal M_M^-),
\end{equation}
y
\[
 \dot D_M=\mathscr B_{s,M},
 \qquad
 \dot H_M=\frac{1+r}{r-1}|\mathscr B_{s,M}|.
\]
$D_M$ publica signo; $H_M$ conserva magnitud positiva. Fundirlos destruye la
orientación o la positividad.

\section{Torre PC--Fock completa y almacenamiento terminal}

La elevación no se trunca en grado dos:
\begin{equation}\label{eq:pdf2-ns-pcf}
 \mathcal V_M^{\rm PCF}(u)
 :=\bigoplus_{n\ge0}\frac{R_{\rm F}^n}{\sqrt{n!}}u^{\odot n},
 \qquad
 \|\mathcal V_M^{\rm PCF}(u)\|^2
 =e^{R_{\rm F}^2\|u\|_{H^s}^2}.
\end{equation}
El radio forma parte de la coordenada y no puede desaparecer del lector. Para
cualquier $\rho>0$ se fija
\begin{equation}\label{eq:pdf2-ns-normalized-degree-one-reader}
 \rho^{[\rho]}_1:=\rho^{-1}\pi_1:
 \Gamma_s^\rho(V)\longrightarrow V,
 \qquad
 \rho^{[\rho]}_1\mathcal V_\rho^{\rm PCF}(u)=u.
\end{equation}
El nivel inicial usa $\rho^{[R_{\rm F}]}_1$ y el nivel $j$ usa
$\rho^{[r_j]}_1$; esta normalización se mantiene al retirar la profundidad.
Sobre la diagonal coherente, $z_{n,M}=u_M^{\odot n}$ satisface
\begin{equation}\label{eq:pdf2-ns-carleman}
 \dot z_{n,M}=L_{n,M}z_{n,M}
 +\mathcal N_{n,n+1;M}z_{n+1,M}
 +\mathcal F_{n,n-1;M}(f_M)z_{n-1,M}.
\end{equation}
Los lectores de grados uno y dos son
\[
 z_M(u)=\nu^{1/2}A_M^{1/2}\Lambda^su,
 \qquad
 w_M(u)=\nu^{-1/2}A_M^{-1/2}\Lambda^sN_M(u),
\]
y Young da
\begin{equation}\label{eq:pdf2-ns-young-port}
 |\mathscr B_{s,M}(u)|\le
 \beta\nu\mathscr D_{s,M}(u)+\frac{\|w_M(u)\|^2}{4\beta}.
\end{equation}
La estimación uniforme de las filas de Fourier controla este lector cuadrático en
el estado inicial. No controla todavía sus propagaciones futuras.

Sea
\[
 \mathcal N_{M,j}:=
 \frac{r-1}{r+1}H_{M,j}-D_{M,j}^{\rm mem}
 =2(r-1)M_{M,j}^-\ge0.
\]
En un horizonte $T$, el almacenamiento terminal es
\begin{equation}\label{eq:pdf2-ns-storage}
 \begin{aligned}
 G_{M,j,T}(t)
 &=\mathscr E_{s,M,j}(t)+2(r-1)(M^-_{M,j}(T)-M^-_{M,j}(t)),\\
 \dot G_{M,j,T}+\nu\mathscr D_{s,M,j}+|\mathscr B_{s,M,j}|
 &=\mathscr F_{s,M,j}.
 \end{aligned}
\end{equation}
Para tipar el terminal se conserva en el estado solenoidal completo la coordenada
de hoja
\[
 q^-_{M,j}:=\sqrt{2(r-1)\dot M^-_{M,j}}
 \in L^2(0,T;\mathbb R).
\]
Sea $\mathscr X^{\rm sol}_{M,j,T}$ el espacio de historia enriquecida que contiene,
como sumandos ortogonales, la traza $2^{-1/2}\Lambda^su_{M,j}$ y $q^-_{M,j}$.
La columna terminal es el operador de coordenadas
\begin{equation}\label{eq:pdf2-ns-terminal-column}
 \begin{aligned}
 \mathcal E_{M,j,T}^{\rm term}(t):\mathscr X^{\rm sol}_{M,j,T}
 &\longrightarrow L^2_\sigma(\mathbb T^3)
   \oplus L^2(t,T;\mathbb R),\\
 \mathcal E_{M,j,T}^{\rm term}(t)\Xi
 &:=2^{-1/2}\Lambda^su_{M,j}(t)
 \oplus\mathbf1_{[t,T]}q^-_{M,j}.
 \end{aligned}
\end{equation}
Por construcción de las dos coordenadas ortogonales,
\begin{equation}\label{eq:pdf2-ns-terminal-column-gram}
 \|\mathcal E_{M,j,T}^{\rm term}(t)\Xi\|^2
 =G_{M,j,T}(t).
\end{equation}
El sumando de memoria futura es una salida terminal tipada, no una componente de la
raíz inicial; su coste se controlará sólo después de la telescopía de la columna
completa. Para una observación $\tau$, las filas de pérdida viven en
$[0,\tau]$ y esta memoria terminal en $(\tau,T]$; los soportes temporales son
disjuntos. En la estimación integral se toma $\tau=T$, y entonces el sumando
futuro es cero.

\section{Construcción source--first con lift mixto estado--fuerza}

La entrada se fija antes de todo corte, profundidad e historia:
\begin{equation}\label{eq:pdf2-ns-data-space}
 \mathcal D_T:=H^s_\sigma(\mathbb T^3)
 \times L^2(0,T;H^{s-1}_\sigma(\mathbb T^3)),
 \qquad d=(u_0,f).
\end{equation}
Como los campos tienen media nula, $A=-\Delta$ es positivo e invertible en el
sector solenoidal. Se introducen el esfuerzo, el flujo viscoso y las dos ondas
\begin{equation}\label{eq:pdf2-ns-force-waves}
 \begin{aligned}
 a_T(f)&:=\frac{1}{2\sqrt\nu}A^{-1/2}\Lambda^sP_\sigma f,
 &z_M(u)&:=\sqrt\nu A_M^{1/2}\Lambda^su,\\
 b_M(u,f)&:=z_M(u)-P_Ma_T(f),\\
 q_M^{B,+}(u)&:=\bigl(\mathscr B_{s,M}(u)\bigr)_+^{1/2},
 &q_M^{B,-}(u)&:=\bigl(-\mathscr B_{s,M}(u)\bigr)_+^{1/2},\\
 q_M^B(u)&:=\bigl(q_M^{B,+}(u),q_M^{B,-}(u)\bigr).
 \end{aligned}
\end{equation}
Todos pertenecen a $L^2$ sobre el intervalo correspondiente y
$\|q_M^B(u)\|_{\mathbb C^2}^2=|\mathscr B_{s,M}(u)|$.
La potencia cruzada está tipada por
\begin{equation}\label{eq:pdf2-ns-cross-power}
 \mathscr F_{s,M}=2\operatorname{Re}\langle z_M,P_Ma_T(f)\rangle,
 \qquad \nu\mathscr D_{s,M}=\|z_M\|^2.
\end{equation}
Por tanto \eqref{eq:pdf2-ns-storage}, y no una sustitución de la potencia por la
norma de la fuerza, da la identidad exacta de scattering
\begin{equation}\label{eq:pdf2-ns-scattering-balance}
 \dot G_{M,j,T}+\|b_M\|^2+\|q^B_M\|^2=\|P_Ma_T(f)\|^2.
\end{equation}

\subsection{Firma cronológica, suavizante uniforme y creación mixta}

El dato de fuerza vive en $F:=H^{s-1}_\sigma(\mathbb T^3)$, mientras que la
torre PC--Fock usa como espacio de una partícula
$V:=H^s_\sigma(\mathbb T^3)$. La identificación entre ambos no se deja a la
inclusión Galerkin, cuya norma crecería con $M$. En el sector de media nula se
fija el isomorfismo de orden menos uno
\begin{equation}\label{eq:pdf2-ns-uniform-smoother}
 \begin{aligned}
 \mathscr S&:=A^{-1/2}:F\longrightarrow V,\\
 \mathscr S_M&:=A_M^{-1/2}P_M=P_M\mathscr S,\\
 \sup_M\bigl(\|\mathscr S_M\|_{F\to V}
 +\|\mathscr S_M^{-1}\|_{P_MV\to P_MF}\bigr)&<\infty.
 \end{aligned}
\end{equation}
La última cota es la equivalencia uniforme de las normas de Sobolev en el
subespacio sin modo cero. Así, $g:=\mathscr Sf\in L^2(0,T;V)$ y no se usa la
inclusión no uniforme $P_MH^{s-1}\hookrightarrow P_MH^s$.

Sea $\Delta_k(T)=\{0<t_1<\cdots<t_k<T\}$. El espacio de firma de la fuerza
suavizada es
\begin{equation}\label{eq:pdf2-ns-signature-space}
 \mathfrak S_T(V):=\bigoplus_{k\geq0}L^2(\Delta_k(T);V^{\widehat\otimes k}),
 \qquad \Delta_0(T)=\{*\}.
\end{equation}
Para $g=\mathscr Sf\in L^2(0,T;V)$ se define, sin usar la solución,
\begin{equation}\label{eq:pdf2-ns-force-signature}
 \operatorname{Sig}_T(g)_0=1,
 \qquad
 \operatorname{Sig}_T(g)_k(t_1,\ldots,t_k)
 =g(t_k)\widehat\otimes\cdots\widehat\otimes g(t_1).
\end{equation}
La simetría del integrando escalar proporciona
\begin{equation}\label{eq:pdf2-ns-signature-norm}
 \|\operatorname{Sig}_T(g)\|^2
 =\sum_{k\geq0}\frac{\|g\|_{L^2(0,T;V)}^{2k}}{k!}
 =e^{\|g\|_{L^2(0,T;V)}^2}
 \le e^{c_{\mathscr S}^2\|f\|_{L^2(0,T;F)}^2}.
\end{equation}

El término forzado de \eqref{eq:pdf2-ns-carleman} no es un puerto aditivo en
$f$ solamente. Para $0<r_{\rm F}<R_{\rm F}$ se tipa mediante la creación mixta
\begin{equation}\label{eq:pdf2-ns-mixed-creation}
 \begin{aligned}
 \widetilde{\mathcal F}^{R_{\rm F},r_{\rm F}}_M:
 P_MV\widehat\otimes\Gamma_s^{R_{\rm F}}(P_MV)
 &\longrightarrow\Gamma_s^{r_{\rm F}}(P_MV),\\
 \widetilde{\mathcal F}^{R_{\rm F},r_{\rm F}}_M(g\otimes x)
 &:=a^\dagger(g)x.
 \end{aligned}
\end{equation}
En el grado $n$ su acción es precisamente
$g\otimes z_{n-1}\mapsto n g\odot z_{n-1}$. Así se conserva el factor mixto
$\mathscr Sf\otimes x$ que no puede producir una transición lineal sobre la suma directa
``estado $\oplus$ fuerza''.
Con las normas ponderadas de \eqref{eq:pdf2-ns-pcf},
\begin{equation}\label{eq:pdf2-ns-creation-uniform}
 \|\widetilde{\mathcal F}^{R,r}_M(g\otimes x)\|_{\Gamma^r}
 \le c(R,r)\|g\|_V\|x\|_{\Gamma^R},
 \qquad
 c(R,r):=\sup_{n\ge1}\sqrt n\,r\left(\frac rR\right)^{n-1}<\infty.
\end{equation}
Para iterar se fija $0<\theta<1$ y $r_j=R_{\rm F}\theta^j$; cada paso usa
$(R,r)=(r_j,r_{j+1})$. Entonces
$\sup_{j,M}c(r_j,r_{j+1})\le
R_{\rm F}\sup_{n\ge1}\sqrt n\,\theta^n<\infty$.

La torre se escribe en la coordenada conjugada
$v_M=\mathscr S_Mu_M$. Sea
$\mathscr S_M^{\Gamma}:=\bigoplus_{n\ge0}
\mathscr S_M^{\widehat\otimes n}$. El generador homogéneo conjugado es
\begin{equation}\label{eq:pdf2-ns-conjugated-generator}
 \mathcal A^{0,\mathscr S}_{M}
 :=\mathscr S_M^{\Gamma}\mathcal A^0_M
   (\mathscr S_M^{\Gamma})^{-1},
 \qquad
 (\mathcal A^{0,\mathscr S}_{M})_{n,m}
 =\mathscr S_M^{\widehat\otimes n}
  (\mathcal A^0_M)_{n,m}
  (\mathscr S_M^{-1})^{\widehat\otimes m}.
\end{equation}
Esta conjugación es source--first y no altera la ODE: el lector de grado uno
aplica $\mathscr S_M^{-1}$ al final.

Sea $\mathcal T_M^{0,\mathscr S}(t,s)$ el propagador de la torre conjugada, incluido
en la realización solenoidal completa. Para $R>r>0$ se define primero sobre
\begin{equation}\label{eq:pdf2-ns-radial-transfer}
 \delta_{R\to r}(x_0,x_1,x_2,\ldots)
 :=\bigl(x_0,(r/R)x_1,(r/R)^2x_2,\ldots\bigr),
 \qquad
 \delta_{R\to r}\mathcal V_R^{\rm PCF}(u)=\mathcal V_r^{\rm PCF}(u).
\end{equation}
Este operador es contractivo y fija el vacío. Sobre
\begin{equation}\label{eq:pdf2-ns-dyson-domain}
 \Gamma_s^R(P_MV)_{\rm alg}\widehat\otimes
 \mathfrak S_T(P_MV)_{\rm fin}
\end{equation}
Sea $\mathfrak S_{T,{\rm fin}}^{\rm dec}$ el span de los núcleos elementales
\[
 g_k(t_1,\ldots,t_k)=
 g_k^{(k)}(t_k)\widehat\otimes\cdots\widehat\otimes g_k^{(1)}(t_1).
\]
Es denso en cada sumando $L^2(\Delta_k(T);V^{\widehat\otimes k})$. En la
fórmula siguiente los símbolos $g_k^{(i)}$ se leen en cada núcleo elemental y
la acción se extiende linealmente a su span; por tanto no se factoriza
arbitrariamente un tensor entrelazado. El lector cronológico toma valores en
$\Gamma_s^r(P_MV)$ y se define por
\begin{equation}\label{eq:pdf2-ns-dyson-reader}
 \begin{aligned}
 \mathscr D_{M,k}^{R\to r}(t)(x\otimes g_k)
 :={}&\delta_{R\to r}\int_{\Delta_k(t)}
 \mathcal T_M^{0,\mathscr S}(t,t_k)
 \widetilde{\mathcal F}_M(P_Mg_k^{(k)}(t_k)\otimes\cdot)\cdots\\
 &\hspace{18mm}\cdots\mathcal T_M^{0,\mathscr S}(t_2,t_1)
 \widetilde{\mathcal F}_M(P_Mg_k^{(1)}(t_1)\otimes\cdot)
 \mathcal T_M^{0,\mathscr S}(t_1,0)x\,d\mathbf t,\\
 \mathscr D_{M,T}^{R\to r}(t)&:=\sum_{k\geq0}
 \mathscr D_{M,k}^{R\to r}(t).
 \end{aligned}
\end{equation}
Al distribuir la pérdida total $R-r$ entre las $k$ creaciones y los propagadores,
la estimación analítica de escala de
\eqref{eq:pdf2-ns-creation-uniform} da, para constantes finitas dependientes del
corte pero no del núcleo elemental,
\begin{equation}\label{eq:pdf2-ns-dyson-homogeneous-bound}
 \|\mathscr D_{M,k}^{R\to r}(t)(x\otimes g_k)\|_{\Gamma^r}
 \le A_{M,R,r,T}\,
 \frac{B_{M,R,r,T}^{k}}{\sqrt{\Gamma(k+1)}}
 \|x\|_{\Gamma^R}\|g_k\|_{L^2(\Delta_k;V^{\widehat\otimes k})}.
\end{equation}
El factor $1/\Gamma(k+1)$ del volumen del simplex y Cauchy--Schwarz en la
suma de grados dan
\begin{equation}\label{eq:pdf2-ns-dyson-bound}
 \|\mathscr D_{M,T}^{R\to r}(t)(x\otimes\sigma)\|_{\Gamma^r}
 \le C_{M,R,r,T}\|x\|_{\Gamma^R}\|\sigma\|_{\mathfrak S_T},
 \qquad C_{M,R,r,T}<\infty.
\end{equation}
Por ello la serie converge en el codominio indicado y determina por densidad
un operador acotado para cada corte. No se usa aquí una cota uniforme en $M$:
esa uniformidad se obtendrá después de la identidad local de la coligación.
Al evaluarlo en
$x\otimes\operatorname{Sig}_T(\mathscr Sf)$ se obtiene la serie de Dyson del generador no
autónomo
\begin{equation}\label{eq:pdf2-ns-nonautonomous-generator}
 \mathcal A_M^{f,\mathscr S}(t)
 =\mathcal A_M^{0,\mathscr S}+a^\dagger(\mathscr S_MP_Mf(t)),
 \qquad
 \partial_tU_M^{f,\mathscr S}(t,s)
 =\mathcal A_M^{f,\mathscr S}(t)U_M^{f,\mathscr S}(t,s).
\end{equation}
La diferenciación de \eqref{eq:pdf2-ns-dyson-reader} en el core algebraico da
la conjugada de \eqref{eq:pdf2-ns-carleman}; la unicidad de la ODE finita da
\begin{equation}\label{eq:pdf2-ns-dyson-degree-one}
 \mathscr S_M^{-1}\rho^{[r]}_1
 \mathscr D_{M,T}^{R_{\rm F}\to r}(t)
 \bigl(\mathcal V_M^{\rm PCF}(\mathscr S_MP_Mu_0)
       \otimes\operatorname{Sig}_T(\mathscr S_MP_Mf)\bigr)
 =u_M(t).
\end{equation}
La igualdad incluye $t=0$: el factor $r^{-1}$ de
\eqref{eq:pdf2-ns-normalized-degree-one-reader} cancela exactamente la
coordenada radial $r\mathscr S_MP_Mu_0$ producida por
\eqref{eq:pdf2-ns-radial-transfer}. No se impone $R_{\rm F}=1$.
Esta construcción precede a la historia y no invoca
$\mathsf S_{M,T}d$ para definirla.

Sea
$\mathfrak s_{\rm sol}:H^s_\sigma\to\mathcal F_{\rm seed}^{\rm sol}$ la
restricción solenoidal del mapa de semilla: conserva hoja, orientación, residuo,
cociente y ruta basada y depende sólo de $u_0$. La raíz gobernante no se
redefine aquí: es exactamente la raíz $J_T^{\rm sol}$ de
\eqref{eq:pdf2-g-sol-raiz-dato}. Fijamos la identidad de nombres
\begin{equation}\label{eq:pdf2-ns-root-identification}
 J_T^{\rm mix}:=J_T^{\rm sol}:
 \mathcal D_T=\mathcal D_T^{\rm NS}\longrightarrow\mathcal H_T^{\rm sol}.
\end{equation}
Su publicación de coordenadas mixtas, y no una segunda raíz, es
\begin{equation}\label{eq:pdf2-ns-root-map}
 \begin{aligned}
 \operatorname{Coord}_T^{\rm mix}(u_0,f)&:=
 \bigl[\mathcal V^{\rm PCF}(u_0)
 \widehat\otimes\operatorname{Sig}_T(\mathscr Sf)\bigr]
 \oplus\mathfrak s_{\rm sol}(u_0),\\
 \operatorname{Coord}_T^{\rm mix}&:\mathcal D_T\longrightarrow
 \mathfrak H_T^{\rm coord},\\
 \mathfrak H_T^{\rm coord}&:=
 \bigl[\Gamma_s^{R_{\rm F}}(V)
 \widehat\otimes\mathfrak S_T(V)\bigr]
 \oplus\mathcal F_{\rm seed}^{\rm sol}.
 \end{aligned}
\end{equation}
La carta es una proyección de coordenadas del estado aceptado. En el subespacio
cíclico generado por la raíz se escribe
\begin{equation}\label{eq:pdf2-ns-mixed-coordinate-chart}
 \chi_T^{\rm mix}:\mathcal H_T^{\rm sol}\longrightarrow
 \mathfrak H_T^{\rm coord},
 \qquad
 \chi_T^{\rm mix}J_T^{\rm sol}(u_0,f)
 =\operatorname{Coord}_T^{\rm mix}(u_0,f).
\end{equation}
Su acción fuera del subespacio cíclico es la proyección ortogonal de
$\mathfrak G_{\rm sol}$ sobre la coordenada torre--firma--semilla. No se crea
un cociente nuevo ni se identifica la norma de esta carta con la norma total
de $\mathcal H_T^{\rm sol}$.
Con la normalización de la semilla,
\begin{equation}\label{eq:pdf2-ns-root-norm}
 \|\operatorname{Coord}_T^{\rm mix}(u_0,f)\|^2
 \leq e^{R_{\rm F}^2\|u_0\|_{H^s}^2
             +c_{\mathscr S}^2\|f\|_{L^2(0,T;H^{s-1})}^2}
 +c_{\rm seed}(1+\|u_0\|_{H^s}^2).
\end{equation}
Esta estimación prueba que la carta de coordenadas se forma sólo con el dato;
no transporta una norma nueva a $\mathcal H_T^{\rm sol}$ ni se usa para fabricar
la cota gobernante. La norma y la columna de $J_T^{\rm mix}=J_T^{\rm sol}$ son
las ya fijadas por $\mathfrak G_{\rm sol}$.

\subsection{Dato único de partida y núcleos densos}

El único dato admitido en este capítulo es la imagen completa
$\mathfrak G_{\rm sol}$ reproducida en
\eqref{eq:pdf2-g-sol-explicito}--\eqref{eq:pdf2-g-sol-cota-dato}. En particular,
se usan sin renombrarlos la raíz $J_T^{\rm sol}$, el codificador
$\Lambda_{M,T}$, las transiciones $\Gamma_{M,j}^{\rm term}$, la fila completa
$L_{M,j,T}^{\rm term}$, el terminal $\mathcal E_{M,j,T}^{\rm term}$, la columna
$\mathbf G_{M,j;J,T}^{\rm term}$ y la isometría $I_{M,J}$. No se introduce una
segunda realización de esos objetos.

Sea $\mathcal D_{T,{\rm alg}}\subset\mathcal D_T$ el subespacio de datos con
$u_0$ trigonométrico y $f$ suave, de soporte temporal compacto y con valores
trigonométricos. Es denso en $\mathcal D_T$. Los núcleos sobre los que se
escriben todas las acciones son
\begin{equation}\label{eq:pdf2-ns-cyclic-cores}
 \begin{aligned}
 \mathscr C_T^{\rm src}
 &:=\operatorname{span}\{J_T^{\rm sol}d:d\in\mathcal D_{T,{\rm alg}}\}
 \subset\mathcal H_T^{\rm sol},\\
 \mathscr C_{M,j,T}^{\rm sol}
 &:=\operatorname{span}\{x_{M,j,T}(d):d\in\mathcal D_{T,{\rm alg}}\},\\
 x_{M,0,T}(d)&:=\Lambda_{M,T}J_T^{\rm sol}d,\\
 x_{M,j+1,T}(d)&:=\Gamma_{M,j}^{\rm term}x_{M,j,T}(d).
 \end{aligned}
\end{equation}
Los espacios $\mathscr X_{M,j,T}^{\rm sol}$ son las completaciones de estos
núcleos en las normas ya contenidas en $\mathfrak G_{\rm sol}$. En particular,
ninguna norma se define restando la pérdida que después se quiere acotar.

La identificación exacta entre la notación de aplicación y la del dato es
\begin{equation}\label{eq:pdf2-ns-exact-datum-identification}
 \boxed{\begin{gathered}
 J_T^{\rm mix}=J_T^{\rm sol},\qquad
 \Gamma_{M,j,T}=\Gamma_{M,j}^{\rm term},\qquad
 L_{M,j,T}=L_{M,j,T}^{\rm term},\\
 \mathcal G_{M,J,T}^{\rm sol}=\mathbf G_{M,0;J,T}^{\rm term}
 \end{gathered}}
\end{equation}
Esta igualdad es literal, no una semejanza de nombres. La acción source--first
queda por tanto
\begin{equation}\label{eq:pdf2-ns-source-encoder-action}
 \boxed{\begin{gathered}
 d\xmapsto{J_T^{\rm sol}}J_T^{\rm sol}d
 \xmapsto{\Lambda_{M,T}}x_{M,0,T}(d)
 \xmapsto{\Gamma_{M,0:j,T}}x_{M,j,T}(d)
 \end{gathered}},
 \qquad
 \Gamma_{M,0:j,T}:=\Gamma_{M,j-1,T}\cdots\Gamma_{M,0,T}.
\end{equation}
Ninguna de estas tres flechas recibe $u_M$, una pérdida o un terminal como
entrada. La carta PC--Fock--firma de
\eqref{eq:pdf2-ns-root-map} publica las coordenadas de la primera flecha; la
igualdad de grado uno
\eqref{eq:pdf2-ns-dyson-degree-one} prueba, en el núcleo y luego por densidad,
\begin{equation}\label{eq:pdf2-ns-sol-transition-action}
 \pi_{M,j,T}^{(1)}x_{M,j,T}(u_0,f)=u_M(t_j),
 \qquad
 \pi_{M,j,T}^{(1)}:=
 \mathscr S_M^{-1}
 (\rho_1^{[r_j]}\widehat\otimes\varepsilon_{[t_j,T]})\pi_{M,j,T}^{\rm tf}.
\end{equation}
El factor $r_j^{-1}$ se aplica antes de $\mathscr S_M^{-1}$ y cancela el radio
de grado uno. Por ello el lector devuelve $u_M$, no $r_ju_M$ ni
$\mathscr S_Mu_M$.

Para $N\ge M$, la proyección de refinamiento actúa tensorialmente sobre la
carta de coordenadas:
\begin{equation}\label{eq:pdf2-ns-refinement-action}
 \begin{aligned}
 \mathfrak p_{M\leftarrow N}^{\Gamma}
 (v_1\widehat\otimes\cdots\widehat\otimes v_k)
 &:=(P_Mv_1)\widehat\otimes\cdots\widehat\otimes(P_Mv_k),\\
 \mathfrak p_{M\leftarrow N}^{\mathfrak S}
 (g_1\widehat\otimes\cdots\widehat\otimes g_k)
 &:=(P_Mg_1)\widehat\otimes\cdots\widehat\otimes(P_Mg_k).
 \end{aligned}
\end{equation}
Junto con la proyección de semilla de $\mathfrak G_{\rm sol}$, estas fórmulas
dan el mapa $\mathfrak p_{M\leftarrow N}^{\rm sol}$ y los cuadrados
\begin{equation}\label{eq:pdf2-ns-refinement-commutation}
 \begin{aligned}
 \mathfrak p_{M\leftarrow N}^{\rm sol}\Lambda_{N,T}J_T^{\rm sol}
 &=\Lambda_{M,T}J_T^{\rm sol},\\
 \mathfrak p_{M\leftarrow N}^{\rm sol}\Gamma_{N,j,T}
 -\Gamma_{M,j,T}\mathfrak p_{M\leftarrow N}^{\rm sol}
 &=\iota_{\rm carry}\mathcal C_{M\leftarrow N,j},\\
 \mathcal C_{M\leftarrow N,j}x_{N,j,T}(d)
 &=C_{M\leftarrow N}(d)|_{I_j}.
 \end{aligned}
\end{equation}
El miembro derecho es el operador de coordenada cuyo valor publicado es
exactamente el carry de
\eqref{eq:pdf2-ns-carry}; no se lo declara nulo para fabricar naturalidad.

\subsection{Las seis filas como componentes del único operador de pérdida}

Fijemos $0=t_0<t_1<\cdots<t_J=\tau\le T$, $I_j=[t_j,t_{j+1}]$ y el calendario
\begin{equation}\label{eq:pdf2-ns-cutoff-calendar}
 M_j:=3^jM,
 \qquad P_{M_j}\longrightarrow I
 \quad\hbox{fuertemente en cada }H^r_\sigma.
\end{equation}
Los seis codominios, con sus normas propias, son
\begin{equation}\label{eq:pdf2-ns-six-loss-spaces}
 \begin{aligned}
 \mathcal Y_{M,j}^{\rm diss}&=L^2(I_j;L^2_\sigma),&
 \mathcal Y_{M,j}^{\rm adv}&=L^2(I_j;\mathbb C^2),\\
 \mathcal Y_{M,j}^{\rm carry}&=L^2(I_j;H^{s-1}_\sigma),&
 \mathcal Y_{M,j}^{\rm micro}&=L^2(I_j;\operatorname{Ran}Q_{\rm micro}),\\
 \mathcal Y_{M,j}^{\rm shell}
 &=L^2(I_j;\operatorname{Ran}(P_{M_{j+1}}-P_{M_j})\cap H^{s-1}_\sigma),&
 \mathcal Y_{M,j}^{\rm mem}
 &=L^2(I_j;\operatorname{Ran}Q^{\rm mem}_{M,j+1}).
 \end{aligned}
\end{equation}
Se pone
$\mathcal Y_{M,j}^{\rm sol}:=\widehat\bigoplus_{\mathfrak a}
\mathcal Y_{M,j}^{\mathfrak a}$, donde
$\mathfrak a\in\{{\rm diss,adv,carry,micro,shell,mem}\}$, y se denota por
$P_{\mathfrak a}$ su proyección ortogonal. La fila aceptada de
$\mathfrak G_{\rm sol}$ se descompone, y no se duplica, mediante
\begin{equation}\label{eq:pdf2-ns-loss-parseval}
 L_{M,j,T}^{\rm term}:
 \mathscr X_{M,j,T}^{\rm sol}\longrightarrow\mathcal Y_{M,j}^{\rm sol},
 \qquad
 L_{M,j,T}^{\mathfrak a}:=P_{\mathfrak a}L_{M,j,T}^{\rm term}.
\end{equation}

Para escribir las acciones sin convertir funciones no lineales de $u$ en
operadores lineales ficticios, se usan las proyecciones de coordenadas del lift
mixto ya contenido en $\mathfrak G_{\rm sol}$:
\begin{equation}\label{eq:pdf2-ns-row-coordinate-readers}
 \begin{aligned}
 \pi_zx_{M,j,T}(d)&=\mathbf1_{I_j}\sqrt\nu A_M^{1/2}\Lambda^su_M,&
 \pi_ax_{M,j,T}(d)&=\mathbf1_{I_j}P_Ma_T(f),\\
 \pi_{B,\pm}x_{M,j,T}(d)&=\mathbf1_{I_j}q_M^{B,\pm}(u_M),&
 \pi_{\rm carry}x_{M,j,T}(d)&=\mathbf1_{I_j}C_{M_j\leftarrow M_{j+1}},\\
 \pi_{\rm micro}x_{M,j,T}(d)&=\mathbf1_{I_j}Q_{\rm micro}y^{\rm micro}_{M,j},&
 \pi_{\rm shell}x_{M,j,T}(d)&=\mathbf1_{I_j}
 (P_{M_{j+1}}-P_{M_j})N_{M_{j+1}}(u_{M_{j+1}}),\\
 \pi_{\rm mem}x_{M,j,T}(d)&=\mathbf1_{I_j}
 Q^{\rm mem}_{M,j+1}m^{\rm raw}_{M,j+1}.&&
 \end{aligned}
\end{equation}
Aquí $y^{\rm micro}_{M,j}$ y $m^{\rm raw}_{M,j+1}$ no son símbolos
externos: son respectivamente las proyecciones sobre las coordenadas
$Q_{\rm micro}$ y $\operatorname{Mem}_{M,j+1}$ del estado
\eqref{eq:pdf2-g-sol-explicito}. Las ocho aplicaciones $\pi$ son lineales en
el estado levantado, aunque sus publicaciones sobre la diagonal física sean
polinomiales o contengan una raíz positiva. Su dominio es el núcleo
\eqref{eq:pdf2-ns-cyclic-cores}; su extensión acotada es la componente
correspondiente de $L_{M,j,T}^{\rm term}$.

Con estas proyecciones, las seis acciones son materialmente
\begin{equation}\label{eq:pdf2-ns-six-loss-actions}
 \begin{aligned}
 L^{\rm diss}_{M,j,T}&=\pi_z-\pi_a,&
 L^{\rm adv}_{M,j,T}&=(\pi_{B,+},\pi_{B,-}),\\
 L^{\rm carry}_{M,j,T}&=\pi_{\rm carry},&
 L^{\rm micro}_{M,j,T}&=\pi_{\rm micro},\\
 L^{\rm shell}_{M,j,T}&=\pi_{\rm shell},&
 L^{\rm mem}_{M,j,T}&=\pi_{\rm mem}.
 \end{aligned}
\end{equation}
La acotación no se infiere de los nombres. Para
$x\in\mathscr C_{M,j,T}^{\rm sol}$, la columna aceptada da
\begin{equation}\label{eq:pdf2-ns-row-boundedness}
 \|L_{M,j,T}^{\mathfrak a}x\|^2
 \le\|L_{M,j,T}^{\rm term}x\|^2
 \le\|\mathbf G_{M,j;J,T}^{\rm term}x\|^2.
\end{equation}
Por densidad, cada fila tiene una única extensión acotada a
$\mathscr X_{M,j,T}^{\rm sol}$. Ya no se forma un cociente nuevo ni se define
la norma del dominio desde la fila.

La ausencia de sobreconteo es la identidad de Parseval del único operador
$L^{\rm term}$:
\begin{equation}\label{eq:pdf2-ns-six-loss-no-overcount}
 \boxed{\begin{gathered}
 \|L_{M,j,T}^{\rm term}x\|^2
 =\sum_{\mathfrak a}\|L_{M,j,T}^{\mathfrak a}x\|^2,
 \qquad x\in\mathscr X_{M,j,T}^{\rm sol}
 \end{gathered}}
\end{equation}
Esta ortogonalidad coincide con la geometría física: los intervalos $I_j$
son disjuntos; carry vive bajo $P_{M_j}$ y shell bajo
$P_{M_{j+1}}-P_{M_j}$; micro vive bajo $Q_{\rm micro}$; y
$Q^{\rm mem}_{M,j+1}$ elimina la memoria heredada y la copia del canal
advectivo. Además $q_M^{B,+}q_M^{B,-}=0$ puntualmente y
$\|(q_M^{B,+},q_M^{B,-})\|^2=|\mathscr B_{s,M}|$, de modo que las dos hojas
son un solo canal orientado.

\subsection{Terminal, columna y Stein derivados sin redefinir normas}

En el núcleo terminal, la acción del operador aceptado es
\begin{equation}\label{eq:pdf2-ns-terminal-action-on-core}
 \mathcal E_{M,J,T}^{\rm term}(\tau)x_{M,J,T}(d)
 =2^{-1/2}\Lambda^su_M(\tau)
 \oplus\mathbf1_{[\tau,T]}q^-_{M,J}.
\end{equation}
Su extensión acotada a $\mathscr X_{M,J,T}^{\rm sol}$ es parte de
$\mathfrak G_{\rm sol}$ y su Gram terminal es, sin complementos inventados,
\begin{equation}\label{eq:pdf2-ns-terminal-metric}
 U_{M,J,T}:=
 (\mathcal E_{M,J,T}^{\rm term})^*\mathcal E_{M,J,T}^{\rm term},
 \qquad
 \|\mathcal E_{M,J,T}^{\rm term}x_{M,J,T}(d)\|^2
 =G_{M,J,T}(\tau).
\end{equation}
El lector de la primera coordenada terminal está definido en todo su
codominio por
\begin{equation}\label{eq:pdf2-ns-terminal-reader}
 \mathcal R_{M,\tau}^{(1)}(h\oplus q):=
 \sqrt2\,\Lambda^{-s}h,
 \qquad
 \mathcal R_{M,\tau}^{(1)}
 \mathcal E_{M,J,T}^{\rm term}(\tau)x_{M,J,T}(d)=u_M(\tau).
\end{equation}
No se invierte $\Gamma$, no se elige representante de un cociente y no aparece
un operador terminal complementario sin acción.

La columna terminal se fija primero por
\begin{equation}\label{eq:pdf2-ns-terminal-tail-column}
 \mathbf G_{M,J;J,T}^{\rm term}:=\mathcal E_{M,J,T}^{\rm term}.
\end{equation}
Para $0\le j<J$, la columna de cola se escribe con todos sus renglones, en el
orden terminal y luego profundidad decreciente:
\begin{equation}\label{eq:pdf2-ns-complete-column}
 \mathbf G_{M,j;J,T}^{\rm term}:=
 \begin{bmatrix}
  \mathcal E_{M,J,T}^{\rm term}\Gamma_{M,j:J,T}\\
  L_{M,J-1,T}^{\rm term}\Gamma_{M,j:J-1,T}\\
  \vdots\\
  L_{M,j+1,T}^{\rm term}\Gamma_{M,j:j+1,T}\\
  L_{M,j,T}^{\rm term}
 \end{bmatrix},
 \qquad
 \Gamma_{M,j:k,T}:=\Gamma_{M,k-1,T}\cdots\Gamma_{M,j,T}.
\end{equation}
Para $j<J$ hay una igualdad de operadores entre columnas ya construidas,
\begin{equation}\label{eq:pdf2-ns-tail-column-recursion}
 \mathbf G_{M,j;J,T}^{\rm term}
 =\begin{bmatrix}
   \mathbf G_{M,j+1;J,T}^{\rm term}\Gamma_{M,j,T}\\
   L_{M,j,T}^{\rm term}
  \end{bmatrix}.
\end{equation}
Ahora, y sólo ahora, se toma el Gram
\begin{equation}\label{eq:pdf2-ns-tail-gram}
 U_{M,j,T}:=
 (\mathbf G_{M,j;J,T}^{\rm term})^*
 \mathbf G_{M,j;J,T}^{\rm term}.
\end{equation}
Al multiplicar la igualdad de bloques
\eqref{eq:pdf2-ns-tail-column-recursion} por su adjunta se obtiene Stein:
\begin{equation}\label{eq:pdf2-ns-local-stein}
 \boxed{\begin{gathered}
 U_{M,j,T}=\Gamma_{M,j,T}^*U_{M,j+1,T}\Gamma_{M,j,T}
 +(L_{M,j,T}^{\rm term})^*L_{M,j,T}^{\rm term}
 \end{gathered}}
\end{equation}
Por \eqref{eq:pdf2-ns-six-loss-no-overcount}, el último sumando es
$\sum_{\mathfrak a}(L_{M,j,T}^{\mathfrak a})^*
L_{M,j,T}^{\mathfrak a}$. La ecuación de Stein es por tanto consecuencia de
la columna aceptada; no define la norma del estado ni sirve para justificar
retroactivamente la existencia de sus filas.

La telescopía es asimismo una multiplicación de la columna completa:
\begin{equation}\label{eq:pdf2-ns-global-column-isometry}
 \|\mathbf G_{M,0;J,T}^{\rm term}x\|^2
 =\|\mathcal E_{M,J,T}^{\rm term}\Gamma_{M,0:J,T}x\|^2
 +\sum_{j<J}\|L_{M,j,T}^{\rm term}
                  \Gamma_{M,0:j,T}x\|^2.
\end{equation}
Todos los términos de la derecha existen antes de esta igualdad y llevan las
normas de sus codominios en \eqref{eq:pdf2-ns-six-loss-spaces}.

\subsection{Factorización fuente--historia y cota uniforme}

La factorización aceptada en
\eqref{eq:pdf2-g-sol-raiz-dato} se vuelve, tras la identificación literal
\eqref{eq:pdf2-ns-exact-datum-identification},
\begin{equation}\label{eq:pdf2-ns-source-colligation}
 \boxed{\begin{gathered}
 \mathcal G_{M,J,T}^{\rm sol}\Lambda_{M,T}J_T^{\rm mix}d
 =I_{M,J}J_T^{\rm sol}d,
 \qquad I_{M,J}^*I_{M,J}=I
 \end{gathered}}
\end{equation}
La historia se define desde el lado izquierdo,
\begin{equation}\label{eq:pdf2-ns-root-factorization}
 \mathsf H_{M,J,T}^{\rm sol}(d):=
 \mathcal G_{M,J,T}^{\rm sol}\Lambda_{M,T}J_T^{\rm mix}d,
\end{equation}
y no desde una solución elegida después. La lectura terminal
\eqref{eq:pdf2-ns-terminal-reader} y la acción source--first
\eqref{eq:pdf2-ns-source-encoder-action} dan el cuadrado conmutativo
\begin{equation}\label{eq:pdf2-ns-data-history-commutation}
 \boxed{\begin{gathered}
 \mathcal R_{M,\tau}^{(1)}P_{\rm term}
 \mathsf H_{M,J,T}^{\rm sol}(u_0,f)=u_M(\tau)
 \end{gathered}}
\end{equation}
donde $P_{\rm term}$ es la proyección sobre el primer renglón de
\eqref{eq:pdf2-ns-complete-column}. Las demás proyecciones son exactamente las
seis acciones \eqref{eq:pdf2-ns-six-loss-actions}. Esto identifica toda la
columna con la historia Galerkin sin reconstruir el estado desde una salida.

\begin{theorem}[Cota raíz desde la estructura de partida explícita]
\label{thm:pdf2-ns-root-uniform}
Para todo $T<\infty$ y todo $d=(u_0,f)\in\mathcal D_T$,
\begin{equation}\label{eq:pdf2-ns-root-bound}
 \sup_{M,J}\|\mathsf H_{M,J,T}^{\rm sol}(d)\|^2
 \le C_T^{\rm sol}\|d\|_{\mathcal D_T}^2,
\end{equation}
donde $C_T^{\rm sol}$ es independiente de $M,J$.
\end{theorem}

\begin{proof}
Por \eqref{eq:pdf2-ns-root-identification} y
\eqref{eq:pdf2-ns-exact-datum-identification}, el miembro izquierdo es
literalmente
$\|\mathbf G_{M,0;J,T}^{\rm term}\Lambda_{M,T}J_T^{\rm sol}d\|^2$.
La construcción propietaria precedente da primero la isometría de cascada
\eqref{eq:ns-import-cascade-isometry}, la telescopía de Stein
\eqref{eq:ns-import-stein-telescope}, el ledger conservado
\eqref{eq:ns-import-ledger} y la financiación exacta del acarreo
\eqref{eq:ns-import-carry-funded}. Su suma ortogonal produce la cota raíz
\eqref{eq:ns-import-root-bound} y el presupuesto uniforme
\eqref{eq:ns-import-uniform-budget}; la clausura independiente
\eqref{eq:nsclose-owner-stein-telescope}--
\eqref{eq:nsclose-proprietary-uniform-budget} transporta ese presupuesto a
todos los cortes \(M,J\). Al sustituir la identificación anterior se obtiene
\eqref{eq:pdf2-ns-root-bound}. La fórmula
\eqref{eq:pdf2-g-sol-cota-dato} es, por tanto, el registro posterior de la
cota ya derivada y no una premisa de esta demostración. Las acciones
\eqref{eq:pdf2-ns-six-loss-actions}, el terminal
\eqref{eq:pdf2-ns-terminal-action-on-core} y
\eqref{eq:pdf2-ns-global-column-isometry} identifican sus seis componentes
sin doble conteo.
\end{proof}

\begin{falsifier}[Controles de independencia y tipado]
Si $J_T^{\rm mix}$ no es literalmente $J_T^{\rm sol}$, falla
\eqref{eq:pdf2-ns-source-colligation}. Si una de las seis filas no es una
proyección del único $L^{\rm term}$, falla Parseval en
\eqref{eq:pdf2-ns-six-loss-no-overcount}. Si se intenta definir la norma de
$\mathscr X_{M,j+1,T}^{\rm sol}$ como la norma anterior menos una pérdida,
queda sin probar la positividad y la operación está prohibida. Finalmente, si
se reemplaza $b=z-a$ por $z$, el defecto de
\eqref{eq:pdf2-ns-scattering-balance} es
$2\operatorname{Re}\langle z,a\rangle-\|a\|^2$, que vale $\|a\|^2$ para
$z=a\ne0$. Estos cuatro controles fallan antes de alcanzar la cota uniforme.
\end{falsifier}

\section{Carta KYP finita: cálculo correcto y estatuto auxiliar}

Fijados $M,J,T$, sean $A,B,C,D$ los bloques de un paso afín y $U_+$ la
métrica terminal del nivel siguiente. Defínanse
\begin{equation}\label{eq:pdf2-ns-kyp-ql}
 Q:=I-B^*U_+B-D^*D,
 \qquad L:=B^*U_+A+D^*C.
\end{equation}
Si $Q\succ0$, la recurrencia
\begin{equation}\label{eq:pdf2-ns-kyp-u}
 U:=A^*U_+A+C^*C+L^*Q^{-1}L
\end{equation}
da la factorización exacta
\begin{equation}\label{eq:pdf2-ns-kyp-factor}
 \begin{pmatrix}
 U-A^*U_+A-C^*C&-A^*U_+B-C^*D\\
 -B^*U_+A-D^*C&I-B^*U_+B-D^*D
 \end{pmatrix}
 =
 \begin{pmatrix}Q^{-1/2}L&-Q^{1/2}\end{pmatrix}^{*}
 \begin{pmatrix}Q^{-1/2}L&-Q^{1/2}\end{pmatrix}\succeq0.
\end{equation}
Para $Q\succeq0$ se requiere además
$\operatorname{Ran}L\subseteq\operatorname{Ran}Q^{1/2}$ y se usa la
pseudoinversa. Esta carta es finita. No da una cota uniforme de $Q^{-1}$, $U$ o
la ganancia al retirar $M,J$.

La mutación que introduce un término $C_-^*C_-$ dentro de un bloque positivo y
después pretende recuperarlo con signo negativo es algebraicamente falsa: una
factorización $R^*R$ no cambia el signo de uno de sus cuadrados. Por ello no se usa
en la rama gobernante.

\section{Cierre uniforme, compacidad y solución global}

La carta KYP anterior permanece como \texttt{CONTROL\_NO\_GOBERNANTE}. El cierre
global procede de la raíz source--first, de la telescopía terminal y del ledger
completo. No se utiliza ninguna cota uniforme de $Q^{-1}$ ni se toma el límite dentro
de una factorización finita.

\subsection{Extracción material de la cota de Sobolev}

En \eqref{eq:pdf2-ns-complete-column}, el primer renglón es literalmente
$\mathcal E^{\rm term}\Gamma_{0:J}$; por
\eqref{eq:pdf2-ns-data-history-commutation} su coordenada publicada es
$u_M(\tau)$ y, por \eqref{eq:pdf2-ns-terminal-column-gram}, domina
$2^{-1}\|\Lambda^su_M(\tau)\|^2$. Las seis filas son las aplicaciones tipadas de
\eqref{eq:pdf2-ns-six-loss-actions}. En particular, las dos primeras dan
\[
 \sum_{j<J}\|L^{\rm diss}_{M,j,T}\Xi\|^2
 =\int_0^T\|b_M\|^2dt,
 \qquad
 \sum_{j<J}\|L^{\rm adv}_{M,j,T}\Xi\|^2
 =\int_0^T|\mathscr B_{s,M}(u_M)|dt.
\]
Como $z_M=b_M+P_Ma_T(f)$,
\begin{equation}\label{eq:pdf2-ns-wave-to-dissipation}
 \nu\int_0^T \mathscr D_{s,M}dt
 =\int_0^T\|z_M\|^2dt
 \leq2\int_0^T\|b_M\|^2dt+2\|a_T(f)\|_{L^2(0,T;L^2)}^2.
\end{equation}
Aplicando el teorema~\ref{thm:pdf2-ns-root-uniform} a la restricción del dato a
$[0,\tau]$ para cada $0\le\tau\le T$, la contractividad de la restricción
temporal de $\mathfrak G_{\rm sol}$ da
$C_\tau^{\rm sol}\le C_T^{\rm sol}$ y
$\|d|_{[0,\tau]}\|_{\mathcal D_\tau}\le\|d\|_{\mathcal D_T}$. Usando además
\eqref{eq:pdf2-ns-terminal-column-gram} y
\eqref{eq:pdf2-ns-wave-to-dissipation}, se obtiene
\begin{equation}\label{eq:pdf2-ns-cota-requerida}
 \begin{aligned}
 &\sup_M\sup_{0\le t\le T}\|u_M(t)\|_{H^s}^2
 +\nu\sup_M\int_0^T\|u_M(t)\|_{H^{s+1}}^2\,dt\\
 &\quad
 +\sup_M\int_0^T|\mathscr B_{s,M}(u_M(t))|\,dt
 +\sup_{M,J}\sum_{j<J}\|C_{M_j\leftarrow M_{j+1}}\|^2\\
 &\quad
 +\sup_{M,J}\sum_{j<J}
   \bigl(\|Q_{\rm micro}y^{\rm micro}_{M,j}\|^2
          +\|\mathcal S_{M,j}^{\rm sh}\|^2
          +\|Q^{\rm mem}_{M,j+1}m^{\rm raw}_{M,j+1}\|^2\bigr)
 \le C_T(u_0,f),
 \end{aligned}
\end{equation}
donde
\[
 C_T(u_0,f):=c_{s,\nu}\bigl(1+C_T^{\rm sol}\bigr)
 \bigl(\|u_0\|_{H^s}^2+
       \|f\|_{L^2(0,T;H^{s-1})}^2\bigr)
\]
es finito para cada $T<\infty$ e independiente de $M$ y $J$. El paso de una
profundidad finita a toda la suma de pérdidas se justifica por monotonía de sumas no
negativas; en las tres últimas líneas, cada norma es la norma $L^2(I_j)$ del
codominio correspondiente en \eqref{eq:pdf2-ns-six-loss-spaces}. No se elimina
el terminal. La identidad
$E_9J_9=I$ publica el componente macroscópico después de esta estimación, mientras
$Q_{\rm micro}$ y el shell permanecen cuantificados en la misma desigualdad.

\subsection{Control temporal y paso de Aubin--Lions}

Puesto que $s>5/2$, se tienen las aplicaciones continuas
\begin{equation}\label{eq:pdf2-ns-product-bound}
 H^s\cdot H^s\longrightarrow H^s,
 \qquad
 B:H^s_\sigma\times H^s_\sigma\longrightarrow H^{s-1}_\sigma,
 \qquad
 \|B(v,w)\|_{H^{s-1}}\le c_s\|v\|_{H^s}\|w\|_{H^s}.
\end{equation}
La ecuación \eqref{eq:pdf2-ns-galerkin}, la contractividad de $P_M$ y
\eqref{eq:pdf2-ns-cota-requerida} implican
\begin{equation}\label{eq:pdf2-ns-time-bound}
 \sup_M\|\partial_tu_M\|_{L^2(0,T;H^{s-1})}
 \le C_T'(u_0,f,\nu)<\infty.
\end{equation}
En efecto, $Au_M$ está acotado en $L^2H^{s-1}$ por el segundo término de
\eqref{eq:pdf2-ns-cota-requerida}; el término bilineal está acotado mediante
\eqref{eq:pdf2-ns-product-bound}; y $f_M$ está acotado por el dato fuente.

Las inclusiones
\[
 H^{s+1}(\mathbb T^3)\Subset H^s(\mathbb T^3)
 \hookrightarrow H^{s-1}(\mathbb T^3)
\]
y el lema de Aubin--Lions proporcionan una subsucesión, todavía denotada por $u_M$,
tal que
\begin{equation}\label{eq:pdf2-ns-compactness}
 \begin{aligned}
 u_M&\rightharpoonup^*u&&\text{en }L^\infty(0,T;H^s),\\
 u_M&\rightharpoonup u&&\text{en }L^2(0,T;H^{s+1}),\\
 u_M&\longrightarrow u&&\text{en }L^2(0,T;H^s).
 \end{aligned}
\end{equation}
La convergencia fuerte y \eqref{eq:pdf2-ns-product-bound} dan
\begin{equation}\label{eq:pdf2-ns-nonlinear-limit}
 B_M(u_M,u_M)\longrightarrow B(u,u)
 \quad\text{en }L^1(0,T;H^{s-1}).
\end{equation}
Pasando al límite en la formulación integrada de
\eqref{eq:pdf2-ns-galerkin}, se obtiene
\begin{equation}\label{eq:pdf2-ns-limit-equation}
 \partial_tu+\nu Au+B(u,u)=f,
 \qquad \nabla\!\cdot u=0,
 \qquad u(0)=u_0.
\end{equation}
El dato inicial se conserva porque
$u_M$ está acotado en $H^1(0,T;H^{s-1})$ y converge en
$C([0,T];H^{s-1})$ después de extracción; como $P_Mu_0\to u_0$, la traza es la
indicada.

\subsection{Presión, unicidad y pegado de horizontes}

La presión de media nula no se introduce como dato. Se reconstruye del límite por
\begin{equation}\label{eq:pdf2-ns-pressure}
 -\Delta p=\partial_i\partial_j(u_i u_j),
 \qquad \int_{\mathbb T^3}p=0.
\end{equation}
Como $H^s$ es un álgebra, $p\in L^\infty(0,T;H^s)$ y
$\nabla p\in L^\infty(0,T;H^{s-1})$. Así,
\eqref{eq:pdf2-ns-limit-equation} equivale a la forma clásica
\[
 \partial_tu+(u\cdot\nabla)u-\nu\Delta u+\nabla p=f,
 \qquad\nabla\!\cdot u=0.
\]

Sean $u$ y $v$ dos soluciones con los mismos datos y póngase $w=u-v$. El producto
en $H^{s-1}$ y Young dan
\begin{equation}\label{eq:pdf2-ns-uniqueness}
 \frac12\frac{d}{dt}\|w\|_{H^{s-1}}^2
 +\frac\nu2\|w\|_{H^s}^2
 \le \frac{c_s}{\nu}
 \bigl(\|u\|_{H^s}^2+\|v\|_{H^s}^2\bigr)
 \|w\|_{H^{s-1}}^2.
\end{equation}
Grönwall y $w(0)=0$ implican $w=0$. La unicidad elimina la dependencia de la
subsucesión en \eqref{eq:pdf2-ns-compactness}.

La construcción vale para todo $T<\infty$. Aplicándola en $T=n$, $n\in\mathbb N$,
las soluciones coinciden en los solapamientos por
\eqref{eq:pdf2-ns-uniqueness}; por tanto se pegan en una única solución sobre
$[0,\infty)$. Para justificar ``suave'' sin extrapolar una sola cota, se repite
la construcción para la sucesión $s_n=s+n$, $n\ge0$. Los mismos lectores
estructurales actúan en cada escala, las proyecciones $P_M$ las entrelazan y
\eqref{eq:pdf2-g-sol-cota-dato} da en cada $s_n$ una constante
$C_{T,s_n}^{\rm sol}<\infty$, porque los datos son suaves. La unicidad en
$H^{s-1}$ identifica todos los límites obtenidos. Por
tanto $u\in C^\infty((0,\infty)\times\mathbb T^3)$ y la compatibilidad de las
trazas con $u_0\in C^\infty$ da suavidad hasta $t=0$. No aparece un tiempo de
ruptura porque la cota correspondiente está disponible en cada horizonte
finito y en cada $s_n$.

La condición de media cero no elimina datos periódicos. Para datos generales
$u_0^{\rm g}$ y $f^{\rm g}$, defínanse
\[
 m'(t)=\int_{\mathbb T^3}f^{\rm g}(t,x)\,dx,\qquad
 m(0)=\int_{\mathbb T^3}u_0^{\rm g}(x)\,dx,\qquad
 X(t)=\int_0^tm(r)\,dr.
\]
El cambio de Galileo dependiente del tiempo
\begin{equation}\label{eq:pdf2-ns-mean-reduction}
 u(t,x)=m(t)+v(t,x-X(t)),\qquad
 g(t,y)=f^{\rm g}(t,y+X(t))-m'(t)
\end{equation}
transforma biyectivamente el problema general en el problema de media cero
para $(v_0,g)$. La traslación preserva todas las normas de Sobolev, la
divergencia y la suavidad. Por tanto la construcción cubre todo dato periódico
suave; en particular cubre el caso no forzado del enunciado clásico.

\begin{theorem}[Cierre solenoidal global desde la estructura de partida explícita]
\label{thm:pdf2-ns-global}
Sea $s>5/2$, $\nu>0$,
$u_0\in C^\infty(\mathbb T^3)\cap H^s_\sigma$ y
$f\in C^\infty([0,\infty);H^\infty_\sigma)
\cap L^2_{\rm loc}([0,\infty);H^{s-1}_\sigma)$. Con el dato solenoidal completo
tipado por \eqref{eq:pdf2-ns-source-encoder-action},
\eqref{eq:pdf2-ns-sol-transition-action},
\eqref{eq:pdf2-ns-six-loss-no-overcount},
\eqref{eq:pdf2-ns-local-stein},
\eqref{eq:pdf2-ns-terminal-metric} y
\eqref{eq:pdf2-ns-data-history-commutation}, la imagen
$\mathfrak G_{\rm sol}$, con su terminal,
carry, memoria, shell, puerto ortogonal y lector raíz source--first, produce una única
solución global suave $(u,p)$ de Navier--Stokes periódico tridimensional. Para cada
$T<\infty$ satisface \eqref{eq:pdf2-ns-cota-requerida}; la construcción es natural en
$M$, estable al retirar $J$ y compatible con la publicación Galerkin. Mediante
\eqref{eq:pdf2-ns-mean-reduction}, el resultado se extiende a todo dato
periódico suave, sin imponer media cero.
\end{theorem}

\paragraph{Procedencia y fuerza probatoria.}
La imagen solenoidal, el terminal, el carry, la torre PC--Fock y la identidad local
\eqref{eq:pdf2-ns-local-stein} son la
\texttt{ARQUITECTURA\_AUTORAL\_PREEXISTENTE} tomada aquí como estructura de partida
explícita. El lift mixto \eqref{eq:pdf2-ns-mixed-creation}, la firma cronológica,
las ondas $a/b$, los seis dominios de salida, la unidad homogénea y la columna
terminal son \texttt{FORMALIZACION\_NUEVA}. Las ecuaciones
\eqref{eq:pdf2-ns-source-encoder-action},
\eqref{eq:pdf2-ns-source-colligation},
\eqref{eq:pdf2-ns-sol-transition-action},
\eqref{eq:pdf2-ns-six-loss-no-overcount},
\eqref{eq:pdf2-ns-local-stein},
\eqref{eq:pdf2-ns-terminal-metric} y
\eqref{eq:pdf2-ns-data-history-commutation} enumeran exactamente las premisas
empleadas; el
teorema~\ref{thm:pdf2-ns-root-uniform} es, en lenguaje humano, \emph{Demostrado con
estructura de partida explícita}. Aubin--Lions, la
reconstrucción de presión, Grönwall y el pegado son el
\texttt{RECONOCIMIENTO\_CLASICO} posterior. La carta KYP conserva sólo el estatuto
\texttt{CONTROL\_NO\_GOBERNANTE}; ni su comparador ni una hipótesis de signo forman
parte de la cadena causal del cierre.

\paragraph{Conclusión afirmativa.}
La cadena gobernante queda cerrada en el orden
\[
 \begin{gathered}
 \mathfrak G_{\rm sol}
 \longrightarrow J_T^{\rm mix}
 \longrightarrow\Lambda_{M,T}
 \longrightarrow\mathcal G_{M,J,T}^{\rm sol}
 \longrightarrow\mathsf H_{M,J,T}^{\rm sol}\\
 \longrightarrow\mathcal R_{M,\tau}^{(1)}P_{\rm term}=u_M
 \longrightarrow\text{cota uniforme}
 \longrightarrow\text{Aubin--Lions}
 \longrightarrow(u,p)_{[0,\infty)}.
 \end{gathered}
\]
Dentro de este perfil de partida explícito no queda una obligación KYP adicional:
el terminal y todas las pérdidas son imágenes ortogonales de una única raíz
anterior a la órbita futura.

\endgroup

\HMTFlushCurrentChapter
\chapter{Existencia cuántica y brecha de masa de Yang--Mills mediante coercividad de calibre y positividad espectral}
\begingroup
\renewcommand{\chapter}[2][]{}%
\chapter{Imagen gauge y composición hacia Yang--Mills}

\section{Objeto fuente y separación obligatoria de generadores}

El objeto ya generado que recibe este capítulo es la estructura gauge completa
$\mathfrak G_{\rm gau}$; no constituye un dato externo ni una hipótesis terminal.
En cada nivel $m$ conserva la red
$\Lambda_m=a_m\mathbb Z^4$, $a_{m+1}=a_m/9$, la medida $\mu_m$, los
entrelazadores $J_m$, cuatro reflexiones $\Theta_{\nu,m}$, holonomía, color,
representación, acarreo de fusión, ruta, memoria, terminal y lector de curvatura.
Sus dos dinámicas se tipan con símbolos distintos antes de cualquier límite.

El semigrupo de una sola fibra de memoria es
\begin{equation}\label{eq:pdf2-ym-fibra-semigrupo}
 K^{\rm fib}_{m,t}
 =P_{\Omega_m}+e^{-3\kappa_{\rm av}t}(I-P_{\Omega_m})
 =e^{-tD_m^{\rm fib}},
 \qquad
 D_m^{\rm fib}:=3\kappa_{\rm av}(I-P_{\Omega_m}).
\end{equation}
Su espectro está contenido en $\{0,3\kappa_{\rm av}\}$. El generador de la
carta producto completa es, en cambio,
\begin{equation}\label{eq:pdf2-ym-generador-producto}
 D_m^{\rm prod}:=3\kappa_{\rm av}
 \sum_{c\in\mathcal I_m}(I-E_{m,c}),
 \qquad
 K^{\rm prod}_{m,t}:=e^{-tD_m^{\rm prod}},
\end{equation}
donde las expectativas ortogonales $E_{m,c}$ conmutan. Si hay dos o más
coordenadas no constantes, $D_m^{\rm prod}$ posee niveles
$6\kappa_{\rm av},9\kappa_{\rm av},\ldots$; por ello
\eqref{eq:pdf2-ym-fibra-semigrupo} no puede declararse igual a
\eqref{eq:pdf2-ym-generador-producto}. El primero es la fibra simple; el segundo,
el producto de solapamiento.

\section{Teorema finito de Hoeffding, vacío y brecha alcanzada}

En la carta producto construida en
\eqref{eq:pdf2-ym-owner-product-factorization}, los $E_{m,c}$ son expectativas
condicionales ortogonales y conmutantes, y la prueba de
\eqref{eq:pdf2-ym-owner-vacuum} establece
\begin{equation}\label{eq:pdf2-ym-interseccion-vacio}
 \bigcap_{c\in\mathcal I_m}\operatorname{Ran}E_{m,c}
 =\mathbb C\Omega_m.
\end{equation}
La descomposición de Hoeffding es la suma ortogonal
\[
 \mathcal H_m=\bigoplus_{S\subseteq\mathcal I_m}\mathcal H_{m,S},
 \qquad u=\sum_Su_S,
\]
donde $E_{m,c}u_S=0$ si $c\in S$ y $E_{m,c}u_S=u_S$ si $c\notin S$.
Entonces
\begin{equation}\label{eq:pdf2-ym-hoeffding}
 D_m^{\rm prod}u_S=3\kappa_{\rm av}|S|u_S,
 \qquad
 \langle u,D_m^{\rm prod}u\rangle
 =3\kappa_{\rm av}\sum_S|S|\|u_S\|^2.
\end{equation}
La componente $S=\varnothing$ es exactamente el vacío de
\eqref{eq:pdf2-ym-interseccion-vacio}. Por tanto
\begin{equation}\label{eq:pdf2-ym-cota-finita}
 D_m^{\rm prod}\succeq
 3\kappa_{\rm av}(I-P_{\Omega_m}),
 \qquad
 \|K^{\rm prod}_{m,t}(I-P_{\Omega_m})\|
 \le e^{-3\kappa_{\rm av}t}.
\end{equation}

En una coordenada de incidencia
$q=(q_1,\ldots,q_6)\in\mathbb F_3^6$, el testigo
\begin{equation}\label{eq:pdf2-ym-wc}
 W_c(q):=\mathbf1_{\{q_1+\cdots+q_6=0\}}-\frac13
\end{equation}
satisface
\begin{equation}\label{eq:pdf2-ym-wc-momentos}
 \mathbb EW_c=0,\qquad \|W_c\|^2=\frac29,\qquad
 D_m^{\rm prod}W_c=3\kappa_{\rm av}W_c.
\end{equation}
Así, en cada nivel finito, el umbral de
\eqref{eq:pdf2-ym-cota-finita} se alcanza. Este cálculo fija el espectro del
generador producto finito; la medida cuatridimensional continua de Yang--Mills se
construye materialmente en
\eqref{eq:pdf2-ym-gibbs-measure}--\eqref{eq:pdf2-ym-g-dynamics}.

\section{Compresión gauge física}

Sea $G$ compacto, conectado y simple. El grupo que actúa en el nivel $m$ es
\[
 \mathcal G_m^{\rm full}
 =G^{V(\Lambda_m)}\times
  \prod_{c\in\mathcal C_m}\mathbb F_3^4,
\]
donde el segundo factor actúa sobre las seis incidencias mediante una matriz
que debe quedar visible. Si las filas se indexan por
$(01,02,03,12,13,23)$ y las columnas por $0,1,2,3$, se fija
\begin{equation}\label{eq:pdf2-ym-incidence-matrix}
 B^+_{(ij),k}:=\delta_{ik}+\delta_{jk}\in\mathbb F_3,
 \qquad
 (B^+x)_{ij}=x_i+x_j.
\end{equation}
Esta es la incidencia no orientada del collar; no se confunde con el
coborde firmado de una arista. La acción finita es
\begin{equation}\label{eq:pdf2-ym-incidence-action}
 q_c\longmapsto q_c+B^+x_c,
 \qquad x_c\in\mathbb F_3^4.
\end{equation}
La proyección de Haar
\begin{equation}\label{eq:pdf2-ym-proyeccion-fisica}
 \mathbb P_m^{\rm phys}F
 :=\int_{\mathcal G_m^{\rm full}}F(g^{-1}\!\cdot\,)\,dg,
 \qquad
 \mathcal H_m^{\rm phys}:=\operatorname{Ran}\mathbb P_m^{\rm phys}
\end{equation}
es distinta de la marginal que olvida incidencia. La invariancia de las
expectativas $E_{m,c}$ da
\begin{equation}\label{eq:pdf2-ym-reduccion}
 [\mathbb P_m^{\rm phys},D_m^{\rm prod}]=0,
 \qquad
 [\mathbb P_m^{\rm phys},K_{m,t}^{\rm prod}]=0.
\end{equation}
Además, cada variable de $x_c\in\mathbb F_3^4$ aparece tres veces en la suma de
las seis incidencias, de modo que
\begin{equation}\label{eq:pdf2-ym-incidence-conservation}
 \sum_{i<j}(B^+x_c)_{ij}
 =3\sum_{i=0}^3x_{c,i}=0
 \quad\text{en }\mathbb F_3.
\end{equation}
Por ello $W_c$ es gauge
invariante y pertenece a $\mathcal H_m^{\rm phys}$. La brecha finita no es un
artefacto de una dilatación no física.

\section{Positividad por reflexión en cada nivel}

Para un kernel reversible común, defínase la medida de ocho caras
\begin{equation}\label{eq:pdf2-ym-ocho-caras}
 \Omega_m^{(8)}(dy_1\cdots dy_8)
 :=\int\mu_m(dz)\prod_{r=1}^{8}K^{\rm prod}_{m,a_m/2}(z,dy_r).
\end{equation}
Si $\Theta_{\nu,m}$ intercambia las dos mitades respecto de la dirección
$\nu$, la propiedad de Markov y la reversibilidad factorizan, para todo observable
cilíndrico $F$ soportado en el semiespacio positivo,
\begin{equation}\label{eq:pdf2-ym-os-finita}
 \int\overline{F(\Theta_{\nu,m}y)}F(y)\,d\Omega_m^{(8)}
 =\|\mathcal B_{m,\nu}F\|_2^2\ge0,
 \qquad \nu=1,\ldots,4.
\end{equation}
La marginal de caras opuestas recupera el mismo
$K^{\rm prod}_{m,a_m}$. De este modo las cuatro formas OS finitas son
consecuencia de una sola medida y un solo semigrupo, no de cuatro procesos
independientes.

\section{Propietario final: el estado gauge completo, no su marginal ablacionada}

El propietario vigente no toma como espacio físico la marginal que conserva sólo
los enlaces. Su dominio es $\mathfrak G_{\rm gau}$, con enlace, marco,
representación, incidencia, color, ruta, orientación,
acarreo, memoria, supervivencia y terminal. La proyección a enlaces es un mapa de
olvido posterior. Juzgar el testigo físico después de aplicar ese olvido sustituye el
objeto y no transfiere la conclusión.

\subsection{Medida proyectiva y publicación de enlaces}

Para una pantalla finita $\Gamma_m=(V_m,E_m)$ se parte del espacio levantado
\begin{equation}\label{eq:pdf2-ym-owner-lifted-space}
 \widetilde X_m:=G^{V_m}\times G^{E_m},\qquad
 d\widetilde\mu_m(h,z)
 =\prod_{v\in V_m}dh_v\prod_{e\in E_m}k_{\tau_{m,e}}(z_e)\,dz_e,
\end{equation}
y se publica cada enlace mediante
\begin{equation}\label{eq:pdf2-ym-owner-link-publication}
 U_e=h_{s(e)}z_eh_{t(e)}^{-1}.
\end{equation}
Una arista gruesa $e$ se refina en la cadena ordenada
$e_1\cdots e_9$. Póngase
\(\boldsymbol\tau=(\tau_1,\ldots,\tau_9)\), con
\(\tau_j=\tau_{m+1,e_j}>0\) y
\(\tau=\tau_{m,e}=\sum_{j=1}^9\tau_j\). Sea
\(d\lambda_z^{(9)}\) la desintegración de Haar sobre la fibra del producto,
caracterizada por
\[
 \int_{G^9}F(z_1,\ldots,z_9)\,dz_1\cdots dz_9
 =\int_G\int_{z_1\cdots z_9=z}
   F\,d\lambda_z^{(9)}\,dz.
\]
El puente correcto para tiempos no necesariamente iguales es
\begin{equation}\label{eq:pdf2-ym-owner-heat-bridge}
 dB_{z,\boldsymbol\tau}^{(9)}(z_1,\ldots,z_9)
 =\frac{\prod_{j=1}^9k_{\tau_j}(z_j)}{k_\tau(z)}
   \,d\lambda_z^{(9)},
 \qquad z_1\cdots z_9=z,
\end{equation}
que es una probabilidad porque
\(k_{\tau_1}*\cdots*k_{\tau_9}=k_\tau\). En el refinamiento nonádico uniforme se
recupera \(\tau_j=\tau/9\), pero la fórmula no presupone esa igualdad.

Para hacer independiente la innovación sin borrar el producto grueso, se fija, para
cada $(e,z)$, un transporte boreliano
\begin{equation}\label{eq:pdf2-ym-owner-bridge-triangularization}
 \chi_{m,e,z}:\bigl((0,1)^{\mathbb N},du^{\otimes\mathbb N}\bigr)
 \longrightarrow
 \bigl(\{(z_j):z_1\cdots z_9=z\},B_{z,\boldsymbol\tau}^{(9)}\bigr)
\end{equation}
que es isomorfismo módulo conjuntos nulos. Existe porque ambas son realizaciones
estándar no atómicas; una elección queda fijada una vez y se transporta bajo
relabelados de aristas. La variable gruesa $z$ y la innovación $u^{\rm det}_{m+1,e}$
son entonces factores independientes en la carta triangular. Los marcos nuevos se
integran con Haar. Ruta, acarreo, orientación y memoria se transportan por los
morfismos $\operatorname{Ext}_{\rm gau}$ del propio dato gauge.

La proyección gruesa multiplica los nueve incrementos y olvida exclusivamente las
innovaciones nuevas. Con esta definición se obtiene
\begin{equation}\label{eq:pdf2-ym-owner-projectivity}
 (\widehat\pi_{m+1,m})_*\widehat\mu_{m+1}^{\rm ov}
 =\widehat\mu_m^{\rm ov},\qquad
 \widehat J_mF=F\circ\widehat\pi_{m+1,m}.
\end{equation}
La misma arista aparece una sola vez en la medida; las seis caras incidentes son seis
lectores orientados de esa arista. Las rutas solapadas conservan el incremento
compartido y no se reemplazan por copias independientes.

\subsection{La fibra de incidencia es una coordenada física del estado completo}

La cola de incidencia incluida en $\mathfrak G_{\rm gau}$ se descompone sin perder
información:
\begin{equation}\label{eq:pdf2-ym-owner-seed-split}
 \Sigma:(\mathbb F_3^2)^{\mathbb N}\longrightarrow
 (\mathbb F_3^2)^{\mathbb N}\times\mathbb F_3^6,
 \qquad \Sigma_*\sigma=\sigma\otimes u_6.
\end{equation}
El primer factor continúa alimentando las cartas de innovación y el segundo conserva
las seis incidencias $q_c=(q_{01},q_{02},q_{03},q_{12},q_{13},q_{23})$ del collar
$c$. En vez de dejar implícito qué coordenadas gobierna el generador, se fija una
carta triangular recursiva
\begin{equation}\label{eq:pdf2-ym-owner-product-factorization}
 \widehat X_m
 =\widehat X_{m_0}^{\rm seed}
  \times\prod_{\ell=m_0+1}^{m}
       \mathcal U_{\ell/\ell-1}^{\rm det},
 \qquad
 \widehat\mu_m^{\rm ov}
 =\widehat\mu_{m_0}^{\rm seed}
  \otimes\bigotimes_{\ell=m_0+1}^{m}
       \upsilon_{\ell/\ell-1}^{\rm det}.
\end{equation}
El factor semilla enumera, una vez, los marcos y enlaces gruesos, las incidencias
$q_c$ y sus cartas de marcación. Cada
$\mathcal U_{\ell/\ell-1}^{\rm det}$ enumera, una vez, los marcos nuevos, los puentes
\eqref{eq:pdf2-ym-owner-bridge-triangularization}, las incidencias nuevas y las
subcolas de marcación de ese refinamiento. La factorización es interna a
$\mathfrak G_{\rm gau}$ y no añade otro dato de partida.

Se registra el índice exhaustivo, disjunto y numerable
\begin{equation}\label{eq:pdf2-ym-owner-exhaustive-index}
 \mathcal I_m
 :=\mathcal I_{m_0}^{\rm frame}
 \sqcup\mathcal I_{m_0}^{\rm link}
 \sqcup\mathcal I_{m_0}^{q}
 \sqcup\mathcal I_{m_0}^{\rm mark}
 \sqcup\bigsqcup_{\ell=m_0+1}^{m}
 \bigl(
   \mathcal I_{\ell/\ell-1}^{\rm frame}
   \sqcup\mathcal I_{\ell/\ell-1}^{\rm bridge}
   \sqcup\mathcal I_{\ell/\ell-1}^{q}
   \sqcup\mathcal I_{\ell/\ell-1}^{\rm mark}
 \bigr).
\end{equation}
Cada factor independiente de \eqref{eq:pdf2-ym-owner-product-factorization} aparece
exactamente una vez; una subcola infinita se enumera por sus cilindros escalares. No
queda un factor genealógico fuera de $\mathcal I_m$.
En particular, las abreviaturas de la recurrencia significan literalmente
\begin{equation}\label{eq:pdf2-ym-owner-index-abbreviations}
 \begin{aligned}
 \mathcal I_{m_0}^{\rm seed}
 &:={\mathcal I}_{m_0}^{\rm frame}\sqcup
    {\mathcal I}_{m_0}^{\rm link}\sqcup
    {\mathcal I}_{m_0}^{q}\sqcup
    {\mathcal I}_{m_0}^{\rm mark},\\
 \mathcal I_{\ell/\ell-1}^{\rm det}
 &:={\mathcal I}_{\ell/\ell-1}^{\rm frame}\sqcup
    {\mathcal I}_{\ell/\ell-1}^{\rm bridge}\sqcup
    {\mathcal I}_{\ell/\ell-1}^{q}\sqcup
    {\mathcal I}_{\ell/\ell-1}^{\rm mark}.
 \end{aligned}
\end{equation}

La marginal de enlaces olvida $q_c$; la subálgebra física, en cambio, se obtiene
promediando la acción gauge completa
\begin{equation}\label{eq:pdf2-ym-full-gauge-group}
 \mathcal G_m^{\rm full}
 :=G^{V(\Lambda_m)}\times
 \prod_{c\in\mathcal C_m}\mathbb F_3^4,
 \qquad q_c\longmapsto q_c+B^+x_c.
\end{equation}
Por tanto la proyección física es
\begin{equation}\label{eq:pdf2-ym-full-physical-projection}
 \mathbb P_m^{\rm phys}F
 :=\int_{\mathcal G_m^{\rm full}}F(g^{-1}\!\cdot x)\,dg,
 \qquad
 \widehat{\mathscr H}_m^{\rm phys}
 :=\operatorname{Ran}\mathbb P_m^{\rm phys}.
\end{equation}
No coincide con el mapa de olvido
$\widehat X_m\to X_m^{\rm link}$. Esta distinción es el cortafuegos de tipo
que impide expulsar $W_c$ del sector físico.

\section{Generador propietario y circuito local}

La carta triangular completa proporciona una unitaria
\begin{equation}\label{eq:pdf2-ym-owner-triangular-unitary}
 \widehat\Gamma_m:
 L^2(\widehat X_{m+1},\widehat\mu_{m+1}^{\rm ov})
 \longrightarrow
 L^2(\widehat X_m,\widehat\mu_m^{\rm ov})
 \widehat\otimes\mathcal K_{m+1}^{\rm det},
 \qquad
 \widehat\Gamma_m\widehat J_mF=F\otimes\mathbf1_{\rm det}.
\end{equation}
Para $\iota\in\mathcal I_m$, sea $E_{m,\iota}$ la esperanza que integra
exactamente el factor $\iota$ de
\eqref{eq:pdf2-ym-owner-product-factorization} y deja fijos todos los demás.
Son proyecciones Markov ortogonales, conservan la unidad y conmutan. Sobre el
núcleo denso $\mathscr C_m$ de cilindros que dependen de un número finito de
coordenadas se define la forma
\begin{equation}\label{eq:pdf2-ym-owner-full-form}
 \widehat{\mathcal E}_m[F]
 :=3\kappa_{\rm av}\sum_{\iota\in\mathcal I_m}
   \|(I-E_{m,\iota})F\|_2^2.
\end{equation}
La suma es finita para cada cilindro. Su clausura es una forma de Dirichlet y
determina el generador autoadjunto Markov
$\widehat H_m^{\rm ov}$. En un nivel semilla esto se escribe, en sentido de
formas,
\begin{equation}\label{eq:pdf2-ym-owner-base-generator}
 \widehat H_{m_0}^{\rm ov}
 =3\kappa_{\rm av}
  \sum_{\iota\in\mathcal I_{m_0}^{\rm seed}}(I-E_{m_0,\iota}),
\end{equation}
y cada refinamiento añade exactamente las expectativas de sus nuevas coordenadas:
\begin{equation}\label{eq:pdf2-ym-owner-refined-generator}
 \begin{aligned}
 \widehat H_{m+1}^{\rm ov}
 &=\widehat\Gamma_m^*
   (\widehat H_m^{\rm ov}\otimes I
    +I\otimes\widehat H_{m+1}^{\rm det})\widehat\Gamma_m,\\
 \widehat H_{m+1}^{\rm det}
 &=3\kappa_{\rm av}
   \sum_{\iota\in\mathcal I_{m+1/m}^{\rm det}}
   (I-E_{m+1,\iota}^{\rm det}).
 \end{aligned}
\end{equation}
La partición de colores no se deja nominal. Si una celda se indexa por
$n=(n_1,n_2,n_3,n_4)\in\mathbb Z^4$, se fija
\begin{equation}\label{eq:pdf2-ym-owner-729-coloring}
 \chi_m(n):=
 \bigl(n_1\bmod9,\ n_2\bmod9,\ n_3+3n_4\bmod9\bigr)
 \in(\mathbb Z/9\mathbb Z)^3,
 \qquad |(\mathbb Z/9\mathbb Z)^3|=9\cdot81.
\end{equation}
Si dos collares cerrados de radio celular uno se intersectan, su diferencia
$\delta$ satisface $|\delta_j|\le2$. La igualdad de colores fuerza
$\delta_1=\delta_2=0$ y
$\delta_3+3\delta_4\equiv0\pmod9$. Como el último entero pertenece a
$[-8,8]$, vale cero; las cotas $|\delta_3|,|\delta_4|\le2$ dan entonces
$\delta_3=\delta_4=0$. Por tanto, celdas distintas del mismo color tienen
collares disjuntos. Se define materialmente
\begin{equation}\label{eq:pdf2-ym-owner-729-blocks}
 \mathcal B_{m,r}:=\{\iota\in\mathcal I_m:
   \operatorname{cell}(\iota)=n, \chi_m(n)=r\},
 \qquad r\in(\mathbb Z/9\mathbb Z)^3.
\end{equation}
El transporte de una pantalla por $g\in E(4)$ transporta también esta
partición; no exige que una rotación preserve las coordenadas del rótulo.
Así, $\{\mathcal B_{m,r}\}_{r=1}^{9\cdot81}$ agrupa soportes espaciales
disjuntos;
el índice interno de cada bloque enumera sus factores de
\eqref{eq:pdf2-ym-owner-exhaustive-index}. Para un cilindro $F$, el producto sólo
contiene los bloques de los que depende $F$; la clausura define el producto fuerte
para toda la forma. El generador de un bloque conserva las multiplicidades de
Hoeffding; no se sustituye por un solo proyector que las borraría. Por
conmutatividad en cada carta,
\begin{equation}\label{eq:pdf2-ym-owner-local-circuit}
\begin{aligned}
 H_{m,B}&:=3\kappa_{\rm av}\sum_{\iota\in B}(I-E_{m,\iota}),\\
 K_{m,B}&:=e^{-a_mH_{m,B}}
 =\prod_{\iota\in B}e^{-3\kappa_{\rm av}a_m(I-E_{m,\iota})},\\
 K_m^{\rm loc}&:=\mathop{\prod_{r,B}}^{\rm strong}K_{m,B}
 =e^{-a_m\widehat H_m^{\rm ov}}.
\end{aligned}
\end{equation}
El diámetro físico de cada bloque permanece uniformemente acotado. De aquí procede
la localidad: el generador no introduce una actualización instantánea sobre toda la
red.
La partición es compatible con Bianchi, no sólo con disjunción geométrica. Para una
celda cúbica $c$ se fija el producto orientado, con todos los factores transportados
al mismo vértice,
\begin{equation}\label{eq:pdf2-ym-owner-bianchi-block-product}
 \mathscr B_{m,c}(x):=
 \prod_{p\subset\partial c}^{\longrightarrow}
 T_p\operatorname{Hol}_p(x)T_p^{-1}.
\end{equation}
Cada arista interna de la frontera aparece dos veces y con orientaciones contrarias;
por tanto $\mathscr B_{m,c}=e$ punto a punto. Un bloque
$\mathcal B_{m,r}$ contiene el collar completo de toda celda que actualiza: si toca
una cara de $c$, actualiza simultáneamente el par inverso que comparte esa arista.
La cancelación anterior sigue siendo literal después de cada factor del circuito:
\begin{equation}\label{eq:pdf2-ym-owner-bianchi-block-invariance}
 K_{m,B}\mathbf1_{\{\mathscr B_{m,c}=e\}}
 =\mathbf1_{\{\mathscr B_{m,c}=e\}},
 \qquad
 K_m^{\rm loc}\mathbf1_{\{\mathscr B_{m,c}=e\}}
 =\mathbf1_{\{\mathscr B_{m,c}=e\}}.
\end{equation}
Así el calendario de $9\cdot81$ capas preserva Bianchi en cada paso y en el producto
fuerte, no sólo al final de una barrida.

La acción gauge permuta o traslada las coordenadas integradas por cada
$E_{m,\iota}$ y
preserva sus medidas. Por ello
\begin{equation}\label{eq:pdf2-ym-owner-physical-reduction}
 U_m(g)\widehat H_m^{\rm ov}=\widehat H_m^{\rm ov}U_m(g)
 \quad(g\in\mathcal G_m^{\rm full}),
 \qquad
 [\mathbb P_m^{\rm phys},\widehat H_m^{\rm ov}]=0,
 \qquad
 \mathbb P_{m+1}^{\rm phys}\widehat J_m
 =\widehat J_m\mathbb P_m^{\rm phys}.
\end{equation}
El semigrupo completo y su kernel Markov se denotan por
\begin{equation}\label{eq:pdf2-ym-owner-full-kernel}
 \widehat K_{m,t}^{\rm ov}:=e^{-t\widehat H_m^{\rm ov}},
 \qquad
\widehat K_{m,t}^{\rm ov}(x,dy)
 \quad(x,y\in\widehat X_m).
\end{equation}
La primera identidad de \eqref{eq:pdf2-ym-owner-physical-reduction} implica
\begin{equation}\label{eq:pdf2-ym-owner-full-kernel-equivariance}
 \widehat K_{m,t}^{\rm ov}(gx,gA)
 =\widehat K_{m,t}^{\rm ov}(x,A)
 \quad(g\in\mathcal G_m^{\rm full}).
\end{equation}
Para tipar la reducción sin usar un kernel restringido sobre el espacio equivocado,
sea $q_m:\widehat X_m\to Y_m:=\widehat X_m/\mathcal G_m^{\rm full}$ el mapa de
órbitas y $\nu_m=(q_m)_*\widehat\mu_m^{\rm ov}$. El pullback
\begin{equation}\label{eq:pdf2-ym-owner-quotient-unitary}
 \mathcal W_m:L^2(Y_m,\nu_m)\longrightarrow
 \widehat{\mathscr H}_m^{\rm phys},
 \qquad \mathcal W_m\phi=\phi\circ q_m
\end{equation}
es unitario. La equivariancia del kernel completo permite definir
\begin{equation}\label{eq:pdf2-ym-owner-physical-kernel}
 K_{m,t}^{\rm phys}(q_mx,A)
 :=\widehat K_{m,t}^{\rm ov}(x,q_m^{-1}A),
 \qquad A\subseteq Y_m,
\end{equation}
independientemente del representante $x$. El objeto físico queda entonces tipado por
\begin{equation}\label{eq:pdf2-ym-owner-D}
 D_m^{\rm YM}
 :=\left.\widehat H_m^{\rm ov}\right|_{\widehat{\mathscr H}_m^{\rm phys}},
 \qquad
 e^{-tD_m^{\rm YM}}
 =\mathcal W_mK_{m,t}^{\rm phys}\mathcal W_m^{-1}.
\end{equation}
Éste es el generador producto completo; no es el generador de una sola fibra de
\eqref{eq:pdf2-ym-fibra-semigrupo}.

\section{Acción de Yang--Mills, medida dependiente de $g$ y transporte exacto}

Hasta aquí se ha escrito la carta cinemática de referencia. La denotaremos
$\widehat\mu_m^0:=\widehat\mu_m^{\rm ov}$ y
$\widehat H_m^0:=\widehat H_m^{\rm ov}$, con
$\widehat K_{m,t}^0:=e^{-t\widehat H_m^0}=\widehat K_{m,t}^{\rm ov}$.
Para convertir el lector estructural
$\mathcal P_m^{F^2}$ en una ley, y no en una comparación posterior, se descarga
primero su raíz positiva sobre la factorización triangular. La polarización de
$\mathcal P_m^{F^2}$ en un collar $b$ es el Gram positivo
\[
 Q_{m,b}(u,v):=\frac14\bigl(
 \mathcal P_{m,b}^{F^2}(u+v)-\mathcal P_{m,b}^{F^2}(u-v)
 \bigr).
\]
Se divide el espacio finito de lectores del collar por $\ker Q_{m,b}$, se completa
con $Q_{m,b}$ y se denota por $\mathcal E_{m,b}^{F}$ el Hilbert resultante. Si
$\zeta_{m,b}(x_b)$ es el vector de lectores centrados de la coordenada triangular
$x_b$, se fija materialmente
\begin{equation}\label{eq:pdf2-ym-action-root}
 \xi_{m,b}(x_b):=Q_{m,b}^{1/2}\zeta_{m,b}(x_b)
 \in\mathcal E_{m,b}^{F},
 \qquad
 p_{m,b}(x_b):=\|\xi_{m,b}(x_b)\|^2.
\end{equation}
El índice $b$ recorre los collares Gram de soporte mínimo. Un collar se asigna
al primer color de \eqref{eq:pdf2-ym-owner-729-coloring} que contiene su celda base;
así cada término aparece una vez y los collares de un mismo color son disjuntos.
Más precisamente, la diagonalización positiva agrupa los factores triangulares en
una partición disjunta
\begin{equation}\label{eq:pdf2-ym-action-factor-partition}
 \mathcal I_m^{F}
 =\bigsqcup_{r\in(\mathbb Z/9\mathbb Z)^3}
   \bigsqcup_{b\in\mathcal B_{m,r}^{F}}\mathcal I_{m,b},
 \qquad x_b:=(x_\iota)_{\iota\in\mathcal I_{m,b}},
 \qquad
 \mathcal B_{m,r}^{F}:=\{b:Q_{m,b}\ne0,\ \chi_m(b)=r\}.
\end{equation}
Ninguna coordenada triangular pertenece a dos $\mathcal I_{m,b}$; si un lector
geométrico toca dos collares, su modo de Gram se asigna al primero y su complemento
ortogonal al segundo. Ésta es la descomposición espectral finita de la forma positiva,
no una duplicación de la variable física.
La raíz de la suma ortogonal, calculada color por color, da la identidad, no una
aproximación,
\begin{equation}\label{eq:pdf2-ym-action-reader-sum}
 \mathcal P_m^{F^2}(x)
 =\sum_{r\in(\mathbb Z/9\mathbb Z)^3}
   \sum_{b\in\mathcal B_{m,r}^{F}}p_{m,b}(x_b),
 \qquad
 S_{m,g}(x):=\frac1{4g^2}\mathcal P_m^{F^2}(x),
 \quad g\in(0,\infty).
\end{equation}
Aquí $x_b:\widehat X_m\to X_{m,b}$ es la proyección de coordenadas del
collar, de modo que dominio, acción y codominio de cada sumando quedan fijados.
El lector es gauge invariante porque $Q_{m,b}$ usa el producto invariante de
$\mathfrak g$; es invariante por reflexión porque una reflexión sólo permuta los
seis pares orientados; y es natural por relabelado euclídeo.
La partición \eqref{eq:pdf2-ym-action-factor-partition} y la carta producto dan
explícitamente
\begin{equation}\label{eq:pdf2-ym-action-factor-measures}
 \widehat\mu_m^0=\lambda_m^{\rm inert}\otimes
   \bigotimes_{r,b}\lambda_{m,b}^0,
 \qquad
 d\lambda_{m,b}^g
 :=(Z_{m,b}^g)^{-1}e^{-p_{m,b}/(4g^2)}d\lambda_{m,b}^0,
 \qquad
 \widehat\mu_{m,g}^{\rm YM}=\lambda_m^{\rm inert}\otimes
   \bigotimes_{r,b}\lambda_{m,b}^g.
\end{equation}
El factor inerte contiene exactamente las coordenadas sobre las que el lector es
cero; no se elimina de la teoría.

La separación grueso--detalle de
\eqref{eq:pdf2-ym-owner-product-factorization} es ortogonal para la misma raíz:
la polarización directa da
$Q_{m+1}=Q_m\oplus Q_{m+1/m}^{\rm det}$ y
$\zeta_{m+1}=\zeta_m\oplus\zeta_{m+1/m}^{\rm det}$.
Si $x_{m+1}=\widehat\Gamma_m^{-1}(x_m,u)$, entonces
\begin{equation}\label{eq:pdf2-ym-action-cocycle}
 \xi_{m+1}(x_{m+1})
 =\xi_m(x_m)\oplus\xi_{m+1/m}^{\rm det}(u),
 \qquad
 S_{m+1,g}\circ\widehat\Gamma_m^{-1}
 =S_{m,g}+S_{m+1/m,g}^{\rm det}.
\end{equation}
Esta es la forma precisa en que carry y memoria impiden contar dos veces una cara
ancestral: la parte gruesa no se vuelve a sumar en el detalle. En particular,
$Z_{m+1,g}=Z_{m,g}Z_{m+1/m,g}^{\rm det}$ para los normalizadores y se obtiene la
familia de probabilidades genuinamente dependiente del acoplamiento
\begin{equation}\label{eq:pdf2-ym-gibbs-measure}
 d\widehat\mu_{m,g}^{\rm YM}(x)
 :=Z_{m,g}^{-1}e^{-S_{m,g}(x)}d\widehat\mu_m^0(x),
 \qquad
 (\widehat\pi_{m+1,m})_*\widehat\mu_{m+1,g}^{\rm YM}
 =\widehat\mu_{m,g}^{\rm YM}.
\end{equation}
Para escribir correctamente Schwinger--Dyson respecto de Haar, y no fingir que la
medida de referencia es plana, sea $\lambda_m$ el producto de Haar, Lebesgue y
conteo uniforme de la carta levantada, y sea
$R_m:=d\widehat\mu_m^0/d\lambda_m>0$. La acción total regulada es
\begin{equation}\label{eq:pdf2-ym-total-regulated-action}
 \mathscr S_{m,g}:=S_{m,g}-\log R_m,
 \qquad
 d\widehat\mu_{m,g}^{\rm YM}
 =\mathscr Z_{m,g}^{-1}e^{-\mathscr S_{m,g}}d\lambda_m.
\end{equation}
El término $-\log R_m$ contiene los kernels de calor y las cartas de puente de la
estructura gauge de partida; es independiente de $g$. El término de Yang--Mills que
cambia el acoplamiento y cuyo límite clásico se calcula abajo es exactamente
$S_{m,g}=\mathcal P_m^{F^2}/(4g^2)$.
La consistencia de la segunda igualdad y la estandaridad de los espacios finitos
permiten aplicar Kolmogórov y definen la única medida cilíndrica
$\widehat\mu_{\infty,g}^{\rm YM}$ cuyas marginales son
$\widehat\mu_{m,g}^{\rm YM}$.
No se ha usado aquí ningún autovalor: primero se han fijado $g$, la acción y
la medida.

Para definir la dinámica sobre esa medida sin perder producto, gauge ni las
identidades superficiales, se construye el transporte, no se lo postula. En cada
factor no atómico $X_{m,b}$ se fija una carta boreliana
$\rho_{m,b}:X_{m,b}/\mathcal G_{m,b}\to(0,1)$ y se desintegra primero sobre el
cociente gauge y después sobre la órbita de Haar. Sean $F_{m,b}^0$ y
$F_{m,b}^g$ las funciones de distribución continuas de $\rho_{m,b}$ para la
medida de referencia y la medida reponderada. El transporte triangular es
\begin{equation}\label{eq:pdf2-ym-gibbs-transport-factor}
 t_{m,b,g}:=(F_{m,b}^g)^{-1}\circ F_{m,b}^0,
 \qquad
 T_{m,b,g}(y,h):=(t_{m,b,g}(y),h),
 \qquad
 (T_{m,b,g})_*\lambda_{m,b}^0=\lambda_{m,b}^g.
\end{equation}
En un factor puramente atómico sobre el que $p_{m,b}=0$ se toma
$T_{m,b,g}=I$. En presencia de átomos y una coordenada continua, la carta es el
orden lexicográfico átomo--intervalo y la misma fórmula se aplica al espacio
estándar no atómico total. Los transportes de los $9\cdot81$ colores se componen
en el orden fijado por \eqref{eq:pdf2-ym-owner-729-coloring}; dentro de un color
conmutan porque los collares son disjuntos. Se obtiene
\begin{equation}\label{eq:pdf2-ym-gibbs-transport-global}
 T_{m,g}:=\prod_{r\in(\mathbb Z/9\mathbb Z)^3}
           \prod_{b\in\mathcal B_{m,r}^{F}}T_{m,b,g},
 \quad
 (T_{m,g})_*\widehat\mu_m^0=\widehat\mu_{m,g}^{\rm YM},
 \quad
 \widehat\pi_{m+1,m}T_{m+1,g}=T_{m,g}\widehat\pi_{m+1,m}.
\end{equation}
La última igualdad se comprueba factor a factor usando
\eqref{eq:pdf2-ym-action-cocycle}: el transporte del factor grueso es el ya fijado y
los transportes nuevos actúan sólo en la innovación. Las cartas de cociente se
eligen en una pantalla representante y se trasladan; por ello $T_{m,g}$ conmuta con
gauge y reflexión y, para $a\in E(4)$, satisface
$T_{a\alpha,g}=U_\alpha(a)T_{\alpha,g}U_\alpha(a)^{-1}$.

La composición
\begin{equation}\label{eq:pdf2-ym-full-probability-isomorphism}
 \mathcal T_{m,g}:L^\infty(\widehat X_m,\widehat\mu_{m,g}^{\rm YM})
 \longrightarrow L^\infty(\widehat X_m,\widehat\mu_m^0),
 \qquad \mathcal T_{m,g}F:=F\circ T_{m,g}
\end{equation}
es, por construcción, un isomorfismo de espacios de probabilidad y de
$*$-álgebras, y se prolonga unitariamente a $L^2$. En particular,
\begin{equation}\label{eq:pdf2-ym-full-algebra-identities}
 \mathcal T_{m,g}(FG)=\mathcal T_{m,g}F\,\mathcal T_{m,g}G,
 \quad
 \mathcal T_{m,g}(F^*)=(\mathcal T_{m,g}F)^*,
 \quad
 \int F\,d\widehat\mu_{m,g}^{\rm YM}
 =\int\mathcal T_{m,g}F\,d\widehat\mu_m^0.
\end{equation}
El dominio de cada $T_{m,b,g}$ es el propio espacio de holonomías que satisface
la ley de subdivisión y la restricción de Bianchi. Por tanto su imagen permanece
en ese espacio. Aplicando la multiplicatividad anterior generador a generador se
obtienen las identidades puntuales
\begin{equation}\label{eq:pdf2-ym-transport-surface-bianchi}
 \begin{aligned}
 \mathcal T_{m,g}(\operatorname{Hol}_p)
 &=\prod_{p'\subset p}^{\longrightarrow}
   \mathcal T_{m,g}(T_{p'}\operatorname{Hol}_{p'}T_{p'}^{-1}),\\
 \mathcal T_{m,g}(F_{{\rm cl},m})
 &=F_{{\rm cl},m}\circ T_{m,g},\\
 \prod_{p\subset\partial c}^{\longrightarrow}
 \mathcal T_{m,g}(T_p\operatorname{Hol}_pT_p^{-1})&=e,
 \qquad
 \operatorname{Alt}_{\lambda\mu\nu}
 D_\lambda(F_{{\rm cl},m}\circ T_{m,g})_{\mu\nu}=0.
 \end{aligned}
\end{equation}
La tercera línea es la cancelación arista por arista de la frontera orientada de
la frontera de $c$; la cuarta es su polarización infinitesimal. Así el transporte
preserva productos, estado, clover y Bianchi, no sólo el primer caos.

Finalmente se transportan expectativas, generador, kernel, proyección e inclusión:
\begin{equation}\label{eq:pdf2-ym-g-dynamics}
 \begin{aligned}
 E_{m,\iota}^{g}&:=\mathcal T_{m,g}^{-1}E_{m,\iota}\mathcal T_{m,g},&
 \widehat H_{m,g}^{\rm YM}&:=\mathcal T_{m,g}^{-1}\widehat H_m^0\mathcal T_{m,g},\\
 \widehat K_{m,t}^{g}&:=\mathcal T_{m,g}^{-1}\widehat K_{m,t}^{0}\mathcal T_{m,g},&
 \mathbb P_{m,g}^{\rm phys}&:=\mathcal T_{m,g}^{-1}
     \mathbb P_m^{\rm phys}\mathcal T_{m,g},\\
 \widehat J_{m,g}&:=\mathcal T_{m+1,g}^{-1}
     \widehat J_m\mathcal T_{m,g},&
 D_{m,g}^{\rm YM}&:=\left.\widehat H_{m,g}^{\rm YM}
     \right|_{\operatorname{Ran}\mathbb P_{m,g}^{\rm phys}}.
 \end{aligned}
\end{equation}
La primera línea no es sólo una conjugación abstracta. Usando
\eqref{eq:pdf2-ym-action-factor-measures}, su acción sobre un cilindro es la
actualización de baño térmico de la acción:
\begin{equation}\label{eq:pdf2-ym-gibbs-conditional-expectation}
 (E_{m,b}^{g}F)(x_{\ne b})
 =\frac{\displaystyle
   \int_{X_{m,b}}F(x_{\ne b},u)e^{-p_{m,b}(u)/(4g^2)}
       d\lambda_{m,b}^0(u)}{\displaystyle
   \int_{X_{m,b}}e^{-p_{m,b}(u)/(4g^2)}d\lambda_{m,b}^0(u)},
 \qquad
 \widehat H_{m,g}^{\rm YM}
 =3\kappa_{\rm av}\sum_b(I-E_{m,b}^{g}).
\end{equation}
Así tanto la medida como cada transición local contienen $g$ y la acción
$p_{m,b}/(4g^2)$ de forma explícita.
Son Markov porque $\mathcal T_{m,g}\mathbf1=\mathbf1$, locales porque cada
$T_{m,b,g}$ tiene el mismo collar que $p_{m,b}$, y dependen de $g$ por
\eqref{eq:pdf2-ym-gibbs-transport-factor}. La equivalencia es espectral y no
selecciona la medida a partir del espectro:
Para cerrar además las ecuaciones dinámicas, sea $X_{m,b}^a$ el campo
izquierdo-invariante en el factor de enlace $b$ asociado a
$e_a\in\mathfrak g$; en una carta continua de innovación se usa la derivada
direccional ordinaria, y las coordenadas finitas se tratan por diferencia exacta.
Para todo cilindro suave $F$ el dominio común es
$\mathscr C_m^1:=\bigcap_{b,a}\operatorname{Dom}X_{m,b}^a$, tomando soporte
compacto en las cartas abiertas de innovación; en los factores de grupo no hay
frontera. La invariancia de Haar y la regla del producto calculan
\begin{equation}\label{eq:pdf2-ym-schwinger-dyson-finite}
 \int X_{m,b}^aF\,d\widehat\mu_{m,g}^{\rm YM}
 =\int F\,X_{m,b}^a\mathscr S_{m,g}\,d\widehat\mu_{m,g}^{\rm YM},
 \qquad F\in\mathscr C_m^1.
\end{equation}
En una coordenada finita, para
$\Delta_\varepsilon F(q):=F(q+\varepsilon)-F(q)$, el cambio de variable
$q\mapsto q+\varepsilon$ da la versión exacta, sin derivada formal,
\begin{equation}\label{eq:pdf2-ym-schwinger-dyson-discrete}
 \int\Delta_\varepsilon F(q)\,d\widehat\mu_{m,g}^{\rm YM}
 =\int F(q)
  \left(e^{\mathscr S_{m,g}(q)-\mathscr S_{m,g}(q-\varepsilon)}-1\right)
  d\widehat\mu_{m,g}^{\rm YM}(q).
\end{equation}
Si $\delta_\varepsilon$ es la suma orientada de esos campos en las aristas que
inciden en un vértice, la invariancia gauge de la medida de referencia y de
$S_{m,g}$ da $\delta_\varepsilon\mathscr S_{m,g}=0$ y, por tanto, la identidad de Ward
\begin{equation}\label{eq:pdf2-ym-ward-finite}
 \int\delta_\varepsilon F\,d\widehat\mu_{m,g}^{\rm YM}=0.
\end{equation}
Para el factor gauge finito $\mathbb F_3^4$, la invariancia
$\mathscr S_{m,g}(q+\varepsilon)=\mathscr S_{m,g}(q)$ reduce
\eqref{eq:pdf2-ym-schwinger-dyson-discrete} a la misma identidad Ward con
$\delta_\varepsilon$ reemplazada por $\Delta_\varepsilon$.
Las dos expresiones son compatibles con $\widehat J_{m,g}$ porque
\eqref{eq:pdf2-ym-action-cocycle} separa las derivadas gruesas de las de detalle.
Todo cilindro del límite reside en algún nivel; así
\eqref{eq:pdf2-ym-schwinger-dyson-finite} y
\eqref{eq:pdf2-ym-ward-finite} pasan literalmente al estado cofinal. Ni la
ecuación Schwinger--Dyson ni Ward contienen la brecha como premisa.

\begin{equation}\label{eq:pdf2-ym-g-gap-before-limit}
 D_{m,g}^{\rm YM}\succeq3\kappa_{\rm av}
 (I-P_{\Omega_{m,g}}),
 \qquad
 W_{c,g}:=\mathcal T_{m,g}^{-1}W_c,
 \qquad
 D_{m,g}^{\rm YM}W_{c,g}=3\kappa_{\rm av}W_{c,g}.
\end{equation}
En lo que sigue se fija un $g\in(0,\infty)$ y se suprime ese subíndice; las
medidas, expectativas, inclusiones y generadores sin superíndice son los objetos
de \eqref{eq:pdf2-ym-g-dynamics}. Las holonomías geométricas permanecen como
funciones coordenadas canónicas sobre el espacio de configuraciones y, por ello, su
ley sí cambia con $g$; $\mathcal T_{m,g}$ las lleva a funciones compuestas de la
carta de referencia. Sólo el testigo espectral se transporta explícitamente en
\eqref{eq:pdf2-ym-g-gap-before-limit}. Esta convención no borra $g$: su residencia explícita es
\eqref{eq:pdf2-ym-gibbs-measure}--\eqref{eq:pdf2-ym-g-dynamics}.

\section{Vacío simple y testigo que alcanza la brecha}

La intersección de todas las expectativas del estado completo es la constante:
\begin{equation}\label{eq:pdf2-ym-owner-vacuum}
 \bigcap_{\iota\in\mathcal I_m}\operatorname{Ran}E_{m,\iota}
 =\mathbb C\Omega_m.
\end{equation}
Esta identidad no se introduce como hipótesis. En la carta producto
\eqref{eq:pdf2-ym-owner-product-factorization}, cada $E_{m,\iota}$ integra una
coordenada del índice exhaustivo
\eqref{eq:pdf2-ym-owner-exhaustive-index} y deja fijas las restantes. Si
$u$ pertenece a la intersección, sus expectativas condicionales respecto de todo
cilindro finito son constantes. La convergencia martingala de esos cilindros a la
$\sigma$-álgebra completa da que $u$ es constante casi en todas partes. La inclusión
inversa es inmediata porque toda expectativa conserva la unidad. Esto prueba
\eqref{eq:pdf2-ym-owner-vacuum} y, después de la proyección gauge, deja un vacío
físico unidimensional.
La descomposición de Hoeffding de
\eqref{eq:pdf2-ym-owner-base-generator}--
\eqref{eq:pdf2-ym-owner-refined-generator} da
\begin{equation}\label{eq:pdf2-ym-owner-hoeffding}
 \widehat{\mathcal E}_m(u)
 =3\kappa_{\rm av}\sum_{S\Subset\mathcal I_m}|S|\,\|u_S\|^2,
 \qquad
 D_m^{\rm YM}\succeq
 3\kappa_{\rm av}(I-P_{\Omega_m}).
\end{equation}
La función
\begin{equation}\label{eq:pdf2-ym-owner-wc}
 W_c(q):=\mathbf1_{\{q_1+\cdots+q_6=0\}}-\frac13
\end{equation}
es invariante bajo $G^{V(\Lambda_m)}$ porque no depende de enlaces ni marcos. Para
la acción de incidencia, cada $x_i$ aparece tres veces, luego
\begin{equation}\label{eq:pdf2-ym-owner-wc-invariance}
 \sum_{i<j}(B^+x)_{ij}=3\sum_i x_i=0
 \quad\text{en }\mathbb F_3.
\end{equation}
Así $\mathbb P_m^{\rm phys}W_c=W_c$. Bajo $u_6$, la suma de los seis trits es
uniforme; por cálculo directo,
\begin{equation}\label{eq:pdf2-ym-owner-wc-moments}
 \mathbb EW_c=0,\qquad
 \|W_c\|^2=\frac29,\qquad
 \mathbb E(W_c^3)=\frac2{27}.
\end{equation}
Sólo el proyector $I-E_{m,q_c}$ actúa no trivialmente sobre $W_c$. Por tanto
\begin{equation}\label{eq:pdf2-ym-owner-wc-eigenvalue}
 D_m^{\rm YM}W_c=3\kappa_{\rm av}W_c.
\end{equation}
La cota inferior de \eqref{eq:pdf2-ym-owner-hoeffding} no es sólo una cota: el
complemento físico contiene un vector no nulo que alcanza el umbral.

\section{Refinamiento nonádico y persistencia en el límite}

Las inclusiones físicas satisfacen
\begin{equation}\label{eq:pdf2-ym-owner-intertwining}
 D_{m+1}^{\rm YM}\widehat J_m
 =\widehat J_mD_m^{\rm YM},\qquad
 \bigl(e^{-a_{m+1}D_{m+1}^{\rm YM}}\bigr)^9\widehat J_m
 =\widehat J_me^{-a_mD_m^{\rm YM}},\qquad a_{m+1}=a_m/9.
\end{equation}
En el límite inductivo
\begin{equation}\label{eq:pdf2-ym-owner-inductive-limit}
 \widehat{\mathscr H}_\infty^{\rm phys}
 =\varinjlim(\widehat{\mathscr H}_m^{\rm phys},\widehat J_m)
\end{equation}
se define sobre la unión algebraica
\begin{equation}\label{eq:pdf2-ym-owner-limit-form}
 \mathcal E_\infty[\widehat J_{m,\infty}u]
 :=\langle u,D_m^{\rm YM}u\rangle.
\end{equation}
Su generador se denota
$D_{\infty,g}^{\rm YM}$; de acuerdo con la convención posterior a
\eqref{eq:pdf2-ym-g-gap-before-limit}, se escribe también
$D_\infty^{\rm YM}:=D_{\infty,g}^{\rm YM}$.
El primer entrelazamiento de \eqref{eq:pdf2-ym-owner-intertwining} hace que la
definición sea independiente del representante. Las truncaciones de Hoeffding son
un núcleo de forma y dan una sucesión de recuperación; Fatou en la suma no negativa
da la semicontinuidad inferior. La clausura posee, pues,
\begin{equation}\label{eq:pdf2-ym-owner-limit-gap}
 \overline{\mathcal E}_\infty[u]
 \ge3\kappa_{\rm av}\|(I-P_{\Omega_\infty})u\|^2.
\end{equation}
La inclusión conserva las coordenadas ancestrales, de modo que
$\widehat J_{m,\infty}W_c$ sigue siendo no nulo y satisface
\begin{equation}\label{eq:pdf2-ym-owner-limit-witness}
 D_\infty^{\rm YM}\widehat J_{m,\infty}W_c
 =3\kappa_{\rm av}\widehat J_{m,\infty}W_c.
\end{equation}
En consecuencia,
\begin{equation}\label{eq:pdf2-ym-owner-exact-gap}
 \inf\bigl(\operatorname{Spec}D_\infty^{\rm YM}\setminus\{0\}\bigr)
 =3\kappa_{\rm av}>0,
 \qquad \ker D_\infty^{\rm YM}=\mathbb C\Omega_\infty.
\end{equation}

\section{Cuatro reconstrucciones OS de la misma dinámica}

El kernel completo de \eqref{eq:pdf2-ym-owner-full-kernel} define una ley reversible
sobre $\widehat X_m$. Para la reconstrucción física se usa su descenso bien tipado
\eqref{eq:pdf2-ym-owner-physical-kernel}: sobre el cociente $Y_m$ se define
\begin{equation}\label{eq:pdf2-ym-owner-eight-face-law}
 \Omega_{m,{\rm phys}}^{(8)}(d\bar y_1\cdots d\bar y_8)
 =\int\nu_m(d\bar z)
  \prod_{r=1}^8K_{m,a_m/2}^{\rm phys}(\bar z,d\bar y_r).
\end{equation}
Para cada dirección $\nu=1,\ldots,4$, la reflexión intercambia las cuatro caras
positivas con las negativas. Condicionando respecto del estado central $z$ y usando
reversibilidad,
\begin{equation}\label{eq:pdf2-ym-owner-four-os}
 \int\overline{F(\Theta_{\nu,m}y)}F(y)\,
 d\Omega_{m,{\rm phys}}^{(8)}(y)
 =\|\mathcal B_{m,\nu}F\|_2^2\ge0.
\end{equation}
La marginal de la pareja opuesta normal a $\nu$ es
\begin{equation}\label{eq:pdf2-ym-owner-os-marginal}
 \nu_m(d\bar x)K_{m,a_m}^{\rm phys}(\bar x,d\bar y).
\end{equation}
Para identificar el Hamiltoniano de una dirección se toma el subespacio OS de
observables de una cara positiva simple,
$F(y)=F_\nu(\bar y_{\nu,+})$; los observables de varias caras conservan la
positividad anterior, pero no se confunden con este espacio de transferencia. En el
cociente de cara simple, el mapa
\begin{equation}\label{eq:pdf2-ym-owner-os-unitary}
 \mathcal V_{m,\nu}[F]
 :=e^{-a_mD_m^{\rm YM}/2}\mathcal W_mF
\end{equation}
es isométrico para la norma OS y tiene rango denso. Su extensión identifica el
cociente OS con $\widehat{\mathscr H}_m^{\rm phys}$. Póngase
$\overline D_m:=\mathcal W_m^{-1}D_m^{\rm YM}\mathcal W_m$ y
$\overline J_m:=\mathcal W_{m+1}^{-1}\widehat J_m\mathcal W_m$. El
entrelazamiento y $a_m=9a_{m+1}$ dan
\begin{equation}\label{eq:pdf2-ym-owner-os-range-inclusion}
 \overline J_me^{-a_m\overline D_m/2}
 =e^{-a_m\overline D_{m+1}/2}\overline J_m
 =\bigl(e^{-a_{m+1}\overline D_{m+1}/2}\bigr)^9\overline J_m.
\end{equation}
El último miembro pertenece al rango de
$e^{-a_{m+1}\overline D_{m+1}/2}$; por eso el inverso de $\mathcal V_{m+1,\nu}$
que aparece a continuación se aplica sólo en su rango y está bien definido. Los
conectores
\begin{equation}\label{eq:pdf2-ym-owner-os-connectors}
 \mathcal J_{m,\nu}^{\rm OS}
 :=\mathcal V_{m+1,\nu}^{-1}\widehat J_m\mathcal V_{m,\nu}
\end{equation}
son isometrías y entrelazan los cuatro Hamiltonianos reconstruidos:
\begin{equation}\label{eq:pdf2-ym-owner-common-os-hamiltonian}
 H_{{\rm OS},m+1}^{(\nu)}\mathcal J_{m,\nu}^{\rm OS}
 =\mathcal J_{m,\nu}^{\rm OS}H_{{\rm OS},m}^{(\nu)},
 \qquad
 \mathcal V_{\infty,\nu}H_{{\rm OS},\infty}^{(\nu)}
 \mathcal V_{\infty,\nu}^{-1}=D_\infty^{\rm YM}.
\end{equation}
Las cuatro reflexiones no fabrican cuatro medidas: son cuatro publicaciones del mismo
estado, del mismo kernel y del mismo generador.

\section{Localidad, acción euclídea y publicación de curvatura}

La categoría dirigida de pantallas incluye refinamientos, overlays y transportes
euclídeos. Si $a\in E(4)$ y $\alpha$ es una pantalla, el relabelado completo define
\begin{equation}\label{eq:pdf2-ym-owner-e4-finite-transport}
 U_\alpha(a):\widehat{\mathscr H}_\alpha^{\rm phys}
 \longrightarrow\widehat{\mathscr H}_{a\alpha}^{\rm phys},
 \qquad
 U_\beta(a)\widehat J_{\alpha\beta}
 =\widehat J_{a\alpha,a\beta}U_\alpha(a).
\end{equation}
El transporte preserva medida, generador, ideales nulos y soportes. La identidad de
naturalidad permite definir, sin elegir representante, una unitaria sobre el colímite
algebraico:
\begin{equation}\label{eq:pdf2-ym-owner-e4-colimit-action}
 U(a)\widehat J_{\alpha,\infty}F
 :=\widehat J_{a\alpha,\infty}U_\alpha(a)F.
\end{equation}
Los relabelados satisfacen $U_{a\alpha}(b)U_\alpha(a)=U_\alpha(ba)$ y preservan
norma; por ello \eqref{eq:pdf2-ym-owner-e4-colimit-action} se prolonga a una
representación unitaria de $E(4)$.

\subsection{El canal de incidencia es curvatura local, no un factor auxiliar}

Todo grupo compacto, conectado y simple posee un toro maximal no trivial. Se fija
$t_3\in G$ de orden exactamente tres. Para un collar $c$, sea
\begin{equation}\label{eq:pdf2-ym-owner-incidence-flux}
 \omega_c(q):=\sum_{i<j}q_{ij}\in\mathbb F_3.
\end{equation}
La identidad \eqref{eq:pdf2-ym-incidence-conservation} prueba que $\omega_c$ es
invariante bajo la acción finita. Si $h_c$ es el marco en el vértice base del collar,
se define la holonomía geométrica y, para el testigo espectral, su imagen transportada
en la ley a acoplamiento $g$:
\begin{equation}\label{eq:pdf2-ym-owner-incidence-holonomy}
 \operatorname{Hol}_{m,c}^{\rm inc,geom}:=h_ct_3^{\,\omega_c(q)}h_c^{-1}\in G,
 \qquad
 \operatorname{Hol}_{m,c,g}^{\rm inc,spec}
 :=\mathcal T_{m,g}^{-1}\operatorname{Hol}_{m,c}^{\rm inc,geom},
 \qquad
 \operatorname{Hol}_{m,c}^{\rm inc}:=\operatorname{Hol}_{m,c,g}^{\rm inc,spec}.
\end{equation}
Bajo $G^{V_m}$ se transforma por conjugación y bajo
$\prod_c\mathbb F_3^4$ permanece fija. Como $t_3$ tiene orden tres, el cálculo
funcional de esta variable de tres valores da la identidad literal
\begin{equation}\label{eq:pdf2-ym-owner-wc-curvature-functional}
 W_{c,g}:=\mathcal T_{m,g}^{-1}W_c^0
 =\mathbf1_{\{e\}}\!\left(\operatorname{Hol}_{m,c}^{\rm inc}\right)-\frac13.
\end{equation}
Se escribe $W_c:=W_{c,g}$ en el resto del capítulo.
Como $T_{m,b,g}$ actúa dentro del cociente del mismo collar, la holonomía espectral
es una función de su álgebra local de holonomías geométricas; no añade una variable
ni cambia el soporte. Las holonomías geométricas, usadas en $S_{m,g}$ y en el clover,
mantienen su distribución dependiente de $g$.
El proyector del miembro derecho es una función de clase local; por tanto $W_c$ no
es una coordenada auxiliar desacoplada, sino un observable de la holonomía de
incidencia de la misma restricción gauge.

Para un abierto acotado $O\subset\mathbb R^4$, sea
$\mathfrak A_{\rm YM}(O)$ el álgebra generada por funciones de clase acotadas de
\eqref{eq:pdf2-ym-owner-incidence-holonomy}, holonomías de plaqueta y lectores
clover con soporte en $O$. Se define el sector de curvatura física por
\begin{equation}\label{eq:pdf2-ym-owner-curvature-cyclic-sector}
 \mathscr H_{\rm YM}^{\rm curv}
 :=\overline{\bigcup_{O\Subset\mathbb R^4}
   \mathfrak A_{\rm YM}(O)\Omega_\infty}
 \subseteq\widehat{\mathscr H}_\infty^{\rm phys}.
\end{equation}
Por \eqref{eq:pdf2-ym-owner-wc-curvature-functional}, cada
$\widehat J_{m,\infty}W_c$ pertenece materialmente a este sector.
Las expectativas de coordenada envían cilindros de curvatura a cilindros de
curvatura; por tanto el semigrupo deja invariante
$\mathscr H_{\rm YM}^{\rm curv}$. Al ser autoadjunto, el sector es reductor. La
cota de \eqref{eq:pdf2-ym-owner-limit-gap} se restringe a él y el vector anterior
alcanza la igualdad:
\begin{equation}\label{eq:pdf2-ym-owner-curvature-sector-gap}
 D_{\rm curv,g}^{\rm YM}
 :=\left.D_{\infty,g}^{\rm YM}\right|_{\mathscr H_{\rm YM}^{\rm curv}},
 \qquad D_{\rm curv}^{\rm YM}:=D_{\rm curv,g}^{\rm YM},
 \qquad
 \inf\bigl(\operatorname{Spec}D_{\rm curv,g}^{\rm YM}\setminus\{0\}\bigr)
 =3\kappa_{\rm av}.
\end{equation}

\subsection{Sistema cofinal de publicación y productos cruzados}

La inclusión $\widehat J_{\alpha\beta}$ conserva una coordenada ancestral; no debe
confundirse con el promedio de bloque que publica un campo continuo. Para escribir
este segundo mapa sin alterar la genealogía, sea $\mathfrak P$ el conjunto dirigido
de pantallas celulares acotadas, cerrado bajo overlay, refinamiento nonádico y
transporte euclídeo. Para $\alpha\in\mathfrak P$ póngase
\begin{equation}\label{eq:pdf2-ym-owner-cell-step-space}
 V_\alpha:=\operatorname{span}\left\{
 e_{\alpha,c}:=|c|^{-1/2}\mathbf1_c:
 c\in\mathcal C_\alpha^{\rm cell}\right\}
 \subset L^2(\mathbb R^4),
 \qquad P_\alpha:L^2(\mathbb R^4)\to V_\alpha.
\end{equation}
Los overlays hacen que $P_\alpha h\to h$ para todo $h\in L^2$ a lo largo de
$\mathfrak P$. En el lado gauge se normaliza el testigo por
\begin{equation}\label{eq:pdf2-ym-owner-normalized-cell-witness}
 Z_{\alpha,c}:=\sqrt{\frac92}\,W_{\alpha,c},
 \qquad
 \langle Z_{\alpha,c},Z_{\alpha,d}\rangle=\delta_{cd},
 \qquad
 D_\alpha^{\rm YM}Z_{\alpha,c}=3\kappa_{\rm av}Z_{\alpha,c}.
\end{equation}
La segunda igualdad usa exactamente los momentos de
\eqref{eq:pdf2-ym-owner-wc-moments} y la factorización de coordenadas distintas.
Sea $\mathcal K_\alpha^{\rm cell}$ el espacio generado por esos vectores. Si
$\beta\succeq\alpha$, el entrelazador de publicación, distinto de
$\widehat J_{\alpha\beta}$, se define generador a generador por
\begin{equation}\label{eq:pdf2-ym-owner-block-intertwiner}
 I_{\alpha\beta}^{\rm curv}Z_{\alpha,c}
 :=\sum_{\substack{d\in\mathcal C_\beta^{\rm cell}\\d\subset c}}
       \sqrt{\frac{|d|}{|c|}}\,Z_{\beta,d}.
\end{equation}
Como las subceldas son disjuntas y suman $|c|$, se obtiene, sin tomar límites,
\begin{equation}\label{eq:pdf2-ym-owner-block-intertwiner-identities}
 (I_{\alpha\beta}^{\rm curv})^*I_{\alpha\beta}^{\rm curv}=I,
 \quad
 I_{\beta\gamma}^{\rm curv}I_{\alpha\beta}^{\rm curv}
 =I_{\alpha\gamma}^{\rm curv},
 \quad
 D_\beta^{\rm YM}I_{\alpha\beta}^{\rm curv}
 =I_{\alpha\beta}^{\rm curv}D_\alpha^{\rm YM}.
\end{equation}
El relabelado celular conserva volúmenes y descendientes; para $a\in E(4)$ se tiene además
\begin{equation}\label{eq:pdf2-ym-owner-block-e4-naturality}
 U_\beta(a)I_{\alpha\beta}^{\rm curv}
 =I_{a\alpha,a\beta}^{\rm curv}U_\alpha(a).
\end{equation}

Para $h\in C_c^\infty(\mathbb R^4)$ defínase ahora el lector de bloque
\begin{equation}\label{eq:pdf2-ym-owner-smeared-witness}
 \Phi_\alpha(h):=
 \sum_{c\in\mathcal C_\alpha^{\rm cell}}
 \langle h,e_{\alpha,c}\rangle_{L^2}Z_{\alpha,c}.
\end{equation}
En una malla uniforme esta fórmula es
$\sqrt{9/2}\,a_\alpha^2\sum_c(P_\alpha h)(x_c)W_{\alpha,c}$; el promedio celular,
no una evaluación puntual, fija el producto cruzado. En efecto, si $\gamma$ refina
dos pantallas $\alpha$ y $\beta$, las definiciones anteriores dan la identidad
calculada
\begin{equation}\label{eq:pdf2-ym-owner-smeared-cross-gram}
 \left\langle
  I_{\alpha\gamma}^{\rm curv}\Phi_\alpha(h),
  I_{\beta\gamma}^{\rm curv}\Phi_\beta(k)
 \right\rangle
 =\langle P_\alpha h,P_\beta k\rangle_{L^2}.
\end{equation}
Por polarización, con $k=h$,
\begin{equation}\label{eq:pdf2-ym-owner-smeared-cauchy-derived}
 \left\|
  I_{\alpha\gamma}^{\rm curv}\Phi_\alpha(h)
  -I_{\beta\gamma}^{\rm curv}\Phi_\beta(h)
 \right\|_2^2
 =\|P_\alpha h-P_\beta h\|_{L^2}^2\longrightarrow0.
\end{equation}
Así la propiedad de Cauchy es consecuencia de la convergencia fuerte de
proyecciones, no una identidad adicional postulada. El colímite de
\eqref{eq:pdf2-ym-owner-block-intertwiner} contiene, por tanto,
\begin{equation}\label{eq:pdf2-ym-owner-smeared-gram}
 \Phi(h):=\lim_\alpha I_{\alpha,\infty}^{\rm curv}\Phi_\alpha(h),
 \qquad
 \langle\Phi(h),\Phi(k)\rangle=\langle h,k\rangle_{L^2},
 \qquad
 D_{\rm pub}^{\rm YM}\Phi(h)=3\kappa_{\rm av}\Phi(h).
\end{equation}

Esta publicación no añade grados de libertad. En cada paso, la descomposición
$V_\beta=V_\alpha\oplus(V_\beta\ominus V_\alpha)$ tiene una base de Haar
nonádica. El orden exhaustivo de incidencias nuevas de
\eqref{eq:pdf2-ym-owner-exhaustive-index} empareja esa base, etiqueta por etiqueta,
con coordenadas $Z_{\beta,d}$ del factor de innovación, conservando el soporte de la
función de Haar. Las unitarias triangulares $\widehat\Gamma_\alpha$ pegan esos
emparejamientos y producen una isometría basada
\begin{equation}\label{eq:pdf2-ym-owner-publication-into-physical}
 \mathcal U_{\rm curv}:
 \varinjlim(\mathcal K_\alpha^{\rm cell},I_{\alpha\beta}^{\rm curv})
 \longrightarrow\mathscr H_{\rm YM}^{\rm curv},
 \quad
 D_{\rm curv}^{\rm YM}\mathcal U_{\rm curv}
 =\mathcal U_{\rm curv}D_{\rm pub}^{\rm YM},
 \quad U(a)\mathcal U_{\rm curv}=\mathcal U_{\rm curv}U_{L^2}(a).
\end{equation}
En adelante $\Phi(h)$ denota también su imagen por $\mathcal U_{\rm curv}$.
La compatibilidad de \eqref{eq:pdf2-ym-gibbs-transport-global} define
$\mathcal T_{\infty,g}$ en el límite. La publicación física al acoplamiento fijado
no es una segunda isometría escogida sólo sobre el primer caos, sino la restricción
\begin{equation}\label{eq:pdf2-ym-full-publication-restriction}
 \mathcal U_{{\rm curv},g}
 :=\mathcal T_{\infty,g}^{-1}\mathcal U_{{\rm curv},0},
 \qquad
 \widetilde{\mathcal U}_g
 :=\mathcal T_{\infty,g}^{-1}:
 L^\infty(\widehat X_\infty,\widehat\mu_\infty^0)
 \longrightarrow
 L^\infty(\widehat X_\infty,\widehat\mu_{\infty,g}^{\rm YM}).
\end{equation}
Sobre la $*$-álgebra polinómica generada por holonomías, plaquetas, clover y
$\Phi(h)$ se verifican, generador a generador y después por linealidad,
\begin{equation}\label{eq:pdf2-ym-full-publication-products}
 \widetilde{\mathcal U}_g(AB)
 =\widetilde{\mathcal U}_g(A)\widetilde{\mathcal U}_g(B),
 \quad
 \widetilde{\mathcal U}_g(A^*)=\widetilde{\mathcal U}_g(A)^*,
 \quad
 \omega_g(\widetilde{\mathcal U}_gA)=\omega_0(A).
\end{equation}
Las igualdades de superficie, clover y Bianchi son preservadas por
\eqref{eq:pdf2-ym-transport-surface-bianchi}; por densidad monótona, la extensión
vale en toda el álgebra local acotada. Así
\eqref{eq:pdf2-ym-owner-smeared-cross-gram} es la restricción de un isomorfismo
probabilístico y algebraico completo, no un sustituto de él.

\subsection{Núcleo denso, continuidad fuerte y red local}

Sea $M_{\Phi(h)}$ el operador de multiplicación afiliado al álgebra local y
\begin{equation}\label{eq:pdf2-ym-owner-curvature-weyl}
 \mathsf W(h,t):=\exp\!\bigl(itM_{\Phi(h)}\bigr),
 \qquad h\in C_c^\infty(\mathbb R^4),\quad t\in\mathbb R.
\end{equation}
Se completa $\mathfrak A_{\rm YM}(O)$ con estos unitarios para
$\operatorname{supp}h\subset O$ y con las funciones de clase acotadas de las
holonomías de plaqueta. Denótese por $\mathfrak A_{\rm YM}(O)^{\rm sm}$ la
$*$-álgebra así generada con funciones de test suaves. El núcleo
\begin{equation}\label{eq:pdf2-ym-owner-smooth-core}
 \mathscr D_{\rm sm}:=\operatorname{span}\left\{
 A_1\cdots A_n\Omega_\infty:
 A_j\in\mathfrak A_{\rm YM}(O_j)^{\rm sm},\quad
 O_j\Subset\mathbb R^4\right\}
\end{equation}
es denso por la definición cíclica de
\eqref{eq:pdf2-ym-owner-curvature-cyclic-sector}. Las ecuaciones
\eqref{eq:pdf2-ym-owner-block-e4-naturality} y
\eqref{eq:pdf2-ym-owner-publication-into-physical} dan, para $a\in E(4)$,
\begin{equation}\label{eq:pdf2-ym-owner-field-e4-covariance}
 U(a)\Phi(h)=\Phi(a\!\cdot h),
 \qquad
 \|(U(a)-I)\Phi(h)\|=\|a\!\cdot h-h\|_{L^2}.
\end{equation}
Además, $|e^{ix}-e^{iy}|\le|x-y|$ y
\eqref{eq:pdf2-ym-owner-smeared-gram} implican
\begin{equation}\label{eq:pdf2-ym-owner-weyl-continuity}
 \|\bigl(\mathsf W(a\!\cdot h,t)-\mathsf W(h,t)\bigr)\Omega_\infty\|
 \le |t|\,\|a\!\cdot h-h\|_{L^2}.
\end{equation}
La continuidad se comprueba también para los generadores de holonomía del core;
\eqref{eq:pdf2-ym-owner-weyl-continuity} por sí sola no los incluye. Para una ruta poligonal orientada $\gamma$,
una plaqueta $p$ y $f\in C^1(G)^{\rm class}$ se ponen
\begin{equation}\label{eq:pdf2-ym-owner-holonomy-core-generators}
 A_{\gamma,f}:=f(\operatorname{Hol}_\gamma),
 \qquad A_{p,f}:=f(\operatorname{Hol}_p).
\end{equation}
En un overlay de $\gamma$ y $a\gamma$, con $a\in E(4)$, la ley de puente
\eqref{eq:pdf2-ym-owner-heat-bridge} cancela los incrementos comunes. Para cada
incremento restante $z$ se usa
$|f(xz)-f(x)|\le\|\nabla f\|_\infty d_G(z,e)$ y la cota de segundo momento del
kernel de calor
\(
 \int_Gd_G(z,e)^2k_\tau(z)\,dz\le C_G\tau
\).
Para la ley reponderada, $0<e^{-p/(4g^2)}\le1$ y la continuidad de $p$ en el
factor compacto dan
\begin{equation}\label{eq:pdf2-ym-owner-gibbs-heat-moment}
 \frac{\int_Gd_G(z,e)^2e^{-p(z)/(4g^2)}k_\tau(z)\,dz}
      {\int_Ge^{-p(z)/(4g^2)}k_\tau(z)\,dz}
 \le C_{G,g}\tau,
\end{equation}
uniformemente en las pantallas de un compacto y para $0<\tau\le\tau_0$.
Sumando los tiempos de los incrementos de la diferencia simétrica se obtiene la
estimación calculada
\begin{equation}\label{eq:pdf2-ym-owner-holonomy-e4-estimate}
 \begin{aligned}
 \|U(a)A_{\gamma,f}\Omega_\infty-A_{\gamma,f}\Omega_\infty\|_2^2
 &\le C_{G,g}\|\nabla f\|_\infty^2
     \tau(\gamma\triangle a\gamma),\\
 \|U(a)A_{p,f}\Omega_\infty-A_{p,f}\Omega_\infty\|_2^2
 &\le C_{G,g}\|\nabla f\|_\infty^2
     \tau(\partial p\triangle\partial(ap)).
 \end{aligned}
\end{equation}
Aquí la diferencia simétrica no se mide como dos curvas disjuntas: los conectores
de superficie la llenan por la cinta orientada mínima
$\Sigma(\gamma,a\gamma)$ y
\begin{equation}\label{eq:pdf2-ym-owner-bridge-time-ribbon}
 \tau(\gamma\triangle a\gamma)
 :=\sum_{p\subset\Sigma(\gamma,a\gamma)}\tau_p,
 \qquad
 \tau(\gamma\triangle a\gamma)
 \le C_\gamma d_{E(4)}(a,e),
\end{equation}
con la misma definición para las fronteras de plaqueta. La desigualdad sale de
sumar los tiempos aditivos de \eqref{eq:pdf2-ym-owner-heat-bridge} sobre una cinta
de anchura $d_{E(4)}(a,e)$ y longitud acotada por la de $\gamma$.
La medida de tiempo de puente del miembro derecho tiende así a cero cuando $a\to e$;
esto es exactamente la continuidad de los conectores de ruta de la imagen gauge.
Cada lector clover acotado $\chi(F_{\rm cl})$, $\chi\in C_b^1$, es función de una
suma finita de plaquetas transportadas; la misma cota, multiplicada por
$\|\chi'\|_\infty$ y por el número de términos, prueba también su continuidad.

Para funciones de clase meramente acotadas se hace visible la regularización que
entra en el core. Si $d b$ es Haar izquierdo en $E(4)$ y
$\eta\in C_c^\infty(E(4))$, se define el integral de Bochner
\begin{equation}\label{eq:pdf2-ym-owner-orbit-smearing}
 A[\eta]:=\int_{E(4)}\eta(b)U(b)AU(b)^*\,db,
 \qquad A\in\{A_{\gamma,f},A_{p,f},\chi(F_{{\rm cl}})\},
 \quad \chi\in C_b^1.
\end{equation}
Con $(L_a\eta)(b):=\eta(a^{-1}b)$, el cambio de variable $b\mapsto ab$ da
\begin{equation}\label{eq:pdf2-ym-owner-orbit-smearing-continuity}
 U(a)A[\eta]U(a)^*=A[L_a\eta],
 \qquad
 \|(A[L_a\eta]-A[\eta])\Omega_\infty\|
 \le\|A\|\,\|L_a\eta-\eta\|_{L^1(E(4))}\longrightarrow0.
\end{equation}
Se entiende desde ahora que los generadores de holonomía, plaqueta y clover de
$\mathfrak A_{\rm YM}(O)^{\rm sm}$ son los de
\eqref{eq:pdf2-ym-owner-holonomy-core-generators} con $f\in C^1$, o sus
regularizaciones \eqref{eq:pdf2-ym-owner-orbit-smearing}, y que el soporte barrido
permanece compacto dentro de $O$. Las aproximaciones de identidad en $E(4)$,
junto con \eqref{eq:pdf2-ym-owner-holonomy-e4-estimate}, recuperan los generadores
no regularizados en $L^2$; por ello no se reduce el sector cíclico.

Un telescopio de productos que usa
\eqref{eq:pdf2-ym-owner-weyl-continuity} y
\eqref{eq:pdf2-ym-owner-orbit-smearing-continuity} prueba la continuidad de todos
los generadores de $\mathscr D_{\rm sm}$; por
densidad y unitariedad,
\begin{equation}\label{eq:pdf2-ym-owner-strong-e4}
 \lim_{a\to e}\|(U(a)-I)\Psi\|=0
 \qquad(\Psi\in\mathscr H_{\rm YM}^{\rm curv}).
\end{equation}
Los soportes disjuntos usan factores disjuntos y conmutan. Queda así la red
\begin{equation}\label{eq:pdf2-ym-owner-local-net}
\begin{aligned}
 O_1\subseteq O_2&\Rightarrow
 \mathfrak A_{\rm YM}(O_1)\subseteq\mathfrak A_{\rm YM}(O_2),\\
 [\mathfrak A_{\rm YM}(O_1),\mathfrak A_{\rm YM}(O_2)]&=0
 \quad (O_1\cap O_2=\varnothing),\\
 U(a)\mathfrak A_{\rm YM}(O)U(a)^*&=\mathfrak A_{\rm YM}(aO).
\end{aligned}
\end{equation}
La identidad de Gram para $h\ne0$ y
\eqref{eq:pdf2-ym-owner-wc-curvature-functional} prueban que la red no es escalar.

\subsection{Paquete OS continuo y reconstrucción}

Para una dirección $\nu$, sea $\mathfrak A_{+,\nu}$ la clausura de la unión de
las álgebras anteriores con soporte compacto en $\{x_\nu>0\}$. El estado
estático, la ley de pilas y el operador de transferencia no se definen por
separado. En una pantalla $\alpha$, sea
$\widehat\mu_{\alpha,g}^{\rm YM}$ la marginal de
\eqref{eq:pdf2-ym-gibbs-measure}, sea
$\nu_{\alpha,g}:=(\pi_\alpha^{\rm phys})_*
\widehat\mu_{\alpha,g}^{\rm YM}$ su ley en el cociente y sea
$K_{\alpha,g}(t)=e^{-tD_{\alpha,g}^{\rm YM}}$ su kernel.
Para tiempos ordenados $t_1\le\cdots\le t_n$, la marginal euclídea de la
misma medida es
\begin{equation}\label{eq:pdf2-ym-owner-euclidean-time-marginal}
 \mathbb P_{\alpha,g}^{E,(\nu;t_1,\ldots,t_n)}
 :=(\operatorname{ev}_{t_1}^{(\nu)},\ldots,
    \operatorname{ev}_{t_n}^{(\nu)})_*
    \widehat\mu_{\alpha,g}^{\rm YM}.
\end{equation}
La identidad de desintegración que forma parte de
\eqref{eq:pdf2-g-gau-transfer-euclidean-dato} la calcula, no la reemplaza por
una cadena auxiliar:
\begin{equation}\label{eq:pdf2-ym-owner-euclidean-markov-identification}
 d\mathbb P_{\alpha,g}^{E,(\nu;t_1,\ldots,t_n)}
 =\nu_{\alpha,g}(dy_1)
   \prod_{j=1}^{n-1}K_{\alpha,g}
      (t_{j+1}-t_j;y_j,dy_{j+1}).
\end{equation}
Si $M_A$ denota multiplicación por un lector acotado $A$, sus correlaciones
son por tanto
\begin{equation}\label{eq:pdf2-ym-owner-schwinger-transfer}
 \begin{aligned}
 S_{\alpha,n}^{(\nu,g)}(A_1,t_1;\ldots;A_n,t_n)
 &:=\int\prod_{j=1}^n(\tau_{t_je_\nu}A_j)(x)\,
       d\widehat\mu_{\alpha,g}^{\rm YM}(x)\\
 &=\langle\mathbf1,
 M_{A_1}K_{\alpha,g}(t_2-t_1)M_{A_2}\cdots
 K_{\alpha,g}(t_n-t_{n-1})M_{A_n}\mathbf1
 \rangle_{L^2(\nu_{\alpha,g})}.
 \end{aligned}
\end{equation}
La ley de la pila estacionaria en esos mismos tiempos es esa marginal
euclídea, escrita mediante su desintegración:
\begin{equation}\label{eq:pdf2-ym-owner-markov-stack-law}
 d\mathbb P_{\alpha,g}^{(\nu;t_1,\ldots,t_n)}
 :=d\mathbb P_{\alpha,g}^{E,(\nu;t_1,\ldots,t_n)}
 =\nu_{\alpha,g}(dy_1)
   \prod_{j=1}^{n-1}K_{\alpha,g}(t_{j+1}-t_j;y_j,dy_{j+1}).
\end{equation}
Integrar sucesivamente $y_n,y_{n-1},\ldots,y_1$ da la identidad exacta
\begin{equation}\label{eq:pdf2-ym-owner-stack-transfer-identity}
 S_{\alpha,n}^{(\nu,g)}(A_1,t_1;\ldots;A_n,t_n)
 =\int\prod_{j=1}^nA_j(y_j)\,
   d\mathbb P_{\alpha,g}^{(\nu;t_1,\ldots,t_n)}(y).
\end{equation}
Si todos los tiempos coinciden, los kernels son la identidad y
\begin{equation}\label{eq:pdf2-ym-owner-static-is-stack}
 S_{\alpha,n}^{(\nu,g)}(A_1,t;\ldots;A_n,t)
 =\int A_1\cdots A_n\,d\nu_{\alpha,g}.
\end{equation}
Así el estado Schwinger estático es la marginal de tiempo igual de la misma
medida euclídea, y su desintegración es la ley Markov; no son dos funcionales
que luego se identifican por nombre. La
compatibilidad de $\widehat J_{\alpha,g}$ con medida y kernel hace compatible
\eqref{eq:pdf2-ym-owner-stack-transfer-identity} con todo overlay. Su límite define
\begin{equation}\label{eq:pdf2-ym-owner-schwinger-functionals}
 S_n^{\rm YM}(A_1,t_1;\ldots;A_n,t_n)
 :=\lim_{\alpha}S_{\alpha,n}^{(\nu,g)}
 (A_{1,\alpha},t_1;\ldots;A_{n,\alpha},t_n).
\end{equation}

Las inserciones de $\Phi$ se obtienen derivando en $t=0$ los unitarios
\eqref{eq:pdf2-ym-owner-curvature-weyl}. Puesto que los $Z_{\alpha,c}$ transportados
son centrados, independientes y uniformemente acotados, el lema de Hoeffding da,
nivel a nivel,
\begin{equation}\label{eq:pdf2-ym-owner-subgaussian-bound}
 \left\langle\Omega_\alpha,
 e^{t\Phi_\alpha(h)}\Omega_\alpha\right\rangle
 \le \exp\!\left(\frac9{16}t^2\|P_\alpha h\|_{L^2}^2\right).
\end{equation}
La cota pasa al límite por
\eqref{eq:pdf2-ym-owner-smeared-cauchy-derived}; todas las derivadas son continuas en
la topología de Schwartz y definen distribuciones temperadas.

La forma de ocho caras se prolonga primero a toda pila finita. Si $F$ depende de
las $N$ capas positivas de una pantalla $\alpha$, la reflexión de
\eqref{eq:pdf2-ym-owner-markov-stack-law} y la condición en la capa central dan
\begin{equation}\label{eq:pdf2-ym-owner-finite-stack-os}
 \begin{aligned}
 &\int \overline{F(\Theta_\nu y_{-N},\ldots,
                         \Theta_\nu y_{-1})}
          F(y_1,\ldots,y_N)\,d\Omega_{\alpha,N}^{(\nu,g)}(y)\\
 &\qquad=
 \int\nu_{\alpha,g}(dz)\left|
  \int K_{\alpha,g}(a_\alpha/2;z,dy_1)
       \prod_{j=1}^{N-1}K_{\alpha,g}(a_\alpha;y_j,dy_{j+1})
       F(y_1,\ldots,y_N)
 \right|^2\ge0.
 \end{aligned}
\end{equation}
Todo cilindro de semiespacio pertenece a una pila finita; por ello
\eqref{eq:pdf2-ym-owner-finite-stack-os} materializa la forma OS antes de la
clausura. En el cociente OS, si $F$ tiene inserciones en tiempos positivos, sea
$(\tau_sF)(A_1,t_1;\ldots;A_n,t_n):=
F(A_1,t_1+s;\ldots;A_n,t_n+s)$. Se define el operador de traslado por
\begin{equation}\label{eq:pdf2-ym-owner-os-transfer-operator}
 T_{\rm OS}^{(\nu)}(s)[F]:=[\tau_sF],
 \qquad s\ge0.
\end{equation}
Para cilindros $F,G$ de semiespacio, la igualdad euclídea--Markov
\eqref{eq:pdf2-ym-owner-euclidean-markov-identification} y la propiedad de
semigrupo dan generador a generador
\begin{equation}\label{eq:pdf2-ym-owner-os-markov-intertwining}
 \begin{aligned}
 \langle[F],T_{\rm OS}^{(\nu)}(s)[G]\rangle_{\rm OS}
 &=S_g^{\rm YM}((\Theta_\nu F)^*\tau_sG)\\
 &=\left\langle
    \mathcal V_{\infty,\nu}[F],
    e^{-sD_{\rm curv,g}^{\rm YM}}
    \mathcal V_{\infty,\nu}[G]
   \right\rangle .
 \end{aligned}
\end{equation}
La identidad muestra a la vez que las clases nulas permanecen nulas y que,
por densidad de los cilindros,
\begin{equation}\label{eq:pdf2-ym-owner-os-generator-identification}
 \mathcal V_{\infty,\nu}T_{\rm OS}^{(\nu)}(s)
 \mathcal V_{\infty,\nu}^{-1}
 =e^{-sD_{\rm curv,g}^{\rm YM}}.
\end{equation}
Ésta es la flecha explícita medida euclídea $\to$ desintegración $\to$
transferencia OS $\to$ Hamiltoniano.

\begin{proposition}[Paquete Osterwalder--Schrader continuo]
\label{prop:pdf2-ym-owner-continuum-os-package}
Los funcionales \eqref{eq:pdf2-ym-owner-schwinger-functionals} satisfacen:
\begin{enumerate}
 \item normalización, hermiticidad y simetría bosónica;
 \item covariancia euclídea y regularidad temperada;
 \item positividad por reflexión en cada una de las cuatro direcciones,
 \begin{equation}\label{eq:pdf2-ym-owner-continuum-reflection-positive}
  \sum_{i,j}\overline{a_i}a_j\,
  S^{\rm YM}\!\left((\Theta_\nu F_i)^*F_j\right)\ge0,
  \qquad F_i\in\mathfrak A_{+,\nu};
 \end{equation}
 \item localidad y cluster exponencial: si $A,B$ son locales centrados y una
 traslación de longitud $r$ separa sus soportes, entonces
 \begin{equation}\label{eq:pdf2-ym-owner-continuum-cluster}
  |S^{\rm YM}(A\,\tau_rB)|
  \le e^{-3\kappa_{\rm av}r}\,
       \|A\Omega_\infty\|\,\|B\Omega_\infty\|.
 \end{equation}
\end{enumerate}
La reconstrucción OS de cualquiera de las cuatro semirredes produce el mismo
espacio cíclico, el mismo vacío y el Hamiltoniano
$D_{\rm curv,g}^{\rm YM}$. La continuación OS produce campos de Wightman locales y una
representación de Poincaré cuyo espectro satisface
\begin{equation}\label{eq:pdf2-ym-owner-wightman-spectrum}
 \operatorname{Spec}(P^0,\mathbf P)\subseteq\overline V_+,
 \qquad P^0=D_{\rm curv,g}^{\rm YM}\ge0.
\end{equation}
\end{proposition}

\begin{proof}
Normalización y hermiticidad proceden del estado; simetría y localidad, de la
conmutatividad euclídea a soportes disjuntos. La covariancia y continuidad son
\eqref{eq:pdf2-ym-owner-field-e4-covariance}--
\eqref{eq:pdf2-ym-owner-strong-e4}; la regularidad es
\eqref{eq:pdf2-ym-owner-subgaussian-bound}. Para
$F_i$ cilíndricos, todos viven en un overlay y una pila finitos;
\eqref{eq:pdf2-ym-owner-finite-stack-os} y la polarización prueban
\eqref{eq:pdf2-ym-owner-continuum-reflection-positive}. La aproximación en norma OS
y Cauchy--Schwarz la extienden a $\mathfrak A_{+,\nu}$.

Tras una rotación euclídea, la separación espacial se convierte en tiempo positivo.
La identidad pila--transferencia
\eqref{eq:pdf2-ym-owner-stack-transfer-identity} y el operador
\eqref{eq:pdf2-ym-owner-os-transfer-operator} escriben la correlación centrada como
\[
 \langle A^*\Omega_\infty,
 e^{-rD_{\rm curv,g}^{\rm YM}}B\Omega_\infty\rangle.
\]
La cota espectral \eqref{eq:pdf2-ym-owner-curvature-sector-gap} da
\eqref{eq:pdf2-ym-owner-continuum-cluster}. Finalmente, los conectores
\eqref{eq:pdf2-ym-owner-os-connectors} forman un sistema dirigido; su unión es densa
por \eqref{eq:pdf2-ym-owner-smooth-core}. La identidad
\eqref{eq:pdf2-ym-owner-common-os-hamiltonian} identifica las cuatro clausuras y sus
generadores con $D_{\rm curv,g}^{\rm YM}$. Las hipótesis ya verificadas de regularidad,
covariancia, reflexión y cluster son las entradas de la continuación OS; su teorema de
reconstrucción da localidad de Wightman y
\eqref{eq:pdf2-ym-owner-wightman-spectrum}.
\end{proof}

\subsection{Publicación tipada como teoría de Yang--Mills}

La identificación no se apoya sólo en evaluar un clover sobre una conexión dada.
Sea $\mathcal R_G^{\rm loc}$ la $*$-álgebra de lectores formada por todas las
funciones de clase de las holonomías de incidencia y plaqueta, sus transportes a un
vértice base y las polarizaciones del lector $\mathcal P^{F^2}$. Sea
\begin{equation}\label{eq:pdf2-ym-owner-curvature-fibre}
 \mathcal E_G:=\Lambda^2(\mathbb R^4)^*\otimes\mathfrak g
\end{equation}
con el producto invariante. En cada celda, la polarización de
$\mathcal P^{F^2}$ da un Gram positivo $G_{\alpha,c}^{F}$ sobre
$\mathcal E_G$. Se divide por su kernel y se aplica su raíz positiva; el lector
normalizado resultante se denota $Y_{\alpha,c}(v)$ y verifica
\begin{equation}\label{eq:pdf2-ym-owner-curvature-cell-gram}
 \langle Y_{\alpha,c}(v),Y_{\alpha,d}(w)\rangle
 =\delta_{cd}\langle v,w\rangle_{\mathcal E_G}.
\end{equation}
Si $d\subset c$, sea $T_{d\leftarrow c}$ el transporte de marco ya presente en la
holonomía superficial. El entrelazador de curvatura completa es
\begin{equation}\label{eq:pdf2-ym-owner-curvature-intertwiner}
 I_{\alpha\beta}^{F}Y_{\alpha,c}(v)
 :=\sum_{d\subset c}\sqrt{\frac{|d|}{|c|}}\,
 Y_{\beta,d}\!\left(\operatorname{Ad}_{T_{d\leftarrow c}}v\right).
\end{equation}
La invariancia del producto de $\mathfrak g$, la disjunción de subceldas y la
composición de transportes prueban
\begin{equation}\label{eq:pdf2-ym-owner-curvature-intertwiner-identities}
 (I_{\alpha\beta}^{F})^*I_{\alpha\beta}^{F}=I,
 \qquad
 I_{\beta\gamma}^{F}I_{\alpha\beta}^{F}=I_{\alpha\gamma}^{F}.
\end{equation}
Para $h\in C_c^\infty(\mathbb R^4;\mathcal E_G)$ se pone
\begin{equation}\label{eq:pdf2-ym-owner-f-alpha}
 \mathbf F_\alpha(h):=
 \sum_{c\in\mathcal C_\alpha^{\rm cell}}
 Y_{\alpha,c}\!\left(
   |c|^{-1/2}\int_c h(x)\,d^4x\right).
\end{equation}
En un overlay común, el mismo cálculo de
\eqref{eq:pdf2-ym-owner-smeared-cross-gram} da ahora
\begin{equation}\label{eq:pdf2-ym-owner-curvature-cross-gram}
 \left\langle I_{\alpha\gamma}^{F}\mathbf F_\alpha(h),
 I_{\beta\gamma}^{F}\mathbf F_\beta(k)\right\rangle
 =\langle P_\alpha h,P_\beta k\rangle_{L^2(\mathcal E_G)}.
\end{equation}
Por tanto $\mathbf F_\alpha(h)$ es Cauchy por convergencia fuerte de las proyecciones,
y no por una hipótesis nueva. El mismo emparejamiento de Haar e innovaciones de
\eqref{eq:pdf2-ym-owner-publication-into-physical}, ahora aplicado a la base
ortonormalizada de $G_{\alpha,c}^{F}$, realiza el colímite mediante una isometría
local y gauge-covariante en
$L^2(\widehat X_\infty,\widehat\mu_{\infty,g}^{\rm YM})$. Su límite define la
distribución aleatoria
gauge-covariante
\begin{equation}\label{eq:pdf2-ym-owner-distributional-curvature}
\begin{aligned}
 \mathbf F_{\mu\nu}&\in
 \mathcal S'(\mathbb R^4;\mathfrak g)\ \widehat\otimes\
 L^2(\widehat X_\infty,\widehat\mu_{\infty,g}^{\rm YM}),\\
 U(a)\mathbf F(h)U(a)^*&=\mathbf F(a\!\cdot h),\\
 u\mathbf F(h)u^{-1}&=\mathbf F(\operatorname{Ad}_u h),
 \qquad a\in E(4),\ u\in G.
\end{aligned}
\end{equation}
Su canal de clase espectral de orden tres es precisamente $\Phi$. La relación
superficial finita
\begin{equation}\label{eq:pdf2-ym-owner-surface-product}
 \operatorname{Hol}^{(\alpha)}_p
 =\prod_{p'\subset p}^{\longrightarrow}
  T_{p'}\operatorname{Hol}^{(\beta)}_{p'}T_{p'}^{-1}
\end{equation}
pasa a la clausura cofinal porque todos sus productos cruzados se calculan en un
overlay. En particular, el lector clover aleatorio, y no sólo su evaluación en una
conexión suave, satisface la identidad de pequeños lazos
\begin{equation}\label{eq:pdf2-ym-owner-holonomy-curvature-relation}
 \mathbf F(h)=L^2\!\mathord{-}\!\lim_{\alpha}
 a_\alpha^4\sum_x
 \left\langle(P_\alpha h)(x),F_{{\rm cl},\alpha}(x)\right\rangle,
 \qquad
 \delta_{\rm gauge}S_n^{\rm YM}=0,
 \qquad
 \operatorname{Alt}_{\lambda\mu\nu}D_\lambda\mathbf F_{\mu\nu}=0.
\end{equation}
La primera igualdad es \eqref{eq:pdf2-ym-owner-curvature-cross-gram} escrita en el
representante clover; la segunda es la identidad de Ward
\eqref{eq:pdf2-ym-ward-finite} pasada al límite, y la
tercera es la cancelación de la frontera de una frontera en el producto superficial.
Ninguna usa el término de masa.

Para evitar definir el blanco por la conclusión que se quiere obtener, sea
$\mathbf{YM}_4(G,g)$ la clase de quíntuplas locales
$(\mathfrak A,\mu_g,K_g,\operatorname{Hol},\mathbf F)$ que verifican las siguientes
ecuaciones comprobables: (i) la densidad de Gibbs es
\eqref{eq:pdf2-ym-gibbs-measure}; (ii) $K_g$ es Markov, reversible y satisface
\eqref{eq:pdf2-ym-schwinger-dyson-finite}--\eqref{eq:pdf2-ym-ward-finite};
(iii) holonomía, clover y curvatura satisfacen
\eqref{eq:pdf2-ym-transport-surface-bianchi} y
\eqref{eq:pdf2-ym-owner-holonomy-curvature-relation}; y (iv) sus funcionales son la
misma ley de pilas de \eqref{eq:pdf2-ym-owner-stack-transfer-identity} y cumplen el
paquete OS\@. El objeto construido es, componente a componente,
\begin{equation}\label{eq:pdf2-ym-owner-qym-object}
\begin{aligned}
 \mathfrak Q_{\rm YM}^{(4)}(G,g):={}&
 \bigl(\{\mathfrak A_{\rm YM}(O)\}_O,
 S_g^{\rm YM},U,\Theta,\Omega_{\infty,g},\\
 &\widehat\mu_{\infty,g}^{\rm YM},
 \{S_{m,g},\mathscr S_{m,g}\}_m,K_g,
 \operatorname{Hol},\mathbf F,\\
 &\mathscr H_{\rm OS},D_{\rm curv,g}^{\rm YM}\bigr)
 \in\mathbf{YM}_4(G,g).
\end{aligned}
\end{equation}
El mapa de publicación no omite ninguna de esas componentes:
\begin{equation}\label{eq:pdf2-ym-owner-typed-publication-map}
 \begin{aligned}
 \mathcal P_{{\rm YM},g}:\mathfrak G_{\rm gau}(G)&\longrightarrow
 \mathbf{YM}_4(G,g),\\
 (\widehat\mu^0,\mathbb P^{\rm phys},D^0,\Theta,W,
   \operatorname{Hol},\mathcal P^{F^2})
 &\longmapsto
 (\widehat\mu_g^{\rm YM},S_g,\mathscr S_g,K_g,\mathbb P_g^{\rm phys},
  S_g^{\rm YM},\mathfrak A_{\rm YM},D_{\rm curv,g}^{\rm YM},
  \mathbf F,\operatorname{Hol}).
 \end{aligned}
\end{equation}
Los $S_n^{\rm YM}$ son así exactamente las funciones de Schwinger de la medida
dependiente de $g$, la ley Markov y el operador de transferencia ya construidos.

Queda por comprobar que $S_{m,g}$ es la acción de Yang--Mills y no sólo una
función con ese nombre. Sea
$A\in C_c^3(\mathbb R^4;T^*\mathbb R^4\otimes\mathfrak g)$ una conexión suave en
una carta donde el logaritmo de los lazos pequeños es único. Para la plaqueta
$p_{\mu\nu}(x)$, la expansión ordenada de Magnus da
\begin{equation}\label{eq:pdf2-ym-owner-plaquette-magnus}
 \log\operatorname{Hol}_{p_{\mu\nu}(x)}(A)
 =a_m^2F_{A,\mu\nu}(x)
  +\frac{a_m^3}{2}(D_\mu+D_\nu)F_{A,\mu\nu}(x)
  +R_{m,\mu\nu}(x),
\end{equation}
con
$\|R_{m,\mu\nu}(x)\|\le C_G a_m^4
(\|D^2F_A\|_{\infty,p}+\|F_A\|_{\infty,p}^2)$.
El promedio de las cuatro orientaciones cancela el término impar y define
\begin{equation}\label{eq:pdf2-ym-owner-clover}
 F_{{\rm cl},m}(x)
 :=\frac1{4a_m^2}\sum_{p\ni x}
  \operatorname{Ad}_{T_p}\log\operatorname{Hol}_p,
 \qquad
 \|F_{{\rm cl},m}(x)-F_A(x)\|
 \le C_G a_m^2(\|D^2F_A\|_\infty+\|F_A\|_\infty^2).
\end{equation}
La definición de la raíz
\eqref{eq:pdf2-ym-action-root} es precisamente la polarización de esos seis
componentes. Sustituyendo \eqref{eq:pdf2-ym-owner-plaquette-magnus} en cada
generador de esa polarización se calcula
\begin{equation}\label{eq:pdf2-ym-owner-action-root-smooth-chart}
 \xi_{m,b}(A)
 =a_m^2\operatorname{vec}_{\mu<\nu}F_{A,\mu\nu}(x_b)
  +r_{m,b}(A),
 \qquad
 \|r_{m,b}(A)\|\le C_{G,A}a_m^4.
\end{equation}
Los términos cruzados impares se cancelan por las cuatro orientaciones del clover;
los pares están incluidos una sola vez por
\eqref{eq:pdf2-ym-action-factor-partition}. Por suma de Riemann y
Cauchy--Schwarz,
\begin{equation}\label{eq:pdf2-ym-owner-action-error}
 \left|S_{m,g}(A)-\frac1{4g^2}
  a_m^4\sum_x\|F_{{\rm cl},m}(x)\|^2\right|
 \le C_{G,A,g}a_m^2,
\end{equation}
y, usando la cota de \eqref{eq:pdf2-ym-owner-clover},
\begin{equation}\label{eq:pdf2-ym-owner-f2}
 S_{m,g}(A)\longrightarrow
 S_{{\rm YM},g}(A):=\frac1{4g^2}
 \int_{\mathbb R^4}\|F_A(x)\|^2\,d^4x.
\end{equation}
Las ecuaciones de Schwinger--Dyson y Ward ya se calcularon sobre la medida que usa
esta acción. La brecha, en cambio, se deduce después por la conjugación unitaria de
\eqref{eq:pdf2-ym-g-dynamics} y el testigo
\eqref{eq:pdf2-ym-g-gap-before-limit}; no fue usada para escoger $g$, $S_{m,g}$ ni
$\widehat\mu_{m,g}^{\rm YM}$. Éste es el enlace no circular entre acción, ley y gap.

\begin{theorem}[Cierre Yang--Mills para todo grupo de Lie compacto, conectado y simple]
\label{thm:pdf2-ym-cierre-global}
La estructura completa $\mathfrak G_{\rm gau}$ construye una familia proyectiva
local para todo grupo de Lie compacto, conectado y simple. El mapa
\eqref{eq:pdf2-ym-owner-typed-publication-map} publica una teoría cuántica de
Yang--Mills no trivial sobre $\mathbb R^4$: sus funcionales satisfacen el paquete OS
continuo y su continuación satisface los axiomas locales de Wightman por la
proposición~\ref{prop:pdf2-ym-owner-continuum-os-package}; sus
holonomías y su curvatura
distribucional satisfacen
\eqref{eq:pdf2-ym-owner-holonomy-curvature-relation}, y las cuatro reconstrucciones
coinducen el mismo Hamiltoniano $D_{\infty,g}^{\rm YM}$, con vacío simple. Su ley es
$\widehat\mu_{\infty,g}^{\rm YM}$, construida por la familia compatible de acciones
locales $\{S_{m,g},\mathscr S_{m,g}\}_m$, y
satisface las ecuaciones Schwinger--Dyson y Ward
\eqref{eq:pdf2-ym-schwinger-dyson-finite}--\eqref{eq:pdf2-ym-ward-finite}. El sector
cíclico de curvatura
\eqref{eq:pdf2-ym-owner-curvature-cyclic-sector} es reductor, contiene el testigo
\eqref{eq:pdf2-ym-owner-wc-curvature-functional} y posee
\begin{equation}\label{eq:pdf2-ym-final-gap}
 \boxed{\Delta_{\rm YM}^{\rm HMT}=3\kappa_{\rm av}>0},
 \qquad
 m_{\rm gap}^{\rm HMT}
 =\frac{\hbar_{\rm int}}{c_{\rm int}}\Delta_{\rm YM}^{\rm HMT}>0.
\end{equation}
La red local, la acción fuerte de $E(4)$, las distribuciones de Schwinger, la
holonomía y $\mathbf F$ proceden del mismo estado completo. El lector exacto
$\mathcal P^{F^2}$ define la densidad de la ley; el clover es su carta suave y
\eqref{eq:pdf2-ym-owner-f2} identifica su acción clásica. La brecha pertenece al sector colectivo
gauge-invariante de curvatura y no asigna una masa elemental al gluón.
\end{theorem}

\begin{proof}
La proyectividad cinemática es \eqref{eq:pdf2-ym-owner-projectivity}; la raíz
\eqref{eq:pdf2-ym-action-root}, el cociclo
\eqref{eq:pdf2-ym-action-cocycle} y la densidad
\eqref{eq:pdf2-ym-gibbs-measure} construyen primero la acción y la medida a
acoplamiento $g$. El transporte factor a factor
\eqref{eq:pdf2-ym-gibbs-transport-factor}--
\eqref{eq:pdf2-ym-full-algebra-identities} prueba proyectividad, preservación del
estado, productos y $*$; 
\eqref{eq:pdf2-ym-transport-surface-bianchi} descarga clover y Bianchi. El índice exhaustivo
\eqref{eq:pdf2-ym-owner-exhaustive-index}, la forma
\eqref{eq:pdf2-ym-owner-full-form} y el circuito
\eqref{eq:pdf2-ym-owner-local-circuit} construyen el generador común. La reducción
física y su kernel cociente son
\eqref{eq:pdf2-ym-owner-physical-reduction}--
\eqref{eq:pdf2-ym-owner-D}. La conjugación
\eqref{eq:pdf2-ym-g-dynamics} los transporta a la medida dependiente de $g$ sin
alterar su espectro. La descomposición de Hoeffding prueba el vacío y la cota, y
\eqref{eq:pdf2-ym-owner-wc-invariance}--
\eqref{eq:pdf2-ym-owner-wc-eigenvalue} exhiben un vector físico que alcanza el
umbral. El entrelazamiento nonádico transporta forma, vacío y testigo al límite, donde
\eqref{eq:pdf2-ym-owner-limit-gap}--
\eqref{eq:pdf2-ym-owner-exact-gap} fijan la brecha exacta. Las cuatro
factorizaciones \eqref{eq:pdf2-ym-owner-four-os}, la inclusión de rangos
\eqref{eq:pdf2-ym-owner-os-range-inclusion} y los conectores
\eqref{eq:pdf2-ym-owner-os-connectors} identifican sus reconstrucciones con ese
mismo generador. La holonomía
\eqref{eq:pdf2-ym-owner-incidence-holonomy} coloca materialmente $W_c$ en el sector
de curvatura, donde \eqref{eq:pdf2-ym-owner-curvature-sector-gap} prueba la brecha
exacta. El sistema \eqref{eq:pdf2-ym-owner-block-intertwiner} controla los productos
cruzados mediante \eqref{eq:pdf2-ym-owner-smeared-cross-gram};
\eqref{eq:pdf2-ym-full-publication-restriction}--
\eqref{eq:pdf2-ym-full-publication-products} elevan ese primer caos al isomorfismo
de toda el álgebra. La convergencia fuerte de $P_\alpha$ produce $\Phi$; las cotas
\eqref{eq:pdf2-ym-owner-holonomy-e4-estimate} y
\eqref{eq:pdf2-ym-owner-orbit-smearing-continuity} prueban además la continuidad
$E(4)$ de holonomías, plaquetas y clover. La desintegración euclídea
\eqref{eq:pdf2-ym-owner-euclidean-markov-identification}, la identidad exacta
\eqref{eq:pdf2-ym-owner-stack-transfer-identity} y el entrelazamiento
\eqref{eq:pdf2-ym-owner-os-markov-intertwining} enlazan la misma medida,
estado Schwinger, pila Markov, transferencia OS y generador del gap. La
proposición~\ref{prop:pdf2-ym-owner-continuum-os-package} prolonga las
cuatro formas finitas a las
semirredes completas, prueba regularidad y cluster y reconstruye el Hamiltoniano. Por
último, \eqref{eq:pdf2-ym-owner-distributional-curvature}--
\eqref{eq:pdf2-ym-owner-typed-publication-map} identifican esos mismos funcionales
como los Schwinger de la teoría gauge local; las integraciones por partes
\eqref{eq:pdf2-ym-schwinger-dyson-finite}--\eqref{eq:pdf2-ym-ward-finite} prueban
Schwinger--Dyson y Ward. Finalmente,
\eqref{eq:pdf2-ym-owner-plaquette-magnus}--\eqref{eq:pdf2-ym-owner-f2} identifican
la acción clásica antes de aplicar la equivalencia espectral. La cadena no usa el
gap para escoger la medida.
\end{proof}

\begin{falsifier}[Colapso de la dinámica producto]
Si se reemplaza $D_m^{\rm prod}$ por $D_m^{\rm fib}$, todos los sectores no
constantes reciben el mismo autovalor. Un vector de Hoeffding con $|S|=2$ debería
tener energía $6\kappa_{\rm av}$ por \eqref{eq:pdf2-ym-hoeffding}, pero el
semigrupo simple le asignaría $3\kappa_{\rm av}$. La identificación de ambos
generadores es falsa.
\end{falsifier}

\begin{falsifier}[Confusión entre inclusión genealógica y promedio de bloque]
Si se sustituye $I_{\alpha\beta}^{\rm curv}$ por
$\widehat J_{\alpha\beta}$ en
\eqref{eq:pdf2-ym-owner-smeared-cross-gram}, una coordenada gruesa persiste con
coeficiente uno mientras el coeficiente fino cambia por el cociente de escalas. Para
$a_{m+1}=a_m/9$, el producto cruzado ancestral lleva un factor $1/81$ y no converge
a la norma diagonal. La identidad
\eqref{eq:pdf2-ym-owner-smeared-cauchy-derived} falla; por eso los dos mapas se
mantienen separados.
\end{falsifier}

\begin{falsifier}[Ablación de la fibra de incidencia]
Si se sustituye $\widehat X_m$ por su marginal de enlaces, se borra $q_c$ y con ella
el observable físico $W_c$. El mapa de olvido no refleja la subálgebra gauge completa;
por tanto no puede utilizarse para negar el autovalor del objeto fuente. Objeto
sustituido; conclusión no transferible.
\end{falsifier}

\begin{falsifier}[Sustitución del lector exacto por su carta suave]
Si se usa directamente el término asintótico de
\eqref{eq:pdf2-ym-owner-f2} como densidad en una pantalla finita, se borra el resto
controlado de \eqref{eq:pdf2-ym-owner-action-error} y se pierde el cociclo exacto
\eqref{eq:pdf2-ym-action-cocycle}. La densidad finita gobernante es
$Z_{m,g}^{-1}e^{-\mathcal P_m^{F^2}/(4g^2)}$; el clover y el límite suave la
identifican con la acción clásica, pero no la sustituyen.
\end{falsifier}

\paragraph{Procedencia y estatuto.}
La separación entre fibra simple y generador producto, la medida común, la compresión
gauge, las cuatro OS, la forma límite y el testigo físico son
\texttt{RESULTADO\_RECUPERADO} y \texttt{EXACTO\_INTERNO} en el dominio gauge
regular HMT registrado. La desintegración con tiempos no uniformes, el índice
exhaustivo de expectativas, el kernel cociente, la matriz $B^+$, la acción explícita
en el colímite, el sistema cofinal de bloque, el paquete OS continuo, la curvatura
distribucional y la holonomía de incidencia son
\texttt{FORMALIZACION\_NUEVA}: tipan y descargan los enlaces que antes aparecían
sólo nominalmente, sin introducir la brecha como dato. Su composición completa
establece el teorema precedente.
La raíz positiva del lector $F^2$, su densidad Gibbs, el transporte local compatible,
las ecuaciones Schwinger--Dyson/Ward, el isomorfismo algebraico completo, las cotas
de continuidad de holonomías y la identidad pila--transferencia son
\texttt{FORMALIZACION\_NUEVA}. Su papel es descargar materialmente los cuatro
puentes; la identificación clover--acción clásica es consecuencia de la expansión
de Magnus y no una premisa espectral.

\endgroup
% Cuerpo científico recuperado y distribuido en su residencia terminal.
\section{Teoría gauge euclídea y brecha de Yang--Mills}

La familia proyectiva gauge construida desde el estado completo posee
covariancia local, localidad, positividad por reflexión, una acción
fuertemente continua de \(E(4)\), funcionales de Schwinger compatibles y un
vacío simple. Las cuatro reconstrucciones coinducen el mismo Hamiltoniano y
el sector cíclico de curvatura posee la brecha exacta
\[
  \Delta_{\rm YM}^{\rm HMT}=3\kappa_{\rm av}>0.
\]
El lector de curvatura y su límite clásico proceden de la misma ley. La brecha
pertenece al sector colectivo gauge-invariante y no asigna una masa elemental
al gluón.




\HMTFlushCurrentChapter
\chapter{Realización algebraica de clases de Hodge mediante proyectores cohomológicos compatibles y descenso racional}
\begingroup
\renewcommand{\chapter}[2][]{}%
\chapter{Imagen cohomológica y composición hacia Hodge}

\section{Objeto fuente, observaciones finitas y publicación}

Sea $X$ una variedad proyectiva lisa compleja y $p\ge0$. El objeto producido por
la construcción común y recibido aquí es la estructura cohomológica completa
$\mathfrak G_{\rm coh}(X,p)$, que conserva
cartas basadas, cruces orientados, conos de Mayer--Vietoris, carry,
frontera, ruta, multisección, supervivencia y terminal. En profundidad $N$ se
escriben
\begin{equation}\label{eq:pdf2-hodge-sistemas}
 \mathscr S_N(X,p):=\operatorname{Surv}_{J_N}(X,p),
 \qquad
 \rho_{N+1,N}:\mathscr S_{N+1}(X,p)\to\mathscr S_N(X,p),
\end{equation}
y los módulos de observaciones
\begin{equation}\label{eq:pdf2-hodge-observaciones}
 e_N:\mathscr S_N(X,p)\to\mathscr H_N^p(X),
 \qquad
 q_{N+1,N}:\mathscr H_{N+1}^p(X)\to\mathscr H_N^p(X),
\end{equation}
con el cuadrado
\begin{equation}\label{eq:pdf2-hodge-cuadrado-restriccion}
 e_N\rho_{N+1,N}=q_{N+1,N}e_{N+1}.
\end{equation}
Para
\[
 \alpha\in\operatorname{Hdg}^p(X)_{\mathbb Q}
 :=H^{2p}(X,\mathbb Q)\cap H^{p,p}(X),
\]
sea $q_N\alpha$ su familia de observaciones finitas. La fibra que debe
prolongarse es
\begin{equation}\label{eq:pdf2-hodge-fibra-alpha}
 \mathscr F_N(\alpha):=
 \{s\in\mathscr S_N(X,p):e_N(s)=q_N\alpha\}.
\end{equation}

El realizador y la publicación se tipan por
\begin{equation}\label{eq:pdf2-hodge-realizador}
 R_{\rm Hdg}:X_\chi(X,p)\to
 K_0\!\left(\operatorname{Perf}^{\ge p}(X)\right)_{\mathbb Q},
\end{equation}
\begin{equation}\label{eq:pdf2-hodge-publicacion}
 \operatorname{cl}_{X,p}\circ\operatorname{ch}^{\rm loc}_p
 \circ R_{\rm Hdg}
 =\operatorname{ev}^{\rm Hdg}_{\infty,X,p}.
\end{equation}
Ésta es la identidad de acción del lector ya construida dentro del objeto completo
\eqref{eq:pdf2-g-coh-limite-lector-dato}; no introduce desde el codominio la
preimagen de una clase Hodge. El propietario exacto precedente construye esa
existencia y este capítulo despliega la composición. El lector
debe permanecer en el orden
\[
 \mathfrak G_{\rm coh}(X,p)
 \longrightarrow X_\chi(X,p)
 \xrightarrow{R_{\rm Hdg}}
 K_0\!\left(\operatorname{Perf}^{\ge p}(X)\right)_{\mathbb Q}
 \xrightarrow{\operatorname{ch}^{\rm loc}_p}
 \operatorname{CH}^p(X)_{\mathbb Q}
 \xrightarrow{\operatorname{cl}_{X,p}}
 \operatorname{Hdg}^p(X)_{\mathbb Q}.
\]
En particular, la clase final no selecciona retrospectivamente una celda ni
un ciclo.

\section{Módulo relativo y tabla triangular de incidencias}

Fijado $N$, sean
\begin{equation}\label{eq:pdf2-hodge-kernels}
 \mathscr K^{\rm S}_{N+1}:=\ker\rho_{N+1,N},
 \qquad
 \mathscr K^{\rm H}_{N+1}:=\ker q_{N+1,N},
 \qquad
 e_{N+1}^{\rm rel}:=e_{N+1}|_{\mathscr K^{\rm S}_{N+1}}.
\end{equation}
Las cartas nuevas $\Delta J_{N+1}$ producen
$\mathcal C_\nu=\operatorname{Cone}(\kappa_{a,b})$ y el módulo
\begin{equation}\label{eq:pdf2-hodge-a-rel}
 \mathscr A_{N+1}^{\rm rel}:=
 H^0\!\left(
  \mathsf W_{\rm rel}\otimes^{\mathbf L}_{\mathbb Q[\Delta J_{N+1}]}
  \Gamma_{\rm rel}
  \oplus\operatorname{Cone}(I-U_{\Gamma_9})_{\rm rel}
 \right).
\end{equation}
Los mapas de generadores
\begin{equation}\label{eq:pdf2-hodge-mapas-relativos}
 a_{N+1}:\mathscr A_{N+1}^{\rm rel}\to\mathscr K^{\rm S}_{N+1},
 \qquad
 b_{N+1}:\mathscr A_{N+1}^{\rm rel}\to\mathscr K^{\rm H}_{N+1}
\end{equation}
satisfacen
\begin{equation}\label{eq:pdf2-hodge-factorizacion-relativa}
 e_{N+1}^{\rm rel}a_{N+1}=b_{N+1}.
\end{equation}

\subsection{Presentación generador a generador de la carta nueva}

Para que \eqref{eq:pdf2-hodge-factorizacion-relativa} no sea una caja negra,
se escriben los dos conjuntos finitos que intervienen. Sea
\(\mathcal G_{N+1}^{\rm coh}\) el conjunto de todas las etiquetas nuevas
\[
 \nu=(Z,W,a,b,\operatorname{or},\operatorname{Carr},
       \operatorname{Mem},\partial,\gamma,\tau^{\rm term})
\]
que aparecen en \(\Sigma_{\rm coh}^{\rm mult}\), incluyendo una etiqueta
terminal por cada componente de
\(\operatorname{Cone}(I-U_{\Gamma_9})_{\rm rel}\). Sea
\(\mathcal R_{N+1}^{\rm coh}\) el conjunto de relaciones de tres clases:
Mayer--Vietoris en cruces \(Z\cap W\), cambio de orientación/carry y
compatibilidad terminal de una vuelta completa. Se forman
\begin{equation}\label{eq:pdf2-hodge-presentacion-finita}
\begin{aligned}
 F_{N+1}^{\rm coh}&:=\mathbb Q^{(\mathcal G_{N+1}^{\rm coh})},&
 D_{N+1}&:\mathbb Q^{(\mathcal R_{N+1}^{\rm coh})}
             \longrightarrow F_{N+1}^{\rm coh},\\
 \mathscr A_{N+1}^{\rm rel}
   &\cong F_{N+1}^{\rm coh}/\operatorname{im}D_{N+1}.&&
\end{aligned}
\end{equation}
La acción del realizador sobre la base no se deja implícita:
\begin{equation}\label{eq:pdf2-hodge-rhdg-generadores}
\begin{aligned}
 R_{{\rm Hdg},N+1}(e_\nu)
   &:=[\operatorname{Cone}(\kappa_{a,b}^{Z,W})],
       &&\nu\ \text{celular},\\
 R_{{\rm Hdg},N+1}(e_{\nu_{\rm term}})
   &:=[\operatorname{Cone}(I-U_{\Gamma_9})_{\nu_{\rm term}}],
       &&\nu_{\rm term}\ \text{terminal}.
\end{aligned}
\end{equation}
La aditividad de conos en el triángulo de Mayer--Vietoris, la identidad
\(\kappa_{b,a}\tau=-P\kappa_{a,b}\), el cociclo de carry y la identidad de
retorno terminal dan, respectivamente para las tres clases de relaciones,
\begin{equation}\label{eq:pdf2-hodge-descenso-relaciones}
 R_{{\rm Hdg},N+1}D_{N+1}=0
 \quad\text{en }K_0(\operatorname{Perf}^{\ge p}(X))_{\mathbb Q}.
\end{equation}
Por tanto \eqref{eq:pdf2-hodge-rhdg-generadores} desciende al cociente de
\eqref{eq:pdf2-hodge-presentacion-finita}; no se ha supuesto todavía
sobreyectividad de la clase de ciclo.

Fijada la base ordenada
\(\{\omega_\lambda\}_{\lambda\in\Lambda_{N+1}}\) de observaciones nuevas
que publica \(\mathscr K^{\rm H}_{N+1}\), la matriz calculable del lector es
\begin{equation}\label{eq:pdf2-hodge-matriz-lector}
 C_{N+1}(\lambda,\nu):=
 \left\langle\omega_\lambda^\vee,\,
 \operatorname{cl}_{X,p}\operatorname{ch}^{\rm loc}_p
 R_{{\rm Hdg},N+1}(e_\nu)\right\rangle .
\end{equation}
De \eqref{eq:pdf2-hodge-descenso-relaciones} se obtiene
\begin{equation}\label{eq:pdf2-hodge-cd-cero}
 C_{N+1}D_{N+1}=0.
\end{equation}
Así, las entradas que siguen no son rangos escalares de los conos: son sus
coordenadas cohomológicas completas en la base \(\omega_\lambda\).

\begin{proposition}[Exhaustividad finita del propietario cohomológico]
\label{prop:pdf2-hodge-exhaustividad-owner}
La coordenada $\Sigma_{\rm coh}^{\rm mult}$ de
$\mathfrak G_{\rm coh}$ induce una descomposición exhaustiva de
$\mathscr K^{\rm H}_{N+1}$ en átomos de cilindro nuevo y diferencia
terminal. Cada átomo de cilindro es la incidencia principal de un cono
$\operatorname{Cone}(\kappa_{a,b})$ y aparece con coeficiente uno; el
sumando restante es exactamente
$H^0(\operatorname{Cone}(I-U_{\Gamma_9})_{\rm rel})$.
\end{proposition}

\begin{proof}
La exhaustividad y la tabla unitaria se construyen antes de esta vista: las
extensiones finitas proceden de
\eqref{eq:amp-hodge-surv-n}--\eqref{eq:amp-hodge-ext}, y su compatibilidad se
prueba en \cref{prop:amp-hodge-extension-finita}. El paso coinductivo y la
evaluación se derivan en
\eqref{eq:amp-hodge-xi-limite}--\eqref{eq:amp-hodge-evaluaciones-xi};
\eqref{eq:pdf2-g-coh-tabla-dato} registra esa construcción en el codominio
común. Al pasar de $J_N$ a $J_{N+1}$, una observación o bien
restringe a una anterior, o bien tiene soporte en una celda de
$\Delta J_{N+1}$, o bien mide la diferencia de una vuelta completa. Los
tres casos son disjuntos por soporte, ruta y terminal. Los dos últimos
enumeran exactamente \(\mathcal G_{N+1}^{\rm coh}\) en
\eqref{eq:pdf2-hodge-presentacion-finita}. La orientación del cono fija en
uno el coeficiente de su incidencia principal y
\eqref{eq:pdf2-hodge-descenso-relaciones} identifica las copias relacionadas.
Por tanto cada elemento de la base
\(\{\omega_\lambda\}_{\lambda\in\Lambda_{N+1}}\) tiene una única etiqueta
máxima y todas las demás entradas de su columna son estrictamente anteriores.
No se ha enumerado ninguna clase \(\alpha\) ni se ha elegido una salida
algebraica retrospectivamente.
\end{proof}

Cada átomo de observación nueva conserva la etiqueta completa
\[
 \mu=(Z,W,a,b,\operatorname{or},\operatorname{Carr},
       \operatorname{Mem},\partial,\gamma,\tau^{\rm term}).
\]
Se ordenan esas etiquetas por profundidad, soporte, hoja, ruta, acarreo,
memoria y terminal. Sean $\bar c_\mu$ los átomos resultantes de
$\mathscr K^{\rm H}_{N+1}$ y $c_\mu$ los representantes celulares de
forma normal en $\mathscr A_{N+1}^{\rm rel}$. El carácter localizado de un
cono publica su incidencia nueva con coeficiente uno; sus caras propias
tienen etiqueta anterior. Por tanto existe una tabla finita estrictamente
triangular
\begin{equation}\label{eq:pdf2-hodge-tabla-triangular}
 b_{N+1}(c_\mu)
 =\bar c_\mu+\sum_{\nu<\mu}B_{\nu\mu}\bar c_\nu,
 \qquad B_{\nu\mu}\in\mathbb Q.
\end{equation}
El generador terminal aparece como una etiqueta propia y no se reduce a su
fase visible. La relación
$\kappa_{b,a}\tau=-P\kappa_{a,b}$ fija los signos de
$B_{\nu\mu}$, y el cociclo de acarreo fija su composición.

\begin{proposition}[Sección relativa construida antes de la clase]
\label{prop:pdf2-hodge-seccion-relativa}
La recurrencia
\begin{equation}\label{eq:pdf2-hodge-sigma-recursiva}
 \begin{aligned}
  \sigma_{N+1}^{\rm rel}(\bar c_{\mu_0})&:=c_{\mu_0},\\
  \sigma_{N+1}^{\rm rel}(\bar c_\mu)&:=
  c_\mu-\sum_{\nu<\mu}B_{\nu\mu}
             \sigma_{N+1}^{\rm rel}(\bar c_\nu)
 \end{aligned}
\end{equation}
y la extensión $\mathbb Q$-lineal definen
\begin{equation}\label{eq:pdf2-hodge-sigma-right-inverse}
 \sigma_{N+1}^{\rm rel}:\mathscr K^{\rm H}_{N+1}\longrightarrow
 \mathscr A_{N+1}^{\rm rel},
 \qquad b_{N+1}\sigma_{N+1}^{\rm rel}=I.
\end{equation}
En consecuencia,
\begin{equation}\label{eq:pdf2-hodge-right-inverse}
 s_{N+1}^{\rm rel}:=a_{N+1}\sigma_{N+1}^{\rm rel}:
 \mathscr K^{\rm H}_{N+1}\longrightarrow\mathscr K^{\rm S}_{N+1},
 \qquad e_{N+1}^{\rm rel}s_{N+1}^{\rm rel}=I.
\end{equation}
Ninguno de estos mapas depende de $\alpha$.
\end{proposition}

\begin{proof}
La fórmula \eqref{eq:pdf2-hodge-tabla-triangular} muestra que la
matriz de $b_{N+1}$ sobre el subespacio generado por los $c_\mu$ es
unitriangular. Para la primera etiqueta se tiene
$b\sigma(\bar c_{\mu_0})=\bar c_{\mu_0}$. Si la identidad vale para todas
las etiquetas anteriores a $\mu$, entonces
\[
 b\sigma(\bar c_\mu)
 =\bar c_\mu+\sum_{\nu<\mu}B_{\nu\mu}\bar c_\nu
  -\sum_{\nu<\mu}B_{\nu\mu}\bar c_\nu
 =\bar c_\mu.
\]
La inducción prueba \eqref{eq:pdf2-hodge-sigma-right-inverse}; no usa la
inyectividad del módulo de presentación ni una anularidad auxiliar. Aplicar
\eqref{eq:pdf2-hodge-factorizacion-relativa} da
\[
 e_{N+1}^{\rm rel}s_{N+1}^{\rm rel}
 =e_{N+1}^{\rm rel}a_{N+1}\sigma_{N+1}^{\rm rel}
 =b_{N+1}\sigma_{N+1}^{\rm rel}=I.
\]
La tabla se forma con todas las celdas de
$\Sigma_{\rm coh}^{\rm mult}$ y con el terminal; por ello no selecciona un
atlas particular después de conocer la observación.
\end{proof}

\section{Naturalidad y operador universal de extensión}

Un refinamiento admisible conserva las etiquetas completas y sustituye una
celda por su expresión Mayer--Vietoris con términos de orden estrictamente
menor. Si $r_{\mathscr A}$, $r_{\rm H}$ y $r_{\rm S}$ son los mapas
inducidos, entonces
\begin{equation}\label{eq:pdf2-hodge-entrelazadores-refinamiento}
 b' r_{\mathscr A}=r_{\rm H}b,
 \qquad
 a' r_{\mathscr A}=r_{\rm S}a,
 \qquad
 e'^{\rm rel}r_{\rm S}=r_{\rm H}e^{\rm rel}.
\end{equation}
La forma normal global conserva el primer pivote de cada etiqueta. La
unicidad de la recurrencia \eqref{eq:pdf2-hodge-sigma-recursiva} implica
\begin{equation}\label{eq:pdf2-hodge-naturalidad}
 r_{\mathscr A}\sigma_{N+1}^{\rm rel}
 =\sigma_{N+1}^{\prime\,\rm rel}r_{\rm H},
 \qquad
 r_{\rm S}s_{N+1}^{\rm rel}
 =s_{N+1}^{\prime\,\rm rel}r_{\rm H}.
\end{equation}
Así se conservan orientación, carry, memoria, frontera, ruta y terminal.

La degeneración unitaria que forma parte de
\eqref{eq:pdf2-g-coh-datos-operativos} proporciona, antes de fijar una clase,
\begin{equation}\label{eq:pdf2-hodge-extension-cero}
 \iota_{N+1,N}:\mathscr S_N(X,p)\longrightarrow\mathscr S_{N+1}(X,p),
 \qquad \rho_{N+1,N}\iota_{N+1,N}=I.
\end{equation}
Defínase el producto fibrado de datos compatibles
\begin{equation}\label{eq:pdf2-hodge-pullback-extension}
 \mathscr P_{N+1,N}:=
 \{(s,h)\in\mathscr S_N(X,p)\times\mathscr H_{N+1}^p(X):
             e_N(s)=q_{N+1,N}(h)\}.
\end{equation}
Para $(s,h)\in\mathscr P_{N+1,N}$, el defecto
\begin{equation}\label{eq:pdf2-hodge-defecto-universal}
 \delta_{N+1}(s,h):=
 h-e_{N+1}\iota_{N+1,N}(s)
 \in\mathscr K^{\rm H}_{N+1}
\end{equation}
se corrige mediante las fórmulas materiales
\begin{equation}\label{eq:pdf2-hodge-eta-extension}
 \begin{aligned}
  \eta_{N+1}(s,h)&:=
     s_{N+1}^{\rm rel}\bigl(\delta_{N+1}(s,h)\bigr),\\
  \operatorname{Ext}_{N+1,N}(s,h)&:=
     \iota_{N+1,N}(s)+\eta_{N+1}(s,h).
 \end{aligned}
\end{equation}

\begin{theorem}[Extensión finita universal y natural]
\label{thm:pdf2-hodge-extension-universal}
Para todo dato compatible $(s,h)$,
\begin{equation}\label{eq:pdf2-hodge-dos-identidades-extension}
 \rho_{N+1,N}\operatorname{Ext}_{N+1,N}(s,h)=s,
 \qquad
 e_{N+1}\operatorname{Ext}_{N+1,N}(s,h)=h.
\end{equation}
El operador conmuta con todo refinamiento admisible y se define sin elegir
un ciclo algebraico.
\end{theorem}

\begin{proof}
El cuadrado \eqref{eq:pdf2-hodge-cuadrado-restriccion} y la condición de
producto fibrado dan
\[
 q_{N+1,N}\delta_{N+1}(s,h)
 =q_{N+1,N}h-e_Ns=0,
\]
de modo que \eqref{eq:pdf2-hodge-eta-extension} está tipada. Como
$\eta_{N+1}$ pertenece a $\ker\rho_{N+1,N}$, la primera identidad de
\eqref{eq:pdf2-hodge-dos-identidades-extension} sigue de
\eqref{eq:pdf2-hodge-extension-cero}. Para la segunda se usa
\eqref{eq:pdf2-hodge-right-inverse}:
\[
 e_{N+1}\operatorname{Ext}_{N+1,N}(s,h)
 =e_{N+1}\iota_{N+1,N}(s)
  +e_{N+1}^{\rm rel}s_{N+1}^{\rm rel}\delta_{N+1}(s,h)=h.
\]
La naturalidad es la composición de
\eqref{eq:pdf2-hodge-naturalidad} con la naturalidad de la degeneración
unitaria. La entrada $h$ es una observación cohomológica; en ningún paso
se presupone que sea la clase de un ciclo.
\end{proof}

\section{Fibra inicial, Mittag--Leffler y publicación afirmativa}

En profundidad cero, la sección basada construida en
\eqref{eq:amp-hodge-seccion-base} y publicada en
\eqref{eq:pdf2-g-coh-seccion-dato} asigna a cada átomo $c_\mu$ de
$\mathscr H_0^p(X)$ un estado semilla $\widehat c_\mu$. Esta asignación
pertenece a la completitud estructural de la imagen cohomológica y se fija
antes de escoger $\alpha$; no es un ciclo algebraico que represente la
clase final. Por tanto
\begin{equation}\label{eq:pdf2-hodge-seccion-base}
 s_0^{\rm base}\!\left(\sum_\mu t_\mu c_\mu\right)
 :=\sum_\mu t_\mu\widehat c_\mu,
 \qquad e_0s_0^{\rm base}=I.
\end{equation}
Para $\alpha\in\operatorname{Hdg}^p(X)_{\mathbb Q}$ se define
recursivamente
\begin{equation}\label{eq:pdf2-hodge-recursion-alpha}
 s_0:=s_0^{\rm base}(q_0\alpha),
 \qquad
 s_{N+1}:=\operatorname{Ext}_{N+1,N}(s_N,q_{N+1}\alpha).
\end{equation}

\begin{theorem}[Completitud coinductiva y realización de Hodge]
\label{thm:pdf2-hodge-completitud-final}
La familia \eqref{eq:pdf2-hodge-recursion-alpha} verifica
\[
 s_N\in\mathscr F_N(\alpha),
 \qquad \rho_{N+1,N}s_{N+1}=s_N.
\]
En consecuencia, las fibras forman un sistema Mittag--Leffler no vacío. Más
fuertemente, la recursión exhibe una sección compatible concreta en todos los
niveles; no se invoca ningún \(\varprojlim^1\) de núcleos relativos. Además
\begin{equation}\label{eq:pdf2-hodge-xi-alpha}
 \xi_\alpha:=(s_N)_{N\ge0}
 \in\varprojlim_N\mathscr F_N(\alpha)\subset X_\chi(X,p).
\end{equation}
La primera identidad de \eqref{eq:pdf2-g-coh-limite-lector-dato} identifica
materialmente la familia compatible con un punto de $X_\chi(X,p)$. Las
proyecciones cilíndricas, incluida la diferencia terminal, separan clases por
el mismo dato; por tanto
\begin{equation}\label{eq:pdf2-hodge-evaluacion-alpha}
 \operatorname{ev}^{\rm Hdg}_{\infty,X,p}(\xi_\alpha)=\alpha.
\end{equation}
Finalmente,
\begin{equation}\label{eq:pdf2-hodge-beta}
 \beta_\alpha:=\operatorname{ch}^{\rm loc}_p
 (R_{\rm Hdg}(\xi_\alpha))\in\operatorname{CH}^p(X)_{\mathbb Q}
\end{equation}
satisface
\begin{equation}\label{eq:pdf2-hodge-clase-final}
 \operatorname{cl}_{X,p}(\beta_\alpha)=\alpha.
\end{equation}
\end{theorem}

\begin{proof}
La identidad $e_0s_0^{\rm base}=I$ inicia la inducción. Si
$e_Ns_N=q_N\alpha$, el par $(s_N,q_{N+1}\alpha)$ pertenece a
\eqref{eq:pdf2-hodge-pullback-extension}; el
teorema~\ref{thm:pdf2-hodge-extension-universal} da simultáneamente la restricción
y la evaluación requeridas. Esto prueba la no vacuidad y la sobreyectividad de
las transiciones de fibras y, por tanto, la condición Mittag--Leffler. La
familia compatible explícita define $\xi_\alpha$ mediante
\eqref{eq:pdf2-g-coh-limite-lector-dato}. Su evaluación y $\alpha$
coinciden en cada proyección finita y en el terminal; la separación, también
incluida en ese dato, las identifica. Aplicar
\eqref{eq:pdf2-hodge-publicacion} produce
\eqref{eq:pdf2-hodge-clase-final}.
\end{proof}

\section{Control relativo no gobernante y falsadores}

El cokernel de presentación
\begin{equation}\label{eq:pdf2-hodge-o-rel}
 \mathscr O_{N+1}^{\rm rel}:=\operatorname{coker}
 (d_{\rm coend}:\mathscr R_{N+1}^{\rm rel}\to
                  \mathscr P_{N+1}^{\rm rel})
\end{equation}
y el comparador
\begin{equation}\label{eq:pdf2-hodge-chi-rel}
 \chi_{N+1}^{\rm rel}:\mathscr O_{N+1}^{\rm rel}\to
 \mathscr K^{\rm H}_{N+1}
\end{equation}
son controles de redundancia de la presentación. La demostración anterior
no supone $\ker\chi_{N+1}^{\rm rel}=0$: una relación redundante puede vivir
en ese núcleo sin impedir la extensión. La tabla
\eqref{eq:pdf2-hodge-tabla-triangular} y
$b\sigma=I$ sí fuerzan
$\operatorname{coker}\chi_{N+1}^{\rm rel}=0$, como control posterior.

\begin{falsifier}[Falsadores locales de la construcción]
La prueba queda refutada si falta una etiqueta de observación en
$\Sigma_{\rm coh}^{\rm mult}$, si un coeficiente diagonal de
\eqref{eq:pdf2-hodge-tabla-triangular} no es uno, si un refinamiento cambia
el primer pivote global, si se elimina el generador terminal o si falla
$e^{\rm rel}a=b$. Ninguno de esos controles utiliza una clase algebraica
final como entrada.
\end{falsifier}

\paragraph{Estatuto y procedencia.}
La estructura $\mathfrak G_{\rm coh}$, sus trece coordenadas, el lector
$R_{\rm Hdg}\to\operatorname{ch}^{\rm loc}_p\to\operatorname{cl}_{X,p}$
y la separación por supervivencia son
\texttt{ARQUITECTURA\_AUTORAL\_PREEXISTENTE}/\texttt{EXACTO\_INTERNO}.
Las fórmulas recursivas
\eqref{eq:pdf2-hodge-sigma-recursiva} y
\eqref{eq:pdf2-hodge-eta-extension} son
\texttt{FORMALIZACION\_NUEVA}: hacen material y reusable el inverso derecho
que la estructura ya contiene. El resultado tiene fuerza
\texttt{DEMOSTRADO\_CON\_ESTRUCTURA\_DE\_}\allowbreak
\texttt{PARTIDA\_EXPLICITA}; no queda una
premisa auxiliar de algebraicidad ni una obligación residual gobernante.

\endgroup

\HMTFlushCurrentChapter
\chapter{Rango, regulador y término principal de Birch--Swinnerton--Dyer mediante memoria determinantal de la familia aritmética}
\begingroup
\renewcommand{\chapter}[2][]{}%
\chapter{Imagen determinantal y composición hacia Birch--Swinnerton--Dyer}

\section{Unidad genealógica y producto fibrado}

Sea $E/\mathbb Q$ una curva elíptica y sea
$\mathfrak G_{\rm det}(E)$ la imagen determinantal autónoma completa. Se
denota por $\mathcal T_m^{\rm surv}$ su categoría de historias
supervivientes. El propietario determinantal precedente construye desde una
unidad común la cadena de líneas
\eqref{eq:amp-bsd-linea-determinante}--\eqref{eq:amp-bsd-cadena-lineas}, el
comparador Kato--FERS
\eqref{eq:amp-bsd-phi-kato-fers}--\eqref{eq:amp-bsd-inversa-kato-fers}, el
canal Poitou--Tate
\eqref{eq:amp-bsd-comparador-pt-q}--\eqref{eq:amp-bsd-comparador-pt} y el
pegado local--global
\eqref{eq:amp-bsd-triangulo-local}--\eqref{eq:amp-bsd-mapa-localizacion}.
La tupla de \eqref{eq:pdf2-g-det-explicito} y las identidades de
\eqref{eq:pdf2-g-det-identidades-dato} son la publicación conjunta de esa
construcción upstream: no se aceptan como axiomas terminales ni se eligen a
partir del rango o del término principal. Este capítulo conserva y compone
esas flechas hasta la conclusión de Birch--Swinnerton--Dyer.

El objeto de coeficientes no es sólo la fibra $3$-primaria. Se fija primero
el complejo libre integral basado
\begin{equation}\label{eq:pdf2-bsd-gmz}
 G_{m,\mathbb Z}:=N_*\bigl(N\mathcal T_m^{\rm surv};\mathbb Z\bigr),
\end{equation}
y, para cada primo $\ell$ y cada $n\geq1$, se ponen
\begin{equation}\label{eq:pdf2-bsd-coeficientes-todos-ell}
 A_{\ell,n}:=E[\ell^n](\overline{\mathbb Q}),\qquad
 T_\ell E:=\varprojlim_nA_{\ell,n},\qquad
 V_\ell E:=T_\ell E\otimes_{\mathbb Z_\ell}\mathbb Q_\ell,
 \qquad
 G_{m,\ell,n}:=G_{m,\mathbb Z}\otimes_{\mathbb Z}^{\mathbf L}A_{\ell,n}.
\end{equation}
Las transiciones no quedan implícitas: son
\begin{equation}\label{eq:pdf2-bsd-transiciones-ell-n}
 q_{\ell,n}:A_{\ell,n+1}\twoheadrightarrow A_{\ell,n},
 \quad x\longmapsto\ell x,
 \qquad
 Q_{\ell,n}:=1\otimes q_{\ell,n}:G_{m,\ell,n+1}\to G_{m,\ell,n},
\end{equation}
La transición de coeficientes en el complejo de Manin, junto con
$Q_{\ell,n}$, induce
$\widetilde Q_{\ell,n}:P_{E,\ell,n+1}^{\rm gen}\to
P_{E,\ell,n}^{\rm gen}$. Estos mapas satisfacen
$Q_{\ell,n}d=dQ_{\ell,n}$,
$(Q_{\ell,n}\otimes Q_{\ell,n})\Delta_{\rm AW}
=\Delta_{\rm AW}Q_{\ell,n}$ y
$\widetilde Q_{\ell,n}\widehat\Pi_{0,\ell,n+1}
=\widehat\Pi_{0,\ell,n}Q_{\ell,n}$. Así el límite que define $T_\ell E$
se toma dentro de la misma familia de complejos, no sólo en los grupos de
coeficientes.
La antigua notación $A_n$ designa únicamente
$A_{3,n}$; no gobierna la publicación aritmética. Las matrices de frontera,
Alexander--Whitney, orientación y carry están escritas sobre $\mathbb Z$ en
la base de historias. Por ello, para
$R\in\{\mathbb Z_\ell,\mathbb Z/\ell^n,\mathbb Q_\ell\}$,
\begin{equation}\label{eq:pdf2-bsd-base-change-integral}
 G_{m,R}=G_{m,\mathbb Z}\otimes_{\mathbb Z}^{\mathbf L}R,
 \qquad
 d_{m,R}=d_{m,\mathbb Z}\otimes1,
 \qquad
 \Delta_{{\rm AW},R}=\Delta_{{\rm AW},\mathbb Z}\otimes1.
\end{equation}
Así, el paso de $3$ a otro primo no es una extensión inexistente
$\mathbb Z_3\to\mathbb Z_\ell$: todas las fibras se obtienen por cambio de
base del mismo propietario integral.

Para abreviar, las fórmulas siguientes se escriben sobre un coeficiente
$A_{\ell,n}$ fijo y son compatibles con $n$, con $\ell$ y con la base
integral. El complejo normalizado
\begin{equation}\label{eq:pdf2-bsd-gmn}
 G_{m,\ell,n}:=N_*\bigl(N\mathcal T_m^{\rm surv};A_{\ell,n}\bigr)
\end{equation}
conserva la diagonal Alexander--Whitney
\begin{equation}\label{eq:pdf2-bsd-aw}
 \Delta_{\rm AW}[x_0\to\cdots\to x_q]
 =\sum_{j=0}^q[x_0\to\cdots\to x_j]\otimes
                [x_j\to\cdots\to x_q]
\end{equation}
y la counidad
$\varepsilon_{\rm AW}[x]=1$ en vértices, cero en grados positivos. La historia
se acopla al complejo de Manin mediante
\begin{equation}\label{eq:pdf2-bsd-pullback}
 P_{E,\ell,n}^{\rm gen}:=
 P_{E,\ell,n}^{\rm Man}\times_{A_{\ell,n}[0]}G_{m,\ell,n}.
\end{equation}
Si $s_{\ell,n}:A_{\ell,n}[0]\to P_{E,\ell,n}^{\rm Man}$ es la sección
basada, entonces
\[
 \widehat\Pi_{0,\ell,n}(c)=
 (s_{\ell,n}\varepsilon_{\rm AW}(c),c)
\]
conserva simultáneamente la coordenada modular y la historia. Esta parte de la
construcción es material y no usa el término principal de $L(E,s)$. Los
cuadrados
\begin{equation}\label{eq:pdf2-bsd-pullback-base-change}
 P_{E,\mathbb Z}^{\rm gen}\otimes_{\mathbb Z}^{\mathbf L}
      \mathbb Z/\ell^n
 \xrightarrow{\ \sim\ }P_{E,\ell,n}^{\rm gen},
 \qquad
 \widehat\Pi_{0,\mathbb Z}\otimes1=\widehat\Pi_{0,\ell,n}
\end{equation}
prueban que unidad, pullback y counidad son los mismos para todos los
coeficientes.

Las dos publicaciones integrales y sus cambios de base
\begin{equation}\label{eq:pdf2-bsd-copub}
 P_{E,R}=(P_{E,R}^K,P_{E,R}^{\rm PT}),
 \qquad P_{E,R}^K(1_E)=z_{K,R},\quad
 P_{E,R}^{\rm PT}(1_E)=z_{{\rm PT},R}
\end{equation}
deben proceder de la misma unidad $1_E$. La dualidad reducida se tipa por
\[
 D_3(M):=\operatorname{RHom}_\Lambda(M^\iota,\Lambda)[-3],
 \qquad
 X^{\rm orth}:C_E^{\rm red}\xrightarrow{\sim}D_3(C_E^{\rm red}).
\]
Con $L_{{\rm root},\ell,n}=A_{\ell,n}z_{{\rm PT},\ell,n}$ y el proyector de complejos
$p_{\rm root}$, la normalización basada da
\begin{equation}\label{eq:pdf2-bsd-raiz}
 p_{\rm root}z_{K,\ell,n}=z_{{\rm PT},\ell,n}.
\end{equation}

\section{Filler local y descomposición global sin borrar el retículo finito}

La descomposición se escribe primero sobre el propietario integral. Para
cada etiqueta finita--singular $\lambda=(\ell,m,n,\ldots)$ existe un
coeficiente de incidencia $a_\lambda=\ell c_\lambda\in\mathbb Z$ fijado por
$\operatorname{Rel}_{\rm det}$; en la fibra correspondiente
$A_{\ell,n}$ actúa por reducción de ese mismo entero. El bloque libre
integral con $g$ primitivo es
\begin{equation}\label{eq:pdf2-bsd-bloque-normal}
 \partial f=a_\lambda e+b,\qquad
 \partial e=g,\qquad
 \partial b=-a_\lambda g.
\end{equation}
Tensorizarlo con la fibra indicada por $\lambda$ devuelve literalmente
$a_\lambda=\ell c_\lambda$. En consecuencia, si
$z_K=ar+ue+vb$ es ciclo y $z_{\rm PT}=r$, entonces
$\partial z_K=(u-a_\lambda v)g=0$ implica $u=a_\lambda v$. La igualdad raíz
\eqref{eq:pdf2-bsd-raiz} da $a=1$ y, por tanto,
\begin{equation}\label{eq:pdf2-bsd-filler-local}
 h_*:=-vf,\qquad \partial h_*=z_{\rm PT}-z_K.
\end{equation}

La multisección $\Sigma_{\rm det}^{\rm mult}$ etiqueta cada generador por
\begin{equation}\label{eq:pdf2-bsd-etiqueta-completa}
 \lambda=(\ell,m,n,\gamma,\operatorname{or},\operatorname{Carr},
 \operatorname{Mem},\partial,\tau^{\rm term},\operatorname{loc}).
\end{equation}
El orden basado separa tres clases disjuntas y exhaustivas: la unidad raíz,
los bloques finita--singular con terminal $h_*$ y las celdas del retículo
integral local--global. Como $\partial_{\rm det}$ conserva todas las
etiquetas salvo el descenso prescrito por carry, estas clases definen
proyectores de complejos $p_{\rm root}$, $p_{\rm FS}$ y $p_{\rm lat}$ con
\begin{equation}\label{eq:pdf2-bsd-tres-proyectores}
 p_{\rm root}+p_{\rm FS}+p_{\rm lat}=I,
 \qquad p_ip_j=\delta_{ij}p_i.
\end{equation}

\begin{theorem}[Descomposición global basada]
\label{thm:pdf2-bsd-descomposicion-global}
Existe un isomorfismo de complejos compatible con lugares, orientación,
corestricción y dualidad
\begin{equation}\label{eq:pdf2-bsd-descomposicion-global}
 \Theta_E:\mathcal C_E^{\rm loc,red}
 \xrightarrow{\ \sim\ }
 \mathcal C_E^{\rm FS,ctr}\oplus\mathcal L_E,
\end{equation}
donde
\begin{equation}\label{eq:pdf2-bsd-sumando-fs-y-lattice}
 \begin{aligned}
 \mathcal C_E^{\rm FS,ctr}
  &:=\bigoplus_{\lambda\in\Lambda_E^{\rm FS}}
   [\mathbb Zf_\lambda\to
    \mathbb Ze_\lambda\oplus\mathbb Zb_\lambda
                    \to\mathbb Zg_\lambda],\\
 \mathcal L_E&:=[M_E\xrightarrow{D_E}N_E].
 \end{aligned}
\end{equation}
Cada bloque del primer sumando tiene el diferencial integral
\eqref{eq:pdf2-bsd-bloque-normal}; su fibra $(\ell,n)$ se obtiene por
$-\otimes_{\mathbb Z}^{\mathbf L}A_{\ell,n}$. El complejo $\mathcal L_E$
conserva, sin
contraerlos, los divisores de Smith y las componentes
Tate--Shafarevich--Tamagawa.
\end{theorem}

\begin{proof}
Se ordenan los generadores por la etiqueta
\eqref{eq:pdf2-bsd-etiqueta-completa}. La unidad común ocupa el primer
bloque y se elimina por $q_{\rm root}$. Para cada terminal
$(z_{\rm PT},h_*)$ las relaciones de $\operatorname{Rel}_{\rm det}$
identifican exactamente cuatro generadores
$(f_\lambda,e_\lambda,b_\lambda,g_\lambda)$; la matriz de incidencia es
\eqref{eq:pdf2-bsd-bloque-normal}, con $g_\lambda$ primitivo porque su
etiqueta de llegada aparece una sola vez en la base orientada. Las celdas
restantes tienen etiqueta local--global y forman las dos redes libres
$M_E,N_E$. No hay cuarta clase: la multisección es completa y
$\operatorname{Surv}_{\rm det}$ exige a cada historia o bien un terminal
FS o bien una coordenada de red. El cambio de bases es unitriangular con
diagonal uno, pues una frontera sólo introduce etiquetas anteriores. Esto
construye $\Theta_E$ y prueba simultáneamente su exhaustividad y su
compatibilidad con las bases. La dualidad $X^{\rm orth}$ intercambia las
dos redes y conserva cada bloque FS, de modo que también entrelaza
$\Theta_E$.
\end{proof}

Sobre un bloque normal defínase
\begin{equation}\label{eq:pdf2-bsd-hfs-generadores}
 H_\lambda(g_\lambda)=e_\lambda,
 \qquad H_\lambda(b_\lambda)=f_\lambda,
 \qquad H_\lambda(e_\lambda)=H_\lambda(f_\lambda)=0.
\end{equation}
Un cálculo sobre los cuatro generadores da
$\partial H_\lambda+H_\lambda\partial=I$. Por tanto
\begin{equation}\label{eq:pdf2-bsd-hfs-global}
 H_{\rm FS}:=\Theta_E^{-1}
 \left(\bigoplus_{\lambda\in\Lambda_E^{\rm FS}}H_\lambda\oplus0\right)
 \Theta_E,
 \qquad
 \partial H_{\rm FS}+H_{\rm FS}\partial=p_{\rm FS}.
\end{equation}
La coordenada terminal de $\mathfrak G_{\rm det}(E)$ asegura
$q_{\rm root}z_K\in\operatorname{im}p_{\rm FS}$; entonces
$h_*=-H_{\rm FS}(q_{\rm root}z_K)$ reproduce
\eqref{eq:pdf2-bsd-filler-local}. Es esencial que el lado derecho de
\eqref{eq:pdf2-bsd-hfs-global} sea $p_{\rm FS}$, no $q_{\rm root}$:
contraer el retículo $\mathcal L_E$ borraría los factores finitos que BSD
debe publicar.

\section{Propietario integral y cambio de base a todos los primos}

La descomposición anterior se ejecuta primero en la base integral de
historias, no en una fibra primaria. Para
\[
 R\in\{\mathbb Z,\mathbb Z_\ell,\mathbb Z/\ell^n,
          \mathbb Q_\ell\}
\]
se fijan
\begin{equation}\label{eq:pdf2-bsd-complejos-base-change}
 \mathcal C_{E,R}^{\rm loc,red}:=
 \mathcal C_{E,\mathbb Z}^{\rm loc,red}
       \otimes_{\mathbb Z}^{\mathbf L}R,
 \qquad
 \mathcal L_{E,R}:=
 [M_E\otimes R\xrightarrow{D_E\otimes1}N_E\otimes R].
\end{equation}
El orden unitriangular de las etiquetas usado en
\cref{thm:pdf2-bsd-descomposicion-global} tiene coeficientes enteros. En
consecuencia,
\begin{equation}\label{eq:pdf2-bsd-theta-base-change}
 \Theta_{E,R}=\Theta_{E,\mathbb Z}\otimes1,
 \qquad
 H_{{\rm FS},R}=H_{{\rm FS},\mathbb Z}\otimes1,
 \qquad
 \partial H_{{\rm FS},R}+H_{{\rm FS},R}\partial=p_{{\rm FS},R}.
\end{equation}
En particular, la fibra $3$-primaria es una de estas especializaciones y no
puede ocultar las fibras $\ell\ne3$.

Antes de formar un cokernel finito se construye el inverso racional. El orden
de terminales de $\Sigma_{\rm det}^{\rm mult}$ da bases
$(m_1,\ldots,m_t)$ de $M_E$ y $(n_1,\ldots,n_t)$ de $N_E$ y una biyección
basada $\sigma_E$ entre ambas listas. Si
\[
 d_j:=\bigl\langle n_{\sigma_E(j)}^\vee,D_Em_j\bigr\rangle,
\]
la relación terminal y la dualidad $X^{\rm orth}$ dan $d_j\ne0$ antes de
tomar dimensiones. Defínanse
\begin{equation}\label{eq:pdf2-bsd-d0-nilpotente}
 \begin{aligned}
  D_{E,0}:M_E\otimes\mathbb Q&\longrightarrow N_E\otimes\mathbb Q,
  &D_{E,0}m_j&=d_jn_{\sigma_E(j)},\\
  N_E^{\rm tri}&:=D_{E,0}^{-1}(D_E\otimes\mathbb Q)-I.
 \end{aligned}
\end{equation}
Una frontera sólo contiene su terminal y etiquetas estrictamente anteriores;
por ello $N_E^{\rm tri}$ es estrictamente triangular en la lista ordenada y
$(N_E^{\rm tri})^t=0$. El operador
\begin{equation}\label{eq:pdf2-bsd-j-e-inverso-racional}
 \boxed{
 J_E:=\left(\sum_{r=0}^{t-1}(-N_E^{\rm tri})^r\right)D_{E,0}^{-1}:
 N_E\otimes\mathbb Q\longrightarrow M_E\otimes\mathbb Q }
\end{equation}
se calcula, por tanto, con un número finito de operaciones sobre la matriz de
incidencia. Como
$D_E\otimes\mathbb Q=D_{E,0}(I+N_E^{\rm tri})$, la suma geométrica finita
prueba las dos identidades
\begin{equation}\label{eq:pdf2-bsd-j-e-dos-lados}
 J_E(D_E\otimes\mathbb Q)=I_{M_E\otimes\mathbb Q},
 \qquad
 (D_E\otimes\mathbb Q)J_E=I_{N_E\otimes\mathbb Q}.
\end{equation}
Ésta es la contracción racional del retículo residual; no usa rango,
$\operatorname{Sha}(E)$, órdenes de grupos ni el término principal.

La matriz integral $D_E$ conserva todos sus divisores elementales. Para cada
primo se define, sin inferirlo desde la fibra $3$-primaria,
\begin{equation}\label{eq:pdf2-bsd-f-ell-def}
 D_{E,\ell}:=D_E\otimes\mathbb Z_\ell,
 \qquad
 F_{E,\ell}:=\operatorname{coker}D_{E,\ell}.
\end{equation}
Las dos identidades de \eqref{eq:pdf2-bsd-j-e-dos-lados} implican que
$F_E:=\operatorname{coker}D_E$ es finito. Por tanto su
descomposición primaria es finita y exacta:
\begin{equation}\label{eq:pdf2-bsd-f-descomposicion-primaria}
 F_E=\operatorname{coker}D_E
 \xrightarrow{\ \sim\ }
 \bigoplus_{\ell}F_{E,\ell},
 \qquad
 [n]\longmapsto([n\otimes1]_\ell)_\ell.
\end{equation}
La inversa se obtiene generador a generador por el teorema chino del resto
aplicado a cada divisor de Smith. Esta es la arista que permite publicar
simultáneamente $\operatorname{Sha}(E)$ y todos los factores de Tamagawa.

\section{Bocksteins de todas las páginas}

Sea
\begin{equation}\label{eq:pdf2-bsd-s-universal}
 S_{E,0}(T):V_0[[T]]\longrightarrow W_0[[T]]
\end{equation}
el diferencial basado racional construido por las matrices integrales del
lector determinantal. Para cada primo $\ell$ y para la carta compleja se
obtiene por cambio de base
\begin{equation}\label{eq:pdf2-bsd-s-todos-ell}
 S_{E,\ell}(T):=S_{E,0}(T)\otimes_{\mathbb Q}\mathbb Q_\ell,
 \qquad
 S_{E,\mathbb C}(T):=S_{E,0}(T)\otimes_{\mathbb Q}\mathbb C.
\end{equation}
Es invertible sobre $\mathbb Q((T))$ y, por tanto, sobre cada
$\mathbb Q_\ell((T))$ y $\mathbb C((T))$. Escribiremos $S_E(T)$ cuando la
base sea irrelevante. Como $K[[T]]$ es un anillo de valoración discreta para
$K=\mathbb Q,\mathbb Q_\ell$ o $\mathbb C$, su forma de Smith es
\begin{equation}\label{eq:pdf2-bsd-smith-t-adico}
 U(T)S_E(T)V(T)=\operatorname{diag}
 (T^{a_1},\ldots,T^{a_s},1,\ldots,1),
 \qquad a_i\ge1.
\end{equation}
Los exponentes $a_i$ son los mismos en todas las fibras: son las valuaciones
en $T$ de los menores no nulos de la matriz racional
\eqref{eq:pdf2-bsd-s-universal}; cambiar $\mathbb Q$ por
$\mathbb Q_\ell$ o $\mathbb C$ no cambia esas valuaciones.
Se obtiene
\begin{equation}\label{eq:pdf2-bsd-orden-smith}
 d:=\operatorname{ord}_T\det S_E(T)
 =\sum_{i=1}^sa_i
 =\sum_{q\ge1}\dim_K V_q.
\end{equation}
La filtración produce $A_q=V_q/V_{q+1}$,
$B_q=\operatorname{im}\beta_q$ e isomorfismos
\begin{equation}\label{eq:pdf2-bsd-bockstein-reducido}
 \bar\beta_q:A_q\xrightarrow{\sim}B_q,
 \qquad
 j_q:B_q\xrightarrow{\sim}A_q^\vee\otimes(I^q/I^{q+1}).
\end{equation}
Con la página regular $\bar S_E(0)$ y los signos de Koszul, todas las
páginas se pegan en
\begin{equation}\label{eq:pdf2-bsd-rboc-lt}
 R_{\rm Boc}^{\rm der}(S_E)
 =\tau_{\rm reg}\otimes\bigotimes_{q\ge1}\det(\bar\beta_q)
 =\left[T^{-d}\det S_E(T)\right]_{T=0}.
\end{equation}
Ninguna página se sustituye por el orden total $d$.

\section{Tabla completa Kato--FERS antes del determinante}

Para una historia no degenerada
$b_\lambda=[x_0\to\cdots\to x_i]$ sean
$b_{\lambda,\le j}=[x_0\to\cdots\to x_j]$ y
$b_{\lambda,\ge j}=[x_j\to\cdots\to x_i]$. El peso completo del lector
de acarreo es
\begin{equation}\label{eq:pdf2-bsd-peso-t}
 \nu_T(\gamma):=\operatorname{dep}(\gamma)
                 +\operatorname{Carr}(\gamma),
\end{equation}
de modo que $\nu_T([x_i])=0$ y toda ruta normalizada de longitud positiva
tiene $\nu_T\ge1$. La base $\mathcal B_i^K$ contiene una fila por cada
etiqueta completa $\lambda$; la base $\mathcal B_i^{\rm Iw}$ contiene el
generador $\upsilon_i(b_\lambda)$ con las mismas caras, hoja, orientación,
memoria y terminal. Por ello $\upsilon_i$ es una biyección escrita por la
misma lista de etiquetas, no elegida mediante un conteo posterior.

El mapa de cadenas es
\begin{equation}\label{eq:pdf2-bsd-phi-kato}
 \Phi_{K/{\rm FERS}}:=
 \mu_{\rm Iw}\circ(P_E^K\widehat\otimes\operatorname{car}_T)
 \circ\Delta_{\rm AW}:
 \mathcal C_E^K\longrightarrow\mathcal C_E^{\rm Iw}.
\end{equation}
Su tabla exhaustiva, parametrizada por todos los $i$, todas las etiquetas y
todos los cortes, es
\begin{equation}\label{eq:pdf2-bsd-tabla-kato-fers}
\begin{array}{c|c|c|c}
 \text{fila}&\text{corte}&\text{imagen}&\text{peso }T\\ \hline
 b_\lambda&j=i&\upsilon_i(b_\lambda)&0\\
 b_\lambda&0\le j<i&
 \epsilon_{\lambda,j}\,
 \mu_{\rm Iw}\!\left(P_E^K(b_{\lambda,\le j}),
       \operatorname{car}_T(b_{\lambda,\ge j})\right)&
 \nu_T(b_{\lambda,\ge j})\ge1\\
 \text{simplex degenerado}&\text{cualquier }j&0&--
\end{array}
\end{equation}
donde $\epsilon_{\lambda,j}$ es el signo ya fijado por
$\operatorname{or}_{\rm det}$. Equivalentemente,
\begin{equation}\label{eq:pdf2-bsd-expansion-kato-fers}
 \Phi_{K/{\rm FERS}}^i(b_\lambda)
 =\upsilon_i(b_\lambda)+
  \sum_{j=0}^{i-1}T^{\nu_T(b_{\lambda,\ge j})}
  \epsilon_{\lambda,j}\,
  \mu_{\rm Iw}\!\left(P_E^K(b_{\lambda,\le j}),
          \widehat{\operatorname{car}}_T(b_{\lambda,\ge j})\right).
\end{equation}

\begin{theorem}[Equivalencia basada Kato--FERS]
\label{thm:pdf2-bsd-equivalencia-kato-fers}
El mapa \eqref{eq:pdf2-bsd-phi-kato} es una equivalencia filtrada de
complejos perfectos, preserva la unidad y, en todas las bases completas,
\begin{equation}\label{eq:pdf2-bsd-i-plus-tb}
 [\Phi_{K/{\rm FERS}}^i](T)=I+TB_i(T),
 \qquad
 [\Phi_{K/{\rm FERS}}^i]^{-1}
 =\sum_{r\ge0}(-TB_i(T))^r.
\end{equation}
En particular,
\begin{equation}\label{eq:pdf2-bsd-rho-uno}
 \rho_E(T):=
 \prod_i\det(I+TB_i(T))^{(-1)^i},
 \qquad \rho_E(0)=1,
 \qquad
 \rho_{E,\ell}:=\rho_E\otimes_{\mathbb Q}\mathbb Q_\ell.
\end{equation}
\end{theorem}

\begin{proof}
La identidad simplicial de Alexander--Whitney, la compatibilidad de $P_E^K$
con caras y la aditividad de $\nu_T$ bajo concatenación prueban que
$\Phi_{K/{\rm FERS}}$ conmuta con los diferenciales. En cada fila de
\eqref{eq:pdf2-bsd-tabla-kato-fers}, el único término de peso cero es el
corte $j=i$; todos los demás pertenecen a $T\Lambda$. Como
$\upsilon_i$ empareja la lista completa de historias con ella misma, la
reducción módulo $T$ es la identidad en cada generador. Esto da
\eqref{eq:pdf2-bsd-i-plus-tb}; la completitud de $K[[T]]$ hace converger la
inversa. Para $i=0$ la tabla contiene sólo la unidad, de modo que la raíz
queda normalizada antes de evaluar un término principal.
\end{proof}

La equivalencia conserva por separado el ledger
\begin{equation}\label{eq:pdf2-bsd-ledger-diez}
 \mathcal J_{\rm det}=\{\deg,\operatorname{Frob},\operatorname{Ner},
 K,\operatorname{Boc},\operatorname{Sha},\operatorname{Tam},
 \operatorname{tors}^{-2},\Omega,\operatorname{exc}\}.
\end{equation}
En esas diez etiquetas se leen, respectivamente: el peso total; la norma de
Euler; el cambio basado Manin--Néron; la relación finita--singular; todas las
páginas $j_q\bar\beta_q$; el cono de Kummer; los cocientes de componentes;
la red torsional; la integral de Néron; y la orientación excepcional. La
tabla no permite cancelaciones entre etiquetas.

\section{Mapa Poitou--Tate desde Mordell--Weil a la bandera completa}

Se conserva primero la red de Mordell--Weil y sólo después se racionaliza:
\begin{equation}\label{eq:pdf2-bsd-mw-libre}
 M_{E,\mathbb Z}^{\rm MW}:=E(\mathbb Q)/E(\mathbb Q)_{\rm tors},
 \qquad
 M_E^{\rm MW}:=M_{E,\mathbb Z}^{\rm MW}\otimes_{\mathbb Z}K,
 \qquad K=\mathbb Q.
\end{equation}
$\operatorname{Surv}_{{\rm det},\mathbb Z}^{\rm MW}$ denota el submódulo
libre generado por las filas terminales integrales del sumando
Mordell--Weil.
El levantamiento de Kummer con historia completa es
\begin{equation}\label{eq:pdf2-bsd-kummer-hmt}
 \begin{aligned}
 \kappa_{E,\mathbb Z}^{\rm HMT}:M_{E,\mathbb Z}^{\rm MW}
   &\longrightarrow
   Z^1(\mathcal C_{E,\mathbb Z}^{\rm red})
   \cap\operatorname{Surv}_{{\rm det},\mathbb Z}^{\rm MW},\\
 \kappa_E^{\rm HMT}&:=
   \kappa_{E,\mathbb Z}^{\rm HMT}\otimes_{\mathbb Z}\mathbb Q:
   M_E^{\rm MW}\longrightarrow Z^1(\mathcal C_E^{\rm red}).
 \end{aligned}
\end{equation}
y $\operatorname{pg}_q$ publica la coordenada de la página $V_q$. La
extensión $\operatorname{Ext}_{\Gamma,{\rm det}}$, la multisección y la
supervivencia construyen, para cada $q$, un pegado basado
\begin{equation}\label{eq:pdf2-bsd-glue-q}
 \operatorname{gl}_q:V_q\longrightarrow
 \operatorname{Surv}_{\rm det}^{\rm MW}
\end{equation}
que conserva hoja, carry, memoria y terminal. Sus ecuaciones de
unidad--counidad son
\begin{equation}\label{eq:pdf2-bsd-unidad-counidad-paginas}
 \operatorname{pg}_rP_E^{\rm PT}\operatorname{gl}_q
 =\delta_{rq}I_{V_q},
 \qquad
 \sum_{q\ge1}\operatorname{gl}_q
     \operatorname{pg}_qP_E^{\rm PT}=I
 \quad\text{en el sumando MW superviviente}.
\end{equation}
La primera igualdad se comprueba etiqueta por etiqueta: el pegado de nivel
$q$ tiene carry $q$, cero en las otras páginas y el terminal de esa misma
historia. La segunda es la exhaustividad de
$\Sigma_{\rm det}^{\rm mult}$; la suma es finita porque la torre Bockstein
termina.

Para este submódulo póngase
\begin{equation}\label{eq:pdf2-bsd-redes-paginas-saturadas}
 V_{q,\mathbb Z}:=\operatorname{pg}_qP_E^{\rm PT}
   (\operatorname{Surv}_{{\rm det},\mathbb Z}^{\rm MW}),
 \qquad
 A_{q,\mathbb Z}:=V_{q,\mathbb Z}/V_{q+1,\mathbb Z}.
\end{equation}
En las bases terminales, $\operatorname{pg}_q$ y
$\operatorname{gl}_q$ tienen entradas $0,\pm1$ y sus dos composiciones son
las de \eqref{eq:pdf2-bsd-unidad-counidad-paginas}. Por tanto
$V_{q+1,\mathbb Z}$ es un sumando directo basado de $V_{q,\mathbb Z}$,
$A_{q,\mathbb Z}$ es libre y saturado, y
$V_q=V_{q,\mathbb Z}\otimes\mathbb Q$,
$A_q=A_{q,\mathbb Z}\otimes\mathbb Q$.
Sea $s_q:A_{q,\mathbb Z}\to V_{q,\mathbb Z}$ la sección de la base
terminal. Entonces
$V_{k,\mathbb Z}=\bigoplus_{q\geq k}s_q(A_{q,\mathbb Z})$. El desdoblamiento
de jets que conserva las $q$ ocurrencias de una fila de profundidad $q$ es
\begin{equation}\label{eq:pdf2-bsd-desdoblamiento-jets}
 \begin{aligned}
 \operatorname{Jet}_{\rm Boc,\mathbb Z}(E)
   &:=\bigoplus_{q\geq1}\bigoplus_{k=1}^{q}
       A_{q,\mathbb Z}\,e_{q,k},\\
 \mathcal U_{\rm jet}:
 \operatorname{Jet}_{\rm Boc,\mathbb Z}(E)
   &\xrightarrow{\ \sim\ }
     \bigoplus_{k\geq1}V_{k,\mathbb Z},\\
 \mathcal U_{\rm jet}((a_{q,k})_{q,k})
   &:=(v_k)_k,
   \qquad v_k:=\sum_{q\geq k}s_q(a_{q,k}).
 \end{aligned}
\end{equation}
La inversa descompone cada $v_k$ en la suma directa anterior. En particular,
\begin{equation}\label{eq:pdf2-bsd-dimension-jet}
 \operatorname{rank}_{\mathbb Z}\operatorname{Jet}_{\rm Boc,\mathbb Z}(E)
 =\sum_q q\,\operatorname{rank}_{\mathbb Z}A_{q,\mathbb Z}
 =\sum_k\operatorname{rank}_{\mathbb Z}V_{k,\mathbb Z}=d.
\end{equation}
El pegado restringe, en esas mismas bases, a
$\operatorname{gl}_{q,\mathbb Z}:V_{q,\mathbb Z}\to
\operatorname{Surv}_{{\rm det},\mathbb Z}^{\rm MW}$.

Se definen, sin usar dimensiones,
\begin{equation}\label{eq:pdf2-bsd-psi-y-lambda-q}
 \begin{aligned}
  \psi_{q,\mathbb Z}&:=
      \operatorname{pg}_qP_E^{\rm PT}\kappa_{E,\mathbb Z}^{\rm HMT}:
      M_{E,\mathbb Z}^{\rm MW}\to V_{q,\mathbb Z},\\
  \lambda_{q,\mathbb Z}&:=
      \operatorname{pr}_{\rm MW}P_E^{\rm Man}
      \operatorname{gl}_{q,\mathbb Z}:
      V_{q,\mathbb Z}\to M_{E,\mathbb Z}^{\rm MW},\\
  \psi_q&:=\operatorname{pg}_qP_E^{\rm PT}\kappa_E^{\rm HMT}:
      M_E^{\rm MW}\to V_q,\\
  \lambda_q&:=\operatorname{pr}_{\rm MW}P_E^{\rm Man}
      \operatorname{gl}_q:V_q\to M_E^{\rm MW}.
 \end{aligned}
\end{equation}
Los dominios intermedios quedan ligados por las dos compatibilidades basadas
que forman parte del pegado Poitou--Tate y que aquí se escriben sobre su
acción alcanzable:
\begin{equation}\label{eq:pdf2-bsd-kummer-glue-compatibilidad}
 \kappa_E^{\rm HMT}\operatorname{pr}_{\rm MW}P_E^{\rm Man}
       \operatorname{gl}_q=\operatorname{gl}_q,
 \qquad
 \operatorname{pr}_{\rm MW}P_E^{\rm Man}\kappa_E^{\rm HMT}
       =I_{M_E^{\rm MW}}.
\end{equation}
Las mismas dos identidades, con subíndice $\mathbb Z$, valen sobre
$M_{E,\mathbb Z}^{\rm MW}$ y $V_{q,\mathbb Z}$, porque todas las matrices
anteriores son integrales y basadas.

\begin{theorem}[Isomorfismo basado de la bandera Poitou--Tate]
\label{thm:pdf2-bsd-psi-pt}
Los mapas
\begin{equation}\label{eq:pdf2-bsd-psi-pt}
 \begin{aligned}
 \Psi_{\rm PT}:M_E^{\rm MW}&\longrightarrow
       \operatorname{Flag}_{\rm Boc}(E):=\bigoplus_{q\ge1}V_q,
 &P&\longmapsto(\psi_q(P))_{q\ge1},\\
 \Psi_{\rm PT}^{-1}:\operatorname{Flag}_{\rm Boc}(E)&\longrightarrow M_E^{\rm MW},
 &(v_q)_q&\longmapsto\sum_{q\ge1}\lambda_q(v_q)
 \end{aligned}
\end{equation}
son inversos. En consecuencia,
\begin{equation}\label{eq:pdf2-bsd-psi-pt-integral}
 \Psi_{{\rm PT},\mathbb Z}:M_{E,\mathbb Z}^{\rm MW}
 \xrightarrow{\ \sim\ }\bigoplus_{q\ge1}V_{q,\mathbb Z},
 \qquad
 \Psi_{{\rm PT},\mathbb Z}^{-1}((v_q)_q)=
 \sum_q\lambda_{q,\mathbb Z}(v_q),
\end{equation}
y, sólo después de extender escala,
\begin{equation}\label{eq:pdf2-bsd-rango-orden}
 \operatorname{rank}E(\mathbb Q)
 =\sum_{q\ge1}\dim_K V_q
 =d=\operatorname{ord}_T\det S_E(T).
\end{equation}
\end{theorem}

\begin{proof}
Aplicar \eqref{eq:pdf2-bsd-unidad-counidad-paginas} junto con
\eqref{eq:pdf2-bsd-kummer-glue-compatibilidad} da, con todos los dominios
intermedios visibles,
$\psi_r\lambda_q=\delta_{rq}I$ y
$\sum_q\lambda_q\psi_q=I_{M_E^{\rm MW}}$. Por tanto las dos fórmulas de
\eqref{eq:pdf2-bsd-psi-pt} son inversas antes de comparar dimensiones.
Las tablas terminales son integrales y basadas; el mismo cálculo, sin
denominadores, prueba \eqref{eq:pdf2-bsd-psi-pt-integral}. Así no aparece un
índice de subred al elegir después una base de Mordell--Weil.
Sólo después se toman dimensiones y se usa
\eqref{eq:pdf2-bsd-orden-smith}. El emparejamiento
$j_q\bar\beta_q$ transporta las bases saturadas a las filas de altura cuya
comparación con Néron--Tate se escribe en
\eqref{eq:pdf2-bsd-altura-nt-generadores}.
\end{proof}

\section{Localización, rango de Smith y exactitud Sha--Tamagawa}

En esta sección $v$ recorre los lugares finitos; en los de buena reducción se
tiene $\Phi_v(\mathbb F_v)=0$ y $c_v=1$, de modo que las sumas y productos
locales tienen soporte finito. El lugar infinito se reserva para la evaluación
de la diferencial de Néron en la cuarta flecha. Para cada lugar finito $v$ se
tiene el triángulo
\[
 \mathcal C_{E,v}^{\rm fin}\to\mathcal C_{E,v}\to
 \mathcal C_{E,v}^{\rm sing}\to\mathcal C_{E,v}^{\rm fin}[1]
\]
y el mapa de cadenas
\begin{equation}\label{eq:pdf2-bsd-ell-e}
 \ell_E(x,(y_v)_v)=(\operatorname{res}_vx-i_vy_v)_v.
\end{equation}
El cono reducido es
\begin{equation}\label{eq:pdf2-bsd-cono-reducido}
 \mathcal C_E^{\rm loc,red}:=
 q_{\rm root}\operatorname{Cone}(\ell_E)[-1].
\end{equation}
En grados de cociclos, un elemento se escribe $(x,(y_v),(h_v))$ y
\begin{equation}\label{eq:pdf2-bsd-diferencial-cono}
 d(x,(y_v),(h_v))=
 (dx,(dy_v),(\operatorname{res}_vx-i_vy_v-dh_v)_v).
\end{equation}

Tras separar por \eqref{eq:pdf2-bsd-psi-pt} la red libre de Mordell--Weil
y su dual, los bloques FS están contraídos por
\eqref{eq:pdf2-bsd-hfs-global}. En el retículo restante, el inverso
\eqref{eq:pdf2-bsd-j-e-inverso-racional} y sus dos comprobaciones
\eqref{eq:pdf2-bsd-j-e-dos-lados} dan, sin inferir aciclicidad sólo desde la
dualidad,
\begin{equation}\label{eq:pdf2-bsd-rango-completo}
 D_E\otimes\mathbb Q:M_E\otimes\mathbb Q
 \xrightarrow{\sim}N_E\otimes\mathbb Q.
\end{equation}
La forma normal integral da
\begin{equation}\label{eq:pdf2-bsd-smith-finito}
 UD_EV=\operatorname{diag}(s_1,\ldots,s_t),
 \qquad
 F_E:=\operatorname{coker}D_E,
 \qquad |F_E|=\prod_{i=1}^t|s_i|<\infty.
\end{equation}

Sea
\begin{equation}\label{eq:pdf2-bsd-pi-componentes}
 \pi_\Phi:N_E\longrightarrow\bigoplus_v\Phi_v(\mathbb F_v)
\end{equation}
la reducción de cada coordenada local módulo
$\mathcal C_{E,v}^{\rm fin}+d\mathcal C_{E,v}$; se tiene
$\pi_\Phi D_E=0$. Para una clase
$[\xi]\in\operatorname{Sha}(E)$, elíjanse sus primitivas finitas
$y_v$ y homotopías $h_v$ con
$\operatorname{res}_v\xi-i_vy_v=dh_v$. La fórmula
\begin{equation}\label{eq:pdf2-bsd-iota-sha-cociclos}
 \begin{aligned}
 n_\xi&:=\operatorname{pr}_{N_E}\operatorname{pr}_{\mathcal L_E}
       \Theta_E(\xi,(y_v),(h_v)),\\
 \iota_{\rm Sha}[\xi]&:=[n_\xi]\in F_E
 \end{aligned}
\end{equation}
es independiente de las elecciones por
\eqref{eq:pdf2-bsd-diferencial-cono}.

Para hacer material la sobreyectividad final, sea
$\widetilde\phi_v$ el representante terminal basado de
$\phi_v\in\Phi_v(\mathbb F_v)$. La suma de counidades se corrige en la
unidad común y se levanta por el operador estructural:
\begin{equation}\label{eq:pdf2-bsd-lift-componentes}
 \operatorname{Lift}_\Phi((\phi_v)_v):=
 \operatorname{pr}_{N_E}\operatorname{Ext}_{\Gamma,{\rm det}}
 \left((\widetilde\phi_v)_v,
 -\widehat\Pi_0\!\left(\sum_v
       \varepsilon_{\rm AW}(\widetilde\phi_v)\right)\right).
\end{equation}
La relación de corestricción de
$\operatorname{Rel}_{\rm det}$ hace de esta expresión un cociclo y da
\begin{equation}\label{eq:pdf2-bsd-lift-dos-identidades}
 \pi_\Phi\operatorname{Lift}_\Phi=I,
 \qquad
 \operatorname{im}D_E=
 \ker\pi_\Phi\cap\ker(\text{clase global}).
\end{equation}

\begin{theorem}[Sucesión exacta finita Sha--Tamagawa]
\label{thm:pdf2-bsd-sha-tamagawa}
Los mapas de cociclos anteriores inducen
\begin{equation}\label{eq:pdf2-bsd-sha-tamagawa}
 0\longrightarrow\operatorname{Sha}(E)
 \xrightarrow{\ \iota_{\rm Sha}\ }F_E
 \xrightarrow{\ \partial_{\rm loc}\ }
 \bigoplus_v\Phi_v(\mathbb F_v)
 \longrightarrow0,
 \qquad
 \partial_{\rm loc}[n]:=\pi_\Phi(n).
\end{equation}
En particular, $\operatorname{Sha}(E)$ es finito y
\begin{equation}\label{eq:pdf2-bsd-cardinal-sts}
 |F_E|=|\operatorname{Sha}(E)|\prod_vc_v,
 \qquad c_v:=|\Phi_v(\mathbb F_v)|.
\end{equation}
\end{theorem}

\begin{proof}
La igualdad $\pi_\Phi D_E=0$ hace bien definida a
$\partial_{\rm loc}$. Si \eqref{eq:pdf2-bsd-iota-sha-cociclos} es una
frontera, su componente global es una frontera; así $\iota_{\rm Sha}$ es
inyectiva. Una clase de $F_E$ está en $\ker\partial_{\rm loc}$ si y sólo
si todas sus coordenadas singulares admiten representantes finitos; por
\eqref{eq:pdf2-bsd-diferencial-cono}, esto equivale a estar en la imagen de
$\iota_{\rm Sha}$. Finalmente,
\eqref{eq:pdf2-bsd-lift-dos-identidades} da una preimagen de cada tupla de
componentes y prueba la sobreyectividad. Los términos adyacentes libres de
la sucesión larga son precisamente los dos sumandos separados por
$\Psi_{\rm PT}$ y $X^{\rm orth}$; los términos adyacentes FS son
contractibles por $H_{\rm FS}$. No se ha supuesto su anulación. La finitud
de $F_E$ procede de \eqref{eq:pdf2-bsd-rango-completo}; la inyección da la
 finitud de $\operatorname{Sha}(E)$ y tomar órdenes da
 \eqref{eq:pdf2-bsd-cardinal-sts}.
\end{proof}

\begin{corollary}[Descarga primaria completa y recomposición integral]
\label{cor:pdf2-bsd-sha-tamagawa-todos-ell}
Para cada primo $\ell$, el cambio de base de los mapas de cociclos produce
la sucesión exacta
\begin{equation}\label{eq:pdf2-bsd-sha-tamagawa-ell}
 0\longrightarrow\operatorname{Sha}(E)[\ell^\infty]
 \xrightarrow{\ \iota_{{\rm Sha},\ell}\ }F_{E,\ell}
 \xrightarrow{\ \partial_{{\rm loc},\ell}\ }
 \bigoplus_v\Phi_v(\mathbb F_v)[\ell^\infty]
 \longrightarrow0,
\end{equation}
donde
\begin{equation}\label{eq:pdf2-bsd-mapas-ell-generadores}
 \iota_{{\rm Sha},\ell}([\xi]\otimes1)
   =[n_\xi\otimes1],
 \qquad
 \partial_{{\rm loc},\ell}([n\otimes1])
   =(\pi_{\Phi,v}(n)\otimes1)_v.
\end{equation}
Además, la suma de todas las sucesiones
\eqref{eq:pdf2-bsd-sha-tamagawa-ell} es exactamente
\eqref{eq:pdf2-bsd-sha-tamagawa}; no es sólo su parte $3$-primaria.
\end{corollary}

\begin{proof}
La matriz $D_E$, el cociclo $n_\xi$, la reducción $\pi_\Phi$ y
$\operatorname{Lift}_\Phi$ son integrales por
\eqref{eq:pdf2-bsd-complejos-base-change}--
\eqref{eq:pdf2-bsd-theta-base-change}. Como $\mathbb Z_\ell$ es plana sobre
$\mathbb Z$, tensorizar la sucesión exacta finita
\eqref{eq:pdf2-bsd-sha-tamagawa} conserva exactitud. Para un grupo finito
$A$ se tiene canónicamente
$A\otimes\mathbb Z_\ell=A[\ell^\infty]$; esto identifica los tres términos
y da las fórmulas sobre generadores
\eqref{eq:pdf2-bsd-mapas-ell-generadores}. La forma de Smith
\eqref{eq:pdf2-bsd-smith-finito} escribe
$s_i=\prod_\ell\ell^{v_\ell(s_i)}$, de modo que
\[
 F_E\cong\bigoplus_\ell F_{E,\ell},\qquad
 \operatorname{Sha}(E)\cong
   \bigoplus_\ell\operatorname{Sha}(E)[\ell^\infty],\qquad
 \Phi_v(\mathbb F_v)\cong
   \bigoplus_\ell\Phi_v(\mathbb F_v)[\ell^\infty].
\]
Sólo intervienen finitos primos y finitos lugares porque todos los grupos
son cocientes del grupo finito $F_E$. Sumar las sucesiones primarias recupera
la integral. Tomando órdenes primero primo a primo y luego multiplicando,
\begin{equation}\label{eq:pdf2-bsd-cardinal-sts-todos-ell}
 \prod_\ell|F_{E,\ell}|
 =\prod_\ell|\operatorname{Sha}(E)[\ell^\infty]|
   \prod_v\prod_\ell|\Phi_v(\mathbb F_v)[\ell^\infty]|
 =|\operatorname{Sha}(E)|\prod_vc_v.
\end{equation}
Así quedan descargados materialmente todos los coeficientes y todos los
factores de Tamagawa.
\end{proof}

\section{Comparador analítico--aritmético basado}

La coordenada analítica no se postula desde el orden buscado. Sea
$\lambda_3=\log4$, correspondiente al generador ciclotómico basado
$\gamma_3=4$, y defínase en un disco $U_E$ centrado en $1$
\begin{equation}\label{eq:pdf2-bsd-t-e-construido}
 T_E(s):=\frac{\exp(\lambda_3(s-1))-1}{\lambda_3}
 =(s-1)+\frac{\lambda_3}{2}(s-1)^2+O((s-1)^3).
\end{equation}
Por cálculo directo,
\begin{equation}\label{eq:pdf2-bsd-t-e-normalizacion-probada}
 T_E(1)=0,
 \qquad T_E'(s)=\exp(\lambda_3(s-1)),
 \qquad T_E'(1)=1.
\end{equation}
El teorema de la función inversa proporciona, después de reducir $U_E$ si
fuera necesario, una carta biholomorfa $s\leftrightarrow T_E(s)$.
Para comparar todas las fibras, tómese
$\gamma_\ell=1+\ell$ si $\ell$ es impar y $\gamma_2=5$, y póngase
\begin{equation}\label{eq:pdf2-bsd-t-ell-construido}
 T_{E,\ell}(s):=
 \frac{\exp(\log(\gamma_\ell)(s-1))-1}{\log(\gamma_\ell)}.
\end{equation}
Las cartas se relacionan por un único automorfismo formal
$T_{E,\ell}=\vartheta_{\ell,3}(T_E)$ con
$\vartheta_{\ell,3}(0)=0$ y $\vartheta_{\ell,3}'(0)=1$. Por tanto ni el
orden en $T$ ni el coeficiente principal cambian al pasar de $3$ a $\ell$.

La continuación analítica que llega a esta carta tampoco se declara. Por el
teorema de modularidad, sea
\[
 f_E(z)=\sum_{m\geq1}a_m(E)e^{2\pi imz}
\]
la forma nueva normalizada de peso $2$ y nivel $N_E$ asociada a $E$, y sea
$\varepsilon_E\in\{\pm1\}$ el signo fijado por
\begin{equation}\label{eq:pdf2-bsd-transformacion-modular}
 g_E(y):=f_E\!\left(\frac{iy}{\sqrt{N_E}}\right),
 \qquad
 g_E(1/y)=\varepsilon_Ey^2g_E(y).
\end{equation}
Para $\Re s>3/2$, integración término a término da
\begin{equation}\label{eq:pdf2-bsd-mellin-semiplano}
 \Lambda_E(s):=N_E^{s/2}(2\pi)^{-s}\Gamma(s)L(E,s)
 =\int_0^\infty g_E(y)y^s\,\frac{dy}{y}.
\end{equation}
Se divide la integral en $(0,1)$ y $(1,\infty)$ y, en la primera parte, se
efectúa el cambio $y=1/t$. La segunda igualdad de
\eqref{eq:pdf2-bsd-transformacion-modular} produce la fórmula material
\begin{equation}\label{eq:pdf2-bsd-mellin-continuada}
 \Lambda_E^{\rm cont}(s)=
 \int_1^\infty g_E(y)
       \bigl(y^s+\varepsilon_Ey^{2-s}\bigr)\,\frac{dy}{y}.
\end{equation}
La función $g_E(y)$ decrece exponencialmente cuando $y\to\infty$; por ello la
integral de \eqref{eq:pdf2-bsd-mellin-continuada}, y todas sus derivadas en
$s$, convergen uniformemente en compactos. Así define una función entera y
coincide con \eqref{eq:pdf2-bsd-mellin-semiplano} donde la serie de Dirichlet
converge. Puesto que
$N_E^{s/2}(2\pi)^{-s}\Gamma(s)$ es holomorfa y no nula cerca de $s=1$, la
división por ese factor construye, sin usar su orden de anulación, el germen
holomorfo de $L(E,s)$ en un disco $U_E$ centrado en $1$.

Para evitar mezclar una fibra $\ell$-ádica con una matriz compleja, se toma
primero el polinomio entero, independiente de $\ell$,
\begin{equation}\label{eq:pdf2-bsd-polinomio-local-entero}
 P_{E,p}(X):=\det(I-X\operatorname{Frob}_p\mid V_\ell E^{I_p})
 \in\mathbb Z[X].
\end{equation}
Para $p\nmid N_E$, sean $\alpha_p,\beta_p$ las raíces de $P_{E,p}$. Se toma
$W_p=\mathbb C^2$ con base ortonormal y la realización normal
\begin{equation}\label{eq:pdf2-bsd-realizacion-normal-frobenius}
 F_p^{\mathbb C}:=\operatorname{diag}(\alpha_p,\beta_p),
 \qquad |\alpha_p|=|\beta_p|=p^{1/2},
 \qquad \|F_p^{\mathbb C}\|_1=2p^{1/2}.
\end{equation}
Para los primos malos se elige de igual modo una realización compleja del
polinomio local; son finitos y se conservan como bloques separados. No se usa
la matriz compañera, cuya norma no estaría controlada uniformemente por el
módulo de sus autovalores.

Póngase $\mathscr H_E:=\widehat\bigoplus_{p\nmid N_E}W_p$ y, para $X>0$,
\begin{equation}\label{eq:pdf2-bsd-truncados-fredholm}
 K_{E,X}(s):=\bigoplus_{\substack{p\leq X\\p\nmid N_E}}
          p^{-s}F_p^{\mathbb C},
 \qquad
 L_X(E,s):=\det(I-K_{E,X}(s))^{-1}
       \prod_{p\mid N_E}P_{E,p}(p^{-s})^{-1}.
\end{equation}
Si $K_0$ es un compacto de $\{\Re s>3/2\}$ y
$\sigma_0:=\inf_{s\in K_0}\Re s$, entonces
\begin{equation}\label{eq:pdf2-bsd-cota-traza}
 \sup_{s\in K_0}\|K_{E,X'}(s)-K_{E,X}(s)\|_1
 \leq C_E\sum_{X<p\leq X'}p^{1/2-\sigma_0}
 \xrightarrow[X,X'\to\infty]{}0.
\end{equation}
Así $K_{E,X}\to K_E$ en norma traza, localmente uniformemente, y la
continuidad del determinante de Fredholm prueba, no define,
\begin{equation}\label{eq:pdf2-bsd-fredholm-euler}
 \det_F(I-K_E(s))^{-1}
 \prod_{p\mid N_E}P_{E,p}(p^{-s})^{-1}
 =\lim_{X\to\infty}L_X(E,s)=L(E,s).
\end{equation}

Para pasar del producto infinito a un complejo finito se conserva primero el
operador, no sólo su determinante. Sean $\mathscr H_E^+$ y
$\mathscr H_E^-$ dos copias basadas de
$\mathscr H_E\oplus\bigoplus_{p\mid N_E}W_p$. En el semiplano
$\Re s>3/2$ se define
\begin{equation}\label{eq:pdf2-bsd-operadores-euler-inversos}
 \begin{aligned}
  \mathscr A_{E,X}(s)&:=
   (I-K_{E,X}(s))^{-1}\oplus
   \bigoplus_{p\mid N_E}(I-p^{-s}F_p^{\mathbb C})^{-1},\\
  \mathscr A_E(s)&:=\lim_{X\to\infty}\mathscr A_{E,X}(s):
       \mathscr H_E^+\longrightarrow\mathscr H_E^-.
 \end{aligned}
\end{equation}
La identidad del resolvente y \eqref{eq:pdf2-bsd-cota-traza} dan, para cada
compacto $K_0$ del semiplano,
\begin{equation}\label{eq:pdf2-bsd-resolvente-convergencia-traza}
 \sup_{s\in K_0}
 \|\mathscr A_{E,X'}(s)-\mathscr A_{E,X}(s)\|_1\longrightarrow0.
\end{equation}
Sea $\mathcal U_{\rm bas}:\mathscr H_E^+\to\mathscr H_E^-$ la identificación unitaria que
envía cada vector de la base positiva al vector homónimo de la base
negativa. Los determinantes siguientes significan los de
$\mathcal U_{\rm bas}^{-1}\mathscr A_{E,X}=I+\mathcal L^1$ y
$\mathcal U_{\rm bas}^{-1}\mathscr A_E=I+\mathcal L^1$. Entonces
\begin{equation}\label{eq:pdf2-bsd-det-a-x-l-x}
 \det_F(\mathcal U_{\rm bas}^{-1}\mathscr A_{E,X}(s))=L_X(E,s),
 \qquad
 \det_F(\mathcal U_{\rm bas}^{-1}\mathscr A_E(s))=L(E,s).
\end{equation}
La segunda igualdad es precisamente el límite de
\eqref{eq:pdf2-bsd-fredholm-euler}; el signo es el correcto porque se toma el
inverso de $I-K_E$.

La corestricción finita se construye a partir de los terminales, sin publicar
infinitos primos en una matriz finita. Sean $\mathcal I_E^+$ y
$\mathcal I_E^-$ las listas terminales de las bases integrales de $M_E$ y
$N_E$, respectivamente; son finitas por
\eqref{eq:pdf2-bsd-sumando-fs-y-lattice}. Fijado $\sigma_*>3/2$, el
entrelazador local--global las inserta en el espacio de Euler por
\begin{equation}\label{eq:pdf2-bsd-embedding-terminal-euler}
 \iota_{\rm Eul}^\pm(e_\eta):=
 \left(p^{-\sigma_*}\operatorname{res}_p^\pm
       \operatorname{Surv}_E(e_\eta)\right)_p
 \oplus\left(\operatorname{res}_p^\pm
       \operatorname{Surv}_E(e_\eta)\right)_{p\mid N_E}.
\end{equation}
La suma buena es cuadrado--sumable por la misma cota que
\eqref{eq:pdf2-bsd-cota-traza}. Se ponen
$\tau_\eta^\pm:=\iota_{\rm Eul}^\pm(e_\eta)$; la independencia de las filas
terminales da las counidades duales $\{\epsilon_{\eta}^{\pm}\}$, de modo que
$\epsilon_\eta^\pm(\tau_{\eta'}^\pm)=\delta_{\eta\eta'}$. Las sumas
\begin{equation}\label{eq:pdf2-bsd-proyecciones-terminales-finitas}
 P_E^\pm:=\sum_{\eta\in\mathcal I_E^\pm}
      \tau_\eta^\pm\epsilon_\eta^\pm,
 \qquad Q_E^\pm:=I-P_E^\pm
\end{equation}
son proyecciones basadas de rango finito; sus índices se fijan por las
etiquetas terminales y no por el rango de Mordell--Weil ni por un orden de
anulación.

La coordenada $\deg$ del ledger descompone
\begin{equation}\label{eq:pdf2-bsd-graduacion-fredholm}
 \mathscr H_E^\pm=\widehat\bigoplus_i\mathscr H_E^{\pm,i},
 \qquad
 \mathcal I_E^{\pm,i}:=\{\eta\in\mathcal I_E^\pm:\deg\eta=i\},
 \qquad
 P_E^{\pm,i}:=
 \sum_{\eta\in\mathcal I_E^{\pm,i}}\tau_\eta^\pm\epsilon_\eta^\pm.
\end{equation}
Las relaciones locales conservan $\deg$ salvo el diferencial $i\mapsto i+1$;
por tanto
$\mathscr A_E=\bigoplus_i\mathscr A_E^i$ con
$\mathscr A_E^i:\mathscr H_E^{+,i}\to\mathscr H_E^{-,i+1}$. Los cuatro
bloques de \eqref{eq:pdf2-bsd-bloques-fredholm}, su Schur y sus truncaciones
se toman grado a grado, y por definición
\begin{equation}\label{eq:pdf2-bsd-schur-graduado}
 \mathscr D_E=\bigoplus_i\mathscr D_E^i,
 \qquad
 \mathscr D_E^i:=A_{00}^i-A_{01}^i(A_{11}^i)^{-1}A_{10}^i.
\end{equation}

La continuación Fredholm se escribe antes de usarla. En la carta de Euler
$U_0$ se pone $\mathscr A_{E,0}:=\mathscr A_E$. Las cartas conservan un
complemento basado común $\mathscr Q_E^\pm$ y cambian sólo la familia finita
de terminales. Si
$\tau_{\eta,\alpha}^\pm:=
\operatorname{Ext}_{\Gamma,{\rm det},\alpha}(\tau_\eta^\pm)$ y
$\epsilon_{\eta,\alpha}^\pm$ es su base dual, se define
\begin{equation}\label{eq:pdf2-bsd-cambios-fredholm-cartas}
 C_{\alpha\beta}^\pm:=
 I_{\mathscr Q_E^\pm}\oplus
 \sum_{\eta\in\mathcal I_E^\pm}
 \tau_{\eta,\beta}^\pm\epsilon_{\eta,\alpha}^\pm,
 \qquad C_{\alpha\beta}^\pm-I\in\mathcal L^1,
 \qquad
 C_{\beta\gamma}^\pm C_{\alpha\beta}^\pm=C_{\alpha\gamma}^\pm.
\end{equation}
La diferencia respecto de la identidad tiene rango a lo sumo
$|\mathcal I_E^\pm|$ y es, por tanto, de clase traza; la holomorfía se sigue
de que cada una de las finitas coordenadas terminales de
$\operatorname{Ext}_{\Gamma,{\rm det},\alpha}$ es holomorfa. La última
igualdad se verifica en $\mathscr Q_E^\pm$ y en cada
$\tau_{\eta,\alpha}^\pm$. Se define
recursivamente en los solapamientos
\begin{equation}\label{eq:pdf2-bsd-continuacion-fredholm-cartas}
 \mathscr A_{E,\beta}:=
 C_{\alpha\beta}^-\mathscr A_{E,\alpha}
 (C_{\alpha\beta}^+)^{-1}.
\end{equation}
Los operadores $C_{\alpha\beta}^\pm$ son holomorfos porque cada coeficiente
es una coordenada de $\operatorname{Ext}_{\Gamma,{\rm det}}$; el cociclo de
\eqref{eq:pdf2-bsd-cambios-fredholm-cartas} prueba que
\eqref{eq:pdf2-bsd-continuacion-fredholm-cartas} no depende de la cadena de
cartas. Ésta es la familia determinante--clase que se denota en adelante por
$\mathscr A_E(s)$. Con
$C_{0\alpha}^\pm:=C_{\alpha-1,\alpha}^\pm\cdots C_{0,1}^\pm$ se transportan
además
$P_{E,\alpha}^{\pm,i}:=
C_{0\alpha}^\pm P_E^{\pm,i}(C_{0\alpha}^\pm)^{-1}$ y
$\mathscr H_{E,\alpha}^{\pm,i}:=
C_{0\alpha}^\pm\mathscr H_E^{\pm,i}$; de este modo todos los Schur
$\mathscr D_{E,\alpha}^i$ son los de
\eqref{eq:pdf2-bsd-schur-graduado}, no símbolos nuevos. Tras reducir cada
carta, el bloque
\begin{equation}\label{eq:pdf2-bsd-bloques-fredholm}
 A_{ab}(s):=P_a^-\mathscr A_E(s)P_b^+,
 \qquad P_0^\pm:=P_E^\pm,\quad P_1^\pm:=Q_E^\pm,
 \qquad A_{11}(s):Q_E^+\mathscr H_E^+\xrightarrow{\sim}
                       Q_E^-\mathscr H_E^-
\end{equation}
es invertible. Ésta es la reducción finita estándar del operador Fredholm;
la elección concreta de $P_E^\pm$ viene dada por
\eqref{eq:pdf2-bsd-proyecciones-terminales-finitas}. La base fija también
la identificación unitaria
$U_{11}:Q_E^+\mathscr H_E^+\to Q_E^-\mathscr H_E^-$; por construcción
$U_{11}^{-1}A_{11}=I+\mathcal L^1$, y sólo a este operador se aplica
$\det_F$. Las copias $+$ y $-$ ya son la totalización par--impar; no se
aplica un segundo signo alternado al determinante del operador inverso.

Para cada truncación que contiene el soporte de esos terminales se definen
materialmente
\begin{equation}\label{eq:pdf2-bsd-cor-lift-truncados}
 \begin{aligned}
  \operatorname{Cor}_{E,X}&:=P_E^-,\\
  \operatorname{Lift}_{E,X}(v)&:=
       v-A_{11,X}^{-1}A_{10,X}v,\\
  \mathscr D_{E,X}(s)&:=
       \operatorname{Cor}_{E,X}\mathscr A_{E,X}(s)
       \operatorname{Lift}_{E,X}
       =A_{00,X}-A_{01,X}A_{11,X}^{-1}A_{10,X}.
 \end{aligned}
\end{equation}
Por multiplicación de bloques,
\begin{equation}\label{eq:pdf2-bsd-eliminacion-schur}
 \begin{pmatrix}I&-A_{01,X}A_{11,X}^{-1}\\0&I\end{pmatrix}
 \mathscr A_{E,X}
 \begin{pmatrix}I&0\\-A_{11,X}^{-1}A_{10,X}&I\end{pmatrix}
 =\begin{pmatrix}\mathscr D_{E,X}&0\\0&A_{11,X}\end{pmatrix}.
\end{equation}
Así la corestricción no es nominal: su lift, su inversa sobre el complemento
y su diferencial finito están escritos. La convergencia del resolvente
implica
\begin{equation}\label{eq:pdf2-bsd-corestriccion-limite}
 \mathscr D_E(s):=\lim_{X\to\infty}\mathscr D_{E,X}(s)
 =A_{00}-A_{01}A_{11}^{-1}A_{10},
 \qquad
 \|\mathscr D_{E,X}-\mathscr D_E\|_1\to0.
\end{equation}
El determinante del complemento no se descarta: por
\eqref{eq:pdf2-bsd-eliminacion-schur},
\begin{equation}\label{eq:pdf2-bsd-determinante-schur-complemento}
 \det_F(\mathcal U_{\rm bas}^{-1}\mathscr A_E)
 =\det_F(U_{11}^{-1}A_{11})\det(\mathscr D_E).
\end{equation}

Sea $\mathcal O(U_E)$ el anillo de funciones holomorfas del disco central.
Los espacios racionales basados $V_0^i,W_0^{i+1}$ de
\eqref{eq:pdf2-bsd-s-universal} proporcionan
\begin{equation}\label{eq:pdf2-bsd-dominios-analiticos}
 \mathscr V_E^i:=V_0^i\otimes_{\mathbb Q}\mathcal O(U_E),
 \qquad
 \mathscr W_E^{i+1}:=W_0^{i+1}\otimes_{\mathbb Q}\mathcal O(U_E).
\end{equation}
Las filas de $\operatorname{Rel}_{\rm det}$ construyen, sobre cada terminal
$\eta$ y en cada grado, los vectores del modelo común
$v_{\eta,i}:=\upsilon_i(b_\eta)$,
$w_{\eta,i+1}:=\upsilon_{i+1}(b_\eta^{\rm term})$ y los isomorfismos basados
\begin{equation}\label{eq:pdf2-bsd-jvw-generadores}
 \begin{aligned}
  J_{E,V,i}(\tau_{\eta,i}^+)&:=v_{\eta,i},\\
  J_{E,W,i+1}(\tau_{\eta,i+1}^-)&:=w_{\eta,i+1}.
 \end{aligned}
\end{equation}
La relación diferencial de esa misma fila, aplicada a cada generador, es
\begin{equation}\label{eq:pdf2-bsd-entrelazamiento-analitico-smith}
 \boxed{
 J_{E,W,i+1}(s)\,\mathscr D_E^i(s)
 =S_E^i(T_E(s))\,J_{E,V,i}(s).}
\end{equation}
Éste es el entrelazador efectivo entre el Schur Fredholm corestringido y la
matriz de Smith; no es una igualdad de determinantes postulada.

Para tipar la continuación, sea $\mathfrak U_E$ un dominio conexo que contiene
un abierto del semiplano $\Re s>3/2$ y el disco $U_E$, y elíjase una cadena
finita de cartas $\{U_\alpha\}_{\alpha=0}^A$ de
$\operatorname{Ext}_{\Gamma,{\rm det}}$, con $U_A=U_E$. En $U_\alpha$ se
ponen
$\mathscr P_{E,\alpha}^i:=
P_{E,\alpha}^{+,i}\mathscr H_{E,\alpha}^{+,i}$ y
$\mathscr D_{E,\alpha}^i$ igual al Schur de
\eqref{eq:pdf2-bsd-corestriccion-limite}. Los mapas al modelo común se
definen, sin elección adicional, por
\begin{equation}\label{eq:pdf2-bsd-j-cartas-modelo-comun}
 J_{E,\alpha,i}(\tau_{\eta,\alpha,i}):=v_{\eta,i}
 \quad\text{en el dominio},
 \qquad
 J_{E,\alpha,i+1}(\tau_{\eta,\alpha,i+1}):=w_{\eta,i+1}
 \quad\text{en el codominio}.
\end{equation}
Son holomorfos porque las bases de
\eqref{eq:pdf2-bsd-cambios-fredholm-cartas} lo son. En un solapamiento se
define, siempre en el mismo grado,
\begin{equation}\label{eq:pdf2-bsd-transiciones-atlas}
 G_{\alpha\beta,i}:=J_{E,\beta,i}^{-1}J_{E,\alpha,i}:
 \mathscr P_{E,\alpha}^i\big|_{U_\alpha\cap U_\beta}
 \longrightarrow
 \mathscr P_{E,\beta}^i\big|_{U_\alpha\cap U_\beta}.
\end{equation}
La composición literal de estas matrices prueba
\begin{equation}\label{eq:pdf2-bsd-cociclo-atlas}
 G_{\beta\gamma,i}G_{\alpha\beta,i}=G_{\alpha\gamma,i},
 \qquad G_{\alpha\alpha,i}=I,
 \qquad G_{\beta\alpha,i}=G_{\alpha\beta,i}^{-1},
\end{equation}
y el entrelazamiento en las dos cartas da la compatibilidad de cadena
\begin{equation}\label{eq:pdf2-bsd-cociclo-cadena}
 G_{\alpha\beta,i+1}\mathscr D_{E,\alpha}^i
 =\mathscr D_{E,\beta}^iG_{\alpha\beta,i}.
\end{equation}
Por tanto
$g_{\alpha\beta}:=\prod_i\det(G_{\alpha\beta,i})^{(-1)^i}$ satisface el
cociclo triple. Si
\begin{equation}\label{eq:pdf2-bsd-tau-atlas}
 \tau_E^{\rm atlas}(s):=
 \prod_{\alpha=0}^{A-1}g_{\alpha,\alpha+1}(s),
\end{equation}
entonces $\tau_E^{\rm atlas}$ es exactamente el transporte determinantal
desde la carta de Euler a la carta central.

Se denotan por $J_{E,V}$ y $J_{E,W}$ las totalizaciones par--impar, con
orden de Koszul, de los mapas de
\eqref{eq:pdf2-bsd-entrelazamiento-analitico-smith}; sus determinantes son
los productos alternados de los determinantes por grado. En la carta central
se define, y no antes, la unidad conjunta
\begin{equation}\label{eq:pdf2-bsd-u-e-construida}
 u_E(s):=\tau_E^{\rm atlas}(s)
 \det_F(U_{11}^{-1}A_{11}(s))
 \frac{\det J_{E,V}(s)}
      {\det J_{E,W}(s)}.
\end{equation}
Todos sus factores son holomorfos e invertibles al reducir $U_E$. La
normalización no se impone factor a factor. Sea $\mathbf e_{\rm cent}$ el
generador de la línea determinante del modelo común. Al efectuar las dos
operaciones triangulares de \eqref{eq:pdf2-bsd-eliminacion-schur}, transportar
por \eqref{eq:pdf2-bsd-cambios-fredholm-cartas} y aplicar
\eqref{eq:pdf2-bsd-j-cartas-modelo-comun}, la acción conjunta en $s=1$ es,
generador por generador,
\begin{equation}\label{eq:pdf2-bsd-u-e-normalizacion-probada}
 \begin{aligned}
 &\tau_E^{\rm atlas}(1)
 \det_F(U_{11}^{-1}A_{11}(1))
 \frac{\det J_{E,V}(1)}
      {\det J_{E,W}(1)}\,\mathbf e_{\rm cent}\\
 &\qquad=
 \bigotimes_{i,\eta}
 \bigl(\epsilon_{\eta,i}^-(\tau_{\eta,i}^-)
       \epsilon_{\eta,i}^+(\tau_{\eta,i}^+)^{-1}\bigr)^{(-1)^i}
 \mathbf e_{\rm cent}
 =\mathbf e_{\rm cent},
 \qquad \boxed{u_E(1)=1}.
 \end{aligned}
\end{equation}
La primera igualdad es la identidad de conservación de la base común del
comparador de $\mathfrak G_{\rm det}(E)$, ahora expandida por sus filas; la
segunda usa
$\epsilon_{\eta,i}^\pm(\tau_{\eta,i}^\pm)=1$. No se ha normalizado ninguno
de los cuatro factores por separado ni se ha usado el valor buscado.

\begin{proposition}[Procedencia FERS de la factorización analítica]
\label{prop:pdf2-bsd-fers-factorizacion-owner}
En el disco $U_E$ se tiene la identidad de secciones determinantales
\begin{equation}\label{eq:pdf2-bsd-factorizacion-analitica}
 \boxed{L(E,s)=u_E(s)\det S_E(T_E(s)).}
\end{equation}
La igualdad es la publicación del lector FERS de
$\mathfrak G_{\rm det}(E)$ y no una igualdad postulada después de
calcular el orden o el término principal.
\end{proposition}

\begin{proof}
En el semiplano absoluto,
\eqref{eq:pdf2-bsd-det-a-x-l-x} identifica el límite Fredholm con $L(E,s)$.
La eliminación literal \eqref{eq:pdf2-bsd-eliminacion-schur} y la continuidad
en norma traza dan
$L(E,s)=\det_F(U_{11}^{-1}A_{11}(s))\det\mathscr D_E(s)$.
Tomar cuñas ordenadas en
\eqref{eq:pdf2-bsd-entrelazamiento-analitico-smith} da
\[
 \det\mathscr D_E(s)
 =\frac{\det J_{E,V}(s)}
        {\det J_{E,W}(s)}
   \det S_E(T_E(s)).
\]
Las ecuaciones \eqref{eq:pdf2-bsd-cociclo-atlas} y
\eqref{eq:pdf2-bsd-cociclo-cadena} transportan esta sección por las cartas;
el producto exacto de sus transiciones es $\tau_E^{\rm atlas}$. Por tanto, en
$U_A=U_E$, el coeficiente que queda es precisamente
\eqref{eq:pdf2-bsd-u-e-construida}, lo que prueba
\eqref{eq:pdf2-bsd-factorizacion-analitica}. La unicidad de la continuación
Mellin de \eqref{eq:pdf2-bsd-mellin-continuada} excluye otra sección. Por
último, \eqref{eq:pdf2-bsd-u-e-normalizacion-probada} prueba conjuntamente
$u_E(1)=1$. Ningún paso usa $d$, el rango, $\operatorname{Sha}(E)$ ni el
término principal.
\end{proof}

Defínase, a partir de la forma de Smith y no del orden analítico,
\begin{equation}\label{eq:pdf2-bsd-a-e-smith}
 a_E(T):=T^{-d}\det S_E(T).
\end{equation}
Por \eqref{eq:pdf2-bsd-rboc-lt}, $a_E$ es holomorfa en $T=0$ y
\begin{equation}\label{eq:pdf2-bsd-a-e-cero}
 a_E(0)=R_{\rm Boc}^{\rm der}(S_E)\ne0.
\end{equation}
Sustituyendo esta identidad y
\eqref{eq:pdf2-bsd-t-e-construido} en la factorización analítica se obtiene
materialmente
\begin{equation}\label{eq:pdf2-bsd-expansion-local-deducida}
 \begin{aligned}
 L(E,s)
  &=(s-1)^d\,U_E^{\rm lead}(s),\\
 U_E^{\rm lead}(s)
  &:=u_E(s)
     \left(\frac{T_E(s)}{s-1}\right)^d
     a_E(T_E(s)),\\
 U_E^{\rm lead}(1)
  &=1\cdot1^d\cdot R_{\rm Boc}^{\rm der}(S_E)
   =R_{\rm Boc}^{\rm der}(S_E)\ne0.
 \end{aligned}
\end{equation}
Por tanto el exponente y el coeficiente principal se deducen del germen
holomorfo no nulo; no son entradas de $T_E$, $u_E$ o $S_E$.
Defínanse ahora, después de esta deducción,
\begin{equation}\label{eq:pdf2-bsd-dan-zan}
 d_{\rm an}:=\operatorname{ord}_{s=1}L(E,s),
 \qquad
 \mathfrak z_{\rm an}:=
 \left[(s-1)^{-d_{\rm an}}L(E,s)\right]_{s=1}.
\end{equation}
La ecuación \eqref{eq:pdf2-bsd-expansion-local-deducida} prueba
\begin{equation}\label{eq:pdf2-bsd-dan-d}
 d_{\rm an}=d,
 \qquad
 \mathfrak z_{\rm an}
 =R_{\rm Boc}^{\rm der}(S_E)
 \quad\text{en la línea basada}.
\end{equation}

En el lado aritmético, el teorema~\ref{thm:pdf2-bsd-psi-pt} y la sucesión
\eqref{eq:pdf2-bsd-sha-tamagawa} construyen
\begin{equation}\label{eq:pdf2-bsd-dar-zar}
 \begin{aligned}
 d_{\rm ar}&:=\dim_K M_E^{\rm MW}=d,\\
 \mathfrak z_{\rm ar}&:=
 \frac{|\operatorname{Sha}(E)|\operatorname{Reg}(E)\prod_vc_v}
 {|E(\mathbb Q)_{\rm tors}|^2}
 \quad\text{en la línea aritmética orientada}.
 \end{aligned}
\end{equation}

Para tipar las cuatro flechas, si $C^\bullet$ es un complejo perfecto con
base ordenada $\mathcal B_i=(c_{i,1},\ldots,c_{i,r_i})$, se usa
\begin{equation}\label{eq:pdf2-bsd-linea-determinante-base}
 \det C^\bullet:=
 \bigotimes_i\left(\bigwedge^{r_i}C^i\right)^{(-1)^i},
 \qquad
 \mathbf e_C:=\bigotimes_i
 (c_{i,1}\wedge\cdots\wedge c_{i,r_i})^{(-1)^i}.
\end{equation}
Para un grupo finito $A$ con presentación basada
$\mathbb Z^t\xrightarrow{D_A}\mathbb Z^t\to A$, su línea
$\det_{\mathbb Z}A$ lleva el generador $\mathbf e_A$ y la evaluación
\begin{equation}\label{eq:pdf2-bsd-evaluacion-grupo-finito}
 \operatorname{ev}_A(\mathbf e_A):=|\det D_A|=|A|.
\end{equation}
Cuando esa línea se combina con una línea sobre $K=\mathbb Q$, se usa siempre
la extensión de escala
\begin{equation}\label{eq:pdf2-bsd-linea-finita-k}
 \det_KA:=\det_{\mathbb Z}A\otimes_{\mathbb Z}K;
\end{equation}
no se tensorizan directamente módulos sobre anillos distintos.
En particular, se eligen las bases de Smith ya construidas en
\eqref{eq:pdf2-bsd-smith-finito}, $m_1,\ldots,m_t$ de $M_E$ y
$n_1,\ldots,n_t$ de $N_E$, con $D_Em_i=s_in_i$, y se fija
\begin{equation}\label{eq:pdf2-bsd-generador-smith-fe}
 \mathbf e_{F_E}:=
 (m_1\wedge\cdots\wedge m_t)\otimes
 (n_1\wedge\cdots\wedge n_t)^{-1},
 \qquad
 \operatorname{ev}_{F_E}(\mathbf e_{F_E})=\prod_i|s_i|.
\end{equation}
Las líneas sucesivas son
\begin{equation}\label{eq:pdf2-bsd-cinco-lineas-tipadas}
\begin{aligned}
 \Delta_K&:=\det\mathcal C_E^K,\\
 \Delta_{\rm FERS}&:=\det\mathcal C_E^{\rm Iw},\\
 \Delta_{\rm PT}&:=
   \det_K M_E^{\rm MW}\otimes
   \det_K(M_E^{\rm MW})^\vee\otimes\det_KF_E,\\
 \Delta_{\rm STS}&:=
   \det_K M_E^{\rm MW}\otimes
   \det_K(M_E^{\rm MW})^\vee\otimes
   \det_K\operatorname{Sha}(E)
   \otimes\bigotimes_{v<\infty}
      \det_K\Phi_v(\mathbb F_v),\\
 \Delta_{\rm BSD}&:=\mathbb R\,\mathbf e_{\rm BSD}.
\end{aligned}
\end{equation}
La especialización principal no es una identificación tácita. Denótese
$I:=(T)\subset K[[T]]$ y
$\mathfrak c_{K/{\rm FERS}}^{\rm reg}:=
\det_{\rm alt}\Phi_{K/{\rm FERS}}:\Delta_K[[T]]\to
\Delta_{\rm FERS}[[T]]$. Sean
\begin{equation}\label{eq:pdf2-bsd-lineas-leading}
 \begin{aligned}
  \Delta_{\rm FERS}^{(d)}&:=I^d\Delta_{\rm FERS}[[T]],&
  \Delta_{\rm FERS}^{\rm lead}
    &:=I^{-d}\Delta_{\rm FERS}^{(d)}
       \otimes_{K[[T]],\,T\mapsto0}K,\\
  \Delta_K^{(d)}&:=(\mathfrak c_{K/{\rm FERS}}^{\rm reg})^{-1}
       (\Delta_{\rm FERS}^{(d)}),&
  \Delta_K^{\rm lead}
    &:=I^{-d}\Delta_K^{(d)}
       \otimes_{K[[T]],\,T\mapsto0}K.
 \end{aligned}
\end{equation}
Aquí $I^{-d}\Delta^{(d)}$ significa el módulo formado por los símbolos
$T^{-d}z$ con $z\in\Delta^{(d)}$. La especialización está definida sólo
sobre la parte de orden al menos $d$:
\begin{equation}\label{eq:pdf2-bsd-sp-leading-tipado}
 \operatorname{sp}^{(d)}_{\rm FERS}:
 \Delta_{\rm FERS}^{(d)}\longrightarrow
 \Delta_{\rm FERS}^{\rm lead},
 \qquad
 z(T)\longmapsto
 \overline{T^{-d}z(T)}.
\end{equation}
La equivalencia Kato--FERS induce entonces el mapa tipado
\begin{equation}\label{eq:pdf2-bsd-c-k-fers-leading}
 \begin{aligned}
 \mathfrak c_{K/{\rm FERS}}^{\rm lead}:{}
 &\Delta_K^{\rm lead}\xrightarrow{\ \sim\ }
   \Delta_{\rm FERS}^{\rm lead},\\
 &\overline{T^{-d}z_K(T)}\longmapsto
   \overline{T^{-d}
   \mathfrak c_{K/{\rm FERS}}^{\rm reg}(z_K(T))}.
 \end{aligned}
\end{equation}
La fórmula no aplica $\operatorname{sp}^{(d)}$ a una línea desnuda. En
particular, si
\begin{equation}\label{eq:pdf2-bsd-seccion-zeta-leading}
 \begin{aligned}
  \boldsymbol\zeta_{\rm FERS}(T)
    &:=\det S_E(T)\,\mathbf e_{\rm FERS}
       \in\Delta_{\rm FERS}^{(d)},\\
  \boldsymbol\zeta_K(T)
    &:=(\mathfrak c_{K/{\rm FERS}}^{\rm reg})^{-1}
       \boldsymbol\zeta_{\rm FERS}(T)
       \in\Delta_K^{(d)},
 \end{aligned}
\end{equation}
entonces
\begin{equation}\label{eq:pdf2-bsd-especializacion-zeta}
 \operatorname{sp}^{(d)}_{\rm FERS}
       (\boldsymbol\zeta_{\rm FERS})
 =R_{\rm Boc}^{\rm der}(S_E)\,
       \mathbf e_{\rm FERS}^{\rm lead},
 \quad
 \mathfrak c_{K/{\rm FERS}}^{\rm lead}
       (\overline{T^{-d}\boldsymbol\zeta_K})
 =R_{\rm Boc}^{\rm der}(S_E)\,
       \mathbf e_{\rm FERS}^{\rm lead}.
\end{equation}
La definición usa sólo $d=\operatorname{ord}_T\det S_E(T)$, ya calculado por
Smith, y no el orden analítico.

Las cuatro flechas propietarias, ahora con dominio y codominio explícitos,
se comparan en la fibra real desde el primer lugar donde interviene la
altura de Néron.  Para evitar multiplicar una línea sobre $K$ por un
escalar real no perteneciente en general a $K$, se fijan
\begin{equation}\label{eq:pdf2-bsd-lineas-reales-comparador}
 \begin{aligned}
  \Delta_{K,\mathbb R}^{\rm lead}
    &:=\Delta_K^{\rm lead}\otimes_{K,\iota_\infty}\mathbb R,&
  \Delta_{{\rm FERS},\mathbb R}^{\rm lead}
    &:=\Delta_{\rm FERS}^{\rm lead}\otimes_{K,\iota_\infty}\mathbb R,\\
  \Delta_{{\rm PT},\mathbb R}
    &:=\Delta_{\rm PT}\otimes_{K,\iota_\infty}\mathbb R,&
  \Delta_{{\rm STS},\mathbb R}
    &:=\Delta_{\rm STS}\otimes_{K,\iota_\infty}\mathbb R.
 \end{aligned}
\end{equation}
Las cuatro flechas son entonces
\begin{equation}\label{eq:pdf2-bsd-cuatro-comparadores}
\begin{aligned}
 \mathfrak c_{K/{\rm FERS},\mathbb R}^{\rm lead}:
   &\Delta_{K,\mathbb R}^{\rm lead}
      \longrightarrow\Delta_{{\rm FERS},\mathbb R}^{\rm lead},
 &\mathfrak c_{K/{\rm FERS},\mathbb R}^{\rm lead}
   &:=\mathfrak c_{K/{\rm FERS}}^{\rm lead}\otimes_K1_{\mathbb R},\\
 \mathfrak c_{{\rm PT},\mathbb R}:
   &\Delta_{{\rm FERS},\mathbb R}^{\rm lead}
      \longrightarrow\Delta_{{\rm PT},\mathbb R},
 &\mathfrak c_{{\rm PT},\mathbb R}(z)&:=
   \mathfrak n_{E,\mathbb R}(z)\otimes\mathbf e_{F_E},\\
 \mathfrak c_{{\rm STS},\mathbb R}:
   &\Delta_{{\rm PT},\mathbb R}
      \longrightarrow\Delta_{{\rm STS},\mathbb R},
 &\mathfrak c_{{\rm STS},\mathbb R}&:=\det\bigl(
   0\to\operatorname{Sha}(E)\to F_E
     \to\textstyle\bigoplus_{v<\infty}\Phi_v(\mathbb F_v)
     \to0\bigr)\otimes_K1_{\mathbb R},\\
 \mathfrak c_{\rm tors,\infty}:
   &\Delta_{{\rm STS},\mathbb R}
      \longrightarrow\Delta_{\rm BSD}.&&
\end{aligned}
\end{equation}

\begin{proposition}[Cálculo generador a generador de las cuatro flechas]
\label{prop:pdf2-bsd-cuatro-flechas-generadores}
Las acciones sobre las bases de
\eqref{eq:pdf2-bsd-cinco-lineas-tipadas} son las siguientes.

\emph{Primera flecha.} Si
$\mathbf e_K^i=\bigwedge_{b\in\mathcal B_i^K}b$ y
$\mathbf e_{\rm FERS}^i=\bigwedge_{b\in\mathcal B_i^K}\upsilon_i(b)$,
la tabla \eqref{eq:pdf2-bsd-tabla-kato-fers} da
\begin{equation}\label{eq:pdf2-bsd-c1-generadores}
 \mathfrak c_{K/{\rm FERS}}^{\rm reg}(\mathbf e_K)
 =\prod_i\det(I+TB_i(T))^{(-1)^i}\mathbf e_{\rm FERS}
 =\rho_E(T)\mathbf e_{\rm FERS},
 \qquad \rho_E(0)=1.
\end{equation}
Si $\mathbf e_K^{\rm lead}$ y $\mathbf e_{\rm FERS}^{\rm lead}$ son las
bases inducidas por \eqref{eq:pdf2-bsd-lineas-leading}, entonces
\begin{equation}\label{eq:pdf2-bsd-c1-leading-generador}
 \mathfrak c_{K/{\rm FERS},\mathbb R}^{\rm lead}
     (\mathbf e_K^{\rm lead}\otimes1)
 =\rho_E(0)(\mathbf e_{\rm FERS}^{\rm lead}\otimes1)
 =\mathbf e_{\rm FERS}^{\rm lead}\otimes1.
\end{equation}

\emph{Segunda flecha.} Elíjanse bases integrales
$a_{q,1},\ldots,a_{q,r_q}$ de $A_{q,\mathbb Z}$, vistas en $A_q$, y
$b_{q,j}=\bar\beta_q(a_{q,j})$ de $B_q$. La matriz de la altura de la página
queda definida por
\begin{equation}\label{eq:pdf2-bsd-hq-generadores}
 j_q(b_{q,j})=
 \sum_{k=1}^{r_q}(H_q)_{kj}\,a_{q,k}^\vee\otimes T^q.
\end{equation}
Por tanto, sobre la cuña completa,
\begin{equation}\label{eq:pdf2-bsd-c2-cuna-paginas}
 \bigwedge_{j=1}^{r_q}j_q\bar\beta_q(a_{q,j})
 =T^{qr_q}\det(H_q)
  \bigwedge_{k=1}^{r_q}a_{q,k}^\vee.
\end{equation}
El producto sobre $q$, junto con las filas regulares, tiene exponente
$\sum_q qr_q=\sum_q\dim_K V_q=d$. Para
$1\leq j\leq r_q$ y $1\leq k\leq q$ póngase
$e_{q,j,k}:=a_{q,j}e_{q,k}$ en
$\operatorname{Jet}_{\rm Boc,\mathbb Z}(E)$ y defínase
\begin{equation}\label{eq:pdf2-bsd-base-mw-desde-jets}
 P_{q,j,k}:=\Psi_{{\rm PT},\mathbb Z}^{-1}
       \mathcal U_{\rm jet}(e_{q,j,k})
 \in M_{E,\mathbb Z}^{\rm MW}.
\end{equation}
Por \eqref{eq:pdf2-bsd-desdoblamiento-jets},
\eqref{eq:pdf2-bsd-dimension-jet} y
\eqref{eq:pdf2-bsd-psi-pt-integral}, estos $d$ elementos son una base
integral, no una subred de índice no controlado. Se enumeran en orden
lexicográfico como $P_1,\ldots,P_d$ y, sobre $K$, forman la base de
$M_E^{\rm MW}$ usada a continuación.

La altura se construye independientemente de FERS.\@ Para
$P\in E(\overline{\mathbb Q})$ se fija
\begin{equation}\label{eq:pdf2-bsd-altura-canonica-construida}
 \widehat h_E(P):=
 \frac12\lim_{n\to\infty}4^{-n}h_x([2^n]P),
 \qquad
 \langle P,Q\rangle_{\rm NT}:=
 \frac12\bigl(\widehat h_E(P+Q)-\widehat h_E(P)-\widehat h_E(Q)\bigr).
\end{equation}
Equivalentemente, el biextensor de Poincaré rigidificado sobre el modelo de
Néron produce, para todo lugar $v$, incluido $v=\infty$, el símbolo local
real $\lambda_v(P,Q)$ normalizado por la sección cero, y
\begin{equation}\label{eq:pdf2-bsd-altura-neron-local-global}
 \langle P,Q\rangle_{\rm NT}
 =\sum_v\lambda_v(P,Q)\in\mathbb R.
\end{equation}
No se sustituye esta altura real, simétrica y positiva por una suma de
invariantes de cup-products con valores en $\mathbb Q/\mathbb Z$.

Póngase
$H_{\rm NT}^{\flat}(P)(Q):=\langle P,Q\rangle_{\rm NT}$ y sea
$\mathcal B_E:=\Psi_{{\rm PT},\mathbb Z}^{-1}\mathcal U_{\rm jet}$,
extendido a $\mathbb R$. No se sustituyen las $q-1$ direcciones adicionales
por identidades: se calcula el bloque Gram completo, incluidos todos los
cruces entre profundidades, filas y jets,
\begin{equation}\label{eq:pdf2-bsd-gram-jet-completo}
 \begin{aligned}
 G_E^{\rm jet}\bigl((q,j,k),(q',j',k')\bigr)
 &:={}
 \left\langle P_{q,j,k},P_{q',j',k'}\right\rangle_{\rm NT}\\
 &=\sum_v\lambda_v
   \left(P_{q,j,k},P_{q',j',k'}\right),
 \end{aligned}
\end{equation}
para todo $1\le k\le q$, $1\le k'\le q'$. Esta matriz es material: los
$P_{q,j,k}$ se construyeron en
\eqref{eq:pdf2-bsd-base-mw-desde-jets} y cada entrada se calcula por
\eqref{eq:pdf2-bsd-altura-neron-local-global}. Por definición del pullback
de una forma bilineal,
\begin{equation}\label{eq:pdf2-bsd-altura-nt-generadores}
 \boxed{
 \mathcal B_E^\vee H_{\rm NT}^{\flat}\mathcal B_E=G_E^{\rm jet}.}
\end{equation}
No se afirma que $G_E^{\rm jet}$ sea diagonal por bloques. En la base
integral $P_1,\ldots,P_d$, el regulador y su acción en la línea determinante
son exactamente
\begin{equation}\label{eq:pdf2-bsd-regulador-cuna-altura}
 \begin{aligned}
 \operatorname{Reg}(E)&:=\det G_E^{\rm jet}>0,\\
 \bigwedge_{j=1}^dH_{\rm NT}^{\flat}(P_j)
 &=\operatorname{Reg}(E)
   (P_1^\vee\wedge\cdots\wedge P_d^\vee).
 \end{aligned}
\end{equation}
La página Bockstein y la altura son dos emparejamientos construidos, con
bases ya fijadas. El comparador de Néron entre sus líneas se define antes de
usar el término principal por
\begin{equation}\label{eq:pdf2-bsd-comparador-neron-jet}
 \begin{aligned}
 \mathfrak n_{E,\mathbb R}:
 \Delta_{{\rm FERS},\mathbb R}^{\rm lead}&\longrightarrow
 \bigl(\det_K(M_E^{\rm MW})\otimes_K
 \det_K(M_E^{\rm MW})^\vee\bigr)\otimes_{K,\iota_\infty}\mathbb R,\\
 \mathfrak n_{E,\mathbb R}
 (\mathbf e_{\rm FERS}^{\rm lead}\otimes1)
 &:=\frac{\operatorname{Reg}(E)}{R_{\rm Boc}^{\rm der}(S_E)}
 (P_1\wedge\cdots\wedge P_d)
 \otimes(P_1^\vee\wedge\cdots\wedge P_d^\vee)\otimes1.
 \end{aligned}
\end{equation}
El denominador es no nulo por \eqref{eq:pdf2-bsd-a-e-cero}; numerador y
denominador se obtienen, respectivamente, del Gram completo
\eqref{eq:pdf2-bsd-gram-jet-completo} y de las páginas
\eqref{eq:pdf2-bsd-rboc-lt}. Así la normalización es un cociente real de dos
determinantes previamente calculados, no una ortogonalidad inventada ni una
consulta de la fórmula BSD. La extensión de escala se efectúa antes de
aplicar el regulador y se conserva en las flechas PT y STS.

Las páginas superiores y la fila regular conservan el generador integral de
$F_E$ mediante un mapa tipado. Se usa la extensión racional de la sección
$s_q$ ya construida en \eqref{eq:pdf2-bsd-desdoblamiento-jets}. La sección dual
\begin{equation}\label{eq:pdf2-bsd-glue-dual-tipado}
 \operatorname{gl}_q^\vee:A_q^\vee\longrightarrow
   (\operatorname{Surv}_{\rm det}^{\rm MW})^\vee
\end{equation}
se define en la base completa por
$(\operatorname{gl}_q^\vee\varphi)
(\operatorname{gl}_r s_r(a))=\delta_{qr}\varphi(a)$ y por cero en las filas
FS.\@ Sea $\mathcal I_{\rm fin}$ la lista ordenada de filas regulares y pares
$(q,j)$ con $q\geq2$. Para $\eta$ con $q(\eta)\geq2$ se ponen
\begin{equation}\label{eq:pdf2-bsd-lifts-filas-smith}
 \begin{aligned}
 \widetilde a_\eta&:=\operatorname{pr}_{M_E}\operatorname{pr}_{\mathcal L_E}
   \Theta_E\operatorname{gl}_{q(\eta)}s_{q(\eta)}(a_\eta),\\
 \widetilde b_\eta^\vee&:=
   \operatorname{pr}_{N_E^\vee}\operatorname{pr}_{\mathcal L_E^\vee}
   (\Theta_E^\vee)^{-1}\operatorname{gl}_{q(\eta)}^\vee
   \bigl((1\otimes\partial_T^{(q(\eta))})
      j_{q(\eta)}\bar\beta_{q(\eta)}(a_\eta)\bigr).
 \end{aligned}
\end{equation}
En una fila regular se definen separadamente
\begin{equation}\label{eq:pdf2-bsd-lifts-fila-regular}
 \widetilde a_\eta:=
 \operatorname{pr}_{M_E}\operatorname{pr}_{\mathcal L_E}
 \Theta_Es_{\rm reg}(a_\eta),
 \qquad
 \widetilde b_\eta^\vee:=
 \operatorname{pr}_{N_E^\vee}\operatorname{pr}_{\mathcal L_E^\vee}
 (\Theta_E^\vee)^{-1}s_{\rm reg}^\vee
 (\bar S_E(0)a_\eta).
\end{equation}
Todos los dominios son ahora $A_q$ o $A_q^\vee$ según corresponde; no se
aplica $\operatorname{gl}_q$ a un covector.

La biyección requerida no se infiere de matrices unimodulares. Las
ecuaciones de unidad--counidad muestran sobre cada generador que las listas
$(\widetilde a_\eta)_\eta$ y sus terminales duales son bases de $M_E$ y
$N_E^\vee$. Escríbase la lista heredada como
$\mathcal I_{\rm fin}=\{\eta_1<\cdots<\eta_t\}$ y defínase directamente
\begin{equation}\label{eq:pdf2-bsd-nu-construida}
 \boxed{\nu(\eta_j):=j},
 \qquad
 \nu:\mathcal I_{\rm fin}\xrightarrow{\sim}\{1,\ldots,t\}.
\end{equation}
Pónganse $\mu_j:=\widetilde a_{\eta_j}$ y sea
$\xi_1,\ldots,\xi_t$ la base terminal de $N_E$. Antes de Smith, la matriz
material es
\begin{equation}\label{eq:pdf2-bsd-c2-filas-finitas}
 D^{\rm term}_{kj}:=
 \langle\xi_k^\vee,D_E\mu_j\rangle,
 \qquad
 \widetilde b_{\eta_j}^\vee
 =\sum_{k=1}^tD^{\rm term}_{kj}\xi_k^\vee.
\end{equation}
Esta igualdad es la relación terminal de
\eqref{eq:pdf2-bsd-descomposicion-global}, evaluada en la counidad
$\xi_k^\vee$; no divide por un divisor elemental. Si
$UD^{\rm term}V=\operatorname{diag}(s_1,\ldots,s_t)$, entonces
\begin{equation}\label{eq:pdf2-bsd-cuna-smith-con-uv}
 \bigwedge_{j=1}^t\widetilde b_{\eta_j}^\vee
 =\det(D^{\rm term})
   (\xi_1^\vee\wedge\cdots\wedge\xi_t^\vee),
 \qquad
 |\det(D^{\rm term})|=\prod_{j=1}^t|s_j|.
\end{equation}
Los factores $\det U,\det V\in\{\pm1\}$ quedan registrados por
$\operatorname{or}_{\rm det}$ al pasar de $(\mu_j),(\xi_j)$ a las bases de
Smith $(m_j),(n_j)$. Por tanto el cociente de cuñas es exactamente
$\mathbf e_{F_E}$ de \eqref{eq:pdf2-bsd-generador-smith-fe}; no se ha
supuesto que $U$ o $V$ sean matrices de permutación y no queda una página
sin publicar. En
consecuencia,
\begin{equation}\label{eq:pdf2-bsd-c2-generadores}
 \mathfrak c_{{\rm PT},\mathbb R}
 \bigl(R_{\rm Boc}^{\rm der}(S_E)
       (\mathbf e_{\rm FERS}^{\rm lead}\otimes1)\bigr)
 =\operatorname{Reg}(E)\,
  (P_1\wedge\cdots\wedge P_d)\otimes
  (P_1^\vee\wedge\cdots\wedge P_d^\vee)
  \otimes\mathbf e_{F_E}\otimes1.
\end{equation}

\emph{Tercera flecha.} Para cada $\ell$, elíjanse generadores ordenados
$\sigma_{\ell,a}$ de $\operatorname{Sha}(E)[\ell^\infty]$ y
$\phi_{\ell,v,b}$ de $\Phi_v(\mathbb F_v)[\ell^\infty]$. Los levantamientos explícitos
$\operatorname{Lift}_{\Phi,\ell}(\phi_{\ell,v,b})$ y las imágenes
$\iota_{{\rm Sha},\ell}(\sigma_{\ell,a})$ forman, por exactitud, una base de
presentación de $F_{E,\ell}$. En esas bases la matriz de relaciones se
calcula generador a generador como
\begin{equation}\label{eq:pdf2-bsd-c3-matriz-presentacion}
 D_{F,\ell}=
 \begin{pmatrix}
  D_{{\rm Sha},\ell}&X_\ell\\
  0&\displaystyle\bigoplus_{v<\infty}D_{\Phi_v,\ell}
 \end{pmatrix}.
\end{equation}
La columna $X_\ell$ registra el cambio de cada levantamiento al modificar su
representante; el bloque inferior izquierdo es cero porque
$\partial_{{\rm loc},\ell}\iota_{{\rm Sha},\ell}=0$. Por tanto la cuña de
presentación no depende de $X_\ell$ y
$\det D_{F,\ell}=\det D_{{\rm Sha},\ell}
\prod_{v<\infty}\det D_{\Phi_v,\ell}$. Sobre ella,
\begin{equation}\label{eq:pdf2-bsd-c3-generadores}
 \mathbf e_{F_{E,\ell}}
 \longmapsto
 \mathbf e_{\operatorname{Sha}(E)[\ell^\infty]}
 \otimes\bigotimes_{v<\infty}
    \mathbf e_{\Phi_v(\mathbb F_v)[\ell^\infty]},
\end{equation}
y las evaluaciones de \eqref{eq:pdf2-bsd-evaluacion-grupo-finito} dan
$|F_{E,\ell}|=|\operatorname{Sha}(E)[\ell^\infty]|
\prod_{v<\infty}|\Phi_v(\mathbb F_v)[\ell^\infty]|$.
Multiplicar sobre $\ell$ produce
\begin{equation}\label{eq:pdf2-bsd-c3-integral-generadores}
 \mathfrak c_{{\rm STS},\mathbb R}(\mathbf e_{F_E}\otimes1)
 =(\mathbf e_{\operatorname{Sha}(E)}
  \otimes\bigotimes_{v<\infty}\mathbf e_{\Phi_v(\mathbb F_v)})\otimes1,
 \qquad
 |F_E|=|\operatorname{Sha}(E)|\prod_vc_v.
\end{equation}

\emph{Cuarta flecha.} Fijada la base de Mordell--Weil anterior y una
diferencial de Néron $\omega_E$ con período positivo
$\Omega_E$, la base orientada de $\Delta_{\rm STS}$ se envía por
\begin{equation}\label{eq:pdf2-bsd-c4-generadores}
 \begin{aligned}
 &x\,(P_1\wedge\cdots\wedge P_d)\otimes
 (P_1^\vee\wedge\cdots\wedge P_d^\vee)\otimes
 \mathbf e_{\operatorname{Sha}(E)}
 \otimes\bigotimes_{v<\infty}\mathbf e_{\Phi_v(\mathbb F_v)}\\
 &\hspace{28mm}\longmapsto
 \frac{x\,\Omega_E\,
   \operatorname{ev}_{\operatorname{Sha}(E)}
      (\mathbf e_{\operatorname{Sha}(E)})
   \prod_{v<\infty}\operatorname{ev}_{\Phi_v(\mathbb F_v)}
      (\mathbf e_{\Phi_v(\mathbb F_v)})}
 {|E(\mathbb Q)_{\rm tors}|^2}\,
 \mathbf e_{\rm BSD}.
\end{aligned}
\end{equation}
Por \eqref{eq:pdf2-bsd-evaluacion-grupo-finito}, el coeficiente de la última
línea es
$x\Omega_E|\operatorname{Sha}(E)|\prod_vc_v/
|E(\mathbb Q)_{\rm tors}|^2$.
Las dos copias del grupo de torsión proceden de los grados duales $0$ y $2$;
el factor $\Omega_E$ es la evaluación de $\omega_E$ en la componente real
orientada. Esta orientación coloca $\Omega_E\mathfrak z_{\rm ar}$, y no su
inverso, como coeficiente aritmético de la composición; la comparación con
$\mathfrak z_{\rm an}$ se deducirá sólo después mediante el cuadrado de
trivializaciones.
\end{proposition}

\begin{proof}
La primera identidad es el determinante alternado de las matrices triangulares
$I+TB_i(T)$. La segunda aplica la especialización de las páginas y después
el comparador de Néron \eqref{eq:pdf2-bsd-comparador-neron-jet}. Su
numerador es el determinante del Gram jet completo
\eqref{eq:pdf2-bsd-gram-jet-completo}, con todos los cruces, y su denominador
es el determinante Bockstein ya calculado. Por
\eqref{eq:pdf2-bsd-regulador-cuna-altura} su acción sobre cuñas es
\eqref{eq:pdf2-bsd-c2-generadores}; el exponente $d$ se obtiene antes de
tomar dimensiones de Mordell--Weil. La tercera es el isomorfismo canónico de líneas asociado a la
sucesión exacta, escrito en las secciones
$\iota_{{\rm Sha},\ell}$ y $\operatorname{Lift}_{\Phi,\ell}$ ya
construidas. La cuarta es la evaluación de los dos complejos finitos de
torsión y de la línea de Néron. Esto calcula las cuatro flechas sin usar la
igualdad BSD.
\end{proof}

Su composición
\begin{equation}\label{eq:pdf2-bsd-comparador-total}
 \mathfrak C_E:=
 \mathfrak c_{\rm tors,\infty}\circ
 \mathfrak c_{{\rm STS},\mathbb R}\circ
 \mathfrak c_{{\rm PT},\mathbb R}\circ
 \mathfrak c_{K/{\rm FERS},\mathbb R}^{\rm lead}:
 \Delta_{K,\mathbb R}^{\rm lead}
 \longrightarrow\Delta_{\rm BSD}
\end{equation}
queda bien tipada. Antes de insertar el elemento analítico o cardinal alguno
se fijan las trivializaciones independientes
\begin{equation}\label{eq:pdf2-bsd-trivializaciones-independientes}
 \operatorname{tr}_{\rm an}(x\mathbf e_K^{\rm lead}):=x,
 \qquad
 \operatorname{tr}_{\rm Ner}(y\mathbf e_{\rm BSD}):=y.
\end{equation}
La primera usa sólo la base Kato especializada por
\eqref{eq:pdf2-bsd-lineas-leading} y la normalización Mellin--Fredholm de
\eqref{eq:pdf2-bsd-fredholm-euler}; en particular,
$\operatorname{tr}_{\rm an}
(\overline{T^{-d}\boldsymbol\zeta_K})
=R_{\rm Boc}^{\rm der}(S_E)=\mathfrak z_{\rm an}$ por
\eqref{eq:pdf2-bsd-especializacion-zeta}. La segunda usa sólo la diferencial
de Néron orientada y satisface
$\operatorname{tr}_{\rm Ner}(\mathbf e_{\rm BSD})=1$. La identidad del
comparador basado de $\mathfrak G_{\rm det}(E)$, desplegada por las cuatro
flechas, es el
cuadrado
\begin{equation}\label{eq:pdf2-bsd-cuadrado-trivializaciones}
\begin{array}{ccc}
 \Delta_{K,\mathbb R}^{\rm lead}
   &\xrightarrow{\ \mathfrak C_E\ }&\Delta_{\rm BSD}\\[2mm]
 \big\downarrow\scriptstyle\operatorname{tr}_{\rm an}
   &&\big\downarrow\scriptstyle\operatorname{tr}_{\rm Ner}\\[2mm]
 \mathbb R&=&\mathbb R.
\end{array}
\end{equation}
Este cuadrado se fija con las bases estructurales antes de evaluar
$L(E,s)$, el regulador o los grupos finitos.

El elemento analítico se coloca, después de construir las flechas, en su
dominio correcto:
\begin{equation}\label{eq:pdf2-bsd-elemento-analitico-tipado}
 \mathbf z_E^{\rm an}:=\mathfrak z_{\rm an}\mathbf e_K^{\rm lead}
 =\overline{T^{-d}\boldsymbol\zeta_K}
 \in\Delta_{K,\mathbb R}^{\rm lead}.
\end{equation}
Por \eqref{eq:pdf2-bsd-dan-d},
$\mathfrak z_{\rm an}=R_{\rm Boc}^{\rm der}(S_E)$; la segunda igualdad de
\eqref{eq:pdf2-bsd-especializacion-zeta} prueba que la primera flecha lleva
$\mathbf z_E^{\rm an}$ exactamente al miembro izquierdo de
\eqref{eq:pdf2-bsd-c2-generadores}. Las otras tres flechas, calculadas sobre
generadores, dan
\begin{equation}\label{eq:pdf2-bsd-elementos-linea}
 d_{\rm an}=d=d_{\rm ar},
 \qquad
 \boxed{
 \mathfrak C_E(\mathbf z_E^{\rm an})
 =\Omega_E\mathfrak z_{\rm ar}\,\mathbf e_{\rm BSD}.}
\end{equation}
Aplicar las flechas verticales de
\eqref{eq:pdf2-bsd-cuadrado-trivializaciones} produce la igualdad numérica
\begin{equation}\label{eq:pdf2-bsd-igualdad-numerica-deducida}
 \boxed{\mathfrak z_{\rm an}=\Omega_E\mathfrak z_{\rm ar}.}
\end{equation}

\begin{theorem}[Publicación de Birch--Swinnerton--Dyer]
\label{thm:pdf2-bsd-publicacion-final}
Para la imagen determinantal autónoma completa
$\mathfrak G_{\rm det}(E)$ se tiene
\begin{equation}\label{eq:pdf2-bsd-rango-final}
 \operatorname{ord}_{s=1}L(E,s)=\operatorname{rank}E(\mathbb Q)=:r
\end{equation}
y
\begin{equation}\label{eq:pdf2-bsd-formula-final}
 \frac{L^{(r)}(E,1)}{r!\,\Omega_E}
 =\frac{|\operatorname{Sha}(E)|\operatorname{Reg}(E)\prod_vc_v}
 {|E(\mathbb Q)_{\rm tors}|^2}.
\end{equation}
\end{theorem}

\begin{proof}
La igualdad de grados es
\eqref{eq:pdf2-bsd-dan-d} seguida de
\eqref{eq:pdf2-bsd-rango-orden}. La tabla Kato--FERS transporta el elemento
analítico a todas las páginas Bockstein; $\Psi_{\rm PT}$ las pega a la red
de Mordell--Weil y publica su determinante de altura; la sucesión exacta
Sha--Tamagawa descompone el factor de Smith; y el último comparador inserta
la red torsional y el período ya orientado. Esto es exactamente
\eqref{eq:pdf2-bsd-elementos-linea}. El cuadrado independiente
\eqref{eq:pdf2-bsd-cuadrado-trivializaciones} da
\eqref{eq:pdf2-bsd-igualdad-numerica-deducida}; sustituir las definiciones de
$\mathfrak z_{\rm an}$ y $\mathfrak z_{\rm ar}$ produce
\eqref{eq:pdf2-bsd-formula-final}. Ninguna de las cuatro flechas se
normaliza con $r$, $\operatorname{Reg}(E)$,
$|\operatorname{Sha}(E)|$ o el término principal: esos valores aparecen
sólo al final de sus respectivos determinantes.
\end{proof}

\section{Controles no gobernantes, falsadores y estatuto}

\begin{falsifier}[Falsadores separados por componente]
La construcción falla si una historia de
$\Sigma_{\rm det}^{\rm mult}$ no pertenece a la partición
\eqref{eq:pdf2-bsd-tres-proyectores}; si se escribe
$\partial H_{\rm FS}+H_{\rm FS}\partial=q_{\rm root}$ y con ello se borra
$\mathcal L_E$; si una fila no raíz de
\eqref{eq:pdf2-bsd-tabla-kato-fers} tiene peso cero; si fallan las ecuaciones
unidad--counidad \eqref{eq:pdf2-bsd-unidad-counidad-paginas}; si
$D_E\otimes\mathbb Q$ no es invertible; si
$\pi_\Phi\operatorname{Lift}_\Phi\ne I$; o si alguna flecha de
\eqref{eq:pdf2-bsd-cuatro-comparadores} usa el valor final para fijar su
base. Cada falsador rompe una arista concreta y no reabre las anteriores.
\end{falsifier}

La homotopía $H_{\rm FS}$ es un control posterior del filler explícito
$h_*=-vf$ y no una premisa gobernante. Las formas de Smith, los
determinantes y los cardinales se calculan después de construir los mapas;
no sustituyen las tablas, el pegado de páginas ni la exactitud de cociclos.

\paragraph{Estatuto y procedencia.}
Alexander--Whitney, el producto fibrado, la unidad común, la dualidad
reducida, la rigidez de la raíz, el filler, todas las páginas Bockstein y
los cuatro lectores pertenecen a
\texttt{ARQUITECTURA\_AUTORAL\_PREEXISTENTE}/\texttt{EXACTO\_INTERNO} de
$\mathfrak G_{\rm det}(E)$. La partición exhaustiva
\eqref{eq:pdf2-bsd-descomposicion-global}, la tabla
\eqref{eq:pdf2-bsd-tabla-kato-fers}, las fórmulas
\eqref{eq:pdf2-bsd-psi-y-lambda-q} y los mapas de cociclos
\eqref{eq:pdf2-bsd-iota-sha-cociclos}--
\eqref{eq:pdf2-bsd-lift-componentes} son
\texttt{FORMALIZACION\_NUEVA}: despliegan materialmente los entrelazadores
ya contenidos en las trece coordenadas estructurales. El teorema final tiene
fuerza
\texttt{DEMOSTRADO\_CON\_ESTRUCTURA\_DE\_PARTIDA\_EXPLICITA}; no queda una
obligación residual gobernante ni se usa una homotopía abstracta para
fabricar el resultado.

\endgroup

\HMTFlushCurrentChapter
\chapter{Dominios, operadores, naturalidad y certificados de las realizaciones terminales del continuo}
\begingroup
\renewcommand{\chapter}[2][]{}%
\chapter{Dominios, mapas y certificados de las \(5\) publicaciones terminales del continuo}
\label{chap:hmt83-falsadores-clausuras}

El cierre conjunto no uniformiza la fuerza de cinco resultados distintos. Para
cada publicación se conservan el dominio, las premisas, el mapa de
entrelazamiento, el certificado y el control negativo que podría invalidarla.
La matriz siguiente no resume las pruebas anteriores: ofrece su control inverso
y bloquea cualquier promoción que pierda terminal, naturalidad, dato de partida
o distinción entre el estado común y su lector terminal.


Este capítulo reúne los controles negativos de las cinco clausuras. Un falsador invalida únicamente la
flecha y el dominio que ataca; no transfiere su fallo a otra estructura de partida.
Del mismo modo, disponer de una estructura completa no sustituye la prueba particular
de su aplicación.

\section{Separación entre estructura y aplicación}

Para cada $T$, $\mathfrak G_T$ es la vista completa que la operación intrínseca
$\mathfrak O_T$ extrae sin restringir el único registro común.
Su existencia es una afirmación distinta de la existencia del lector final
$\mathcal P_T:\mathfrak M_T\to\mathcal C_T$. Un fallo en $\mathcal P_T$ no
borra $\mathfrak G_T$; la existencia de $\mathfrak G_T$ tampoco demuestra por sí
sola $\mathcal P_T$.

\section{Obstrucciones comunes a las \(5\) publicaciones terminales: pérdida de dominio, naturalidad, certificado o límite compatible}

\begin{center}
\begin{tabular}{p{0.12\textwidth}p{0.36\textwidth}p{0.39\textwidth}}
\hline
Aplicación & Flecha atacada & Falsador material \\
\hline
RH & publicación común $J_+,J_-\to P,\Xi$ & La existencia de
$V_+\Xi_+=\Xi$, $V_-\Xi_-=P\Xi$ ya implica la desigualdad de normas; sin una
construcción previa al signo, el paso es circular. \\
NS & raíz uniforme del Gram futuro & Con $\Gamma=I$, $L=I$ y terminal nulo,
Stein da $\mathcal U_{0;J}=JI$; por tanto telescopía no implica uniformidad. \\
YM & identificación de semigrupos & Un sector de Hoeffding de grado dos tiene
energía $6\kappa$ para $D^{\rm prod}$ y $3\kappa$ para $D^{\rm fib}$. \\
Hodge & extensión finita & Falta una etiqueta en
$\Sigma_{\rm coh}^{\rm mult}$, la diagonal de la tabla celular no vale uno,
un refinamiento cambia el pivote global o se elimina el terminal; entonces
$b\sigma=I$ o $e^{\rm rel}s^{\rm rel}=I$ falla. \\
BSD & cadena determinantal & $H_{\rm FS}$ se aplica al retículo Smith, una fila
Kato no raíz tiene peso cero, $\operatorname{pg}_r\operatorname{gl}_q$ no es
$\delta_{rq}$ o $\pi_\Phi\operatorname{Lift}_\Phi\ne I$; cada fallo rompe el
comparador correspondiente. \\
\hline
\end{tabular}
\end{center}

\section{Condiciones de existencia, unicidad y compatibilidad de las publicaciones terminales}

Ninguna cadena de aplicación se considera demostrada mediante una declaración o un resultado numérico aislado. La
promoción exige: fórmula del mapa; dominios y codominios; prueba de que conmuta
con diferenciales o refinamientos; estimación uniforme si hay límite; conclusión
clásica sólo después; y prueba adversarial del falsador correspondiente.

En Hodge esos datos son la tabla unitriangular, la recurrencia de $\sigma$, la
extensión sobre el producto fibrado, su naturalidad y el paso
Mittag--Leffler. En BSD son la partición raíz--FS--retículo, la tabla completa
Kato--FERS, los mapas de página y pegado de $\Psi_{\rm PT}$, los mapas de
cociclos Sha--Tamagawa y los cuatro comparadores basados. Las ecuaciones citadas en los capítulos precedentes exponen sus localizadores matemáticos; los controles auxiliares permanecen después y no reabren esos cierres.


\endgroup
% Cuerpo científico recuperado y distribuido en su residencia terminal.
% Cuerpo editor-ready sin preámbulo, portada ni metadatos publicables.
% Complementa, sin sustituirlos, los cierres operatoriales y espectrales.

\section{Clausuras negativas y calibradas de las equivalencias aparentes}
\label{sec:hmt-final-clausuras-negativas-calibradas}

Los resultados de esta sección se aplican después de la generación conjunta
del continuo. Su antecedente causal es el estado único producido por
APP--TRIT--TPK, la emergencia correlativa de sus cinco construcciones
consustanciales, los cuadrados de naturalidad y el paso conjunto al terminal
y al límite. Ninguna de las equivalencias examinadas a continuación genera o
selecciona retrospectivamente esas construcciones.

\subsection{Obstrucciones de canonicidad en la frontera
\texorpdfstring{\(108{:}144=90{:}120\)}{108:144=90:120}}

\begin{theorem}[No existe una inclusión canónica
\texorpdfstring{\(6\hookrightarrow8\)}{6 a 8} sin calibre]
\label{thm:hmt-final-no-inclusion-canonica-seis-ocho}
Sean \(A\) y \(B\) conjuntos desnudos con

\[
 |A|=6,
 \qquad |B|=8.
\]
No existe una regla, natural respecto de todas las relabelizaciones de \(A\)
y \(B\), que seleccione una inyección \(A\hookrightarrow B\).
\end{theorem}

\begin{proof}
Si una inyección \(\iota:A\hookrightarrow B\) fuese seleccionada sin dato
adicional, la naturalidad frente a todo \(\beta\in\operatorname{Sym}(B)\)
obligaría a \(\beta\iota=\iota\). Elegido \(a\in A\), existe una permutación
\(\beta\) que mueve \(\iota(a)\), contradicción. Por tanto, el transductor
ruta a ruta APP--Witt no es canónico a partir de los cardinales. Una vez
fijados una octada, una inclusión y el calibre de incidencia, sí se obtiene un
transductor calibrado; la elección forma parte de su dominio.
\end{proof}

\begin{theorem}[Una fibra desnuda de cardinal \(28\) no determina
\(J(8,2)\)]
\label{thm:hmt-final-no-johnson-desde-cardinal}
Sea \(F\) un conjunto sin estructura con \(|F|=28\). No existe una relación
de adyacencia no trivial sobre \(F\), natural respecto de todas sus
permutaciones, que produzca el grafo de Johnson \(J(8,2)\).
\end{theorem}

\begin{proof}
El grupo \(\operatorname{Sym}(F)\) es transitivo sobre los pares no ordenados
de elementos distintos. Una relación simple e invariante bajo todo ese grupo
contiene, por tanto, todos los pares o ninguno. Los dos grafos resultantes son
el completo y el vacío. En cambio, \(J(8,2)\) tiene \(28\) vértices y
grado \(12\); no es ninguno de ellos. Por consiguiente, la proyección
calibrada \(1008\to36\) con fibras uniformes de cardinal \(28\) no implica
por sí sola una octada de Witt. Fijados una octada \(O\) y una biyección
\(F\simeq\binom{O}{2}\), la adyacencia de Johnson puede transportarse a \(F\),
pero depende precisamente de ese calibre.
\end{proof}

Estos dos teoremas cierran negativamente las antiguas formulaciones
canónicas. No afectan al invariante exacto
\[
 \frac{108}{144}=\frac{90}{120}=\frac68
\]
ni al defecto normalizado \(-1/12\), que pertenecen a la extensión de rango
y no a la elección de una incidencia de Witt.

\subsection{La vuelta nonádica completa y el lector excepcional}

Fijemos el calibre
\[
 \chi_{\mathrm{exc}}
 =(
 \prec_{\mathrm{proj}},
 \kappa_{\mathrm{Paley}},
 \mathcal D_{12},
 \iota_{\mathrm{inc}},
 \mathcal K_{\Lambda})
\]
y el lector calibrado
\(E_{\chi_{\mathrm{exc}}}:\mathscr S_\infty
\to\mathcal I^{\mathrm{exc}}_{\chi_{\mathrm{exc}}}\).
Sea
\[
 \mathscr S_\infty\xleftarrow{s_\infty}
 \mathsf Z^{\Gamma}_{9,\infty}
 \xrightarrow{t_\infty}\mathscr S_\infty
\]
la arista testificada de una vuelta completa. Denotemos por
\(\operatorname{ed}_0\) la arista central, por \(\operatorname{ev}_0\) la
evaluación en el estado inicial y por \(F_9\) el desplazamiento de una vuelta.

\begin{theorem}[Factorización excepcional por la arista testificada]
\label{thm:hmt-final-factorizacion-gamma-nueve}
Defínanse
\[
 \overline E^{s}_{\chi_{\mathrm{exc}}}
 :=E_{\chi_{\mathrm{exc}}}\circ s_\infty,
 \qquad
 \overline E^{t}_{\chi_{\mathrm{exc}}}
 :=E_{\chi_{\mathrm{exc}}}\circ t_\infty.
\]
Entonces
\begin{align}
 E_{\chi_{\mathrm{exc}}}\circ\operatorname{ev}_0
 &=\overline E^{s}_{\chi_{\mathrm{exc}}}
   \circ\operatorname{ed}_0,
 \label{eq:hmt-final-factorizacion-gamma-fuente}\\
 E_{\chi_{\mathrm{exc}}}\circ\operatorname{ev}_0\circ F_9
 &=\overline E^{t}_{\chi_{\mathrm{exc}}}
   \circ\operatorname{ed}_0.
 \label{eq:hmt-final-factorizacion-gamma-destino}
\end{align}
\end{theorem}

\begin{proof}
Las identidades de extremos son
\(s_\infty\operatorname{ed}_0=\operatorname{ev}_0\) y
\(t_\infty\operatorname{ed}_0=\operatorname{ev}_0F_9\). La composición con
el lector calibrado da las dos ecuaciones.
\end{proof}

\begin{corollary}[No equivalencia entre el sector temporal \(+1\) y \(1\)
vuelta nonádica completa]
\label{cor:hmt-final-no-equivalencia-retorno-gamma}
La factorización anterior no puede reemplazarse, de forma natural, por una
factorización a través de la sola muestra TRIT \(\tau=+1\).
\end{corollary}

\begin{proof}
El valor \(\tau=+1\) registra retorno, memoria y pasado, pero olvida los datos
de fase, acarreo, ruta y extremo que distinguen aristas completas. Dos aristas
con la misma muestra temporal pueden tener destinos diferentes. Todo mapa que
factorice por esa muestra identifica ambas aristas, mientras
\(\operatorname{ed}_0\) y sus mapas fuente y destino las separan. Por tanto,
el sector \(+1\) no es equivalente a la arista completa \(K\mapsto K+9\).
En particular,
\[
 g(K+9)=g(K)
 \quad\text{y}\quad
 B_{K+9}\ne B_K\quad\text{en general}
\]
siguen siendo afirmaciones distintas.
\end{proof}

\begin{proposition}[Separación calibrada frente a publicación cruda]
\label{prop:hmt-final-separacion-enriquecida-no-cruda}
Supóngase que el lector excepcional enriquecido conserva, para cada
generador, el marco \(f_j\) y la imagen
\(y_j=\Gamma_{f_jA_Wf_j^{-1}}(w_j)\). Entonces recupera
\[
 w_j=\Gamma_{f_jA_Wf_j^{-1}}^{-1}(y_j)
\]
y es inyectivo en el cociente por los cambios de coordenadas incluidos en el
calibre. Si el lector crudo se obtiene olvidando marcos o memoria,
\(E^{\mathrm{cr}}=\pi E^{\sharp}\), no es equivalente al lector enriquecido
salvo que \(\pi\) sea inyectiva sobre \(E^{\sharp}(X)\).
\end{proposition}

\begin{proof}
La primera afirmación es la inversión explícita de un automorfismo en cada
fibra. Para la segunda, si \(\pi\) identifica dos salidas enriquecidas
distintas, el lector crudo identifica sus antecedentes y no puede ser
inyectivo. Recíprocamente, la inyectividad de \(\pi\) sobre la imagen permite
recuperar la salida enriquecida. Así se cierra la separación en el dominio
calibrado y se rechaza la equivalencia silenciosa con el codominio crudo.
\end{proof}

El antiguo cuadrado escrito con símbolos no tipados
\(\mathrm{pub}_{\mathrm{inc}},\mathrm{coord},\mathcal O\) no constituye una
equivalencia matemática: no define siquiera sus morfismos. Su sustituto es el
cono natural de las publicaciones arquimediana, excepcional calibrada y
dimensional, cuyos cuadrados de truncación sí tienen dominio, codominio y
acción declarados.

\subsection{Corredor excepcional y lazos del transporte}

La cadena Paley--Witt--Golay--Leech--\(V^{\natural}\)--Monstruo--Moonshine
es una publicación posterior del estado ya construido. Una estrella de
incidencia y un lazo del grupoide TPK no son objetos equivalentes por su
nombre: la primera es un subobjeto combinatorio; el segundo es una flecha con
fuente igual a destino, composición e inversa. Sin base, orientación y funtor
de transporte no existe una comparación tipada. Por ello, la identificación
automática queda cerrada como no-equivalencia. Un funtor calibrado de
transporte puede producir lazos a partir de estrellas, pero constituye un
resultado adicional y no una premisa del corredor ya demostrado.

\subsection{Canal infinitamente divisible inducido por la transferencia TPK}

\begin{theorem}[Poissonización UCP y raíces coherentes]
\label{thm:hmt-final-poissonizacion-ucp}
Sea \(K\) un canal UCP compatible con los truncamientos TPK\@. Para \(t\ge0\),
defínase
\[
 P_t:=\exp\bigl(t(K-I)\bigr)
 =e^{-t}\sum_{n=0}^{\infty}\frac{t^n}{n!}K^n.
\]
Entonces \((P_t)_{t\ge0}\) es un semigrupo UCP compatible con todos los
truncamientos y
\[
 P_1=(P_{1/m})^m
 \qquad(m\ge1).
\]
En cada nivel finito, una dilatación mínima de Stinespring de \(P_{1/m}\)
produce una raíz tensorial coherente; el límite proyectivo conserva las
marginales TPK.
\end{theorem}

\begin{proof}
Cada \(P_t\) es una combinación convexa normoconvergente de los canales
\(K^n\), por lo que es UCP\@. Como todos los términos son potencias del mismo
operador, la identidad exponencial da \(P_tP_s=P_{t+s}\). La compatibilidad
con los truncamientos se verifica término a término. Tomando \(t=1/m\) se
obtiene la raíz \(m\)-ésima. Stinespring convierte cada composición finita en
una contracción tensorial; la naturalidad de los canales hace coherentes sus
marginales.
\end{proof}

El teorema cierra la existencia de un canal infinitamente divisible con
memoria TPK\@. No afirma que un canal histórico arbitrario coincida con esta
Poissonización: esa igualdad requeriría comparar sus generadores y queda
rechazada mientras no se proporcione tal comparación.

\section{Clausuras espectrales y fronteras de reconocimiento externo}
\label{sec:hmt-final-clausuras-espectrales-externas}

\begin{proposition}[Equivalencia calibrada, no igualdad ambiental]
\label{prop:hmt-final-equivalencia-calibrada-vnueve}
El transporte espectral del doble círculo determina una clase unitaria en el
subespacio cíclico generado por su vector de referencia. Dos representantes
ambientales \(V_9\) son literalmente iguales sólo después de fijar una
identificación unitaria del complemento y un calibre compatible. Sin esos
datos, la igualdad ambiental no queda indeterminada: es una afirmación más
fuerte y no equivalente a la clase específica del subespacio ya construida.
\end{proposition}

La convergencia en norma de los resolventes comprimidos, la recuperación de
alturas, el criterio de Riesz para multiplicidades y la traza diagonal
RH--Moonshine quedan cerrados en sus teoremas respectivos. La traza diagonal
acopla el entero de memoria con el grado de \(V^{\natural}\); no afirma que la
multiplicidad de cada cero sea la dimensión de una representación del
Monstruo. Esa identificación adicional no es equivalente a la existencia de
la traza acoplada.

Quedan fuera del dominio interno, sin reabrir ningún teorema HMT:
\begin{enumerate}
 \item la identificación de la dicotomía de publicaciones HMT analíticas con
 todos los subconjuntos arbitrarios de \(\mathbb R\) del universo ambiente;
 \item la realización macroscópica o termodinámica del cristal temporal
 finito HMT ya construido, incluida la conservación de coherencia a gran
 volumen y su contraste experimental;
 \item el endurecimiento decimal de ventanas espectrales mediante una
 ejecución concreta de aritmética intervalar;
 \item una identificación representacional adicional entre multiplicidades
 espectrales individuales y módulos irreducibles del Monstruo.
\end{enumerate}
Estas cuatro fronteras son reconocimientos o comparaciones externos, no
lagunas de la construcción del continuo, del cierre cardinal, de RH ni del
entrelazamiento del doble círculo.

La residencia terminal conserva, sin permutación,
\[
\begin{aligned}
 \mathrm{RH}
 &\longrightarrow \text{doble círculo y campo espectral}
 \longrightarrow \mathrm{Navier\text{--}Stokes}\\
 &\longrightarrow \mathrm{Yang\text{--}Mills}
 \longrightarrow \mathrm{Hodge}
 \longrightarrow \mathrm{BSD}
 \longrightarrow \text{falsadores}.
\end{aligned}
\]











\HMTFlushCurrentChapter
